% book.tex -- generated by tools/assemble_book.py; edit the
% chapter files, not this file, then re-run the assembler.
\documentclass[11pt]{book}
\usepackage[paperwidth=145mm,paperheight=200mm,top=18mm,bottom=20mm,inner=20mm,outer=16mm]{geometry}
\usepackage{brumfiel}
% figures01.tex -- Chapter 1 (INTRODUCTION) figures for Brumfiel, _Geometry_.
% Macro names: \FIG + I (chapter 1) + spelled-out figure number(s), letters only.
% Every constrained point is computed, never eyeballed:
%   bisectors  -> equal-angle construction / (intersection of ...)
%   altitudes  -> ($(A)!(P)!(B)$)
%   midpoints  -> ($(A)!0.5!(B)$)
%   parallels  -> shared direction vector / equal ratios
% Constraint audit: tools/constraints/ch01.py

%--------------------------------------------------------------- Fig 1-1
% Plane figures. Freeform illustrative shapes traced from scan01 p24 (book p.7).
\newcommand{\FIGIONE}{%
\begin{bkfigure}[\linewidth]
\begin{tikzpicture}[scale=0.92, every node/.style={font=\scriptsize}]
  % --- row 1
  \draw[given] (0,0) -- (2.0,0);
  \fill (0,0) circle (1.1pt) (2.0,0) circle (1.1pt);
  \node[vlab, above] at (0,0.02) {$A$};
  \node[vlab, above] at (2.0,0.02) {$B$};
  \node[below=5pt] at (1.0,-0.05) {Line segment};

  \draw[given] (3.3,0) -- (5.3,0);
  \fill (3.3,0) circle (1.1pt);
  \node[vlab, above] at (3.3,0.02) {$O$};
  \node[below=5pt] at (4.3,-0.05) {Ray from $O$};

  \draw[given] (6.6,-0.12) -- (8.7,-0.12);
  \draw[given] (6.6,-0.12) -- (8.25,0.95);
  \node[below=5pt] at (7.6,-0.17) {Angle};

  % --- row 2
  \begin{scope}[yshift=-2.35cm]
    \draw[given] (0.05,0.22) -- (0.95,1.05) -- (1.95,0) -- cycle;
    \node[below=5pt] at (1.0,-0.05) {Triangle};

    \draw[given] (2.85,0) rectangle (4.15,1.3);
    \node[below=5pt] at (3.5,-0.05) {Square};

    \draw[given] (4.95,0) -- (6.75,0) -- (7.35,1.15) -- (5.55,1.15) -- cycle;
    \node[below=5pt] at (6.15,-0.05) {Parallelogram};

    \draw[given] (8.6,0.6) circle (0.62);
    \node[below=5pt] at (8.6,-0.05) {Circle};
  \end{scope}

  % --- row 3
  \begin{scope}[yshift=-4.7cm]
    \draw[given] (0,0) -- (2.0,0) -- (1.65,1.15) -- (0.35,1.15) -- cycle;
    \node[below=5pt] at (1.0,-0.05) {Trapezoid};

    \draw[given] (3.05,0.62) -- (3.85,1.35) -- (5.05,0) -- (3.25,0.28) -- cycle;
    \node[align=center, below=5pt] at (4.0,-0.05)
         {Quadrilateral, or\\4-sided polygon};

    \draw[given] (6.5,0.55) -- (7.05,1.35) -- (8.35,0.62) -- (7.75,0) -- (6.95,0.1) -- cycle;
    \node[align=center, below=5pt] at (7.5,-0.05)
         {Pentagon, or\\5-sided polygon};
  \end{scope}

  % --- row 4
  \begin{scope}[yshift=-7.35cm]
    \draw[given] (3.3,0) rectangle (5.5,1.25);
    \node[below=5pt] at (4.4,-0.05) {Rectangle};
  \end{scope}
\end{tikzpicture}
\figcap{1--1}
\end{bkfigure}}

%--------------------------------------------------------------- Fig 1-2
% Solid figures. Hidden edges dashed, as the book draws them.
\newcommand{\FIGITWO}{%
\begin{bkfigure}[\linewidth]
\begin{tikzpicture}[scale=0.92, every node/.style={font=\scriptsize},
                    hid/.style={fig, dashed, dash pattern=on 2pt off 1.6pt}]
  % ---- cube
  \begin{scope}
    \def\d{0.42}
    \draw[given] (0,0) rectangle (1.25,1.25);
    \draw[given] (\d,\d) -- (\d+1.25,\d) -- (\d+1.25,\d+1.25) -- (\d,\d+1.25);
    \draw[hid]   (\d,\d) -- (\d,\d+1.25);
    \draw[given] (0,1.25) -- (\d,\d+1.25);
    \draw[given] (1.25,1.25) -- (\d+1.25,\d+1.25);
    \draw[given] (1.25,0) -- (\d+1.25,\d);
    \draw[hid]   (0,0) -- (\d,\d);
    \node[below=5pt] at (0.85,-0.05) {Cube};
  \end{scope}
  % ---- box
  \begin{scope}[xshift=3.15cm]
    \def\d{0.42}
    \draw[given] (0,0) rectangle (1.95,1.05);
    \draw[given] (\d,\d) -- (\d+1.95,\d) -- (\d+1.95,\d+1.05) -- (\d,\d+1.05);
    \draw[hid]   (\d,\d) -- (\d,\d+1.05);
    \draw[given] (0,1.05) -- (\d,\d+1.05);
    \draw[given] (1.95,1.05) -- (\d+1.95,\d+1.05);
    \draw[given] (1.95,0) -- (\d+1.95,\d);
    \draw[hid]   (0,0) -- (\d,\d);
    \node[align=center, below=5pt] at (1.2,-0.05)
         {Rectangular\\parallelepiped or box};
  \end{scope}
  % ---- circular cylinder
  \begin{scope}[xshift=6.95cm]
    \draw[given] (0,0) -- (0,1.6);
    \draw[given] (1.3,0) -- (1.3,1.6);
    \draw[given] (0.65,1.6) ellipse (0.65 and 0.24);
    \draw[given] (0.65,0) ++(180:0.65) arc (180:360:0.65 and 0.24);
    \draw[hid]   (0.65,0) ++(180:0.65) arc (180:0:0.65 and 0.24);
    \node[align=center, below=5pt] at (0.65,-0.26) {Circular\\cylinder};
  \end{scope}

  % ---- tetrahedron: as the book draws it, the near vertex is at the front and
  % the single hidden element is the back base edge L--R
  \begin{scope}[yshift=-3.1cm]
    \coordinate (tL) at (0,0.402);
    \coordinate (tR) at (1.517,0.674);
    \coordinate (tF) at (0.817,0);
    \coordinate (tT) at (0.596,1.750);
    \draw[given] (tT) -- (tL) -- (tF) -- (tR) -- cycle;
    \draw[given] (tT) -- (tF);
    \draw[hid]   (tL) -- (tR);
    \node[below=5pt] at (0.76,-0.1) {Tetrahedron};
  \end{scope}
  % ---- pyramid
  \begin{scope}[xshift=3.15cm, yshift=-3.1cm]
    \coordinate (p1) at (0,0.15);
    \coordinate (p2) at (1.45,0);
    \coordinate (p3) at (2.05,0.55);
    \coordinate (p4) at (0.6,0.72);
    \coordinate (pa) at (1.05,1.75);
    \draw[given] (p1) -- (p2) -- (p3);
    \draw[hid]   (p3) -- (p4) -- (p1);
    \draw[given] (p1) -- (pa) -- (p2);
    \draw[given] (p3) -- (pa);
    \draw[hid]   (p4) -- (pa);
    \node[below=5pt] at (1.0,-0.1) {Pyramid};
  \end{scope}
  % ---- prism: upright triangular prism, near vertical edge facing the reader,
  % top face fully visible, only the back bottom edge hidden (as in the book)
  \begin{scope}[xshift=6.95cm, yshift=-3.1cm]
    \coordinate (uL) at (0,1.732);        \coordinate (dL) at (0,0.426);
    \coordinate (uR) at (1.205,1.656);    \coordinate (dR) at (1.205,0.350);
    \coordinate (uF) at (0.266,1.306);    \coordinate (dF) at (0.266,0);
    \draw[given] (uL) -- (uR) -- (uF) -- cycle;   % top face
    \draw[given] (uL) -- (dL);
    \draw[given] (uR) -- (dR);
    \draw[given] (uF) -- (dF);
    \draw[given] (dL) -- (dF) -- (dR);
    \draw[hid]   (dL) -- (dR);
    \node[below=5pt] at (0.60,-0.1) {Prism};
  \end{scope}

  % ---- sphere
  \begin{scope}[yshift=-6.2cm]
    \draw[given] (0.85,0.85) circle (0.78);
    \draw[given] (0.85,0.85) ++(180:0.78) arc (180:360:0.78 and 0.27);
    \draw[hid]   (0.85,0.85) ++(180:0.78) arc (180:0:0.78 and 0.27);
    \node[below=5pt] at (0.85,-0.05) {Sphere};
  \end{scope}
  % ---- torus
  \begin{scope}[xshift=3.15cm, yshift=-6.2cm]
    \draw[given] (1.05,0.85) circle (0.82);
    \draw[given] (1.12,0.92) circle (0.40);
    \draw[given] (0.72,0.92) arc (180:300:0.40 and 0.16);
    \node[align=center, below=5pt] at (1.05,-0.05)
         {Torus, or doughnut,\\or anchor ring};
  \end{scope}
  % ---- cone
  \begin{scope}[xshift=6.95cm, yshift=-6.2cm]
    \draw[given] (0,0.35) -- (0.78,1.85) -- (1.56,0.35);
    \draw[given] (0.78,0.35) ++(180:0.78) arc (180:360:0.78 and 0.27);
    \draw[hid]   (0.78,0.35) ++(180:0.78) arc (180:0:0.78 and 0.27);
    \node[below=5pt] at (0.78,-0.05) {Cone};
  \end{scope}
\end{tikzpicture}
\figcap{1--2}
\end{bkfigure}}

%--------------------------------------------------------------- Fig 1-3
% "Strange" sets of points. ORGANIC -- traced against scan02 p2 (book p.9):
% damped zigzag, chain of shrinking circles, recursively quartered square,
% face built from circles.
\newcommand{\FIGITHREE}{%
\begin{bkfigure}[0.92\linewidth]
\begin{tikzpicture}[scale=1.0]
  % ---- damped zigzag: amplitude falls geometrically left to right
  \begin{scope}[shift={(0,3.05)}]
    % the path must OPEN with a coordinate -- starting it with "--" emits a
    % lineto with no current point and some viewers drop the whole zigzag
    \draw[given] (0,0.42)
      \foreach \i [evaluate={\x=0.30*\i; \s=ifthenelse(mod(\i,2)==0,1,-1);
                            \a=0.42*pow(0.80,\i)} ] in {1,...,15} { -- (\x,{\s*\a}) }
      -- (4.85,0);
  \end{scope}
  % ---- chain of shrinking circles converging to a point
  \begin{scope}[shift={(0.15,1.55)}]
    \foreach \i [evaluate={\r=0.235*pow(0.80,\i); \c=0.44*(1-pow(0.80,\i))/0.20}] in {0,...,9}
      \draw[given] ({\c+\r},0) ellipse ({\r} and {1.28*\r});
    \draw[given] (2.30,0.05) -- (2.62,0) -- (2.30,-0.05);
  \end{scope}
  % ---- square quartered again and again in its upper-right corner
  \begin{scope}[shift={(3.15,0.60)}]
    \def\S{2.0}
    \draw[given] (0,0) rectangle (\S,\S);
    \draw[given] (0,{\S/2}) -- (\S,{\S/2});
    \draw[given] ({\S/2},0) -- ({\S/2},\S);
    \draw[given] ({\S/2},{3*\S/4}) -- (\S,{3*\S/4});
    \draw[given] ({3*\S/4},{\S/2}) -- ({3*\S/4},\S);
    \draw[given] ({3*\S/4},{7*\S/8}) -- (\S,{7*\S/8});
    \draw[given] ({7*\S/8},{3*\S/4}) -- ({7*\S/8},\S);
    \draw[given] ({7*\S/8},{15*\S/16}) -- (\S,{15*\S/16});
    \draw[given] ({15*\S/16},{7*\S/8}) -- ({15*\S/16},\S);
  \end{scope}
  % ---- face made of circles
  \begin{scope}[shift={(1.55,-1.25)}]
    \draw[given] (0,0) circle (0.80 and 0.86);      % head
    \draw[given] (-0.86,0.02) circle (0.10);        % ears
    \draw[given] ( 0.86,0.02) circle (0.10);
    \draw[given] (-0.30,0.26) circle (0.15);        % eyes
    \draw[given] ( 0.30,0.26) circle (0.15);
    \draw[given] (-0.30,0.26) circle (0.06);
    \draw[given] ( 0.30,0.26) circle (0.06);
    \draw[given] (0,-0.06) circle (0.09 and 0.14);  % nose
    \draw[given] (0,-0.48) circle (0.17);           % mouth
  \end{scope}
\end{tikzpicture}
\figcap{1--3}
\end{bkfigure}}

%--------------------------------------------------------------- Fig 1-4
% Angle: vertex, interior, exterior.
\newcommand{\FIGIFOUR}{%
\begin{bkfigure}[0.60\linewidth]
\begin{tikzpicture}[scale=1.0]
  \coordinate (A) at (0,0);
  \draw[given] (A) -- (2.35,1.20);
  \draw[given] (A) -- (2.35,-1.20);
  \node[vlab] at (-0.20,0) {$A$};
  \node[slab] at (-0.42,0.42) {Vertex};
  \node[slab] at (1.02,0) {Interior};
  \node[slab] at (0.62,-0.88) {Exterior};
\end{tikzpicture}
\figcap{1--4}
\end{bkfigure}}

%--------------------------------------------------------------- Figs 1-5, 1-6
% 1-5: ray AC bisects angle BAD -- C placed at exactly half the angle of B.
% 1-6: acute (< 90) and obtuse (> 90) angles.
\newcommand{\FIGIFIVESIX}{%
\begin{bkfigure}[\linewidth]
\begin{minipage}[t]{0.48\linewidth}\centering
\begin{tikzpicture}[scale=1.0]
  \coordinate (A) at (0,0);
  \coordinate (D) at (2.30,0);
  \coordinate (B) at (58:2.30);
  \coordinate (C) at (29:2.15);       % exact bisector of angle BAD
  \draw[given] (A) -- (D);
  \draw[given] (A) -- (B);
  \draw[given] (A) -- (C);
  \node[vlab, left]  at (A) {$A$};
  \node[vlab, right] at (D) {$D$};
  \node[vlab, above] at (B) {$B$};
  \node[vlab, right] at (C) {$C$};
\end{tikzpicture}
\figcap{1--5}
\end{minipage}\hfill
\begin{minipage}[t]{0.48\linewidth}\centering
\begin{tikzpicture}[scale=1.0]
  % acute: vertex bottom-left, 32 degrees
  \draw[given] (0,0) -- (1.55,0);
  \draw[given] (0,0) -- (32:1.55);
  \node[slab, above] at (0.72,0.85) {Acute angle};
  % obtuse: vertex bottom-right; leftward ray at 180 deg, second ray at 47 deg,
  % so the angle at the vertex is 180 - 47 = 133 deg
  \draw[given] (3.35,0) -- (1.95,0);
  \draw[given] (3.35,0) -- ($(3.35,0)+(47:1.55)$);
  \node[slab, above] at (2.85,1.15) {Obtuse angle};
\end{tikzpicture}
\figcap{1--6}
\end{minipage}
\end{bkfigure}}

%--------------------------------------------------------------- Figs 1-7, 1-8
% 1-7: right triangle, right angle at the lower-left vertex.
% 1-8: angle AOB is a right angle; C interior.
\newcommand{\FIGISEVENEIGHT}{%
\begin{bkfigure}[\linewidth]
\begin{minipage}[t]{0.48\linewidth}\centering
\begin{tikzpicture}[scale=1.0]
  \coordinate (R) at (0,0);
  \coordinate (T) at (0,1.75);
  \coordinate (S) at (2.15,0);
  \draw[given] (R) -- (T) -- (S) -- cycle;
  \draw[fig] (0,0.17) -- (0.17,0.17) -- (0.17,0);
  \node[slab] at (0.56,0.30) {90\textdegree};
\end{tikzpicture}
\figcap{1--7}
\end{minipage}\hfill
\begin{minipage}[t]{0.48\linewidth}\centering
\begin{tikzpicture}[scale=1.0]
  \coordinate (O) at (0,0);
  \coordinate (A) at (0,1.75);
  \coordinate (B) at (1.85,0);
  \coordinate (C) at (33:1.62);
  \draw[given] (O) -- (A);
  \draw[given] (O) -- (B);
  \draw[given] (O) -- (C);
  \fill (A) circle (1.2pt) (B) circle (1.2pt) (C) circle (1.2pt);
  \node[vlab, above] at (A) {$A$};
  \node[vlab, right] at (B) {$B$};
  \node[vlab, right] at (C) {$C$};
  \node[vlab, below left, xshift=-1pt, yshift=-0pt] at (O) {$O$};
\end{tikzpicture}
\figcap{1--8}
\end{minipage}
\end{bkfigure}}

%--------------------------------------------------------------- Figs 1-9, 1-10
% 1-9: A, O, C collinear on l (so angle AOB and angle BOC are supplementary).
% 1-10: straight angle A-O-B.
\newcommand{\FIGININETEN}{%
\begin{bkfigure}[\linewidth]
\begin{minipage}[t]{0.48\linewidth}\centering
\begin{tikzpicture}[scale=1.0]
  \coordinate (A) at (-1.55,0);
  \coordinate (O) at (0,0);
  \coordinate (C) at (1.55,0);
  \coordinate (B) at (38:1.62);
  \draw[given] (A) -- (C);
  \draw[given] (O) -- (B);
  \fill (A) circle (1.2pt) (O) circle (1.2pt) (C) circle (1.2pt) (B) circle (1.2pt);
  \node[vlab, left]  at (A) {$A$};
  \node[vlab, right] at (C) {$C$};
  \node[vlab, above right] at (B) {$B$};
  \node[vlab, below] at (O) {$O$};
  \node[vlab, below] at (0.80,-0.02) {$l$};
\end{tikzpicture}
\figcap{1--9}
\end{minipage}\hfill
\begin{minipage}[t]{0.48\linewidth}\centering
\begin{tikzpicture}[scale=1.0]
  \coordinate (A) at (-1.45,0);
  \coordinate (O) at (0,0);
  \coordinate (B) at (1.45,0);
  \draw[given] (A) -- (B);
  \fill (A) circle (1.2pt) (O) circle (1.2pt) (B) circle (1.2pt);
  \node[vlab, above] at (A) {$A$};
  \node[vlab, above] at (B) {$B$};
  \node[vlab, below] at (O) {$O$};
\end{tikzpicture}
\figcap{1--10}
\end{minipage}
\end{bkfigure}}

%--------------------------------------------------------------- Fig 1-11
\newcommand{\FIGIELEVEN}{%
\begin{bkfigure}[0.52\linewidth]
\begin{tikzpicture}[scale=1.0]
  % vertex ratios measured off the book's own quadrilateral
  \draw[given] (0,0.310) -- (0.558,1.565) -- (2.200,0.558) -- (1.270,0) -- cycle;
\end{tikzpicture}
\figcap{1--11}
\end{bkfigure}}

%--------------------------------------------------------------- Figs 1-12, 1-13
% 1-12: equilateral triangle by the two-arc construction (Euclid I.1).
% 1-13: an angle and its copy -- both drawn at exactly 52 degrees.
\newcommand{\FIGITWELVETHIRTEEN}{%
\begin{bkfigure}[\linewidth]
\begin{minipage}[t]{0.46\linewidth}\centering
\begin{tikzpicture}[scale=1.0]
  \def\s{2.05}
  \coordinate (A) at (0,0);
  \coordinate (B) at (\s,0);
  \coordinate (C) at (60:\s);          % exact equilateral apex
  \draw[given] (A) -- (B) -- (C) -- cycle;
  \draw[fig] (A) ++(-18:\s) arc (-18:78:\s);
  \draw[fig] (B) ++(198:\s) arc (198:102:\s);
  \draw[fig] ($(C)+(-0.09,-0.09)$) -- ($(C)+(0.09,0.09)$);
  \draw[fig] ($(C)+(-0.09,0.09)$) -- ($(C)+(0.09,-0.09)$);
  % compass-point marks where each arc passes through the opposite base vertex
  \draw[fig] ($(A)+(135:0.11)$) -- ($(A)+(-45:0.11)$);
  \draw[fig] ($(B)+(45:0.11)$)  -- ($(B)+(225:0.11)$);
\end{tikzpicture}
\figcap{1--12}
\end{minipage}\hfill
\begin{minipage}[t]{0.50\linewidth}\centering
\begin{tikzpicture}[scale=1.0]
  \foreach \x/\lo/\lb/\la in {0/{O}/{B}/{A}, 2.75/{O\mkern2mu'}/{B\mkern2mu'}/{A\mkern2mu'}} {
    \begin{scope}[xshift=\x cm]
      \coordinate (V) at (0,0);
      \draw[given] (V) -- (1.45,0);
      \draw[given] (V) -- (60:1.60);
      \draw[fig] (V) ++(0:0.62) arc (0:60:0.62);
      % the book marks where the arc cuts each ray: a tick on OB, a cross on OA
      \draw[fig] ($(V)+(0.62,-0.10)$) -- ($(V)+(0.62,0.10)$);
      \draw[fig] ($(V)+(60:0.62)+(15:0.11)$)  -- ($(V)+(60:0.62)+(195:0.11)$);
      \draw[fig] ($(V)+(60:0.62)+(105:0.11)$) -- ($(V)+(60:0.62)+(285:0.11)$);
      \node[vlab, below left, xshift=-0pt, yshift=-0pt] at (V) {$\lo$};
      \node[vlab, right] at (1.45,0) {$\lb$};
      \node[vlab, above] at (60:1.60) {$\la$};
    \end{scope}}
\end{tikzpicture}
\figcap{1--13}
\end{minipage}
\end{bkfigure}}

%--------------------------------------------------------------- Fig 1-14
% Compass designs. ORGANIC -- traced against scan02 p9 (book p.14):
% (left) four circles of radius r whose centres lie at distance r from a common
% point; (right) the six-petal rosette inscribed in a circle.
\newcommand{\FIGIFOURTEEN}{%
\begin{bkfigure}[0.86\linewidth]
\begin{tikzpicture}[scale=1.0]
  % four-circle design
  \begin{scope}
    \def\r{0.78}
    \foreach \a in {0,90,180,270} \draw[given] (\a:\r) circle (\r);
  \end{scope}
  % six-petal rosette
  \begin{scope}[xshift=4.15cm]
    \def\R{1.12}
    \begin{scope}
      \clip (0,0) circle (\R);
      \foreach \a in {0,60,...,300} \draw[given] (\a:\R) circle (\R);
    \end{scope}
    \draw[given] (0,0) circle (\R);
  \end{scope}
\end{tikzpicture}
\figcap{1--14}
\end{bkfigure}}

%--------------------------------------------------------------- Figs 1-15, 1-16
% 1-15: angle bisected with compass; Z is the exact arc crossing, so VZ bisects.
% 1-16: adjacent angles AOB and BOC sharing ray OB.
\newcommand{\FIGIFIFTEENSIXTEEN}{%
\begin{bkfigure}[\linewidth]
\begin{minipage}[t]{0.48\linewidth}\centering
\begin{tikzpicture}[scale=1.0]
  \coordinate (V) at (0,0);
  \draw[given] (V) -- (26:2.15);
  \draw[given] (V) -- (-26:2.15);
  \coordinate (X) at (26:0.72);
  \coordinate (Y) at (-26:0.72);
  \draw[fig] (V) ++(34:0.72) arc (34:-34:0.72);
  % arcs of radius 0.716170 about X and Y meet at Z on the bisector; as in the
  % book only their crossing is marked, the two small arcs are not drawn
  \coordinate (Z) at (1.29,0);
  \draw[given] (V) -- (2.15,0);
  \draw[fig] ($(Z)+(-0.09,-0.09)$) -- ($(Z)+(0.09,0.09)$);
  \draw[fig] ($(Z)+(-0.09,0.09)$) -- ($(Z)+(0.09,-0.09)$);
\end{tikzpicture}
\figcap{1--15}
\end{minipage}\hfill
\begin{minipage}[t]{0.48\linewidth}\centering
\begin{tikzpicture}[scale=1.0]
  \coordinate (O) at (0,0);
  \coordinate (A) at (36:2.05);
  \coordinate (B) at (-2:2.05);
  \coordinate (C) at (-24:2.05);
  \draw[given] (O) -- (A);
  \draw[given] (O) -- (B);
  \draw[given] (O) -- (C);
  \node[vlab, left]  at (O) {$O$};
  \node[vlab, above right] at (A) {$A$};
  \node[vlab, right] at (B) {$B$};
  \node[vlab, right] at (C) {$C$};
\end{tikzpicture}
\figcap{1--16}
\end{minipage}
\end{bkfigure}}

%--------------------------------------------------------------- Figs 1-17, 1-18
% 1-17: three pairs of angles, none of them adjacent (traced from scan02 p9):
%   (a) two angles at two DIFFERENT vertices;
%   (b) two angles sharing a vertex and a side but overlapping interiors;
%   (c) two angles sharing only the vertex -- no common side.
% 1-18: copying a segment onto a line at A.
\newcommand{\FIGISEVENTEENEIGHTEEN}{%
\begin{bkfigure}[\linewidth]
\begin{minipage}[t]{0.56\linewidth}\centering
\begin{tikzpicture}[scale=0.95]
  % (a) two nearby but distinct vertices
  \begin{scope}
    \coordinate (V) at (0,0);
    \coordinate (W) at (0.16,-0.14);
    \draw[given] (V) -- (34:1.75);
    \draw[given] (V) -- (-9:1.70);
    \draw[given] (W) -- ($(W)+(-6:1.55)$);
    \draw[given] (W) -- ($(W)+(-44:1.30)$);
    \node[vlab] at (12.5:0.80) {$A$};
    \node[vlab] at ($(W)+(-25:0.98)$) {$B$};
  \end{scope}
  % (b) shared vertex and side, interiors overlap
  \begin{scope}[xshift=3.35cm, yshift=0.05cm]
    \coordinate (V) at (1.30,0);
    \draw[given] (V) -- ($(V)+(140:1.45)$);
    \draw[given] (V) -- ($(V)+(218:1.35)$);
    \draw[given] (V) -- ($(V)+(2:1.15)$);
    \node[vlab] at ($(V)+(35:0.52)$) {$A$};
    \node[vlab] at ($(V)+(-35:0.52)$) {$B$};
  \end{scope}
  % (c) shared vertex only
  \begin{scope}[xshift=1.55cm, yshift=-2.55cm]
    \coordinate (V) at (0,0);
    \draw[given] (V) -- (92:1.55);
    \draw[given] (V) -- (62:1.45);
    \draw[given] (V) -- (-10:1.50);
    \draw[given] (V) -- (-38:1.45);
    \node[vlab] at (77:1.15) {$A$};
    \node[vlab] at (-24:1.25) {$B$};
  \end{scope}
\end{tikzpicture}
\figcap{1--17}
\end{minipage}\hfill
\begin{minipage}[t]{0.40\linewidth}\centering
\begin{tikzpicture}[scale=1.0]
  \draw[given] (0.28,0.52) -- (2.05,0.52);        % the given segment
  \draw[given] (0,0) -- (2.30,0);                  % the line
  \coordinate (A) at (0.28,0);
  \fill (A) circle (1.2pt);
  \draw[fig] (2.05,-0.13) -- (2.05,0.13);          % compass mark at B
  \node[vlab, below left, xshift=2pt, yshift=-0pt] at (A) {$A$};
  \node[vlab, above right, xshift=-3pt, yshift=0pt] at (2.05,0.02) {$B$};
\end{tikzpicture}
\figcap{1--18}
\end{minipage}
\end{bkfigure}}

%--------------------------------------------------------------- Figs 1-19, 1-20
% 1-19: perpendicular bisector of AB; C is the exact midpoint, l is exactly
%       perpendicular (vertical through C).
% 1-20: perpendicular to l at a point O on l.
\newcommand{\FIGININETEENTWENTY}{%
\begin{bkfigure}[\linewidth]
\begin{minipage}[t]{0.48\linewidth}\centering
\begin{tikzpicture}[scale=1.0]
  % drawn 1.4x the first draft's size, to the book's own proportions
  \coordinate (A) at (-1.89,0);
  \coordinate (B) at (1.89,0);
  \coordinate (C) at ($(A)!0.5!(B)$);     % exact midpoint
  \draw[given] (A) -- (B);
  \draw[given] (0,-2.17) -- (0,2.17);
  % compass arcs of radius 2.45 about A and B; the sweep runs past each
  % crossing, as the book draws it. Chained 23-degree pieces keep every
  % Bezier tight against its own arc.
  \draw[aux] (A) ++(-46:2.45) arc (-46:-23:2.45) arc (-23:0:2.45)
                              arc (0:23:2.45)    arc (23:46:2.45);
  \draw[aux] (B) ++(134:2.45) arc (134:157:2.45) arc (157:180:2.45)
                              arc (180:203:2.45) arc (203:226:2.45);
  \foreach \y in {1.558974,-1.558974} {   % sqrt(2.45^2 - 1.89^2)
    \draw[fig] (-0.12,\y-0.12) -- (0.12,\y+0.12);
    \draw[fig] (-0.12,\y+0.12) -- (0.12,\y-0.12);}
  % where the arcs cut AB (x = -1.89+2.45 and 1.89-2.45)
  \foreach \x in {-0.56,0.56} \draw[fig] (\x,-0.15) -- (\x,0.15);
  \fill (A) circle (1.2pt) (B) circle (1.2pt);
  \node[vlab, left]  at (A) {$A$};
  \node[vlab, right] at (B) {$B$};
  \node[vlab] at (0.20,-0.33) {$C$};
  \node[vlab, below right, xshift=0pt, yshift=-1pt] at (0,-2.17) {$l$};
\end{tikzpicture}
\figcap{1--19}
\end{minipage}\hfill
\begin{minipage}[t]{0.48\linewidth}\centering
\begin{tikzpicture}[scale=1.0]
  % same 1.4x enlargement as Fig 1-19, so the pair matches
  \coordinate (O) at (0,0);
  \draw[given] (-2.03,0) -- (2.03,0);
  \draw[given] (0,-2.03) -- (0,2.17);
  \foreach \x in {-1.12,1.12} \draw[fig] (\x,-0.18) -- (\x,0.18);
  \draw[aux] (-1.12,0) ++(-58:1.61)
      arc (-58:-38.667:1.61) arc (-38.667:-19.333:1.61) arc (-19.333:0:1.61)
      arc (0:19.333:1.61)    arc (19.333:38.667:1.61)   arc (38.667:58:1.61);
  \draw[aux] ( 1.12,0) ++(122:1.61)
      arc (122:141.333:1.61)   arc (141.333:160.667:1.61) arc (160.667:180:1.61)
      arc (180:199.333:1.61)   arc (199.333:218.667:1.61) arc (218.667:238:1.61);
  \foreach \y in {1.156590,-1.156590} {   % sqrt(1.61^2 - 1.12^2)
    \draw[fig] (-0.12,\y-0.12) -- (0.12,\y+0.12);
    \draw[fig] (-0.12,\y+0.12) -- (0.12,\y-0.12);}
  \node[vlab, above right, xshift=0pt, yshift=0pt] at (0,2.17) {$m$};
  \node[vlab, left] at (-2.03,0) {$l$};
  \node[vlab] at (0.20,-0.33) {$O$};
\end{tikzpicture}
\figcap{1--20}
\end{minipage}
\end{bkfigure}}

%--------------------------------------------------------------- Figs 1-21, 1-22
% 1-21: perpendicular to l from an external point P (m is exactly vertical).
% 1-22: the three medians; O is the exact centroid.
\newcommand{\FIGITWENTYONETWENTYTWO}{%
\begin{bkfigure}[\linewidth]
\begin{minipage}[t]{0.46\linewidth}\centering
\begin{tikzpicture}[scale=1.0]
  \coordinate (P) at (0,1.30);
  \draw[given] (-1.60,0) -- (1.60,0);
  \draw[given] (0,-1.15) -- (0,1.55);
  % arc of radius 1.75 about P cuts l at x = +- sqrt(1.75^2 - 1.30^2) = 1.171538;
  % the crossings are at 227.95 deg and 312.05 deg, so the sweep runs past them
  \draw[fig] (P) ++(222:1.75) arc (222:318:1.75);
  \draw[fig] (-0.09,1.30-0.09) -- (0.09,1.30+0.09);
  \draw[fig] (-0.09,1.30+0.09) -- (0.09,1.30-0.09);
  % arcs of radius 1.45 about those points meet at y = -sqrt(1.45^2 - 1.171538^2)
  \draw[fig] (-0.09,-0.854400-0.09) -- (0.09,-0.854400+0.09);
  \draw[fig] (-0.09,-0.854400+0.09) -- (0.09,-0.854400-0.09);
  \node[vlab, left] at (P) {$P$};
  \node[vlab, left] at (-1.60,0) {$l$};
  \node[vlab, below] at (0,-1.15) {$m$};
\end{tikzpicture}
\figcap{1--21}
\end{minipage}\hfill
\begin{minipage}[t]{0.50\linewidth}\centering
\begin{tikzpicture}[scale=1.0]
  % proportions from the book: apex at 0.64 of the base, height 0.627 of the base
  \coordinate (A) at (0,0);
  \coordinate (B) at (3.40,0);
  \coordinate (C) at (2.176,2.132);
  \coordinate (Mc) at ($(A)!0.5!(B)$);
  \coordinate (Ma) at ($(B)!0.5!(C)$);
  \coordinate (Mb) at ($(C)!0.5!(A)$);
  \coordinate (O)  at ($(A)!0.6666666!(Ma)$);   % exact centroid
  \draw[given] (A) -- (B) -- (C) -- cycle;
  \draw[aux] (A) -- (Ma);
  \draw[aux] (B) -- (Mb);
  \draw[aux] (C) -- (Mc);
  \draw[fig] ($(O)+(-0.08,-0.08)$) -- ($(O)+(0.08,0.08)$);
  \draw[fig] ($(O)+(-0.08,0.08)$) -- ($(O)+(0.08,-0.08)$);
  % 296.33 deg bisects the widest gap between the six median rays at O
  \node[vlab] at ($(O)+(296.33:0.47)$) {$O$};
\end{tikzpicture}
\figcap{1--22}
\end{minipage}
\end{bkfigure}}

%--------------------------------------------------------------- Fig 1-23
% Cardboard triangle balancing on a fingertip. ORGANIC (pictorial) -- traced
% against scan02 p12 (book p.16). The medians are exact; the hand is stylised.
\newcommand{\FIGITWENTYTHREE}{%
\begin{bkfigure}[0.50\linewidth]
\begin{tikzpicture}[scale=1.0]
  % triangle seen almost edge-on, balanced at its centroid
  \coordinate (A) at (-1.55,0.18);
  \coordinate (B) at (1.62,0.30);
  \coordinate (C) at (-0.28,0.92);
  \coordinate (Ma) at ($(B)!0.5!(C)$);
  \coordinate (Mb) at ($(C)!0.5!(A)$);
  \coordinate (Mc) at ($(A)!0.5!(B)$);
  \coordinate (O)  at ($(A)!0.6666666!(Ma)$);
  \draw[given] (A) -- (B) -- (C) -- cycle;
  \draw[aux] (A) -- (Ma);
  \draw[aux] (B) -- (Mb);
  \draw[aux] (C) -- (Mc);
  % stylised finger and hand: the rounded fingertip touches the centroid and
  % the finger runs down to meet the top of the fist (it must not float free)
  \draw[fig] ($(O)+(-0.075,-1.14)$) -- ($(O)+(-0.075,-0.02)$)
             arc (180:0:0.075) -- ($(O)+(0.075,-1.14)$);
  \draw[fig] (-0.30,-0.62) .. controls (-0.34,-0.98) and (-0.30,-1.24) .. (-0.24,-1.42);
  \draw[fig] ( 0.16,-0.62) .. controls (0.22,-0.98) and (0.24,-1.24) .. (0.20,-1.42);
  \draw[fig] (-0.30,-0.62) .. controls (-0.20,-0.70) and (0.06,-0.70) .. (0.16,-0.62);
  \foreach \k in {0,1,2} {
    \draw[fig] (-0.30,-0.78-0.17*\k)
       .. controls (-0.40,-0.84-0.17*\k) and (-0.42,-0.92-0.17*\k) ..
          (-0.32,-0.96-0.17*\k);}
  \draw[fig] (-0.30,-1.42) rectangle (0.22,-1.60);
\end{tikzpicture}
\figcap{1--23}
\end{bkfigure}}

%--------------------------------------------------------------- Fig 1-24
% Altitudes. (a) acute, foot inside; (b) obtuse, foot on the extension of the
% base; (c) the three altitude lines concurrent at O. All feet are exact
% projections ($(X)!(P)!(Y)$).
\newcommand{\FIGITWENTYFOUR}{%
\begin{bkfigure}[\linewidth]
\begin{tikzpicture}[scale=1.0]
  % Proportions measured off the book: all three panels are about the same width
  % (100 : 106 : 110) and grow in height (64 : 92 : 113).
  % ---- (a) apex left of centre, foot inside, right-angle box right of the altitude
  \begin{scope}
    \coordinate (A) at (0,0);
    \coordinate (B) at (2.15,0);
    \coordinate (C) at (0.7256,1.376);
    \coordinate (F) at ($(A)!(C)!(B)$);
    \draw[given] (A) -- (B) -- (C) -- cycle;
    \draw[aux] (C) -- (F);
    \draw[fig] ($(F)+(0.16,0)$) -- ($(F)+(0.16,0.16)$) -- ($(F)+(0,0.16)$);
    \node[slab, below=6pt] at (1.075,0) {(a)};
  \end{scope}
  % ---- (b) obtuse at B: the foot falls beyond B
  \begin{scope}[xshift=3.05cm]
    \coordinate (A) at (0,0);
    \coordinate (B) at (1.3975,0);
    \coordinate (C) at (2.2844,1.9834);
    \coordinate (F) at ($(A)!(C)!(B)$);
    \draw[given] (A) -- (B) -- (C) -- cycle;
    \draw[given] (B) -- (F);
    \draw[aux] (C) -- (F);
    \draw[fig] ($(F)+(-0.16,0)$) -- ($(F)+(-0.16,0.16)$) -- ($(F)+(0,0.16)$);
    \node[slab, below=6pt] at (1.142,0) {(b)};
  \end{scope}
  % ---- (c) three altitudes concurrent; near-isosceles, apex at 0.584 of the base
  \begin{scope}[xshift=6.10cm]
    \coordinate (A) at (0,0);
    \coordinate (B) at (2.365,0);
    \coordinate (C) at (1.3812,2.4295);
    \coordinate (Fc) at ($(A)!(C)!(B)$);
    \coordinate (Fa) at ($(B)!(A)!(C)$);
    \coordinate (Fb) at ($(A)!(B)!(C)$);
    \coordinate (H)  at (intersection of C--Fc and A--Fa);
    \draw[given] (A) -- (B) -- (C) -- cycle;
    \draw[aux] (C) -- (Fc);
    \draw[aux] (A) -- (Fa);
    \draw[aux] (B) -- (Fb);
    \draw[fig] ($(H)+(-0.08,-0.08)$) -- ($(H)+(0.08,0.08)$);
    \draw[fig] ($(H)+(-0.08,0.08)$) -- ($(H)+(0.08,-0.08)$);
    % 176.21 deg bisects the gap between the altitude rays at 150.38 and 202.05
    \node[vlab] at ($(H)+(176.21:0.65)$) {$O$};
    \node[slab, below=6pt] at (1.1825,0) {(c)};
  \end{scope}
\end{tikzpicture}
\figcap{1--24}
\end{bkfigure}}

%--------------------------------------------------------------- Fig 1-25
% The three angle bisectors meet at the incentre. Each bisector is drawn from a
% vertex through the exact incentre to the opposite side.
\newcommand{\FIGITWENTYFIVE}{%
\begin{bkfigure}[0.52\linewidth]
\begin{tikzpicture}[scale=1.0]
  % proportions from the book: apex at 0.775 of the base, height 0.667 of the base
  \coordinate (A) at (0,0);
  \coordinate (B) at (2.95,0);
  \coordinate (C) at (2.28625,1.96765);
  % incentre = (a A + b B + c C)/(a+b+c), a=|BC|, b=|CA|, c=|AB|
  % |BC| = 2.077003, |CA| = 3.016421, |AB| = 2.95  ->  I = (1.944900, 0.721694)
  \coordinate (I) at (1.944900,0.721694);
  \draw[given] (A) -- (B) -- (C) -- cycle;
  \draw[aux] (A) -- (intersection of A--I and B--C);
  \draw[aux] (B) -- (intersection of B--I and C--A);
  \draw[aux] (C) -- (intersection of C--I and A--B);
  \draw[fig] ($(I)+(-0.08,-0.08)$) -- ($(I)+(0.08,0.08)$);
  \draw[fig] ($(I)+(-0.08,0.08)$) -- ($(I)+(0.08,-0.08)$);
  % 227.52 deg bisects the gap between the bisector rays at 200.36 and 254.68
  \node[vlab] at ($(I)+(227.52:0.58)$) {$O$};
\end{tikzpicture}
\figcap{1--25}
\end{bkfigure}}

%--------------------------------------------------------------- Figs 1-26, 1-27
% 1-26: l parallel to m (same direction vector), transversal n; angles 1 and 2
%       are corresponding angles, hence equal.
% 1-27: copying angle 2 at P produces a line through P parallel to l.
\newcommand{\FIGITWENTYSIXTWENTYSEVEN}{%
\begin{bkfigure}[\linewidth]
\begin{minipage}[t]{0.48\linewidth}\centering
\begin{tikzpicture}[scale=1.0]
  % l and m share the direction (3.2, 0.9); n is the transversal
  \coordinate (l1) at (0.10,0.95);  \coordinate (l2) at (3.30,1.85);
  \coordinate (m1) at (0.10,0.30);  \coordinate (m2) at (3.30,1.20);
  \coordinate (n1) at (0.55,2.05);  \coordinate (n2) at (2.05,-0.10);
  \draw[given] (l1) -- (l2);
  \draw[given] (m1) -- (m2);
  \draw[given] (n1) -- (n2);
  \coordinate (P1) at (intersection of l1--l2 and n1--n2);
  \coordinate (P2) at (intersection of m1--m2 and n1--n2);
  % the book marks the two ALTERNATE INTERIOR angles: 1 below l and right of n,
  % 2 above m and left of n. Each label sits on the bisector of its own sector
  % (l runs at 15.71 deg, n descends at 304.90 deg).
  \node[slab] at ($(P1)+(340.30:0.42)$) {1};
  \node[slab] at ($(P2)+(160.30:0.42)$) {2};
  \node[vlab, right] at (l2) {$l$};
  \node[vlab, right] at (m2) {$m$};
  \node[vlab, below right, xshift=-1pt, yshift=-0pt] at (n2) {$n$};
\end{tikzpicture}
\figcap{1--26}
\end{minipage}\hfill
\begin{minipage}[t]{0.48\linewidth}\centering
\begin{tikzpicture}[scale=1.0]
  \coordinate (l1) at (0,1.35);  \coordinate (l2) at (3.10,1.35);
  \coordinate (n1) at (0.72,0);  \coordinate (n2) at (2.02,2.10);
  \draw[given] (l1) -- (l2);
  \draw[given] (n1) -- (n2);
  \coordinate (Q) at (intersection of l1--l2 and n1--n2);
  \coordinate (P) at ($(n1)!0.30!(n2)$);
  \draw[aux] ($(P)+(-0.66,0)$) -- ($(P)+(1.22,0)$);   % the copied parallel
  \fill (P) circle (1.2pt);
  % as in the book: 2 is the lower-left angle at Q, 1 the upper-right angle at P
  % (alternate interior angles). n rises at 58.24 deg.
  \node[slab] at ($(Q)+(209.12:0.55)$) {2};
  \node[slab] at ($(P)+(29.12:0.55)$) {1};
  \node[vlab] at ($(P)+(299.12:0.34)$) {$P$};
  \node[vlab, left] at (l1) {$l$};
  \node[vlab, above] at (n2) {$n$};
\end{tikzpicture}
\figcap{1--27}
\end{minipage}
\end{bkfigure}}

%--------------------------------------------------------------- Fig 1-28
% Right triangle: right angle at C, hypotenuse c, legs a and b.
\newcommand{\FIGITWENTYEIGHT}{%
\begin{bkfigure}[0.52\linewidth]
\begin{tikzpicture}[scale=1.0]
  \coordinate (A) at (0,0);
  \coordinate (C) at (2.45,0);
  \coordinate (B) at (2.45,1.70);
  \draw[given] (A) -- (C) -- (B) -- cycle;
  \draw[fig] (2.29,0) -- (2.29,0.16) -- (2.45,0.16);
  \node[vlab, left]  at (A) {$A$};
  \node[vlab, below right, xshift=-1pt, yshift=-0pt] at (C) {$C$};
  \node[vlab, above] at (B) {$B$};
  \node[vlab, below] at (1.22,0) {$b$};
  \node[vlab, right] at (2.45,0.85) {$a$};
  \node[vlab, above left, xshift=-0pt, yshift=-1pt] at (1.22,0.85) {$c$};
\end{tikzpicture}
\figcap{1--28}
\end{bkfigure}}

% figures03.tex -- Chapter 3 ("Concerning Lines") figures redrawn in TikZ.
%
% Fig 3-1 / 3-2 geometry is DIGITIZED from scan04.pdf p4 (book p.44): the six
% straight lines of Fig 3-1 were recovered by a Hough fit over the thresholded
% scan crop, and Fig 3-2 is the same six lines carrying the printed hand-wobble
% (measured perpendicular deviation 8-15 px at 170 dpi ~= 0.04-0.075 cm).
% Remaining figures are reconstructions at printed proportion; constrained
% points are placed by construction (never by eye) -- see tools/constraints/ch03.py.

% ============================================================ Fig 3-1, 3-2
% Six lines. L1 baseline carries A,B; L3 carries A',B'.
% AB and A'B' are drawn the SAME length in 3-1 (the text's whole point).
\newcommand{\FIGIIIONETWO}{\begin{bkfigure}[\linewidth]
\begin{minipage}{0.48\linewidth}\centering
\begin{tikzpicture}[scale=0.62]
  % digitized endpoints (scan px / 200, y flipped)
  \coordinate (p1) at (0.18,0.985);  \coordinate (q1) at (8.69,0.80);
  \coordinate (p2) at (0.905,3.195); \coordinate (q2) at (7.945,2.315);
  \coordinate (p3) at (4.06,4.075);  \coordinate (q3) at (7.60,0.81);
  \coordinate (p4) at (4.64,4.08);   \coordinate (q4) at (1.81,0.625);
  \coordinate (p5) at (2.005,3.565); \coordinate (q5) at (3.25,0.005);
  \coordinate (p6) at (7.35,3.64);   \coordinate (q6) at (5.74,0.765);
  \foreach \a/\b in {p1/q1,p2/q2,p3/q3,p4/q4,p5/q5,p6/q6} \draw[fig] (\a)--(\b);
  % A,B on the baseline; A',B' on L3 -- equal separations by construction
  \coordinate (A) at ($(p1)!0.085!(q1)$);  \coordinate (B) at ($(p1)!0.150!(q1)$);
  % A'B' drawn the SAME LENGTH as AB: t-step scaled by |L1|/|L3|.
  % Both slid down L3 past its crossing with L6 (t=0.658) so the labels have
  % clear air on the right-hand side of the line.
  \coordinate (Ap) at ($(p3)!0.83489!(q3)$); \coordinate (Bp) at ($(p3)!0.72000!(q3)$);
  \foreach \p in {A,B,Ap,Bp} {\dt{\p}}
  \node[vlab,above] at (A) {$A$};   \node[vlab,above] at (B) {$B$};
  \node[vlab,right,yshift=4pt] at (Ap) {$A'$};
  \node[vlab] at ($(Bp)+(47.3:0.45)$) {$B'$};
\end{tikzpicture}
\figcap{3--1}
\end{minipage}\hfill
\begin{minipage}{0.48\linewidth}\centering
\begin{tikzpicture}[scale=0.62]
  % same six lines seen "through uneven glass": wobble traced from the scan
  \draw[fig] plot[smooth,tension=0.85] coordinates
    {(0.245,0.745)(1.30,0.79)(2.35,0.755)(3.40,0.815)(4.45,0.80)(5.50,0.86)(6.45,0.90)(7.40,0.945)};
  \draw[fig] plot[smooth,tension=0.85] coordinates
    {(1.005,2.98)(1.90,2.90)(2.80,2.90)(3.70,2.79)(4.60,2.68)(5.50,2.62)(6.45,2.46)(7.39,2.355)};
  \draw[fig] plot[smooth,tension=0.85] coordinates
    {(4.05,3.87)(4.52,3.42)(4.95,2.98)(5.40,2.55)(5.85,2.09)(6.30,1.68)(6.75,1.30)(7.185,0.905)};
  \draw[fig] plot[smooth,tension=0.85] coordinates
    {(4.515,3.98)(4.14,3.53)(3.73,3.04)(3.38,2.55)(2.99,2.03)(2.63,1.53)(2.29,1.03)(1.945,0.515)};
  \draw[fig] plot[smooth,tension=0.85] coordinates
    {(1.97,3.355)(2.11,2.92)(2.28,2.49)(2.40,2.06)(2.56,1.64)(2.72,1.21)(2.93,0.79)(3.09,0.37)};
  \draw[fig] plot[smooth,tension=0.85] coordinates
    {(6.835,3.555)(6.62,3.11)(6.42,2.66)(6.24,2.20)(6.02,1.75)(5.83,1.30)(5.63,0.85)(5.43,0.40)};
  % same two pairs as Fig 3-1, at the same fractions along their curves
  \coordinate (A) at (0.83,0.775);  \coordinate (B) at (1.28,0.79);
  \coordinate (Ap) at (6.681,1.359); \coordinate (Bp) at (6.319,1.664);
  \foreach \p in {A,B,Ap,Bp} {\dt{\p}}
  \node[vlab,above] at (A) {$A$};   \node[vlab,above] at (B) {$B$};
  \node[vlab,right,yshift=4pt] at (Ap) {$A'$};
  \node[vlab] at ($(Bp)+(49.8:0.50)$) {$B'$};
\end{tikzpicture}
\figcap{3--2}
\end{minipage}
\end{bkfigure}}

% ======================================================= Fig 3-3, 3-4, 3-5
\newcommand{\FIGIIITHREEFOURFIVE}{\begin{bkfigure}[\linewidth]
\begin{minipage}{0.30\linewidth}\centering
\begin{tikzpicture}[scale=0.62]
  \draw[fig] (0,0) -- (5.0,0.55); \node[vlab,right] at (5.0,0.55) {$l$};
  \coordinate (P) at (2.0,0.22); \dt{P}
  \node[vlab,above] at (P) {$P$};
\end{tikzpicture}
\figcap{3--3}
\end{minipage}\hfill
\begin{minipage}{0.30\linewidth}\centering
\begin{tikzpicture}[scale=0.62]
  \coordinate (la) at (0,0.10); \coordinate (lb) at (4.8,1.42);
  \coordinate (ma) at (0,1.55); \coordinate (mb) at (4.8,0.02);
  \draw[fig] (la) -- (lb); \node[vlab,right] at (lb) {$l$};
  \draw[fig] (ma) -- (mb); \node[vlab,right] at (mb) {$m$};
  \coordinate (P) at (intersection of la--lb and ma--mb);
  \dt{P} \node[vlab,below] at (P) {$P$};
\end{tikzpicture}
\figcap{3--4}
\end{minipage}\hfill
\begin{minipage}{0.30\linewidth}\centering
\begin{tikzpicture}[scale=0.62]
  \coordinate (A) at (2.35,2.55); \coordinate (B) at (0.75,0.32);
  \coordinate (C) at (4.25,0.32);
  \draw[fig] ($(A)!-0.16!(B)$) -- ($(B)!-0.16!(A)$);
  \draw[fig] ($(A)!-0.16!(C)$) -- ($(C)!-0.16!(A)$);
  \draw[fig] ($(B)!-0.20!(C)$) -- ($(C)!-0.20!(B)$);
  % labels on the bisector of the free wedge at each vertex (the book puts
  % A above the crossing, B and C just inside the base)
  \node[vlab] at ($(A)+(92.5:0.75)$) {$A$};
  \node[vlab] at ($(B)+(27:0.90)$)   {$B$};
  \node[vlab] at ($(C)+(155:0.95)$)  {$C$};
\end{tikzpicture}
\figcap{3--5}
\end{minipage}
\end{bkfigure}}

% ============================================================ Fig 3-6, 3-7
% 3-6: two CURVED "lines" l and m meeting at A and B (the impossible picture).
\newcommand{\FIGIIISIXSEVEN}{\begin{bkfigure}[\linewidth]
\begin{minipage}{0.48\linewidth}\centering
\begin{tikzpicture}[scale=0.714]
  \coordinate (A) at (0.45,0.55); \coordinate (B) at (4.65,1.15);
  \draw[fig] (A) .. controls (1.7,1.55) and (3.3,1.75) .. (B);
  \draw[fig] (A) .. controls (1.7,0.10) and (3.3,0.35) .. (B);
  % BOTH curves run on past BOTH meeting points, so A and B each read as a
  % crossing (as in the book), not as a cusp
  \draw[fig] ($(A)!-0.09!(1.7,1.55)$) -- (A);
  \draw[fig] ($(A)!-0.09!(1.7,0.10)$) -- (A);
  \draw[fig] (B) -- ($(B)!-0.09!(3.3,1.75)$);
  \draw[fig] (B) -- ($(B)!-0.09!(3.3,0.35)$);
  \dt{A} \dt{B}
  \node[vlab,left] at (A) {$A$}; \node[vlab] at ($(B)+(100:0.50)$) {$B$};
  \node[vlab,above right] at (3.92,1.50) {$l$};
  \node[vlab,below right] at (4.55,0.88) {$m$};
\end{tikzpicture}
\figcap{3--6}
\end{minipage}\hfill
\begin{minipage}{0.48\linewidth}\centering
\begin{tikzpicture}[scale=0.714]
  \coordinate (C) at (2.55,2.60); \coordinate (A) at (0.55,0.35);
  \coordinate (B) at (4.35,0.62);
  \draw[fig] ($(A)!-0.18!(B)$) -- ($(B)!-0.16!(A)$);
  \draw[fig] ($(C)!-0.20!(A)$) -- ($(A)!-0.10!(C)$);
  \draw[fig] ($(C)!-0.20!(B)$) -- ($(B)!-0.10!(C)$);
  % m = line BC (extends up-LEFT above C); n = line CA (extends up-RIGHT).
  % C sits in the wedge ABOVE the crossing, between m and n, as in the book.
  \node[vlab] at ($(C)+(90:0.72)$) {$C$};
  \node[vlab,above left] at (A) {$A$}; \node[vlab,below right] at (B) {$B$};
  \node[vlab,left]  at ($(C)!-0.20!(B)$) {$m$};
  \node[vlab,right] at ($(C)!-0.20!(A)$) {$n$};
  \node[vlab,right] at ($(B)!-0.16!(A)$) {$l$};
\end{tikzpicture}
\figcap{3--7}
\end{minipage}
\end{bkfigure}}

% ============================================================ Fig 3-8, 3-9
\newcommand{\FIGIIIEIGHT}{\begin{bkfigure}[0.62\linewidth]
\begin{tikzpicture}[scale=0.806]
  \coordinate (P) at (0.40,0.35); \coordinate (Q) at (4.55,1.35);
  \draw[fig] (P) .. controls (1.7,1.35) and (3.2,1.70) .. (Q);
  \draw[fig] (P) .. controls (1.8,0.05) and (3.4,0.55) .. (Q);
  \draw[fig] ($(P)!-0.10!(1.7,1.35)$) -- (P);
  \draw[fig] ($(P)!-0.10!(1.8,0.05)$) -- (P);
  \draw[fig] (Q) -- ($(Q)!-0.10!(3.2,1.70)$);
  \draw[fig] (Q) -- ($(Q)!-0.10!(3.4,0.55)$);
  \dt{P} \dt{Q}
  \node[vlab,left] at (P) {$P$}; \node[vlab,above left] at (Q) {$Q$};
  \node[vlab,above right] at (4.30,1.60) {$a$};
  \node[vlab,below right] at (4.55,1.12) {$b$};
\end{tikzpicture}
\figcap{3--8}
\end{bkfigure}}

% 3-9: figure-eight -- two curved "lines" l and m crossing twice, A..E marked.
\newcommand{\FIGIIININE}{\begin{bkfigure}[0.62\linewidth]
\begin{tikzpicture}[scale=0.806]
  \coordinate (A) at (0.30,1.05); \coordinate (B) at (2.45,1.05);
  \coordinate (C) at (4.60,0.92);
  % D and E sit ON the right lobe: Bezier midpoints of the two arcs B..C
  \coordinate (D) at (3.5625,1.74625);
  \coordinate (E) at (3.5625,0.32125);
  \draw[fig] (A) .. controls (0.95,2.05) and (1.85,2.00) .. (B)
                 .. controls (3.05,2.05) and (4.10,1.95) .. (C);
  \draw[fig] (A) .. controls (0.95,0.05) and (1.85,0.10) .. (B)
                 .. controls (3.05,0.05) and (4.10,0.15) .. (C);
  \draw[fig] ($(A)!-0.12!(0.95,2.05)$) -- (A);
  \draw[fig] ($(A)!-0.12!(0.95,0.05)$) -- (A);
  \draw[fig] (C) -- ($(C)!-0.14!(4.10,1.95)$);
  \draw[fig] (C) -- ($(C)!-0.14!(4.10,0.15)$);
  \foreach \p in {A,B,C,D,E} {\dt{\p}}
  \node[vlab,above left] at (A) {$A$}; \node[vlab] at ($(B)+(90:0.58)$) {$B$};
  \node[vlab,below] at (C) {$C$};
  \node[vlab,above] at (D) {$D$}; \node[vlab,below] at (E) {$E$};
  \node[vlab,above right] at (4.80,1.36) {$l$};
  \node[vlab,right] at (5.02,0.62) {$m$};
\end{tikzpicture}
\figcap{3--9}
\end{bkfigure}}

% =============================================================== Fig 3-10
\newcommand{\FIGIIITEN}{\begin{bkfigure}[0.55\linewidth]
\begin{tikzpicture}[scale=0.806]
  \coordinate (A) at (2.45,2.85); \coordinate (B) at (0.35,0.30);
  \coordinate (C) at (4.55,0.55);
  \draw[fig] (A)--(B)--(C)--cycle;
  % D on AB, E on AC; the drawn line DE extended both ways
  \coordinate (D) at ($(A)!0.52!(B)$); \coordinate (E) at ($(A)!0.52!(C)$);
  \draw[fig] ($(D)!-0.30!(E)$) -- ($(E)!-0.22!(D)$);
  \foreach \p in {A,B,C,D,E} {\dt{\p}}
  \node[vlab,above] at (A) {$A$}; \node[vlab,below left] at (B) {$B$};
  \node[vlab,right] at (C) {$C$};
  \node[vlab,above left] at (D) {$D$}; \node[vlab,above right] at (E) {$E$};
\end{tikzpicture}
\figcap{3--10}
\end{bkfigure}}

% ========================================================= Fig 3-11, 3-12
\newcommand{\FIGIIIELEVENTWELVE}{\begin{bkfigure}[\linewidth]
\begin{minipage}{0.48\linewidth}\centering
\begin{tikzpicture}[scale=0.762]
  \draw[fig] (0,0) -- (4.6,0.62);
  \coordinate (A) at (0.85,0.115); \coordinate (B) at (2.95,0.398);
  \coordinate (C) at (3.55,0.478);
  \foreach \p in {A,B,C} {\dt{\p}}
  \node[vlab,above] at (A) {$A$}; \node[vlab,above] at (B) {$B$};
  \node[vlab,above] at (C) {$C$};
\end{tikzpicture}
\figcap{3--11}
\end{minipage}\hfill
\begin{minipage}{0.48\linewidth}\centering
\begin{tikzpicture}[scale=0.762]
  \draw[fig] (0,0) -- (4.6,0.62);
  \coordinate (A) at (0.60,0.081); \coordinate (D) at (2.05,0.276);
  \coordinate (B) at (2.60,0.350); \coordinate (C) at (4.05,0.546);
  \foreach \p in {A,D,B,C} {\dt{\p}}
  \node[vlab,above] at (A) {$A$}; \node[vlab,above] at (D) {$D$};
  \node[vlab,above] at (B) {$B$}; \node[vlab,above] at (C) {$C$};
\end{tikzpicture}
\figcap{3--12}
\end{minipage}
\end{bkfigure}}

% ========================================================= Fig 3-13, 3-14
\newcommand{\FIGIIITHIRTEENFOURTEEN}{\begin{bkfigure}[\linewidth]
\begin{minipage}{0.48\linewidth}\centering
\begin{tikzpicture}[scale=0.741]
  \draw[fig] (0,0) -- (4.6,1.55); \node[vlab,right] at (4.6,1.55) {$l$};
  \coordinate (P) at (0.55,1.95); \coordinate (Q) at (2.25,1.95);
  \draw[fig] (P) -- (Q);
  \dt{P} \dt{Q}
  \node[vlab,above] at (P) {$P$}; \node[vlab,above] at (Q) {$Q$};
  \node[slab,below] at ($(P)!0.5!(Q)$) {Same side};
\end{tikzpicture}
\figcap{3--13}
\end{minipage}\hfill
\begin{minipage}{0.48\linewidth}\centering
\begin{tikzpicture}[scale=0.741]
  \coordinate (P) at (1.25,2.10); \coordinate (Q) at (2.85,0.20);
  \coordinate (la) at (0.30,0.05); \coordinate (lb) at (4.30,2.30);
  \draw[fig] (la) -- (lb); \node[vlab,right] at (lb) {$l$};
  \draw[fig] (P) -- (Q);
  % R is where PQ meets l -- by construction, never by eye
  \coordinate (R) at (intersection of la--lb and P--Q);
  \foreach \p in {P,Q,R} {\dt{\p}}
  \node[vlab,above] at (P) {$P$}; \node[vlab,below] at (Q) {$Q$};
  \node[vlab] at ($(R)+(170:0.80)$) {$R$};
  \node[slab,right,align=left] at (3.05,1.14) {Opposite\\sides};
\end{tikzpicture}
\figcap{3--14}
\end{minipage}
\end{bkfigure}}

% =============================================================== Fig 3-15
% Halftone-style line illustration of three men standing in a line: PLATE.
\newcommand{\FIGIIIFIFTEEN}{\begin{bkfigure}[0.72\linewidth]
\setlength{\fboxsep}{10pt}%
\fbox{\begin{minipage}{0.86\linewidth}\centering\footnotesize
\textsc{Plate.}\ Illustration: three men standing in a line, marked
$A$, $B$, $C$, with the sight line from $A$ blocked at $B$.\\[2pt]
\emph{Original artwork -- not redrawn.}
\end{minipage}}
\figcap{3--15}
\end{bkfigure}}

% ========================================================= Fig 3-16, 3-17
\newcommand{\FIGIIISIXTEENSEVENTEEN}{\begin{bkfigure}[\linewidth]
\begin{minipage}{0.48\linewidth}\centering
\begin{tikzpicture}[scale=0.768]
  \coordinate (A) at (0.20,0.15); \coordinate (B) at (3.55,2.20);
  \draw[fig] (A) -- (B);
  \coordinate (P) at ($(A)!0.46!(B)$);
  \foreach \p in {A,B,P} {\dt{\p}}
  \node[vlab,below left] at (A) {$A$}; \node[vlab,above] at (B) {$B$};
  \node[vlab,above left] at (P) {$P$};
\end{tikzpicture}
\figcap{3--16}
\end{minipage}\hfill
\begin{minipage}{0.48\linewidth}\centering
\begin{tikzpicture}[scale=0.768]
  \coordinate (A) at (2.05,2.45); \coordinate (B) at (0.25,0.20);
  \coordinate (C) at (3.75,0.05);
  \draw[fig] (A)--(B)--(C)--cycle;
  \node[vlab,above] at (A) {$A$}; \node[vlab,left] at (B) {$B$};
  \node[vlab,below right] at (C) {$C$};
\end{tikzpicture}
\figcap{3--17}
\end{minipage}
\end{bkfigure}}

% =============================================================== Fig 3-18
% l meets side AB and misses BC; C is the apex. l passes through no vertex.
\newcommand{\FIGIIIEIGHTEEN}{\begin{bkfigure}[0.55\linewidth]
\begin{tikzpicture}[scale=0.806]
  \coordinate (C) at (1.85,2.60); \coordinate (B) at (0.20,0.30);
  \coordinate (A) at (4.30,0.35);
  \draw[fig] (C)--(B)--(A)--cycle;
  % l: enters through AB (the base) and leaves through AC
  \coordinate (X) at ($(B)!0.42!(A)$);      % on AB (base)
  \coordinate (Y) at ($(C)!0.52!(A)$);      % on AC
  % l runs from just below the base up past AC to about C's height, as printed
  \coordinate (ltop) at ($(Y)!-1.05!(X)$);
  \draw[fig] ($(X)!-0.35!(Y)$) -- (ltop);
  \node[vlab,above] at (C) {$C$}; \node[vlab,left] at (B) {$B$};
  \node[vlab,right] at (A) {$A$};
  \node[vlab,above right] at (ltop) {$l$};
\end{tikzpicture}
\figcap{3--18}
\end{bkfigure}}

% =============================================================== Fig 3-19
\newcommand{\FIGIIININETEEN}{\begin{bkfigure}[0.62\linewidth]
\begin{tikzpicture}[scale=0.806]
  \draw[fig] (0,0) -- (5.4,0); \node[vlab,right] at (5.4,0) {$l$};
  \coordinate (O) at (2.55,0); \dt{O}
  \node[vlab,below] at (O) {$O$};
  \draw[fig,line width=0.5pt] (0.35,0.30) -- (0.35,0.46) -- (2.45,0.46) -- (2.45,0.30);
  \draw[fig,line width=0.5pt] (2.65,0.30) -- (2.65,0.46) -- (4.95,0.46) -- (4.95,0.30);
  \node[slab,above] at (1.40,0.46) {One side};
  \node[slab,above] at (3.80,0.46) {Other side};
\end{tikzpicture}
\figcap{3--19}
\end{bkfigure}}

% ========================================================= Fig 3-20, 3-21
\newcommand{\FIGIIITWENTYTWENTYONE}{\begin{bkfigure}[\linewidth]
\begin{minipage}{0.48\linewidth}\centering
\begin{tikzpicture}[scale=0.751]
  \coordinate (la) at (0,0.55);     \coordinate (lb) at (4.35,0.55);
  \coordinate (ma) at (1.05,-0.45); \coordinate (mb) at (2.95,2.40);
  \draw[fig] (la) -- (lb); \node[vlab,right] at (lb) {$l$};
  \draw[fig] (ma) -- (mb); \node[vlab,above] at (mb) {$m$};
  % O is the point l and m have in common; A is a point of m, not of l
  \coordinate (O) at (intersection of la--lb and ma--mb);
  \coordinate (A) at ($(ma)!0.79!(mb)$);
  \dt{O} \dt{A}
  \node[vlab] at ($(O)+(298:0.55)$) {$O$}; \node[vlab,right] at (A) {$A$};
\end{tikzpicture}
\figcap{3--20}
\end{minipage}\hfill
\begin{minipage}{0.48\linewidth}\centering
\begin{tikzpicture}[scale=0.751]
  \coordinate (la) at (0,0.55);     \coordinate (lb) at (4.35,0.55);
  \coordinate (ma) at (1.05,-0.45); \coordinate (mb) at (2.95,2.40);
  \draw[fig] (la) -- (lb); \node[vlab,right] at (lb) {$l$};
  \draw[fig] (ma) -- (mb); \node[vlab,above] at (mb) {$m$};
  \coordinate (O) at (intersection of la--lb and ma--mb);
  \coordinate (A) at ($(ma)!0.79!(mb)$);
  % P and Q are points of l on opposite sides of O:  POQ
  \coordinate (Q) at ($(la)!0.138!(lb)$); \coordinate (P) at ($(la)!0.655!(lb)$);
  \foreach \p in {O,A,P,Q} {\dt{\p}}
  \node[vlab] at ($(O)+(298:0.55)$) {$O$}; \node[vlab,right] at (A) {$A$};
  \node[vlab,below] at (Q) {$Q$}; \node[vlab,below] at (P) {$P$};
\end{tikzpicture}
\figcap{3--21}
\end{minipage}
\end{bkfigure}}

% ========================================================= Fig 3-22, 3-23
\newcommand{\FIGIIITWENTYTWOTWENTYTHREE}{\begin{bkfigure}[\linewidth]
\begin{minipage}{0.48\linewidth}\centering
\begin{tikzpicture}[scale=0.755]
  \coordinate (O) at (1.05,0.35); \coordinate (lend) at (4.35,1.32);
  % the part of l that is NOT in the ray, dashed: the exact opposite ray
  \draw[fig,dashed,dash pattern=on 3pt off 2.4pt] ($(O)!-0.318!(lend)$) -- (O);
  \draw[fig] (O) -- (lend); \node[vlab,right] at (lend) {$l$};
  \coordinate (P) at ($(O)!0.394!(lend)$);
  \dt{O} \dt{P}
  \node[vlab,above] at (O) {$O$}; \node[vlab,above] at (P) {$P$};
  \node[slab,below] at (1.90,0.52) {Ray};
\end{tikzpicture}
\figcap{3--22}
\end{minipage}\hfill
\begin{minipage}{0.48\linewidth}\centering
\begin{tikzpicture}[scale=0.755]
  \coordinate (P) at (2.05,0.30);
  \draw[fig] (0.15,0.30) -- (4.35,0.30); \node[vlab,right] at (4.35,0.30) {$l$};
  \draw[fig] (P) -- (3.35,2.55); \node[vlab,above] at (3.35,2.55) {$r$};
  \dt{P} \node[vlab,above left] at (P) {$P$};
\end{tikzpicture}
\figcap{3--23}
\end{minipage}
\end{bkfigure}}

% ========================================================= Fig 3-24, 3-25
\newcommand{\FIGIIITWENTYFOURTWENTYFIVE}{\begin{bkfigure}[\linewidth]
\begin{minipage}{0.52\linewidth}\centering
\begin{tikzpicture}[scale=0.780]
  % (i) angle at O, upper left
  \coordinate (O1) at (0.30,2.05);
  \draw[fig] (O1) -- (1.85,2.95); \draw[fig] (O1) -- (1.85,2.05);
  \dt{O1} \node[vlab,left] at (O1) {$O$};
  % (ii) angle at O, lower left
  \coordinate (O2) at (0.62,1.35);
  \draw[fig] (O2) -- (2.05,1.05); \draw[fig] (O2) -- (-0.15,0.20);
  \dt{O2} \node[vlab,left] at (O2) {$O$};
  % (iii) angle AOB, right
  \coordinate (O3) at (3.55,2.20); \coordinate (A3) at (2.55,3.10);
  \coordinate (B3) at (2.85,1.25);
  \draw[fig] (O3) -- (A3); \draw[fig] (O3) -- (B3);
  \dt{O3} \dt{A3} \dt{B3}
  \node[vlab,right] at (O3) {$O$}; \node[vlab,above] at (A3) {$A$};
  \node[vlab,below] at (B3) {$B$};
  % (iv) straight angle
  \draw[fig] (0.85,0.00) -- (3.35,0.00);
  \coordinate (O4) at (2.05,0.00); \dt{O4}
  \node[vlab,below] at (O4) {$O$};
\end{tikzpicture}
\figcap{3--24}
\end{minipage}\hfill
\begin{minipage}{0.44\linewidth}\centering
\begin{tikzpicture}[scale=0.780]
  % (a) the arrowheads are on the ENDS OF THE ARC, pointing at the two sides;
  % the sides themselves carry no arrows.  Verified on scans/scan04.pdf p18
  % (book p.54), where both arrowheads are plainly on the arc.
  \coordinate (Oa) at (0.15,2.55);
  \draw[fig] (Oa) -- (1.90,3.15);
  \draw[fig] (Oa) -- (1.75,2.10);
  % every arc is centred on its vertex and stops exactly on the two sides
  \draw[fig,line width=0.5pt,<->,>=stealth]
    ($(Oa)+(18.93:1.05)$) arc (18.93:-15.71:1.05);
  \node[slab] at ($(Oa)+(1.61:1.30)$) {1};
  \node[slab,right] at (2.25,2.55) {(a)};
  % (b) plain sides
  \coordinate (Ob) at (0.15,1.35);
  \draw[fig] (Ob) -- (1.90,1.95); \draw[fig] (Ob) -- (1.75,0.90);
  \draw[fig,line width=0.5pt] ($(Ob)+(18.93:1.05)$) arc (18.93:-15.71:1.05);
  \node[slab] at ($(Ob)+(1.61:1.30)$) {1};
  \node[slab,right] at (2.25,1.35) {(b)};
  % (c) obtuse
  \coordinate (Oc) at (0.35,0.10);
  \draw[fig] (Oc) -- (2.10,0.10); \draw[fig] (Oc) -- (-0.25,0.95);
  \draw[fig,line width=0.5pt] ($(Oc)+(0:0.90)$) arc (0:125.18:0.90);
  \node[slab] at ($(Oc)+(62.59:0.58)$) {2};
  \node[slab,right] at (2.25,0.10) {(c)};
\end{tikzpicture}
\figcap{3--25}
\end{minipage}
\end{bkfigure}}

% =============================================================== Fig 3-26
% Two lines l, m crossing at O; four rays l1,l2,m1,m2; angles 1,2,3,4.
\newcommand{\FIGIIITWENTYSIX}{\begin{bkfigure}[0.66\linewidth]
\begin{tikzpicture}[scale=0.910]
  \coordinate (L1) at (4.85,2.05); \coordinate (L2) at (0.35,0.15);
  \coordinate (M1) at (4.85,0.35); \coordinate (M2) at (0.35,1.95);
  \draw[fig] (L2) -- (L1); \draw[fig] (M2) -- (M1);
  \coordinate (O) at (intersection of L2--L1 and M2--M1);
  \dt{O}
  \node[vlab,right] at (L1) {$l_1$}; \node[vlab,left]  at (L2) {$l$};
  \node[vlab,right] at (M1) {$m_1$}; \node[vlab,left]  at (M2) {$m$};
  % ray labels offset perpendicular to their own line, so they sit beside it
  \node[vlab] at ($($(M2)!0.28!(M1)$)+(70.42:0.42)$) {$m_2$};
  \node[vlab] at ($($(L2)!0.30!(L1)$)+(-67.10:0.42)$) {$l_2$};
  % the four arcs are centred on O and stop exactly on the four rays
  \draw[fig,line width=0.5pt] ($(O)+(160.44:0.62)$) arc (160.44:202.89:0.62);
  \draw[fig,line width=0.5pt] ($(O)+(-19.58:0.62)$) arc (-19.58:22.90:0.62);
  \draw[fig,line width=0.5pt] ($(O)+(22.90:0.62)$)  arc (22.90:160.44:0.62);
  \draw[fig,line width=0.5pt] ($(O)+(202.89:0.62)$) arc (202.89:340.42:0.62);
  % angle numerals on each region's bisector, clear of the arcs
  \node[slab] at ($(O)+(181.67:0.95)$) {1};
  \node[slab] at ($(O)+(1.66:0.95)$)   {2};
  \node[slab] at ($(O)+(91.67:0.95)$)  {3};
  \node[slab] at ($(O)+(271.66:0.95)$) {4};
  % the book sets O just below-left of the crossing, inside the ring of arcs;
  % at this type size the label padding will not clear the two lines there, so
  % it stays outside the ring, below-right (see ch03-UNCERTAIN.md)
  \node[vlab] at ($(O)+(303:1.00)$) {$O$};
\end{tikzpicture}
\figcap{3--26}
\end{bkfigure}}

% =============================================================== Fig 3-27
% Three drawings of a pair of angles sharing a vertex and one side.
\newcommand{\FIGIIITWENTYSEVEN}{\begin{bkfigure}[0.82\linewidth]
\begin{tikzpicture}[scale=0.780]
  % (i) narrow pair
  % the narrowest of the three pairs; opened from 18/24 deg to 23/25 so the
  % numerals clear both rays at this type size (the book's own fan is tight)
  \coordinate (O) at (0.15,0.55);
  \draw[fig] (O) -- (1.891,1.724); \draw[fig] (O) -- (2.261,0.960);
  \draw[fig] (O) -- (2.091,0.066);
  \draw[fig,line width=0.5pt] ($(O)+(34:1.25)$) arc (34:11:1.25);
  \draw[fig,line width=0.5pt] ($(O)+(11:1.25)$) arc (11:-14:1.25);
  \node[slab] at ($(O)+(22.5:1.62)$) {1};
  \node[slab] at ($(O)+(-1.5:1.62)$) {2};
\end{tikzpicture}\hspace{1.6em}
\begin{tikzpicture}[scale=0.780]
  % (ii) the book's middle drawing: ONE ray straight up, TWO rays down-left and
  % down-right (an inverted Y), the vertical being the common side.  The two
  % arcs bulge out either side of it -- measured off the printed figure.
  \coordinate (O) at (1.15,1.35);
  \draw[fig] (O) -- (1.15,2.75);            % common side, straight up
  \draw[fig] (O) -- (0.30,0.25);            % noncommon side, down-left
  \draw[fig] (O) -- (2.05,0.25);            % noncommon side, down-right
  \draw[fig,line width=0.5pt] ($(O)+(232.30:0.80)$) arc (232.30:90:0.80);
  \draw[fig,line width=0.5pt] ($(O)+(90:0.80)$) arc (90:-50.71:0.80);
  \node[slab] at ($(O)+(161.15:1.14)$) {1};
  \node[slab] at ($(O)+(19.65:1.14)$)  {2};
\end{tikzpicture}\hspace{1.6em}
\begin{tikzpicture}[scale=0.780]
  % (iii) wide pair over a base line
  \coordinate (O) at (1.25,0.25);
  \draw[fig] (O) -- (0.05,1.35); \draw[fig] (O) -- (1.55,2.05);
  \draw[fig] (O) -- (2.85,0.90);
  \draw[fig,line width=0.5pt] ($(O)+(137.48:1.35)$) arc (137.48:80.54:1.35);
  \draw[fig,line width=0.5pt] ($(O)+(80.54:1.35)$) arc (80.54:22.11:1.35);
  \node[slab] at ($(O)+(109.01:1.80)$) {1};
  \node[slab] at ($(O)+(51.33:1.80)$)  {2};
\end{tikzpicture}
\figcap{3--27}
\end{bkfigure}}

% =============================================================== Fig 3-28
\newcommand{\FIGIIITWENTYEIGHT}{\begin{bkfigure}[0.55\linewidth]
\begin{tikzpicture}[scale=0.858]
  \coordinate (O) at (2.05,0.20);
  \draw[fig] (0.20,0.20) -- (4.15,0.20);
  \draw[fig] (O) -- (3.55,1.70);          % common side, at 45 degrees
  % both arcs run from a side of the angle to the common side, no gap
  \draw[fig,line width=0.5pt] ($(O)+(180:1.20)$) arc (180:45:1.20);
  \node[slab] at ($(O)+(112.5:0.85)$) {1};
  \draw[fig,line width=0.5pt] ($(O)+(0:0.67)$) arc (0:45:0.67);
  \node[slab] at ($(O)+(22.5:1.00)$) {2};
\end{tikzpicture}
\figcap{3--28}
\end{bkfigure}}

% ========================================================= Fig 3-29, 3-30
% 3-29: lines l and m crossing at O; the four regions cut by them, two of
% which are hatched (vertical / horizontal) so their overlap reads doubly hatched.
\newcommand{\FIGIIITWENTYNINETHIRTY}{\begin{bkfigure}[\linewidth]
\begin{minipage}{0.48\linewidth}\centering
\begin{tikzpicture}[scale=0.730]
  % The book's picture: A B C D at the corners of a rectangle, l = AC and
  % m = BD crossing at O.  VERTICAL rules shade the side of l that contains D;
  % HORIZONTAL rules shade the side of m that contains C.  The two families
  % therefore overlap in triangle OCD, and that doubly-shaded wedge IS the
  % interior of angle COD -- which is the whole point of the figure.
  \coordinate (B) at (0.35,2.55); \coordinate (Cc) at (3.95,2.55);
  \coordinate (A) at (0.35,0.35); \coordinate (D) at (3.95,0.35);
  \coordinate (O) at (intersection of A--Cc and B--D);
  % Each rule is drawn at its exact length rather than clipped: a clipped rule
  % still carries its full path in the PDF, so the collision audit would see
  % ink where the picture shows none.
  % D's side of l = triangle A D C: vertical rules from AD up to line AC
  \foreach \x in {0.45,0.55,...,3.90}
    \draw[fig,line width=0.35pt] (\x,0.35) -- (\x,{0.35+(\x-0.35)*0.611111});
  % C's side of m = triangle B C D: horizontal rules from line BD across to DC
  \foreach \y in {0.45,0.55,...,2.50}
    \draw[fig,line width=0.35pt] ({0.35+(2.55-\y)*1.636364},\y) -- (3.95,\y);
  % the two lines: carried past the far corner to take the label, and past the
  % near point as well -- in the book A and B are POINTS ON the lines, not ends
  \draw[fig] ($(A)!-0.07!(Cc)$) -- ($(A)!1.10!(Cc)$);
  \node[vlab,right] at ($(A)!1.10!(Cc)$) {$l$};
  \draw[fig] ($(B)!-0.07!(D)$) -- ($(B)!1.10!(D)$);
  \node[vlab,right] at ($(B)!1.10!(D)$) {$m$};
  \foreach \p in {A,B,Cc,D,O} {\dt{\p}}
  % A and B sit clear of the tails the two lines now carry past them
  \node[vlab] at ($(A)+(150:0.50)$) {$A$};  \node[vlab] at ($(B)+(95:0.58)$) {$B$};
  \node[vlab,above left] at (Cc) {$C$}; \node[vlab,below left] at (D) {$D$};
  \node[vlab] at ($(O)+(180:0.85)$) {$O$};   % the one unshaded wedge, as in the book
\end{tikzpicture}
\figcap{3--29}
\end{minipage}\hfill
\begin{minipage}{0.48\linewidth}\centering
\begin{tikzpicture}[scale=0.730]
  % r_1 lies on l_1 and r_2 on l_2, so each dashed piece is the ray OPPOSITE
  % its solid one: built as a reflection through O, never placed by eye.
  \coordinate (O) at (2.30,1.35);
  \coordinate (R1) at (4.30,2.55); \coordinate (R2) at (4.10,0.10);
  \coordinate (L1) at ($(O)!-0.90!(R1)$);   % opposite ray of r_1
  \coordinate (L2) at ($(O)!-1.00!(R2)$);   % opposite ray of r_2
  \draw[fig] (O) -- (R1); \draw[fig] (O) -- (R2);
  \draw[aux] (O) -- (L2); \draw[aux] (O) -- (L1);
  \node[vlab,right] at (R1) {$r_1$}; \node[vlab,right] at (R2) {$r_2$};
  \node[vlab,left] at (L2) {$l_2$};  \node[vlab,left] at (L1) {$l_1$};
  \dt{O} \node[vlab] at ($(O)+(90:0.62)$) {$O$};
\end{tikzpicture}
\figcap{3--30}
\end{minipage}
\end{bkfigure}}

% =============================================================== Fig 3-31
\newcommand{\FIGIIITHIRTYONE}{\begin{bkfigure}[0.70\linewidth]
\begin{tikzpicture}[scale=0.806]
  \coordinate (O) at (2.05,1.35);
  \coordinate (R1) at (4.05,2.45); \coordinate (R2) at (3.95,0.15);
  \coordinate (L1) at ($(O)!-0.85!(R1)$);   % opposite ray of r_1
  \coordinate (L2) at ($(O)!-0.95!(R2)$);   % opposite ray of r_2
  \draw[fig] (O) -- (R1); \node[vlab,right] at (R1) {$r_1$};
  \draw[fig] (O) -- (R2); \node[vlab,right] at (R2) {$r_2$};
  % The book draws l_1 and l_2 SOLID here, unlike Fig 3-30 where the rays that
  % are not sides of the angle are dashed -- checked on scans/scan04.pdf p20.
  \draw[fig] (O) -- (L2); \node[vlab,left] at (L2) {$l_2$};
  \draw[fig] (O) -- (L1); \node[vlab,left] at (L1) {$l_1$};
  \dt{O} \node[vlab] at ($(O)+(90:0.62)$) {$O$};
  % the separate segment PQ with A between; it runs a little past Q
  \coordinate (P) at (5.15,0.60); \coordinate (Q) at (7.15,0.60);
  \draw[fig] (P) -- ($(P)!1.14!(Q)$);
  \coordinate (Aa) at ($(P)!0.5!(Q)$);
  \foreach \p in {P,Q,Aa} {\dt{\p}}
  \node[vlab,above] at (P) {$P$}; \node[vlab,above] at (Q) {$Q$};
  \node[vlab,above] at (Aa) {$A$};
\end{tikzpicture}
\figcap{3--31}
\end{bkfigure}}

% =============================================================== Fig 3-32
\newcommand{\FIGIIITHIRTYTWO}{\begin{bkfigure}[0.60\linewidth]
\begin{tikzpicture}[scale=0.806]
  \node[vlab,left] at (0.20,2.45) {$r_2'$};
  \node[vlab,left] at (0.25,0.20) {$r_1'$};
  % the two full lines through O; r1' / r2' are the rays opposite r1 / r2
  \coordinate (ta) at (0.20,2.45); \coordinate (tb) at (4.05,0.30);
  \coordinate (ua) at (0.25,0.20); \coordinate (ub) at (4.15,2.35);
  \draw[fig] (ta) -- (tb);
  \draw[fig] (ua) -- (ub);
  \coordinate (O) at (intersection of ta--tb and ua--ub);
  \draw[fig] (O) -- (4.20,1.34); \node[vlab,right] at (4.20,1.34) {$r$};
  \node[vlab,right] at (4.15,2.35) {$r_1$}; \node[vlab,right] at (4.05,0.30) {$r_2$};
  \dt{O} \node[vlab] at ($(O)+(90:0.52)$) {$O$};
  % the two arcs mark angle(r_1,r) and angle(r,r_2): centred on O, ending on
  % the rays (O = (2.2524,1.3039); r_1 at 28.88, r at 1.06, r_2 at -29.18 deg)
  \draw[fig,line width=0.5pt] ($(O)+(28.88:1.40)$) arc (28.88:1.06:1.40);
  \draw[fig,line width=0.5pt] ($(O)+(1.06:1.40)$)  arc (1.06:-29.18:1.40);
\end{tikzpicture}
\figcap{3--32}
\end{bkfigure}}

% =============================================================== Fig 3-33
\newcommand{\FIGIIITHIRTYTHREE}{\begin{bkfigure}[0.72\linewidth]
\begin{tikzpicture}[scale=0.806]
  % l2 : the shallow line, carrying r2 to the right and r2' to the left
  \coordinate (L2R) at (5.75,1.95); \coordinate (L2L) at (0.30,1.75);
  \draw[fig] (L2L) -- (L2R);
  \node[vlab,right] at (L2R) {$l_2$}; \node[vlab,left] at (L2L) {$r_2'$};
  % l1 : the steep line; O is where they meet
  \coordinate (L1T) at (4.15,3.35); \coordinate (L1B) at (1.95,0.15);
  \draw[fig] (L1B) -- (L1T);
  \coordinate (O) at (intersection of L1B--L1T and L2L--L2R);
  \node[vlab,right] at (L1T) {$l_1$};
  \draw[aux] (L1B) -- (1.70,-0.20); \node[vlab,below] at (1.70,-0.20) {$r_1'$};
  \node[vlab,left]  at ($(O)!0.72!(L1T)$) {$r_1$};   % ray r1: above O
  \node[vlab,below] at ($(O)!0.80!(L2R)$) {$r_2$};   % ray r2: right of O
  % P lies in the interior of the angle; Q_3 is across BOTH lines from it, so
  % the book's segment Q_3 P has to meet l_1 in a point R and l_2 in a point S.
  % R and S are taken as intersections -- the proof's whole content.
  \coordinate (P) at (5.15,2.45);  \coordinate (Qc) at (1.85,1.15);
  \draw[fig] (Qc) -- (P);
  \coordinate (R) at (intersection of Qc--P and L1B--L1T);
  \coordinate (S) at (intersection of Qc--P and L2L--L2R);
  \foreach \p in {O,P,S,R} {\dt{\p}}
  \node[vlab] at ($(O)+(118.8:0.55)$) {$O$};
  \node[vlab,above right] at (P) {$P$};
  \node[vlab] at ($(S)+(60:0.75)$) {$S$};
  \node[vlab] at ($(R)+(308.5:0.48)$) {$R$};
  % the five sample positions.  Q_1, Q_3, Q_5 lie in the three regions the
  % lines cut off; Q_2 lies ON the ray r_2' and Q_4 ON the ray r_1' -- that
  % distinction is what Exercises *12-*14 turn on.
  \coordinate (Qa) at (2.55,3.05);                    % Q_1
  \coordinate (Qb) at ($(L2L)!0.22!(L2R)$);           % Q_2, on r_2'
  \coordinate (Qd) at ($(L1B)!0.12!(L1T)$);           % Q_4, on r_1'
  \coordinate (Qe) at (4.30,0.80);                    % Q_5
  \foreach \p in {Qa,Qb,Qc,Qd,Qe} {\dt{\p}}
  \node[vlab,above] at (Qa) {$Q_1$}; \node[vlab,above] at (Qb) {$Q_2$};
  \node[vlab] at ($(Qc)+(235:0.55)$) {$Q_3$};
  \node[vlab,right] at (Qd) {$Q_4$}; \node[vlab,above] at (Qe) {$Q_5$};
\end{tikzpicture}
\figcap{3--33}
\end{bkfigure}}

% ========================================================= Fig 3-34, 3-35
\newcommand{\FIGIIITHIRTYFOURTHIRTYFIVE}{\begin{bkfigure}[\linewidth]
\begin{minipage}{0.48\linewidth}\centering
\begin{tikzpicture}[scale=0.749]
  \coordinate (B) at (1.75,2.35); \coordinate (A) at (0.20,0.25);
  \coordinate (C) at (3.75,0.25);
  \draw[fig] (A)--(B)--(C)--cycle;
  \node[vlab,above] at (B) {$B$}; \node[vlab,left] at (A) {$A$};
  \node[vlab,right] at (C) {$C$};
\end{tikzpicture}
\figcap{3--34}
\end{minipage}\hfill
\begin{minipage}{0.48\linewidth}\centering
\begin{tikzpicture}[scale=0.749]
  \coordinate (B) at (1.55,2.35); \coordinate (A) at (0.20,0.25);
  \coordinate (C) at (3.55,0.25);
  \draw[fig] (A)--(B)--(C)--cycle;
  \coordinate (P) at (1.45,1.05); \coordinate (Q) at (3.55,2.20);
  \draw[fig] (P) -- (Q);
  \dt{P} \dt{Q}
  \node[vlab,above] at (B) {$B$}; \node[vlab,left] at (A) {$A$};
  \node[vlab,right] at (C) {$C$};
  \node[vlab,above,yshift=2pt] at (P) {$P$}; \node[vlab,above right] at (Q) {$Q$};
\end{tikzpicture}
\figcap{3--35}
\end{minipage}
\end{bkfigure}}

% =============================================================== Fig 3-36
\newcommand{\FIGIIITHIRTYSIX}{\begin{bkfigure}[0.50\linewidth]
\begin{tikzpicture}[scale=0.806]
  \coordinate (O) at (0.75,0.35); \coordinate (A) at (3.85,0.15);
  \coordinate (Ot) at (2.55,2.85); \coordinate (At) at (1.55,2.75);
  \draw[fig,dashed,dash pattern=on 3pt off 2.4pt] ($(O)!-0.23!(A)$) -- (O);
  \draw[fig] (O) -- (A);
  \draw[fig] (O) -- (Ot);
  \draw[fig] (A) -- (At);
  % B is where the ray from O meets the other side of the angle at A
  \coordinate (B) at (intersection of O--Ot and A--At);
  \dt{O} \dt{A} \dt{B}
  \node[vlab,above left] at (O) {$O$}; \node[vlab,right] at (A) {$A$};
  \node[vlab] at ($(B)+(93:0.80)$) {$B$};
\end{tikzpicture}
\figcap{3--36}
\end{bkfigure}}

% =============================================================== Fig 3-37
\newcommand{\FIGIIITHIRTYSEVEN}{\begin{bkfigure}[0.52\linewidth]
\begin{tikzpicture}[scale=0.806]
  \coordinate (V) at (1.35,0.65);
  \coordinate (M1) at (3.35,2.75);
  \draw[fig] (V) -- (3.95,0.65); \node[vlab,right] at (3.95,0.65) {$l_1$};
  \draw[fig] (V) -- (M1); \node[vlab,above] at (M1) {$m_1$};
  % l and m are the REST of the two lines: exact opposite rays through V
  \draw[fig,dashed,dash pattern=on 3pt off 2.4pt] (V) -- (0.15,0.65);
  \node[vlab,left] at (0.15,0.65) {$l$};
  \coordinate (Mo) at ($(V)!-0.45!(M1)$);
  \draw[fig,dashed,dash pattern=on 3pt off 2.4pt] (V) -- (Mo);
  \node[vlab,below] at (Mo) {$m$};
  \draw[fig,line width=0.5pt] ($(V)+(0:1.10)$) arc (0:46.40:1.10);
\end{tikzpicture}
\figcap{3--37}
\end{bkfigure}}

% ========================================================= Fig 3-38, 3-39
\newcommand{\FIGIIITHIRTYEIGHTTHIRTYNINE}{\begin{bkfigure}[\linewidth]
\begin{minipage}{0.48\linewidth}\centering
\begin{tikzpicture}[scale=0.675]
  % sides of the angle solid, the opposite rays r_1' and r_2' dashed
  \coordinate (l2L) at (0.30,1.05); \coordinate (l2R) at (4.55,1.25);
  \coordinate (ka) at (1.65,0.15);  \coordinate (kb) at (3.35,2.85);
  \coordinate (O) at (intersection of l2L--l2R and ka--kb);
  \draw[fig] (O) -- (l2R); \node[vlab,right] at (l2R) {$l_2$};
  \draw[aux] (O) -- (l2L); \node[vlab,left]  at (l2L) {$r_2'$};
  \draw[fig] (O) -- (kb);  \node[vlab,right] at (kb) {$l_1$};
  \draw[aux] (O) -- ($(O)!1.18!(ka)$);
  \node[vlab,below] at ($(O)!1.18!(ka)$) {$r_1'$};
  \node[vlab,below] at ($(O)!0.62!(l2R)$) {$r_2$};
  \node[vlab] at ($($(O)!0.78!(kb)$)+(147.8:0.42)$) {$r_1$};
  % Q_1 is on the same side of l_2 as r_1, the far side of l_1 from r_2;
  % R is where segment Q_1 P meets l_1 -- it has to land on the ray r_1.
  \coordinate (P) at (4.05,2.10); \coordinate (Q) at (1.35,2.45);
  \draw[fig] (Q) -- (P);
  \coordinate (R) at (intersection of Q--P and ka--kb);
  \foreach \p in {P,Q,R} {\dt{\p}}
  \node[vlab,right] at (P) {$P$}; \node[vlab,left] at (Q) {$Q_1$};
  \node[vlab] at ($(R)+(299.5:0.55)$) {$R$};
\end{tikzpicture}
\figcap{3--38}
\end{minipage}\hfill
\begin{minipage}{0.48\linewidth}\centering
\begin{tikzpicture}[scale=0.675]
  \coordinate (l2L) at (0.30,1.05); \coordinate (l2R) at (4.55,1.25);
  \coordinate (ka) at (1.65,0.15);  \coordinate (kb) at (3.35,2.85);
  \coordinate (O) at (intersection of l2L--l2R and ka--kb);
  \draw[fig] (O) -- (l2R); \node[vlab,right] at (l2R) {$l_2$};
  \draw[aux] (O) -- (l2L); \node[vlab,left]  at (l2L) {$r_2'$};
  \draw[fig] (O) -- (kb);  \node[vlab,right] at (kb) {$l_1$};
  \draw[aux] (O) -- ($(O)!1.18!(ka)$);
  \node[vlab,below] at ($(O)!1.18!(ka)$) {$r_1'$};
  \node[vlab,below] at ($(O)!0.62!(l2R)$) {$r_2$};
  \node[vlab] at ($($(O)!0.78!(kb)$)+(147.8:0.42)$) {$r_1$};
  % the one thing that separates this figure from 3-38: Q_2 lies ON the ray
  % r_2' itself, not merely on that side of it
  \coordinate (P) at (4.05,2.10); \coordinate (Q) at ($(l2L)!0.16!(l2R)$);
  \draw[fig] (Q) -- (P);
  \coordinate (R) at (intersection of Q--P and ka--kb);
  \foreach \p in {P,Q,R} {\dt{\p}}
  \node[vlab,right] at (P) {$P$}; \node[vlab,above] at (Q) {$Q_2$};
  \node[vlab] at ($(R)+(126:0.60)$) {$R$};
\end{tikzpicture}
\figcap{3--39}
\end{minipage}
\end{bkfigure}}

% figures04.tex -- Chapter 4 figures redrawn in TikZ.
%
% Label clearance, chapter-local (2026-08-17).  Every figure in this chapter
% is a point-on-a-line diagram whose labels sit directly over their point.
% The global `outer sep' was tightened to 2.2pt in the same pass that turned
% the point dots into fixed 3.2pt nodes -- and a 3.2pt dot reaches 1.6pt out
% of its point, so at 2.2pt the label box lands ON the dot it labels.
% check_labels bears this out exactly: the collision count for this chapter
% falls off a cliff at 3.2pt (39 -> 2) and does not improve above it.  3.2pt
% is therefore the tightest clearance that still clears the new dots, and is
% still tighter than the 3.4pt this book used before the dots existed.
% Coordinates are reconstructions; every congruence the printed figure
% asserts holds EXACTLY here (see tools/constraints/ch04.py).
% Self-contained: congruence ticks are drawn as explicit strokes, sized so
% that each prints at about 2.2pt half-length at the figure's own scale
% (Figs. 4-7/4-8/4-9, split out and enlarged, carry the book's own
% tick-to-segment ratio instead and so print nearer 3pt).
\tikzset{vlab/.append style={outer sep=3.2pt},
         slab/.append style={outer sep=3.2pt}}

% ------------------------------------------------------------- Fig 4-1
% Pictorial vignette (boy beside a marked spear stuck in hatched ground).
% Figurative line art, not a geometric diagram -> placeholder, per the
% plate rule. See ch04-UNCERTAIN.md.
\newcommand{\FIGIVONE}{\begin{bkfigure}[0.40\linewidth]
\includegraphics[width=0.40\linewidth]{fig-4-1}%
\figcap{4--1}\end{bkfigure}}

% -------------------------------------------------------- Figs 4-2, 4-3
% 4-2: III-1, the existence postulate.  AB congruent to A'B' (exact).
% 4-3: III-3, addition/subtraction.  AB=A'B' and BC=B'C' (exact: the two
%      lines carry identical parameter fractions over equal lengths).
\newcommand{\FIGIVTWOTHREE}{\begin{bkfigure}[\linewidth]
\begin{minipage}{0.47\linewidth}\centering
\begin{tikzpicture}[scale=0.86]
  \coordinate (A)  at (1.00000,0.42000);
  \coordinate (B)  at (4.20000,1.76000);
  \coordinate (Ap) at (1.55000,-0.53479);
  \coordinate (Bp) at (5.01280,-0.74609);
  \draw[fig] (0.12371,0.05306) -- (5.21465,2.18488);
  \draw[fig] (0.83242,-0.49100) -- (5.62546,-0.78347);
  \foreach \p in {A,B,Ap,Bp} {\dt{\p}}
  \node[vlab,above left]  at (A)  {$A$};
  \node[vlab,above]       at (B)  {$B$};
  \node[vlab,above left]  at (Ap) {$A'$};
  \node[vlab,above]       at (Bp) {$B'$};
  \node[vlab,right]       at (5.21465,2.18488) {$l$};
  \node[vlab,right]       at (5.62546,-0.78347) {$l'$};
\end{tikzpicture}
\figcap{4--2}
\end{minipage}\hfill
\begin{minipage}{0.47\linewidth}\centering
\begin{tikzpicture}[scale=0.86]
  \coordinate (A)  at (0.86972,0.09602);
  \coordinate (B)  at (2.70855,0.29902);
  \coordinate (C)  at (3.41923,0.37748);
  \coordinate (Ap) at (1.09193,-1.07678);
  \coordinate (Bp) at (2.93545,-0.92196);
  \coordinate (Cp) at (3.64794,-0.86212);
  % Left-hand stubs pulled in (2026-08-17): both lines ran 0.47 of AB to the
  % left of A / A', where the plate runs only about 0.22.  New start points
  % are exactly on the same lines (slope unchanged), so the incidence checks
  % in tools/constraints/ch04.py still hold -- their anchors move with them.
  \draw[fig] (0.46519,0.05136) -- (4.96981,0.54866);
  \draw[fig] (0.68636,-1.11084) -- (5.20247,-0.73157);
  \foreach \p in {A,B,C,Ap,Bp,Cp} {\dt{\p}}
  \node[vlab,above]      at (A)  {$A$};
  \node[vlab,above]      at (B)  {$B$};
  \node[vlab,above]      at (C)  {$C$};
  \node[vlab,above]      at (Ap) {$A'$};
  \node[vlab,above]      at (Bp) {$B'$};
  \node[vlab,above]      at (Cp) {$C'$};
  \node[vlab,right]      at (4.96981,0.54866)  {$l$};
  \node[vlab,right]      at (5.20247,-0.73157) {$l'$};
\end{tikzpicture}
\figcap{4--3}
\end{minipage}
\end{bkfigure}}

% ------------------------------------------------------------- Fig 4-4
% Exercise 4-2 #7.  AB=BC, CD=CE=ED (equilateral CDE), EF=EB.
% Construction: A,C,D collinear; E is the equilateral apex on CD below AD;
% B on ray EC beyond C with AB=BC (which forces ABC equilateral); F on ray
% ED beyond D with EF=EB.
\newcommand{\FIGIVFOUR}{\begin{bkfigure}[0.62\linewidth]
\begin{tikzpicture}[scale=0.689]
  \coordinate (A) at (0,0);
  \coordinate (C) at (3.60000,0);
  \coordinate (D) at (6.50000,0);
  \coordinate (E) at (5.05000,-2.51147);
  \coordinate (B) at (1.80000,3.11769);
  \coordinate (F) at (8.30000,3.11769);
  \draw[fig] (A) -- (D);
  \draw[fig] (A) -- (B) -- (E) -- (F);
  \node[vlab,below left]  at (A) {$A$};
  \node[vlab,above]       at (B) {$B$};
  \node[vlab,below left]  at (C) {$C$};
  \node[vlab,above left]  at (D) {$D$};
  \node[vlab,below]       at (E) {$E$};
  \node[vlab,above]       at (F) {$F$};
\end{tikzpicture}
\figcap{4--4}
\end{bkfigure}}

% ------------------------------------------------------------- Fig 4-5
% Exercise 4-2 #15/#16.  (a) two congruent copies: AB=A'B' (double tick),
% AC=A'C' (single tick on the arc).  (b) CA=BD.  (c) AC=BD.
\newcommand{\FIGIVFIVE}{\begin{bkfigure}[\linewidth]
\begin{tikzpicture}[scale=1.040]
  % ---- (a) left copy
  \begin{scope}[shift={(0.60,2.90)}]
    \coordinate (A) at (0,0); \coordinate (B) at (1.95000,0);
    \coordinate (C) at (2.85000,0);
    \draw[fig] (A) -- (C);
    \draw[fig] (A) .. controls +(0,-0.82667) and +(0,-0.82667) .. (C);
    \draw[tick] (0.90000,-0.09700) -- (0.90000,0.09700);
    \draw[tick] (1.05000,-0.09700) -- (1.05000,0.09700);
    \draw[tick] (1.42500,-0.71700) -- (1.42500,-0.52300);
    \foreach \p in {A,B,C} {\dt{\p}}
    \node[vlab,above] at (A) {$A$};
    \node[vlab,above] at (B) {$B$};
    \node[vlab,above] at (C) {$C$};
  \end{scope}
  % ---- (a) right copy
  \begin{scope}[shift={(5.20,2.90)}]
    \coordinate (Ap) at (0,0); \coordinate (Bp) at (1.95000,0);
    \coordinate (Cp) at (2.85000,0);
    \draw[fig] (Ap) -- (Cp);
    \draw[fig] (Ap) .. controls +(0,-0.82667) and +(0,-0.82667) .. (Cp);
    \draw[tick] (0.90000,-0.09700) -- (0.90000,0.09700);
    \draw[tick] (1.05000,-0.09700) -- (1.05000,0.09700);
    \draw[tick] (1.42500,-0.71700) -- (1.42500,-0.52300);
    \foreach \p in {Ap,Bp,Cp} {\dt{\p}}
    \node[vlab,above] at (Ap) {$A'$};
    \node[vlab,above] at (Bp) {$B'$};
    \node[vlab,above] at (Cp) {$C'$};
  \end{scope}
  \node[slab,left] at (0.20,2.90) {(a)};
  % ---- (b)
  \begin{scope}[shift={(0.60,1.45)}]
    \coordinate (Cb) at (1.51900,0); \coordinate (Ab) at (2.26900,0);
    \coordinate (Bb) at (4.67600,0); \coordinate (Db) at (5.42600,0);
    \draw[fig] (0,0) -- (7.00000,0);
    \draw[tick] (1.89400,-0.09700) -- (1.89400,0.09700);
    \draw[tick] (5.05100,-0.09700) -- (5.05100,0.09700);
    \foreach \p in {Cb,Ab,Bb,Db} {\dt{\p}}
    \node[vlab,below] at (Cb) {$C$};
    \node[vlab,below] at (Ab) {$A$};
    \node[vlab,below] at (Bb) {$B$};
    \node[vlab,below] at (Db) {$D$};
  \end{scope}
  \node[slab,left] at (0.20,1.45) {(b)};
  % ---- (c)
  \begin{scope}[shift={(0.60,0)}]
    \coordinate (Ac) at (0.75000,0); \coordinate (Bc) at (3.07000,0);
    \coordinate (Cc) at (4.07000,0); \coordinate (Dc) at (6.39000,0);
    \draw[fig] (0,0) -- (7.00000,0);
    \draw[fig] (Ac) .. controls +(0,0.96000) and +(0,0.96000) .. (Cc);
    \draw[fig] (Bc) .. controls +(0,0.96000) and +(0,0.96000) .. (Dc);
    \draw[tick] (2.41000,0.62300) -- (2.41000,0.81700);
    \draw[tick] (4.73000,0.62300) -- (4.73000,0.81700);
    \foreach \p in {Ac,Bc,Cc,Dc} {\dt{\p}}
    \node[vlab,below] at (Ac) {$A$};
    \node[vlab,below] at (Bc) {$B$};
    \node[vlab,below] at (Cc) {$C$};
    \node[vlab,below] at (Dc) {$D$};
  \end{scope}
  \node[slab,left] at (0.20,0) {(c)};
\end{tikzpicture}
\figcap{4--5}
\end{bkfigure}}

% ------------------------------------------------------------- Fig 4-6
% Definition 4-1.  AB congruent to A'B'' (exact); B' between A' and B''.
\newcommand{\FIGIVSIX}{\begin{bkfigure}[0.72\linewidth]
\begin{tikzpicture}[scale=1.170]
  \coordinate (A)   at (0.90000,1.32000);
  \coordinate (B)   at (4.55000,2.28000);
  \coordinate (Ap)  at (0.85000,0);
  \coordinate (Bp)  at (3.32023,0.09490);
  \coordinate (Bpp) at (4.62134,0.14489);
  \coordinate (Le)  at (5.47071,0.17752);
  \draw[fig] (A) -- (B);
  \draw[fig] (Ap) -- (Le);
  \foreach \p in {A,B,Ap,Bp,Bpp,Le} {\dt{\p}}
  \node[vlab,above left]  at (A)   {$A$};
  \node[vlab,above]       at (B)   {$B$};
  \node[vlab,above left]  at (Ap)  {$A'$};
  \node[vlab,above]       at (Bp)  {$B'$};
  \node[vlab,above]       at (Bpp) {$B''$};
  \node[vlab,above right] at (Le)  {$l'$};
\end{tikzpicture}
\figcap{4--6}
\end{bkfigure}}

% ------------------------------------------------ Figs 4-7, 4-8, 4-9
% 4-7 (Ex 4-3 #1): AB congruent to CB', with C-B'-D.
% 4-8 (Ex 4-3 #2): RS congruent to TS', with T-U-S'.
% 4-9 (Ex 4-3 #3): AB=DE (single tick), BC=EC' (double tick).
%
% SPLIT 2026-08-17: these three were one side-by-side block, printed after
% the whole of Exercise Group 4-3.  The book sets each one immediately after
% the exercise that cites it (4-7 after ex. 1 and 4-8 after ex. 2 in the left
% column of p.68, 4-9 above ex. 3 at the head of the right column), so each
% is now its own full-measure figure and ch04.tex places it inline.  Point
% coordinates are unchanged -- only the wrapper and `scale' moved -- so
% tools/constraints/ch04.py still describes them exactly.  The three scales
% are chosen so all three print about 4.4 cm wide, as the book draws them at
% one size; ticks ride the scale, keeping the book's tick-to-segment ratio.
\newcommand{\FIGIVSEVEN}{\begin{bkfigure}[0.56\linewidth]
\begin{tikzpicture}[scale=1.00]
  \coordinate (A)  at (1.20000,0.95000);
  \coordinate (B)  at (2.40000,0.95000);
  \coordinate (C)  at (0.92000,0);
  \coordinate (Bp) at (2.12000,0);
  \coordinate (D)  at (3.34000,0);
  \draw[fig] (0.30000,0.95000) -- (3.35000,0.95000);
  \draw[fig] (0,0) -- (4.40000,0);
  \draw[tick] (1.80000,0.84200) -- (1.80000,1.05800);
  \draw[tick] (1.52000,-0.10800) -- (1.52000,0.10800);
  \foreach \p in {A,B,C,Bp,D} {\dt{\p}}
  \node[vlab,above] at (A)  {$A$};
  \node[vlab,above] at (B)  {$B$};
  \node[vlab,below] at (C)  {$C$};
  \node[vlab,below] at (Bp) {$B'$};
  \node[vlab,below] at (D)  {$D$};
\end{tikzpicture}
\figcap{4--7}
\end{bkfigure}}

\newcommand{\FIGIVEIGHT}{\begin{bkfigure}[0.56\linewidth]
\begin{tikzpicture}[scale=1.20]
  \coordinate (R)  at (0.65000,1.30000);
  \coordinate (S)  at (3.00000,1.30000);
  \coordinate (T)  at (0.62000,0);
  \coordinate (U)  at (2.32000,0);
  \coordinate (Sp) at (2.97000,0);
  \draw[fig] (0,1.30000) -- (3.60000,1.30000);
  \draw[fig] (0,0) -- (3.65000,0);
  \draw[fig] (R) .. controls +(0,0.64000) and +(0,0.64000) .. (S);
  \draw[fig] (T) .. controls +(0,0.64000) and +(0,0.64000) .. (Sp);
  % tick half-length 0.090 here (not 0.108): this picture is drawn at
  % scale 1.20, so the shorter coordinate prints the same ink as 4-7/4-9.
  \draw[tick] (1.82500,1.69000) -- (1.82500,1.87000);
  \draw[tick] (1.79500,0.39000) -- (1.79500,0.57000);
  \foreach \p in {R,S,T,U,Sp} {\dt{\p}}
  \node[vlab,above left]  at (R)  {$R$};
  \node[vlab,above right] at (S)  {$S$};
  \node[vlab,below] at (T)  {$T$};
  \node[vlab,below] at (U)  {$U$};
  \node[vlab,below] at (Sp) {$S'$};
\end{tikzpicture}
\figcap{4--8}
\end{bkfigure}}

\newcommand{\FIGIVNINE}{\begin{bkfigure}[0.56\linewidth]
\begin{tikzpicture}[scale=1.02]
  \coordinate (A)  at (0.75000,0.95000);
  \coordinate (B)  at (2.60000,0.95000);
  \coordinate (C)  at (3.40000,0.95000);
  \coordinate (D)  at (0.75000,0);
  \coordinate (E)  at (2.60000,0);
  \coordinate (Cp) at (3.40000,0);
  \draw[fig] (0,0.95000) -- (4.30000,0.95000);
  \draw[fig] (0,0) -- (4.30000,0);
  \draw[tick] (1.67500,0.84200) -- (1.67500,1.05800);
  \draw[tick] (2.91700,0.84200) -- (2.91700,1.05800);
  \draw[tick] (3.08300,0.84200) -- (3.08300,1.05800);
  \draw[tick] (1.67500,-0.10800) -- (1.67500,0.10800);
  \draw[tick] (2.91700,-0.10800) -- (2.91700,0.10800);
  \draw[tick] (3.08300,-0.10800) -- (3.08300,0.10800);
  \foreach \p in {A,B,C,D,E,Cp} {\dt{\p}}
  \node[vlab,above] at (A)  {$A$};
  \node[vlab,above] at (B)  {$B$};
  \node[vlab,above] at (C)  {$C$};
  \node[vlab,below] at (D)  {$D$};
  \node[vlab,below] at (E)  {$E$};
  \node[vlab,below] at (Cp) {$C'$};
\end{tikzpicture}
\figcap{4--9}
\end{bkfigure}}

% ------------------------------------------------------ Figs 4-10, 4-11
% 4-10 (Thm 4-2): AB=A'C' (double tick on the arcs), AC=A'B' (single tick).
% 4-11: AC=A'B' (single tick), CX=B'C' (double tick).
\newcommand{\FIGIVTENELEVEN}{\begin{bkfigure}[\linewidth]
\begin{minipage}{0.47\linewidth}\centering
\begin{tikzpicture}[scale=0.901]
  \coordinate (A)  at (0.30000,1.15000);
  \coordinate (C)  at (2.75000,1.15000);
  \coordinate (B)  at (3.45000,1.15000);
  \coordinate (Ap) at (0.30000,0);
  \coordinate (Bp) at (2.75000,0);
  \coordinate (Cp) at (3.45000,0);
  % Right-hand stubs trimmed 4.30 -> 3.80 (2026-08-17).  Measured off the
  % plate, each line runs about 0.10 of AB past B / C'; ours ran 0.27 past,
  % so both lines were over-drawn to the right.  The left stubs (0.30, also
  % 0.10 of AB) already matched.
  \draw[fig] (0,1.15000) -- (3.80000,1.15000);
  \draw[fig] (0,0) -- (3.80000,0);
  \draw[fig] (A)  .. controls +(0,0.76000) and +(0,0.76000) .. (B);
  \draw[fig] (Ap) .. controls +(0,-0.76000) and +(0,-0.76000) .. (Cp);
  \draw[tick] (1.80000,1.62300) -- (1.80000,1.81700);
  \draw[tick] (1.95000,1.62300) -- (1.95000,1.81700);
  \draw[tick] (1.80000,-0.66700) -- (1.80000,-0.47300);
  \draw[tick] (1.95000,-0.66700) -- (1.95000,-0.47300);
  \draw[tick] (1.52500,1.05300) -- (1.52500,1.24700);
  \draw[tick] (1.52500,-0.09700) -- (1.52500,0.09700);
  % The two correspondence arrows.  In the book they bow AWAY from each
  % other, so that "then (" and ") if" read as a facing pair of parentheses:
  % the B'->C arrow bows left (west), the B->C' arrow bows right (east).
  % Both were drawn with `bend right' here, which mirrored each one and made
  % the pair read ")(" instead.  For an upward path `bend left' is the
  % westward bow; for a downward path it is the eastward one.
  \draw[fig,->,line width=0.55pt] (2.72000,0.14000)
        to[bend left=28] (2.72000,1.01000);
  \draw[fig,->,line width=0.55pt] (3.45000,1.01000)
        to[bend left=28] (3.45000,0.14000);
  \foreach \p in {A,C,B,Ap,Bp,Cp} {\dt{\p}}
  \node[vlab,above left]  at (A)  {$A$};
  \node[vlab,above]       at (C)  {$C$};
  % B and C' are the far ends of the two arcs, so the arc arrives into the
  % label's own corner.  The plate sets both in the free quadrant on the
  % outside of the arc -- B above-RIGHT of its dot, C' below-RIGHT -- which
  % is also what keeps them clear of the arc ends here.
  \node[vlab,above right] at (B)  {$B$};
  \node[vlab,below left]  at (Ap) {$A'$};
  \node[vlab,below]       at (Bp) {$B'$};
  \node[vlab,below right] at (Cp) {$C'$};
  \node[slab,left]  at (2.42000,0.57500) {then};
  \node[slab,right] at (3.78000,0.57500) {if};
\end{tikzpicture}
\figcap{4--10}
\end{minipage}\hfill
\begin{minipage}{0.47\linewidth}\centering
\begin{tikzpicture}[scale=0.901]
  \coordinate (A)  at (0.60000,0.95000);
  \coordinate (C)  at (2.55000,0.95000);
  \coordinate (X)  at (3.35000,0.95000);
  \coordinate (Ap) at (0.60000,0);
  \coordinate (Bp) at (2.55000,0);
  \coordinate (Cp) at (3.35000,0);
  \draw[fig] (0,0.95000) -- (4.20000,0.95000);
  \draw[fig] (0,0) -- (4.20000,0);
  \draw[tick] (1.57500,0.85300) -- (1.57500,1.04700);
  \draw[tick] (2.87500,0.85300) -- (2.87500,1.04700);
  \draw[tick] (3.02500,0.85300) -- (3.02500,1.04700);
  \draw[tick] (1.57500,-0.09700) -- (1.57500,0.09700);
  \draw[tick] (2.87500,-0.09700) -- (2.87500,0.09700);
  \draw[tick] (3.02500,-0.09700) -- (3.02500,0.09700);
  \foreach \p in {A,C,X,Ap,Bp,Cp} {\dt{\p}}
  \node[vlab,above] at (A)  {$A$};
  \node[vlab,above] at (C)  {$C$};
  \node[vlab,above] at (X)  {$X$};
  \node[vlab,below] at (Ap) {$A'$};
  \node[vlab,below] at (Bp) {$B'$};
  \node[vlab,below] at (Cp) {$C'$};
\end{tikzpicture}
\figcap{4--11}
\end{minipage}
\end{bkfigure}}

% ------------------------------------------------------ Figs 4-12, 4-13
% 4-12 (Thm 4-3): AB=A'B' (double tick), AC=A'C' (single tick on the arcs).
% 4-13 (Thm 4-4): AD'=CD=EF.
\newcommand{\FIGIVTWELVETHIRTEEN}{\begin{bkfigure}[\linewidth]
\begin{minipage}{0.47\linewidth}\centering
\begin{tikzpicture}[scale=0.833]
  \coordinate (A)  at (0.30000,1.25000);
  \coordinate (B)  at (2.75000,1.25000);
  \coordinate (C)  at (3.50000,1.25000);
  \coordinate (Ap) at (0.30000,0);
  \coordinate (Bp) at (2.75000,0);
  \coordinate (Cp) at (3.50000,0);
  % Right-hand stubs trimmed 4.10 -> 3.90: the plate runs about 0.12 of AC
  % past C / C', ours ran 0.19 past.
  \draw[fig] (0,1.25000) -- (3.90000,1.25000);
  \draw[fig] (0,0) -- (3.90000,0);
  \draw[fig] (A)  .. controls +(0,-0.60000) and +(0,-0.60000) .. (C);
  \draw[fig] (Ap) .. controls +(0,-0.60000) and +(0,-0.60000) .. (Cp);
  \draw[tick] (1.45000,1.15300) -- (1.45000,1.34700);
  \draw[tick] (1.60000,1.15300) -- (1.60000,1.34700);
  \draw[tick] (1.45000,-0.09700) -- (1.45000,0.09700);
  \draw[tick] (1.60000,-0.09700) -- (1.60000,0.09700);
  \draw[tick] (1.90000,0.70300) -- (1.90000,0.89700);
  \draw[tick] (1.90000,-0.54700) -- (1.90000,-0.35300);
  % The B'->B correspondence arrow.  The plate draws a full arrow from just
  % above B' up to just under B, crossing the upper arc on the way; ours was
  % a short stub hanging off B and pointing the wrong way (down into empty
  % space), so the figure said nothing about which point matched which.
  % Upward path + `bend right' = the plate's slight eastward bow.
  \draw[fig,->,line width=0.55pt] (2.79000,0.15000)
        to[bend right=14] (2.75000,1.11000);
  \foreach \p in {A,B,C,Ap,Bp,Cp} {\dt{\p}}
  \node[vlab,above] at (A)  {$A$};
  \node[vlab,above] at (B)  {$B$};
  \node[vlab,above] at (C)  {$C$};
  \node[vlab,below left] at (Ap) {$A'$};
  \node[vlab,below]      at (Bp) {$B'$};
  \node[vlab,below]      at (Cp) {$C'$};
\end{tikzpicture}
\figcap{4--12}
\end{minipage}\hfill
\begin{minipage}{0.47\linewidth}\centering
\begin{tikzpicture}[scale=0.833]
  \coordinate (A)  at (0,1.45000);
  \coordinate (Dp) at (3.34012,1.70701);
  \coordinate (B)  at (5.63333,1.88347);
  \coordinate (C)  at (0,0.72000);
  \coordinate (D)  at (3.33824,1.00033);
  \coordinate (E)  at (0,0);
  \coordinate (F)  at (3.34776,0.12274);
  \draw[fig] (A) -- (B);
  \draw[fig] (C) -- (D);
  \draw[fig] (E) -- (F);
  \foreach \p in {A,Dp,B,C,D,E,F} {\dt{\p}}
  \node[vlab,above left] at (A)  {$A$};
  \node[vlab,above]      at (Dp) {$D'$};
  \node[vlab,above]      at (B)  {$B$};
  \node[vlab,above left] at (C)  {$C$};
  \node[vlab,above]      at (D)  {$D$};
  \node[vlab,above left] at (E)  {$E$};
  \node[vlab,above]      at (F)  {$F$};
\end{tikzpicture}
\figcap{4--13}
\end{minipage}
\end{bkfigure}}

% figures05.tex -- Chapter 5, "Measurement of Segments".
% All figures in this chapter are number-line / segment diagrams: labelled
% points on straight lines, with the unit segment shown below.  Every point
% position below is exact arithmetic on the unit u -- no point is placed by
% eye.  Proportions read off the iPad scans (scan05 pp.18-24, scan06 pp.6-7)
% and the photo PDF (PDF pp.87-97); see tools/constraints/ch05.py.
%
% Fig. 5-1 (a pictorial vignette of a man pacing off a building) is NOT
% redrawn -- see chapters/ch05-UNCERTAIN.md.
%
% Loaded here rather than in ch05.tex so that this file stands alone for
% tools/check_labels.py, whose driver inputs figures05.tex by itself.
\usetikzlibrary{decorations.pathreplacing}

% ---------------------------------------------------------------- Fig. 5-2
% Unit CD laid off from A: A1,A2,A3,A4 at 1..4 units; B between A3 and A4
% (measured at 0.36 of that gap, giving AB about 3.36 units).
\newcommand{\FIGVTWO}{%
\begin{bkfigure}[0.92\linewidth]
\begin{tikzpicture}[x=1.75cm,y=1cm,scale=1.257]
  \draw[given] (-0.20,0) -- (4.40,0);
  \foreach \x/\lb in {0/A,1/{A_1},2/{A_2},3/{A_3},3.36/B,4/{A_4}}{%
    \dt{\x,0}
    \node[vlab,above] at (\x,0) {$\lb$};}
  \node[vlab,right] at (4.40,0) {$l$};
  % unit segment CD, drawn below, spanning 0.52u to 1.52u.  Drop measured
  % from the book plate: gap between the two lines is 0.35 of the unit.
  \draw[given] (0.52,-0.61) -- (1.52,-0.61);
  \dt{0.52,-0.61} \dt{1.52,-0.61}
  \node[vlab,above,xshift=-2pt] at (0.52,-0.61) {$C$};
  \node[vlab,above,xshift=2pt] at (1.52,-0.61) {$D$};
  \node[slab,above] at (1.02,-0.63) {1};
\end{tikzpicture}
\figcap{5--2}
\end{bkfigure}}

% ---------------------------------------------------------------- Fig. 5-3
% The unit is too large: B lies between A and A1.  Re-measured off the book
% plate: B sits at 0.32 of AA1 (not 0.45), and the line runs on to 1.71 u.
\newcommand{\FIGVTHREE}{%
\begin{bkfigure}[0.62\linewidth]
\begin{tikzpicture}[x=1.75cm,y=1cm,scale=1.300]
  \draw[given] (-0.26,0) -- (1.71,0);
  \foreach \x/\lb in {0/A,0.32/B,1/{A_1}}{%
    \dt{\x,0}
    \node[vlab,above] at (\x,0) {$\lb$};}
\end{tikzpicture}
\figcap{5--3}
\end{bkfigure}}

% ----------------------------------------------------------- Figs. 5-4/5-5
% Archimedes' integer n = 4.  5-4: B = A4 exactly.  5-5: B between A3 and A4
% (measured at 0.18 of that gap).  In both plates the drawn line STARTS at A
% -- there is no stub to the left of A, unlike Fig. 5-2.
\newcommand{\FIGVFOURFIVE}{%
\begin{bkfigure}[0.94\linewidth]
\begin{tikzpicture}[x=1.75cm,y=1cm,scale=1.274]
  \draw[given] (0,0) -- (4.40,0);
  \foreach \x/\lb in {0/{A=A_0},1/{A_1},2/{A_2},3/{A_3},4/{A_4=B}}{%
    \dt{\x,0}
    \node[vlab,above] at (\x,0) {$\lb$};}
  \draw[given] (0.52,-0.61) -- (1.52,-0.61);
  \dt{0.52,-0.61} \dt{1.52,-0.61}
  \node[vlab,above,xshift=-2pt] at (0.52,-0.61) {$C$};
  \node[vlab,above,xshift=2pt] at (1.52,-0.61) {$D$};
\end{tikzpicture}
\figcap{5--4}
\end{bkfigure}%
\begin{bkfigure}[0.94\linewidth]
\begin{tikzpicture}[x=1.75cm,y=1cm,scale=1.274]
  \draw[given] (0,0) -- (4.40,0);
  \foreach \x in {0,1,2,3,3.18,4}{\dt{\x,0}}
  \foreach \x/\lb in {0/A,1/{A_1},2/{A_2},4/{A_4}}{%
    \node[vlab,above] at (\x,0) {$\lb$};}
  % A3 and B sit 0.18 u apart, as in the book; the labels are nudged apart
  % by 1.5pt each so they read as "A_3 B" rather than colliding.  The DOTS
  % stay at the exact positions 3 and 3.18 -- only the text moves.
  \node[vlab,above,xshift=-1.5pt] at (3,0) {$A_3$};
  \node[vlab,above,xshift=1.5pt] at (3.18,0) {$B$};
\end{tikzpicture}
\figcap{5--5}
\end{bkfigure}}

% ----------------------------------------------------------- Figs. 5-6/5-7
% 5-6: the unit interval CD divided into tenths (9 interior marks).
% 5-7: AB measured as about 2.3 units; the tenth-scale unit CD alongside.
\newcommand{\FIGVSIXSEVEN}{%
\begin{bkfigure}[0.90\linewidth]
\begin{tikzpicture}[x=7.5cm,y=1cm,scale=1.300]
  \draw[given] (0,0) -- (1,0);
  \dt{0,0} \dt{1,0}
  \node[vlab,above,xshift=-2pt] at (0,0) {$C$};
  \node[vlab,above,xshift=2pt] at (1,0) {$D$};
  \foreach \k in {1,...,9}{\draw[tick] ({\k/10},-0.085) -- ({\k/10},0.085);}
\end{tikzpicture}
\figcap{5--6}
\end{bkfigure}%
\begin{bkfigure}[0.90\linewidth]
\begin{tikzpicture}[x=1.9cm,y=1cm,scale=1.300]
  \draw[given] (0,0) -- (2.30,0);
  \foreach \x/\lb in {0/A,1/{A_1},2/{A_2},2.30/B}{%
    \dt{\x,0}
    \node[vlab,above] at (\x,0) {$\lb$};}
  % the unit CD alongside, itself divided into tenths
  \draw[given] (2.64,0) -- (3.64,0);
  \dt{2.64,0} \dt{3.64,0}
  \node[vlab,above,xshift=-2pt] at (2.64,0) {$C$};
  \node[vlab,above,xshift=2pt] at (3.64,0) {$D$};
  \foreach \k in {1,...,9}{\draw[tick] ({2.64+\k/10},-0.075) -- ({2.64+\k/10},0.075);}
\end{tikzpicture}
\figcap{5--7}
\end{bkfigure}}

% ---------------------------------------------------------------- Fig. 5-8
% Theorem 5-2: B'' chosen on the same side of A as B with AB'' = A'B'.
% B'' drawn between A and B (measured at 0.77 of AB).  AB is drawn the same
% length as A'B' (the theorem's given |AB| = |A'B'|), A sits directly under
% A', and the two rows are 0.17 of the segment length apart -- all measured
% off the book plate.
\newcommand{\FIGVEIGHT}{%
\begin{bkfigure}[0.72\linewidth]
\begin{tikzpicture}[x=4.6cm,y=1cm,scale=1.300]
  \draw[given] (0,0.80) -- (1,0.80);
  \dt{0,0.80} \dt{1,0.80}
  \node[vlab,above] at (0,0.80) {$A'$};
  \node[vlab,above] at (1,0.80) {$B'$};
  \draw[given] (0,0) -- (1,0);
  \foreach \x/\lb in {0/A,1/B}{%
    \dt{\x,0}
    \node[vlab,above] at (\x,0) {$\lb$};}
  \dt{0.77,0}
  \node[vlab,above,yshift=1.2pt] at (0.77,0) {$B''$};
\end{tikzpicture}
\figcap{5--8}
\end{bkfigure}}

% ---------------------------------------------------------------- Fig. 5-9
% Theorem 5-3: E on the same side of C as D with AB = CE; D between C and E
% (measured at 0.66 of CE).
\newcommand{\FIGVNINE}{%
\begin{bkfigure}[0.72\linewidth]
\begin{tikzpicture}[x=4.6cm,y=1cm,scale=1.300]
  \draw[given] (0,0.75) -- (1,0.75);
  \dt{0,0.75} \dt{1,0.75}
  \node[vlab,above] at (0,0.75) {$A$};
  \node[vlab,above] at (1,0.75) {$B$};
  \draw[given] (0,0) -- (1,0);
  \foreach \x/\lb in {0/C,0.66/D,1/E}{%
    \dt{\x,0}
    \node[vlab,above] at (\x,0) {$\lb$};}
\end{tikzpicture}
\figcap{5--9}
\end{bkfigure}}

% --------------------------------------------------------------- Fig. 5-10
% Change of scale: E is the midpoint of the unit CD, so AB = 3 with unit CD
% and AB = 6 with unit ED.  Horizontal layout re-measured off the book plate
% (all in CD units from C): A at 1.55, B at 4.55, the half-unit ED alongside
% at 5.15..5.65.  The earlier 2.24/5.24/5.92/6.42 spread the plate too wide.
\newcommand{\FIGVTEN}{%
\begin{bkfigure}[0.96\linewidth]
\begin{tikzpicture}[x=1.25cm,y=1cm,scale=1.300]
  % unit CD with midpoint E, braced and labelled 1
  \draw[given] (0,0) -- (1,0);
  \dt{0,0} \dt{1,0}
  \draw[tick] (0.5,-0.09) -- (0.5,0.09);
  \node[vlab,below] at (0,-0.02) {$C$};
  \node[vlab,below] at (0.5,-0.02) {$E$};
  \node[vlab,below] at (1,-0.02) {$D$};
  \draw[fig,line width=0.5pt,decorate,decoration={brace,amplitude=3pt}]
        (0,0.16) -- (1,0.16);
  \node[slab,above,yshift=2.6pt] at (0.5,0.19) {1};
  % AB, three units long; the "1" sits under its first third
  \draw[given] (1.55,0.55) -- (4.55,0.55);
  \dt{1.55,0.55} \dt{4.55,0.55}
  \node[vlab,above] at (1.55,0.55) {$A$};
  \node[vlab,above] at (4.55,0.55) {$B$};
  \draw[tick] (2.55,0.46) -- (2.55,0.64);
  \draw[tick] (3.55,0.46) -- (3.55,0.64);
  \node[slab,below,yshift=-2.2pt] at (2.05,0.53) {1};
  % the half-unit ED alongside
  \draw[given] (5.15,0.30) -- (5.65,0.30);
  \dt{5.15,0.30} \dt{5.65,0.30}
  \node[vlab,above] at (5.15,0.30) {$E$};
  \node[vlab,above] at (5.65,0.30) {$D$};
\end{tikzpicture}
\figcap{5--10}
\end{bkfigure}}

% figures06.tex -- Chapter 6 figures (Congruence of Angles and Triangles).
%
% Built 2026-08-17 from sources/Geometry.pdf pp. 98-129 (book pp. 84-115) with
% structure checked against the iPad scans (scan06 p11-p24, scan07 p1-p24,
% scan08 p1-p10; see sources/scan-index/).
%
% Constrained points are DERIVED, never eyeballed:
%   parallels  -> equal ratios along the two sides   ($(X)!r!(Y)$)
%   altitudes / perpendicular feet                   ($(X)!(P)!(Y)$)
%   crossings  -> (intersection of A--B and C--D)
%   midpoints  -> ($(X)!0.5!(Y)$)
% so the hypothesis the book states in words holds exactly in the drawing.
% The same hypotheses are re-asserted numerically in tools/constraints/ch06.py.
%
% Macro names are chapter-prefixed English number words: \FIGVI<NUMBERWORD>.
% A macro spanning several consecutive figures concatenates the words
% (e.g. \FIGVITWOTHREEFOUR holds Figs 6-2, 6-3 and 6-4).
%
% Line treatment follows the book: solid where the book prints solid (given),
% dashed where the book prints dashed (the constructed D' in the congruence
% proofs, and the "moved" angle B' in the inequality figures).  This chapter is
% otherwise plain `given' line work -- the book draws these diagrams uniformly.

\usetikzlibrary{decorations.pathreplacing,arrows.meta,angles,quotes}

% Primed labels: a math prime sits high enough that the collision audit reads it
% as a separate word overlapping its own letter; a small drop keeps the pair in
% one word and is visually indistinguishable.  (Same fix as ch09.)
\newcommand{\VIpr}{\raisebox{-1.3pt}{\ensuremath{{}^{\prime}}}}
\newcommand{\VIppr}{\VIpr\VIpr}
\newcommand{\VIpppr}{\VIpr\VIpr\VIpr}

% Chapter-local label clearance.  A labelled point lying ON a line needs a
% little more room than brumfiel.sty's default or the line clips the box.
\tikzset{
  vlab/.append style={outer sep=2.9pt},
  slab/.append style={outer sep=2.9pt},
  klab/.append style={outer sep=2.9pt},
}

% ---------------------------------------------------------------- house marks
% The congruence apparatus this chapter runs on, defined once.
%
% All helpers are chapter-prefixed: \aa and friends collide with LaTeX
% built-ins (\aa is the Scandinavian ring-a), so every one carries the VI code.
%
% \VItk{A}{B}{n} -- n tick marks across the middle of segment AB, perpendicular
%                 to it.  The book uses one tick for one pair of congruent
%                 segments, two for the second pair.
\newcommand{\VItk}[3]{%
  \foreach \tki in {1,...,#3}{%
    \pgfmathsetmacro{\tkd}{(\tki-(#3+1)/2)*2.0}%
    \coordinate (tkm) at ($(#1)!0.5!(#2)$);
    \coordinate (tkt) at ($(tkm)!\tkd pt!(#2)$);
    \draw[tick] ($(tkt)!2.6pt!90:(#2)$) -- ($(tkt)!2.6pt!270:(#2)$);}}
%
% The angle marks are struck by hand rather than through the `angles' pic
% library, because that library always sweeps COUNTERCLOCKWISE from its first
% ray to its second.  Whenever a call happened to name the two sides the other
% way round, the pic drew the REFLEX angle -- a 331-degree sweep that renders
% as a near-complete circle swallowing the whole figure (this is what put the
% spurious circle in Fig. 6-13, and it silently mis-marked several others).
% Striking the arc from the two bearings and normalising the sweep into
% (-180,180] makes every call draw the minor angle no matter which side is
% named first, so the mark can never depend on argument order again.
%
% \VIarc{P}{V}{Q}{r} -- single arc of radius r at vertex V, from side VP to VQ.
\newcommand{\VIarc}[4]{%
  \begingroup
  \pgfmathanglebetweenpoints{\pgfpointanchor{#2}{center}}{\pgfpointanchor{#1}{center}}%
  \edef\VIaP{\pgfmathresult}%
  \pgfmathanglebetweenpoints{\pgfpointanchor{#2}{center}}{\pgfpointanchor{#3}{center}}%
  \edef\VIaQ{\pgfmathresult}%
  \pgfmathparse{\VIaQ-\VIaP-360*round((\VIaQ-\VIaP)/360)}%
  \edef\VIdA{\pgfmathresult}%
  \draw[fig, line width=0.55pt] ([shift=(\VIaP:#4)]#2)
       arc[start angle=\VIaP, end angle=\VIaP+\VIdA, radius=#4];
  \endgroup}
% \VIarcb{P}{V}{Q}{r} -- double arc (the book's second congruent-angle mark).
\newcommand{\VIarcb}[4]{%
  \VIarc{#1}{#2}{#3}{#4}%
  \VIarc{#1}{#2}{#3}{#4+2.1pt}}
% \VIarcvia{P}{V}{Q}{W}{r} -- arc at V from side VP to side VQ, sweeping the way
%   that passes through ray VW.  \VIarc always takes the MINOR angle, which is
%   right for a numbered angle but cannot draw the marks Fig. 6-94 needs: there
%   the book sweeps one arc from OA all the way round to OA' on the side B lies
%   on -- a straight angle in case 2 and a reflex one in case 3.  Naming the ray
%   the arc must pass through picks the branch without hard-coding bearings, so
%   the mark still follows the points if they move.
\newcommand{\VIarcvia}[5]{%
  \begingroup
  \pgfmathanglebetweenpoints{\pgfpointanchor{#2}{center}}{\pgfpointanchor{#1}{center}}%
  \edef\VIvP{\pgfmathresult}%
  \pgfmathanglebetweenpoints{\pgfpointanchor{#2}{center}}{\pgfpointanchor{#3}{center}}%
  \edef\VIvQ{\pgfmathresult}%
  \pgfmathanglebetweenpoints{\pgfpointanchor{#2}{center}}{\pgfpointanchor{#4}{center}}%
  \edef\VIvW{\pgfmathresult}%
  \pgfmathparse{mod(\VIvQ-\VIvP+720,360)}\edef\VIdQ{\pgfmathresult}%
  \pgfmathparse{mod(\VIvW-\VIvP+720,360)}\edef\VIdW{\pgfmathresult}%
  \pgfmathparse{\VIdW<\VIdQ ? \VIdQ : \VIdQ-360}\edef\VIsw{\pgfmathresult}%
  \draw[fig, line width=0.55pt] ([shift=(\VIvP:#5)]#2)
       arc[start angle=\VIvP, end angle=\VIvP+\VIsw, radius=#5];
  \endgroup}
% \VIrt{P}{V}{Q}{r} -- the small square marking a right angle.  Built from the
% two side directions directly, so it always lands inside the angle.
\newcommand{\VIrt}[4]{%
  \begingroup
  \coordinate (rt@p) at ($(#2)!#4!(#1)$);
  \coordinate (rt@q) at ($(#2)!#4!(#3)$);
  \draw[fig, line width=0.55pt] (rt@p) -- ($(rt@p)+(rt@q)-(#2)$) -- (rt@q);
  \endgroup}
% \VInum{V}{ang}{d}{txt} -- a small numeral label (the book's "1", "2", ...
%   inside an angle), placed d from vertex V along direction ang.
\newcommand{\VInum}[4]{\node[slab] at ($(#1)+({#2}:#3)$) {#4};}
%
% \VInumb{P}{V}{Q}{d}{txt} -- the same numeral, but placed on the BISECTOR of
%   angle P-V-Q at distance d.  Placing by bisector rather than by a
%   hand-picked compass angle is what keeps the numeral off both rays, which
%   is where most label/geometry collisions in this chapter came from.
\newcommand{\VInumb}[5]{%
  \begingroup
  \coordinate (nb@1) at ($(#2)!1cm!(#1)$);
  \coordinate (nb@2) at ($(#2)!1cm!(#3)$);
  \coordinate (nb@m) at ($(nb@1)!0.5!(nb@2)$);
  \node[slab] at ($(#2)!#4!(nb@m)$) {#5};
  \endgroup}

% ================================================================== Fig 6-1
% Book p.84.  The given angle A (a free "V"), and below it the line l' with A'
% on it: r' is the ray of l' to the right of A', s' the constructed ray.  The
% dashed stub to the left of A' carries the caption "Given side of l'".
\newcommand{\FIGVIONE}{\begin{bkfigure}[\linewidth]
\begin{tikzpicture}[scale=1.235]
  % given angle A, drawn free above
  \coordinate (A)  at (2.55,1.32);
  \draw[given] (A) -- (3.12,2.14);
  \draw[given] (A) -- (3.95,0.92);
  \node[vlab, above left] at (A) {$A$};
  % line l' with A' on it
  \coordinate (Ap) at (1.72,0.16);
  \draw[given] (Ap) -- (3.62,-0.22);           % ray r'
  \draw[aux]   (Ap) -- (0.62,0.38);            % given side of l', dashed stub
  \draw[given] (Ap) -- (2.62,1.24);            % ray s'
  \node[vlab, below left] at (Ap) {$A\VIpr$};
  \node[slab, right] at (3.62,-0.22) {$r\VIpr$};
  \node[slab] at (2.10,1.02) {$s\VIpr$};
  \node[slab, below]  at (0.62,0.38) {$l\VIpr$};
  % The book runs the caption straight through the ray s', with the ray
  % passing in the word space.  Setting it as two nodes either side of the ray
  % reproduces that exactly and leaves the ray unobscured.
  \node[slab, anchor=east] at (2.01,0.60) {Given side};
  \node[slab, anchor=west] at (2.17,0.60) {of $l\VIpr$};
\end{tikzpicture}
\figcap{6--1}
\end{bkfigure}}

% =========================================================== Figs 6-2, 6-3, 6-4
% Book p.85, Exercise Group 6-1.
%  6-2: angle A with sides r (right) and s (up-left); ray r' from A' below.
%  6-3: angle A with sides r and s; a detached ray r' ending at A'.
%  6-4: a straight angle A on a line, with A' on it and ray r'.
\newcommand{\FIGVITWOTHREEFOUR}{\begin{bkfigure}[\linewidth]
\begin{minipage}{0.46\linewidth}\centering
\begin{tikzpicture}[scale=1.008]
  \coordinate (A)  at (1.35,0.52);
  \coordinate (Ap) at (0.22,-0.62);
  \draw[given] (A) -- (2.62,0.44);            % r
  \draw[given] (A) -- (0.52,1.02);            % toward r'/s corner
  \draw[given] (0.52,1.02) -- (0.30,1.62);    % s
  \draw[given] (Ap) -- (0.52,1.02);
  \node[vlab, below] at (A) {$A$};
  \node[vlab, below left] at (Ap) {$A\VIpr$};
  \node[slab, right] at (2.62,0.44) {$r$};
  \node[slab, above] at (0.30,1.62) {$s$};
  \node[slab, left]  at (0.52,1.02) {$r\VIpr$};
\end{tikzpicture}
\figcap{6--2}
\end{minipage}\hfill
\begin{minipage}{0.46\linewidth}\centering
\begin{tikzpicture}[scale=1.008]
  \coordinate (A) at (1.05,1.05);
  \draw[given] (A) -- (0.18,0.05);            % s
  \draw[given] (A) -- (1.72,0.12);            % r
  \node[vlab, above] at (A) {$A$};
  \node[slab, below left]  at (0.18,0.05) {$s$};
  \node[slab, below right] at (1.72,0.12) {$r$};
  % detached ray r' ending at A'
  \draw[given] (2.18,1.00) -- (3.30,1.00);
  \node[slab, above] at (2.30,1.00) {$r\VIpr$};
  \node[vlab, right] at (3.30,1.00) {$A\VIpr$};
\end{tikzpicture}
\figcap{6--3}

\medskip
\begin{tikzpicture}[scale=1.008]
  \coordinate (A)  at (0.95,0.72);
  \coordinate (Ap) at (1.42,0.72);
  \draw[given] (0.12,0.28) -- (Ap);
  \draw[given] (Ap) -- (3.05,0.16);
  \dt{A}
  \node[vlab, above] at (A) {$A$};
  \node[vlab, above right] at (Ap) {$A\VIpr$};
  \node[slab, right] at (3.05,0.16) {$r\VIpr$};
\end{tikzpicture}
\figcap{6--4}
\end{minipage}
\end{bkfigure}}

% ================================================================== Fig 6-5
% Book p.86.  THREE rays from each vertex, not four.  The common side splits
% the whole angle A into the two parts B and C.  The book marks this with two
% concentric arcs struck from the vertex: a SMALL one spanning the whole
% opening (that is angle A) and a LARGE one, also spanning the whole opening,
% which the middle ray cuts into the two arcs belonging to B and C.  Each
% letter is written just BEYOND its own arc -- A outside the small arc, B and
% C outside the large one, in their respective sectors.  The common side is
% drawn a little longer than the other two, as in the book.
% Ray directions and arc radii measured off the plate (book p.86, left fan):
% sides at 146/94/52 degrees, large arc at 0.89 and small arc at 0.35 of the
% side length.
\newcommand{\FIGVIFIVE}{\begin{bkfigure}[\linewidth]
\begin{tikzpicture}[scale=1.300]
 \foreach \dx/\lA/\lB/\lC in {0/A/B/C, 4.27/A\VIpr/B\VIpr/C\VIpr}{
  \begin{scope}[xshift=\dx cm]
   \coordinate (V) at (0,0);
   \draw[given] (V) -- ++(146:2.30);        % outer side, left
   \draw[given] (V) -- ++(94:2.38);         % common side of B and C
   \draw[given] (V) -- ++(52:2.20);         % outer side, right
   % large arc -- the middle ray cuts it into the marks for B and C
   \draw[tick] ([shift=(146:2.04)]V)
        arc[start angle=146, end angle=52, radius=2.04];
   % small arc near the vertex -- the mark for the whole angle A
   \draw[tick] ([shift=(146:0.80)]V)
        arc[start angle=146, end angle=52, radius=0.80];
   \node[vlab] at ($(V)+(116:1.14)$) {$\lA$};
   \node[vlab] at ($(V)+(120:2.36)$) {$\lB$};
   \node[vlab] at ($(V)+(71:2.33)$)  {$\lC$};
  \end{scope}}
\end{tikzpicture}
\figcap{6--5}
\end{bkfigure}}

% =========================================================== Figs 6-6, 6-7
% Book p.86, Exercise Group 6-3.
%  6-6: same two fans as 6-5 but labelled A,B,C and A',C',B' (the primed fan
%       has C' inside B', which is what makes part (b) of the exercise work).
%  6-7: two separate three-ray fans, P Q R from O and R V U from S T U -- the
%       book draws them as two "V" clusters sharing a baseline.
% 6-6 is the same three-ray idea as 6-5, but here each of the three angles
% carries its OWN arc at its own radius: the whole angle A on the outside, and
% the two parts on smaller arcs.  The unprimed fan opens up-left to
% horizontal-right with B the upper part and C the lower; the primed fan is
% the mirror-ish case with C' the part next to the horizontal side and B' the
% part next to the slanting one -- that swap is what the exercise turns on.
\newcommand{\FIGVISIXSEVEN}{\begin{bkfigure}[\linewidth]
\begin{minipage}{0.60\linewidth}\centering
\begin{tikzpicture}[scale=1.040]
   % ---- unprimed fan: sides at 134 / 89 / 0 degrees
   \coordinate (V) at (0,0);
   \draw[given] (V) -- ++(134:1.84);
   \draw[given] (V) -- ++(89:1.87);
   \draw[given] (V) -- ++(0:1.86);
   \draw[tick] ([shift=(134:1.61)]V)
        arc[start angle=134, end angle=0, radius=1.61];    % whole angle A
   \draw[tick] ([shift=(134:0.50)]V)
        arc[start angle=134, end angle=89, radius=0.50];   % part B
   \draw[tick] ([shift=(89:1.06)]V)
        arc[start angle=89, end angle=0, radius=1.06];     % part C
   \node[vlab] at ($(V)+(64:1.94)$)  {$A$};
   \node[vlab] at ($(V)+(111:0.96)$) {$B$};
   \node[vlab] at ($(V)+(45:0.62)$)  {$C$};
   % ---- primed fan: sides at 163 / 89 / 48 degrees
   \coordinate (W) at (4.24,0);
   \draw[given] (W) -- ++(163:1.91);
   \draw[given] (W) -- ++(89:2.06);
   \draw[given] (W) -- ++(48:1.93);
   \draw[tick] ([shift=(163:1.58)]W)
        arc[start angle=163, end angle=48, radius=1.58];   % whole angle A'
   \draw[tick] ([shift=(163:1.06)]W)
        arc[start angle=163, end angle=89, radius=1.06];   % part C'
   \draw[tick] ([shift=(89:0.50)]W)
        arc[start angle=89, end angle=48, radius=0.50];    % part B'
   \node[vlab] at ($(W)+(119:2.01)$) {$A\VIpr$};
   \node[vlab] at ($(W)+(126:0.62)$) {$C\VIpr$};
   \node[vlab] at ($(W)+(68:1.10)$)  {$B\VIpr$};
\end{tikzpicture}
\figcap{6--6}
\end{minipage}\hfill
\begin{minipage}{0.36\linewidth}\centering
\begin{tikzpicture}[scale=0.832]
  \coordinate (O) at (0.72,0);
  \draw[given] (O) -- (0.10,1.62);   \node[vlab, above] at (0.10,1.62) {$P$};
  \draw[given] (O) -- (0.62,1.62);   \node[vlab, above] at (0.62,1.62) {$Q$};
  \draw[given] (O) -- (1.46,1.62);   \node[vlab, above] at (1.46,1.62) {$R$};
  \node[vlab, below] at (O) {$O$};
  \coordinate (S) at (1.72,0);
  \coordinate (T) at (2.10,0);
  \coordinate (U) at (2.72,0);
  \coordinate (Vv) at (2.16,1.62);
  \draw[given] (S) -- (Vv);
  \draw[given] (T) -- (Vv);
  \draw[given] (U) -- (Vv);
  \node[vlab, below] at (S) {$S$};
  \node[vlab, below] at (T) {$T$};
  \node[vlab, below] at (U) {$U$};
  \node[vlab, above] at (Vv) {$V$};
\end{tikzpicture}
\figcap{6--7}
\end{minipage}
\end{bkfigure}}

% =========================================================== Figs 6-8, 6-9
% Book p.87.  6-8: fan at A with rays to B, C, D, E.
%             6-9: fan at P with rays to R, S, T, U.
\newcommand{\FIGVIEIGHTNINE}{\begin{bkfigure}[\linewidth]
\begin{minipage}{0.48\linewidth}\centering
\begin{tikzpicture}[scale=1.086]
  \coordinate (A) at (1.30,0);
  \draw[given] (A) -- (0.10,1.28);  \node[vlab, above left]  at (0.10,1.28) {$B$};
  \draw[given] (A) -- (1.02,1.62);  \node[vlab, above] at (1.02,1.62) {$C$};
  \draw[given] (A) -- (1.42,1.62);  \node[vlab, above] at (1.42,1.62) {$D$};
  \draw[given] (A) -- (2.42,1.10);  \node[vlab, right] at (2.42,1.10) {$E$};
  \node[vlab, below] at (A) {$A$};
\end{tikzpicture}
\figcap{6--8}
\end{minipage}\hfill
\begin{minipage}{0.48\linewidth}\centering
\begin{tikzpicture}[scale=1.086]
  \coordinate (P) at (1.30,0);
  \draw[given] (P) -- (0.28,1.42);  \node[vlab, above left] at (0.28,1.42) {$R$};
  \draw[given] (P) -- (0.92,1.62);  \node[vlab, above] at (0.92,1.62) {$S$};
  \draw[given] (P) -- (1.52,1.62);  \node[vlab, above] at (1.52,1.62) {$T$};
  \draw[given] (P) -- (2.16,1.44);  \node[vlab, above right] at (2.16,1.44) {$U$};
  \node[vlab, below] at (P) {$P$};
\end{tikzpicture}
\figcap{6--9}
\end{minipage}
\end{bkfigure}}

% ==================================================== Figs 6-10, 6-11, 6-12
% Book p.87.  6-10: two supplement pictures (angle A with its adjacent
%   supplement B on a line; the primed copy).  6-11: two crossing lines with
%   vertical angles A and A', and B one of the others.  6-12: a right angle.
\newcommand{\FIGVITENELEVENTWELVE}{\begin{bkfigure}[\linewidth]
\begin{minipage}{0.50\linewidth}\centering
% In the book the ray rises to the RIGHT from a point on the line, so the
% OBTUSE angle (marked by the big arc sweeping over the top from the
% left-hand ray) is A and the small acute one tucked between the ray and the
% right-hand part of the line is B.  An earlier draft mirrored the ray and
% then put both letters in the same wedge, so the picture no longer showed a
% pair of supplements at all.
\begin{tikzpicture}
 \foreach \dx/\lA/\lB in {0/A/B, 3.30/A\VIpr/B\VIpr}{
  \begin{scope}[xshift=\dx cm]
   \coordinate (O)  at (0,0);
   \coordinate (Lt) at (-1.05,0);
   \coordinate (Rt) at (1.55,0);
   \coordinate (St) at (32:1.55);
   \draw[given] (Lt) -- (Rt);
   \draw[given] (O) -- (St);
   \VIarc{Lt}{O}{St}{0.62cm}   % the obtuse angle A
   \VIarc{St}{O}{Rt}{0.40cm}   % the acute angle B
   \node[slab] at ($(O)+(118:0.86)$) {$\lA$};
   \node[slab] at ($(O)+(16:1.02)$)  {$\lB$};
  \end{scope}}
\end{tikzpicture}
\figcap{6--10}
\end{minipage}\hfill
\begin{minipage}{0.30\linewidth}\centering
% Two crossing lines: A and A' are the vertical pair (left and right), B is
% one of the other two (the top one).  Each letter sits on the BISECTOR of
% its own wedge -- off the bisector a letter drifts into one of the four rays,
% which is what the collision audit was catching here.
\begin{tikzpicture}
  \coordinate (O)  at (0,0);
  \coordinate (P1) at (-1.30,-0.44);   \coordinate (P2) at (1.30,0.44);
  \coordinate (Q1) at (-1.30,0.44);    \coordinate (Q2) at (1.30,-0.44);
  \draw[given] (P1) -- (P2);
  \draw[given] (Q1) -- (Q2);
  \VIarc{Q1}{O}{P1}{0.60cm}   % left vertical angle  A
  \VIarc{P2}{O}{Q1}{0.90cm}   % top angle            B
  \VIarc{P2}{O}{Q2}{0.60cm}   % right vertical angle A'
  \node[slab] at ($(O)+(208:1.45)$) {$A$};
  \node[slab] at ($(O)+(90:0.70)$)  {$B$};
  \node[slab] at ($(O)+(-32:1.56)$)   {$A\VIpr$};
\end{tikzpicture}
\figcap{6--11}
\end{minipage}\hfill
\begin{minipage}{0.16\linewidth}\centering
\begin{tikzpicture}[scale=0.80]
  \coordinate (O) at (0.62,0);
  \coordinate (R) at (1.24,0);
  \coordinate (U) at (0.62,1.02);
  \draw[given] (0.02,0) -- (R);
  \draw[given] (O) -- (U);
  \VIrt{R}{O}{U}{6pt}
\end{tikzpicture}
\figcap{6--12}
\end{minipage}
\end{bkfigure}}

% ==================================================== Figs 6-13, 6-14, 6-15
% Book p.88, Exercise Group 6-4.
%  6-13: fan at A with rays to E, D, C, B; angle 1 = DAE, angle 2 = CAB.
%  6-14: two parallel lines cut by a transversal through A, B, C, D, E.
%  6-15: three concurrent lines through O: A--F, C--D, E--B.
\newcommand{\FIGVITHIRTEENFOURTEENFIFTEEN}{\begin{bkfigure}[\linewidth]
\begin{minipage}{0.46\linewidth}\centering
\begin{tikzpicture}[scale=1.009]
  \coordinate (A) at (1.42,0);
  \coordinate (E) at (0.05,1.02);
  \coordinate (D) at (0.72,1.52);
  \coordinate (C) at (2.02,1.52);
  \coordinate (B) at (2.62,1.02);
  \draw[given] (A) -- (E);  \draw[given] (A) -- (D);
  \draw[given] (A) -- (C);  \draw[given] (A) -- (B);
  \node[vlab, left]  at (E) {$E$};
  \node[vlab, above] at (D) {$D$};
  \node[vlab, above] at (C) {$C$};
  \node[vlab, right] at (B) {$B$};
  \node[vlab, below] at (A) {$A$};
  \VIarc{E}{A}{D}{0.86cm}
  \VIarc{C}{A}{B}{0.86cm}
  % Both wedges here are under 30 degrees wide, so a numeral set at a
  % hand-picked bearing lands on a ray; take the bisector instead.
  \VInumb{E}{A}{D}{1.10cm}{1}
  \VInumb{C}{A}{B}{1.10cm}{2}
\end{tikzpicture}
\figcap{6--13}

\medskip
% Rebuilt against scan06 p17 / photo p.88.  A, B, C, D, E are ANGLES, not
% points: the book draws THREE parallels cut by one transversal and strikes an
% arc on five of the angles so formed -- A below-left and B above-right at the
% lowest crossing, C below-right at the middle one, D below-right and E
% above-right at the top one.  (An earlier draft drew two parallels and
% lettered A..E as five POINTS, which left Exercise 3 -- "angle B is congruent
% to the supplement of angle C, and angle C is congruent to angle D; prove
% angle A is congruent to angle E" -- with nothing in the picture to name.)
\begin{tikzpicture}[scale=1.009]
  \coordinate (L1a) at (0.02,0.05);  \coordinate (L1b) at (3.20,0.05);
  \coordinate (L2a) at (0.02,1.00);  \coordinate (L2b) at (3.20,1.00);
  \coordinate (L3a) at (0.02,1.95);  \coordinate (L3b) at (3.20,1.95);
  \draw[given] (L1a) -- (L1b);
  \draw[given] (L2a) -- (L2b);
  \draw[given] (L3a) -- (L3b);
  \coordinate (P) at (1.02,-0.28);
  \coordinate (Q) at (2.86,2.24);
  \draw[given] (P) -- (Q);
  \coordinate (X1) at (intersection of P--Q and L1a--L1b);
  \coordinate (X2) at (intersection of P--Q and L2a--L2b);
  \coordinate (X3) at (intersection of P--Q and L3a--L3b);
  % arcs: the angle each letter names
  \VIarc{L1a}{X1}{P}{0.30cm}    % A -- below-left at the lowest crossing
  \VIarc{L1b}{X1}{Q}{0.34cm}    % B -- above-right there
  \VIarc{L2b}{X2}{P}{0.34cm}    % C -- below-right at the middle crossing
  \VIarc{L3b}{X3}{P}{0.34cm}    % D -- below-right at the top crossing
  \VIarc{L3b}{X3}{Q}{0.34cm}    % E -- above-right there
  \VInumb{L1a}{X1}{P}{0.56cm}{$A$}
  \VInumb{L1b}{X1}{Q}{0.60cm}{$B$}
  \VInumb{L2b}{X2}{P}{0.60cm}{$C$}
  \VInumb{L3b}{X3}{P}{0.60cm}{$D$}
  \VInumb{L3b}{X3}{Q}{0.62cm}{$E$}
\end{tikzpicture}
\figcap{6--14}
\end{minipage}\hfill
\begin{minipage}{0.46\linewidth}\centering
\begin{tikzpicture}[scale=1.009]
  % The six ends carry names so the constraint suite can read them straight
  % out of the drawing; the three lines really do meet at O.
  \coordinate (A) at (0.10,1.62);
  \coordinate (F) at (2.74,0.10);
  \coordinate (E) at (0.10,0.10);
  \coordinate (B) at (2.74,1.62);
  \coordinate (C) at (1.42,1.72);
  \coordinate (D) at (1.42,0.02);
  \coordinate (O) at (intersection of A--F and E--B);
  \draw[given] (A) -- (F);
  \draw[given] (E) -- (B);
  \draw[given] (C) -- (D);
  \node[vlab, above left]  at (A) {$A$};
  \node[vlab, above] at (C) {$C$};
  \node[vlab, above right] at (B) {$B$};
  \node[vlab, below left]  at (E) {$E$};
  \node[vlab, below] at (D) {$D$};
  \node[vlab, below right] at (F) {$F$};
  % Three lines meet at O, so six rays leave it; the letter goes on the
  % bisector of one of the 60-degree gaps rather than beside a ray.
  \node[vlab] at ($(O)+(68:0.64)$) {$O$};
\end{tikzpicture}
\figcap{6--15}
\end{minipage}
\end{bkfigure}}

% ================================================================= Fig 6-16
% Book p.89.  The S.A.S. postulate picture: triangles ABC and A'B'C' with one
% tick on AB/A'B', two on AC/A'C', and a double arc at A and A'.
\newcommand{\FIGVISIXTEEN}{\begin{bkfigure}[\linewidth]
\begin{tikzpicture}[scale=1.235]
 \foreach \dx/\lA/\lB/\lC in {0/A/B/C, 3.85/A\VIpr/B\VIpr/C\VIpr}{
  \begin{scope}[xshift=\dx cm]
   \coordinate (A) at (0,0);
   \coordinate (C) at (2.72,0);
   \coordinate (B) at (1.18,1.42);
   \draw[given] (A) -- (C) -- (B) -- cycle;
   \VItk{A}{B}{1}
   \VItk{A}{C}{2}
   \VIarcb{C}{A}{B}{0.52cm}
   \node[vlab, below left]  at (A) {$\lA$};
   \node[vlab, above] at (B) {$\lB$};
   \node[vlab, below right] at (C) {$\lC$};
  \end{scope}}
\end{tikzpicture}
\figcap{6--16}
\end{bkfigure}}

% =========================================================== Figs 6-17, 6-18
% Book p.89, Exercise Group 6-5.
%  6-17: triangles ABC (flat, opening right) and DEF.
%  6-18: triangles PQR and STU.
\newcommand{\FIGVISEVENTEENEIGHTEEN}{\begin{bkfigure}[\linewidth]
\begin{minipage}{0.50\linewidth}\centering
\begin{tikzpicture}[scale=0.894]
  \coordinate (A) at (0,0.62);
  \coordinate (B) at (1.02,0);
  \coordinate (C) at (1.92,0.68);
  \draw[given] (A) -- (B) -- (C) -- cycle;
  \node[vlab, left]  at (A) {$A$};
  \node[vlab, below] at (B) {$B$};
  \node[vlab, right] at (C) {$C$};
  % The two triangles are separate figures; keep a clear gap between C and E
  % so their letters cannot run together.
  \coordinate (E) at (2.62,0);
  \coordinate (D) at (3.28,1.42);
  \coordinate (F) at (4.24,1.10);
  \draw[given] (E) -- (D) -- (F) -- cycle;
  \node[vlab, below] at (E) {$E$};
  \node[vlab, above] at (D) {$D$};
  \node[vlab, right] at (F) {$F$};
\end{tikzpicture}
\figcap{6--17}
\end{minipage}\hfill
\begin{minipage}{0.46\linewidth}\centering
\begin{tikzpicture}[scale=0.894]
  \coordinate (P) at (0,0);
  \coordinate (R) at (1.32,0);
  \coordinate (Q) at (0.86,1.30);
  \draw[given] (P) -- (R) -- (Q) -- cycle;
  \node[vlab, left]  at (P) {$P$};
  \node[vlab, right] at (R) {$R$};
  \node[vlab, above] at (Q) {$Q$};
  % R and U were close enough that their letters merged into one word.
  \coordinate (U) at (2.14,0);
  \coordinate (S) at (2.04,1.30);
  \coordinate (T) at (3.38,1.30);
  \draw[given] (U) -- (S) -- (T) -- cycle;
  \node[vlab, below] at (U) {$U$};
  \node[vlab, above left]  at (S) {$S$};
  \node[vlab, above right] at (T) {$T$};
\end{tikzpicture}
\figcap{6--18}
\end{minipage}
\end{bkfigure}}

% =========================================================== Figs 6-19, 6-20
% Book p.89.  6-19: A above, B below it, with C and D at the ends -- the
%   "arrowhead" of Exercise 5 (AC = AD, angle CAB = angle BAD).
%  6-20: segments D--B and A--C crossing at O.
\newcommand{\FIGVININETEENTWENTY}{\begin{bkfigure}[\linewidth]
\begin{minipage}{0.46\linewidth}\centering
\begin{tikzpicture}[scale=0.966]
  \coordinate (A) at (1.42,1.52);
  \coordinate (B) at (1.42,0.52);
  \coordinate (C) at (0.02,0);
  \coordinate (D) at (2.82,0);
  \draw[given] (C) -- (A) -- (D);
  \draw[given] (C) -- (B) -- (D);
  \draw[given] (A) -- (B);
  \node[vlab, above] at (A) {$A$};
  \node[vlab, below] at (B) {$B$};
  \node[vlab, left]  at (C) {$C$};
  \node[vlab, right] at (D) {$D$};
\end{tikzpicture}
\figcap{6--19}
\end{minipage}\hfill
\begin{minipage}{0.50\linewidth}\centering
\begin{tikzpicture}[scale=0.966]
  \coordinate (D) at (0,0.52);
  \coordinate (B) at (3.02,0.42);
  \coordinate (A) at (0.86,0);
  \coordinate (C) at (2.42,1.02);
  \draw[given] (D) -- (B);
  \draw[given] (A) -- (C);
  \coordinate (O) at (intersection of D--B and A--C);
  \node[vlab, left]  at (D) {$D$};
  \node[vlab, right] at (B) {$B$};
  \node[vlab, below] at (A) {$A$};
  \node[vlab, above] at (C) {$C$};
  \node[vlab, below right] at (O) {$O$};
\end{tikzpicture}
\figcap{6--20}
\end{minipage}
\end{bkfigure}}

% =========================================================== Figs 6-21, 6-22
% Book p.90.  6-21: rectangle A B D C with the diagonal C--B.
%             6-22: triangle ACB with the cevian CD to the base.
\newcommand{\FIGVITWENTYONETWENTYTWO}{\begin{bkfigure}[\linewidth]
\begin{minipage}{0.50\linewidth}\centering
\begin{tikzpicture}[scale=0.991]
  \coordinate (A) at (0,0.92);
  \coordinate (B) at (2.92,0.92);
  \coordinate (C) at (0,0);
  \coordinate (D) at (2.92,0);
  \draw[given] (A) -- (B) -- (D) -- (C) -- cycle;
  \draw[given] (C) -- (B);
  \node[vlab, above left]  at (A) {$A$};
  \node[vlab, above right] at (B) {$B$};
  \node[vlab, below left]  at (C) {$C$};
  \node[vlab, below right] at (D) {$D$};
\end{tikzpicture}
\figcap{6--21}
\end{minipage}\hfill
\begin{minipage}{0.46\linewidth}\centering
\begin{tikzpicture}[scale=0.991]
  \coordinate (A) at (0,0);
  \coordinate (B) at (2.32,0);
  \coordinate (C) at (1.12,1.62);
  \coordinate (D) at ($(A)!0.48!(B)$);
  \draw[given] (A) -- (B) -- (C) -- cycle;
  \draw[given] (C) -- (D);
  \node[vlab, left]  at (A) {$A$};
  \node[vlab, right] at (B) {$B$};
  \node[vlab, above] at (C) {$C$};
  \node[vlab, below] at (D) {$D$};
\end{tikzpicture}
\figcap{6--22}
\end{minipage}
\end{bkfigure}}

% =========================================================== Figs 6-23, 6-24
% Book p.90.  6-23: A--B--C on a line with E and D above (two triangles meeting
%   at B).  6-24: quadrilateral A D C with B at the right, angles 1..4.
\newcommand{\FIGVITWENTYTHREETWENTYFOUR}{\begin{bkfigure}[\linewidth]
\begin{minipage}{0.52\linewidth}\centering
\begin{tikzpicture}[scale=0.940]
  \coordinate (A) at (0,0);
  \coordinate (B) at (1.62,0);
  \coordinate (C) at (3.12,0);
  \coordinate (E) at (0.86,1.22);
  \coordinate (D) at (2.02,1.28);
  \draw[given] (A) -- (C);
  \draw[given] (A) -- (E) -- (B);
  \draw[given] (B) -- (D) -- (C);
  \node[vlab, left]  at (A) {$A$};
  \node[vlab, below] at (B) {$B$};
  \node[vlab, right] at (C) {$C$};
  \node[vlab, above] at (E) {$E$};
  \node[vlab, above] at (D) {$D$};
\end{tikzpicture}
\figcap{6--23}
\end{minipage}\hfill
\begin{minipage}{0.44\linewidth}\centering
\begin{tikzpicture}[scale=1.490]
  \coordinate (A) at (0.72,1.52);
  \coordinate (C) at (0.72,0);
  % Exercise 10 gives AD = DC, so D must sit on the perpendicular bisector of
  % AC -- i.e. at the height of AC's midpoint.
  % Measured off scan06 p20 (book p.90): the book's quadrilateral is A D C B
  % with the single diagonal D--B.  There is NO segment A--C: D sits only a
  % little to the left of the line AC (offset 0.12 of AC's length), so the
  % angle ADC is nearly straight and its one big arc sweeps almost 180 deg.
  % An earlier draft drew A--C as well and pushed D out to 0.46 of AC, which
  % turned the plate into a diamond and buried the numerals 1 and 2 under the
  % spurious chord.
  \coordinate (D) at (0.54,0.76);
  \coordinate (B) at (2.62,0.76);
  \draw[given] (A) -- (D) -- (C);
  \draw[given] (A) -- (B) -- (C);
  \draw[given] (D) -- (B);
  \node[vlab, above] at (A) {$A$};
  \node[vlab, below] at (C) {$C$};
  \node[vlab, left]  at (D) {$D$};
  \node[vlab, right] at (B) {$B$};
  % D carries one big arc across the whole of angle ADC, which the segment DB
  % cuts into the parts 2 (upper) and 1 (lower); B carries a small arc for
  % each of 3 and 4.  Every numeral sits on its own bisector.
  \VIarc{A}{D}{C}{0.40cm}
  \VInumb{A}{D}{B}{0.62cm}{2}
  \VInumb{B}{D}{C}{0.62cm}{1}
  \VIarc{A}{B}{D}{0.52cm}   \VInumb{A}{B}{D}{1.20cm}{3}
  \VIarc{D}{B}{C}{0.52cm}   \VInumb{D}{B}{C}{1.20cm}{4}
\end{tikzpicture}
\figcap{6--24}
\end{minipage}
\end{bkfigure}}

% ==================================================== Figs 6-25, 6-26, 6-27
% Book p.90.
%  6-25: the zigzag B G C F H D A -- C and D are the upper vertices, G and H
%        the lower ones, with F between C and D on the upper line.
%  6-26: A B C on a line with peaks E and D; angles 1 and 2 at the ends.
%  6-27: triangle A M B with apex C, cevian CM, angles 1..4.
\newcommand{\FIGVITWENTYFIVETWENTYSIXTWENTYSEVEN}{\begin{bkfigure}[\linewidth]
\begin{tikzpicture}[scale=0.984]
  \coordinate (B) at (0,0);
  \coordinate (C) at (0.62,1.12);
  \coordinate (G) at (1.42,0);
  \coordinate (F) at (2.12,1.12);
  \coordinate (H) at (2.82,0);
  \coordinate (D) at (3.62,1.12);
  \coordinate (A) at (4.24,0);
  \draw[given] (B) -- (C) -- (G) -- (F) -- (H) -- (D) -- (A);
  \node[vlab, left]  at (B) {$B$};
  \node[vlab, above] at (C) {$C$};
  \node[vlab, below] at (G) {$G$};
  \node[vlab, above] at (F) {$F$};
  \node[vlab, below] at (H) {$H$};
  \node[vlab, above] at (D) {$D$};
  \node[vlab, right] at (A) {$A$};
  \VItk{C}{F}{2}
  \VItk{F}{D}{2}
  % Exercise 11 gives CG = HD as well; the book marks that pair with a single
  % tick, which an earlier draft left off.
  \VItk{C}{G}{1}
  \VItk{H}{D}{1}
  % Angle 3 is the one BETWEEN the top line and FG (i.e. angle CFG), not the
  % whole wedge GFH: the proof subtracts angle 1 from angle BCF to get angle
  % GCF, and it is the pair CFG / DFH that S.A.S. then makes congruent.
  \VIarc{B}{C}{G}{0.34cm}  \VInumb{B}{C}{G}{0.50cm}{1}
  \VIarc{C}{F}{G}{0.34cm}  \VInumb{C}{F}{G}{0.56cm}{3}
  \VIarc{H}{F}{D}{0.34cm}  \VInumb{H}{F}{D}{0.56cm}{4}
  \VIarc{H}{D}{A}{0.34cm}  \VInumb{H}{D}{A}{0.50cm}{2}
\end{tikzpicture}
\figcap{6--25}

\medskip
\begin{minipage}{0.52\linewidth}\centering
\begin{tikzpicture}[scale=0.938]
  % Exercise 12 gives AE = CD and AB = BC, so the two peaks are mirror
  % images about B and B is the midpoint of AC.  The draft had AB = 1.72
  % against BC = 1.70 and AE = 1.46 against CD = 1.60.
  \coordinate (A) at (0,0);
  \coordinate (C) at (3.42,0);
  \coordinate (B) at ($(A)!0.5!(C)$);
  \coordinate (E) at (0.62,1.32);
  \coordinate (D) at (2.80,1.32);
  \draw[given] (-0.28,0) -- (3.70,0);
  \draw[given] (A) -- (E) -- (B) -- (D) -- (C);
  \node[vlab] at ($(A)+(150:0.66)$) {$A$};
  \node[vlab, below] at (B) {$B$};
  \node[vlab] at ($(C)+(30:0.66)$) {$C$};
  \node[vlab, above] at (E) {$E$};
  \node[vlab, above] at (D) {$D$};
  % 1 and 2 are the EXTERIOR angles at the ends of the line -- between the
  % part of the line beyond A (resp. beyond C) and the side AE (resp. CD) --
  % and the book strikes an arc across each.  Those arcs were missing.
  \coordinate (Lx) at (-0.28,0);
  \coordinate (Rx) at (3.70,0);
  \VIarc{Lx}{A}{E}{0.30cm}   \VInum{A}{210}{0.52cm}{1}
  \VIarc{Rx}{C}{D}{0.30cm}   \VInum{C}{330}{0.52cm}{2}
\end{tikzpicture}
\figcap{6--26}
\end{minipage}\hfill
\begin{minipage}{0.44\linewidth}\centering
\begin{tikzpicture}[scale=0.938]
  \coordinate (A) at (0,0);
  \coordinate (B) at (2.62,0);
  \coordinate (C) at (1.32,1.72);
  \coordinate (M) at ($(A)!0.5!(B)$);
  % The book carries the base BEYOND A and beyond B, and angles 1 and 2 are
  % the EXTERIOR angles struck against those stubs -- not the interior base
  % angles an earlier draft marked (which, with AC = BC given, would have been
  % congruent by inspection and left Exercise 13 with nothing to prove).
  \coordinate (Lx) at (-0.52,0);
  \coordinate (Rx) at (3.14,0);
  \draw[given] (Lx) -- (Rx);
  \draw[given] (A) -- (C) -- (B);
  \draw[given] (C) -- (M);
  \node[vlab, below left]  at (A) {$A$};
  \node[vlab, below right] at (B) {$B$};
  \node[vlab, above] at (C) {$C$};
  \node[vlab, below] at (M) {$M$};
  \VIarc{A}{C}{M}{0.46cm}  \VInumb{A}{C}{M}{0.64cm}{3}
  \VIarc{M}{C}{B}{0.46cm}  \VInumb{M}{C}{B}{0.64cm}{4}
  \VIarc{Lx}{A}{C}{0.34cm} \VInumb{Lx}{A}{C}{0.60cm}{1}
  \VIarc{Rx}{B}{C}{0.34cm} \VInumb{Rx}{B}{C}{0.60cm}{2}
\end{tikzpicture}
\figcap{6--27}
\end{minipage}
\end{bkfigure}}

% =========================================================== Figs 6-28, 6-29
% Book p.91.  6-28: parallelogram-shaped quadrilateral D C B A with diagonal
%   D--B, angles 1 (at D) and 2 (at B).  6-29: triangle A B with apex B and
%   cevian BD to the base, angles 1 and 2 at D.
\newcommand{\FIGVITWENTYEIGHTTWENTYNINE}{\begin{bkfigure}[\linewidth]
\begin{minipage}{0.52\linewidth}\centering
\begin{tikzpicture}[scale=0.975]
  % Proportions off scan06 p21 (book p.91): the book's parallelogram is
  % width 1.8 x height, leaning right by 0.17 of its height.  The draft was
  % 2.3 x height and leaned 0.29, which flattened the plate enough to push
  % the two numerals into one another in the middle of the figure.
  \coordinate (A) at (0,0);
  \coordinate (B) at (2.62,0);
  \coordinate (D) at (0.25,1.45);
  \coordinate (C) at (2.87,1.45);
  \draw[given] (A) -- (B) -- (C) -- (D) -- cycle;
  \draw[given] (D) -- (B);
  \node[vlab, below left]  at (A) {$A$};
  \node[vlab, below right] at (B) {$B$};
  \node[vlab, above right] at (C) {$C$};
  \node[vlab, above left]  at (D) {$D$};
  \VIarc{C}{D}{B}{0.42cm}  \VInumb{C}{D}{B}{0.92cm}{1}
  \VIarc{A}{B}{D}{0.42cm}  \VInumb{A}{B}{D}{0.92cm}{2}
\end{tikzpicture}
\figcap{6--28}
\end{minipage}\hfill
\begin{minipage}{0.44\linewidth}\centering
\begin{tikzpicture}[scale=0.975]
  \coordinate (A) at (0,0);
  \coordinate (C) at (2.72,0);
  % Exercise 15 gives AD = DC and angle 1 = angle 2.  Angles 1 and 2 are the
  % two angles the downward stub of BD makes with AC, and they are
  % supplementary, so congruent means each is a right angle: B must stand
  % directly OVER the midpoint D, exactly as the book draws it.  The draft put
  % B at 0.34 of the base, which made the plate contradict its own hypothesis.
  \coordinate (B) at (1.36,1.62);
  % Exercise 15 gives AD = DC: D is the MIDPOINT of AC, not a point 0.42
  % along it, which is what the figure used to show.
  \coordinate (D) at ($(A)!0.5!(C)$);
  \draw[given] (A) -- (C) -- (B) -- cycle;
  \draw[given] (B) -- (D);
  \node[vlab, left]  at (A) {$A$};
  \node[vlab, right] at (C) {$C$};
  \node[vlab, above] at (B) {$B$};
  \node[vlab, above right] at (D) {$D$};
  % In the book BD is carried a little way BELOW the base, and the two
  % numbered angles are the ones under the line, spanned by a single arc.
  \draw[given] (D) -- ($(D)!-0.34!(B)$);
  % ADC is a straight angle, so which way \VIarc would sweep is a coin toss:
  % strike this one explicitly, below the base, as the book prints it.
  \draw[fig, line width=0.55pt] ([shift=(180:0.40cm)]D)
       arc[start angle=180, end angle=360, radius=0.40cm];
  \VInum{D}{225}{0.62cm}{1}
  \VInum{D}{315}{0.62cm}{2}
\end{tikzpicture}
\figcap{6--29}
\end{minipage}
\end{bkfigure}}

% =========================================================== Figs 6-30, 6-31
% Book p.91.  6-30: quadrilateral A B C D with E and F interior on the
%   diagonals; angles 1..6.  6-31: quadrilateral A B C D with centre O.
\newcommand{\FIGVITHIRTYTHIRTYONE}{\begin{bkfigure}[\linewidth]
\begin{minipage}{0.58\linewidth}\centering
% The book draws ONE diagonal here, A--C, and hangs two segments off it:
% B--F and E--D, with F and E on AC.  (An earlier draft drew the second
% diagonal B--D instead and left BF and ED out altogether, which made the
% six numbered angles label things that were not in the picture.)
% Exercise *16 gives AE = CF, so E and F sit symmetrically about the midpoint
% of AC -- placed here by ratio, not by eye.
\begin{tikzpicture}[scale=1.046]
  \coordinate (A) at (0,0);
  \coordinate (D) at (3.30,0);
  \coordinate (B) at (0.75,1.45);
  \coordinate (C) at (4.05,1.45);
  \draw[given] (A) -- (D) -- (C) -- (B) -- cycle;
  \draw[given] (A) -- (C);
  \coordinate (F) at ($(A)!0.33!(C)$);
  \coordinate (E) at ($(A)!0.67!(C)$);
  \draw[given] (B) -- (F);
  \draw[given] (E) -- (D);
  \node[vlab, below left]  at (A) {$A$};
  \node[vlab, below right] at (D) {$D$};
  \node[vlab, above left]  at (B) {$B$};
  \node[vlab, above right] at (C) {$C$};
  % E above-left of its point as the book sets it; F sits low on AC, close
  % to the edge AD, so its letter goes up into the wedge between AC and
  % the segment F--B.
  \node[vlab, above left] at (E) {$E$};
  \node[vlab] at ($(F)+(60:0.38)$) {$F$};
  % Radii are set per wedge: the two slivers at A and C (only about 20 degrees
  % wide, because BC is parallel to AD) need the numeral pushed well out
  % before it clears both sides, while the open wedges at F and E do not.
  \VIarc{B}{A}{C}{0.36cm}   \VInumb{B}{A}{C}{0.78cm}{3}
  \VIarc{C}{A}{D}{0.52cm}   \VInumb{C}{A}{D}{1.30cm}{6}
  \VIarc{A}{F}{B}{0.26cm}   \VInumb{A}{F}{B}{0.44cm}{1}
  \VIarc{C}{E}{D}{0.24cm}   \VInumb{C}{E}{D}{0.42cm}{2}
  \VIarc{B}{C}{A}{0.42cm}   \VInumb{B}{C}{A}{1.30cm}{5}
  \VIarc{A}{C}{D}{0.30cm}   \VInumb{A}{C}{D}{0.55cm}{4}
\end{tikzpicture}
\figcap{6--30}
\end{minipage}\hfill
\begin{minipage}{0.42\linewidth}\centering
\begin{tikzpicture}[scale=1.360]
  % Exercise *17 gives AB = AD and BO = OD, so B and D must sit symmetrically
  % about the line AC: B and D share an x, and AC passes through their midpoint.
  % Proportions off the plate: AO is about 2.5 times OC, and OB = OD is about
  % 0.37 of AO -- the book's kite is much longer in the nose than a first
  % draft suggests.
  \coordinate (A) at (0,0.68);
  \coordinate (B) at (1.85,1.36);
  \coordinate (C) at (2.60,0.68);
  \coordinate (D) at (1.85,0);
  \draw[given] (A) -- (B) -- (C) -- (D) -- cycle;
  \draw[given] (B) -- (D);
  \coordinate (O) at (intersection of A--C and B--D);
  \draw[given] (A) -- (C);
  \node[vlab, left]  at (A) {$A$};
  \node[vlab, above] at (B) {$B$};
  \node[vlab, right] at (C) {$C$};
  \node[vlab, below] at (D) {$D$};
  \node[vlab, left] at ($(O)+(0.02,-0.20)$) {$O$};
  \VItk{A}{B}{1}
  \VItk{A}{D}{1}
  \VItk{O}{B}{2}
  \VItk{O}{D}{2}
  % The pair of congruent angles at A, marked as the book marks them.  These
  % arcs used to be omitted because the old pic-based mark wrapped to a full
  % circle whenever the sides were named the other way round; the hand-struck
  % \VIarc has no such failure mode.
  \VIarc{B}{A}{C}{0.85cm}
  \VIarc{C}{A}{D}{0.85cm}
  % OC is only 0.75 long, so a numeral 0.68 out along the 45-degree bisector
  % lands OUTSIDE the kite, across BC.  The book keeps 1 and 2 tucked in
  % beside the double ticks, about 0.19 of AO from the crossing.
  \VIarc{B}{O}{C}{0.18cm}   \VInum{O}{35}{0.36cm}{1}
  \VIarc{C}{O}{D}{0.18cm}   \VInum{O}{-35}{0.36cm}{2}
\end{tikzpicture}
\figcap{6--31}
\end{minipage}
\end{bkfigure}}

% ================================================================= Fig 6-32
% Book p.92.  The isosceles triangle labelled two ways at once, which is the
% whole trick of the Theorem 6-3 proof.
\newcommand{\FIGVITHIRTYTWO}{\begin{bkfigure}[\linewidth]
\begin{tikzpicture}[scale=1.235]
  \coordinate (A) at (1.42,1.62);
  \coordinate (B) at (0,0);
  \coordinate (C) at (2.84,0);
  \draw[given] (A) -- (B) -- (C) -- cycle;
  \node[vlab, above] at (A) {$A = A\VIpr$};
  \node[vlab, below left]  at (B) {$B = C\VIpr$};
  \node[vlab, below right] at (C) {$C = B\VIpr$};
\end{tikzpicture}
\figcap{6--32}
\end{bkfigure}}

% ==================================================== Figs 6-33, 6-34, 6-35
% Book p.92, Exercise Group 6-6.
%  6-33: triangle P Q R.
%  6-34: the kite A B C D with both diagonals drawn.
%  6-35: triangle A D B with apex C and cevian CD, angles 1 and 2.
\newcommand{\FIGVITHIRTYTHREETHIRTYFOURTHIRTYFIVE}{\begin{bkfigure}[\linewidth]
\begin{minipage}{0.46\linewidth}\centering
\begin{tikzpicture}[scale=0.961]
  \coordinate (Q) at (0,0);
  \coordinate (R) at (2.62,0);
  \coordinate (P) at (1.22,1.12);
  \draw[given] (Q) -- (R) -- (P) -- cycle;
  \node[vlab, left]  at (Q) {$Q$};
  \node[vlab, right] at (R) {$R$};
  \node[vlab, above] at (P) {$P$};
\end{tikzpicture}
\figcap{6--33}

\medskip
\begin{tikzpicture}[scale=0.961]
  \coordinate (A) at (1.32,2.02);
  \coordinate (B) at (0,1.32);
  \coordinate (C) at (2.64,1.32);
  \coordinate (D) at (1.32,0);
  \draw[given] (A) -- (B) -- (D) -- (C) -- cycle;
  \draw[given] (B) -- (C);
  \draw[given] (A) -- (D);
  \node[vlab, above] at (A) {$A$};
  \node[vlab, left]  at (B) {$B$};
  \node[vlab, right] at (C) {$C$};
  \node[vlab, below] at (D) {$D$};
\end{tikzpicture}
\figcap{6--34}
\end{minipage}\hfill
\begin{minipage}{0.46\linewidth}\centering
\begin{tikzpicture}[scale=0.961]
  \coordinate (A) at (0,0);
  \coordinate (B) at (2.82,0);
  \coordinate (C) at (1.32,1.72);
  \coordinate (D) at ($(A)!0.47!(B)$);
  \draw[given] (A) -- (B) -- (C) -- cycle;
  \draw[given] (C) -- (D);
  \node[vlab, left]  at (A) {$A$};
  \node[vlab, right] at (B) {$B$};
  \node[vlab, above] at (C) {$C$};
  \node[vlab, below] at (D) {$D$};
  \VIarc{A}{C}{D}{0.46cm}  \VInumb{A}{C}{D}{0.76cm}{1}
  \VIarc{D}{C}{B}{0.46cm}  \VInumb{D}{C}{B}{0.72cm}{2}
\end{tikzpicture}
\figcap{6--35}
\end{minipage}
\end{bkfigure}}

% =========================================================== Figs 6-36, 6-37
% Book p.93.  6-36: the "dart" A B C E D -- AD is the long spine, B interior.
%             6-37: A B C D E F -- two triangles meeting at C/D, angles 1, 2.
\newcommand{\FIGVITHIRTYSIXTHIRTYSEVEN}{\begin{bkfigure}[\linewidth]
\begin{minipage}{0.50\linewidth}\centering
\begin{tikzpicture}[scale=0.972]
  \coordinate (A) at (0,0.86);
  \coordinate (C) at (1.02,1.62);
  \coordinate (E) at (1.02,0);
  \coordinate (D) at (2.72,0.86);
  \coordinate (B) at ($(A)!0.42!(D)$);
  \draw[given] (A) -- (D);
  \draw[given] (C) -- (B) -- (E);
  \draw[given] (C) -- (D) -- (E);
  \node[vlab, left]  at (A) {$A$};
  \node[vlab, above] at (C) {$C$};
  \node[vlab, below] at (E) {$E$};
  \node[vlab, right] at (D) {$D$};
  \node[vlab, below left] at (B) {$B$};
\end{tikzpicture}
\figcap{6--36}
\end{minipage}\hfill
\begin{minipage}{0.46\linewidth}\centering
\begin{tikzpicture}[scale=0.972]
  \coordinate (A) at (0,0.72);
  \coordinate (B) at (0.92,1.82);
  \coordinate (F) at (2.92,0.72);
  \coordinate (E) at (1.92,-0.22);
  % AC = DF is given, so C and D sit symmetrically on A--F.
  \coordinate (C) at ($(A)!0.35!(F)$);
  \coordinate (D) at ($(A)!0.65!(F)$);
  \draw[given] (A) -- (F);
  \draw[given] (A) -- (B) -- (D);
  \draw[given] (C) -- (E) -- (F);
  \node[vlab, left]  at (A) {$A$};
  \node[vlab, above] at (B) {$B$};
  \node[vlab, right] at (F) {$F$};
  \node[vlab, below] at (E) {$E$};
  \node[vlab, below left]  at (C) {$C$};
  \node[vlab, above right] at (D) {$D$};
  % Angle 2 is at D (between DB and DA), NOT at B: the S.A.S. pair the
  % exercise builds is (AD, DB, angle 2) against (FC, CE, angle 1), so both
  % marked angles sit on the line A--F.  The book strikes an arc on each; the
  % draft put 2 at vertex B and left both arcs off.
  \VIarc{A}{D}{B}{0.34cm}   \VInumb{A}{D}{B}{0.80cm}{2}
  \VIarc{F}{C}{E}{0.34cm}   \VInumb{F}{C}{E}{0.80cm}{1}
\end{tikzpicture}
\figcap{6--37}
\end{minipage}
\end{bkfigure}}

% =========================================================== Figs 6-38, 6-39
% Book p.93.  6-38: triangle A B C with E and D on the base and cevians from B.
%             6-39: triangle R S T with U inside, joined to S and T.
\newcommand{\FIGVITHIRTYEIGHTTHIRTYNINE}{\begin{bkfigure}[\linewidth]
\begin{minipage}{0.50\linewidth}\centering
\begin{tikzpicture}[scale=0.963]
  \coordinate (A) at (0,0);
  \coordinate (C) at (3.22,0);
  % Exercise 7 gives AD = EC and BE = BD.  The first puts E and D
  % symmetrically about the midpoint of AC; the second then puts B over that
  % midpoint.  The draft had E at 0.36, D at 0.56 and B at 0.41 of the base,
  % so neither hypothesis held in the drawing.
  \coordinate (B) at (1.61,1.62);
  \coordinate (E) at ($(A)!0.42!(C)$);
  \coordinate (D) at ($(A)!0.58!(C)$);
  \draw[given] (A) -- (C) -- (B) -- cycle;
  \draw[given] (B) -- (E);
  \draw[given] (B) -- (D);
  \node[vlab, left]  at (A) {$A$};
  \node[vlab, right] at (C) {$C$};
  \node[vlab, above] at (B) {$B$};
  \node[vlab, below] at (E) {$E$};
  \node[vlab, below] at (D) {$D$};
\end{tikzpicture}
\figcap{6--38}
\end{minipage}\hfill
\begin{minipage}{0.46\linewidth}\centering
\begin{tikzpicture}[scale=0.963]
  \coordinate (S) at (0,0);
  \coordinate (T) at (2.72,0);
  % RS = RT and US = UT are both given, so R and U must stand on the
  % perpendicular bisector of ST -- x = 1.36, not 1.22.
  \coordinate (R) at (1.36,1.72);
  \coordinate (U) at (1.36,0.92);
  \draw[given] (S) -- (T) -- (R) -- cycle;
  \draw[given] (S) -- (U) -- (T);
  \node[vlab, left]  at (S) {$S$};
  \node[vlab, right] at (T) {$T$};
  \node[vlab, above] at (R) {$R$};
  \node[vlab, below] at (U) {$U$};
\end{tikzpicture}
\figcap{6--39}
\end{minipage}
\end{bkfigure}}

% ==================================================== Figs 6-40, 6-41, 6-42
% Book p.93.
%  6-40: triangle A B with apex C; D and E on the sides, cevians AE and BD,
%        angles 1 and 2 at D and E.
%  6-41: triangle A D with apex E; B, F, C on the base, cevians from E.
%  6-42: triangle A B with apex C; G, D, O, E, F interior, angles 1..4.
\newcommand{\FIGVIFORTYFORTYONEFORTYTWO}{\begin{bkfigure}[\linewidth]
\begin{minipage}{0.46\linewidth}\centering
% Measured off the plate on book p.93 (photo PDF p.107): base 645 px, apex over
% the midpoint (AC = BC, as Exercise 9 requires), height/base = 0.674, D and E
% about 0.60 of the way up their sides.  The book sets D level with its point
% on the LEFT and E level with its point on the RIGHT -- not above them -- and
% leaves a clear digit's width between the numerals 1 and 2.
% Drawn at 1.05 rather than 0.80: D and E are only 1.17 units apart, and the
% two bisector-placed numerals aim at each other, so at the old scale the
% numerals had to sit so close to their vertices that both grazed the cevians
% (0.27 pt).  Widening the plate buys the clearance without moving either
% numeral off its own bisector.
\begin{tikzpicture}[scale=0.927]
  \coordinate (A) at (0,0);
  \coordinate (B) at (2.92,0);
  \coordinate (C) at (1.46,1.97);
  \coordinate (D) at ($(A)!0.60!(C)$);
  \coordinate (E) at ($(B)!0.60!(C)$);
  \draw[given] (A) -- (B) -- (C) -- cycle;
  \draw[given] (A) -- (E);
  \draw[given] (B) -- (D);
  \node[vlab, left]  at (A) {$A$};
  \node[vlab, right] at (B) {$B$};
  \node[vlab, above] at (C) {$C$};
  \node[vlab, left]  at (D) {$D$};
  \node[vlab, right] at (E) {$E$};
  % Angle 1 is CDB and angle 2 is CEA -- the pair the exercise asks about.
  % Both go on their own bisector, which is what keeps them off DB and EA.
  % The two bisectors point at EACH OTHER, so distance along them is what sets
  % the gap between the numerals: at 0.48/0.46 they closed to 1.2pt and read as
  % the single number "12".  Pulled back to 0.40 for the daylight the book
  % leaves (about one digit width); still clear of the 0.26 arcs.
  \VIarc{C}{D}{B}{0.26cm}   \VInumb{C}{D}{B}{0.40cm}{1}
  \VIarc{C}{E}{A}{0.26cm}   \VInumb{C}{E}{A}{0.40cm}{2}
\end{tikzpicture}
\figcap{6--40}

\medskip
% Rebuilt against scan07 p1 (book p.93) at 150 dpi, cross-read with the photo
% plate.  The earlier draft fanned every cevian from the APEX E, which cannot
% satisfy Exercise 10 at all: its hypotheses are AO = OD, AF = FD, BO = OC,
% BF = FC, so O and F are both distinct interior points joined to A, B, C, D.
% The plate agrees -- at A and at D two lines leave the vertex (one to E, one
% to O), and six lines meet at O (up to E, down to A, B, F, C, D).
% Structure: base A-B-F-C-D collinear; F the MIDPOINT of AD (AF = FD) and of
% BC (BF = FC); E directly above F, so the whole plate is mirror-symmetric
% about FE and AO = OD, BO = OC hold exactly rather than by eye.
% Proportions off the scan: base 825 px, F at 0.488 (drawn 0.5, as the
% exercise demands), B at 0.228 and C at 0.737 (drawn symmetric at
% 0.245/0.755), O 0.40 of the base above it, E 0.70 of the base above it.
\begin{tikzpicture}[scale=1.131]
  \coordinate (A) at (0,0);
  \coordinate (D) at (3.22,0);
  \coordinate (F) at ($(A)!0.5!(D)$);      % AF = FD
  \coordinate (B) at ($(A)!0.245!(D)$);
  \coordinate (C) at ($(A)!0.755!(D)$);    % BF = FC by symmetry about F
  \coordinate (E) at (1.61,2.254);         % directly above F
  % Re-measured on the plate: O sits 0.636 of the way from F up to E, not 0.593.
  \coordinate (O) at ($(F)!0.636!(E)$);    % on FE, so AO = OD and BO = OC
  \draw[given] (A) -- (D) -- (E) -- cycle;
  \draw[given] (E) -- (F);
  \draw[given] (O) -- (A);
  \draw[given] (O) -- (B);
  \draw[given] (O) -- (C);
  \draw[given] (O) -- (D);
  \node[vlab, left]  at (A) {$A$};
  \node[vlab, right] at (D) {$D$};
  \node[vlab, above] at (E) {$E$};
  \node[vlab, below] at (B) {$B$};
  \node[vlab, below] at (F) {$F$};
  \node[vlab, below] at (C) {$C$};
  % Six lines meet at O; the only clear quadrant is up-and-right, between OE
  % and OD, which is exactly where the book sets the letter.  That quadrant is
  % a wedge closing toward E -- bounded on the left by the vertical EF and on
  % the right by the side ED -- so the letter has to hug the vertical and stay
  % LOW, which is exactly what the plate does (its O starts about 2.5% of the
  % base to the right of the axis).  The default 4.6pt outer sep pushed the
  % box out along the 45-degree diagonal instead, straight into ED.
  \node[vlab, above right, inner sep=1pt, outer sep=2.4pt] at (O) {$O$};
\end{tikzpicture}
\figcap{6--41}
\end{minipage}\hfill
\begin{minipage}{0.50\linewidth}\centering
% Rebuilt against scan07 p1 (book p.93) at 200 dpi.  The cevians the book
% draws are A--F and B--G (F on CB, G on CA), crossing at O.  D and E are NOT
% on the base -- an earlier draft put them there, which is the one structural
% error the plate contradicts outright: the book marks D on the cevian A--F
% and E on the cevian B--G, and joins each back to C.  Measured off the plate,
% each sits 0.40 of the way from its base vertex along that cevian.
% C is placed exactly over the midpoint of AB so that AC = BC holds, and G, F
% are taken at the same ratio down the two sides so that CG = CF holds -- both
% are hypotheses of Exercise 11.
% Angles 1 and 2 sit at G and F; 3 and 4 where the cevians CD and CE cut
% BG and AF, which is where the book strikes their arcs.
% Drawn a little larger than the neighbouring plates: four numbered angles
% share the narrow band between G, P, Q and F, and at the smaller scale their
% bisector labels ran into one another in the middle of the triangle.
\begin{tikzpicture}[scale=1.148]
  \coordinate (A) at (0,0);
  \coordinate (B) at (4.30,0);
  \coordinate (C) at (2.15,2.64);
  \coordinate (G) at ($(C)!0.45!(A)$);
  \coordinate (F) at ($(C)!0.45!(B)$);
  \draw[given] (A) -- (B) -- (C) -- cycle;
  \draw[given] (A) -- (F);
  \draw[given] (B) -- (G);
  % Measured off scan07 p1: D and E sit 0.47 of the way along their cevians,
  % not 0.40.
  \coordinate (D) at ($(A)!0.47!(F)$);
  \coordinate (E) at ($(B)!0.47!(G)$);
  \draw[given] (C) -- (D);
  \draw[given] (C) -- (E);
  \coordinate (O) at (intersection of A--F and B--G);
  \coordinate (P) at (intersection of C--D and B--G);
  \coordinate (Q) at (intersection of C--E and A--F);
  \node[vlab, left]  at (A) {$A$};
  \node[vlab, right] at (B) {$B$};
  \node[vlab, above] at (C) {$C$};
  \node[vlab, below]  at (D) {$D$};
  \node[vlab, below] at (E) {$E$};
  \node[vlab, above left]  at (G) {$G$};
  \node[vlab, above right] at (F) {$F$};
  \node[vlab, below] at (O) {$O$};
  \VIarc{C}{G}{B}{0.20cm}   \VInumb{C}{G}{B}{0.38cm}{1}
  \VIarc{C}{F}{A}{0.20cm}   \VInumb{C}{F}{A}{0.38cm}{2}
  % Same trap as 6-40, and worse here because scale=1.30 multiplies the cm
  % distance too: the bisectors at P and Q aim straight at one another, so
  % 0.48cm drove both numerals onto the figure's axis and they printed as "34"
  % right on top of O.  0.36cm restores the gap the plate shows between them.
  \VIarc{C}{P}{B}{0.22cm}   \VInumb{C}{P}{B}{0.36cm}{3}
  \VIarc{C}{Q}{A}{0.22cm}   \VInumb{C}{Q}{A}{0.36cm}{4}
\end{tikzpicture}
\figcap{6--42}
\end{minipage}
\end{bkfigure}}

% ================================================================= Fig 6-43
% Book p.94.  The Theorem 6-4 proof picture: triangle ABC, and A'B'C' with the
% constructed point D' on ray C'B' (dashed, since it is added in the proof).
% The plate sets all four base-vertex letters BELOW the base line -- A and C
% under the first triangle, A' and C' under the second.  The draft lifted C and
% A' ABOVE it, which put them at the same height in the gap between the two
% triangles, where they read as one floating pair "C A'" belonging to neither.
% Panels also pulled 0.4cm further apart for the daylight the plate shows.
\newcommand{\FIGVIFORTYTHREE}{\begin{bkfigure}[\linewidth]
\begin{tikzpicture}[scale=1.235]
  \coordinate (A) at (0,0);
  \coordinate (C) at (2.52,0);
  \coordinate (B) at (1.02,1.42);
  \draw[given] (A) -- (C) -- (B) -- cycle;
  \VItk{A}{B}{1}  \VItk{A}{C}{2}  \VIarc{C}{A}{B}{0.46cm}
  \node[vlab, below left]  at (A) {$A$};
  \node[vlab, above] at (B) {$B$};
  \node[vlab, below right] at (C) {$C$};
  \begin{scope}[xshift=4.05cm]
   \coordinate (Ap) at (0,0);
   \coordinate (Cp) at (2.52,0);
   \coordinate (Bp) at (1.02,1.42);
   \coordinate (Dp) at ($(Cp)!0.30!(Bp)$);
   \draw[given] (Ap) -- (Cp) -- (Bp) -- cycle;
   \draw[aux]   (Ap) -- (Dp);
   \VItk{Ap}{Bp}{1}  \VItk{Ap}{Cp}{2}  \VIarc{Cp}{Ap}{Bp}{0.46cm}
   \node[vlab, below left]  at (Ap) {$A\VIpr$};
   \node[vlab, above] at (Bp) {$B\VIpr$};
   \node[vlab, below right] at (Cp) {$C\VIpr$};
   \node[vlab, above right] at (Dp) {$D\VIpr$};
  \end{scope}
\end{tikzpicture}
\figcap{6--43}
\end{bkfigure}}

% =========================================================== Figs 6-44, 6-45
% Book p.95, Exercise Group 6-8.
%  6-44: the rhombus A B C D with diagonal BD, angles 1 and 2 at B.
%  6-45: the parallelogram E F H G with diagonal E--H, angles 1 and 2.
\newcommand{\FIGVIFORTYFOURFORTYFIVE}{\begin{bkfigure}[\linewidth]
\begin{minipage}{0.46\linewidth}\centering
\begin{tikzpicture}[scale=0.982]
  % Off scan07 p3 the book's rhombus is very nearly square-on-its-corner
  % (half-width 117 against half-height 130); the draft was half again as
  % wide as it was tall.
  \coordinate (B) at (1.17,2.28);
  \coordinate (A) at (0,1.08);
  \coordinate (C) at (2.34,1.08);
  \coordinate (D) at (1.17,0);
  \draw[given] (A) -- (B) -- (C) -- (D) -- cycle;
  \draw[given] (B) -- (D);
  \node[vlab, above] at (B) {$B$};
  \node[vlab, left]  at (A) {$A$};
  \node[vlab, right] at (C) {$C$};
  \node[vlab, below] at (D) {$D$};
  \VInumb{A}{B}{D}{0.56cm}{1}
  \VInumb{D}{B}{C}{0.54cm}{2}
\end{tikzpicture}
\figcap{6--44}
\end{minipage}\hfill
\begin{minipage}{0.50\linewidth}\centering
\begin{tikzpicture}[scale=0.982]
  \coordinate (G) at (0,0);
  \coordinate (H) at (2.62,0);
  \coordinate (E) at (0.52,1.12);
  \coordinate (F) at (3.14,1.12);
  \draw[given] (G) -- (H) -- (F) -- (E) -- cycle;
  \draw[given] (E) -- (H);
  \node[vlab, below left]  at (G) {$G$};
  \node[vlab, below right] at (H) {$H$};
  \node[vlab, above left]  at (E) {$E$};
  \node[vlab, above right] at (F) {$F$};
  \VInumb{G}{E}{H}{0.38cm}{1}
  \VInumb{F}{H}{E}{0.38cm}{2}
\end{tikzpicture}
\figcap{6--45}
\end{minipage}
\end{bkfigure}}

% =========================================================== Figs 6-46, 6-47
% Book p.95.  6-46: triangle C D with apex A and cevian AB to the base.
%             6-47: triangle B D with apex A and cevian AC to the base.
\newcommand{\FIGVIFORTYSIXFORTYSEVEN}{\begin{bkfigure}[\linewidth]
\begin{minipage}{0.48\linewidth}\centering
\begin{tikzpicture}[scale=0.983]
  \coordinate (C) at (0,0);
  \coordinate (D) at (2.72,0);
  \coordinate (A) at (1.36,1.62);
  \coordinate (B) at ($(C)!0.5!(D)$);
  \draw[given] (C) -- (D) -- (A) -- cycle;
  \draw[given] (A) -- (B);
  \node[vlab, left]  at (C) {$C$};
  \node[vlab, right] at (D) {$D$};
  \node[vlab, above] at (A) {$A$};
  \node[vlab, below] at (B) {$B$};
\end{tikzpicture}
\figcap{6--46}
\end{minipage}\hfill
\begin{minipage}{0.48\linewidth}\centering
\begin{tikzpicture}[scale=0.983]
  \coordinate (B) at (0,0);
  \coordinate (D) at (2.72,0);
  % Exercise 4 gives AB = AD and BC = CD, so C is the MIDPOINT of BD and A
  % stands directly over it -- which is how the book draws it.  The draft had
  % C at 0.46 and A off to the left of it.
  \coordinate (C) at ($(B)!0.5!(D)$);
  \coordinate (A) at (1.36,1.62);
  \draw[given] (B) -- (D) -- (A) -- cycle;
  \draw[given] (A) -- (C);
  \node[vlab, left]  at (B) {$B$};
  \node[vlab, right] at (D) {$D$};
  \node[vlab, above] at (A) {$A$};
  \node[vlab, below] at (C) {$C$};
\end{tikzpicture}
\figcap{6--47}
\end{minipage}
\end{bkfigure}}

% =========================================================== Figs 6-48, 6-49
% Book p.96.  6-48: triangle A B C with D, E on the sides and both cevians
%             BE, CD drawn.  6-49: the kite B A D C with both diagonals.
\newcommand{\FIGVIFORTYEIGHTFORTYNINE}{\begin{bkfigure}[\linewidth]
\begin{minipage}{0.52\linewidth}\centering
\begin{tikzpicture}[scale=1.027]
  \coordinate (B) at (0,0);
  \coordinate (C) at (3.22,0);
  \coordinate (A) at (1.61,1.62);
  % Exercise 5 makes D and E the MIDPOINTS of AB and AC.
  \coordinate (D) at ($(B)!0.5!(A)$);
  \coordinate (E) at ($(C)!0.5!(A)$);
  \draw[given] (B) -- (C) -- (A) -- cycle;
  \draw[given] (B) -- (E);
  \draw[given] (C) -- (D);
  \draw[given] (D) -- (E);
  \node[vlab, left]  at (B) {$B$};
  \node[vlab, right] at (C) {$C$};
  \node[vlab, above] at (A) {$A$};
  \node[vlab, above left]  at (D) {$D$};
  \node[vlab, above right] at (E) {$E$};
\end{tikzpicture}
\figcap{6--48}
\end{minipage}\hfill
\begin{minipage}{0.42\linewidth}\centering
\begin{tikzpicture}[scale=1.027]
  % Off scan07 p3 the book's kite is TALL and narrow: half-width 0.85 of the
  % short (upper) axis, against a lower axis twice the upper.  The draft was
  % nearly twice as wide as the upper axis was long.
  \coordinate (A) at (1.22,1.92);
  \coordinate (B) at (0.68,1.28);
  \coordinate (D) at (1.76,1.28);
  \coordinate (C) at (1.22,0);
  \draw[given] (B) -- (A) -- (D) -- (C) -- cycle;
  \draw[given] (B) -- (D);
  \draw[given] (A) -- (C);
  \node[vlab, above] at (A) {$A$};
  \node[vlab, left]  at (B) {$B$};
  \node[vlab, right] at (D) {$D$};
  \node[vlab, below] at (C) {$C$};
\end{tikzpicture}
\figcap{6--49}
\end{minipage}
\end{bkfigure}}

% =========================================================== Figs 6-50, 6-51
% Book p.96.  6-50: triangle A B C, D and E midpoints of AB and BC, with the
%   crossing cevians AE and CD.  6-51: A B C on a line with peaks E and D.
\newcommand{\FIGVIFIFTYFIFTYONE}{\begin{bkfigure}[\linewidth]
\begin{minipage}{0.48\linewidth}\centering
\begin{tikzpicture}[scale=0.960]
  \coordinate (A) at (0,0);
  \coordinate (C) at (2.92,0);
  % Exercise 7 gives AB = BC, so B stands over the midpoint of AC; the draft
  % had it at 0.36 of the base and the drawing was not isosceles at all.
  \coordinate (B) at (1.46,1.62);
  \coordinate (D) at ($(A)!0.5!(B)$);
  \coordinate (E) at ($(B)!0.5!(C)$);
  \draw[given] (A) -- (C) -- (B) -- cycle;
  \draw[given] (A) -- (E);
  \draw[given] (C) -- (D);
  \node[vlab, left]  at (A) {$A$};
  \node[vlab, right] at (C) {$C$};
  \node[vlab, above] at (B) {$B$};
  \node[vlab, left]  at (D) {$D$};
  \node[vlab, above right] at (E) {$E$};
\end{tikzpicture}
\figcap{6--50}
\end{minipage}\hfill
\begin{minipage}{0.48\linewidth}\centering
\begin{tikzpicture}[scale=0.960]
  \coordinate (A) at (0,0);
  \coordinate (C) at (3.02,0);
  % Exercise 8 makes B the midpoint of AC and gives EB = BD, so E and D must
  % sit symmetrically about B.  The draft's E and D were 0.76 and 0.50 out.
  \coordinate (B) at ($(A)!0.5!(C)$);
  \coordinate (E) at (0.88,1.22);
  \coordinate (D) at (2.14,1.22);
  \draw[given] (A) -- (C);
  \draw[given] (A) -- (E) -- (B);
  \draw[given] (B) -- (D) -- (C);
  \node[vlab, left]  at (A) {$A$};
  \node[vlab, below] at (B) {$B$};
  \node[vlab, right] at (C) {$C$};
  \node[vlab, above] at (E) {$E$};
  \node[vlab, above] at (D) {$D$};
\end{tikzpicture}
\figcap{6--51}
\end{minipage}
\end{bkfigure}}

% ==================================================== Figs 6-52, 6-53, 6-54
% Book p.96.
%  6-52: quadrilateral D A B C with the crossing segments DB and AC.
%  6-53: triangle A B C with E, D on the base and cevians from B.
%  6-54: the farmer/barn word problem -- rectangle "Barn" with the dashed
%        sight-line A..B above and the crossed lines A--E, B--D through C.
\newcommand{\FIGVIFIFTYTWOFIFTYTHREEFIFTYFOUR}{\begin{bkfigure}[\linewidth]
\begin{minipage}{0.46\linewidth}\centering
\begin{tikzpicture}[scale=0.923]
  % Exercise 9 gives angle DAB = angle CBA and AD = BC, so the quadrilateral
  % is an ISOSCELES trapezoid; the draft's D and C were 0.42 and 0.50 in from
  % the ends, which broke both hypotheses.
  \coordinate (A) at (0,0);
  \coordinate (B) at (3.02,0);
  \coordinate (D) at (0.46,1.22);
  \coordinate (C) at (2.56,1.22);
  \draw[given] (A) -- (B);
  \draw[given] (D) -- (C);
  \draw[given] (D) -- (A);
  \draw[given] (C) -- (B);
  \draw[given] (D) -- (B);
  \draw[given] (A) -- (C);
  \node[vlab, below left]  at (A) {$A$};
  \node[vlab, below right] at (B) {$B$};
  \node[vlab, above left]  at (D) {$D$};
  \node[vlab, above right] at (C) {$C$};
\end{tikzpicture}
\figcap{6--52}

\medskip
\begin{tikzpicture}[scale=0.923]
  % Exercise 10 gives AD = EC and BE = BD, the same pair of hypotheses as
  % Fig. 6-38: E and D symmetric about the midpoint, B over that midpoint.
  \coordinate (A) at (0,0);
  \coordinate (C) at (3.02,0);
  \coordinate (B) at (1.51,1.62);
  \coordinate (E) at ($(A)!0.42!(C)$);
  \coordinate (D) at ($(A)!0.58!(C)$);
  \draw[given] (A) -- (C) -- (B) -- cycle;
  \draw[given] (B) -- (E);
  \draw[given] (B) -- (D);
  \node[vlab, left]  at (A) {$A$};
  \node[vlab, right] at (C) {$C$};
  \node[vlab, above] at (B) {$B$};
  \node[vlab, below] at (E) {$E$};
  \node[vlab, below] at (D) {$D$};
\end{tikzpicture}
\figcap{6--53}
\end{minipage}\hfill
\begin{minipage}{0.50\linewidth}\centering
\begin{tikzpicture}[scale=0.923]
  \coordinate (A) at (0,1.72);
  \coordinate (B) at (3.42,1.72);
  \coordinate (D) at (0.42,0);
  \coordinate (E) at (3.02,0);
  \draw[given] (A) -- (E);
  \draw[given] (B) -- (D);
  \coordinate (C) at (intersection of A--E and B--D);
  % the book's barn is drawn round the word with clear space; at the old
  % height the label was as tall as the box and touched both long sides.
  \draw[given] (1.02,1.42) rectangle (2.42,2.02);
  \node[slab] at (1.72,1.72) {Barn};
  \draw[aux] (A) -- (1.02,1.72);
  \draw[aux] (2.42,1.72) -- (B);
  \node[vlab, left]  at (A) {$A$};
  \node[vlab, right] at (B) {$B$};
  \node[vlab, left]  at (D) {$D$};
  \node[vlab, right] at (E) {$E$};
  \node[vlab] at ($(C)+(0:0.66)$) {$C$};
\end{tikzpicture}
\figcap{6--54}
\end{minipage}
\end{bkfigure}}

% ================================================================= Fig 6-55
% Book p.97.  The A.S.A. proof picture -- triangles ABC and A'B'C' with
% the constructed D' on ray A'B' (dashed), arcs at A/A' and C/C', tick on
% AC/A'C'.
% Fig 6-56 used to be bundled into this block and emitted beside Theorem 6-5.
% It belongs to Exercise Group 6-9 #1, further down p.97, and is now its own
% macro so it can be placed there.
\newcommand{\FIGVIFIFTYFIVE}{\begin{bkfigure}[\linewidth]
\begin{tikzpicture}[scale=1.196]
  \coordinate (A) at (0,0);
  \coordinate (C) at (2.62,0);
  \coordinate (B) at (0.62,1.42);
  \draw[given] (A) -- (C) -- (B) -- cycle;
  \VItk{A}{C}{1}
  \VIarc{C}{A}{B}{0.46cm}
  \VIarcb{B}{C}{A}{0.46cm}
  \node[vlab, below left]  at (A) {$A$};
  \node[vlab, above] at (B) {$B$};
  % C and A' face each other across the gap between the two triangles.  Set
  % above the base line and with only 1.34cm of canvas between the panels they
  % closed up and printed as one word, "CA'".  Both base vertices now take the
  % same below-the-line placement A and C' already use, and the panels are
  % 4.10cm apart, which is the daylight the book leaves.
  \node[vlab, below right] at (C) {$C$};
  \begin{scope}[xshift=4.10cm]
   \coordinate (Ap) at (0,0);
   \coordinate (Cp) at (2.62,0);
   \coordinate (Bp) at (0.62,1.42);
   \coordinate (Dp) at ($(Ap)!0.72!(Bp)$);
   \draw[given] (Ap) -- (Cp) -- (Bp) -- cycle;
   \draw[aux]   (Dp) -- (Cp);
   \VItk{Ap}{Cp}{1}
   \VIarc{Cp}{Ap}{Bp}{0.46cm}
   \VIarcb{Bp}{Cp}{Ap}{0.46cm}
   \node[vlab, below left]  at (Ap) {$A\VIpr$};
   \node[vlab, above] at (Bp) {$B\VIpr$};
   \node[vlab, below right] at (Cp) {$C\VIpr$};
   \node[vlab, left] at (Dp) {$D\VIpr$};
  \end{scope}
\end{tikzpicture}
\figcap{6--55}
\end{bkfigure}}

% ================================================================= Fig 6-56
% Book p.97, Exercise Group 6-9 #1.  The kite A B C D with the spine B--D:
% angles 1 and 2 under B, angles 3 and 4 over D.  Checked against scan07 p6
% (book p.97): the book strikes an ARC across each of the four angles, which
% the earlier draft left off, and every numeral sits inside its own arc on
% that angle's bisector.
\newcommand{\FIGVIFIFTYSIX}{\begin{bkfigure}[\linewidth]
\begin{tikzpicture}[scale=1.118]
  \coordinate (B) at (1.32,1.62);
  \coordinate (D) at (1.32,0);
  \coordinate (A) at (0,1.02);
  \coordinate (C) at (2.64,1.02);
  \draw[given] (A) -- (B) -- (C);
  \draw[given] (A) -- (D) -- (C);
  \draw[given] (B) -- (D);
  \node[vlab, above] at (B) {$B$};
  \node[vlab, below] at (D) {$D$};
  \node[vlab, left]  at (A) {$A$};
  \node[vlab, right] at (C) {$C$};
  \VIarc{A}{B}{D}{0.34cm}   \VInumb{A}{B}{D}{0.54cm}{1}
  \VIarc{D}{B}{C}{0.34cm}   \VInumb{D}{B}{C}{0.54cm}{2}
  \VIarc{A}{D}{B}{0.34cm}   \VInumb{A}{D}{B}{0.62cm}{3}
  \VIarc{B}{D}{C}{0.34cm}   \VInumb{B}{D}{C}{0.62cm}{4}
\end{tikzpicture}
\figcap{6--56}
\end{bkfigure}}

% =========================================================== Figs 6-57, 6-58
% Book p.98.  6-57: parallelogram A D C B with diagonal B--D, angles 1..4.
%             6-58: triangle A C with apex B, E and D on the base, angles 1, 2.
\newcommand{\FIGVIFIFTYSEVENFIFTYEIGHT}{\begin{bkfigure}[\linewidth]
\begin{minipage}{0.50\linewidth}\centering
% Rebuilt against scan07 p8 (book p.98).  The book draws ONE diagonal, B--D,
% and all four numbered angles are the parts it cuts off: 3 = ABD and
% 1 = DBC at B, 2 = ADB and 4 = BDC at D, each carrying its own arc.  The
% earlier draft drew a second diagonal A--C that is not in the plate, put 3 at
% vertex A, and stacked 2 and 4 on the SAME bisector at two radii -- which is
% why the collision audit read them as one word "42".
\begin{tikzpicture}[scale=1.289]
  \coordinate (B) at (0,0);
  \coordinate (C) at (2.92,0);
  \coordinate (A) at (0.42,1.22);
  \coordinate (D) at (3.34,1.22);
  \draw[given] (B) -- (C) -- (D) -- (A) -- cycle;
  \draw[given] (B) -- (D);
  \node[vlab, below left]  at (B) {$B$};
  \node[vlab, below right] at (C) {$C$};
  \node[vlab, above left]  at (A) {$A$};
  \node[vlab, above right] at (D) {$D$};
  \VIarc{A}{B}{D}{0.52cm}   \VInumb{A}{B}{D}{0.86cm}{3}
  \VIarc{D}{B}{C}{0.34cm}   \VInumb{D}{B}{C}{1.05cm}{1}
  \VIarc{A}{D}{B}{0.34cm}   \VInumb{A}{D}{B}{1.05cm}{2}
  \VIarc{B}{D}{C}{0.52cm}   \VInumb{B}{D}{C}{0.86cm}{4}
\end{tikzpicture}
\figcap{6--57}
\end{minipage}\hfill
\begin{minipage}{0.46\linewidth}\centering
\begin{tikzpicture}[scale=0.881]
  \coordinate (A) at (0,0);
  \coordinate (C) at (3.02,0);
  \coordinate (B) at (1.32,1.72);
  \coordinate (E) at ($(A)!0.30!(C)$);
  \coordinate (D) at ($(A)!0.62!(C)$);
  \draw[given] (A) -- (C) -- (B) -- cycle;
  \draw[given] (B) -- (E);
  \draw[given] (B) -- (D);
  \node[vlab, left]  at (A) {$A$};
  \node[vlab, right] at (C) {$C$};
  \node[vlab, above] at (B) {$B$};
  \node[vlab, below] at (E) {$E$};
  \node[vlab, below] at (D) {$D$};
  \VInumb{C}{B}{A}{0.88cm}{1}
  \VInumb{D}{B}{C}{1.10cm}{2}
\end{tikzpicture}
\figcap{6--58}
\end{minipage}
\end{bkfigure}}

% ==================================================== Figs 6-59, 6-60
% Book p.98.  6-59: the "dart" A B C E D with the spine through A and B.
%             6-60: triangle A B C with E, D on the sides and crossing cevians.
\newcommand{\FIGVIFIFTYNINESIXTY}{\begin{bkfigure}[\linewidth]
\begin{minipage}{0.48\linewidth}\centering
\begin{tikzpicture}[scale=1.273]
  \coordinate (A) at (0.32,0.92);
  \coordinate (B) at (1.32,0.92);
  \coordinate (C) at (0.92,1.72);
  \coordinate (E) at (0.92,0.10);
  \coordinate (D) at (2.62,0.92);
  \draw[given] (0,0.92) -- (D);
  \draw[given] (C) -- (B) -- (E);
  \draw[given] (C) -- (D) -- (E);
  \node[vlab, above left]  at (A) {$A$};
  % B's letter used to sit up-and-left, straight along C--B; up-and-right
  % runs into C--D, so it drops into the band between the spine and E--D.
  \node[vlab] at ($(B)+(300:0.32)$) {$B$};
  \node[vlab, above] at (C) {$C$};
  \node[vlab, below] at (E) {$E$};
  \node[vlab, right] at (D) {$D$};
\end{tikzpicture}
\figcap{6--59}
\end{minipage}\hfill
\begin{minipage}{0.48\linewidth}\centering
\begin{tikzpicture}[scale=0.926]
  \coordinate (A) at (0,0);
  \coordinate (C) at (3.12,0);
  \coordinate (B) at (1.42,1.62);
  \coordinate (E) at ($(A)!0.52!(B)$);
  \coordinate (D) at ($(C)!0.52!(B)$);
  \draw[given] (A) -- (C) -- (B) -- cycle;
  \draw[given] (A) -- (D);
  \draw[given] (C) -- (E);
  \node[vlab, left]  at (A) {$A$};
  \node[vlab, right] at (C) {$C$};
  \node[vlab, above] at (B) {$B$};
  \node[vlab, above left]  at (E) {$E$};
  \node[vlab, above right] at (D) {$D$};
\end{tikzpicture}
\figcap{6--60}
\end{minipage}
\end{bkfigure}}

% ==================================================== Figs 6-61, 6-62, 6-63
% Book p.98.
%  6-61: segments A--D and E--C crossing at B.
%  6-62: the "kite" A F C D E -- F and C on the spine A--D, E below.
%  6-63: parallelogram-ish B C A D with the base extended past D, angles 3, 4.
\newcommand{\FIGVISIXTYONESIXTYTWOSIXTYTHREE}{\begin{bkfigure}[\linewidth]
\begin{minipage}{0.46\linewidth}\centering
\begin{tikzpicture}[scale=0.914]
  \coordinate (A) at (0,1.22);
  \coordinate (D) at (3.02,0);
  \coordinate (E) at (0,0);
  \coordinate (C) at (3.02,1.22);
  \draw[given] (A) -- (D);
  \draw[given] (E) -- (C);
  \coordinate (B) at (intersection of A--D and E--C);
  \node[vlab, left]  at (A) {$A$};
  \node[vlab, right] at (D) {$D$};
  \node[vlab, left]  at (E) {$E$};
  \node[vlab, right] at (C) {$C$};
  \node[vlab, above] at (B) {$B$};
\end{tikzpicture}
\figcap{6--61}

\medskip
\begin{tikzpicture}[scale=1.256]
  \coordinate (A) at (0,0.72);
  \coordinate (D) at (3.22,0.72);
  \coordinate (B) at (1.12,1.62);
  \coordinate (E) at (1.72,0);
  \coordinate (F) at ($(A)!0.34!(D)$);
  \coordinate (C) at ($(A)!0.66!(D)$);
  \draw[given] (A) -- (D);
  \draw[given] (A) -- (B) -- (C);
  \draw[given] (F) -- (E) -- (D);
  \draw[given] (B) -- (D);
  \node[vlab, left]  at (A) {$A$};
  \node[vlab, right] at (D) {$D$};
  \node[vlab, above] at (B) {$B$};
  \node[vlab, below] at (E) {$E$};
  \node[vlab, below left]  at (F) {$F$};
  % C is the foot of B--C on the spine; below-LEFT lay the letter along
  % that segment, so it goes straight down into the band above E--D.
  \node[vlab, below] at (C) {$C$};
\end{tikzpicture}
\figcap{6--62}
\end{minipage}\hfill
\begin{minipage}{0.48\linewidth}\centering
\begin{tikzpicture}[scale=0.914]
  % Exercise 10 gives AB = AD and "AC bisects angle BAD", so B and D are
  % equidistant from A and C stands on the bisector of that angle.  The
  % draft's AB was 1.33 against AD 2.12 and its AC missed the bisector
  % altogether -- the plate showed neither hypothesis.
  \coordinate (A) at (0,0);
  \coordinate (D) at ($(A)+(0:1.55)$);
  \coordinate (B) at ($(A)+(70:1.55)$);
  \coordinate (C) at ($(A)+(35:2.30)$);
  \coordinate (X) at (2.30,0);          % the base carried on past D
  \draw[given] (A) -- (X);
  \draw[given] (A) -- (B) -- (C);
  \draw[given] (A) -- (C);
  \draw[given] (C) -- (D);
  \node[vlab, left]  at (A) {$A$};
  \node[vlab, below] at (D) {$D$};
  \node[vlab, above left]  at (B) {$B$};
  \node[vlab, above right] at (C) {$C$};
  % 4 is the angle between DC and the STUB of the base beyond D -- the book
  % extends the base for exactly that purpose.  The draft aimed the numeral
  % along DC itself.
  \VIarc{A}{B}{C}{0.34cm}   \VInumb{A}{B}{C}{0.58cm}{3}
  \VIarc{C}{D}{X}{0.30cm}   \VInumb{C}{D}{X}{0.54cm}{4}
\end{tikzpicture}
\figcap{6--63}
\end{minipage}
\end{bkfigure}}

% =========================================================== Figs 6-64, 6-65
% Book p.99.  6-64: triangle A B with apex C, D and E on the sides, crossing
%   cevians AE and BD meeting at F, angles 3 and 4 at D and E.
%  6-65: triangle A C with apex B and D below B, angles 1 and 2.
\newcommand{\FIGVISIXTYFOURSIXTYFIVE}{\begin{bkfigure}[\linewidth]
\begin{minipage}{0.50\linewidth}\centering
\begin{tikzpicture}[scale=0.935]
  \coordinate (A) at (0,0);
  \coordinate (B) at (3.22,0);
  \coordinate (C) at (1.61,1.72);
  \coordinate (D) at ($(A)!0.62!(C)$);
  \coordinate (E) at ($(B)!0.62!(C)$);
  \draw[given] (A) -- (B) -- (C) -- cycle;
  \draw[given] (A) -- (E);
  \draw[given] (B) -- (D);
  \coordinate (F) at (intersection of A--E and B--D);
  \node[vlab, left]  at (A) {$A$};
  \node[vlab, right] at (B) {$B$};
  \node[vlab, above] at (C) {$C$};
  \node[vlab, above left]  at (D) {$D$};
  \node[vlab, above right] at (E) {$E$};
  \node[vlab, below] at (F) {$F$};
  % 3 is angle CDB at D and 4 is angle CEA at E, each on its own bisector and
  % each carrying an arc.  The draft aimed them at fixed bearings that fell
  % straight along DB and EA.
  % The two bisectors aim at one another, so 0.58 drove the numerals together
  % until they printed as "34" in a single heap; 0.44 gives them the book's gap
  % and still clears the 0.26 arcs.
  \VIarc{C}{D}{B}{0.26cm}   \VInumb{C}{D}{B}{0.44cm}{3}
  \VIarc{C}{E}{A}{0.26cm}   \VInumb{C}{E}{A}{0.44cm}{4}
\end{tikzpicture}
\figcap{6--64}
\end{minipage}\hfill
\begin{minipage}{0.46\linewidth}\centering
\begin{tikzpicture}[scale=0.935]
  % Exercise 12 gives AB = BC and AD = DC, so B and D BOTH lie on the
  % perpendicular bisector of AC, which is how the book draws them -- D
  % directly under B, half its height.  The draft put B at 0.35 of the base
  % and D off the axis, so neither hypothesis held.
  \coordinate (A) at (0,0);
  \coordinate (C) at (2.92,0);
  \coordinate (B) at (1.46,1.82);
  \coordinate (D) at (1.46,0.91);
  \draw[given] (A) -- (C) -- (B) -- cycle;
  \draw[given] (A) -- (D) -- (C);
  \draw[given] (B) -- (D);
  \node[vlab, left]  at (A) {$A$};
  \node[vlab, right] at (C) {$C$};
  \node[vlab, above] at (B) {$B$};
  \node[vlab, below] at (D) {$D$};
  \VIarc{A}{B}{D}{0.40cm}   \VInumb{A}{B}{D}{0.72cm}{1}
  \VIarc{D}{B}{C}{0.40cm}   \VInumb{D}{B}{C}{0.72cm}{2}
\end{tikzpicture}
\figcap{6--65}
\end{minipage}
\end{bkfigure}}

% =========================================================== Figs 6-66, 6-67
% Book p.99.  Both are the "cevian web" figures.
%  6-66: triangle A B with apex C; D, E on the sides, H on AB; cevians AE, BD
%        and CH concurrent at G; angles 3, 4 at D, E.
%  6-67: triangle A B with apex C; D, E on the sides; F, G on AB; cevians
%        crossing at H, I, J; angles 1, 2 at the base.
\newcommand{\FIGVISIXTYSIXSIXTYSEVEN}{\begin{bkfigure}[\linewidth]
\begin{minipage}{0.48\linewidth}\centering
% Re-measured on the plate (book p.99, photo PDF p.113) at 4x: base 700 px,
% apex 610 px above it -- height/base 0.871, NOT the 0.565 the draft drew.
% Flattening the triangle is what squeezed the band between D--E and the
% cevian crossing G, and that is the band the letters F and G live in: F was
% pushed 13.7% of the base out from its point to find room and landed on A--E.
% D and E measured at 0.60 of the way up their sides (draft had 0.66).
\begin{tikzpicture}[scale=1.50]
  \coordinate (A) at (0,0);
  \coordinate (B) at (3.22,0);
  \coordinate (C) at (1.61,2.80);
  \coordinate (H) at ($(A)!0.5!(B)$);
  \coordinate (D) at ($(A)!0.60!(C)$);
  \coordinate (E) at ($(B)!0.60!(C)$);
  \draw[given] (A) -- (B) -- (C) -- cycle;
  \draw[given] (A) -- (E);
  \draw[given] (B) -- (D);
  \draw[given] (C) -- (H);
  \coordinate (G) at (intersection of A--E and B--D);
  \draw[given] (D) -- (E);
  % F is where C--H crosses D--E; the book letters it and the draft left it
  % out altogether.
  \coordinate (F) at (intersection of C--H and D--E);
  \node[vlab, left]  at (A) {$A$};
  \node[vlab, right] at (B) {$B$};
  \node[vlab, above] at (C) {$C$};
  \node[vlab, below] at (H) {$H$};
  \node[vlab, above left]  at (D) {$D$};
  \node[vlab, above right] at (E) {$E$};
  % Offsets measured off the plate, as a fraction of the base: F sits about
  % 5% out from its point and G about 9%, both down-and-left.  The draft used
  % 13.7% and 14.3%, which is what carried F across the cevian A--E.
  \node[vlab] at ($(F)+(225:0.227)$) {$F$};
  \node[vlab] at ($(G)+(247:0.33)$) {$G$};
  % 3 and 4 are the two halves of the vertex angle at C, and the book sets
  % them at the SAME radius on either side of C--H.  The draft used 1.30 and
  % 0.34, which threw both numerals onto the sides of the triangle.
  \VIarc{A}{C}{H}{0.30cm}   \VInumb{A}{C}{H}{0.44cm}{3}
  \VIarc{H}{C}{B}{0.30cm}   \VInumb{H}{C}{B}{0.44cm}{4}
\end{tikzpicture}
\figcap{6--66}
\end{minipage}\hfill
\begin{minipage}{0.48\linewidth}\centering
% Rebuilt against the photo of book p.99 (PDF 113).  F and G are on the SIDES
% CA and CB (below D and E), not on the base: the segment F--G runs across the
% triangle, and the two lines drawn are the cevians A--E and B--D.  H is where
% those cevians cross; I and J are where they cut F--G.  The earlier draft put
% F and G on AB, drew D--G and E--F instead, and then placed I and J at
% arbitrary ratios along F--G rather than deriving them, so the exercise's
% triangle HIJ was not the one in the picture.
% Exercise 14 gives AC = BC and DF = EG, and the plate is mirror-symmetric
% about the altitude from C -- which is what makes triangle HIJ isosceles.
\begin{tikzpicture}[scale=1.30]
  \coordinate (A) at (0,0);
  \coordinate (B) at (3.32,0);
  \coordinate (C) at (1.66,2.70);
  \coordinate (D) at ($(A)!0.55!(C)$);
  \coordinate (E) at ($(B)!0.55!(C)$);
  \coordinate (F) at ($(A)!0.30!(C)$);
  \coordinate (G) at ($(B)!0.30!(C)$);
  \draw[given] (A) -- (B) -- (C) -- cycle;
  \draw[given] (F) -- (G);
  \draw[given] (A) -- (E);
  \draw[given] (B) -- (D);
  \coordinate (H) at (intersection of A--E and B--D);
  \coordinate (I) at (intersection of A--E and F--G);
  \coordinate (J) at (intersection of B--D and F--G);
  \node[vlab, left]  at (A) {$A$};
  \node[vlab, right] at (B) {$B$};
  \node[vlab, above] at (C) {$C$};
  \node[vlab, above left]  at (D) {$D$};
  \node[vlab, above right] at (E) {$E$};
  \node[vlab, left]  at (F) {$F$};
  \node[vlab, right] at (G) {$G$};
  \node[vlab, above] at (H) {$H$};
  \node[vlab] at ($(I)+(285:0.32)$) {$I$};
  \node[vlab] at ($(J)+(255:0.32)$) {$J$};
  % 1 is the angle at A between the side AC and the cevian AE; 2 its mirror
  % image at B.
  \VIarc{C}{A}{E}{0.40cm}   \VInumb{C}{A}{E}{0.72cm}{1}
  \VIarc{C}{B}{D}{0.40cm}   \VInumb{C}{B}{D}{0.72cm}{2}
\end{tikzpicture}
\figcap{6--67}
\end{minipage}
\end{bkfigure}}

% ================================================================= Fig 6-68
% Book p.99.  Theorem 6-6: the same triangle labelled two ways, as in 6-32.
\newcommand{\FIGVISIXTYEIGHT}{\begin{bkfigure}[\linewidth]
\begin{tikzpicture}[scale=1.196]
  \coordinate (B) at (1.32,1.52);
  \coordinate (A) at (0,0);
  \coordinate (C) at (2.64,0);
  \draw[given] (A) -- (B) -- (C) -- cycle;
  \node[vlab, above] at (B) {$B = B\VIpr$};
  \node[vlab, below left]  at (A) {$A = C\VIpr$};
  \node[vlab, below right] at (C) {$C = A\VIpr$};
\end{tikzpicture}
\figcap{6--68}
\end{bkfigure}}

% ================================================================= Fig 6-69
% Book p.100.  The S.S.S. proof: (a) triangle ABC, (b) triangle A'B'C' with the
% arc at A', (c) the reflected copy with E and D below AC (ticks on AB, AE).
\newcommand{\FIGVISIXTYNINE}{\begin{bkfigure}[\linewidth]
\begin{tikzpicture}[scale=1.118]
  % (a)
  \coordinate (A) at (0,0);
  \coordinate (C) at (1.92,0);
  \coordinate (B) at (0.72,1.12);
  \draw[given] (A) -- (C) -- (B) -- cycle;
  \node[vlab, left]  at (A) {$A$};
  \node[vlab, below left] at (C) {$C$};
  \node[vlab, above] at (B) {$B$};
  \node[slab, below] at (0.96,-0.28) {(a)};
  % (b)
  \begin{scope}[xshift=2.72cm]
   \coordinate (Ap) at (0,0);
   \coordinate (Cp) at (1.92,0);
   \coordinate (Bp) at (0.72,1.12);
   \draw[given] (Ap) -- (Cp) -- (Bp) -- cycle;
   \VIarc{Cp}{Ap}{Bp}{0.42cm}
   \node[vlab, below]  at (Ap) {$A\VIpr$};
   % Set flush right, C' stacked under panel (c)'s A with 0.35pt between them.
   % Below-right matches the placement A' already takes at the other end of the
   % same base and drops it clear of the next panel.
   \node[vlab, below right] at (Cp) {$C\VIpr$};
   \node[vlab, above] at (Bp) {$B\VIpr$};
   \node[slab, below] at (0.96,-0.28) {(b)};
  \end{scope}
  % (c)
  \begin{scope}[xshift=5.44cm]
   \coordinate (Ac) at (0,0.52);
   \coordinate (Cc) at (1.72,0.52);
   \coordinate (Bc) at (1.22,1.52);
   \coordinate (Ec) at (1.22,-0.48);
   \coordinate (Dc) at (1.46,-0.62);
   \draw[given] (Ac) -- (Cc);
   \draw[given] (Ac) -- (Bc) -- (Cc);
   \draw[given] (Ac) -- (Ec) -- (Cc);
   \draw[given] (Ac) -- (Dc);
   \VItk{Ac}{Bc}{1}
   \VItk{Ac}{Ec}{1}
   \VIarc{Cc}{Ac}{Ec}{0.40cm}
   \node[vlab, left]  at (Ac) {$A$};
   \node[vlab, right] at (Cc) {$C$};
   \node[vlab, above] at (Bc) {$B$};
   \node[vlab, below]  at (Ec) {$E$};
   \node[vlab, right] at (Dc) {$D$};
   \node[slab, below] at (0.62,-0.92) {(c)};
  \end{scope}
\end{tikzpicture}
\figcap{6--69}
\end{bkfigure}}

% ================================================================= Fig 6-70
% Book p.100.  The kite A B C E with both diagonals, angles 1..4 at B and E.
\newcommand{\FIGVISEVENTY}{\begin{bkfigure}[\linewidth]
\begin{tikzpicture}[scale=1.196]
  \coordinate (B) at (1.32,1.72);
  \coordinate (A) at (0,0.86);
  \coordinate (C) at (2.64,0.86);
  \coordinate (E) at (1.32,-0.32);
  \draw[given] (A) -- (B) -- (C) -- (E) -- cycle;
  \draw[given] (A) -- (C);
  \draw[given] (B) -- (E);
  \node[vlab, above] at (B) {$B$};
  \node[vlab, left]  at (A) {$A$};
  \node[vlab, right] at (C) {$C$};
  \node[vlab, below] at (E) {$E$};
  \VInum{B}{-118}{0.42cm}{1}
  \VInum{B}{-62}{0.42cm}{2}
  \VInumb{B}{E}{A}{0.64cm}{3}
  \VInumb{C}{E}{B}{0.64cm}{4}
\end{tikzpicture}
\figcap{6--70}
\end{bkfigure}}

% ================================================================= Fig 6-71
% Book p.101.  The two cases the S.S.S. proof tacitly assumed away: BE meets AC
% inside the segment (a) or outside it (b).
\newcommand{\FIGVISEVENTYONE}{\begin{bkfigure}[\linewidth]
\begin{tikzpicture}[scale=1.144]
  % (a) C between A and the crossing
  \coordinate (A) at (0,0.62);
  \coordinate (B) at (1.42,1.52);
  \coordinate (E) at (1.42,-0.42);
  \coordinate (C) at (0.92,0.62);
  \draw[given] (A) -- (B) -- (E) -- cycle;
  \draw[given] (A) -- (C);
  \draw[given] (B) -- (E);
  \node[vlab, left]  at (A) {$A$};
  \node[vlab, above] at (B) {$B$};
  \node[vlab, below] at (E) {$E$};
  \node[vlab, right] at (C) {$C$};
  \node[slab, below] at (0.86,-0.86) {(a)};
  \begin{scope}[xshift=2.72cm]
   \coordinate (A2) at (0,0.62);
   \coordinate (B2) at (1.42,1.52);
   \coordinate (E2) at (1.42,-0.42);
   % C2 sat close enough to the vertical B2--E2 that its letter, set to the
   % right, ran into that segment.
   \coordinate (C2) at (0.88,0.62);
   \draw[given] (A2) -- (B2) -- (E2) -- cycle;
   \draw[given] (A2) -- (C2);
   \draw[given] (B2) -- (C2) -- (E2);
   \node[vlab, left]  at (A2) {$A$};
   \node[vlab, above] at (B2) {$B$};
   \node[vlab, below] at (E2) {$E$};
   \node[vlab, right] at (C2) {$C$};
   \node[slab, below] at (0.86,-0.86) {(b)};
  \end{scope}
\end{tikzpicture}
\figcap{6--71}
\end{bkfigure}}

% ==================================================== Figs 6-72, 6-73, 6-74
% Book p.101.  6-72: the braced gate (a picture, not a lettered diagram) --
%   post at the left, rectangular gate, diagonal brace, ground hatching.
%  6-73: kite A B C D with the diagonal B--D.
%  6-74: parallelogram A B C D with the diagonal A--C.
\newcommand{\FIGVISEVENTYTWOSEVENTYTHREESEVENTYFOUR}{\begin{bkfigure}[\linewidth]
\begin{minipage}{0.48\linewidth}\centering
\begin{tikzpicture}[scale=0.880]
  % post
  \draw[given] (0.10,-0.12) rectangle (0.42,1.92);
  % gate
  \draw[given] (0.52,0.12) rectangle (3.02,1.62);
  \draw[given] (0.66,0.26) rectangle (2.88,1.48);
  % brace
  \draw[given] (0.66,0.26) -- (2.88,1.48);
  % hinges
  \draw[given] (0.42,1.42) -- (0.66,1.42);
  \draw[given] (0.42,0.36) -- (0.66,0.36);
  % ground
  \draw[given] (-0.10,-0.12) -- (3.32,-0.12);
  \foreach \x in {-0.02,0.16,...,3.20}{
    \draw[tick] (\x,-0.12) -- ($(\x,-0.12)+(-0.10,-0.14)$);}
\end{tikzpicture}
\figcap{6--72}
\end{minipage}\hfill
\begin{minipage}{0.22\linewidth}\centering
\begin{tikzpicture}[scale=0.880]
  \coordinate (B) at (0.92,1.72);
  \coordinate (A) at (0,1.02);
  \coordinate (C) at (1.84,1.02);
  \coordinate (D) at (0.92,0);
  \draw[given] (A) -- (B) -- (C) -- (D) -- cycle;
  \draw[given] (B) -- (D);
  \node[vlab, above] at (B) {$B$};
  \node[vlab, left]  at (A) {$A$};
  \node[vlab, right] at (C) {$C$};
  \node[vlab, below] at (D) {$D$};
\end{tikzpicture}
\figcap{6--73}
\end{minipage}\hfill
\begin{minipage}{0.26\linewidth}\centering
\begin{tikzpicture}[scale=0.880]
  \coordinate (D) at (0,0);
  \coordinate (C) at (2.12,0.22);
  \coordinate (B) at (1.82,1.12);
  \coordinate (A) at (0.32,1.02);
  \draw[given] (D) -- (C) -- (B) -- (A) -- cycle;
  \draw[given] (A) -- (C);
  \node[vlab, below left]  at (D) {$D$};
  \node[vlab, right] at (C) {$C$};
  \node[vlab, above right] at (B) {$B$};
  \node[vlab, above left]  at (A) {$A$};
\end{tikzpicture}
\figcap{6--74}
\end{minipage}
\end{bkfigure}}

% =========================================================== Figs 6-75, 6-76
% Book p.102.  6-75: triangle A C with apex B and cevian BD.
%   6-76: the kite A B C D with spine A--C through E, angles 1 and 2 at A.
\newcommand{\FIGVISEVENTYFIVESEVENTYSIX}{\begin{bkfigure}[\linewidth]
\begin{minipage}{0.44\linewidth}\centering
\begin{tikzpicture}[scale=0.938]
  % Exercise 4 gives AB = BC and AD = DC, so D is the midpoint of AC and B
  % stands over it, as the book draws it.
  \coordinate (A) at (0,0);
  \coordinate (C) at (2.62,0);
  \coordinate (B) at (1.31,1.72);
  \coordinate (D) at ($(A)!0.5!(C)$);
  \draw[given] (A) -- (C) -- (B) -- cycle;
  \draw[given] (B) -- (D);
  \node[vlab, left]  at (A) {$A$};
  \node[vlab, right] at (C) {$C$};
  \node[vlab, above] at (B) {$B$};
  \node[vlab, below] at (D) {$D$};
\end{tikzpicture}
\figcap{6--75}
\end{minipage}\hfill
\begin{minipage}{0.52\linewidth}\centering
\begin{tikzpicture}[scale=0.938]
  % Exercise 5 gives AB = AD and BC = CD, so A and C both lie on the
  % perpendicular bisector of BD: B and D must straddle the line AC evenly.
  % They did not -- BD's midpoint sat 0.05 below AC -- and E was placed by a
  % hand-picked ratio instead of being taken as the actual crossing.
  \coordinate (A) at (0,0.86);
  \coordinate (C) at (3.42,0.86);
  \coordinate (B) at (1.22,1.72);
  \coordinate (D) at (1.22,0);
  \draw[given] (A) -- (B) -- (C) -- (D) -- cycle;
  \draw[given] (A) -- (C);
  \draw[given] (B) -- (D);
  \coordinate (E) at (intersection of A--C and B--D);
  \node[vlab, left]  at (A) {$A$};
  \node[vlab, right] at (C) {$C$};
  \node[vlab, above] at (B) {$B$};
  \node[vlab, below] at (D) {$D$};
  \node[vlab, above right] at (E) {$E$};
  \VInumb{C}{A}{B}{0.86cm}{1}
  \VInumb{D}{A}{C}{0.84cm}{2}
\end{tikzpicture}
\figcap{6--76}
\end{minipage}
\end{bkfigure}}

% =========================================================== Figs 6-77, 6-78
% Book p.102.  6-77: parallelogram A B C D with E on BC-side, F interior, and
%   the segments A--C, B--D drawn.  6-78: triangle A C with apex B, E and D on
%   the base.
\newcommand{\FIGVISEVENTYSEVENSEVENTYEIGHT}{\begin{bkfigure}[\linewidth]
\begin{minipage}{0.52\linewidth}\centering
\begin{tikzpicture}[scale=1.179]
  \coordinate (A) at (0,0);
  \coordinate (D) at (2.62,0);
  \coordinate (B) at (0.52,1.12);
  \coordinate (C) at (3.14,1.12);
  \draw[given] (A) -- (D) -- (C) -- (B) -- cycle;
  \draw[given] (B) -- (D);
  \draw[given] (A) -- (C);
  % AF = EC is given, so F and E sit symmetrically on the diagonal.
  \coordinate (F) at ($(A)!0.30!(C)$);
  \coordinate (E) at ($(A)!0.70!(C)$);
  \node[vlab, below left]  at (A) {$A$};
  \node[vlab, below right] at (D) {$D$};
  \node[vlab, above left]  at (B) {$B$};
  \node[vlab, above right] at (C) {$C$};
  % E and F both sit ON the diagonal, so each letter goes on the side away
  % from the long edge it is nearest.
  \node[vlab, below right] at (E) {$E$};
  \node[vlab, above left] at (F) {$F$};
\end{tikzpicture}
\figcap{6--77}
\end{minipage}\hfill
\begin{minipage}{0.44\linewidth}\centering
\begin{tikzpicture}[scale=0.898]
  % Exercise 7: triangle ABC isosceles with vertex B (so B stands over the
  % midpoint of AC) and AD = EC (so E and D are symmetric about it).
  \coordinate (A) at (0,0);
  \coordinate (C) at (2.92,0);
  \coordinate (B) at (1.46,1.72);
  \coordinate (E) at ($(A)!0.32!(C)$);
  \coordinate (D) at ($(A)!0.68!(C)$);
  \draw[given] (A) -- (C) -- (B) -- cycle;
  \draw[given] (B) -- (E);
  \draw[given] (B) -- (D);
  \node[vlab, left]  at (A) {$A$};
  \node[vlab, right] at (C) {$C$};
  \node[vlab, above] at (B) {$B$};
  \node[vlab, below] at (E) {$E$};
  \node[vlab, below] at (D) {$D$};
\end{tikzpicture}
\figcap{6--78}
\end{minipage}
\end{bkfigure}}

% =========================================================== Figs 6-79, 6-80
% Book p.102.  6-79: the re-entrant "dart" A B C D.
%   6-80: triangle A B with apex C, D on the base, E, G, F on the cevian rows,
%   angles 1 and 2 at D.
\newcommand{\FIGVISEVENTYNINEEIGHTY}{\begin{bkfigure}[\linewidth]
\begin{minipage}{0.48\linewidth}\centering
\begin{tikzpicture}[scale=0.938]
  % Exercise 8 gives AB = AD and BC = CD, so A and C lie on the perpendicular
  % bisector of BD -- both at the height of BD's midpoint, 0.81.
  \coordinate (A) at (0,0.81);
  \coordinate (B) at (3.02,1.62);
  \coordinate (D) at (3.02,0);
  \coordinate (C) at (2.12,0.81);
  \draw[given] (A) -- (B) -- (C) -- (D) -- cycle;
  \node[vlab, left]  at (A) {$A$};
  \node[vlab, right] at (B) {$B$};
  \node[vlab, right] at (D) {$D$};
  \node[vlab, left] at (C) {$C$};
\end{tikzpicture}
\figcap{6--79}
\end{minipage}\hfill
\begin{minipage}{0.48\linewidth}\centering
\begin{tikzpicture}[scale=0.938]
  \coordinate (A) at (0,0);
  \coordinate (B) at (3.02,0);
  \coordinate (C) at (1.51,1.72);
  \coordinate (D) at ($(A)!0.5!(B)$);
  \coordinate (E) at ($(A)!0.54!(C)$);
  \coordinate (F) at ($(B)!0.54!(C)$);
  \draw[given] (A) -- (B) -- (C) -- cycle;
  \draw[given] (C) -- (D);
  \draw[given] (E) -- (F);
  % G is where C--D cuts E--F; it was placed as the midpoint of EF instead of
  % being derived, and its letter, set up-and-right, ran into the side CB.
  \coordinate (G) at (intersection of C--D and E--F);
  \node[vlab, left]  at (A) {$A$};
  \node[vlab, right] at (B) {$B$};
  \node[vlab, above] at (C) {$C$};
  \node[vlab, below] at (D) {$D$};
  \node[vlab, above left]  at (E) {$E$};
  \node[vlab, above right] at (F) {$F$};
  \node[vlab] at ($(G)+(-45:0.42)$) {$G$};
  % The book strikes ONE arc across the whole straight angle ADB and writes
  % the two numerals inside it, either side of DC.
  \draw[fig, line width=0.55pt] ([shift=(0:0.36cm)]D)
       arc[start angle=0, end angle=180, radius=0.36cm];
  \VInum{D}{145}{0.92cm}{2}
  \VInum{D}{35}{0.92cm}{1}
\end{tikzpicture}
\figcap{6--80}
\end{minipage}
\end{bkfigure}}

% =========================================================== Figs 6-81, 6-82
% Book p.102.  6-81: triangle A B with D on AB and apex C below, E and F on the
%   line through C; angles 1 and 2 at C.  6-82: segments C--B and D--A crossing.
\newcommand{\FIGVIEIGHTYONEEIGHTYTWO}{\begin{bkfigure}[\linewidth]
\begin{minipage}{0.50\linewidth}\centering
\begin{tikzpicture}[scale=0.958]
  \coordinate (A) at (0,1.42);
  \coordinate (B) at (2.82,1.42);
  \coordinate (C) at (1.41,0);
  \coordinate (D) at ($(A)!0.5!(B)$);
  \draw[given] (A) -- (B);
  \draw[given] (A) -- (C) -- (B);
  \draw[given] (C) -- (D);
  \draw[given] (0.12,0) -- (2.70,0);
  \node[vlab, left]  at (A) {$A$};
  \node[vlab, right] at (B) {$B$};
  \node[vlab, below] at (C) {$C$};
  \node[vlab, above] at (D) {$D$};
  \node[vlab, left]  at (0.12,0) {$E$};
  \node[vlab, right] at (2.70,0) {$F$};
  % One arc across the whole straight angle ECF, numerals inside it either
  % side of CD.  At 0.46 they sat on the line EF itself.
  \draw[fig, line width=0.55pt] ([shift=(0:0.44cm)]C)
       arc[start angle=0, end angle=180, radius=0.44cm];
  \VInum{C}{112}{0.74cm}{1}
  \VInum{C}{68}{0.74cm}{2}
\end{tikzpicture}
\figcap{6--81}
\end{minipage}\hfill
\begin{minipage}{0.46\linewidth}\centering
\begin{tikzpicture}[scale=0.958]
  % Exercise 11 gives AC = BD and AD = BC, so C and D must sit symmetrically
  % over AB; the draft's D was 0.18 off, breaking both.
  \coordinate (A) at (0,0);
  \coordinate (B) at (2.82,0);
  \coordinate (C) at (0.72,1.32);
  \coordinate (D) at (2.10,1.32);
  \draw[given] (A) -- (B);
  \draw[given] (A) -- (C);
  \draw[given] (A) -- (D);
  \draw[given] (C) -- (B);
  \draw[given] (D) -- (B);
  \node[vlab, left]  at (A) {$A$};
  \node[vlab, right] at (B) {$B$};
  \node[vlab, above] at (C) {$C$};
  \node[vlab, above] at (D) {$D$};
\end{tikzpicture}
\figcap{6--82}
\end{minipage}
\end{bkfigure}}

% =========================================================== Figs 6-83, 6-84
% Book p.103.  6-83: A at the top with B, C at the ends and D, E crossing.
%   6-84: triangle A B with apex C, D on the base, O interior.
\newcommand{\FIGVIEIGHTYTHREEEIGHTYFOUR}{\begin{bkfigure}[\linewidth]
\begin{minipage}{0.46\linewidth}\centering
\begin{tikzpicture}[scale=0.987]
  % Exercise 12 gives AB = BC, so A and C must be mirror images across the
  % horizontal through B; the draft had them 0.80 and 1.02 from it.
  \coordinate (A) at (2.02,1.82);
  \coordinate (B) at (0,1.02);
  \coordinate (C) at (2.02,0.22);
  \coordinate (D) at ($(B)!0.52!(A)$);
  \coordinate (E) at ($(B)!0.52!(C)$);
  \draw[given] (B) -- (A);
  \draw[given] (B) -- (C);
  \draw[given] (A) -- (E);
  \draw[given] (C) -- (D);
  \node[vlab, above] at (A) {$A$};
  \node[vlab, left]  at (B) {$B$};
  \node[vlab, below right] at (C) {$C$};
  \node[vlab, above left]  at (D) {$D$};
  \node[vlab, below left]  at (E) {$E$};
\end{tikzpicture}
\figcap{6--83}
\end{minipage}\hfill
\begin{minipage}{0.50\linewidth}\centering
\begin{tikzpicture}[scale=0.987]
  \coordinate (A) at (0,0);
  \coordinate (B) at (3.02,0);
  \coordinate (C) at (1.51,1.72);
  \coordinate (D) at ($(A)!0.5!(B)$);
  \draw[given] (A) -- (B) -- (C) -- cycle;
  \draw[given] (C) -- (D);
  \coordinate (Pb) at ($(B)!0.42!(C)$);
  \coordinate (Pa) at ($(A)!0.42!(C)$);
  \draw[given] (A) -- (Pb);
  \draw[given] (B) -- (Pa);
  \coordinate (O) at (intersection of C--D and A--Pb);
  \node[vlab, left]  at (A) {$A$};
  \node[vlab, right] at (B) {$B$};
  \node[vlab, above] at (C) {$C$};
  \node[vlab, below] at (D) {$D$};
  % Six rays leave O; the only wide gap is the one opening to the right,
  % between the two cevians.  "above right" lay along A--Pb.
  \node[vlab] at ($(O)+(56:0.62)$) {$O$};
\end{tikzpicture}
\figcap{6--84}
\end{minipage}
\end{bkfigure}}

% =========================================================== Figs 6-85, 6-86
% Book p.103.  6-85: triangle A C with B above A, D and E on the cevians,
%   angles 1 and 2 at A/B.  6-86: parallelogram D C A B with E, F on the sides
%   and M the crossing of the two cevians.
\newcommand{\FIGVIEIGHTYFIVEEIGHTYSIX}{\begin{bkfigure}[\linewidth]
\begin{minipage}{0.48\linewidth}\centering
% Drawn at 1.05, not 0.80.  Both numerals hang off the SAME short side AB --
% 1 in the wedge at A, 2 in the wedge at B -- and each wedge is only about 56
% degrees wide, so the numeral has to sit far enough out to clear both its
% rays while staying clear of the other numeral coming the other way.  At 0.80
% there was no such distance: 1 grazed A--E (0.00 pt).  Enlarging the plate
% buys the room without moving either numeral off its bisector.
\begin{tikzpicture}[scale=1.194]
  % Exercise 14 gives AC = BC, so C must sit on the perpendicular bisector of
  % AB.  The draft's C gave AC = 3.08 against BC = 2.91.
  \coordinate (A) at (0,0);
  \coordinate (B) at (0.22,1.42);
  \coordinate (C) at (3.07,0.25);
  \coordinate (D) at ($(A)!0.52!(C)$);
  \coordinate (E) at ($(B)!0.52!(C)$);
  \draw[given] (A) -- (B);
  \draw[given] (A) -- (C);
  \draw[given] (B) -- (C);
  \draw[given] (A) -- (E);
  \draw[given] (B) -- (D);
  \node[vlab, below] at (A) {$A$};
  \node[vlab, above] at (B) {$B$};
  \node[vlab, right] at (C) {$C$};
  \node[vlab, below] at (D) {$D$};
  \node[vlab, above] at (E) {$E$};
  % 1 sits at A (angle EAB) and 2 at B (angle ABD) -- opposite ends of the same
  % short side, each pushed along a bisector that runs INTO the other.  AB is
  % only 1.44 units, and 0.66 + 0.90 overshot it, so the two numerals crossed
  % and printed as "12".  Measured off the plate, the book puts 1 about 0.31 of
  % AB above A and 2 about 0.41 of AB below B, leaving a clear gap between them.
  \VInumb{E}{A}{B}{0.45cm}{1}
  \VInumb{A}{B}{D}{0.59cm}{2}
\end{tikzpicture}
\figcap{6--85}
\end{minipage}\hfill
\begin{minipage}{0.48\linewidth}\centering
\begin{tikzpicture}[scale=0.909]
  \coordinate (A) at (0,0);
  \coordinate (B) at (2.62,0);
  \coordinate (D) at (0.52,1.22);
  \coordinate (C) at (3.14,1.22);
  % AM = MC is given, so M must be the midpoint of the diagonal -- which
  % happens exactly when E and F are symmetric about the parallelogram's
  % centre, i.e. when their two ratios add to 1.  The draft's 0.30 and 0.62
  % did not, so the drawing contradicted its own hypothesis.
  \coordinate (E) at ($(D)!0.38!(C)$);
  \coordinate (F) at ($(A)!0.62!(B)$);
  \draw[given] (A) -- (B) -- (C) -- (D) -- cycle;
  \draw[given] (A) -- (C);
  \draw[given] (E) -- (F);
  \coordinate (M) at (intersection of A--C and E--F);
  \node[vlab, below left]  at (A) {$A$};
  \node[vlab, below right] at (B) {$B$};
  \node[vlab, above left]  at (D) {$D$};
  \node[vlab, above right] at (C) {$C$};
  \node[vlab, above] at (E) {$E$};
  \node[vlab, below] at (F) {$F$};
  \node[vlab, above left] at (M) {$M$};
\end{tikzpicture}
\figcap{6--86}
\end{minipage}
\end{bkfigure}}

% =========================================================== Figs 6-87, 6-88
% Book p.103.  6-87: triangle A B with apex C, D and E on the sides, crossing
%   cevians AE and BD; angles 1..4.  6-88: the kite A B D C with E below.
\newcommand{\FIGVIEIGHTYSEVENEIGHTYEIGHT}{\begin{bkfigure}[\linewidth]
\begin{minipage}{0.52\linewidth}\centering
\begin{tikzpicture}[scale=1.381]
  \coordinate (A) at (0,0);
  \coordinate (B) at (3.32,0);
  \coordinate (C) at (1.66,1.72);
  \coordinate (D) at ($(A)!0.58!(C)$);
  \coordinate (E) at ($(B)!0.58!(C)$);
  \draw[given] (A) -- (B) -- (C) -- cycle;
  \draw[given] (A) -- (E);
  \draw[given] (B) -- (D);
  \node[vlab, left]  at (A) {$A$};
  \node[vlab, right] at (B) {$B$};
  \node[vlab, above] at (C) {$C$};
  \node[vlab, above left]  at (D) {$D$};
  \node[vlab, above right] at (E) {$E$};
  % All four angles sit at the BASE vertices: the cevians split the angle at
  % A into 1 (next to AB) and 3 (next to AC), and the angle at B into 2 and
  % 4.  The draft put 3 at D and 4 at E, which is not what the plate shows
  % and left Exercise 16 marking angles it never names.
  \VIarc{B}{A}{E}{0.66cm}   \VInumb{B}{A}{E}{0.94cm}{1}
  \VIarc{A}{B}{D}{0.66cm}   \VInumb{A}{B}{D}{0.94cm}{2}
  \VIarc{E}{A}{C}{0.40cm}   \VInumb{E}{A}{C}{0.98cm}{3}
  \VIarc{D}{B}{C}{0.40cm}   \VInumb{D}{B}{C}{0.98cm}{4}
\end{tikzpicture}
\figcap{6--87}
\end{minipage}\hfill
\begin{minipage}{0.44\linewidth}\centering
% Proportions remeasured against scan07 p18 (book p.103) at 150 dpi.  The draft
% kite was twice as wide for its height as the book's (height/width 0.96 against
% the plate's 1.76) and put D barely below BC, which is what drove D's letter
% into the numeral 2.  Measured off the plate, with BC as the unit: A sits 0.520
% above BC, D 0.495 below it, E 1.244 below it.
% Everything is mirror-symmetric about the vertical AE, so BE = CE and angle 1 =
% angle 2 (the hypotheses) and angle 3 = angle 4 (the conclusion) all hold
% exactly rather than by eye.
\begin{tikzpicture}[scale=1.062]
  \coordinate (B) at (0,0);
  \coordinate (C) at (2.44,0);
  \coordinate (A) at (1.22,1.269);
  \coordinate (D) at (1.22,-1.208);
  \coordinate (E) at (1.22,-3.035);
  \draw[given] (B) -- (A) -- (C);
  \draw[given] (B) -- (D) -- (C);
  \draw[given] (B) -- (E) -- (C);
  \draw[given] (A) -- (E);
  \node[vlab, above] at (A) {$A$};
  \node[vlab, left]  at (B) {$B$};
  \node[vlab, right] at (C) {$C$};
  % The plate letters D in the wedge between DA and DC, just up and right of the
  % point -- the only quadrant at D that no stroke crosses.  That wedge is only
  % 45 degrees wide (DA runs at 90, DC at 44.7), so "above right" -- the 45
  % degree diagonal -- laid the letter exactly along DC.  Measured on the plate
  % the letter sits at 67.7 degrees, which is the bisector; place it there.
  \node[vlab] at ($(D)+(70:0.58)$) {$D$};
  \node[vlab, below] at (E) {$E$};
  % 3 and 4 are at B and C -- the angles ABD and ACD that Exercise 17 asks
  % about -- not at D, where an earlier draft put them.  1 and 2 stay at E.
  % Distances measured off the plate: numerals at 0.33 of BC from B and C, arcs
  % at 0.21; at E the numerals sit 0.39 of BC out and the arcs at 0.28, so there
  % the numeral is OUTSIDE its arc, as the book draws it.
  \VIarc{A}{B}{D}{0.52cm}   \VInumb{A}{B}{D}{0.80cm}{3}
  \VIarc{A}{C}{D}{0.52cm}   \VInumb{A}{C}{D}{0.80cm}{4}
  \VIarc{B}{E}{A}{0.68cm}   \VInumb{B}{E}{A}{0.96cm}{1}
  \VIarc{A}{E}{C}{0.68cm}   \VInumb{A}{E}{C}{0.96cm}{2}
\end{tikzpicture}
\figcap{6--88}
\end{minipage}
\end{bkfigure}}

% ================================================================= Fig 6-89
% Book p.103.  Parallelogram D C A B with E and F on the diagonal, angles 1, 2.
\newcommand{\FIGVIEIGHTYNINE}{\begin{bkfigure}[\linewidth]
\begin{tikzpicture}[scale=1.118]
  \coordinate (A) at (0,0);
  \coordinate (B) at (2.82,0);
  \coordinate (D) at (0.52,1.22);
  \coordinate (C) at (3.34,1.22);
  % The book draws the diagonal A--C and hangs D--E and B--F off it, with E
  % and F on that diagonal; there is no D--B.  Angle 1 = ADE at D and
  % angle 2 = CBF at B are the pair Exercise 19 gives, and DE = BF is what it
  % asks for -- so both of those segments have to be in the picture.  The
  % draft drew D--B instead and left E and F unjoined to anything.
  \draw[given] (A) -- (B) -- (C) -- (D) -- cycle;
  \draw[given] (A) -- (C);
  \coordinate (E) at ($(A)!0.30!(C)$);
  \coordinate (F) at ($(A)!0.70!(C)$);
  \draw[given] (D) -- (E);
  \draw[given] (B) -- (F);
  \node[vlab, below left]  at (A) {$A$};
  \node[vlab, below right] at (B) {$B$};
  \node[vlab, above left]  at (D) {$D$};
  \node[vlab, above right] at (C) {$C$};
  \node[vlab] at ($(E)+(82:0.46)$) {$E$};
  \node[vlab] at ($(F)+(262:0.46)$) {$F$};
  \VIarc{A}{D}{E}{0.36cm}   \VInumb{A}{D}{E}{0.62cm}{1}
  \VIarc{C}{B}{F}{0.36cm}   \VInumb{C}{B}{F}{0.62cm}{2}
\end{tikzpicture}
\figcap{6--89}
\end{bkfigure}}

% =========================================================== Figs 6-90, 6-91
% Book p.104.  6-90: trapezoid A B C D with both diagonals crossing at E.
%   6-91: parallelogram D C A B with E, F on the diagonal, angles 1..4.
\newcommand{\FIGVININETYNINETYONE}{\begin{bkfigure}[\linewidth]
\begin{minipage}{0.48\linewidth}\centering
\begin{tikzpicture}[scale=0.873]
  % Exercise 20 gives AD = CB and BD = AC, so the trapezoid is isosceles;
  % the draft's A and B were 0.52 and 0.50 in from the ends.
  \coordinate (D) at (0,0);
  \coordinate (C) at (2.92,0);
  \coordinate (A) at (0.51,1.12);
  \coordinate (B) at (2.41,1.12);
  \draw[given] (D) -- (C) -- (B) -- (A) -- cycle;
  \draw[given] (A) -- (C);
  \draw[given] (D) -- (B);
  \coordinate (E) at (intersection of A--C and D--B);
  \node[vlab, below left]  at (D) {$D$};
  \node[vlab, below right] at (C) {$C$};
  \node[vlab, above left]  at (A) {$A$};
  \node[vlab, above right] at (B) {$B$};
  \node[vlab, below] at (E) {$E$};
\end{tikzpicture}
\figcap{6--90}
\end{minipage}\hfill
\begin{minipage}{0.48\linewidth}\centering
\begin{tikzpicture}[scale=1.254]
  \coordinate (A) at (0,0);
  \coordinate (B) at (2.82,0);
  \coordinate (D) at (0.52,1.22);
  \coordinate (C) at (3.34,1.22);
  % Same construction as Fig. 6-89, and checked against the same plate: the
  % diagonal A--C with D--E and B--F hung off it, and NO D--B.  Exercise 21
  % gives AB = CD, AE = CF and angle 3 = angle 4, and asks for angle 1 =
  % angle 2, where 1 = ADE at D, 2 = CBF at B, 3 = BAC at A and 4 = DCA at C.
  % The draft drew D--B, left E and F unjoined, and marked four different
  % angles.
  \draw[given] (A) -- (B) -- (C) -- (D) -- cycle;
  \draw[given] (A) -- (C);
  \coordinate (E) at ($(A)!0.30!(C)$);
  \coordinate (F) at ($(A)!0.70!(C)$);
  \draw[given] (D) -- (E);
  \draw[given] (B) -- (F);
  \node[vlab, below left]  at (A) {$A$};
  \node[vlab, below right] at (B) {$B$};
  \node[vlab, above left]  at (D) {$D$};
  \node[vlab, above right] at (C) {$C$};
  \node[vlab] at ($(E)+(70:0.34)$) {$E$};
  \node[vlab] at ($(F)+(250:0.34)$) {$F$};
  \VIarc{A}{D}{E}{0.36cm}   \VInumb{A}{D}{E}{0.62cm}{1}
  \VIarc{C}{B}{F}{0.36cm}   \VInumb{C}{B}{F}{0.62cm}{2}
  \VIarc{B}{A}{C}{0.36cm}   \VInumb{B}{A}{C}{1.24cm}{3}
  \VIarc{D}{C}{A}{0.36cm}   \VInumb{D}{C}{A}{1.24cm}{4}
\end{tikzpicture}
\figcap{6--91}
\end{minipage}
\end{bkfigure}}

% ================================================================= Fig 6-92
% Book p.104, Exercise 22.  Rebuilt against scan07 p19 (book p.104).
% The book draws the OUTER PENTAGON A B E D C as well as the five chords --
% an earlier draft drew the bare five-pointed star, leaving the outer edges
% A--B, B--E, E--D, D--C, C--A off the page even though the exercise's two
% triangles ABC and ADE are built from them.
% The pentagon is NOT regular: measured off the plate it is markedly wider
% than tall (width 2.84 against height 2.37).  It IS mirror-symmetric about
% the vertical through A, which is what makes AB = AC, AD = AE and
% angle 1 = angle 2 hold exactly -- the hypotheses of Exercise 22.
\newcommand{\FIGVININETYTWO}{\begin{bkfigure}[\linewidth]
\begin{tikzpicture}[scale=1.612]
  \coordinate (A) at (0,0);
  \coordinate (B) at (-1.42,1.214);
  \coordinate (C) at (1.42,1.214);
  \coordinate (E) at (-0.81,2.366);
  \coordinate (D) at (0.81,2.366);
  \draw[given] (A) -- (B) -- (E) -- (D) -- (C) -- cycle;   % outer pentagon
  \draw[given] (B) -- (C);                                 % the five chords
  \draw[given] (B) -- (D);
  \draw[given] (C) -- (E);
  \draw[given] (A) -- (D);
  \draw[given] (A) -- (E);
  % Inner pentagon: every vertex is the crossing of two chords, so each is
  % derived, never placed by hand.
  \coordinate (F) at (intersection of B--D and C--E);
  \coordinate (G) at (intersection of B--D and A--E);
  \coordinate (H) at (intersection of A--D and C--E);
  \coordinate (I) at (intersection of A--E and B--C);
  \coordinate (J) at (intersection of A--D and B--C);
  \node[vlab, above left]  at (E) {$E$};
  \node[vlab, above right] at (D) {$D$};
  \node[vlab, left]  at (B) {$B$};
  \node[vlab, right] at (C) {$C$};
  \node[vlab, below] at (A) {$A$};
  % Label sides re-measured on scan07 p19 (book p.104) at 4x, by solving each
  % crossing from the plate's own vertices rather than eyeballing it:
  %   F point (568,266), letter (565,210) -> directly ABOVE   (draft: below)
  %   G point (369,357), letter (413,383) -> down and RIGHT   (draft: above left)
  %   H point (742,363), letter (680,390) -> down and LEFT    (draft: above right)
  % G and H both lean INWARD, into the wedge under the chord that makes them;
  % the draft had all three on the opposite side of their own point.
  \node[vlab] at ($(F)+(90:0.30)$) {$F$};
  \node[vlab] at ($(G)+(330:0.32)$) {$G$};
  \node[vlab] at ($(H)+(205:0.34)$) {$H$};
  \node[vlab, below left]  at (I) {$I$};
  \node[vlab, below right] at (J) {$J$};
  % Angle 1 = BAE and angle 2 = DAC (the two given angles at the vertex);
  % angle 3 = DBC and angle 4 = ECB (the pair to be proved congruent).
  % Measured on the plate: the numerals at A sit 0.205 of BC out from the
  % vertex, just outside their arcs (0.16 of BC).  The draft put them at 0.35
  % and 0.32 of BC, which carried numeral 1 up into the letter I -- the two
  % then printed as "I_1" and the collision audit merged them into one word
  % and reported nothing.  See the note on stacked labels in ch06-UNCERTAIN.
  \VIarc{B}{A}{E}{0.46cm}   \VInumb{B}{A}{E}{0.78cm}{1}
  \VIarc{D}{A}{C}{0.46cm}   \VInumb{D}{A}{C}{0.78cm}{2}
  % The wedges at B and C are narrow, so a numeral 1.00 out along the
  % bisector still grazed the chords BD and CE; 0.72 keeps it inside.
  \VIarc{D}{B}{C}{0.46cm}   \VInumb{D}{B}{C}{0.72cm}{3}
  \VIarc{E}{C}{B}{0.46cm}   \VInumb{E}{C}{B}{0.72cm}{4}
\end{tikzpicture}
\figcap{6--92}
\end{bkfigure}}

% ================================================================= Fig 6-93
% Book p.104.  Three copies of "a ray OA from a point O on a line l, with B on
% l": the angle AOB acute, right, and obtuse.
\newcommand{\FIGVININETYTHREE}{\begin{bkfigure}[\linewidth]
\begin{tikzpicture}[scale=1.170]
 \foreach \dx/\ax/\ay in {0/0.86/1.22, 2.62/0.32/1.32, 5.24/-0.32/1.28}{
  \begin{scope}[xshift=\dx cm]
   \coordinate (O) at (0.32,0);
   \coordinate (B) at (1.12,0);
   \draw[given] (-0.22,0) -- (1.62,0);
   \draw[given] (O) -- ($(O)+(\ax,\ay)$);
   \dt{O}
   \dt{B}
   \node[vlab, above right] at ($(O)+(\ax,\ay)$) {$A$};
   \node[vlab, below] at (O) {$O$};
   \node[vlab, above right] at (B) {$B$};
   \node[slab, left] at (-0.22,0) {$l$};
  \end{scope}}
\end{tikzpicture}
\figcap{6--93}
\end{bkfigure}}

% ================================================================= Fig 6-94
% Book p.105.  The three cases of the right-angle construction: C on the same
% side of O as B; C = O; C on the opposite side.  A' is the reflection of A in
% l, and AA' meets l at C.  The circular sweep is the book's arc marking the
% two congruent angles at O.
\newcommand{\FIGVININETYFOUR}{\begin{bkfigure}[\linewidth]
\begin{tikzpicture}[scale=1.118]
  % Arcs rebuilt against scan07 p20 (book p.105) at 150 dpi.  The draft drew a
  % full CIRCLE round O in all three panels; the book draws no circle anywhere
  % here.  It strikes one arc from ray OA round to ray OA' on the side B lies
  % on -- the two supplementary angles the proof pairs up -- which is a straight
  % angle in case 2 and a reflex one in case 3, so \VIarc (always minor) cannot
  % express it and \VIarcvia is used instead.  Case 3 carries a SECOND, smaller
  % arc the other way round, through C's side, because there the proof pairs
  % angle AOC with angle A'OC and those now lie opposite B.
  % --- case 1: C between O and B
  \begin{scope}
   \coordinate (O) at (0.32,0);
   % Measured off the scan, OB is 2.25 times OC; the draft had it at 1.86, which
   % is what pushed B's letter up against C's.
   \coordinate (B) at (1.78,0);
   \coordinate (A) at (1.02,1.42);
   \coordinate (Ap) at (1.02,-1.42);
   \coordinate (C) at (intersection of A--Ap and O--B);
   \draw[given] (-0.12,0) -- (1.96,0);
   \draw[given] (O) -- (A);
   \draw[given] (O) -- (Ap);
   \draw[given] (A) -- (Ap);
   \VIarcvia{A}{O}{Ap}{B}{0.60cm}
   \dt{O} \dt{B}
   \node[slab, left] at (-0.12,0) {$l$};
   % Straight up from O is outside the arc's sector (it spans OA round to OA'
   % through B), so the book's own "just above the point" placement is clear.
   \node[vlab, above left, inner sep=0.6pt] at (O) {$O$};
   \node[vlab, above right] at (C) {$C$};
   % Centred, not above-right: set to the right it reached across the panel gap
   % and came within 0.5pt of the NEXT panel's l.  The plate centres it too.
   \node[vlab, above] at (B) {$B$};
   \node[vlab, above left] at (A) {$A$};
   \node[vlab, below left] at (Ap) {$A\VIpr$};
  \end{scope}
  % --- case 2: C = O
  \begin{scope}[xshift=2.72cm]
   \coordinate (O) at (0.42,0);
   \coordinate (B) at (1.52,0);
   \coordinate (A) at (0.42,1.42);
   \coordinate (Ap) at (0.42,-1.42);
   \draw[given] (-0.12,0) -- (1.92,0);
   \draw[given] (A) -- (Ap);
   \VIarcvia{A}{O}{Ap}{B}{0.62cm}
   \dt{O} \dt{B}
   \node[slab, left] at (-0.12,0) {$l$};
   % O and C coincide here, so the book cannot letter the point in place: it
   % sets the name clear of the drawing, up and to the left, and runs a leader
   % arrow down to the point.  The old draft instead anchored this label
   % south-EAST of a point already left of the panel's own left edge, so the
   % text grew leftwards out of the panel and landed on top of Fig. 6-94's
   % first panel -- which is where the "O=CB" and "C" pile-up came from.
   % Kept clear of the B/C letter band the three panels share: at y = 0.72 this
   % label's box still overlapped the neighbouring panel's B by 0.3pt.
   \node[vlab, anchor=east] at ($(O)+(-0.10,1.02)$) {$O = C$};
   \draw[fig, line width=0.5pt, ->] ($(O)+(-0.34,0.78)$) -- ($(O)+(-0.07,0.11)$);
   \node[vlab, above right] at (B) {$B$};
   \node[vlab, above left] at (A) {$A$};
   \node[vlab, below left] at (Ap) {$A\VIpr$};
  \end{scope}
  % --- case 3: C on the far side of O from B
  \begin{scope}[xshift=5.44cm]
   \coordinate (O) at (1.42,0);
   % B must lie OUTSIDE both arcs.  Measured on the plate (book p.105, photo
   % PDF p.119, panel 3 at 5x): OC = 217 px and OB = 270 px, so OB = 1.244*OC,
   % while the outer arc's radius is 218 px -- it passes essentially through C
   % and B sits clear beyond it.  The draft had OB = 0.667*OC, which put B
   % INSIDE the arc and drove its letter onto the arc stroke.
   \coordinate (B) at (2.54,0);
   \coordinate (A) at (0.52,1.42);
   \coordinate (Ap) at (0.52,-1.42);
   \coordinate (C) at (intersection of A--Ap and O--B);
   \draw[given] (0.02,0) -- (2.76,0);
   \draw[given] (O) -- (A);
   \draw[given] (O) -- (Ap);
   \draw[given] (A) -- (Ap);
   % Measured off the scan: the inner arc crosses l at 0.72 of OC from O, the
   % outer at 0.98 of it.
   \VIarcvia{A}{O}{Ap}{C}{0.65cm}
   \VIarcvia{A}{O}{Ap}{B}{0.88cm}
   \dt{O} \dt{B}
   \node[slab, left] at (0.02,0) {$l$};
   \node[vlab, below right, inner sep=0.6pt] at (O) {$O$};
   \node[vlab, above left]  at (C) {$C$};
   \node[vlab, above] at (B) {$B$};
   \node[vlab, above left] at (A) {$A$};
   \node[vlab, below left] at (Ap) {$A\VIpr$};
  \end{scope}
\end{tikzpicture}
\figcap{6--94}
\end{bkfigure}}

% =========================================================== Figs 6-95, 6-96
% Book p.105.  6-95: P at the left with Q, R, S stacked at the right.
%   6-96: the kite B A D C with spine A--C through E.
\newcommand{\FIGVININETYFIVENINETYSIX}{\begin{bkfigure}[\linewidth]
\begin{minipage}{0.50\linewidth}\centering
\begin{tikzpicture}[scale=1.028]
  \coordinate (P) at (0,0.72);
  \coordinate (Q) at (2.62,1.44);
  \coordinate (R) at (2.62,0.72);
  \coordinate (S) at (2.62,0);
  \draw[given] (P) -- (Q) -- (S) -- cycle;
  \draw[given] (P) -- (R);
  \draw[given] (Q) -- (S);
  \node[vlab, left]  at (P) {$P$};
  \node[vlab, right] at (Q) {$Q$};
  \node[vlab, right] at (R) {$R$};
  \node[vlab, right] at (S) {$S$};
\end{tikzpicture}
\figcap{6--95}
\end{minipage}\hfill
\begin{minipage}{0.46\linewidth}\centering
\begin{tikzpicture}[scale=1.028]
  % Off scan07 p20 the book's kite is narrower than the draft's: the
  % half-width is about 0.6 of the upper axis, and the lower axis is about
  % 1.7 times the upper.  E is taken as the actual crossing, not by ratio.
  \coordinate (A) at (1.32,1.82);
  \coordinate (B) at (0.57,0.92);
  \coordinate (D) at (2.07,0.92);
  \coordinate (C) at (1.32,-0.32);
  \draw[given] (B) -- (A) -- (D) -- (C) -- cycle;
  \draw[given] (B) -- (D);
  \draw[given] (A) -- (C);
  \coordinate (E) at (intersection of B--D and A--C);
  \node[vlab, above] at (A) {$A$};
  \node[vlab, left]  at (B) {$B$};
  \node[vlab, right] at (D) {$D$};
  \node[vlab, below] at (C) {$C$};
  \node[vlab] at ($(E)+(250:0.42)$) {$E$};
\end{tikzpicture}
\figcap{6--96}
\end{minipage}
\end{bkfigure}}

% ================================================================= Fig 6-97
% Book p.106.  The comparison picture for "less than": angle B has been moved
% to the position B' with one side along the first side of A, and its second
% side (dashed) falling interior to A.  A and B are drawn side by side.
\newcommand{\FIGVININETYSEVEN}{\begin{bkfigure}[\linewidth]
\begin{tikzpicture}[scale=1.196]
  % angle A, with the moved copy B' inside it
  \coordinate (A) at (2.42,0);
  \draw[given] (0.10,0) -- (A);
  \draw[given] (A) -- ($(A)+(118:1.52)$);
  \draw[aux]   (A) -- ($(A)+(146:1.32)$);
  \node[vlab, below right] at (A) {$A$};
  \node[slab, above] at ($(A)+(146:1.42)$) {$B\VIpr$};
  % angle B, drawn separately to the right
  \coordinate (B) at (5.62,0);
  \draw[given] (2.92,0) -- (B);
  \draw[given] (B) -- ($(B)+(132:1.62)$);
  \node[vlab, below right] at (B) {$B$};
\end{tikzpicture}
\figcap{6--97}
\end{bkfigure}}

% =========================================================== Figs 6-98, 6-99
% Book p.106.  Definition 6-5: (a) second side of B' interior to A  => B < A;
%   (b) second side of B' exterior to A  => B > A.
\newcommand{\FIGVININETYEIGHTNINETYNINE}{\begin{bkfigure}[\linewidth]
\begin{minipage}{0.48\linewidth}\centering
\begin{tikzpicture}[scale=0.885]
  \coordinate (A) at (2.02,0);
  \draw[given] (0.10,0) -- (A);
  \draw[given] (A) -- ($(A)+(112:1.42)$);
  \draw[aux]   (A) -- ($(A)+(142:1.22)$);
  \node[vlab, below right] at (A) {$A$};
  \node[slab, above] at ($(A)+(142:1.32)$) {$B\VIpr$};
  \coordinate (B) at (4.32,0);
  \draw[given] (2.42,0) -- (B);
  \draw[given] (2.42,0) -- ($(2.42,0)+(38:1.62)$);
  \node[vlab, below right] at (B) {$B$};
\end{tikzpicture}
\figcap{6--98}
\end{minipage}\hfill
\begin{minipage}{0.48\linewidth}\centering
\begin{tikzpicture}[scale=0.885]
  \coordinate (A) at (2.02,0);
  \draw[given] (0.10,0) -- (A);
  \draw[given] (A) -- ($(A)+(122:1.42)$);
  \draw[aux]   (A) -- ($(A)+(96:1.52)$);
  \node[vlab, below right] at (A) {$A$};
  \node[slab, left] at ($(A)+(150:0.92)$) {$B\VIpr$};
  \coordinate (B) at (4.42,0);
  \draw[given] (2.52,0) -- (B);
  \draw[given] (2.52,0) -- ($(2.52,0)+(28:1.72)$);
  \node[vlab, below right] at (B) {$B$};
\end{tikzpicture}
\figcap{6--99}
\end{minipage}
\end{bkfigure}}

% ========================================================= Figs 6-100, 6-101
% Book p.107.  6-100 shows B < A (B' interior to A); 6-101 shows the same pair
% with A moved instead, so that A' is exterior to B.
\newcommand{\FIGVIONEHUNDREDONEHUNDREDONE}{\begin{bkfigure}[\linewidth]
\begin{minipage}{0.48\linewidth}\centering
\begin{tikzpicture}[scale=0.878]
  \coordinate (A) at (2.22,0);
  \draw[given] (0.10,0) -- (A);
  \draw[given] (A) -- ($(A)+(118:1.42)$);
  \draw[aux]   (A) -- ($(A)+(150:1.22)$);
  \node[vlab, below right] at (A) {$A$};
  \node[slab, above] at ($(A)+(150:1.34)$) {$B\VIpr$};
  \coordinate (B) at (4.52,0);
  \draw[given] (2.62,0) -- (B);
  \draw[given] (2.62,0) -- ($(2.62,0)+(32:1.62)$);
  \node[vlab, below right] at (B) {$B$};
\end{tikzpicture}
\figcap{6--100}
\end{minipage}\hfill
\begin{minipage}{0.48\linewidth}\centering
\begin{tikzpicture}[scale=0.878]
  \coordinate (A) at (1.92,0);
  \draw[given] (0.10,0) -- (A);
  \draw[given] (A) -- ($(A)+(126:1.42)$);
  \node[vlab, below right] at (A) {$A$};
  \coordinate (B) at (4.42,0);
  \draw[given] (2.32,0) -- (B);
  \draw[given] (B) -- ($(B)+(152:1.62)$);
  \draw[aux]   (B) -- ($(B)+(122:1.52)$);
  \node[vlab, below right] at (B) {$B$};
  \node[slab, above] at ($(B)+(122:1.62)$) {$A\VIpr$};
\end{tikzpicture}
\figcap{6--101}
\end{minipage}
\end{bkfigure}}

% ================================================================ Fig 6-102
% Book p.107.  The Theorem 6-8 proof: on angle A, points D and E on the sides
% with DE cutting the second side of B' at F; on angle A', the matching D', E'
% with F' chosen so that DF = D'F'.  F is DERIVED as the crossing of DE with
% the second side of B'.
\newcommand{\FIGVIONEHUNDREDTWO}{\begin{bkfigure}[\linewidth]
\begin{tikzpicture}[scale=1.170]
  % --- left: angle A with B' interior
  \coordinate (A) at (2.42,0);
  \coordinate (E) at ($(A)+(150:1.62)$);
  \coordinate (D) at ($(A)+(212:1.62)$);
  \draw[given] (A) -- ($(A)+(150:2.02)$);
  \draw[given] (A) -- ($(A)+(212:2.02)$);
  \draw[given] (D) -- (E);
  \coordinate (Bs) at ($(A)+(174:2.02)$);
  \draw[given] (A) -- (Bs);
  \coordinate (F) at (intersection of D--E and A--Bs);
  \node[vlab, below right] at (A) {$A$};
  \node[vlab, above right] at (E) {$E$};
  % D sits on the lower side of the angle, which runs on past it to
  % 212:2.02 -- so the letter goes square off that ray, not along it.
  \node[vlab] at ($(D)+(302:0.34)$) {$D$};
  \node[vlab, above left] at (F) {$F$};
  \node[slab, above right] at ($(A)+(120:1.13)$) {$B\VIpr$};
  % --- right: angle A' with the matching construction
  \begin{scope}[xshift=3.62cm]
   \coordinate (B) at (2.42,0);
   \coordinate (Ep) at ($(B)+(150:1.62)$);
   \coordinate (Dp) at ($(B)+(212:1.62)$);
   \draw[given] (B) -- ($(B)+(150:2.02)$);
   \draw[given] (B) -- ($(B)+(212:2.02)$);
   \draw[aux]   (Dp) -- (Ep);
   \coordinate (As) at ($(B)+(174:2.02)$);
   \draw[given] (B) -- (As);
   \coordinate (Fp) at (intersection of Dp--Ep and B--As);
   \node[vlab, below right] at (B) {$B$};
   \node[vlab, above right] at (Ep) {$E\VIpr$};
   \node[vlab] at ($(Dp)+(302:0.34)$) {$D\VIpr$};
   \node[vlab, above left] at (Fp) {$F\VIpr$};
   \node[slab, above right] at ($(B)+(132:1.22)$) {$A\VIpr$};
  \end{scope}
\end{tikzpicture}
\figcap{6--102}
\end{bkfigure}}

% ================================================================ Fig 6-103
% Book p.109.  Theorem 6-11: (a) right angle A with supplement A'; (b) the same
% for B and B'; (c) the comparison figure with B'' congruent to B, and B''' its
% supplement, all at one vertex.
\newcommand{\FIGVIONEHUNDREDTHREE}{\begin{bkfigure}[\linewidth]
\begin{tikzpicture}[scale=1.118]
  % (a)
  \begin{scope}
   \coordinate (O) at (0.86,0);
   \draw[given] (0.02,0) -- (1.72,0);
   \draw[given] (O) -- (0.86,1.32);
   \draw[tick] ($(O)+(180:0.52)$) arc[start angle=180, end angle=0, radius=0.52];
   \node[slab] at ($(O)+(136:0.76)$) {$A$};
   \node[slab] at ($(O)+(36:0.81)$)  {$A\VIpr$};
   \node[slab, below] at (0.86,-0.28) {(a)};
  \end{scope}
  % (b)
  \begin{scope}[xshift=2.32cm]
   \coordinate (O) at (0.86,0);
   \draw[given] (0.02,0) -- (1.72,0);
   \draw[given] (O) -- (0.86,1.32);
   \draw[tick] ($(O)+(180:0.52)$) arc[start angle=180, end angle=0, radius=0.52];
   \node[slab] at ($(O)+(136:0.76)$) {$B$};
   \node[slab] at ($(O)+(32:0.80)$)  {$B\VIpr$};
   \node[slab, below] at (0.86,-0.28) {(b)};
  \end{scope}
  % (c)
  \begin{scope}[xshift=4.64cm]
   \coordinate (O) at (1.12,0);
   \draw[given] (0.02,0) -- (2.22,0);
   \draw[given] (O) -- ($(O)+(90:1.32)$);
   \draw[given] (O) -- ($(O)+(62:1.32)$);
   \draw[tick] ($(O)+(180:0.72)$) arc[start angle=180, end angle=0, radius=0.72];
   \node[slab] at ($(O)+(142:1.00)$) {$B\VIppr$};
   \node[slab] at ($(O)+(80:1.24)$) {$A$};
   \node[slab] at ($(O)+(0:1.32)$)  {$A\VIpr$};
   \node[slab] at ($(O)+(30:1.23)$)  {$B\VIpppr$};
   \node[slab, below] at (1.12,-0.28) {(c)};
  \end{scope}
\end{tikzpicture}
\figcap{6--103}
\end{bkfigure}}

% ================================================================ Fig 6-104
% Book p.109.  (a) l perpendicular to m; (b) l the perpendicular bisector of AB.
\newcommand{\FIGVIONEHUNDREDFOUR}{\begin{bkfigure}[\linewidth]
\begin{tikzpicture}[scale=1.170]
  % (a)
  \begin{scope}
   \coordinate (O) at (0.92,0);
   \coordinate (R) at (1.82,0);
   \coordinate (U) at (0.92,1.32);
   \draw[given] (0.02,0) -- (R);
   \draw[given] (O) -- (U);
   \VIrt{R}{O}{U}{6pt}
   \node[slab, above left] at (U) {$l$};
   \node[slab, right] at (R) {$m$};
   \node[slab, below] at (0.92,-0.36) {(a)};
  \end{scope}
  % (b)
  \begin{scope}[xshift=2.92cm]
   \coordinate (A) at (0.02,0);
   \coordinate (B) at (1.92,0);
   \coordinate (M) at ($(A)!0.5!(B)$);
   \coordinate (U) at ($(M)+(0,1.32)$);
   \draw[given] (A) -- (B);
   \draw[given] (M) -- (U);
   \VIrt{B}{M}{U}{6pt}
   \node[slab, above left] at (U) {$l$};
   \node[vlab, left]  at (A) {$A$};
   \node[vlab, right] at (B) {$B$};
   \node[slab, below] at (0.96,-0.36) {(b)};
  \end{scope}
\end{tikzpicture}
\figcap{6--104}
\end{bkfigure}}

% ======================================================= Figs 6-105, 6-106
% Book p.110.  6-105: rectangle A B C D with both diagonals.
%   6-106: triangle A C with apex B and cevian BD.
\newcommand{\FIGVIONEHUNDREDFIVEONEHUNDREDSIX}{\begin{bkfigure}[\linewidth]
\begin{minipage}{0.52\linewidth}\centering
\begin{tikzpicture}[scale=0.955]
  \coordinate (D) at (0,0);
  \coordinate (C) at (2.92,0);
  \coordinate (A) at (0,1.22);
  \coordinate (B) at (2.92,1.22);
  \draw[given] (A) -- (B) -- (C) -- (D) -- cycle;
  \draw[given] (A) -- (C);
  \draw[given] (B) -- (D);
  \node[vlab, above left]  at (A) {$A$};
  \node[vlab, above right] at (B) {$B$};
  \node[vlab, below left]  at (D) {$D$};
  \node[vlab, below right] at (C) {$C$};
\end{tikzpicture}
\figcap{6--105}
\end{minipage}\hfill
\begin{minipage}{0.44\linewidth}\centering
\begin{tikzpicture}[scale=0.955]
  \coordinate (A) at (0,0);
  \coordinate (C) at (2.72,0);
  \coordinate (B) at (1.02,1.62);
  \coordinate (D) at ($(A)!0.44!(C)$);
  \draw[given] (A) -- (C) -- (B) -- cycle;
  \draw[given] (B) -- (D);
  \node[vlab, left]  at (A) {$A$};
  \node[vlab, right] at (C) {$C$};
  \node[vlab, above] at (B) {$B$};
  \node[vlab, below] at (D) {$D$};
\end{tikzpicture}
\figcap{6--106}
\end{minipage}
\end{bkfigure}}

% ======================================================= Figs 6-107, 6-108
% Book p.110.  6-107: rectangle A B C D with E on the base and the two
%   segments AE, BE.  6-108: triangle A C with apex B and cevian BD.
\newcommand{\FIGVIONEHUNDREDSEVENONEHUNDREDEIGHT}{\begin{bkfigure}[\linewidth]
\begin{minipage}{0.52\linewidth}\centering
\begin{tikzpicture}[scale=0.947]
  \coordinate (D) at (0,0);
  \coordinate (C) at (3.02,0);
  \coordinate (A) at (0,1.12);
  \coordinate (B) at (3.02,1.12);
  \coordinate (E) at ($(D)!0.5!(C)$);
  \draw[given] (A) -- (B) -- (C) -- (D) -- cycle;
  \draw[given] (A) -- (E);
  \draw[given] (B) -- (E);
  \node[vlab, above left]  at (A) {$A$};
  \node[vlab, above right] at (B) {$B$};
  \node[vlab, below left]  at (D) {$D$};
  \node[vlab, below right] at (C) {$C$};
  \node[vlab, below] at (E) {$E$};
\end{tikzpicture}
\figcap{6--107}
\end{minipage}\hfill
\begin{minipage}{0.44\linewidth}\centering
\begin{tikzpicture}[scale=0.947]
  \coordinate (A) at (0,0);
  \coordinate (C) at (2.82,0);
  \coordinate (B) at (1.22,1.62);
  \coordinate (D) at ($(A)!0.5!(C)$);
  \draw[given] (A) -- (C) -- (B) -- cycle;
  \draw[given] (B) -- (D);
  \node[vlab, left]  at (A) {$A$};
  \node[vlab, right] at (C) {$C$};
  \node[vlab, above] at (B) {$B$};
  \node[vlab, below] at (D) {$D$};
\end{tikzpicture}
\figcap{6--108}
\end{minipage}
\end{bkfigure}}

% ======================================================= Figs 6-109, 6-110
% Book p.110.  6-109: B--D and A--F crossing at C, with E above.
%   6-110: right triangle D A C with B on DC and E on AD.
\newcommand{\FIGVIONEHUNDREDNINEONEHUNDREDTEN}{\begin{bkfigure}[\linewidth]
\begin{minipage}{0.50\linewidth}\centering
\begin{tikzpicture}[scale=0.960]
  \coordinate (B) at (0.32,1.02);
  \coordinate (D) at (2.62,0.62);
  \coordinate (A) at (0,0);
  \coordinate (E) at (2.92,1.42);
  \draw[given] (B) -- (D);
  \draw[given] (A) -- (E);
  \coordinate (C) at (intersection of B--D and A--E);
  \node[vlab, left]  at (B) {$B$};
  \node[vlab, right] at (D) {$D$};
  \node[vlab, left]  at (A) {$A$};
  \node[vlab, right] at (E) {$E$};
  \node[vlab, below] at (C) {$C$};
\end{tikzpicture}
\figcap{6--109}
\end{minipage}\hfill
\begin{minipage}{0.46\linewidth}\centering
\begin{tikzpicture}[scale=0.960]
  \coordinate (A) at (0,0);
  \coordinate (C) at (2.62,0);
  \coordinate (D) at (0.42,1.72);
  \coordinate (B) at ($(D)!0.52!(C)$);
  \coordinate (E) at ($(A)!0.62!(D)$);
  \draw[given] (A) -- (C) -- (D) -- cycle;
  \draw[given] (A) -- (B);
  \draw[given] (E) -- (B);
  \VIrt{D}{B}{A}{5pt}
  \node[vlab, below left]  at (A) {$A$};
  \node[vlab, right] at (C) {$C$};
  \node[vlab, above] at (D) {$D$};
  \node[vlab, above right] at (B) {$B$};
  \node[vlab, left]  at (E) {$E$};
\end{tikzpicture}
\figcap{6--110}
\end{minipage}
\end{bkfigure}}

% ================================================================ Fig 6-111
% Book p.111.  A, C, D, B with the right angles at C (AC perp CB) and at D
% (AD perp DB); the two triangles share the hypotenuse-like segment A--B.
\newcommand{\FIGVIONEHUNDREDELEVEN}{\begin{bkfigure}[\linewidth]
\begin{tikzpicture}[scale=1.196]
  % Exercise 11 gives AC = AD with AC perp CB and AD perp DB.  Two right
  % angles on AB put C and D on the circle with diameter AB; AC = AD then
  % makes them a mirror pair across AB.  These coordinates are that pair
  % (circle centre (1.47,0.71), radius 1.03276, parameter 100 deg), not an
  % eyeballed quadrilateral.
  \coordinate (A) at (0.72,1.42);
  \coordinate (D) at (2.0390,1.5719);
  \coordinate (C) at (0.6405,0.0947);
  \coordinate (B) at (2.22,0);
  \draw[given] (A) -- (C) -- (B);
  \draw[given] (A) -- (D) -- (B);
  \draw[given] (A) -- (B);
  \VIrt{A}{C}{B}{5pt}
  \VIrt{B}{D}{A}{5pt}
  \VItk{A}{C}{1}
  \VItk{A}{D}{1}
  \node[vlab, above left]  at (A) {$A$};
  \node[vlab, above right] at (D) {$D$};
  \node[vlab, left]  at (C) {$C$};
  \node[vlab, right] at (B) {$B$};
\end{tikzpicture}
\figcap{6--111}
\end{bkfigure}}

% ======================================================= Figs 6-112, 6-113
% Book p.112.  6-112: the two angle-fans of Review exercise 3 -- R, S, T at one
%   vertex and A_1, A_2 at the other, each with a sweeping arc.
%  6-113: two parallels cut by two transversals, points A, B, C, D.
\newcommand{\FIGVIONEHUNDREDTWELVEONEHUNDREDTHIRTEEN}{\begin{bkfigure}[\linewidth]
\begin{minipage}{0.48\linewidth}\centering
% Rebuilt against scan08 p5 (book p.112).  Two corrections to an earlier
% draft, both settled by Review Exercise 3 ("if angle A_1 = angle R and
% angle A_2 = angle S, prove angle B = angle T"), which makes R, S the two
% PARTS of T and A_1, A_2 the two parts of B:
%   (a) the left fan is drawn VERTEX-UP in the book, its three rays running
%       downward; the draft had it vertex-down like the right one.
%   (b) every one of the six names is an ANGLE, marked by its own arc -- R
%       and S on small arcs, T on a large one enclosing both, and likewise
%       A_1, A_2 inside B.  The draft lettered B as a POINT at the tip of the
%       middle ray, which is not what the plate shows.
% Each numeral sits on its own angle's BISECTOR, just outside that angle's
% arc, which is where the book puts them.
\begin{tikzpicture}[scale=0.895]
  % ---- left fan, vertex UP: rays at 238 / 267 / 306 degrees
  \coordinate (V) at (0,0);
  \draw[given] (V) -- ++(238:1.95);
  \draw[given] (V) -- ++(267:1.95);
  \draw[given] (V) -- ++(306:1.95);
  \draw[tick] ([shift=(238:0.66)]V) arc[start angle=238, end angle=267, radius=0.66];
  \draw[tick] ([shift=(267:0.58)]V) arc[start angle=267, end angle=306, radius=0.58];
  \draw[tick] ([shift=(238:1.30)]V) arc[start angle=238, end angle=306, radius=1.30];
  \node[slab] at ($(V)+(228.5:1.52)$) {$R$};
  \node[slab] at ($(V)+(286.5:0.96)$) {$S$};
  \node[slab] at ($(V)+(288:1.62)$)   {$T$};
  % ---- right fan, vertex DOWN: rays at 124 / 87 / 54 degrees
  \begin{scope}[xshift=3.20cm, yshift=-1.95cm]
   \coordinate (W) at (0,0);
   \draw[given] (W) -- ++(124:1.95);
   \draw[given] (W) -- ++(87:1.95);
   \draw[given] (W) -- ++(54:1.95);
   \draw[tick] ([shift=(87:0.58)]W)  arc[start angle=87,  end angle=124, radius=0.58];
   \draw[tick] ([shift=(54:0.66)]W)  arc[start angle=54,  end angle=87,  radius=0.66];
   \draw[tick] ([shift=(54:1.30)]W)  arc[start angle=54,  end angle=124, radius=1.30];
   \node[slab] at ($(W)+(77.5:1.59)$) {$A_1$};
   \node[slab] at ($(W)+(42.5:1.61)$)  {$A_2$};
   \node[slab] at ($(W)+(105.5:1.60)$) {$B$};
  \end{scope}
\end{tikzpicture}
\figcap{6--112}
\end{minipage}\hfill
\begin{minipage}{0.48\linewidth}\centering
\begin{tikzpicture}[scale=0.895]
  \draw[given] (0.02,1.12) -- (3.22,1.12);
  \draw[given] (0.02,0.42) -- (3.22,0.42);
  \coordinate (P) at (0.62,-0.12);
  \coordinate (Q) at (2.42,1.72);
  \draw[given] (P) -- (Q);
  \coordinate (A) at (intersection of P--Q and 0.02,0.42--3.22,0.42);
  \coordinate (B) at (intersection of P--Q and 0.02,1.12--3.22,1.12);
  \coordinate (C) at ($(B)+(0.42,0.16)$);
  \coordinate (D) at ($(A)+(-0.30,-0.42)$);
  \draw[given] (B) -- (C);
  \node[vlab, above left]  at (B) {$B$};
  \node[vlab, above left] at (C) {$C$};
  \node[vlab, below right] at (A) {$A$};
  \node[vlab, below] at (D) {$D$};
\end{tikzpicture}
\figcap{6--113}
\end{minipage}
\end{bkfigure}}

% ======================================================= Figs 6-114, 6-115
% Book p.112.  6-114: A O C on a line with OB perpendicular.
%   6-115: triangles E F G and R S T.
\newcommand{\FIGVIONEHUNDREDFOURTEENONEHUNDREDFIFTEEN}{\begin{bkfigure}[\linewidth]
\begin{minipage}{0.46\linewidth}\centering
\begin{tikzpicture}[scale=0.945]
  \coordinate (A) at (0,0);
  \coordinate (C) at (2.62,0);
  \coordinate (O) at ($(A)!0.5!(C)$);
  \coordinate (B) at ($(O)+(0,1.52)$);
  \draw[given] (A) -- (C);
  \draw[given] (O) -- (B);
  \node[vlab, left]  at (A) {$A$};
  \node[vlab, right] at (C) {$C$};
  \node[vlab, below] at (O) {$O$};
  \node[vlab, above] at (B) {$B$};
\end{tikzpicture}
\figcap{6--114}
\end{minipage}\hfill
\begin{minipage}{0.50\linewidth}\centering
\begin{tikzpicture}[scale=0.945]
  \coordinate (F) at (0,0);
  \coordinate (G) at (1.62,0);
  \coordinate (E) at (0.42,1.22);
  \draw[given] (F) -- (G) -- (E) -- cycle;
  \node[vlab, below left]  at (F) {$F$};
  \node[vlab, below right] at (G) {$G$};
  \node[vlab, above] at (E) {$E$};
  \coordinate (S) at (1.82,0);
  \coordinate (T) at (3.32,0.62);
  \coordinate (R) at (2.32,1.42);
  \draw[given] (S) -- (T) -- (R) -- cycle;
  \node[vlab, below left]  at (S) {$S$};
  \node[vlab, right] at (T) {$T$};
  \node[vlab, above] at (R) {$R$};
\end{tikzpicture}
\figcap{6--115}
\end{minipage}
\end{bkfigure}}

% ======================================================= Figs 6-116, 6-117
% Book p.112.  6-116: triangle R S T.
%  6-117: the river-width figure -- P and Q on opposite banks (the banks drawn
%  as the book's wavy lines), R off the line PQ, with U, S on the sight lines.
\newcommand{\FIGVIONEHUNDREDSIXTEENONEHUNDREDSEVENTEEN}{\begin{bkfigure}[\linewidth]
\begin{minipage}{0.44\linewidth}\centering
\begin{tikzpicture}[scale=0.962]
  \coordinate (T) at (0,0);
  \coordinate (S) at (2.72,0.32);
  \coordinate (R) at (0.72,1.32);
  \draw[given] (T) -- (S) -- (R) -- cycle;
  \node[vlab, below left]  at (T) {$T$};
  \node[vlab, right] at (S) {$S$};
  \node[vlab, above] at (R) {$R$};
\end{tikzpicture}
\figcap{6--116}
\end{minipage}\hfill
\begin{minipage}{0.52\linewidth}\centering
\begin{tikzpicture}[scale=0.962]
  \coordinate (R) at (1.42,1.72);
  \coordinate (U) at (0.32,-0.42);
  \coordinate (T) at (0.02,0);
  \coordinate (S) at (0.72,0);
  \coordinate (P) at (1.72,0.42);
  \coordinate (Q) at (2.92,0.12);
  \draw[given] (R) -- (U);
  \draw[given] (T) -- (P);
  \draw[given] (R) -- (P);
  \draw[aux]   (R) -- (Q);
  \draw[given] (P) -- (Q);
  % the river bank, drawn as the book's wavy rule
  \draw[tick] (2.62,0.92) .. controls (2.82,0.72) and (2.62,0.52) ..
              (2.82,0.32) .. controls (3.02,0.12) and (2.82,-0.08) .. (3.02,-0.28);
  \node[vlab, above] at (R) {$R$};
  \node[vlab, below] at (U) {$U$};
  \node[vlab, left]  at (T) {$T$};
  \node[vlab, below] at (S) {$S$};
  \node[vlab, below] at (P) {$P$};
  \node[vlab, right] at (Q) {$Q$};
\end{tikzpicture}
\figcap{6--117}
\end{minipage}
\end{bkfigure}}

% ======================================================= Figs 6-118, 6-119
% Book p.113.  6-118: triangle A C with apex B and D interior joined to all
%   three vertices.  6-119: triangle A C with apex B, D and E on the base.
\newcommand{\FIGVIONEHUNDREDEIGHTEENONEHUNDREDNINETEEN}{\begin{bkfigure}[\linewidth]
\begin{minipage}{0.46\linewidth}\centering
\begin{tikzpicture}[scale=0.951]
  \coordinate (A) at (0,0);
  \coordinate (C) at (2.62,0);
  \coordinate (B) at (1.31,1.72);
  \coordinate (D) at (1.31,0.62);
  \draw[given] (A) -- (C) -- (B) -- cycle;
  \draw[given] (A) -- (D) -- (C);
  \draw[given] (B) -- (D);
  \node[vlab, left]  at (A) {$A$};
  \node[vlab, right] at (C) {$C$};
  \node[vlab, above] at (B) {$B$};
  % D sits low, so its letter has to share the gap with the base; the book
  % crowds it the same way, and a tighter clearance is what keeps it off AC.
  \node[vlab, below, outer sep=2.0pt] at (D) {$D$};
\end{tikzpicture}
\figcap{6--118}
\end{minipage}\hfill
\begin{minipage}{0.50\linewidth}\centering
\begin{tikzpicture}[scale=0.951]
  \coordinate (A) at (0,0);
  \coordinate (C) at (3.12,0);
  % Exercise 13: ABC isosceles and angle ABD = angle EBC, so D and E are
  % symmetric about the foot of the altitude.
  \coordinate (B) at (1.56,1.62);
  \coordinate (D) at ($(A)!0.32!(C)$);
  \coordinate (E) at ($(A)!0.68!(C)$);
  \draw[given] (A) -- (C) -- (B) -- cycle;
  \draw[given] (B) -- (D);
  \draw[given] (B) -- (E);
  \node[vlab, left]  at (A) {$A$};
  \node[vlab, right] at (C) {$C$};
  \node[vlab, above] at (B) {$B$};
  \node[vlab, below] at (D) {$D$};
  \node[vlab, below] at (E) {$E$};
\end{tikzpicture}
\figcap{6--119}
\end{minipage}
\end{bkfigure}}

% ======================================================= Figs 6-120, 6-121
% Book p.113.  6-120: triangle A C with apex B, D on the base, E and F the
%   midpoints of the sides, forming the medial triangle.
%  6-121: Q S R on a line with PS perpendicular at S.
\newcommand{\FIGVIONEHUNDREDTWENTYONEHUNDREDTWENTYONE}{\begin{bkfigure}[\linewidth]
\begin{minipage}{0.52\linewidth}\centering
\begin{tikzpicture}[scale=0.964]
  \coordinate (A) at (0,0);
  \coordinate (C) at (3.02,0);
  \coordinate (B) at (1.51,1.72);
  \coordinate (D) at ($(A)!0.5!(C)$);
  \coordinate (E) at ($(A)!0.5!(B)$);
  \coordinate (F) at ($(B)!0.5!(C)$);
  \draw[given] (A) -- (C) -- (B) -- cycle;
  \draw[given] (E) -- (D) -- (F);
  \draw[given] (E) -- (F);
  \node[vlab, left]  at (A) {$A$};
  \node[vlab, right] at (C) {$C$};
  \node[vlab, above] at (B) {$B$};
  \node[vlab, below] at (D) {$D$};
  \node[vlab, above left]  at (E) {$E$};
  \node[vlab, above right] at (F) {$F$};
\end{tikzpicture}
\figcap{6--120}
\end{minipage}\hfill
\begin{minipage}{0.44\linewidth}\centering
\begin{tikzpicture}[scale=0.964]
  \coordinate (Q) at (0,0);
  \coordinate (R) at (2.42,0);
  \coordinate (S) at ($(Q)!0.5!(R)$);
  \coordinate (P) at ($(S)+(0,1.62)$);
  \draw[given] (Q) -- (R);
  \draw[given] (S) -- (P);
  \node[vlab, left]  at (Q) {$Q$};
  \node[vlab, right] at (R) {$R$};
  \node[vlab, below] at (S) {$S$};
  \node[vlab, above] at (P) {$P$};
\end{tikzpicture}
\figcap{6--121}
\end{minipage}
\end{bkfigure}}

% ======================================================= Figs 6-122, 6-123
% Book p.113.  6-122: A--B and C--D crossing at O (O bisects both).
%   6-123: A with B, C, D, E -- the two isosceles triangles of exercise 19.
\newcommand{\FIGVIONEHUNDREDTWENTYTWOONEHUNDREDTWENTYTHREE}{\begin{bkfigure}[\linewidth]
\begin{minipage}{0.50\linewidth}\centering
\begin{tikzpicture}[scale=0.958]
  \coordinate (A) at (0,1.22);
  \coordinate (B) at (3.02,0);
  % O bisects BOTH segments, so C--D's midpoint must land exactly on A--B's;
  % D was 0.02 out.
  \coordinate (C) at (1.92,1.22);
  \coordinate (D) at (1.10,0);
  \coordinate (O) at (intersection of A--B and C--D);
  \draw[given] (A) -- (B);
  \draw[given] (C) -- (D);
  \draw[given] (A) -- (D);
  \draw[given] (C) -- (B);
  \node[vlab, left]  at (A) {$A$};
  \node[vlab, right] at (B) {$B$};
  \node[vlab, above] at (C) {$C$};
  \node[vlab, below] at (D) {$D$};
  \node[vlab, below right] at (O) {$O$};
\end{tikzpicture}
\figcap{6--122}
\end{minipage}\hfill
\begin{minipage}{0.46\linewidth}\centering
\begin{tikzpicture}[scale=0.958]
  % Exercise 19 gives AB = AD, BC = DE and AC = AE.  Those three hold exactly
  % when the fan is struck with |AB| = |AD| = 1.90 and |AC| = |AE| = 1.35 and
  % the two opening angles BAC and DAE are equal (40 degrees each): the rays
  % run at -5, -45, -85 and -125 degrees from A.  The draft's five vertices
  % were eyeballed and satisfied none of the three.
  \coordinate (A) at (0.80,1.95);
  \coordinate (B) at ($(A)+(-5:1.90)$);
  \coordinate (C) at ($(A)+(-45:1.35)$);
  \coordinate (D) at ($(A)+(-85:1.90)$);
  \coordinate (E) at ($(A)+(-125:1.35)$);
  \draw[given] (A) -- (B) -- (C) -- (D) -- (E) -- cycle;
  \draw[given] (A) -- (C);
  \draw[given] (A) -- (D);
  \node[vlab, above] at (A) {$A$};
  \node[vlab, right] at (B) {$B$};
  \node[vlab, below right] at (C) {$C$};
  \node[vlab, below right] at (D) {$D$};
  \node[vlab, left]  at (E) {$E$};
\end{tikzpicture}
\figcap{6--123}
\end{minipage}
\end{bkfigure}}

% ================================================================ Fig 6-124
% Book p.113.  Equilateral triangle A with base C B, D, E, F the midpoints of
% the sides -- the medial triangle EFD of exercise 20.
\newcommand{\FIGVIONEHUNDREDTWENTYFOUR}{\begin{bkfigure}[\linewidth]
\begin{tikzpicture}[scale=1.196]
  \coordinate (C) at (0,0);
  \coordinate (B) at (2.82,0);
  \coordinate (A) at ($(C)!0.5!(B)+(0,2.44)$);
  \coordinate (D) at ($(C)!0.5!(B)$);
  \coordinate (E) at ($(C)!0.5!(A)$);
  \coordinate (F) at ($(A)!0.5!(B)$);
  \draw[given] (C) -- (B) -- (A) -- cycle;
  \draw[given] (E) -- (F);
  \draw[given] (E) -- (D);
  \draw[given] (F) -- (D);
  \node[vlab, left]  at (C) {$C$};
  \node[vlab, right] at (B) {$B$};
  \node[vlab, above] at (A) {$A$};
  \node[vlab, below] at (D) {$D$};
  \node[vlab, left]  at (E) {$E$};
  \node[vlab, right] at (F) {$F$};
\end{tikzpicture}
\figcap{6--124}
\end{bkfigure}}

% ======================================================= Figs 6-125, 6-126
% Book p.114.  6-125: A D with B, C above and the crossing segments AC, BD.
%  6-126: the cevian web -- triangle A B with apex C; D, E on the sides;
%  J, I, K on the base; F, G, H the interior crossings; angles 1, 2.
\newcommand{\FIGVIONEHUNDREDTWENTYFIVEONEHUNDREDTWENTYSIX}{\begin{bkfigure}[\linewidth]
\begin{minipage}{0.46\linewidth}\centering
\begin{tikzpicture}[scale=0.833]
  \coordinate (A) at (0,0);
  \coordinate (D) at (3.02,0);
  \coordinate (B) at (0.72,1.32);
  \coordinate (C) at (2.02,1.32);
  \draw[given] (A) -- (D);
  \draw[given] (A) -- (B) -- (D);
  \draw[given] (A) -- (C) -- (D);
  \node[vlab, left]  at (A) {$A$};
  \node[vlab, right] at (D) {$D$};
  \node[vlab, above] at (B) {$B$};
  \node[vlab, above] at (C) {$C$};
\end{tikzpicture}
\figcap{6--125}
\end{minipage}\hfill
\begin{minipage}{0.50\linewidth}\centering
% Rebuilt against scan08 p8 (book p.114).  The book joins D--J and E--K, not
% D--E.  That matters twice over: the draft's D--E left the exercise's
% segments GJ and HK undrawn, and it made G and H DEGENERATE -- they were
% asked for as (A--E cap D--E) and (B--D cap D--E), two pairs of lines that
% already share an endpoint, so G landed exactly on E and H exactly on D and
% the two letters printed on top of the other two.
% Exercise 25 gives angle 1 = angle 2, AC = BC and IJ = IK, so the plate is
% mirror-symmetric about CI: I is the MIDPOINT of AB (the draft had it at
% 0.46, which broke IJ = IK), and J, K sit symmetrically either side of it.
\begin{tikzpicture}[scale=1.479]
  \coordinate (A) at (0,0);
  \coordinate (B) at (3.22,0);
  \coordinate (C) at (1.61,1.92);
  \coordinate (I) at ($(A)!0.5!(B)$);
  \coordinate (J) at ($(A)!0.39!(B)$);
  \coordinate (K) at ($(A)!0.61!(B)$);
  \coordinate (D) at ($(A)!0.45!(C)$);
  \coordinate (E) at ($(B)!0.45!(C)$);
  \draw[given] (A) -- (B) -- (C) -- cycle;
  \draw[given] (C) -- (I);
  \draw[given] (A) -- (E);
  \draw[given] (B) -- (D);
  \draw[given] (D) -- (J);
  \draw[given] (E) -- (K);
  \coordinate (F) at (intersection of A--E and B--D);
  \coordinate (G) at (intersection of D--J and A--E);
  \coordinate (H) at (intersection of E--K and B--D);
  \node[vlab, left]  at (A) {$A$};
  \node[vlab, right] at (B) {$B$};
  \node[vlab, above] at (C) {$C$};
  \node[vlab, above left]  at (D) {$D$};
  \node[vlab, above right] at (E) {$E$};
  \node[vlab, below left] at (F) {$F$};
  % G and H sit just above the base between J/I and I/K; "below" laid their
  % letters straight across AB, so they go above their own points instead.
  \node[vlab] at ($(G)+(160:0.30)$) {$G$};
  \node[vlab] at ($(H)+(20:0.30)$) {$H$};
  \node[vlab, below] at (J) {$J$};
  \node[vlab, below] at (I) {$I$};
  \node[vlab, below] at (K) {$K$};
  \VIarc{B}{A}{E}{0.46cm}   \VInumb{B}{A}{E}{0.86cm}{1}
  \VIarc{A}{B}{D}{0.46cm}   \VInumb{A}{B}{D}{0.86cm}{2}
\end{tikzpicture}
\figcap{6--126}
\end{minipage}
\end{bkfigure}}

% ================================================================ Fig 6-127
% Book p.115, Algebra Review 4.  Triangle B C with apex A; D on AB, E on AC,
% and O the crossing of BE and CD.
\newcommand{\FIGVIONEHUNDREDTWENTYSEVEN}{\begin{bkfigure}[\linewidth]
\begin{tikzpicture}[scale=1.196]
  \coordinate (B) at (0,0);
  \coordinate (C) at (2.82,0);
  \coordinate (A) at (1.41,1.92);
  \coordinate (D) at ($(A)!0.46!(B)$);
  \coordinate (E) at ($(A)!0.46!(C)$);
  \draw[given] (B) -- (C) -- (A) -- cycle;
  \draw[given] (B) -- (E);
  \draw[given] (C) -- (D);
  \draw[given] (D) -- (E);
  \coordinate (O) at (intersection of B--E and C--D);
  \node[vlab, left]  at (B) {$B$};
  \node[vlab, right] at (C) {$C$};
  \node[vlab, above] at (A) {$A$};
  \node[vlab, above left]  at (D) {$D$};
  \node[vlab, above right] at (E) {$E$};
  \node[vlab, below] at (O) {$O$};
\end{tikzpicture}
\figcap{6--127}
\end{bkfigure}}

% figures07.tex -- Chapter 7 figures redrawn in TikZ.
% Geometry reconstructed from sources/Geometry.pdf pp.130-146 (book pp.116-132)
% and checked against scans/scan08.pdf pp.11-24 and scans/scan09.pdf pp.1-6.
% Every constrained point is computed, never eyeballed; see tools/constraints/ch07.py.

% ---------------------------------------------------------------- Fig 7-1
\newcommand{\FIGVIIONE}{\begin{bkfigure}[0.60\linewidth]
\begin{tikzpicture}[scale=1.235]
  \coordinate (A) at (0,0); \coordinate (B) at (3.1,0);
  \coordinate (D) at (3.9,0); \coordinate (C) at (1.64,2.1);
  \draw[given] (A) -- (D);
  \draw[given] (A) -- (C) -- (B);
  \foreach \p/\d in {A/below left, B/below, D/below right, C/above}
    {\node[vlab,\d] at (\p) {$\p$};}
\end{tikzpicture}
\figcap{7--1}
\end{bkfigure}}

% ------------------------------------------------------------ Fig 7-2, 7-3
\newcommand{\FIGVIITWOTHREE}{\begin{bkfigure}[\linewidth]
\begin{minipage}{0.48\linewidth}\centering
\begin{tikzpicture}[scale=0.82]
  % BD is chosen congruent to AC (book's own wording in the First Proof),
  % so D sits at distance |AC| beyond B.  |AC| = sqrt(1.10^2+2.10^2).
  \coordinate (A) at (0,0); \coordinate (B) at (2.6,0);
  \coordinate (C) at (1.10,2.10);
  \coordinate (D) at (4.97066,0);
  \draw[given] (A) -- (5.28,0);
  \draw[given] (A) -- (C) -- (B);
  \draw[aux] (C) -- (D);
  % angle ACB at C, angle DBC at B -- the pair the proof assumes congruent
  \draw[tick] ($(C)+(242.35:0.58)$) arc (242.35:305.54:0.58);
  \draw[tick] ($(B)+(0:0.48)$) arc (0:125.54:0.48);
  \foreach \p/\d in {A/below left, B/below, D/below right, C/above}
    {\node[vlab,\d] at (\p) {$\p$};}
\end{tikzpicture}
\figcap{7--2}
\end{minipage}\hfill
\begin{minipage}{0.48\linewidth}\centering
\begin{tikzpicture}[scale=0.82]
  \coordinate (A) at (0,0); \coordinate (B) at (2.4,0);
  \coordinate (D) at (4.29,0); \coordinate (C) at (1.2,2.1);
  \coordinate (F) at (1.05,0);
  \coordinate (E) at ($(C)!1.245!(F)$);
  \draw[given] (A) -- (D);
  \draw[given] (A) -- (C) -- (B);
  \draw[aux] (C) -- (E);
  % angle ECB at C -- the angle the proof makes congruent to angle DBC
  \draw[tick] ($(C)+(265.91:0.72)$) arc (265.91:299.74:0.72);
  \foreach \p/\d in {A/below left, B/below, D/below right, C/above,
                     F/above right, E/below}
    {\node[vlab,\d] at (\p) {$\p$};}
\end{tikzpicture}
\figcap{7--3}
\end{minipage}
\end{bkfigure}}

% ---------------------------------------------------------------- Fig 7-4
\newcommand{\FIGVIIFOUR}{\begin{bkfigure}[0.62\linewidth]
\begin{tikzpicture}[scale=1.235]
  \coordinate (A) at (0,0); \coordinate (B) at (2.6,0);
  \coordinate (D) at (5.05,0); \coordinate (C) at (1.15,2.1);
  \coordinate (G) at ($(C)!0.5!(B)$);
  \coordinate (H) at ($(A)!2!(G)$);
  \draw[given] (A) -- (D);
  \draw[given] (A) -- (C) -- (B);
  \draw[given] (B) -- (H);
  \draw[aux] (A) -- (H);
  \foreach \p/\d in {A/below left, B/below, D/below right, C/above,
                     H/above right}
    {\node[vlab,\d] at (\p) {$\p$};}
  \node[vlab] at (1.957,1.40) {$G$};
\end{tikzpicture}
\figcap{7--4}
\end{bkfigure}}

% ---------------------------------------------------------------- Fig 7-5
\newcommand{\FIGVIIFIVE}{\begin{bkfigure}[0.72\linewidth]
\begin{tikzpicture}[scale=1.235]
  \coordinate (P) at (1.30,1.95);
  \coordinate (Q) at (1.15,0); \coordinate (R) at (1.75,0);
  \draw[given] (0,0) -- (3.2,0);
  \draw[given] (P) -- (Q); \draw[given] (P) -- (R);
  \draw[tick] (1.15,0.16) -- (1.31,0.16) -- (1.31,0);
  \draw[tick] (1.75,0.16) -- (1.59,0.16) -- (1.59,0);
  \node[vlab,above] at (P) {$P$};
  \node[vlab,below] at (Q) {$Q$};
  \node[vlab,below] at (R) {$R$};
  \node[vlab,right] at (3.2,0) {$l$};
\end{tikzpicture}
\figcap{7--5}
\end{bkfigure}}

% ---------------------------------------------------------------- Fig 7-6
\newcommand{\FIGVIISIX}{\begin{bkfigure}[0.58\linewidth]
\begin{tikzpicture}[scale=1.235]
  \coordinate (B) at (0,0); \coordinate (C) at (2.7,0);
  \coordinate (A) at (2.5,1.9);
  % C' on AB with AC' = AC  (t = |AC|/|AB|)
  \coordinate (Cp) at ($(A)!0.60843!(B)$);
  \draw[given] (B) -- (C) -- (A) -- cycle;
  \draw[aux] (Cp) -- (C);
  \node[vlab,above] at (A) {$A$};
  \node[vlab,below left] at (B) {$B$};
  \node[vlab,below right=2pt] at (C) {$C$};
  \node[vlab] at (0.737,1.062) {$C'$};
\end{tikzpicture}
\figcap{7--6}
\end{bkfigure}}

% ------------------------------------------------------------ Fig 7-7, 7-8
\newcommand{\FIGVIISEVENEIGHT}{\begin{bkfigure}[\linewidth]
\begin{minipage}{0.48\linewidth}\centering
\begin{tikzpicture}[scale=1.197]
  % sides 17, 16, 25 to scale (0.1 unit per unit length)
  \coordinate (A) at (0,0); \coordinate (C) at (2.5,0);
  \coordinate (B) at (1.316,1.076);
  \draw[given] (A) -- (C) -- (B) -- cycle;
  \node[vlab,below left] at (A) {$A$};
  \node[vlab,below right] at (C) {$C$};
  \node[vlab,above] at (B) {$B$};
  \node[slab,above left] at ($(A)!0.5!(B)$) {17};
  \node[slab,above right] at ($(B)!0.5!(C)$) {16};
  \node[slab,below] at ($(A)!0.5!(C)$) {25};
\end{tikzpicture}
\figcap{7--7}
\end{minipage}\hfill
\begin{minipage}{0.48\linewidth}\centering
\begin{tikzpicture}[scale=1.197]
  \coordinate (P) at (0,0); \coordinate (Q) at (2.6,0.35);
  \coordinate (R) at (1.10,2.55);
  % QS bisects angle PQR  =>  PS/SR = QP/QR
  \coordinate (S) at ($(P)!0.49629!(R)$);
  \draw[given] (P) -- (Q) -- (R) -- cycle;
  \draw[given] (S) -- (Q);
  % congruence ticks across the two halves of angle PQR
  \draw[tick] ([shift=(65.95:0.09)]$(Q)!0.30!(S)$)
            -- ([shift=(245.95:0.09)]$(Q)!0.30!(S)$);
  \draw[tick] ([shift=(34.30:0.09)]$(Q)!0.30!(R)$)
            -- ([shift=(214.30:0.09)]$(Q)!0.30!(R)$);
  \node[vlab,below left] at (P) {$P$};
  \node[vlab,right] at (Q) {$Q$};
  \node[vlab,above] at (R) {$R$};
  \node[vlab,above left] at (S) {$S$};
\end{tikzpicture}
\figcap{7--8}
\end{minipage}
\end{bkfigure}}

% ----------------------------------------------------------- Fig 7-9, 7-10
\newcommand{\FIGVIININETEN}{\begin{bkfigure}[\linewidth]
\begin{minipage}{0.46\linewidth}\centering
\begin{tikzpicture}[scale=1.175]
  \coordinate (A) at (0,0); \coordinate (C) at (2.9,0);
  \coordinate (B) at (1.0,1.62);
  % B' on AC with AB' = AB   (|AB| = sqrt(1.0^2 + 1.62^2))
  \coordinate (Bp) at (1.903786,0);
  \draw[given] (A) -- (C) -- (B) -- cycle;
  \draw[aux] (B) -- (Bp);
  % the book's two angle arcs: one at B crossed by BB' (angles 1, 4),
  % one at B' crossed by B'B (angles 2, 3)
  \draw[tick] ($(B)+(238.33:0.42)$) arc (238.33:319.54:0.42);
  \draw[tick] ($(Bp)+(0:0.20)$) arc (0:180:0.20);
  \node[vlab,below left] at (A) {$A$};
  \node[vlab,below right] at (C) {$C$};
  \node[vlab,above] at (B) {$B$};
  \node[vlab,below] at (Bp) {$B'$};
  \node[slab] at (0.986,1.000) {1};   % angle ABB'
  \node[slab] at (1.665,0.808) {4};   % angle B'BC
  \node[slab] at (1.429,0.278) {2};   % angle AB'B
  \node[slab] at (2.210,0.300) {3};   % angle BB'C
\end{tikzpicture}
\figcap{7--9}
\end{minipage}\hfill
\begin{minipage}{0.50\linewidth}\centering
\begin{tikzpicture}[scale=0.729]
  % vertex ratios measured off scan08 p.15 (grid overlay); three interior
  % vertices are unlabelled in the book, as here.
  \coordinate (Pone) at (0.73,1.58);
  \coordinate (Ptwo) at (1.61,4.41);
  \coordinate (Pthree) at (1.73,1.61);
  \coordinate (U) at (0.12,2.99);
  \coordinate (R) at (2.78,2.16);
  \coordinate (S) at (2.38,0.83);
  \coordinate (Pn1) at (0.34,0.45);
  \coordinate (Pn) at (2.60,0.10);
  \draw[given] (Pone) -- (Ptwo) -- (Pthree) -- (U) -- (R) -- (S)
               -- (Pn1) -- (Pn);
  \draw[aux] (Pone) -- (Pn);
  \node[vlab,left] at (Pone) {$P_1$};
  \node[vlab,above] at (Ptwo) {$P_2$};
  \node[vlab,right] at (Pthree) {$P_3$};
  \node[vlab,left] at (Pn1) {$P_{n-1}$};
  \node[vlab,right] at (Pn) {$P_n$};
\end{tikzpicture}
\figcap{7--10}
\end{minipage}
\end{bkfigure}}

% --------------------------------------------------------------- Fig 7-11
\newcommand{\FIGVIIELEVEN}{\begin{bkfigure}[0.56\linewidth]
\begin{tikzpicture}[scale=1.300]
  \coordinate (B) at (0,0); \coordinate (C) at (2.6,0);
  \coordinate (A) at (1.85,1.35);
  % D on BA with DB = DC  (D lies on the perpendicular bisector x = 1.3)
  \coordinate (D) at (1.3,0.94865);
  \draw[given] (B) -- (C) -- (A) -- cycle;
  \draw[given] (D) -- (C);
  % DB = DC, so the book marks angle DBC and angle DCB congruent
  % with a DOUBLE arc at each vertex.
  \draw[tick] (0.30,0) arc (0:36.13:0.30);
  \draw[tick] (0.40,0) arc (0:36.13:0.40);
  \draw[tick] ($(C)+(180:0.30)$) arc (180:143.87:0.30);
  \draw[tick] ($(C)+(180:0.40)$) arc (180:143.87:0.40);
  \node[vlab,below left] at (B) {$B$};
  \node[vlab,below right] at (C) {$C$};
  \node[vlab,above] at (A) {$A$};
  \node[vlab,above left] at (D) {$D$};
\end{tikzpicture}
\figcap{7--11}
\end{bkfigure}}

% --------------------------------------------------------------- Fig 7-12
\newcommand{\FIGVIITWELVE}{\begin{bkfigure}[0.62\linewidth]
\begin{tikzpicture}[scale=1.300]
  % C is on the perpendicular bisector of BB', so CB' = CB exactly
  % (the hypothesis BC = B'C' with C' = C).  Q on AB with QB = QB'.
  \coordinate (A) at (0,0); \coordinate (B) at (3.2,0);
  \coordinate (C) at (1.8875,1.43);
  \coordinate (Bp) at (1.30,-0.42);     % CB' = CB
  \coordinate (Q) at (2.2036,0);        % QB = QB'
  \draw[given] (A) -- (B);
  \draw[given] (A) -- (C) -- (B);
  \draw[given] (C) -- (Bp);
  \draw[aux] (Bp) -- (Q);
  \node[vlab,above] at (C) {$C'=C$};
  \node[vlab,left] at (A) {$A'=A$};
  \node[vlab,left] at (Bp) {$B'$};
  \node[vlab,above right] at (Q) {$Q$};
  \node[vlab,right] at (B) {$B$};
\end{tikzpicture}
\figcap{7--12}
\end{bkfigure}}

% ------------------------------------------------------- Fig 7-13, 7-14
\newcommand{\FIGVIITHIRTEENFOURTEEN}{\begin{bkfigure}[\linewidth]
\begin{minipage}{0.48\linewidth}\centering
\begin{tikzpicture}[scale=1.145]
  \coordinate (A) at (0,0); \coordinate (C) at (3.0,0);
  \coordinate (B) at (1.45,1.93);
  \coordinate (E) at ($(B)!0.49!(C)$);
  \coordinate (D) at ($(A)!0.655!(E)$);
  \draw[given] (A) -- (C) -- (B) -- cycle;
  \draw[given] (A) -- (D);
  \draw[aux] (D) -- (E);
  \draw[given] (D) -- (C);
  \node[vlab,below left] at (A) {$A$};
  \node[vlab,below right] at (C) {$C$};
  \node[vlab,above] at (B) {$B$};
  \node[vlab,above] at (D) {$D$};
  \node[vlab,above right] at (E) {$E$};
\end{tikzpicture}
\figcap{7--13}
\end{minipage}\hfill
\begin{minipage}{0.48\linewidth}\centering
\begin{tikzpicture}[scale=1.145]
  \draw[given] (1.2,2.03) -- (1.2,-0.10);
  \draw[given] (0,1.6) -- (2.8,1.6);
  \draw[given] (0,0.2) -- (2.8,0.2);
  % both right-angle boxes sit ABOVE the horizontal, right of m (as printed)
  \draw[tick] (1.2,1.76) -- (1.36,1.76) -- (1.36,1.6);
  \draw[tick] (1.2,0.36) -- (1.36,0.36) -- (1.36,0.2);
  \node[vlab,above] at (1.2,2.03) {$m$};
  \node[vlab,right] at (2.8,1.6) {$l$};
  \node[vlab,right] at (2.8,0.2) {$l'$};
\end{tikzpicture}
\figcap{7--14}
\end{minipage}
\end{bkfigure}}

% --------------------------------------------------------------- Fig 7-15
\newcommand{\FIGVIIFIFTEEN}{\begin{bkfigure}[\linewidth]
\begin{minipage}{0.48\linewidth}\centering
\begin{tikzpicture}[scale=1.184]
  \def\rr{1.15}
  \coordinate (O) at (0,0);
  \draw[given] (O) circle (\rr);
  \coordinate (A) at (40:\rr);          % on the circle
  \coordinate (Bi) at (95:0.80);        % interior: OB < r
  \coordinate (Cx) at (-10:1.87);       % exterior: OC > r
  \coordinate (E) at (215:\rr);
  \coordinate (F) at (293:\rr);
  \draw[given] (O) -- (A);
  \draw[given] (O) -- (Bi);
  \draw[given] (O) -- (Cx);
  \draw[given] (E) -- (F);            % the chord EF
  \draw[aux] (E) -- ($(F)!-0.30!(E)$);  % secant, produced both ways
  \draw[aux] (F) -- ($(E)!-0.30!(F)$);
  \node[vlab,above right] at (A) {$A$};
  \node[vlab,left] at (Bi) {$B$};
  \node[vlab,right] at (Cx) {$C$};
  \node[vlab,below left] at (E) {$E$};
  \node[vlab,below] at (F) {$F$};
  \node[vlab,below left] at (O) {$O$};
\end{tikzpicture}
\end{minipage}\hfill
\begin{minipage}{0.48\linewidth}\centering
\begin{tikzpicture}[scale=1.184]
  \def\rr{1.15}
  \coordinate (O) at (0,0);
  \draw[given] (O) circle (\rr);
  \coordinate (A) at (95:\rr);
  \coordinate (B) at (45:\rr);
  \coordinate (F) at (178:\rr);
  \coordinate (E) at (250:\rr);
  \coordinate (G) at (310:\rr);
  \draw[given] (O) -- (A);            % central angle AOB
  \draw[given] (O) -- (B);
  % inscribed angle EFG: only the two chords FE and FG are drawn.
  % (Checked on scan08 p.18 -- the book does NOT close E--G.)
  \draw[given] (F) -- (E);
  \draw[given] (F) -- (G);
  \node[vlab,above] at (A) {$A$};
  \node[vlab] at (38:1.48) {$B$};
  \node[vlab,left] at (F) {$F$};
  \node[vlab,below] at (E) {$E$};
  \node[vlab,below right] at (G) {$G$};
  \node[vlab,below right] at (O) {$O$};
\end{tikzpicture}
\end{minipage}
\figcap{7--15}
\end{bkfigure}}

% --------------------------------------------------------------- Fig 7-16
\newcommand{\FIGVIISIXTEEN}{\begin{bkfigure}[0.56\linewidth]
\begin{tikzpicture}[scale=1.365]
  % Ex. 7-5 #3 asks what kind of angle each of /A, /B, /C, /D, /E is, so
  % every labelled point must carry two drawn rays.  C is the CENTRE
  % (printed "C = O"), and the two radii CD, CP_3 form the central angle.
  % Chord set read off scan08 p.20; P_1..P_4 are the book's unlabelled
  % points where the other chords meet the circle.
  \def\rr{1.3}
  \coordinate (O) at (0,0);
  \draw[given] (O) circle (\rr);
  \coordinate (A) at (57:\rr);
  \coordinate (B) at (132:\rr);
  \coordinate (D) at (240:\rr);
  \coordinate (E) at (308:\rr);
  \coordinate (Pone) at (-4:\rr);
  \coordinate (Ptwo) at (13:\rr);
  \coordinate (Pthree) at (-13:\rr);
  \coordinate (Pfour) at (-34:\rr);
  \draw[given] (B) -- (D);                       % chord, side of /B and /D
  \draw[given] (B) -- (Pone);                    % chord, other side of /B
  \draw[given] (A) -- (Ptwo);                    % chord, one side of /A
  \draw[given] (A) -- ($(A)+(-19.5:1.05)$);      % ray outside, other side of /A
  \draw[given] (O) -- (D);                       % radius
  \draw[given] (O) -- (Pthree);                  % radius
  \draw[given] (D) -- (Pfour);                   % chord
  \draw[given] (D) -- (E);                       % chord
  \draw[given] (E) -- ($(E)+(9.9:1.20)$);        % the two rays at E, both
  \draw[given] (E) -- ($(E)+(-11.5:1.20)$);      % lying outside the circle
  \node[vlab,above left] at (B) {$B$};
  \node[vlab,above right] at (A) {$A$};
  \node[vlab,below left] at (D) {$D$};
  \node[vlab,below] at (E) {$E$};
  \node[slab] at (-0.30,0.22) {$C=O$};
\end{tikzpicture}
\figcap{7--16}
\end{bkfigure}}

% --------------------------------------------------------------- Fig 7-17
% Line drawing of a drafting compass, brace beneath marking the radius r.
\newcommand{\FIGVIISEVENTEEN}{\begin{bkfigure}[0.40\linewidth]
\begin{tikzpicture}[scale=1.300]
  % handle
  \draw[given] (-0.075,2.62) rectangle (0.075,3.05);
  % knurled collar
  \draw[given] (-0.115,2.30) -- (-0.115,2.62) -- (0.115,2.62) -- (0.115,2.30);
  \draw[given] (-0.115,2.30) arc (180:360:0.115);
  % hinge head
  \draw[given] (-0.16,1.92) .. controls (-0.20,2.16) and (-0.12,2.28) ..
               (0,2.30) .. controls (0.12,2.28) and (0.20,2.16) .. (0.16,1.92);
  % legs
  \draw[given] (-0.13,2.02) -- (-0.58,0.30);
  \draw[given] (0.13,2.02) -- (0.58,0.30);
  \draw[given] (-0.055,2.00) -- (-0.50,0.30);
  \draw[given] (0.055,2.00) -- (0.50,0.30);
  % needle point (left) and pen nib (right)
  \draw[given] (-0.58,0.30) -- (-0.605,0.10) -- (-0.545,0.10) -- (-0.50,0.30);
  \draw[given] (-0.605,0.10) -- (-0.62,0.0);
  \draw[given] (0.50,0.30) -- (0.545,0.06) -- (0.605,0.06) -- (0.58,0.30);
  \draw[given] (0.575,0.06) -- (0.62,0.0);
  % brace showing the radius
  \draw[tick] (-0.62,-0.16) -- (-0.62,-0.24) -- (-0.04,-0.24);
  \draw[tick] (0.0,-0.30) -- (0.04,-0.24) -- (0.62,-0.24) -- (0.62,-0.16);
  \node[vlab,below] at (0,-0.34) {$r$};
\end{tikzpicture}
\figcap{7--17}
\end{bkfigure}}

% ------------------------------------------------------- Fig 7-18, 7-19
\newcommand{\FIGVIIEIGHTEENNINETEEN}{\begin{bkfigure}[\linewidth]
\begin{minipage}{0.48\linewidth}\centering
\begin{tikzpicture}[scale=1.176]
  \coordinate (V) at (0,0);
  \foreach \a in {45,18,-9,-36}{\draw[given] (V) -- (\a:2.35);}
  \draw[tick] (-36:1.75) arc (-36:45:1.75);
  \node[slab] at (31.5:1.20) {$B$};
  \node[slab] at (4.5:1.20) {$C$};
  \node[slab] at (-22.5:1.20) {$D$};
  \node[slab] at (2.020,0.324) {$A$};
\end{tikzpicture}
\figcap{7--18}
\end{minipage}\hfill
\begin{minipage}{0.48\linewidth}\centering
\begin{tikzpicture}[scale=1.176]
  \draw[given] (0,0) circle (0.85);
  \draw[given] (1.6,-0.85) rectangle (3.3,0.85);
\end{tikzpicture}
\figcap{7--19}
\end{minipage}
\end{bkfigure}}

% ------------------------------------------------------- Fig 7-20, 7-21
\newcommand{\FIGVIITWENTYTWENTYONE}{\begin{bkfigure}[\linewidth]
\begin{minipage}{0.48\linewidth}\centering
\begin{tikzpicture}[scale=0.998]
  % angle A
  \coordinate (A) at (0,1.75);
  \draw[given] (A) -- ++(0:3.0);
  \draw[given] (A) -- ++(14:3.0);
  \node[vlab,left] at (A) {$A$};
  % ray from B, with the two possible copied angles dotted in
  \coordinate (B) at (0,0);
  \draw[given] (B) -- ++(0:3.0);
  \draw[aux] (B) -- ++(20:2.5);
  \draw[aux] (B) -- ++(-20:2.5);
  \node[vlab,left] at (B) {$B$};
  \node[slab,right] at (3.0,0) {Ray};
\end{tikzpicture}
\figcap{7--20}
\end{minipage}\hfill
\begin{minipage}{0.48\linewidth}\centering
\begin{tikzpicture}[scale=0.998]
  % copied angle at A
  \coordinate (A) at (0,1.75);
  \coordinate (E) at ($(A)+(0:1.0)$);
  \coordinate (D) at ($(A)+(38:1.0)$);
  \draw[given] (A) -- ++(0:2.9);
  \draw[given] (A) -- ++(38:1.85);
  % the compass arc of radius r about A, cutting both sides (at D and E)
  \draw[aux] ($(A)+(-8:1.0)$) arc (-8:50:1.0);
  \node[vlab,left] at (A) {$A$};
  \node[vlab] at ($(A)+(50:1.35)$) {$D$};
  \node[vlab,below right] at (E) {$E$};
  % construction at B
  \coordinate (B) at (0,0);
  \coordinate (F) at ($(B)+(0:1.0)$);
  \coordinate (G) at ($(B)+(38:1.0)$);
  \draw[given] (B) -- ++(0:2.9);
  \draw[given] (B) -- ($(B)+(38:1.85)$);
  % arc of radius r about B, cutting the ray at F;
  % then the arc of radius DE about F, which meets the first at G.
  \draw[aux] ($(B)+(-8:1.0)$) arc (-8:50:1.0);
  \draw[aux] ($(F)+(96:0.6512)$) arc (96:122:0.6512);
  \node[vlab,left] at (B) {$B$};
  \node[vlab,below] at (F) {$F$};
  \node[vlab] at ($(B)+(50:1.35)$) {$G$};
  \node[vlab,above right] at ($(B)+(38:1.85)$) {$H$};
  \node[slab,right] at (2.9,0) {Ray};
\end{tikzpicture}
\figcap{7--21}
\end{minipage}
\end{bkfigure}}

% ------------------------------------------------------- Fig 7-22, 7-23
\newcommand{\FIGVIITWENTYTWOTWENTYTHREE}{\begin{bkfigure}[\linewidth]
\begin{minipage}{0.48\linewidth}\centering
\begin{tikzpicture}[scale=0.983]
  % given triangle
  \coordinate (A) at (0,2.05); \coordinate (B) at (0.95,3.40);
  \coordinate (C) at (2.35,2.05);
  \draw[given] (A) -- (B) -- (C) -- cycle;
  \node[vlab,left] at (A) {$A$};
  \node[vlab,above] at (B) {$B$};
  \node[vlab,right] at (C) {$C$};
  % copy on a ray
  \coordinate (Ap) at (0,0); \coordinate (Bp) at (0.95,1.35);
  \coordinate (Cp) at (2.35,0);   % A'C' = AC, A'B' = AB (congruent copy)
  \draw[given] (Ap) -- (3.05,0);
  % Construction 7-2 stops here: the copied side A'B', the ray, C' cut off
  % on it, and the arc of radius AB about A' that fixes B'.  The book does
  % NOT close B'C' (scan08 p.24).
  \draw[given] (Ap) -- (Bp);
  \draw[aux] ($(Ap)+(44.9:1.65076)$) arc (44.9:64.9:1.65076);
  \draw[tick] ($(Cp)+(0,-0.13)$) -- ($(Cp)+(0,0.13)$);
  \node[vlab,left] at (Ap) {$A'$};
  \node[vlab] at ($(Bp)+(0.10,0.36)$) {$B'$};
  \node[vlab] at ($(Cp)+(0.26,0.34)$) {$C'$};
\end{tikzpicture}
\figcap{7--22}
\end{minipage}\hfill
\begin{minipage}{0.48\linewidth}\centering
\begin{tikzpicture}[scale=0.983]
  % given triangle
  \coordinate (A) at (0,2.15); \coordinate (B) at (1.15,3.45);
  \coordinate (C) at (2.4,2.15);
  \draw[given] (A) -- (B) -- (C) -- cycle;
  \node[vlab,left] at (A) {$A$};
  \node[vlab,above] at (B) {$B$};
  \node[vlab,right] at (C) {$C$};
  % ray with the two arcs crossing at B'
  \coordinate (Ap) at (0,0); \coordinate (Cp) at (2.4,0);
  \coordinate (Bp) at (1.15,1.30);
  \draw[given] (Ap) -- (3.1,0);
  \dt{Ap}
  \draw[tick] ($(Cp)+(0,-0.13)$) -- ($(Cp)+(0,0.13)$);
  % the two construction circles: radius AB about A', radius BC about C'.
  % They cross at B'.
  \draw[aux] ($(Ap)+(38.5:1.73566)$) arc (38.5:58.5:1.73566);
  \draw[aux] ($(Cp)+(123.9:1.80347)$) arc (123.9:143.9:1.80347);
  \node[vlab,left] at (Ap) {$A'$};
  \node[vlab] at ($(Bp)+(0,0.44)$) {$B'$};
  \node[vlab] at ($(Cp)+(0.26,0.34)$) {$C'$};
  \node[slab,right] at (3.1,0) {Ray};
\end{tikzpicture}
\figcap{7--23}
\end{minipage}
\end{bkfigure}}

% ------------------------------------------------------- Fig 7-24, 7-25
\newcommand{\FIGVIITWENTYFOURTWENTYFIVE}{\begin{bkfigure}[\linewidth]
\begin{minipage}{0.48\linewidth}\centering
\begin{tikzpicture}[scale=0.993]
  \coordinate (A) at (0,0); \coordinate (B) at (3.0,0);
  \coordinate (P) at (1.5,1.0828); \coordinate (Q) at (1.5,-1.0828);
  \draw[given] (A) -- (B);
  \dt{A} \dt{B}
  \draw[aux] ([shift=(-50:1.85)]A) arc (-50:50:1.85);
  \draw[aux] ([shift=(130:1.85)]B) arc (130:230:1.85);
  \node[vlab,left] at (A) {$A$};
  \node[vlab,right] at (B) {$B$};
  \node[vlab] at (2.35,1.30) {$P$};
  \node[vlab] at (2.35,-1.30) {$Q$};
\end{tikzpicture}
\figcap{7--24}
\end{minipage}\hfill
\begin{minipage}{0.48\linewidth}\centering
\begin{tikzpicture}[scale=0.993]
  \coordinate (A) at (0,0); \coordinate (B) at (3.0,0);
  \coordinate (P) at (1.5,1.0828); \coordinate (Q) at (1.5,-1.0828);
  \coordinate (C) at (1.5,0);
  \draw[given] (A) -- (B);
  \dt{A} \dt{B}
  \draw[aux] ([shift=(-50:1.85)]A) arc (-50:50:1.85);
  \draw[aux] ([shift=(130:1.85)]B) arc (130:230:1.85);
  \draw[given] ($(P)!-0.16!(Q)$) -- ($(Q)!-0.16!(P)$);
  \node[vlab,left] at (A) {$A$};
  \node[vlab,right] at (B) {$B$};
  \node[vlab] at (2.35,1.30) {$P$};
  \node[vlab] at (2.35,-1.30) {$Q$};
  \node[vlab] at (2.05,0.72) {$C$};
\end{tikzpicture}
\figcap{7--25}
\end{minipage}
\end{bkfigure}}

% -------------------------------------------------- Fig 7-26, 7-27, 7-28
\newcommand{\FIGVIITWENTYSIXSEVENEIGHT}{\begin{bkfigure}[\linewidth]
\begin{minipage}{0.30\linewidth}\centering
\begin{tikzpicture}[scale=0.80]
  \draw[given] (0,0) -- (3.0,0);
  \coordinate (A) at (0.55,0); \coordinate (P) at (1.5,0);
  \coordinate (B) at (2.45,0);
  \draw[tick] (0.55,-0.16) -- (0.55,0.16);
  \draw[tick] (2.45,-0.16) -- (2.45,0.16);
  \dt{P}
  \node[vlab,above left] at (A) {$A$};
  \node[vlab,above] at (P) {$P$};
  \node[vlab,above right] at (B) {$B$};
\end{tikzpicture}
\figcap{7--26}
\end{minipage}\hfill
\begin{minipage}{0.30\linewidth}\centering
\begin{tikzpicture}[scale=0.80]
  \draw[given] (0,0) -- (3.2,0);
  \coordinate (A) at (0.35,0); \coordinate (B) at (2.85,0);
  \coordinate (P) at (1.6,1.75);
  \draw[tick] (0.35,-0.16) -- (0.35,0.16);
  \draw[tick] (2.85,-0.16) -- (2.85,0.16);
  \draw[aux] (P) -- (1.6,0);
  \node[vlab,above] at (P) {$P$};
  \node[vlab,above right] at (A) {$A$};
  \node[vlab,above left] at (B) {$B$};
\end{tikzpicture}
\figcap{7--27}
\end{minipage}\hfill
\begin{minipage}{0.34\linewidth}\centering
\begin{tikzpicture}[scale=0.80]
  \coordinate (A) at (0,0);
  \coordinate (P) at (1.55,1.8); \coordinate (Pp) at (1.55,-1.8);
  \coordinate (D) at (1.55,0);
  \coordinate (B) at (2.45,0); \coordinate (C) at (2.9,0);
  \draw[given] (-0.4,0) -- (3.4,0);
  \draw[given] (A) -- (P);
  \draw[given] (A) -- ($(A)!1.22!(Pp)$);   % the ray AQ, through P'
  \draw[given] (P) -- (Pp);
  % angle PAC congruent to angle QAC -- the book marks both
  \draw[tick] (0.42,0) arc (0:49.28:0.42);
  \draw[tick] (0.42,0) arc (0:-49.28:0.42);
  \dt{B} \dt{C}
  \node[vlab] at (-0.24,0.46) {$A$};
  \node[vlab,above] at (P) {$P$};
  \node[vlab,below left] at (Pp) {$P'$};
  \node[vlab] at (1.80,0.30) {$D$};
  \node[vlab,above] at (B) {$B$};
  \node[vlab,above] at (C) {$C$};
  \node[vlab,right] at (3.4,0) {$l$};
  \node[vlab,below right] at ($(A)!1.22!(Pp)$) {$Q$};
\end{tikzpicture}
\figcap{7--28}
\end{minipage}
\end{bkfigure}}

% ------------------------------------------------------- Fig 7-29, 7-30
\newcommand{\FIGVIITWENTYNINETHIRTY}{\begin{bkfigure}[\linewidth]
\begin{minipage}{0.48\linewidth}\centering
\begin{tikzpicture}[scale=1.074]
  \coordinate (A) at (0,0);
  \coordinate (P) at (1.6,1.6); \coordinate (Pp) at (1.6,-1.6);
  \draw[given] (-0.3,0) -- (3.4,0);
  \draw[given] (A) -- (P); \draw[given] (A) -- (Pp);
  \draw[aux] (P) -- (Pp);
  \node[vlab] at (-0.18,0.26) {$A$};
  \node[vlab,above] at (P) {$P$};
  \node[vlab,below] at (Pp) {$P'$};
  \node[vlab,right] at (3.4,0) {$l$};
  \node[slab] at (0.66,0.34) {1};
  \node[slab] at (0.66,-0.36) {2};
\end{tikzpicture}
\figcap{7--29}
\end{minipage}\hfill
\begin{minipage}{0.48\linewidth}\centering
\begin{tikzpicture}[scale=1.124]
  \coordinate (A) at (0,2.0);
  \coordinate (C) at (0,1.0);
  \coordinate (B) at (0.643,1.234);
  \coordinate (Q) at (0.5303,0.5435);
  \draw[given] (A) -- (0,0);
  \draw[given] (A) -- (1.672,0.008);
  \draw[tick] (-0.13,1.0) -- (0.13,1.0);
  \draw[tick] ($(B)+(-50:0.13)$) -- ($(B)+(130:0.13)$);
  \draw[aux] ($(Q)+(150:0.26)$) arc (150:205:0.26);
  \draw[aux] ($(Q)+(-25:0.26)$) arc (-25:30:0.26);
  \dt{Q}
  \node[vlab,above] at (A) {$A$};
  \node[vlab,left] at (C) {$C$};
  \node[vlab,above right] at (B) {$B$};
  \node[vlab,below] at (Q) {$Q$};
\end{tikzpicture}
\figcap{7--30}
\end{minipage}
\end{bkfigure}}

% ------------------------------------------------------- Fig 7-31, 7-32
\newcommand{\FIGVIITHIRTYONETHIRTYTWO}{\begin{bkfigure}[\linewidth]
\begin{minipage}{0.48\linewidth}\centering
\begin{tikzpicture}[scale=1.171]
  \coordinate (A) at (0,0);
  \coordinate (C) at (1.5,0);          % r  on the lower side
  \coordinate (Cp) at (2.1,0);         % r' on the lower side
  \coordinate (D) at (30:1.5);         % r  on the upper side
  \coordinate (Dp) at (30:2.1);        % r' on the upper side
  \coordinate (P) at (1.633,0.4373);   % DC' meets D'C on the bisector
  \draw[given] (A) -- (0:2.55);
  \draw[given] (A) -- (30:2.55);
  % the two compass arcs, radius r and r', each cutting BOTH sides of the
  % angle -- at C, D and at C', D' respectively.
  \draw[aux] ([shift=(0:1.5)]A) arc (0:30:1.5);
  \draw[aux] ([shift=(0:2.1)]A) arc (0:30:2.1);
  \draw[aux] (D) -- (Cp);
  \draw[aux] (Dp) -- (C);
  \dt{P}
  \node[vlab,left] at (A) {$A$};
  \node[vlab] at (1.149,1.010) {$D$};
  \node[vlab] at (1.669,1.310) {$D'$};
  \node[vlab,below] at (C) {$C$};
  \node[vlab,below] at (Cp) {$C'$};
  \node[vlab] at (1.227,0.329) {$P$};
\end{tikzpicture}
\figcap{7--31}
\end{minipage}\hfill
\begin{minipage}{0.48\linewidth}\centering
\begin{tikzpicture}[scale=1.076]
  % upper: angle CAD bisected
  \coordinate (A) at (0,1.35);
  \draw[given] (A) -- ++(30:2.7);
  \draw[given] (A) -- ++(15:2.7);
  \draw[given] (A) -- ++(0:2.7);
  \node[vlab,left] at (A) {$A$};
  \node[vlab,above right] at ($(A)+(30:2.7)$) {$C$};
  \node[vlab,below right] at ($(A)+(0:2.7)$) {$D$};
  \node[slab] at ($(A)+(22.5:1.90)$) {1};
  \node[slab] at ($(A)+(7.5:1.90)$) {2};
  % lower: angle EBF bisected
  \coordinate (B) at (0,-0.55);
  \draw[given] (B) -- ++(30:2.7);
  \draw[given] (B) -- ++(15:2.7);
  \draw[given] (B) -- ++(0:2.7);
  \node[vlab,left] at (B) {$B$};
  \node[vlab,above right] at ($(B)+(30:2.7)$) {$E$};
  \node[vlab,below right] at ($(B)+(0:2.7)$) {$F$};
  \node[slab] at ($(B)+(22.5:1.90)$) {3};
  \node[slab] at ($(B)+(7.5:1.90)$) {4};
\end{tikzpicture}
\figcap{7--32}
\end{minipage}
\end{bkfigure}}

% ------------------------------------------------------- Fig 7-33, 7-34
\newcommand{\FIGVIITHIRTYTHREETHIRTYFOUR}{\begin{bkfigure}[\linewidth]
\begin{minipage}{0.48\linewidth}\centering
\begin{tikzpicture}[scale=1.000]
  \coordinate (B) at (0,0); \coordinate (A) at (0,1.85);
  \coordinate (C) at (1.25,0); \coordinate (D) at (2.85,0);
  \draw[given] (B) -- (D);
  \draw[given] (A) -- (B);
  \draw[given] (A) -- (C);
  \draw[given] (A) -- (D);
  \draw[tick] (0,0.20) -- (0.20,0.20) -- (0.20,0);
  \node[vlab,above] at (A) {$A$};
  \node[vlab,below left] at (B) {$B$};
  \node[vlab,below] at (C) {$C$};
  \node[vlab,below right] at (D) {$D$};
\end{tikzpicture}
\figcap{7--33}
\end{minipage}\hfill
\begin{minipage}{0.48\linewidth}\centering
\begin{tikzpicture}[scale=1.305]
  \coordinate (A) at (0,1.35); \coordinate (B) at (2.9,1.1);
  \coordinate (C) at (0.1,0);  \coordinate (D) at (3.0,0.35);
  \coordinate (O) at (1.913,0.7123);
  \draw[given] (A) -- (B);
  \draw[given] (C) -- (D);
  \draw[given] (A) -- (D);
  \draw[given] (C) -- (B);
  \node[vlab,above left] at (A) {$A$};
  \node[vlab,above right] at (B) {$B$};
  \node[vlab,below left] at (C) {$C$};
  \node[vlab,right] at (D) {$D$};
  \node[vlab] at (1.99,0.48) {$O$};
\end{tikzpicture}
\figcap{7--34}
\end{minipage}
\end{bkfigure}}

% ------------------------------------------------------- Fig 7-35, 7-36
\newcommand{\FIGVIITHIRTYFIVETHIRTYSIX}{\begin{bkfigure}[\linewidth]
\begin{minipage}{0.48\linewidth}\centering
\begin{tikzpicture}[scale=1.026]
  \coordinate (A) at (0,0);
  % B, A, F are collinear in the printed figure (158 and -22 degrees), which
  % is what makes Review Ex. 18's conclusion true: AC bisects angle DAB and
  % AE bisects angle DAF, hence angle CAE is right.
  \draw[given] (A) -- (158:1.70);     % B
  \draw[given] (A) -- (100.5:1.75);   % C
  \draw[given] (A) -- (43:1.75);      % D
  \draw[given] (A) -- (10.5:1.75);    % E
  \draw[given] (A) -- (-22:1.75);     % F
  \draw[aux] (A) -- (190.5:0.95);     % extension of AE (the exercise's hint)
  \node[vlab] at (0.10,-0.30) {$A$};
  \node[vlab,left] at (158:1.70) {$B$};
  \node[vlab,above] at (100.5:1.75) {$C$};
  \node[vlab,above right] at (43:1.75) {$D$};
  \node[vlab,right] at (10.5:1.75) {$E$};
  \node[vlab,right] at (-22:1.75) {$F$};
\end{tikzpicture}
\figcap{7--35}
\end{minipage}\hfill
\begin{minipage}{0.48\linewidth}\centering
\begin{tikzpicture}[scale=1.026]
  \coordinate (A) at (0,0);
  % AC bisects angle DAB ; angle CAE is right
  \draw[given] (A) -- (180:1.55);     % B
  \draw[given] (A) -- (127.5:1.75);   % C
  \draw[given] (A) -- (75:1.75);      % D
  \draw[given] (A) -- (37.5:1.75);    % E
  \draw[given] (A) -- (0:1.55);       % F
  \node[vlab] at (0,-0.32) {$A$};
  \node[vlab,left] at (180:1.55) {$B$};
  \node[vlab,above] at (127.5:1.75) {$C$};
  \node[vlab,above] at (75:1.75) {$D$};
  \node[vlab,right] at (37.5:1.75) {$E$};
  \node[vlab,right] at (0:1.55) {$F$};
\end{tikzpicture}
\figcap{7--36}
\end{minipage}
\end{bkfigure}}

% figures08.tex -- Chapter 8 (Parallel Lines) figures redrawn in TikZ.
%
% Local label clearance: ch08 has several figures where a vertex label sits
% at a point that two or three drawn lines pass through, so the shared
% 3.4pt outer sep is not enough.  Raised here only (brumfiel.sty untouched).
\tikzset{
  vlab/.style = {font=\footnotesize\itshape, inner sep=1.2pt, outer sep=5.2pt},
  slab/.style = {font=\scriptsize, inner sep=1.2pt, outer sep=5.2pt},
}
% Geometry read off the photo pages (book pp.133-153 = PDF 147-167) and the
% iPad scans (scan09 pp.7-23 = book pp.133-148; scan10 pp.1-4 = book
% pp.151-153).  Vertex ratios were measured off 400 dpi crops of the source
% figures; constrained points are constructed, never eyeballed: parallels via
% equal ratios, feet via projection, crossings via (intersection of ...).

% ---------------------------------------------------------------- Fig 8-1
% Two rays with the same direction: r1 || r2, both on the same side of the
% line through their vertices (Definition 8-2).
\newcommand{\FIGVIIIONE}{\begin{bkfigure}[0.62\linewidth]
\begin{tikzpicture}[scale=0.9]
  \coordinate (V1) at (0.62,1.15);          % vertex of r1
  \coordinate (V2) at (0.16,0.0);           % vertex of r2
  \draw[aux] ($(V2)!-0.30!(V1)$) -- ($(V2)!1.55!(V1)$);   % line of vertices
  \draw[fig,-latex] (V1) -- ++(4.35,0);
  \draw[fig,-latex] (V2) -- ++(4.35,0);
  \node[slab,above] at ($(V1)+(2.5,0)$) {$r_1$};
  \node[slab,above] at ($(V2)+(2.5,0)$) {$r_2$};
\end{tikzpicture}
\figcap{8--1}
\end{bkfigure}}

% ----------------------------------------------------------- Fig 8-2, 8-3
% 8-2: first proof of Thm 8-1.  n = AB; C on n beyond A; D on m and E on l
% chosen so that angle DAC is congruent to angle EBA, hence m || l, and the
% lines cannot meet at P.  Both congruent angles are arced, as in the book.
% 8-3: the lemma's figure -- the same three lines with m NOT parallel to l,
% and the two corresponding angles numbered 1 and 2.
\newcommand{\FIGVIIITWOTHREE}{\begin{bkfigure}[\linewidth]
\begin{minipage}{0.48\linewidth}\centering
\begin{tikzpicture}[scale=0.88]
  \coordinate (B) at (1.1470,0.0000);      \coordinate (A) at (1.9220,1.4756);
  \coordinate (C) at (2.1993,2.0035);      % on n, beyond A
  \coordinate (D) at (3.0194,1.4756);      % on m
  \coordinate (E) at (2.5110,0.0000);      % on l
  \coordinate (P) at (5.3320,0.7626);
  \draw[fig] (0.4650,0) -- (3.8440,0);                        % l (solid part)
  \draw[fig] (0.9378,-0.3984) -- (2.3295,2.2514);             % n
  \draw[aux] (0.4340,1.4756) -- (3.3790,1.4756);              % m (dashed)
  \draw[aux] (3.3790,1.4756) .. controls +(0.62,-0.06) and +(-0.50,0.38) .. (P);
  \draw[aux] (3.8440,0)      .. controls +(0.72,0.06) and +(-0.66,-0.34) .. (P);
  \draw[fig] ($(P)+(0.16,0.20)$) -- (P) -- ($(P)+(0.16,-0.20)$);
  % the two congruent angles of the construction
  \draw[tick] ($(A)+(0:0.30)$)  arc (0:62.29:0.30);
  \draw[tick] ($(B)+(0:0.36)$)  arc (0:62.29:0.36);
  \foreach \p in {C,D,E}{\fill (\p) circle (1.0pt);}
  \node[vlab] at (1.5100,1.8700) {$A$};
  \node[vlab,right] at (C) {$C$};
  \node[vlab,above] at (D) {$D$};
  \node[vlab,above] at (E) {$E$};   % book sets E above its dot, as with D
  \node[vlab] at (1.5200,-0.3300) {$B$};
  \node[vlab,right] at (P) {$P$};
  \node[vlab,below right] at (D) {$m$};
  \node[vlab,above] at (1.75,0) {$l$};
  \node[vlab] at (1.1200,0.8500) {$n$};
\end{tikzpicture}
\figcap{8--2}
\end{minipage}\hfill
\begin{minipage}{0.48\linewidth}\centering
\begin{tikzpicture}[scale=0.88]
  \coordinate (A) at (1.5190,1.4756);   \coordinate (B) at (0.7750,0.0000);
  \draw[fig] (-0.0250,-0.0594) -- (3.1950,0.1796);            % l through B
  \draw[fig] (0.7990,1.6310)  -- (3.5690,1.0332);             % m through A
  \draw[fig] (0.5769,-0.3929) -- (1.7531,1.9399);             % n through A,B
  \draw[fig,-latex] (2.387,0.620) -- (2.945,0.620);
  \node[vlab,right] at (2.945,0.620) {$P$};
  \draw[tick] ($(A)+(-12.18:0.34)$) arc (-12.18:63.24:0.34);
  \draw[tick] ($(B)+(4.24:0.38)$)   arc (4.24:63.24:0.38);
  \node[slab] at (2.0604,1.7342) {$1$};
  \node[slab] at (1.3238,0.3666) {$2$};
  \node[vlab] at (0.9050,1.1830) {$A$};
  \node[vlab] at (0.1400,-0.4000) {$B$};
  \node[vlab,right] at (3.5690,1.0332) {$m$};
  \node[vlab,above] at (2.1000,0.0983) {$l$};
  \node[vlab] at (0.5000,0.7800) {$n$};
\end{tikzpicture}
\figcap{8--3}
\end{minipage}
\end{bkfigure}}

% ---------------------------------------------------------------- Fig 8-4
% Third proof: n cuts m at A and l at B, with the indicated angles at A and B
% congruent (both arced).  Q is on m with AQ congruent to BP.
% The book draws m and l PARALLEL (both horizontal) and bends each with a
% dashed curve to the hypothetical meeting point P -- the same convention as
% Fig 8-2.  That is what makes the two arced angles at A and B congruent, as
% the third proof requires; m must not be tilted off l.
\newcommand{\FIGVIIIFOUR}{\begin{bkfigure}[0.66\linewidth]
\begin{tikzpicture}[scale=0.86]
  \coordinate (B) at (2.0460,0.0000);  \coordinate (A) at (2.8830,1.6244);
  \coordinate (P) at (4.5880,0.9424);
  % Q on m, with m || l through A and AQ = BP = 2.7111 exactly
  % (both re-checked in tools/constraints/ch08.py)
  \coordinate (Q) at (0.1719,1.6244);
  \draw[fig] (0.0000,0) -- (B);                          % l, drawn up to B
  \draw[fig] (1.9269,-0.2311) -- (3.0342,1.9177);        % n
  \draw[aux] (Q) -- (A);                                 % m
  \draw[fig] (Q) -- (B);
  \draw[aux] (A) .. controls +(1.00,-0.07) and +(-0.52,0.40) .. (P);
  \draw[aux] (B) .. controls +(1.08,0.10) and +(-0.66,-0.42) .. (P);
  \draw[fig] ($(P)+(0.17,0.20)$) -- (P) -- ($(P)+(0.17,-0.20)$);
  % the two congruent angles at A and B (both 62.74 deg, since m || l)
  \draw[tick] ($(A)+(0:0.32)$) arc (0:62.74:0.32);
  \draw[tick] ($(B)+(0:0.36)$) arc (0:62.74:0.36);
  \node[vlab,above left] at (A) {$A$};
  \node[vlab,below right] at (B) {$B$};
  \node[vlab,above] at (Q) {$Q$};
  \node[vlab,right] at (P) {$P$};
  \node[vlab,left] at (Q) {$m$};
  \node[vlab,left] at (0.0000,0) {$l$};
  \node[vlab,above right] at (3.0342,1.9177) {$n$};
\end{tikzpicture}
\figcap{8--4}
\end{bkfigure}}

% ----------------------------------------------------------- Fig 8-5, 8-6
% 8-5: l || m and m || n -- all three drawn parallel -- with l and n bent
% (dashed) to meet at P, the contradiction.
% 8-6: l || m and transversal n through P on l.
\newcommand{\FIGVIIIFIVESIX}{\begin{bkfigure}[\linewidth]
\begin{minipage}{0.48\linewidth}\centering
\begin{tikzpicture}[scale=0.80]
  \draw[fig] (0.30,1.90) -- (2.55,1.90);                  % l
  \draw[fig] (0.30,1.12) -- (2.55,1.12);                  % n
  \draw[fig] (0.30,0.34) -- (2.55,0.34);                  % m
  \coordinate (P) at (3.60,1.50);
  \draw[aux] (2.55,1.90) -- (P);
  \draw[aux] (2.55,1.12) -- (P);
  \draw[fig] ($(P)+(0.17,0.20)$) -- (P) -- ($(P)+(0.17,-0.20)$);
  \node[vlab,left] at (0.30,1.90) {$l$};
  \node[vlab,left] at (0.30,1.12) {$n$};
  \node[vlab,left] at (0.30,0.34) {$m$};
  \node[vlab,right] at (P) {$P$};
\end{tikzpicture}
\figcap{8--5}
\end{minipage}\hfill
\begin{minipage}{0.48\linewidth}\centering
\begin{tikzpicture}[scale=0.80]
  \draw[fig] (0.20,1.70) -- (3.70,1.70);                  % l
  \draw[fig] (0.20,0.42) -- (3.70,0.42);                  % m
  \coordinate (P) at (1.72,1.70);
  \draw[fig] ($(P)!-0.28!(2.86,0.10)$) -- (2.86,0.10);    % n
  \fill (P) circle (1.0pt);
  \node[vlab,above right] at (P) {$P$};
  \node[vlab,right] at (3.70,1.70) {$l$};
  \node[vlab,right] at (3.70,0.42) {$m$};
  \node[vlab,below] at (2.86,0.10) {$n$};
\end{tikzpicture}
\figcap{8--6}
\end{minipage}
\end{bkfigure}}

% ----------------------------------------------------------- Fig 8-7, 8-8
% 8-7: angle ABC congruent to angle BCD, hence AB || CD.
% 8-8: AB = CD, angles at B and C right, hence AD || BC.
\newcommand{\FIGVIIISEVENEIGHT}{\begin{bkfigure}[\linewidth]
\begin{minipage}{0.50\linewidth}\centering
\begin{tikzpicture}[scale=0.80]
  \coordinate (A) at (0,0);    \coordinate (B) at (1.72,0);
  \coordinate (C) at (2.52,0.86);
  \coordinate (D) at ($(C)+(1.72,0)$);
  \draw[fig] (A) -- (B) -- (C) -- (D);
  \node[vlab,left] at (A) {$A$};
  \node[vlab,below right] at (B) {$B$};
  \node[vlab,above left] at (C) {$C$};
  \node[vlab,right] at (D) {$D$};
\end{tikzpicture}
\figcap{8--7}
\end{minipage}\hfill
\begin{minipage}{0.46\linewidth}\centering
\begin{tikzpicture}[scale=0.80]
  \coordinate (B) at (0,0);    \coordinate (C) at (3.05,0);
  \coordinate (A) at (0,1.72); \coordinate (D) at (3.05,1.72);
  \draw[fig] (A) -- (B) -- (C) -- (D) -- cycle;
  % right-angle boxes at B and C
  \draw[fig] ($(B)+(0,0.20)$) -- ($(B)+(0.20,0.20)$) -- ($(B)+(0.20,0)$);
  \draw[fig] ($(C)+(0,0.20)$) -- ($(C)+(-0.20,0.20)$) -- ($(C)+(-0.20,0)$);
  % double ticks on AB and DC
  \foreach \dy in {-0.07,0.07}{
    \draw[tick] ($(A)!0.5!(B)+(-0.13,\dy)$) -- ($(A)!0.5!(B)+(0.13,\dy)$);
    \draw[tick] ($(D)!0.5!(C)+(-0.13,\dy)$) -- ($(D)!0.5!(C)+(0.13,\dy)$);}
  \node[vlab,above left] at (A) {$A$};
  \node[vlab,below left] at (B) {$B$};
  \node[vlab,below right] at (C) {$C$};
  \node[vlab,above right] at (D) {$D$};
\end{tikzpicture}
\figcap{8--8}
\end{minipage}
\end{bkfigure}}

% --------------------------------------------------------- Fig 8-9, 8-10
% 8-9: transversal l cutting l1 and l2; the eight numbered angles.  As in the
% book, l leans the other way from l1/l2 -- its upper crossing is to the LEFT
% of its lower one.  Each number sits on the bisector of its own angle.
% 8-10: proof figure for Thm 8-4, with G chosen so that AG || l2 (drawn just
% off l1 on purpose -- the proof concludes that G must in fact lie on l1).
\newcommand{\FIGVIIININETEN}{\begin{bkfigure}[\linewidth]
\begin{minipage}{0.50\linewidth}\centering
\begin{tikzpicture}[scale=1.30]
  \coordinate (L1a) at (-0.0620,1.1036);  \coordinate (L1b) at (3.0380,1.0745);
  \coordinate (L2a) at ( 0.0620,0.1240);  \coordinate (L2b) at (3.0380,0.3887);
  \coordinate (La)  at ( 1.3640,1.5047);  \coordinate (Lb)  at (2.3355,0.0000);
  \draw[fig] (L1a) -- (L1b);                              % l1
  \draw[fig] (L2a) -- (L2b);                              % l2
  \draw[fig] (La) -- (Lb);                                % l
  \node[slab] at (1.8070,1.4030) {$1$};
  \node[slab] at (1.3180,1.2614) {$2$};
  \node[slab] at (1.4596,0.7724) {$3$};
  \node[slab] at (1.9486,0.9140) {$4$};
  \node[slab] at (2.2943,0.6320) {$5$};
  \node[slab] at (1.8128,0.4665) {$6$};
  \node[slab] at (1.9783,-0.0150) {$7$};
  \node[slab] at (2.4598,0.1505) {$8$};
  \node[vlab,right] at (L1b) {$l_1$};
  \node[vlab,right] at (L2b) {$l_2$};
  \node[vlab,above] at (La) {$l$};
\end{tikzpicture}
\figcap{8--9}
\end{minipage}\hfill
\begin{minipage}{0.46\linewidth}\centering
\begin{tikzpicture}[scale=0.92]
  \coordinate (A) at (1.62,1.42);  \coordinate (B) at (1.34,0.30);
  \draw[fig] (-0.20,1.42) -- (3.60,1.42);                 % l1
  \draw[fig] (-0.20,0.30) -- (3.60,0.30);                 % l2
  \coordinate (E) at ($(B)!1.68!(A)$);
  \draw[fig] ($(B)!-0.42!(A)$) -- (E);                    % l
  \coordinate (C) at (0.42,1.42);
  \coordinate (F) at (3.00,1.42);
  \coordinate (D) at (2.85,0.30);
  % G on the ray from A that the proof compares with l1
  \coordinate (G) at (2.60,1.86);
  \draw[aux] (A) -- (G);
  \foreach \p in {C,F,D,G,E}{\fill (\p) circle (1.0pt);}  % book dots E too
  \node[vlab,above] at (C) {$C$};
  \node[vlab] at (1.374,1.735) {$A$};
  \node[vlab,above] at (F) {$F$};
  \node[vlab] at (1.586,-0.015) {$B$};
  \node[vlab,above] at (D) {$D$};
  \node[vlab,above left] at (E) {$E$};
  \node[vlab,above left] at (G) {$G$};
  \node[vlab,right] at (3.60,1.42) {$l_1$};
  \node[vlab,right] at (3.60,0.30) {$l_2$};
  \node[vlab,below left] at ($(B)!-0.42!(A)$) {$l$};
\end{tikzpicture}
\figcap{8--10}
\end{minipage}
\end{bkfigure}}

% --------------------------------------------------------------- Fig 8-11
% Bowtie: AO = CO and BO = OD, so AB || CD and BC || AD.
\newcommand{\FIGVIIIELEVEN}{\begin{bkfigure}[0.56\linewidth]
\begin{tikzpicture}[scale=0.86]
  \coordinate (O) at (2.05,0.92);
  \coordinate (A) at (1.15,0);      \coordinate (C) at ($(A)!2!(O)$);
  \coordinate (B) at (0,0.72);      \coordinate (D) at ($(B)!2!(O)$);
  \draw[fig] (A) -- (C);
  \draw[fig] (B) -- (D);
  \draw[fig] (B) -- (A);
  \draw[fig] (C) -- (D);
  \node[vlab,below] at (A) {$A$};
  \node[vlab,left]  at (B) {$B$};
  \node[vlab,above] at (C) {$C$};
  \node[vlab,right] at (D) {$D$};
  \node[vlab] at (1.86,1.36) {$O$};
\end{tikzpicture}
\figcap{8--11}
\end{bkfigure}}

% --------------------------------------------------------------- Fig 8-12
% Adjacent angles AOB and BOC: OB interior to angle AOC.
\newcommand{\FIGVIIITWELVE}{\begin{bkfigure}[0.48\linewidth]
\begin{tikzpicture}[scale=0.90]
  \coordinate (O) at (0,0);
  \coordinate (A) at (76:2.15);
  \coordinate (B) at (48:1.80);
  \coordinate (C) at (-9:2.45);
  \draw[fig] (O) -- (A);
  \draw[fig] (O) -- (B);
  \draw[fig] (O) -- (C);
  \node[vlab,above] at (A) {$A$};
  \node[vlab,above right] at (B) {$B$};
  \node[vlab,right] at (C) {$C$};
  \node[vlab,left] at (O) {$O$};
\end{tikzpicture}
\figcap{8--12}
\end{bkfigure}}

% --------------------------------------------------------------- Fig 8-13
% Thm 8-5: l through A parallel to BC; D and E on l, one on each side of AB.
\newcommand{\FIGVIIITHIRTEEN}{\begin{bkfigure}[0.58\linewidth]
\begin{tikzpicture}[scale=0.88]
  \coordinate (B) at (0,0);  \coordinate (C) at (3.35,0);
  \coordinate (A) at (1.675,1.90);
  \draw[fig] (B) -- (C) -- (A) -- cycle;
  % l is the horizontal through A: same direction as BC by construction
  \coordinate (D) at ($(A)+(-1.15,0)$);
  \coordinate (E) at ($(A)+(1.15,0)$);
  \draw[fig] ($(A)+(-1.72,0)$) -- ($(A)+(1.72,0)$);
  \fill (D) circle (1.0pt); \fill (E) circle (1.0pt);
  \node[vlab,above] at (D) {$D$};
  \node[vlab,above] at (A) {$A$};
  \node[vlab,above] at (E) {$E$};
  \node[vlab,left]  at (B) {$B$};
  \node[vlab,right] at (C) {$C$};
  \node[vlab,right] at ($(A)+(1.72,0)$) {$l$};
\end{tikzpicture}
\figcap{8--13}
\end{bkfigure}}

% --------------------------------------------------------------- Fig 8-14
% Two triangles with AB congruent to A'B', angle A congruent to angle A',
% angle C congruent to angle C'.  The book sets the primed triangle up and to
% the right, overlapping the first in x, and mirrored in sense.
\newcommand{\FIGVIIIFOURTEEN}{\begin{bkfigure}[0.66\linewidth]
\begin{tikzpicture}[scale=0.95]
  \coordinate (A) at (0,1.35);  \coordinate (B) at (0.39,0);
  \coordinate (C) at (2.34,0.78);
  \draw[fig] (A) -- (B) -- (C) -- cycle;
  \node[vlab,left] at (A) {$A$};
  \node[vlab,below] at (B) {$B$};
  \node[vlab,below right] at (C) {$C$};
  % A'B'C' solved from the three side lengths of ABC, so it is congruent to
  % ABC exactly (opposite sense, as the book draws it); checked in ch08.py.
  \coordinate (Cp) at (1.2000,2.1000);
  \coordinate (Ap) at (3.4950,2.8302);
  \coordinate (Bp) at (3.1992,1.4565);
  \draw[fig] (Ap) -- (Bp) -- (Cp) -- cycle;
  \node[vlab,above right] at (Ap) {$A'$};
  \node[vlab,below right] at (Bp) {$B'$};
  \node[vlab,left] at (Cp) {$C'$};
\end{tikzpicture}
\figcap{8--14}
\end{bkfigure}}

% -------------------------------------------------------- Fig 8-15, 8-16
% 8-15: angle A right, AD perpendicular to BC (D is the foot).
% 8-16: D the midpoint of BC, with BD congruent to AD.
\newcommand{\FIGVIIIFIFTEENSIXTEEN}{\begin{bkfigure}[\linewidth]
\begin{minipage}{0.48\linewidth}\centering
\begin{tikzpicture}[scale=0.80]
  % right angle at A: put A on the circle with diameter BC
  \coordinate (B) at (0,0);  \coordinate (C) at (3.45,0);
  \coordinate (M) at ($(B)!0.5!(C)$);
  \coordinate (A) at ($(M)+(65:1.725)$);
  \coordinate (D) at ($(B)!(A)!(C)$);       % foot of the altitude
  \draw[fig] (B) -- (C) -- (A) -- cycle;
  \draw[fig] (A) -- (D);
  \node[vlab,above] at (A) {$A$};
  \node[vlab,left]  at (B) {$B$};
  \node[vlab,right] at (C) {$C$};
  \node[vlab,below] at (D) {$D$};
\end{tikzpicture}
\figcap{8--15}
\end{minipage}\hfill
\begin{minipage}{0.48\linewidth}\centering
\begin{tikzpicture}[scale=0.80]
  \coordinate (B) at (0,0);  \coordinate (C) at (3.45,0);
  \coordinate (D) at ($(B)!0.5!(C)$);       % midpoint of BC
  % A on the circle centred D through B, so that BD = AD
  \coordinate (A) at ($(D)+(62:1.725)$);
  \draw[fig] (B) -- (C) -- (A) -- cycle;
  \draw[fig] (A) -- (D);
  \node[vlab,above] at (A) {$A$};
  \node[vlab,left]  at (B) {$B$};
  \node[vlab,right] at (C) {$C$};
  \node[vlab,below] at (D) {$D$};
\end{tikzpicture}
\figcap{8--16}
\end{minipage}
\end{bkfigure}}

% -------------------------------------------------------- Fig 8-17, 8-18
% 8-17: a polygonal path P1 ... Pn.  8-18: a six-sided path that crosses itself.
\newcommand{\FIGVIIISEVENTEENEIGHTEEN}{\begin{bkfigure}[\linewidth]
\begin{minipage}{0.48\linewidth}\centering
\begin{tikzpicture}[scale=0.74]
  \coordinate (P1) at (0,0.30);
  \coordinate (P2) at (1.15,0.92);
  \coordinate (P3) at (0.62,2.62);
  \coordinate (Q1) at (1.72,1.72);
  \coordinate (Q2) at (2.35,2.42);
  \coordinate (Q3) at (2.95,1.05);
  \coordinate (Pn) at (2.28,0);
  \draw[fig] (P1) -- (P2) -- (P3) -- (Q1) -- (Q2) -- (Q3) -- (Pn);
  \node[vlab,left]  at (P1) {$P_1$};
  \node[vlab,below right] at (P2) {$P_2$};
  \node[vlab,above left] at (P3) {$P_3$};
  \node[vlab,below right] at (Pn) {$P_n$};
\end{tikzpicture}
\figcap{8--17}
\end{minipage}\hfill
\begin{minipage}{0.48\linewidth}\centering
\begin{tikzpicture}[scale=0.74]
  \coordinate (P1) at (0,0.62);
  \coordinate (P2) at (3.35,1.42);
  \coordinate (P3) at (1.05,2.62);
  \coordinate (P4) at (0.62,0);
  \coordinate (P5) at (1.95,2.62);
  \coordinate (P6) at (2.05,0);
  \coordinate (P7) at (0.30,1.95);
  \draw[fig] (P1) -- (P2) -- (P3) -- (P4) -- (P5) -- (P6) -- (P7);
  \node[vlab,left]  at (P1) {$P_1$};
  \node[vlab,right] at (P2) {$P_2$};
  \node[vlab,above] at (P3) {$P_3$};
  \node[vlab,below] at (P4) {$P_4$};
  \node[vlab,above] at (P5) {$P_5$};
  \node[vlab,below] at (P6) {$P_6$};
  \node[vlab,left]  at (P7) {$P_7$};
\end{tikzpicture}
\figcap{8--18}
\end{minipage}
\end{bkfigure}}

% --------------------------------------------------------------- Fig 8-19
% Three simple polygons: a quadrilateral, a triangle, and a non-convex
% pentagon (the notch at P2).
\newcommand{\FIGVIIININETEEN}{\begin{bkfigure}[\linewidth]
\begin{minipage}{0.31\linewidth}\centering
\begin{tikzpicture}[scale=0.66]
  % measured off the book crop: both uprights vertical, top edge rising to
  % the right (the earlier version had it falling to the right)
  \coordinate (P1) at (0,0);       \coordinate (P4) at (2.350,0.044);
  \coordinate (P3) at (2.350,1.862); \coordinate (P2) at (0,1.703);
  \draw[fig] (P1) -- (P4) -- (P3) -- (P2) -- cycle;
  \node[vlab,below left] at (P1) {$P_1$};
  \node[vlab,above left] at (P2) {$P_2$};
  \node[vlab,above right] at (P3) {$P_3$};
  \node[vlab,below right] at (P4) {$P_4$};
\end{tikzpicture}
\end{minipage}\hfill
\begin{minipage}{0.31\linewidth}\centering
\begin{tikzpicture}[scale=0.66]
  \coordinate (P1) at (0,0);   \coordinate (P3) at (2.62,0);
  \coordinate (P2) at (1.01,1.82);   % apex at 0.386 of the base, as measured
  \draw[fig] (P1) -- (P3) -- (P2) -- cycle;
  \node[vlab,below left] at (P1) {$P_1$};
  \node[vlab,above] at (P2) {$P_2$};
  \node[vlab,below right] at (P3) {$P_3$};
\end{tikzpicture}
\end{minipage}\hfill
\begin{minipage}{0.31\linewidth}\centering
\begin{tikzpicture}[scale=0.66]
  % measured off the book crop: P1P2 nearly level, notch (reflex vertex) at P2
  \coordinate (P1) at (0,1.275);   \coordinate (P2) at (1.214,1.213);
  \coordinate (P3) at (1.565,2.418); \coordinate (P4) at (2.621,1.803);
  \coordinate (P5) at (1.346,0);
  \draw[fig] (P1) -- (P2) -- (P3) -- (P4) -- (P5) -- cycle;
  \node[vlab,left] at (P1) {$P_1$};
  % book sets P2 below-right of the notch; placed explicitly so the glyph
  % clears side P4P5, which runs close by on that side
  \node[vlab] at (1.250,0.700) {$P_2$};
  \node[vlab,above] at (P3) {$P_3$};
  \node[vlab,right] at (P4) {$P_4$};
  \node[vlab,below] at (P5) {$P_5$};
\end{tikzpicture}
\end{minipage}
\figcap{8--19.\ \ Simple polygons.}
\end{bkfigure}}

% -------------------------------------------------------- Fig 8-20, 8-21
% 8-20: naming parts of a polygon -- a quadrilateral with two consecutive
% angles arced, and (as the book draws it) a PENTAGON carrying two of its
% diagonals, dashed, with a pair of adjacent sides called out.
% 8-21: convex and non-convex.
\newcommand{\FIGVIIITWENTYTWENTYONE}{\begin{bkfigure}[\linewidth]
\begin{minipage}{0.60\linewidth}\centering
\begin{tikzpicture}[scale=0.836]
  % left quadrilateral: two consecutive angles marked with arcs
  \coordinate (a1) at (0,0);       \coordinate (a2) at (2.42,0);
  \coordinate (a3) at (2.42,1.42); \coordinate (a4) at (0.55,1.42);
  \draw[fig] (a1) -- (a2) -- (a3) -- (a4) -- cycle;
  \draw[tick] ($(a3)+(180:0.36)$) arc (180:270:0.36);
  \draw[tick] ($(a2)+(180:0.36)$) arc (180:90:0.36);
  \node[slab,align=center] at (1.15,2.70) {Consecutive\\angles};
  \draw[fig,-latex,line width=0.4pt] (1.20,2.00) -- (2.14,1.19);
  \draw[fig,-latex,line width=0.4pt] (1.34,2.00) -- (2.14,0.29);
  % right pentagon: adjacent sides and two diagonals
  \begin{scope}[shift={(3.10,0)}]
    \coordinate (b1) at (0,0);           % bottom-left
    \coordinate (b2) at (0.373,1.779);   % top-left
    \coordinate (b3) at (2.315,1.779);   % top-right
    \coordinate (b4) at (2.929,0.886);   % right point
    \coordinate (b5) at (2.286,0);       % bottom-right
    \draw[fig] (b1) -- (b2) -- (b3) -- (b4) -- (b5) -- cycle;
    \draw[aux] (b1) -- (b3);
    \draw[aux] (b3) -- (b5);
    \node[slab,align=center] at (0.35,2.70) {Adjacent\\sides};
    \node[slab] at (2.45,2.86) {Diagonals};
    \draw[fig,-latex,line width=0.4pt] (0.24,2.00) -- (0.131,0.623);
    \draw[fig,-latex,line width=0.4pt] (0.62,2.00) -- (0.850,1.779);
    \draw[fig,-latex,line width=0.4pt] (2.18,2.40) -- (1.572,1.208);
    \draw[fig,-latex,line width=0.4pt] (2.62,2.40) -- (2.301,0.943);
  \end{scope}
\end{tikzpicture}
\figcap{8--20}
\end{minipage}\hfill
\begin{minipage}{0.36\linewidth}\centering
\begin{tikzpicture}[scale=0.62]
  \coordinate (c1) at (0,0);   \coordinate (c2) at (1.86,0);
  \coordinate (c3) at (2.05,1.42); \coordinate (c4) at (0.42,1.62);
  \draw[fig] (c1) -- (c2) -- (c3) -- (c4) -- cycle;
  \node[slab] at (1.05,2.20) {Convex};
  \begin{scope}[shift={(2.75,0)}]
    \coordinate (d1) at (0,0.55);  \coordinate (d2) at (1.30,1.05);
    \coordinate (d3) at (2.62,0);  \coordinate (d4) at (1.15,0.62);
    \draw[fig] (d1) -- (d2) -- (d3) -- (d4) -- cycle;
    \node[slab] at (1.30,2.20) {Not convex};
  \end{scope}
\end{tikzpicture}
\figcap{8--21}
\end{minipage}
\end{bkfigure}}

% -------------------------------------------------------- Fig 8-22, 8-23
% 8-22: a pentagon cut into three triangles by the two diagonals (dashed)
% from one vertex, with all five exterior angles extended (dashed) and arced.
% 8-23: the same idea for an n-gon -- the book draws a seven-gon split by the
% four diagonals from one vertex, with two exterior angles suggested.
\newcommand{\FIGVIIITWENTYTWOTHREE}{\begin{bkfigure}[\linewidth]
\begin{minipage}{0.48\linewidth}\centering
\begin{tikzpicture}[scale=0.88]
  % Vertex ratios measured off a 400 dpi crop of the book figure with a 50 px
  % grid (book p.147): side v1v2 horizontal, side v2v3 vertical, v4 nearly
  % above v1.  The earlier coordinates were noticeably taller and narrower.
  \coordinate (v1) at (1.3360,0.0000); \coordinate (v2) at (2.9500,0.0230);
  \coordinate (v3) at (2.9090,1.0060); \coordinate (v4) at (1.0290,1.6310);
  \coordinate (v5) at (0.0000,0.5490);
  \draw[fig] (v1) -- (v2) -- (v3) -- (v4) -- (v5) -- cycle;
  \draw[aux] (v1) -- (v3);
  \draw[aux] (v1) -- (v4);
  % exterior-angle extensions: at each vertex, extend the side arrived on
  \draw[aux] (v1) -- (1.7060,-0.1520);
  \draw[aux] (v2) -- (3.3500,0.0287);
  \draw[aux] (v3) -- (2.8923,1.4057);
  \draw[aux] (v4) -- (0.6494,1.7572);
  \draw[aux] (v5) -- (-0.2757,0.2591);
  \draw[tick] ($(v1)+(337.66:0.26)$) arc (337.66:360.82:0.26);
  \draw[tick] ($(v2)+(0.82:0.26)$)   arc (0.82:92.39:0.26);
  \draw[tick] ($(v3)+(92.39:0.26)$)  arc (92.39:161.61:0.26);
  \draw[tick] ($(v4)+(161.61:0.26)$) arc (161.61:226.44:0.26);
  \draw[tick] ($(v5)+(226.44:0.26)$) arc (226.44:337.66:0.26);
  % each number sits on its exterior angle's bisector, at a radius chosen so
  % the glyph clears both bounding rays (the angle at v1 is only 23 deg, so
  % that one has to sit further out)
  \node[slab] at (2.1711,-0.1587) {$1$};
  \node[slab] at (3.2935,0.3863)  {$2$};
  \node[slab] at (2.6081,1.4053)  {$3$};
  \node[slab] at (0.5245,1.5050)  {$4$};
  \node[slab] at (0.1044,0.0600)  {$5$};
\end{tikzpicture}
\figcap{8--22.\ \ Pentagon\\(three triangles).}
\end{minipage}\hfill
\begin{minipage}{0.48\linewidth}\centering
\begin{tikzpicture}[scale=0.88]
  \coordinate (w1) at (0.0000,1.6150);   % the vertex the diagonals come from
  \coordinate (w2) at (0.3610,2.6980);
  \coordinate (w3) at (0.8740,2.7930);
  \coordinate (w4) at (3.0400,1.4630);
  \coordinate (w5) at (2.8690,0.6840);
  \coordinate (w6) at (1.9190,0.0000);
  \coordinate (w7) at (0.7600,0.0380);
  \draw[fig] (w1) -- (w2) -- (w3) -- (w4) -- (w5) -- (w6) -- (w7) -- cycle;
  \foreach \v in {w3,w4,w5,w6}{\draw[fig] (w1) -- (\v);}
  \draw[aux] (w3) -- (1.4836,2.9059);
  \draw[aux] (w4) -- (3.1729,2.0686);
\end{tikzpicture}
\figcap{8--23.\ \ $n$-gon\\($n{-}2$ triangles).}
\end{minipage}
\end{bkfigure}}

% --------------------------------------------------------------- Fig 8-24
% Regular five-pointed star: the points are the vertices of a regular pentagon.
\newcommand{\FIGVIIITWENTYFOUR}{\begin{bkfigure}[0.42\linewidth]
\begin{tikzpicture}[scale=0.90]
  \foreach \i in {0,...,4}{\coordinate (s\i) at ({90+72*\i}:1.72);}
  \draw[fig] (s0) -- (s2) -- (s4) -- (s1) -- (s3) -- cycle;
\end{tikzpicture}
\figcap{8--24}
\end{bkfigure}}

% --------------------------------------------------------------- Fig 8-25
% Parallelogram, rectangle, square, rhombus.
\newcommand{\FIGVIIITWENTYFIVE}{\begin{bkfigure}[\linewidth]
\begin{tikzpicture}[scale=0.70]
  % parallelogram
  \begin{scope}
    \draw[fig] (0,0) -- (1.42,0) -- (1.86,1.72) -- (0.44,1.72) -- cycle;
    \node[slab,below] at (0.92,-0.12) {Parallelogram};
  \end{scope}
  % rectangle
  \begin{scope}[shift={(3.05,0)}]
    \draw[fig] (0,0) -- (2.05,0) -- (2.05,1.62) -- (0,1.62) -- cycle;
    \node[slab,below] at (1.02,-0.12) {Rectangle};
  \end{scope}
  % square
  \begin{scope}[shift={(6.30,0)}]
    \draw[fig] (0,0) -- (1.62,0) -- (1.62,1.62) -- (0,1.62) -- cycle;
    \node[slab,below] at (0.80,-0.12) {Square};
  \end{scope}
  % rhombus
  \begin{scope}[shift={(8.95,0)}]
    \draw[fig] (0,0) -- (1.62,0) -- (2.06,1.56) -- (0.44,1.56) -- cycle;
    \node[slab,below] at (1.02,-0.12) {Rhombus};
  \end{scope}
\end{tikzpicture}
\figcap{8--25}
\end{bkfigure}}

% -------------------------------------------------------- Fig 8-26, 8-27
% 8-26: DE || AC and DF || AB.   8-27: regular hexagon ABCDEF, diagonal FD.
\newcommand{\FIGVIIITWENTYSIXSEVEN}{\begin{bkfigure}[\linewidth]
\begin{minipage}{0.50\linewidth}\centering
\begin{tikzpicture}[scale=0.80]
  \coordinate (B) at (0,0);  \coordinate (C) at (3.45,0);
  \coordinate (A) at (1.72,1.86);
  \coordinate (D) at ($(B)!0.52!(C)$);
  % E on AB and F on AC with DE || AC, DF || AB  (equal ratios)
  \coordinate (E) at ($(B)!0.52!(A)$);
  \coordinate (F) at ($(C)!0.48!(A)$);
  \draw[fig] (B) -- (C) -- (A) -- cycle;
  \draw[aux] (D) -- (E);
  \draw[aux] (D) -- (F);
  \node[vlab,above] at (A) {$A$};
  \node[vlab,left]  at (B) {$B$};
  \node[vlab,right] at (C) {$C$};
  \node[vlab,below] at (D) {$D$};
  \node[vlab,above left]  at (E) {$E$};
  \node[vlab,above right] at (F) {$F$};
\end{tikzpicture}
\figcap{8--26}
\end{minipage}\hfill
\begin{minipage}{0.46\linewidth}\centering
\begin{tikzpicture}[scale=0.80]
  \foreach \i in {0,...,5}{\coordinate (h\i) at ({-60+60*\i}:1.42);}
  \draw[fig] (h0) -- (h1) -- (h2) -- (h3) -- (h4) -- (h5) -- cycle;
  \draw[aux] (h4) -- (h2);
  \node[vlab,below right] at (h0) {$B$};
  \node[vlab,right] at (h1) {$C$};
  \node[vlab,above right] at (h2) {$D$};
  \node[vlab,above left] at (h3) {$E$};
  \node[vlab,left] at (h4) {$F$};
  \node[vlab,below left] at (h5) {$A$};
\end{tikzpicture}
\figcap{8--27}
\end{minipage}
\end{bkfigure}}

% -------------------------------------------------------- Fig 8-28, 8-29
% 8-28: two cevians AQ and AR from A.  The lettered angles are, exactly as
% the book places them:  x = angle B, y = angle BAQ, s = angle RAC,
% z = angle AQC (right of AQ), t = angle ARB (left of AR), r = angle C.
% With that reading x + y + t = r + s + z, which is what Ex. 9 asks for; it
% is verified numerically in tools/constraints/ch08.py.
% 8-29: AB = BC, BE = BD, angle CBD = 40 degrees; angle B obtuse as drawn.
\newcommand{\FIGVIIITWENTYEIGHTNINE}{\begin{bkfigure}[\linewidth]
\begin{minipage}{0.48\linewidth}\centering
\begin{tikzpicture}[scale=1.05]
  % Ratios measured off a 400 dpi crop of the book figure (book p.151):
  % Q at 0.369 and R at 0.620 of BC, apex almost exactly over Q so that AQ
  % reads as vertical.  A is nudged 0.04 left of Q so that angle AQC is
  % 91.3 deg, as the book draws it, rather than an unstated exact right angle.
  \coordinate (B) at (0,0);  \coordinate (C) at (3.45,0);
  \coordinate (A) at (1.2331,1.7664);
  \coordinate (Q) at ($(B)!0.369!(C)$);
  \coordinate (R) at ($(B)!0.620!(C)$);
  \draw[fig] (B) -- (C) -- (A) -- cycle;
  \draw[fig] (A) -- (Q);
  \draw[fig] (A) -- (R);
  \node[vlab,above] at (A) {$A$};
  \node[vlab,left]  at (B) {$B$};
  \node[vlab,right] at (C) {$C$};
  % the two angles at A carry arcs, as in the book
  \draw[tick] ($(A)+(-124.92:0.45)$) arc (-124.92:-88.70:0.45);
  \draw[tick] ($(A)+(-62.85:0.50)$)  arc (-62.85:-38.55:0.50);
  % each letter on the bisector of the angle it names, far enough out to clear
  % both sides of that angle
  \node[slab] at (0.4079,0.2127) {$x$};
  \node[slab] at (1.0307,1.0963) {$y$};
  \node[slab] at (1.7716,1.1087) {$s$};
  \node[slab] at (1.5527,0.2860) {$z$};
  \node[slab] at (1.7465,0.2398) {$t$};
  \node[slab] at (2.8459,0.2113) {$r$};
\end{tikzpicture}
\figcap{8--28}
\end{minipage}\hfill
\begin{minipage}{0.48\linewidth}\centering
\begin{tikzpicture}[scale=0.80]
  \coordinate (B) at (0,0);  \coordinate (C) at (3.05,0);
  % AB = BC: A on the circle centred B with radius BC, at 110 deg as drawn
  \coordinate (A) at ($(B)+(110:3.05)$);
  % D is where the 40-degree ray from B meets AC.  Solved exactly and written
  % out as literals so no extra tikz library is needed; angle CBD = 40.0000
  % deg is re-checked in tools/constraints/ch08.py.
  \coordinate (D) at (1.38738,1.16417);
  % BE = BD: E on BA at the same distance from B as D  (ratio 0.59381)
  \coordinate (E) at ($(B)!0.59381!(A)$);
  \draw[fig] (B) -- (C) -- (A) -- cycle;
  \draw[fig] (B) -- (D);
  \draw[fig] (E) -- (D);
  \node[vlab,above left] at (A) {$A$};
  \node[vlab,below] at (B) {$B$};
  \node[vlab,below right] at (C) {$C$};
  \node[vlab] at (1.6670,1.4040) {$D$};
  \node[vlab] at (-1.3700,1.4300) {$E$};
  \node[slab] at (1.2592,0.4583) {$40^\circ$};
\end{tikzpicture}
\figcap{8--29}
\end{minipage}
\end{bkfigure}}

% --------------------------------------------------------------- Fig 8-30
% AB congruent to A'B', with B between A and C.
\newcommand{\FIGVIIITHIRTY}{\begin{bkfigure}[0.66\linewidth]
\begin{tikzpicture}[scale=0.90]
  \draw[fig] (-0.20,0) -- (2.20,0);
  \coordinate (A) at (0,0); \coordinate (B) at (1.05,0);
  \coordinate (C) at (1.86,0);
  \foreach \p in {A,B,C}{\fill (\p) circle (1.0pt);}
  \draw[tick] ($(A)!0.5!(B)+(0,-0.13)$) -- ($(A)!0.5!(B)+(0,0.13)$);
  \node[vlab,above] at (A) {$A$};
  \node[vlab,above] at (B) {$B$};
  \node[vlab,above] at (C) {$C$};
  \begin{scope}[shift={(3.05,0)}]
    \draw[fig] (-0.20,0) -- (1.60,0);
    \coordinate (Ap) at (0,0); \coordinate (Bp) at (1.05,0);
    \fill (Ap) circle (1.0pt); \fill (Bp) circle (1.0pt);
    \draw[tick] ($(Ap)!0.5!(Bp)+(0,-0.13)$) -- ($(Ap)!0.5!(Bp)+(0,0.13)$);
    \node[vlab,above] at (Ap) {$A'$};
    \node[vlab,above] at (Bp) {$B'$};
  \end{scope}
\end{tikzpicture}
\figcap{8--30}
\end{bkfigure}}

% -------------------------------------------------------- Fig 8-31, 8-32
% 8-31: angles CAB and ABD right, CA = DB, M and N midpoints of CD and AB.
% 8-32: X and Y on one side of AB; rays AX and BY meet at P.
\newcommand{\FIGVIIITHIRTYONETWO}{\begin{bkfigure}[\linewidth]
\begin{minipage}{0.50\linewidth}\centering
\begin{tikzpicture}[scale=0.80]
  \coordinate (A) at (0,0);      \coordinate (B) at (3.05,0);
  \coordinate (C) at (0,1.62);   \coordinate (D) at (3.05,1.62);
  \coordinate (M) at ($(C)!0.5!(D)$);
  \coordinate (N) at ($(A)!0.5!(B)$);
  \draw[fig] (A) -- (B) -- (D) -- (C) -- cycle;
  \draw[aux] (M) -- (N);
  % right-angle boxes at A and B
  \draw[fig] ($(A)+(0,0.20)$) -- ($(A)+(0.20,0.20)$) -- ($(A)+(0.20,0)$);
  \draw[fig] ($(B)+(0,0.20)$) -- ($(B)+(-0.20,0.20)$) -- ($(B)+(-0.20,0)$);
  % single ticks on CM, MD ; double ticks on AN, NB
  \draw[tick] ($(C)!0.5!(M)+(0,-0.13)$) -- ($(C)!0.5!(M)+(0,0.13)$);
  \draw[tick] ($(M)!0.5!(D)+(0,-0.13)$) -- ($(M)!0.5!(D)+(0,0.13)$);
  \foreach \dx in {-0.06,0.06}{
    \draw[tick] ($(A)!0.5!(N)+(\dx,-0.13)$) -- ($(A)!0.5!(N)+(\dx,0.13)$);
    \draw[tick] ($(N)!0.5!(B)+(\dx,-0.13)$) -- ($(N)!0.5!(B)+(\dx,0.13)$);}
  \node[vlab,above left] at (C) {$C$};
  \node[vlab,above] at (M) {$M$};
  \node[vlab,above right] at (D) {$D$};
  \node[vlab,below left] at (A) {$A$};
  \node[vlab,below] at (N) {$N$};
  \node[vlab,below right] at (B) {$B$};
\end{tikzpicture}
\figcap{8--31}
\end{minipage}\hfill
\begin{minipage}{0.46\linewidth}\centering
\begin{tikzpicture}[scale=0.80]
  \coordinate (A) at (0,1.72);  \coordinate (B) at (0,0);
  \coordinate (P) at (3.05,1.02);
  \coordinate (X) at ($(A)!0.55!(P)$);
  \coordinate (Y) at ($(B)!0.55!(P)$);
  \draw[fig] (A) -- (B);
  \draw[fig] (A) -- (X);
  \draw[fig] (B) -- (Y);
  \draw[aux] (X) -- (P);
  \draw[aux] (Y) -- (P);
  \draw[fig] ($(P)+(0.17,0.20)$) -- (P) -- ($(P)+(0.17,-0.20)$);
  \fill (X) circle (1.0pt); \fill (Y) circle (1.0pt);
  \node[vlab,above] at (A) {$A$};
  \node[vlab,below] at (B) {$B$};
  \node[vlab,above] at (X) {$X$};
  \node[vlab,anchor=north] at ($(Y)+(0,-0.32)$) {$Y$};
  \node[vlab,above right] at (P) {$P$};
\end{tikzpicture}
\figcap{8--32}
\end{minipage}
\end{bkfigure}}

% figures09.tex -- Chapter 9 figures, rebuilt 2026-08-16 from the iPad scans.
%
% Every figure below was measured off a scan crop (50/100 px grid overlay,
% tools/ + scratchpad/ch09/grid.py) rather than eyeballed from the photo PDF.
% Constrained points are DERIVED, never placed by hand:
%   parallels  -> equal ratios along the two sides   ($(X)!r!(Y)$)
%   altitudes  -> foot of perpendicular              ($(X)!(P)!(Y)$)
%   crossings  -> (intersection of A--B and C--D)
%   bisectors  -> ratio from the angle-bisector theorem
% so the hypothesis stated in the book holds exactly in the drawing.
% Hypotheses are re-asserted numerically in tools/constraints/ch09.py.
%
% Macro names are chapter-prefixed: \FIGIX<NUMBERWORD>.
%
% Line treatment follows the book's own convention: solid where the book
% prints solid (given), dashed where the book prints dashed (aux).  The
% heavier `key' weight is GONE from this chapter (2026-08-17): Caleb flagged
% Figs 9-23/9-25/9-14/9-27/9-50 for drawing one segment heavier than the rest
% where the book strikes every line at a single weight.  Do not reintroduce it.

\usetikzlibrary{decorations.pathreplacing,arrows.meta}

% The book's fine dashed rule (Fig 9-24's DE, Fig 9-50's compass arcs, the
% dashed base extension of Fig 9-54): thinner and shorter-dashed than `aux'.
\tikzset{
  finedash/.style={fig, line width=0.55pt, dash pattern=on 2.4pt off 1.9pt},
  longdash/.style={fig, dash pattern=on 5.2pt off 3.0pt},
  lead/.style={\clB, line width=0.5pt, -{Latex[length=3.6pt,width=2.6pt]}},
}

% Primed labels.  The math prime sits ~4pt above the letter, which makes the
% collision audit read it as a separate word overlapping its own letter; a
% 1.4pt drop keeps the pair in one word and is visually indistinguishable.
\newcommand{\IXpr}{\raisebox{-1.3pt}{\ensuremath{{}^{\prime}}}}

% Chapter-local label clearance.  brumfiel.sty sets outer sep = 3.4pt, which
% is not quite enough when a labelled point lies ON a line: the line clips the
% corner of the box.  4.8pt clears it without the labels drifting loose.
\tikzset{
  vlab/.append style={outer sep=2.9pt},
  slab/.append style={outer sep=2.9pt},
  klab/.append style={outer sep=2.9pt},
}

% Book-style dimension brace (Figs 9-14, 9-45, 9-46).
\tikzset{
  bkbrace/.style={fig, line width=0.5pt, decorate,
                  decoration={brace, amplitude=3.2pt, raise=1pt}},
  bkbracer/.style={fig, line width=0.5pt, decorate,
                  decoration={brace, mirror, amplitude=3.2pt, raise=1pt}},
}

% ---------------------------------------------------------------- Fig 9-1
% Three parallels l, l', l''; transversal m leans left going down, m' leans
% right; m'' through A and m''' through B are parallel to m' (dashed in the
% book).  Spacing of the parallels is equal because AB = BC on m.
\newcommand{\FIGIXONE}{\begin{bkfigure}[\linewidth]
\begin{tikzpicture}[scale=1.133]
  \foreach \y/\nm in {2.2/$l$, 1.1/$l\IXpr$, 0/$l\IXpr\IXpr$}
    {\draw[given] (0.35,\y) -- (9.05,\y);
     \node[vlab, right] at (9.05,\y) {\nm};}
  % transversal m : dx/dy = -0.527
  \coordinate (A) at (2.90,2.2);
  \coordinate (B) at (2.32,1.1);
  \coordinate (C) at (1.75,0);
  \draw[given] (3.19,2.75) -- (1.57,-0.35);
  % transversal m' : dx/dy = +0.341
  \coordinate (Ap) at (6.90,2.2);
  \coordinate (Bp) at (7.275,1.1);
  \coordinate (Cp) at (7.65,0);
  \draw[given] (6.71,2.75) -- (7.77,-0.35);
  % m'' through A, m''' through B, both parallel to m'
  \coordinate (D) at (3.275,1.1);
  \coordinate (E) at (2.695,0);
  \draw[aux] (A) -- (3.75,-0.30);
  \draw[aux] (B) -- (2.80,-0.30);
  \foreach \p/\n/\d in {A/A/above left, B/B/above left, C/C/above left,
                        D/D/above right, E/E/above right,
                        Ap/A\IXpr/above right, Bp/B\IXpr/above right, Cp/C\IXpr/above right}
    {\node[vlab,\d] at (\p) {$\n$};}
  \node[vlab, below] at (1.57,-0.35) {$m$};
  \node[vlab, below] at (2.80,-0.30) {$m\IXpr\IXpr\IXpr$};
  \node[vlab, below] at (3.75,-0.30) {$m\IXpr\IXpr$};
  \node[vlab, below] at (7.77,-0.35) {$m\IXpr$};
\end{tikzpicture}
\figcap{9--1}
\end{bkfigure}}

% ------------------------------------------------------------ Fig 9-2, 9-3
% Same construction twice: D, E midpoints of AB, BC.  In 9-3 the line m
% through D parallel to AC is carried past E as a dashed ray.
\newcommand{\FIGIXTWOTHREE}{\begin{bkfigure}[\linewidth]
\begin{minipage}{0.36\linewidth}\centering
\begin{tikzpicture}[scale=0.70]
  \coordinate (A) at (0,0); \coordinate (C) at (3.93,0);
  \coordinate (B) at (1.87,1.75);
  \draw[given] (A) -- (C) -- (B) -- cycle;
  \coordinate (D) at ($(A)!0.5!(B)$);
  \coordinate (E) at ($(B)!0.5!(C)$);
  \draw[given] (D) -- (E);
  \foreach \p/\n/\d in {A/A/below left, C/C/below right, B/B/above,
                        D/D/above left, E/E/above right}
    {\node[vlab,\d] at (\p) {$\n$};}
\end{tikzpicture}
\figcap{9--2}
\end{minipage}\hfill
\begin{minipage}{0.47\linewidth}\centering
\begin{tikzpicture}[scale=0.70]
  \coordinate (A) at (0,0); \coordinate (C) at (5.65,0);
  \coordinate (B) at (1.83,1.73);
  \draw[given] (A) -- (C) -- (B) -- cycle;
  \coordinate (D) at ($(A)!0.5!(B)$);
  \coordinate (E) at ($(B)!0.5!(C)$);
  \draw[given] (D) -- (E);
  \draw[aux] (E) -- ($(D)!1.42!(E)$);
  \node[vlab,right] at ($(D)!1.42!(E)$) {$m$};
  \foreach \p/\n/\d in {A/A/below left, C/C/below right, B/B/above,
                        D/D/above left, E/E/above}
    {\node[vlab,\d] at (\p) {$\n$};}
\end{tikzpicture}
\figcap{9--3}
\end{minipage}
\end{bkfigure}}

% ---------------------------------------------------------------- Fig 9-4
% D, E, G are the midpoints of AB, BC, CA.  m' is the line through E parallel
% to AB; it meets AC in G.
\newcommand{\FIGIXFOUR}{\begin{bkfigure}[0.62\linewidth]
\begin{tikzpicture}[scale=1.667]
  \coordinate (A) at (0,0); \coordinate (C) at (2.39,0);
  \coordinate (B) at (0.69,0.98);
  \draw[given] (A) -- (C) -- (B) -- cycle;
  \coordinate (D) at ($(A)!0.5!(B)$);
  \coordinate (E) at ($(B)!0.5!(C)$);
  \coordinate (G) at ($(A)!0.5!(C)$);
  \coordinate (Mp) at ($(G)!1.60!(E)$);
  \draw[given] (D) -- (E);
  % m' is a LINE: the book runs it past G on the far side as well.
  \draw[given] ($(E)!1.34!(G)$) -- (Mp);
  \node[vlab,above right] at (Mp) {$m\IXpr$};
  \foreach \p/\n/\d in {A/A/below left, C/C/right, B/B/above,
                        D/D/above left, G/G/below right}
    {\node[vlab,\d] at (\p) {$\n$};}
  \node[vlab] at ($(E)+(0.41,0.13)$) {$E$};   % right of m', as in the book
\end{tikzpicture}
\figcap{9--4}
\end{bkfigure}}

% ------------------------------------------------------------ Fig 9-5..9-7
% SPLIT 2026-08-17: Ex. 9-2 cites 9-5 in #2, 9-6 in #4/5/6 and 9-7 in #7, so
% each now stands alone at full width and is set inline after its exercise.
\newcommand{\FIGIXFIVE}{\begin{bkfigure}[0.62\linewidth]
\begin{tikzpicture}[scale=1.55]
  % 9-5: triangle EFG, C the midpoint of FG, D fixed by CD || EG and DG || EF,
  % which forces CD = 1/2 EG exactly (the statement to be proved).
  % Vertices measured off the p.156 figure at 400 dpi (E,F,G,C,D all checked:
  % C is the midpoint of FG and D sits at C + (G-E)/2, as the two parallels
  % require).  Do NOT draw to $(C)+(G)-(E)$ -- that is twice too far out and
  % breaks DG || EF.
  \coordinate (E) at (0,0); \coordinate (G) at (3.88,-0.12);
  \coordinate (F) at (0.60,1.338);
  \coordinate (C) at ($(F)!0.5!(G)$);
  \coordinate (D) at ($(C)!0.5!($(C)+(G)-(E)$)$);
  \draw[given] (E) -- (F) -- (G) -- cycle;
  \draw[given] (G) -- (D);
  \draw[given]   (C) -- (D);
  \foreach \p/\n/\d in {E/E/left, F/F/above, G/G/below right}
    {\node[vlab,\d] at (\p) {$\n$};}
  \node[vlab,above] at (C) {$C$};
  \node[vlab,above right] at (D) {$D$};
\end{tikzpicture}
\figcap{9--5}
\end{bkfigure}}

\newcommand{\FIGIXSIX}{\begin{bkfigure}[0.60\linewidth]
\begin{tikzpicture}[scale=1.45]
  % 9-6: three parallels, unequal spacing (the exercises vary AB : BC),
  % cut by transversals m (leaning left) and m' (leaning right).
  % spacing AB : BC = 1.90 : 1, measured off the p.156 figure at 400 dpi.
  \foreach \y in {0, 0.70, 2.03}{\draw[given] (-0.62,\y) -- (3.35,\y);}
  \coordinate (C) at (0,0); \coordinate (B) at (0.0854,0.70);
  \coordinate (A) at (0.2477,2.03);
  \coordinate (Cp) at (2.85,0); \coordinate (Bp) at (2.7716,0.70);
  \coordinate (Ap) at (2.6226,2.03);
  \draw[given] (-0.0915,-0.75) -- (0.3577,2.95);
  \draw[given] (2.934,-0.75)  -- (2.5136,2.95);
  \node[vlab,above] at (0.3577,2.95) {$m$};
  \node[vlab,above] at (2.5136,2.95) {$m\IXpr$};
  \foreach \p/\n/\d in {A/A/above left, B/B/above left, C/C/below left,
                        Ap/A\IXpr/above right, Bp/B\IXpr/above right, Cp/C\IXpr/below right}
    {\node[vlab,\d] at (\p) {$\n$};}
\end{tikzpicture}
\figcap{9--6}
\end{bkfigure}}

\newcommand{\FIGIXSEVEN}{\begin{bkfigure}[0.55\linewidth]
\begin{tikzpicture}[scale=2.30]
  % 9-7: F on EG and P on EQ at the same ratio 2/5, hence FP || GQ.
  \coordinate (G) at (0,0); \coordinate (Q) at (2.35,0);
  \coordinate (E) at (0.80,1.15);
  \draw[given] (G) -- (Q) -- (E) -- cycle;
  \coordinate (F) at ($(E)!0.4!(G)$);
  \coordinate (P) at ($(E)!0.4!(Q)$);
  \draw[given] (F) -- (P);
  \foreach \p/\n/\d in {G/G/below left, Q/Q/below right, E/E/above,
                        F/F/left, P/P/above right}
    {\node[vlab,\d] at (\p) {$\n$};}
\end{tikzpicture}
\figcap{9--7}
\end{bkfigure}}

% ---------------------------------------------------------------- Fig 9-8
% AD = DE = EF = FG on the ray AC; the parallels to BG through D, E, F cut AB
% at its quarter points H, I, J -- so every crossing is derived, not placed.
\newcommand{\FIGIXEIGHT}{\begin{bkfigure}[0.72\linewidth]
\begin{tikzpicture}[scale=1.479]
  \coordinate (A) at (0,0); \coordinate (C) at (5.30,0);
  \coordinate (B) at (3.02,1.50);
  \coordinate (G) at (4.22,0);
  \draw[given] (A) -- (C);
  \draw[given] (A) -- ($(A)!1.09!(B)$);
  \draw[given] (B) -- (G);
  \foreach \i in {1,2,3}{%
     \coordinate (P\i) at ($(A)!\i/4!(G)$);
     \coordinate (H\i) at ($(A)!\i/4!(B)$);
     \draw[given] (H\i) -- (P\i);}
  \node[vlab,left] at (A) {$A$};   \node[vlab,right] at (C) {$C$};
  \node[vlab,above left] at (B) {$B$};
  \node[vlab,below] at (P1) {$D$}; \node[vlab,below] at (P2) {$E$};
  \node[vlab,below] at (P3) {$F$}; \node[vlab,below] at (G) {$G$};
  \node[vlab,above=2.5pt] at (H1) {$H$}; \node[vlab,above=2.5pt] at (H2) {$I$};
  \node[vlab,above=2.5pt] at (H3) {$J$};
\end{tikzpicture}
\figcap{9--8}
\end{bkfigure}}

% ------------------------------------------------------------- Fig 9-9
% SPLIT 2026-08-17 out of \FIGIXNINEFOURTEEN: Ex. 9-3 cites 9-9 in #5/#6,
% 9-10 in #7/#8, 9-11 in #9, 9-12 in #11, 9-13 in #12 and 9-14 in #13, so the
% six now stand alone and are set inline after the exercise that names them.
% 9-9: D on AB, E on BC with BD/BA = BE/BC, hence DE || AC exactly.
\newcommand{\FIGIXNINE}{\begin{bkfigure}[0.70\linewidth]
\begin{tikzpicture}[scale=1.45]
  \coordinate (A) at (0,0); \coordinate (C) at (4.12,0.42);
  \coordinate (B) at (1.32,2.32);
  \draw[given] (A) -- (C) -- (B) -- cycle;
  \coordinate (D) at ($(A)!0.172!(B)$);
  \coordinate (E) at ($(B)!0.828!(C)$);
  \draw[given] (D) -- (E);
  % The book's "Parallel" note is a FAN: single-headed leaders from one apex
  % at the word, one per line named.  (Caleb: our double-headed spans are
  % wrong -- see 9-10, 9-12, 9-14.)
  \coordinate (PA) at (3.44,1.66);
  \node[klab] at (3.30,2.02) {Parallel};
  \draw[lead] (PA) -- ($(D)!0.70!(E)$);
  \draw[lead] (PA) -- ($(A)!0.80!(C)$);
  \foreach \p/\n/\d in {A/A/below left, C/C/right, B/B/above,
                        D/D/left, E/E/above right}
    {\node[vlab,\d] at (\p) {$\n$};}
\end{tikzpicture}
\figcap{9--9}
\end{bkfigure}}

% ------------------------------------------------------------ Fig 9-10
% REBUILT 2026-08-17 against the book photograph (feedback/small/20260817_102341).
% What was wrong: T1 missed A entirely, T1/T2 were drawn near-vertical instead
% of the book's ~31 deg off vertical, the three parallels were dead horizontal
% and cut off just past the outer transversals, and both "Parallel" notes were
% double-headed arrows.  Measured off the photograph: the parallels slope
% 8 deg down to the right, spacing AB : BC = 1.5 : 1, the two slanted
% transversals run 0.586 across per 1 down, and both notes are fans of
% single-headed leaders.
\newcommand{\FIGIXTEN}{\begin{bkfigure}[0.94\linewidth]
\begin{tikzpicture}[scale=1.32]
  % three parallels  y = c - 0.145 x  for c = 2.00, 0.80, 0
  \foreach \c in {2.00,0.80,0}
    {\draw[given] (-0.30,{\c+0.0435}) -- (5.60,{\c-0.812});}
  % transversal 1 through A, nearly vertical, leaning LEFT as it descends
  \coordinate (A) at (0.20,1.971);
  \coordinate (B) at (0.0817,0.7882);
  \coordinate (C) at (0.00286,-0.0004);
  \draw[given] (0.255,2.521) -- (-0.045,-0.479);
  % transversal 2 through the SAME point A (the book crosses them on l1)
  \coordinate (D) at (0.96850,0.65957);
  \coordinate (X) at (1.48083,-0.21472);
  \coordinate (T2a) at (-0.0051,2.321);  \coordinate (T2b) at (1.7352,-0.649);
  \draw[given] (T2a) -- (T2b);
  % transversal 3, parallel to it, through E on l1 and F on l3
  \coordinate (E) at (3.30,1.5215);
  \coordinate (F) at (4.58083,-0.66422);
  \coordinate (T3a) at (3.0949,1.8715); \coordinate (T3b) at (4.8353,-1.0985);
  \draw[given] (T3a) -- (T3b);
  % "Parallel" (top right): three leaders fanning from one apex, one per line
  \coordinate (PA) at (5.35,2.02);
  \node[klab] at (5.02,2.40) {Parallel};
  \draw[lead] (PA) -- (5.15,1.2533);
  \draw[lead] (PA) -- (5.20,0.0460);
  \draw[lead] (PA) -- (5.42,-0.7859);
  % "Parallel" (bottom): one leader each side of the word, to the two slants
  \node[klab] at (3.18,-0.72) {Parallel};
  \draw[lead] (2.52,-0.66) -- ($(T2b)!0.06!(T2a)$);
  \draw[lead] (3.84,-0.78) -- ($(T3b)!0.06!(T3a)$);
  \node[vlab] at ($(A)+(-0.20,-0.22)$) {$A$};
  \node[vlab] at ($(B)+(-0.20,0.22)$) {$B$};
  \node[vlab] at ($(C)+(-0.20,0.22)$) {$C$};
  \node[vlab] at ($(D)+(0.18,0.20)$) {$D$};
  \node[vlab] at ($(E)+(0.18,0.20)$) {$E$};
  \node[vlab] at ($(F)+(0.18,0.20)$) {$F$};
\end{tikzpicture}
\figcap{9--10}
\end{bkfigure}}

% ------------------------------------------------------------ Fig 9-11
% AB : BC : CD = 4 : 3 : 2 down the left side; E and F sit on AG at the
% matching ratios, so BE || CF || DG exactly.  AE is marked 8.
\newcommand{\FIGIXELEVEN}{\begin{bkfigure}[0.68\linewidth]
\begin{tikzpicture}[scale=1.35]
  \coordinate (A) at (0,3.15); \coordinate (D) at (0,0);
  \coordinate (G) at (4.55,0);
  \coordinate (B) at (0,1.75); \coordinate (C) at (0,0.70);
  \coordinate (E) at ($(A)!4/9!(G)$);
  \coordinate (F) at ($(A)!7/9!(G)$);
  \draw[given] (A) -- (D) -- (G) -- cycle;
  \draw[given] (B) -- (E); \draw[given] (C) -- (F);
  \node[slab,left=0pt]  at (0,2.45) {4};
  \node[slab,left=0pt]  at (0,1.225) {3};
  \node[slab,left=0pt]  at (0,0.35) {2};
  \node[slab,above] at (0.95,2.60) {8};
  \foreach \p/\n/\d in {A/A/above, D/D/below left, G/G/below right}
    {\node[vlab,\d] at (\p) {$\n$};}
  \node[vlab,left=7pt] at (B) {$B$};
  \node[vlab,left=7pt] at (C) {$C$};
  \node[vlab,above right] at (E) {$E$};
  \node[vlab,above right] at (F) {$F$};
\end{tikzpicture}
\figcap{9--11}
\end{bkfigure}}

% ------------------------------------------------------------ Fig 9-12
% parallels l1,l2,l3 with gaps 2 : 3; a zigzag of three transversals.
% 2 and 3 are the pieces of the first, 4 the lower piece of the second,
% 6 the lower piece of the third; AB is the piece asked for.
% The numerals here are SEGMENT LENGTHS, not angle names -- the book strikes
% no arcs at them and neither do we (checked on the p.158 plate at 500 dpi).
% The "Parallel" note is the book's three-leader fan, not a double-headed span.
\newcommand{\FIGIXTWELVE}{\begin{bkfigure}[0.90\linewidth]
\begin{tikzpicture}[scale=1.15]
  \coordinate (T)  at (0,0);           % foot of transversal 1 on l3
  \coordinate (V1) at (1.189,1.886);   % apex, on l1
  \coordinate (V2) at (3.156,-0.303);  % on l3
  \coordinate (B)  at (6.411,1.3845);  % on l1
  \coordinate (P)  at ($(T)!0.6!(V1)$);
  \coordinate (Q)  at ($(V1)!0.4!(V2)$);
  \coordinate (A)  at ($(V2)!0.6!(B)$);
  \draw[given] (V1) -- (6.62,1.365);   % l1
  \draw[given] (P)  -- (6.60,0.566);   % l2
  \draw[given] (T)  -- (6.62,-0.635);  % l3
  \draw[given] (T) -- (V1) -- (V2) -- (B);
  \coordinate (PA) at (6.85,2.05);
  \node[klab] at (6.58,2.42) {Parallel};
  \draw[lead] (PA) -- (6.45,1.3813);
  \draw[lead] (PA) -- (6.45,0.5805);
  \draw[lead] (PA) -- (6.50,-0.6235);
  \node[slab,left=1pt]  at ($(P)!0.5!(V1)$) {2};
  \node[slab,left]  at ($(T)!0.5!(P)$)  {3};
  \node[slab] at ($($(Q)!0.5!(V2)$)+(0.17,0.15)$) {4};
  \node[slab] at ($(V2)!0.5!(A)+(-0.18,0.30)$) {6};
  \node[vlab] at ($(A)+(-0.22,0.28)$) {$A$};
  \node[vlab,above left] at (B) {$B$};
\end{tikzpicture}
\figcap{9--12}
\end{bkfigure}}

% ------------------------------------------------------------ Fig 9-13
% parallelogram ABCD, E and F the midpoints of AB and CD; the only diagonal
% drawn is AC, cut by ED and BF at its trisection points G and H.
% vertices measured off the p.158 figure at 400 dpi: AD runs 1.5 deg BELOW
% the horizontal and AB stands at 77 deg, giving the book's upright shape.
\newcommand{\FIGIXTHIRTEEN}{\begin{bkfigure}[0.66\linewidth]
\begin{tikzpicture}[scale=1.40]
  \coordinate (A) at (0,0); \coordinate (D) at (3.70,-0.10);
  \coordinate (B) at (0.60,2.55); \coordinate (C) at (4.30,2.45);
  \draw[given] (A) -- (D) -- (C) -- (B) -- cycle;
  \coordinate (E) at ($(A)!0.5!(B)$);
  \coordinate (F) at ($(C)!0.5!(D)$);
  \draw[given] (A) -- (C);
  \draw[given] (E) -- (D); \draw[given] (B) -- (F);
  \coordinate (G) at (intersection of A--C and E--D);
  \coordinate (H) at (intersection of A--C and B--F);
  \foreach \p/\n/\d in {A/A/below left, D/D/below right, C/C/above right,
                        B/B/above, E/E/left, F/F/right}
    {\node[vlab,\d] at (\p) {$\n$};}
  \node[vlab] at ($(G)+(0,0.42)$) {$G$};
  \node[vlab] at ($(H)+(0,0.42)$) {$H$};
\end{tikzpicture}
\figcap{9--13}
\end{bkfigure}}

% ------------------------------------------------------------ Fig 9-14
% D on LT and E on LR at the same ratio, so DE || TR; the braces mark
% x = LT, 1 = LD, y = LE, z = LR.  x : 1 = z : y is the point of the problem.
% REBUILT 2026-08-17 from the book photograph (feedback/small/20260817_102437):
% the two brace rows on each side were nearly on top of the side they measure,
% so the nesting could not be read, and the x brace was drawn INSIDE the 1
% brace (the book has the long brace outermost on both sides).  Offsets are
% now measured off the photograph: 0.07 and 0.24 of the side's own length.
% The "Parallel" leaders now fan from ONE apex at the word, as the book draws
% them, and the upper leader crosses the right side on its way in (so does
% the book's).
\newcommand{\FIGIXFOURTEEN}{\begin{bkfigure}[0.88\linewidth]
\begin{tikzpicture}[scale=1.30]
  \coordinate (L) at (0,0); \coordinate (R) at (4.66,0);
  \coordinate (T) at (2.16,2.12);
  \coordinate (D) at ($(L)!0.5!(T)$);
  \coordinate (E) at ($(L)!0.5!(R)$);
  \draw[given] (L) -- (R) -- (T) -- cycle;
  \draw[given] (D) -- (E);
  % outward unit normal of LT is (-0.7003, 0.7136); of LR it is (0,-1)
  % Drawn L -> T (not T -> L) so the brace teeth point AWAY from the side, out
  % towards their labels -- which is the way the base braces below already run,
  % and the way the book strikes all four.
  \draw[bkbrace] ($(L)+(-0.5042,0.5138)$) -- ($(T)+(-0.5042,0.5138)$);
  \draw[bkbrace] ($(L)+(-0.1541,0.1570)$) -- ($(D)+(-0.1541,0.1570)$);
  \draw[bkbracer] ($(L)+(0,-0.24)$) -- ($(E)+(0,-0.24)$);
  \draw[bkbracer] ($(L)+(0,-0.76)$) -- ($(R)+(0,-0.76)$);
  \node[slab] at (0.422,1.731) {$x$};
  \node[slab] at (0.232,0.844) {1};
  \node[slab] at (1.165,-0.52) {$y$};
  \node[slab] at (2.330,-1.04) {$z$};
  \coordinate (PA) at (3.62,1.28);
  \node[klab,right] at (3.74,1.30) {Parallel};
  \draw[lead] (PA) -- ($(D)!0.5!(E)$);
  \draw[lead] (PA) -- ($(T)!0.55!(R)$);
\end{tikzpicture}
\figcap{9--14}
\end{bkfigure}}

% --------------------------------------------------------------- Fig 9-15
% Sides 5, 6, 7 built exactly; the primed triangle is the same shape at 2/5.
\newcommand{\FIGIXFIFTEEN}{\begin{bkfigure}[0.86\linewidth]
\begin{tikzpicture}[scale=1.160]
  \coordinate (E) at (0,0); \coordinate (G) at (3.5,0);
  \coordinate (F) at (1.357,2.099);
  \draw[given] (E) -- (G) -- (F) -- cycle;
  \node[slab,left]  at ($(E)!0.5!(F)$) {5};
  \node[slab] at ($($(F)!0.5!(G)$)+(0.154,0.157)$) {6};
  \node[slab,below] at ($(E)!0.5!(G)$) {7};
  \node[vlab,below left] at (E) {$E$}; \node[vlab,below right] at (G) {$G$};
  \node[vlab,above] at (F) {$F$};
  \begin{scope}[xshift=5.1cm, yshift=0.02cm]
    \coordinate (Ep) at (0,0); \coordinate (Gp) at (1.4,0);
    \coordinate (Fp) at (0.543,0.840);
    \draw[given] (Ep) -- (Gp) -- (Fp) -- cycle;
    \node[slab,left] at ($(Ep)!0.5!(Fp)$) {2};
    \node[vlab,below left] at (Ep) {$E\IXpr$}; \node[vlab,below right] at (Gp) {$G\IXpr$};
    \node[vlab,above] at (Fp) {$F\IXpr$};
  \end{scope}
\end{tikzpicture}
\figcap{9--15}
\end{bkfigure}}

% --------------------------------------------------------------- Fig 9-16
\newcommand{\FIGIXSIXTEEN}{\begin{bkfigure}[0.9\linewidth]
\begin{tikzpicture}[scale=1.377]
  \draw[given] (0,0.62)   -- (2.75,0.62);
  \draw[given] (0.10,0)   -- (2.40,0);
  \draw[given] (4.35,0.62) -- (5.50,0.62);
  \draw[given] (4.65,0)   -- (5.25,0);
  \foreach \x/\y in {0/0.62, 2.75/0.62, 0.10/0, 2.40/0,
                     4.35/0.62, 5.50/0.62, 4.65/0, 5.25/0}
    {\dt{\x,\y}}
  \node[vlab,left]  at (0,0.62)    {$A$};
  \node[vlab,right] at (2.75,0.62) {$B$};
  \node[vlab,left]  at (0.10,0)    {$A\IXpr$};
  \node[vlab,right] at (2.40,0)    {$B\IXpr$};
  \node[vlab,left]  at (4.35,0.62) {$C$};
  \node[vlab,right] at (5.50,0.62) {$D$};
  \node[vlab,left]  at (4.65,0)    {$E$};
  \node[vlab,right] at (5.25,0)    {$F$};
\end{tikzpicture}
\figcap{9--16}
\end{bkfigure}}

% --------------------------------------------------------------- Fig 9-17
% CD cut into four congruent subintervals; EF is one of them.
\newcommand{\FIGIXSEVENTEEN}{\begin{bkfigure}[0.6\linewidth]
\begin{tikzpicture}[scale=1.958]
  \draw[given] (0,0) -- (2.28,0);
  \foreach \x in {0,0.57,1.14,1.71,2.28}{\dt{\x,0}}
  \node[vlab,left]  at (0,0)    {$C$};
  \node[vlab,right] at (2.28,0) {$D$};
  \node[vlab,above=2.6pt] at (0.57,0) {$E$};
  \node[vlab,above=2.6pt] at (1.14,0) {$F$};
\end{tikzpicture}
\figcap{9--17}
\end{bkfigure}}

% --------------------------------------------------------- Fig 9-18, 9-19
\newcommand{\FIGIXEIGHTEENNINETEEN}{\begin{bkfigure}[\linewidth]
\begin{minipage}{0.52\linewidth}\centering
\begin{tikzpicture}[scale=0.55]
  \draw[given] (0,0) -- (3.6,0) -- (1.6,2.2) -- cycle;
  \node[vlab,below left] at (0,0) {$A$}; \node[vlab,below right] at (3.6,0) {$C$};
  \node[vlab,above] at (1.6,2.2) {$B$};
  \begin{scope}[xshift=6.2cm, scale=0.55]
    \draw[given] (0,0) -- (3.6,0) -- (1.0,2.28) -- cycle;
    \node[vlab,below left] at (0,0) {$A\IXpr$}; \node[vlab,below right] at (3.6,0) {$C\IXpr$};
    \node[vlab,above] at (1.6,2.2) {$B\IXpr$};
  \end{scope}
\end{tikzpicture}
\figcap{9--18}
\end{minipage}\hfill
\begin{minipage}{0.42\linewidth}\centering
\begin{tikzpicture}[scale=0.57]
  % 9-19: convex pentagon ABCDE.  Vertices measured off the p.163 figure at
  % 400 dpi.  The cycle is A-B-C-D-E (an earlier draft joined E to B and D to
  % A, i.e. it swapped A and E in the path, which put A almost on top of E).
  \def\pentA{(0,1.38) (0.66,2.47) (2.25,1.67) (1.87,0.31) (0.05,0)}
  \draw[given] plot coordinates {\pentA} -- cycle;
  \node[vlab,left]  at (0,1.38)    {$A$};
  \node[vlab,above] at (0.66,2.47) {$B$};
  \node[vlab,right] at (2.25,1.67) {$C$};
  \node[vlab,below right] at (1.87,0.31) {$D$};
  \node[vlab,below left]  at (0.05,0)    {$E$};
  \begin{scope}[xshift=4.15cm, yshift=0.30cm, scale=0.62]
    \draw[given] plot coordinates {\pentA} -- cycle;
    \node[vlab,left]  at (0,1.38)    {$A\IXpr$};
    \node[vlab,above] at (0.66,2.47) {$B\IXpr$};
    \node[vlab,right] at (2.25,1.67) {$C\IXpr$};
    \node[vlab,below right] at (1.87,0.31) {$D\IXpr$};
    \node[vlab,below left]  at (0.05,0)    {$E\IXpr$};
  \end{scope}
\end{tikzpicture}
\figcap{9--19}
\end{minipage}
\end{bkfigure}}

% --------------------------------------------------------- Fig 9-20, 9-21
% SPLIT 2026-08-17: Ex. 9-8 cites 9-20 in #15 and 9-21 in #16.
\newcommand{\FIGIXTWENTY}{\begin{bkfigure}[0.54\linewidth]
\begin{tikzpicture}[scale=2.00]
  % 9-20: E, F midpoints of AB and BC.
  \coordinate (A) at (0,0); \coordinate (C) at (2.55,0);
  \coordinate (B) at (1.68,1.10);
  \draw[given] (A) -- (C) -- (B) -- cycle;
  \coordinate (E) at ($(A)!0.5!(B)$);
  \coordinate (F) at ($(B)!0.5!(C)$);
  \draw[given] (E) -- (F);
  \foreach \p/\n/\d in {A/A/below left, C/C/below right, B/B/above,
                        E/E/above left, F/F/above right}
    {\node[vlab,\d] at (\p) {$\n$};}
\end{tikzpicture}
\figcap{9--20}
\end{bkfigure}}

\newcommand{\FIGIXTWENTYONE}{\begin{bkfigure}[0.62\linewidth]
\begin{tikzpicture}[scale=1.95]
  % 9-21: A', B', C', D' are the midpoints of OA, OB, OC, OD.
  % Spoke lengths and directions measured off the p.164 figure at 400 dpi:
  % OA : OB : OC : OD = 428 : 190 : 275 : 200 px.  (An earlier draft had OD
  % 30% longer than OB; the book draws them nearly equal.)
  \coordinate (O) at (0,0);
  \coordinate (A) at (-1.72,0.01);  \coordinate (B) at (-0.05,0.76);
  \coordinate (C) at (1.105,0.07);  \coordinate (D) at (0.30,-0.75);
  \draw[given] (A) -- (B) -- (C) -- (D) -- cycle;
  \foreach \p in {A,B,C,D}{\draw[given] (O) -- (\p);
                           \coordinate (\p p) at ($(O)!0.5!(\p)$);}
  \draw[given] (Ap) -- (Bp) -- (Cp) -- (Dp) -- cycle;
  \node[vlab,left]  at (A) {$A$};   \node[vlab,above] at (B) {$B$};
  \node[vlab,right] at (C) {$C$};   \node[vlab,below] at (D) {$D$};
  \node[vlab] at ($(Ap)+(-0.08,0.19)$)  {$A\IXpr$};
  \node[vlab] at ($(Bp)+(-0.24,0.08)$)  {$B\IXpr$};
  \node[vlab] at ($(Cp)+(0.11,0.115)$)  {$C\IXpr$};
  \node[vlab] at ($(Dp)+(-0.22,-0.10)$) {$D\IXpr$};
  \node[vlab] at ($(O)+(0.13,0.14)$)    {$O$};
\end{tikzpicture}
\figcap{9--21}
\end{bkfigure}}

% --------------------------------------------------------------- Fig 9-22
% AB is cut into five unit pieces by F, A'', E, D; the parallels to AC through
% them meet BC in G, C'', H, I.  The primed triangle is the same shape at 3/5.
\newcommand{\FIGIXTWENTYTWO}{\begin{bkfigure}[0.9\linewidth]
\begin{tikzpicture}[scale=1.885]
  \coordinate (A) at (0,0); \coordinate (C) at (1.75,0);
  \coordinate (B) at (1.02,1.87);
  \draw[given] (A) -- (C) -- (B) -- cycle;
  \coordinate (F)   at ($(A)!1/5!(B)$);   \coordinate (G)   at ($(B)!4/5!(C)$);
  \coordinate (App) at ($(A)!2/5!(B)$);   \coordinate (Cpp) at ($(B)!3/5!(C)$);
  \coordinate (E)   at ($(A)!3/5!(B)$);   \coordinate (H)   at ($(B)!2/5!(C)$);
  \coordinate (D)   at ($(A)!4/5!(B)$);   \coordinate (I)   at ($(B)!1/5!(C)$);
  \draw[aux] (F) -- (G);
  \draw[given] (App) -- (Cpp);
  \draw[aux] (E) -- (H);
  \draw[aux] (D) -- (I);
  \node[vlab,left] at (F) {$F$};   \node[vlab,left] at (App) {$A\IXpr\IXpr$};
  \node[vlab,left] at (E) {$E$};   \node[vlab,left] at (D) {$D$};
  \node[vlab,right] at (G) {$G$};  \node[vlab,right] at (Cpp) {$C\IXpr\IXpr$};
  \node[vlab,right] at (H) {$H$};  \node[vlab,right] at (I) {$I$};
  \node[vlab,below left] at (A) {$A$}; \node[vlab,below right] at (C) {$C$};
  \node[vlab,above] at (B) {$B$};
  \begin{scope}[xshift=2.70cm, scale=0.6]
    \draw[given] (0,0) -- (1.75,0) -- (1.02,1.87) -- cycle;
    \node[vlab,below left] at (0,0) {$A\IXpr$};
    \node[vlab,below right] at (1.75,0) {$C\IXpr$};
    \node[vlab,above] at (1.02,1.87) {$B\IXpr$};
  \end{scope}
\end{tikzpicture}
\figcap{9--22}
\end{bkfigure}}

% ---------------------------------------------------------- Fig 9-23..9-29
% SPLIT 2026-08-17: Ex. 9-9 cites 9-23 in #1, 9-24 in #4, 9-25 in #5, 9-26 in
% #7, 9-27 in #10, 9-28 in #11 and 9-29 in #12; the book sets each beside the
% exercise that names it, so each now stands alone and is set inline.
%
% 9-23: AE = 4, ED = 6, so B and E sit at 0.4 along AC and AD; BE || CD.
% The book runs CD almost horizontally right and SLIGHTLY DOWN to D (the
% render had it dead level), puts E level with its own point rather than down
% by D, and strikes BE at the same weight as the sides.
\newcommand{\FIGIXTWENTYTHREE}{\begin{bkfigure}[0.56\linewidth]
\begin{tikzpicture}[scale=2.00]
  \coordinate (C) at (0,0); \coordinate (D) at (2.10,-0.10);
  \coordinate (A) at (2.58,1.32);
  \draw[given] (C) -- (D) -- (A) -- cycle;
  \coordinate (B) at ($(A)!0.4!(C)$);
  \coordinate (E) at ($(A)!0.4!(D)$);
  \draw[given] (B) -- (E);
  \foreach \p/\n/\d in {C/C/left, D/D/below, A/A/above,
                        B/B/above left, E/E/right}
    {\node[vlab,\d] at (\p) {$\n$};}
\end{tikzpicture}
\figcap{9--23}
\end{bkfigure}}

% 9-24: AB and DE both run north-south; D at 0.85 along AC, E directly below.
% The lake is traced from the iPad scan (scan12 p8) and placed against the
% triangle by matching its bounding box to the printed figure.
% 2026-08-17, against the p.166 plate at 500 dpi:
%   * the lake carries THREE contours, not four -- the innermost oval went;
%   * DE is a FINE DASHED segment in the book, not a heavy solid one;
%   * C was not actually ON the base line (the base ran from B to (3.265,0.075)
%     while C sat at y = 0.025), which is the "kink near B" Caleb saw.  The
%     base is now dead horizontal and AB dead vertical, as "run north-south"
%     requires, and C lies on it exactly.
\newcommand{\FIGIXTWENTYFOUR}{\begin{bkfigure}[0.62\linewidth]
\begin{tikzpicture}[scale=1.75]
  \coordinate (B) at (0,0); \coordinate (A) at (0,2.139);
  \coordinate (C) at (2.377,0);
  \coordinate (D) at ($(A)!0.85!(C)$);
  \coordinate (E) at (D |- B);
  \draw[given, line width=0.5pt]
    plot[smooth cycle, tension=0.7] coordinates {
      (1.636,2.336) (2.026,2.242) (2.307,2.055) (2.463,1.868) (2.588,1.727)
      (2.666,1.540) (2.791,1.384) (2.806,1.197) (2.775,0.994) (2.650,0.791)
      (2.494,0.635) (2.260,0.448) (2.026,0.339) (1.714,0.245) (1.480,0.191)
      (1.246,0.229) (1.012,0.339) (0.841,0.526) (0.669,0.822) (0.466,0.947)
      (0.334,1.150) (0.303,1.353) (0.334,1.509) (0.388,1.727) (0.513,1.914)
      (0.700,2.070) (0.934,2.195) (1.246,2.289)};
  \draw[given, line width=0.5pt]
    plot[smooth cycle, tension=0.7] coordinates {
      (1.558,2.070) (1.839,2.016) (2.058,1.922) (2.214,1.789) (2.292,1.634)
      (2.284,1.478) (2.370,1.337) (2.370,1.150) (2.292,0.978) (2.167,0.838)
      (1.995,0.729) (1.792,0.651) (1.574,0.604) (1.356,0.588) (1.153,0.620)
      (0.966,0.697) (0.810,0.822) (0.700,0.978) (0.622,1.150) (0.576,1.321)
      (0.544,1.493) (0.576,1.665) (0.669,1.821) (0.825,1.946) (1.028,2.024)
      (1.293,2.062)};
  \draw[given, line width=0.5pt]
    plot[smooth cycle, tension=0.7] coordinates {
      (1.558,1.556) (1.761,1.509) (1.855,1.384) (1.855,1.228) (1.792,1.087)
      (1.636,0.978) (1.480,0.900) (1.309,0.838) (1.137,0.830) (1.012,0.900)
      (0.934,1.010) (0.888,1.150) (0.864,1.290) (0.903,1.399) (1.012,1.478)
      (1.184,1.524) (1.371,1.553)};
  \draw[given] (A) -- (B) -- (3.30,0);
  \draw[given] (A) -- (C);
  \draw[finedash] (D) -- (E);
  \node[vlab,above left]  at (A) {$A$};
  \node[vlab,below left]  at (B) {$B$};
  \node[vlab,below] at ($(C)+(0.10,0)$) {$C$};
  % D's label sits in the clear band between contour 1 and contour 2, on the
  % same side as the book's.
  \node[vlab] at ($(D)+(0.07,0.27)$) {$D$};
  \node[vlab,below] at (E) {$E$};
\end{tikzpicture}
\figcap{9--24}
\end{bkfigure}}

% 9-25: D on BA and E on BC at the same ratio, so DE || AC.  Uniform weight:
% the book does NOT strike DE heavier than the sides (Caleb, 2026-08-17).
\newcommand{\FIGIXTWENTYFIVE}{\begin{bkfigure}[0.55\linewidth]
\begin{tikzpicture}[scale=2.10]
  \coordinate (A) at (0,0); \coordinate (C) at (2.40,0);
  \coordinate (B) at (1.30,1.42);
  \draw[given] (A) -- (C) -- (B) -- cycle;
  \coordinate (D) at ($(B)!0.53!(A)$);
  \coordinate (E) at ($(B)!0.53!(C)$);
  \draw[given] (D) -- (E);
  \foreach \p/\n/\d in {A/A/below left, C/C/right, B/B/above}
    {\node[vlab,\d] at (\p) {$\n$};}
  \node[vlab] at ($(D)+(-0.147,0.135)$) {$D$};
  \node[vlab] at ($(E)+(0.158,0.122)$)  {$E$};
\end{tikzpicture}
\figcap{9--25}
\end{bkfigure}}

% 9-26: AB _|_ AE at A and DE _|_ AE at E, on opposite sides, so BD crosses AE
% at C with AC/CE = AB/DE.  2026-08-17: AE now LEANS as the book draws it (the
% render had it dead vertical), DE is the short side and AB the long one in the
% book's 2.5 : 1 proportion, and C's label is lifted clear of the crossing --
% it was sitting on the lines.
\newcommand{\FIGIXTWENTYSIX}{\begin{bkfigure}[0.44\linewidth]
\begin{tikzpicture}[scale=2.80]
  \coordinate (A) at (0.18,1.60); \coordinate (E) at (0,0);
  \coordinate (B) at (1.0346,1.5038); \coordinate (D) at (-0.3379,0.0381);
  \draw[given] (A) -- (E); \draw[given] (A) -- (B); \draw[given] (E) -- (D);
  \draw[given] (B) -- (D);
  \coordinate (C) at (intersection of A--E and B--D);
  \foreach \p/\n/\d in {A/A/above left, B/B/above right, E/E/below right,
                        D/D/below left}
    {\node[vlab,\d] at (\p) {$\n$};}
  \node[vlab] at ($(C)+(-0.20,0.02)$) {$C$};
\end{tikzpicture}
\figcap{9--26}
\end{bkfigure}}

% 9-27: l and m are parallel (both of slope 2); EI and FH cross at G, so
% triangle EFG ~ triangle IHG.  2026-08-17: both parallels are drawn
% appreciably longer (they were stubs, so l || m could not be read), and the
% l/m labels now sit AT the ends of their own lines instead of floating past
% them.  The transversals are no longer heavier than the parallels.
\newcommand{\FIGIXTWENTYSEVEN}{\begin{bkfigure}[0.60\linewidth]
\begin{tikzpicture}[scale=1.70]
  \coordinate (F) at (0,0);    \coordinate (E) at (0.40,0.80);
  \coordinate (I) at (1.98,0); \coordinate (H) at (2.454,0.948);
  \draw[given] ($(F)!-0.70!(E)$) -- ($(F)!1.85!(E)$);
  \draw[given] ($(I)!-0.70!(H)$) -- ($(I)!1.80!(H)$);
  \draw[given] (E) -- (I); \draw[given] (F) -- (H);
  \coordinate (G) at (intersection of E--I and F--H);
  \node[vlab,below left]  at ($(F)!-0.70!(E)$) {$l$};
  \node[vlab,above right] at ($(I)!1.80!(H)$)  {$m$};
  \foreach \p/\n/\d in {I/I/below right, H/H/right, G/G/above}
    {\node[vlab,\d] at (\p) {$\n$};}
  \node[vlab] at ($(E)+(-0.42,0.12)$) {$E$};
  \node[vlab,left]       at (F) {$F$};
\end{tikzpicture}
\figcap{9--27}
\end{bkfigure}}

% 9-28: PS bisects angle P, so QS/SR = PQ/PR (ratio 0.468 for these sides).
% 2026-08-17 -- the figure Caleb called the worst offender.  The R label of the
% left triangle and the P' label of the right one printed jammed together; the
% book leaves a clear horizontal gap between the two triangles.  The gap is now
% 0.55 of a left-triangle unit (the book's is 0.17 of its WIDTH; ours needs
% more because our lettering is proportionally larger), the primed triangle is
% at the book's 0.62 rather than 0.583, and the whole block is drawn at nearly
% five times the old area.
\newcommand{\FIGIXTWENTYEIGHT}{\begin{bkfigure}[0.76\linewidth]
\begin{tikzpicture}[scale=1.90]
  \coordinate (P) at (0,0); \coordinate (R) at (1.80,0);
  \coordinate (Q) at (0.96,1.26);
  \draw[given] (P) -- (R) -- (Q) -- cycle;
  \coordinate (S) at ($(Q)!0.468!(R)$);
  \draw[given] (P) -- (S);
  \foreach \p/\n/\d in {P/P/below left, R/R/below right, Q/Q/above, S/S/right}
    {\node[vlab,\d] at (\p) {$\n$};}
  % primed triangle: same shape at 0.62, set 0.55 units clear of R
  \coordinate (Pp) at (2.35,0); \coordinate (Rp) at (3.466,0);
  \coordinate (Qp) at (2.9452,0.7812);
  \draw[given] (Pp) -- (Rp) -- (Qp) -- cycle;
  \coordinate (Sp) at ($(Qp)!0.468!(Rp)$);
  \draw[given] (Pp) -- (Sp);
  \foreach \p/\n/\d in {Pp/P\IXpr/below left, Rp/R\IXpr/below right,
                        Qp/Q\IXpr/above, Sp/S\IXpr/right}
    {\node[vlab,\d] at (\p) {$\n$};}
\end{tikzpicture}
\figcap{9--28}
\end{bkfigure}}

% 9-29: CD and AE are altitudes -- D and E are true perpendicular feet.
% The book strikes no right-angle squares at D or E; none are drawn.
\newcommand{\FIGIXTWENTYNINE}{\begin{bkfigure}[0.54\linewidth]
\begin{tikzpicture}[scale=2.60]
  \coordinate (A) at (0,0); \coordinate (C) at (1.88,0);
  \coordinate (B) at (1.15,1.45);
  \draw[given] (A) -- (C) -- (B) -- cycle;
  \coordinate (D) at ($(A)!(C)!(B)$);
  \coordinate (E) at ($(B)!(A)!(C)$);
  \draw[given] (C) -- (D); \draw[given] (A) -- (E);
  \coordinate (F) at (intersection of C--D and A--E);
  \foreach \p/\n/\d in {A/A/below left, C/C/below right, B/B/above,
                        D/D/above left, E/E/above right}
    {\node[vlab,\d] at (\p) {$\n$};}
  \node[vlab] at ($(F)+(0.00,0.36)$) {$F$};
\end{tikzpicture}
\figcap{9--29}
\end{bkfigure}}

% --------------------------------------------------------------- Fig 9-30
% AD/AB = AE/AC, so DE || BC.  DF is the auxiliary line of the proof, drawn
% to a point F beyond E on AC -- the proof's job is to show F = E.
\newcommand{\FIGIXTHIRTY}{\begin{bkfigure}[0.52\linewidth]
\begin{tikzpicture}[scale=1.522]
  \coordinate (B) at (0,0); \coordinate (C) at (2.30,0.20);
  \coordinate (A) at (1.20,1.30);
  \draw[given] (B) -- (C) -- (A) -- cycle;
  \coordinate (D) at ($(A)!0.50!(B)$);
  \coordinate (E) at ($(A)!0.50!(C)$);
  \coordinate (F) at ($(A)!0.68!(C)$);
  \draw[given] (D) -- (E);
  \draw[aux] (D) -- (F);
  \node[vlab] at ($(D)+(-0.162,0.149)$) {$D$};
  \foreach \p/\n/\d in {B/B/below left, C/C/right, A/A/above,
                        E/E/above right}
    {\node[vlab,\d] at (\p) {$\n$};}
  \node[vlab] at ($(F)+(0.36,0.03)$) {$F$};
\end{tikzpicture}
\figcap{9--30}
\end{bkfigure}}

% --------------------------------------------------------------- Fig 9-31
% B'' on AB and C'' on AC at the same ratio, so B''C'' || BC (Statement 4).
\newcommand{\FIGIXTHIRTYONE}{\begin{bkfigure}[0.84\linewidth]
\begin{tikzpicture}[scale=1.667]
  \coordinate (A) at (0,0); \coordinate (C) at (1.63,0);
  \coordinate (B) at (1.03,1.13);
  \draw[given] (A) -- (C) -- (B) -- cycle;
  \coordinate (Bpp) at ($(A)!0.46!(B)$);
  \coordinate (Cpp) at ($(A)!0.46!(C)$);
  \draw[given] (Bpp) -- (Cpp);
  \foreach \p/\n/\d in {A/A/left, C/C/right, B/B/above,
                        Bpp/B\IXpr\IXpr/above left, Cpp/C\IXpr\IXpr/below}
    {\node[vlab,\d] at (\p) {$\n$};}
  \begin{scope}[xshift=2.85cm, scale=0.54]
    \draw[given] (0,0) -- (1.63,0) -- (1.03,1.13) -- cycle;
    \node[vlab,below left] at (0,0) {$A\IXpr$};
    \node[vlab,below right] at (1.63,0) {$C\IXpr$};
    \node[vlab,above] at (1.03,1.13) {$B\IXpr$};
  \end{scope}
\end{tikzpicture}
\figcap{9--31}
\end{bkfigure}}

% ---------------------------------------------------------- Fig 9-32..9-35
% SPLIT 2026-08-17: Ex. 9-10 cites 9-32 in #2, 9-33 in #4, 9-34 in #5 and
% 9-35 in *#6, so each stands alone and is set inline after its exercise.
\newcommand{\FIGIXTHIRTYTWO}{\begin{bkfigure}[0.58\linewidth]
\begin{tikzpicture}[scale=1.55]
  % 9-32: O divides both AD and BC in the ratio 2 : 1, matching AO = 10,
  % OD = 5, BO = 12, OC = 6.
  % AB : AD = 0.27 in the book (measured p.169 at 400 dpi); an earlier draft
  % drew it at 0.44, which made the figure far too tall.
  \coordinate (O) at (1.35,0);
  \coordinate (A) at (0,-0.06);  \coordinate (B) at (-0.03,-0.60);
  \coordinate (D) at ($(A)!1.5!(O)$);
  \coordinate (C) at ($(B)!1.5!(O)$);
  \draw[given] (A) -- (B); \draw[given] (C) -- (D);
  \draw[given] (A) -- (D); \draw[given] (B) -- (C);
  \foreach \p/\n/\d in {A/A/above left, B/B/below left, C/C/above right,
                        D/D/right, O/O/above}
    {\node[vlab,\d] at (\p) {$\n$};}
\end{tikzpicture}
\figcap{9--32}
\end{bkfigure}}

\newcommand{\FIGIXTHIRTYTHREE}{\begin{bkfigure}[0.55\linewidth]
\begin{tikzpicture}[scale=2.00]
  % 9-33: AE = x.AC and AD = x.AB with the same x, so ED || BC.
  \coordinate (B) at (0,0); \coordinate (C) at (2.50,0);
  \coordinate (A) at (1.03,1.44);
  \draw[given] (B) -- (C) -- (A) -- cycle;
  \coordinate (E) at ($(A)!0.44!(B)$);
  \coordinate (D) at ($(A)!0.44!(C)$);
  \draw[given] (E) -- (D);
  \node[vlab] at ($(D)+(0.140,0.143)$)  {$D$};
  \node[vlab] at ($(E)+(-0.162,0.116)$) {$E$};
  \foreach \p/\n/\d in {B/B/below left, C/C/below right, A/A/above}
    {\node[vlab,\d] at (\p) {$\n$};}
\end{tikzpicture}
\figcap{9--33}
\end{bkfigure}}

\newcommand{\FIGIXTHIRTYFOUR}{\begin{bkfigure}[0.54\linewidth]
\begin{tikzpicture}[scale=2.20]
  % 9-34: D at 0.285 of AB, C at the height the right angle at C forces.
  % NO right-angle marks: in Ex. 9-10 #5 both "angle C is right" and
  % "CD perp AB" are what the reader has to PROVE, and the book draws the
  % figure bare (checked on the p.169 plate at 400 dpi).
  \coordinate (A) at (0,0); \coordinate (B) at (2.25,0);
  \coordinate (D) at (0.64,0); \coordinate (C) at (0.64,1.0156);
  \draw[given] (A) -- (B) -- (C) -- cycle;
  \draw[given] (C) -- (D);
  \foreach \p/\n/\d in {A/A/below left, B/B/below right, C/C/above, D/D/below}
    {\node[vlab,\d] at (\p) {$\n$};}
\end{tikzpicture}
\figcap{9--34}
\end{bkfigure}}

\newcommand{\FIGIXTHIRTYFIVE}{\begin{bkfigure}[0.54\linewidth]
\begin{tikzpicture}[scale=2.20]
  % 9-35: same construction with D at 0.712 of AB.  CD perp AB IS given here
  % (Ex. 9-10 *#6), but the book still prints no right-angle mark -- checked
  % on the p.169 plate at 400 dpi -- so none is drawn.
  \coordinate (A) at (0,0); \coordinate (B) at (2.25,0);
  \coordinate (D) at (1.602,0); \coordinate (C) at (1.602,1.0189);
  \draw[given] (A) -- (B) -- (C) -- cycle;
  \draw[given] (C) -- (D);
  \foreach \p/\n/\d in {A/A/below left, B/B/below right, C/C/above, D/D/below}
    {\node[vlab,\d] at (\p) {$\n$};}
\end{tikzpicture}
\figcap{9--35}
\end{bkfigure}}

% --------------------------------------------------------------- Fig 9-36
% A'' on BA and C'' on BC at the same ratio 0.60, matching the 0.60 scale of
% the primed triangle; hence A''C'' || AC.
\newcommand{\FIGIXTHIRTYSIX}{\begin{bkfigure}[0.84\linewidth]
\begin{tikzpicture}[scale=1.667]
  \coordinate (A) at (0,0); \coordinate (C) at (1.65,0);
  \coordinate (B) at (0.85,1.50);
  \draw[given] (A) -- (C) -- (B) -- cycle;
  \coordinate (App) at ($(B)!0.60!(A)$);
  \coordinate (Cpp) at ($(B)!0.60!(C)$);
  \draw[given] (App) -- (Cpp);
  \foreach \p/\n/\d in {A/A/below left, C/C/below right, B/B/above,
                        App/A\IXpr\IXpr/left, Cpp/C\IXpr\IXpr/right}
    {\node[vlab,\d] at (\p) {$\n$};}
  \begin{scope}[xshift=3.05cm, scale=0.60]
    \draw[given] (0,0) -- (1.65,0) -- (0.85,1.50) -- cycle;
    \node[vlab,below left] at (0,0) {$A\IXpr$};
    \node[vlab,below right] at (1.65,0) {$C\IXpr$};
    \node[vlab,above] at (0.85,1.50) {$B\IXpr$};
  \end{scope}
\end{tikzpicture}
\figcap{9--36}
\end{bkfigure}}

% ---------------------------------------------------------- Fig 9-37..9-39
% SPLIT 2026-08-17: Ex. 9-11 cites 9-37 in #2, 9-38 in #3 and 9-39 in #4.
\newcommand{\FIGIXTHIRTYSEVEN}{\begin{bkfigure}[0.58\linewidth]
\begin{tikzpicture}[scale=2.10]
  % 9-37: EA = 10, AC = 14, EB = 15, BD = 21, so EA/EC = EB/ED = 5/12 and
  % AB || CD.  A and B are placed at that ratio, not by eye.
  \coordinate (E) at (0,0); \coordinate (D) at (2.48,0);
  \coordinate (C) at (1.32,1.00);
  \coordinate (A) at ($(E)!5/12!(C)$);
  \coordinate (B) at ($(E)!5/12!(D)$);
  \draw[given] (E) -- (D) -- (C) -- cycle;
  \draw[given] (A) -- (B);
  \node[slab] at ($($(E)!0.5!(A)$)+(-0.133,0.175)$) {10};
  \node[slab] at ($($(A)!0.5!(C)$)+(-0.133,0.175)$) {14};
  \node[slab,below] at ($(E)!0.5!(B)$) {15};
  \node[slab,below] at ($(B)!0.5!(D)$) {21};
  \node[vlab,below left]  at (E) {$E$};  \node[vlab,right] at (D) {$D$};
  \node[vlab,above] at (C) {$C$};
  \node[vlab,above left=1pt] at (A) {$A$}; \node[vlab,below] at (B) {$B$};
\end{tikzpicture}
\figcap{9--37}
\end{bkfigure}}

\newcommand{\FIGIXTHIRTYEIGHT}{\begin{bkfigure}[0.74\linewidth]
% This paired diagram is wider than a single restored-book exercise column.
% Reduce only this exercise figure enough to clear the neighboring column;
% labels and line weights retain the shared book settings.
\begin{tikzpicture}[scale=1.05]
  % 9-38: one hexagon, then the same hexagon rotated -30 and scaled 0.72,
  % so the two really are similar.
  \def\hexA{(-0.87,0.58) (0.41,0.81) (0.79,0.06) (0.39,-0.65)
            (-0.09,-0.62) (-0.64,-0.18)}
  \draw[given] plot coordinates {\hexA} -- cycle;
  \node[vlab,left]  at (-0.87,0.58) {$A$};
  \node[vlab,above] at (0.41,0.81)  {$B$};
  \node[vlab,right] at (0.79,0.06)  {$C$};
  \node[vlab,right] at (0.39,-0.65) {$D$};
  \node[vlab,below] at (-0.09,-0.62){$E$};
  \node[vlab,left]  at (-0.64,-0.18){$F$};
  \begin{scope}[xshift=3.25cm]
    \begin{scope}[rotate=-30, scale=0.72]
      \draw[given] plot coordinates {\hexA} -- cycle;
      \coordinate (ha) at (-0.87,0.58);  \coordinate (hb) at (0.41,0.81);
      \coordinate (hc) at (0.79,0.06);   \coordinate (hd) at (0.39,-0.65);
      \coordinate (he) at (-0.09,-0.62); \coordinate (hf) at (-0.64,-0.18);
    \end{scope}
    \node[vlab,above left]  at (ha) {$A\IXpr$};
    \node[vlab,above right] at (hb) {$B\IXpr$};
    \node[vlab,below right] at (hc) {$C\IXpr$};
    \node[vlab,below]       at (hd) {$D\IXpr$};
    \node[vlab,below left]  at (he) {$E\IXpr$};
    \node[vlab,left]        at (hf) {$F\IXpr$};
  \end{scope}
\end{tikzpicture}
\figcap{9--38}
\end{bkfigure}}

% 9-39: BE _|_ AC with E on AC produced, CD _|_ AB with D the true foot.
% 2026-08-17, against the p.170 plate at 500 dpi: the right-angle square at D
% belongs in the wedge between DB and DC (it was drawn on the A side), and the
% dashed base runs a short way PAST E (it stopped dead at E).
\newcommand{\FIGIXTHIRTYNINE}{\begin{bkfigure}[0.58\linewidth]
\begin{tikzpicture}[scale=2.00]
  \coordinate (A) at (0,0); \coordinate (E) at (2.29,0);
  \coordinate (C) at (1.55,0); \coordinate (B) at (2.29,1.55);
  \draw[given] (A) -- (C) -- (B) -- cycle;
  \coordinate (D) at ($(A)!(C)!(B)$);
  \draw[given] (C) -- (D);
  \draw[aux] (C) -- (2.62,0); \draw[aux] (B) -- (E);
  \draw[fig, line width=0.5pt] (2.29,0.18) -- (2.11,0.18) -- (2.11,0);
  \draw[fig, line width=0.5pt]
    ($(D)!0.18!(B)$) -- ($($(D)!0.18!(B)$)!0.18!(C)$) -- ($(D)!0.18!(C)$);
  \foreach \p/\n/\d in {A/A/below left, B/B/above right, C/C/below,
                        D/D/above left, E/E/below right}
    {\node[vlab,\d] at (\p) {$\n$};}
\end{tikzpicture}
\figcap{9--39}
\end{bkfigure}}

% --------------------------------------------------------- Fig 9-40, 9-41
% Both are body figures (Theorem 9-11 and its proof), so they keep the book's
% side-by-side placement; only the drawing grows.
\newcommand{\FIGIXFORTYFORTYONE}{\begin{bkfigure}[\linewidth]
\begin{minipage}{0.46\linewidth}\centering
\begin{tikzpicture}[scale=1.32]
  % 9-40: right angle at C, D the foot of the altitude (CD^2 = AD.DB).
  \coordinate (A) at (0,0); \coordinate (B) at (2.67,0);
  \coordinate (D) at (1.95,0); \coordinate (C) at (1.95,1.1849);
  \draw[given] (A) -- (B) -- (C) -- cycle;
  \draw[given] (C) -- (D);
  % Both right angles ARE hypotheses of Theorem 9-11, and the book marks both.
  % The foot's box sits on the A side of CD (checked at 400 dpi), and the mark
  % at C is a full square, not an open chevron.
  \draw[rmark] (1.73,0) -- (1.73,0.22) -- (1.95,0.22);
  \draw[rmark] (1.762,1.0706) -- (1.876,0.8826) -- (2.064,0.9969);
  \foreach \p/\n/\d in {A/A/below left, B/B/below right, C/C/above, D/D/below}
    {\node[vlab,\d] at (\p) {$\n$};}
\end{tikzpicture}
\figcap{9--40}
\end{minipage}\hfill
\begin{minipage}{0.44\linewidth}\centering
\begin{tikzpicture}[scale=1.32]
  % 9-41: the same triangle with l', l, l'' through A, D(=C) and B.
  % The apex is labelled C, NOT C' -- Caleb checked the physical copy
  % 2026-08-17 and the p.171 plate at 500 dpi agrees.  The label sits in the
  % wedge between side AC and the dashed line l, just under the apex, and l is
  % named low down beside D exactly as the book names it.
  \coordinate (A) at (0,0); \coordinate (B) at (2.48,0);
  \coordinate (D) at (1.83,0); \coordinate (C) at (1.83,1.0906);
  \draw[given] (A) -- (B) -- (C) -- cycle;
  \draw[aux] (0,-0.38) -- (0,1.62);      \node[vlab,above] at (0,1.62) {$l\IXpr$};
  \draw[aux] (1.83,-0.38) -- (1.83,1.62);
  \draw[aux] (2.48,-0.38) -- (2.48,1.62); \node[vlab,above] at (2.48,1.62) {$l\IXpr\IXpr$};
  \node[vlab] at (1.70,0.26) {$l$};
  \node[vlab] at ($(C)+(-0.17,0.24)$) {$C$};
  \foreach \p/\n/\d in {A/A/below left, B/B/below right, D/D/below right}
    {\node[vlab,\d] at (\p) {$\n$};}
\end{tikzpicture}
\figcap{9--41}
\end{minipage}
\end{bkfigure}}

% ---------------------------------------------------------- Fig 9-42..9-45
% SPLIT 2026-08-17: Ex. 9-12 cites 9-42 in #1, 9-43 in #2, 9-44 in #3/#4 and
% 9-45 in #5, so each stands alone and is set inline after its exercise.
\newcommand{\FIGIXFORTYTWO}{\begin{bkfigure}[0.55\linewidth]
\begin{tikzpicture}[scale=1.90]
  % 9-42: EF = 8, FG = 5, right angle at F; H the foot, so EH : HG = 64 : 25.
  \coordinate (E) at (0,0); \coordinate (G) at (2.6412,0);
  \coordinate (H) at (1.8996,0); \coordinate (F) at (1.8996,1.1873);
  \draw[given] (E) -- (G) -- (F) -- cycle;
  \draw[given] (F) -- (H);
  \node[slab,above left]  at ($(E)!0.5!(F)$) {8};
  \node[slab,above right] at ($(F)!0.5!(G)$) {5};
  \foreach \p/\n/\d in {E/E/below left, G/G/below right, F/F/above, H/H/below}
    {\node[vlab,\d] at (\p) {$\n$};}
\end{tikzpicture}
\figcap{9--42}
\end{bkfigure}}

\newcommand{\FIGIXFORTYTHREE}{\begin{bkfigure}[0.55\linewidth]
\begin{tikzpicture}[scale=1.85]
  % 9-43: a true 5-12-13 right triangle; the altitude on the hypotenuse is
  % 60/13, its foot 25/13 from the short leg.
  \coordinate (A) at (0,0); \coordinate (B) at (2.73,0);
  \coordinate (D) at (0.4038,0); \coordinate (C) at (0.4038,0.9692);
  \draw[given] (A) -- (B) -- (C) -- cycle;
  \draw[given] (C) -- (D);
  \node[slab,left]  at ($(A)!0.5!(C)$) {5};
  \node[slab,above right] at ($(C)!0.5!(B)$) {12};
  \node[slab,below] at ($(A)!0.5!(B)$) {13};
  \node[slab,right] at (0.42,0.50) {$h$};
\end{tikzpicture}
\figcap{9--43}
\end{bkfigure}}

\newcommand{\FIGIXFORTYFOUR}{\begin{bkfigure}[0.55\linewidth]
\begin{tikzpicture}[scale=2.00]
  % 9-44: PS = 4, SR = 5; right angle at Q, QS _|_ PR, so QS = sqrt(20).
  \coordinate (R) at (0,0); \coordinate (P) at (2.52,0);
  \coordinate (S) at (1.40,0); \coordinate (Q) at (1.40,1.2522);
  \draw[given] (R) -- (P) -- (Q) -- cycle;
  \draw[aux] (Q) -- (S);
  \node[slab,below] at ($(R)!0.5!(S)$) {5};
  \node[slab,below] at ($(S)!0.5!(P)$) {4};
  \foreach \p/\n/\d in {R/R/below left, P/P/below right, Q/Q/above, S/S/below}
    {\node[vlab,\d] at (\p) {$\n$};}
\end{tikzpicture}
\figcap{9--44}
\end{bkfigure}}

\newcommand{\FIGIXFORTYFIVE}{\begin{bkfigure}[0.60\linewidth]
\begin{tikzpicture}[scale=1.90]
  % 9-45: right angle at C, right angle at D; b = 2.05, a = 1.50, c = 2.540,
  % so AD = b^2/c and h = ab/c hold exactly.
  \coordinate (A) at (0,0); \coordinate (B) at (2.540,0);
  \coordinate (D) at (1.6545,0); \coordinate (C) at (1.6545,1.2106);
  \draw[given] (A) -- (B) -- (C) -- cycle;
  \draw[aux] (C) -- (D);
  % The book braces ALL THREE sides here (b along AC, a along CB, c along AB)
  % and puts h on the B side of the dashed altitude.
  \draw[bkbrace]  (A) -- (C);
  \draw[bkbrace]  (C) -- (B);
  \draw[bkbracer] (A) -- (B);
  \node[slab] at (0.414,1.170) {$b$};
  \node[slab] at (2.662,1.018) {$a$};
  \node[slab] at (1.92,0.40) {$h$};
  \node[slab,below=3pt] at (1.00,-0.10) {$c$};
  \foreach \p/\n/\d in {A/A/left, B/B/right, C/C/above}
    {\node[vlab,\d] at (\p) {$\n$};}
  \node[vlab] at ($(D)+(-0.30,0.45)$) {$D$};
\end{tikzpicture}
\figcap{9--45}
\end{bkfigure}}

% --------------------------------------------------------------- Fig 9-46
% Right angle at C; D the foot of the altitude, x = AD, y = DB, x.y = h^2.
\newcommand{\FIGIXFORTYSIX}{\begin{bkfigure}[0.58\linewidth]
\begin{tikzpicture}[scale=1.522]
  \coordinate (A) at (0,0); \coordinate (B) at (2.75,0);
  \coordinate (D) at (1.05,0); \coordinate (C) at (1.05,1.3360);
  \draw[given] (A) -- (B) -- (C) -- cycle;
  \draw[aux] (C) -- (D);
  \node[slab,above left]  at ($(A)!0.5!(C)$) {$b$};
  \node[slab,above right] at ($(C)!0.5!(B)$) {$a$};
  \node[slab,above] at (0.44,0) {$x$};
  \node[slab,above] at (1.94,0) {$y$};
  \draw[bkbracer] (A) -- (B);
  \node[slab,below=3pt] at (1.62,-0.10) {$c$};
  \foreach \p/\n/\d in {A/A/left, B/B/right, C/C/above, D/D/above left}
    {\node[vlab,\d] at (\p) {$\n$};}
\end{tikzpicture}
\figcap{9--46}
\end{bkfigure}}

% --------------------------------------------------------------- Fig 9-47
% Two congruent right triangles, legs solid and hypotenuse dashed as printed.
\newcommand{\FIGIXFORTYSEVEN}{\begin{bkfigure}[0.72\linewidth]
\begin{tikzpicture}[scale=1.522]
  \coordinate (C) at (0,0); \coordinate (B) at (0.94,0);
  \coordinate (A) at (0,1.60);
  \draw[given] (C) -- (B); \draw[given] (C) -- (A); \draw[aux] (A) -- (B);
  \node[slab,left]  at ($(C)!0.5!(A)$) {$b$};
  \node[slab,below] at ($(C)!0.5!(B)$) {$a$};
  \node[slab,right] at (0.62,0.85) {$c$};
  \node[vlab,above] at (A) {$A$}; \node[vlab,below left] at (C) {$C$};
  \node[vlab,below right] at (B) {$B$};
  \begin{scope}[xshift=2.35cm]
    \coordinate (P) at (0,0); \coordinate (Q) at (0.94,0);
    \coordinate (R) at (0,1.60);
    \draw[given] (P) -- (Q); \draw[given] (P) -- (R); \draw[aux] (R) -- (Q);
    \node[slab,left]  at ($(P)!0.5!(R)$) {$b$};
    \node[slab,below] at ($(P)!0.5!(Q)$) {$a$};
    \node[vlab,above] at (R) {$R$}; \node[vlab,below left] at (P) {$P$};
    \node[vlab,below right] at (Q) {$Q$};
  \end{scope}
\end{tikzpicture}
\figcap{9--47}
\end{bkfigure}}

% --------------------------------------------------------- Fig 9-48, 9-49
% SPLIT 2026-08-17: Ex. 9-13 cites 9-48 in *#4 and 9-49 in #8.
%
% 9-48: SCHEMATIC, by Caleb's decision of 2026-08-17.  The earlier draft drew
% the diagonals at the exercise's own BD : AC = 24 : 10, which hands the reader
% the answer; the book draws an ordinary-looking rhombus at about 1.58 : 1 and
% lets the numbers do the work.  The three hypotheses the exercise DOES state
% -- AB || CD, BC || AD, AC _|_ BD -- still hold exactly: the figure is built
% from two perpendicular half-diagonals, which forces all four sides equal.
\newcommand{\FIGIXFORTYEIGHT}{\begin{bkfigure}[0.52\linewidth]
\begin{tikzpicture}[scale=3.20]
  % u = (0.450,-0.893) along AC ; v = (0.893,0.450) along BD ; |AC|/2 = 0.52,
  % |BD|/2 = 0.82, matching the book's proportions on the p.174 plate.
  \coordinate (A) at (-0.234, 0.464);  \coordinate (C) at (0.234,-0.464);
  \coordinate (B) at (-0.732,-0.369);  \coordinate (D) at (0.732, 0.369);
  \draw[given] (A) -- (D) -- (C) -- (B) -- cycle;
  \draw[given] (A) -- (C); \draw[given] (B) -- (D);
  \foreach \p/\n/\d in {A/A/above left, D/D/above right,
                        B/B/below left, C/C/below right}
    {\node[vlab,\d] at (\p) {$\n$};}
\end{tikzpicture}
\figcap{9--48}
\end{bkfigure}}

\newcommand{\FIGIXFORTYNINE}{\begin{bkfigure}[0.56\linewidth]
\begin{tikzpicture}[scale=2.20]
  % 9-49: OA = 21 north, OB = 20 east, right angle at O; AB dashed across the
  % lake.  The four contours are traced from the iPad scan (scan11 p8).
  \coordinate (O) at (0,0); \coordinate (A) at (0,2.10);
  \coordinate (B) at (2.00,0);
  \draw[given, line width=0.5pt]
    plot[smooth cycle, tension=0.7] coordinates {
      (1.025,2.024) (1.424,1.943) (1.833,1.736) (2.132,1.504) (2.309,1.246)
      (2.324,1.029) (2.242,0.781) (2.051,0.562) (1.804,0.452) (1.560,0.397)
      (1.312,0.372) (1.066,0.375) (0.848,0.431) (0.671,0.522) (0.520,0.659)
      (0.398,0.829) (0.306,1.029) (0.276,1.224) (0.306,1.422) (0.368,1.599)
      (0.464,1.777) (0.645,1.928) (0.823,2.002)};
  \draw[given, line width=0.5pt]
    plot[smooth cycle, tension=0.7] coordinates {
      (1.015,1.847) (1.339,1.777) (1.627,1.629) (1.859,1.449) (2.025,1.271)
      (2.107,1.079) (2.092,0.892) (1.996,0.740) (1.804,0.648) (1.586,0.603)
      (1.368,0.589) (1.150,0.611) (0.944,0.670) (0.767,0.765) (0.630,0.902)
      (0.528,1.068) (0.479,1.261) (0.516,1.437) (0.601,1.599) (0.782,1.750)};
  \draw[given, line width=0.5pt]
    plot[smooth cycle, tension=0.7] coordinates {
      (1.025,1.573) (1.312,1.496) (1.560,1.353) (1.738,1.191) (1.833,1.029)
      (1.812,0.877) (1.697,0.765) (1.519,0.710) (1.320,0.703) (1.136,0.748)
      (0.985,0.829) (0.884,0.950) (0.845,1.101) (0.870,1.267) (0.929,1.422)};
  \draw[given, line width=0.5pt]
    plot[smooth cycle, tension=0.7] coordinates {
      (1.150,1.271) (1.310,1.304) (1.428,1.267) (1.482,1.193) (1.457,1.120)
      (1.361,1.075) (1.247,1.056) (1.150,1.095) (1.099,1.164) (1.111,1.230)};
  \draw[given] (A) -- (O) -- (B);
  \draw[aux] (A) -- (B);
  \node[vlab,above] at (A) {$A$};
  \node[vlab,below left] at (O) {$O$};
  \node[vlab,right] at (B) {$B$};
\end{tikzpicture}
\figcap{9--49}
\end{bkfigure}}

% --------------------------------------------------------- Fig 9-50, 9-51
\newcommand{\FIGIXFIFTYFIFTYONE}{\begin{bkfigure}[\linewidth]
\begin{minipage}{0.38\linewidth}\centering
\begin{tikzpicture}[scale=1.45]
  % 9-50: C is the midpoint of AB, P at height sqrt(r^2 - (AB/2)^2); the two
  % dashed arcs are the circles of radius r about A and B.
  % REBUILT 2026-08-17 from the book photograph (feedback/small/20260817_102830):
  %   * the apex was at 0.65 of AB, the book's at 0.50, which made the lens far
  %     too fat (arcs cutting AB at 18%/82% instead of the book's 31%/69%) and
  %     drove three thick strokes into a blob at P that read as an arrowhead;
  %   * the arcs stopped dead at their intersections.  The book OVERSHOOTS each
  %     arc so a small X is struck at P and again at the lower crossing -- the
  %     distinctive mark of the construction.  Each arc now runs to +-57 deg,
  %     i.e. 12 deg past the +-45 deg where the two meet;
  %   * the arcs are drawn finer than the triangle, as in the book, and PC is
  %     no longer heavier than PA and PB.
  \coordinate (A) at (0,0); \coordinate (B) at (1.80,0);
  \coordinate (C) at (0.90,0); \coordinate (P) at (0.90,0.90);
  \draw[finedash] ([shift=(-57:1.27279)]A) arc (-57:57:1.27279);
  \draw[finedash] ([shift=(123:1.27279)]B) arc (123:237:1.27279);
  \draw[given] (A) -- (B);
  \draw[given] (C) -- ($(C)!0.985!(P)$);
  \draw[given] (A) -- (P) -- (B);
  \node[slab] at (0.337,0.563) {$r$};
  \foreach \p/\n/\d in {A/A/below left, B/B/below right, C/C/below}
    {\node[vlab,\d] at (\p) {$\n$};}
  \node[vlab] at (0.90,1.15) {$P$};   % in the gap between the two arc tails
\end{tikzpicture}
\figcap{9--50}
\end{minipage}\hfill
\begin{minipage}{0.55\linewidth}\centering
\begin{tikzpicture}[scale=1.60]
  % 9-51: C on the semicircle, D the foot of the perpendicular from C, with
  % x = AD = r - y, y = DO, z = CD = sqrt(r^2 - y^2), r the radius.
  \def\R{1.40}
  \coordinate (O) at (0,0);
  \coordinate (A) at (-\R,0); \coordinate (B) at (\R,0);
  \coordinate (D) at (-0.637,0); \coordinate (C) at (-0.637,1.2464);
  \draw[given] (A) -- (B);
  \draw[given] (A) arc (180:0:\R);
  \draw[given] (A) -- (C) -- (B);
  \draw[given] (C) -- (D);
  \draw[given] (C) -- (O);
  \dt{O}
  \node[vlab,below left]  at (A) {$A$};
  \node[vlab,below right] at (B) {$B$};
  \node[vlab,above] at (C) {$C$};
  \node[vlab,below] at (D) {$D$};
  \node[vlab,below] at (O) {$O$};
  \node[slab,above=1pt] at (-1.02,0) {$x$};
  \node[slab,above=1pt] at (-0.32,0) {$y$};
  \node[slab] at (-0.80,0.34) {$z$};
  \node[slab,right=1pt] at (-0.3185,0.6232) {$r$};
  \node[slab,above=1pt] at (0.70,0) {$r$};
\end{tikzpicture}
\figcap{9--51}
\end{minipage}
\end{bkfigure}}

% --------------------------------------------------------------- Fig 9-52
% AB a vertical diameter; CD _|_ AB at E, with EO = 12 and radius 15, so the
% half-chord is 9.
\newcommand{\FIGIXFIFTYTWO}{\begin{bkfigure}[0.56\linewidth]
\begin{tikzpicture}[scale=1.75]
  \draw[given] (0,0) circle (1.5);
  \coordinate (A) at (0,1.5); \coordinate (B) at (0,-1.5);
  \coordinate (E) at (0,1.2);
  \coordinate (C) at (-0.9,1.2); \coordinate (D) at (0.9,1.2);
  \draw[given] (A) -- (B);
  \draw[given] (C) -- (D);
  \dt{0,0}
  \node[vlab,above] at (A) {$A$};   \node[vlab,below] at (B) {$B$};
  \node[vlab] at ($(C)+(-0.30,0.18)$) {$C$};
  \node[vlab] at ($(D)+(0.30,0.18)$) {$D$};
  \node[vlab,below right] at (E) {$E$};
  \node[vlab,right] at (0,0) {$O$};
\end{tikzpicture}
\figcap{9--52}
\end{bkfigure}}

% ------------------------------------------ Figs 9-53/9-54, 9-55 (hoisted
% from ch09.tex local block after the Review-Exercise transcription pass)
% Figures 9-53, 9-54, 9-55 belong to the Review Exercises (book pp.178-179) and
% did not exist when figures09.tex was reviewed; that file is owned by the
% figure-review pass, so they are defined here per the "new local macro" rule.
% HOIST: move these three into figures09.tex and add their hypotheses to
% tools/constraints/ch09.py.  See chapters/ch09-UNCERTAIN.md sec. 3.
%
% 9-53: right trapezoid, bases 18 (bottom) and 12, altitude 8; the slant leg is
% 10, so its run is 6 (6-8-10), which is exactly the 18-12 difference.  The
% nonparallel sides extended meet at (18,24): the enclosing triangle is 18-24-30
% with the right angle on the vertical side, which is what Ex. 11 asks students
% to show.  Every vertex is forced by those numbers; nothing is placed by eye.
% SPLIT 2026-08-17: Review Ex. 11 cites 9-53, *Ex. 15 cites 9-54 and Ex. 18
% cites 9-55, so all three are set inline after their own exercise.
\newcommand{\FIGIXFIFTYTHREE}{\begin{bkfigure}[0.46\linewidth]
\begin{tikzpicture}[x=0.165cm,y=0.165cm]
  \coordinate (A) at (0,0);   \coordinate (B) at (18,0);
  \coordinate (P) at (18,24); \coordinate (D) at (6,8);
  \coordinate (C) at (18,8);  \coordinate (H) at (6,0);
  \draw[given] (A) -- (B) -- (P) -- cycle;
  \draw[given] (D) -- (C);
  \draw[given] (D) -- (H);
  \node[slab,left]  at ($($(A)!0.5!(D)$)+(-0.9,0.7)$) {10};
  \node[slab,right] at ($(D)!0.5!(H)$) {8};
  \node[slab,above] at ($(D)!0.5!(C)$) {12};
  \node[slab,below] at ($(A)!0.5!(B)$) {18};
\end{tikzpicture}
\figcap{9--53}
\end{bkfigure}}

% 9-54: Menelaus.  D is fixed at 2/3 of AB and F on ray AC beyond C; E is then
% the true intersection of DF with BC, so AD/DB . BE/EC . CF/FA = 1 holds
% exactly by construction.  X, Y, Z are projection feet, never eyeballed.
% REBUILT 2026-08-17 from the book photograph (feedback/small/20260817_102917).
% The render's triangle was tall and its apex central (apex 53% along AC,
% height 0.89 of the base); the book's is flat and wide with the apex about
% 39% along and 0.67 of the base high.  That mattered: at the old shape the
% dashed AX ran within 5 deg of side AB and the two were indistinguishable for
% a third of their length.  At the book's shape they diverge by 15 deg.  The
% base also rises slightly to the right, as the book's does, the l/X/D labels
% are pulled apart along and across the line, Y and Z now hug their own feet,
% and the dashed extension of AC uses a longer dash so it does not read as a
% fourth perpendicular.
\newcommand{\FIGIXFIFTYFOUR}{\begin{bkfigure}[0.92\linewidth]
\begin{tikzpicture}[scale=1.28]
  \coordinate (A) at (0,0); \coordinate (B) at (1.24,2.21);
  \coordinate (C) at (3.18,0.20);
  \coordinate (F) at ($(A)!1.70!(C)$);
  \coordinate (D) at ($(A)!2/3!(B)$);
  \coordinate (E) at (intersection of D--F and B--C);
  \coordinate (L) at ($(D)!-0.30!(F)$);
  \coordinate (G) at ($(D)!1.10!(F)$);
  \coordinate (X) at ($(D)!(A)!(F)$);
  \coordinate (Y) at ($(D)!(B)!(F)$);
  \coordinate (Z) at ($(D)!(C)!(F)$);
  \draw[given] (A) -- (B) -- (C) -- cycle;
  \draw[longdash] (C) -- ($(A)!1.80!(C)$);
  \draw[given] (L) -- (G);
  \draw[aux]   (A) -- (X);
  \draw[aux]   (B) -- (Y);
  \draw[aux]   (C) -- (Z);
  % Every letter is offset along l's own normal n = (0.2402, 0.9707): the ones
  % the book sets above the line get +n, the ones below get -n, so none of them
  % lands on l or on the perpendicular that meets it there.
  \node[vlab] at ($(L)+(-0.06,0.20)$) {$l$};
  % X and D printed as one blob "XD" in the render Caleb photographed, so they
  % are pushed apart ALONG l as well as across it: 0.43 units centre to centre.
  \node[vlab] at ($(X)+(-0.044,0.238)$) {$X$};
  \node[vlab] at ($(D)+(-0.063,0.384)$) {$D$};
  \node[vlab] at ($(Y)+(-0.058,-0.233)$) {$Y$};
  \node[vlab] at ($(E)+(0.058,0.233)$)  {$E$};
  \node[vlab] at ($(Z)+(0.053,0.214)$) {$Z$};
  \node[vlab,above]      at (B) {$B$};
  \node[vlab,below left] at (A) {$A$};
  \node[vlab,below]      at (C) {$C$};
  \node[vlab,below]      at (F) {$F$};
\end{tikzpicture}
\figcap{9--54}
\end{bkfigure}}

% 9-55: D and E are the feet of the altitudes from C and from B, placed with the
% projection operator, so the two right angles are exact.
% 2026-08-17, against the p.179 plate at 500 dpi: the right-angle square at E
% was struck at 0.16 of |EA| and 0.16 of |EB| -- i.e. proportional to sides
% four times longer than the ones at D -- so it printed far larger than D's and
% hung into the triangle.  Both squares are now 0.17 units on a side, as the
% book draws them, and E's sits in the wedge between EC and EB (the book's
% side).  The triangle is also flattened from 0.85 to the book's 0.72.
\newcommand{\FIGIXFIFTYFIVE}{\begin{bkfigure}[0.72\linewidth]
\begin{tikzpicture}[scale=1.90]
  \coordinate (A) at (0,0); \coordinate (B) at (3.43,0);
  \coordinate (C) at (2.05,2.46);
  \coordinate (D) at ($(A)!(C)!(B)$);
  \coordinate (E) at ($(A)!(B)!(C)$);
  \draw[given] (A) -- (B) -- (C) -- cycle;
  \draw[given] (C) -- (D);
  \draw[given] (B) -- (E);
  % D = (2.05, 0); E = (1.40575, 1.68689) with unit vectors E->C (0.64018,
  % 0.76821) and E->B (0.76821, -0.64019).  Both squares are 0.17 a side.
  \draw[tick] (1.88,0) -- (1.88,0.17) -- (2.05,0.17);
  \draw[tick] (1.51458,1.81749) -- (1.64518,1.70866) -- (1.53635,1.57806);
  \node[vlab,above]       at (C) {$C$};
  \node[vlab,below left]  at (A) {$A$};
  \node[vlab,below right] at (B) {$B$};
  \node[vlab,below]       at (D) {$D$};
  \node[vlab,above left]  at (E) {$E$};
\end{tikzpicture}
\figcap{9--55}
\end{bkfigure}}

% figures10.tex -- Chapter 10 (AREA) figures redrawn in TikZ.
% Proportions read off scans/scan13.pdf p1-p16 (book pp.180-192).
% Every constrained point is constructed, never eyeballed; the stated
% hypotheses are mirrored in tools/constraints/ch10.py.

% ------------------------------------------------------------ Fig 10-1, 10-2
% 10-1: unit-square counting -- 1x1, 2x1, 3x2.
% 10-2: half-unit square inside a unit square; 2.5 x 3 grid.
\newcommand{\FIGXONETWO}{\begin{bkfigure}[\linewidth]
\begin{minipage}{0.52\linewidth}\centering
\begin{tikzpicture}[scale=0.706]
  % (a) single unit square.  Panels (a) and (b) sit on a band 0.5 units above
  % the old one: their "below" labels used to run into the top edge of the
  % 3x2 grid of panel (c), which the two panels overlap in x.
  \draw[fig] (0,3.1) rectangle (1,4.1);
  \node[slab] at (0.5,3.6) {1};
  \node[slab,left] at (0,3.6) {1};  \node[slab,below] at (0.5,3.1) {1};
  % (b) two unit squares side by side
  \draw[fig] (2.4,3.1) rectangle (4.4,4.1);
  \draw[fig] (3.4,3.1) -- (3.4,4.1);
  \node[slab] at (2.9,3.6) {1}; \node[slab] at (3.9,3.6) {1};
  \node[slab,left] at (2.4,3.6) {1}; \node[slab,right] at (4.4,3.6) {1};
  \node[slab,below] at (2.9,3.1) {1}; \node[slab,below] at (3.9,3.1) {1};
  % (c) 3 x 2 grid
  \draw[fig] (0.7,0) rectangle (3.7,2);
  \draw[fig] (1.7,0) -- (1.7,2); \draw[fig] (2.7,0) -- (2.7,2);
  \draw[fig] (0.7,1) -- (3.7,1);
  \foreach \x in {1.2,2.2,3.2}{\node[slab] at (\x,1.5) {1};
                               \node[slab] at (\x,0.5) {1};
                               \node[slab,below] at (\x,0) {1};}
  \node[slab,left] at (0.7,1.5) {1}; \node[slab,left] at (0.7,0.5) {1};
\end{tikzpicture}
\figcap{10--1}
\end{minipage}\hfill
\begin{minipage}{0.46\linewidth}\centering
\begin{tikzpicture}[scale=0.706]
  % (a) half-unit square inside a dashed unit square.  The book quarters the
  % unit square: a dashed mid-vertical and mid-horizontal run right across it,
  % and only the lower-left quarter is inked solid.  The two dashed halves of
  % those medians were missing.
  \draw[aux] (0,1.8) rectangle (2,3.8);
  \draw[aux] (1,2.8) -- (1,3.8);          % upper half of the mid-vertical
  \draw[aux] (1,2.8) -- (2,2.8);          % right half of the mid-horizontal
  \draw[fig] (0,1.8) rectangle (1,2.8);
  \node[slab,left] at (0,3.3) {$\tfrac{1}{2}$};
  \node[slab,left] at (0,2.3) {$\tfrac{1}{2}$};
  % (b) 2.5 wide x 3 tall grid
  \draw[fig] (3.1,0) rectangle (5.6,3);
  \draw[fig] (4.1,0) -- (4.1,3); \draw[fig] (5.1,0) -- (5.1,3);
  \draw[fig] (3.1,1) -- (5.6,1); \draw[fig] (3.1,2) -- (5.6,2);
  \node[slab,left] at (3.1,2.5) {1}; \node[slab,left] at (3.1,1.5) {1};
  \node[slab,left] at (3.1,0.5) {1};
  \node[slab,below] at (3.6,0) {1}; \node[slab,below] at (4.6,0) {1};
  \node[slab,below] at (5.35,0) {$\tfrac{1}{2}$};
\end{tikzpicture}
\figcap{10--2}
\end{minipage}
\end{bkfigure}}

% ---------------------------------------------------------------- Fig 10-3
% 1.4 by 2.3 rectangle, partly chopped into squares 0.1 on a side.
\newcommand{\FIGXTHREE}{\begin{bkfigure}[0.62\linewidth]
\begin{tikzpicture}[scale=2.015]
  \draw[fig] (0,0) rectangle (2.3,1.4);
  \draw[fig] (1,0) -- (1,1.4);            % 1-unit divisions
  \draw[fig] (2,0) -- (2,1.4);
  \draw[fig] (0,1) -- (2.3,1);            % the 0.4 band
  \foreach \y in {1.1,1.2,1.3} \draw[fig] (0,\y) -- (2.3,\y);
  \foreach \x in {2.1,2.2} \draw[fig] (\x,0) -- (\x,1.4);
  \node[slab,left] at (0,1.2) {0.4};  \node[slab,left] at (0,0.5) {1};
  \node[slab,below] at (0.5,0) {1};   \node[slab,below] at (1.5,0) {1};
  \node[slab,below] at (2.15,0) {0.3};
\end{tikzpicture}
\figcap{10--3}
\end{bkfigure}}

% ---------------------------------------------------------------- Fig 10-4
% rectangle sqrt(2) by 1, the part past the whole unit divided into tenths.
% Structure measured off scan13 p3 by tracking every vertical (solid lines
% are 100% inked down the page, dashed ones 76-83%):
%   left edge (solid) -- unit divider at exactly 1 unit (solid) -- five
%   DASHED tenth-marks -- and the right edge, at sqrt(2) units, falls
%   BETWEEN the fourth and fifth marks.  That is the whole point of the
%   figure: it shows 1.4 < sqrt(2) < 1.5, which the text asserts two lines
%   above it.  The old drawing put six dashed marks inside and ended the
%   rectangle on a mark, so the reading was lost.  The top and bottom edges
%   run past the right edge to the last mark, as they do in the book.
\newcommand{\FIGXFOUR}{\begin{bkfigure}[0.62\linewidth]
\begin{tikzpicture}[scale=2.015]
  % one unit = 1.6; the book's tenth-marks are slightly compressed (0.089
  % unit apart) so that the sqrt(2) edge clears them -- keep that.
  \def\uu{1.6}\def\tt{0.1425}\def\ww{2.2627}\def\hh{1.91}% w = sqrt2 * 1.6
  \def\xe{2.3125}%                                        last tenth-mark
  \draw[fig] (0,0) -- (\xe,0);                          % bottom, extended
  \draw[fig] (0,\hh) -- (\xe,\hh);                      % top, extended
  \draw[fig] (0,0) -- (0,\hh);                          % left edge
  \draw[fig] (\ww,0) -- (\ww,\hh);                      % right edge = sqrt2
  \draw[fig] (\uu,0) -- (\uu,\hh);                      % the whole unit
  \foreach \k in {1,2,3,4,5}
    \draw[aux] ({\uu+\k*\tt},0) -- ({\uu+\k*\tt},\hh);
  \node[slab] at (0.81,1.74) {1};
  \node[slab,left] at (0,0.955) {1};
  \node[slab] at (1.42,2.18) {0.1};
  \draw[fig,->] (1.60,2.10) -- (1.70,1.97);             % leader into a strip
  % dimension line under the rectangle; the label sits in a break of it,
  % as the book draws it.
  \draw[fig,|<->|] (0,-0.40) -- node[slab,fill=white,inner sep=1.5pt]
        {$\sqrt{2}$} (\ww,-0.40);
\end{tikzpicture}
\figcap{10--4}
\end{bkfigure}}

% ---------------------------------------------------------------- Fig 10-5
% (a) isosceles triangle ABC and rectangle on the same base;
% (b) parallelogram A'B'C'D' with A'B' = D'E';
% (c) two parallelograms on the base A''D''.
\newcommand{\FIGXFIVE}{\begin{bkfigure}[\linewidth]
\begin{minipage}[b]{0.31\linewidth}\centering
\begin{tikzpicture}[scale=0.62]
  % height read off scan13 p4: h : BC = 570 : 535, i.e. 1.07 (was 0.73).
  \coordinate (B) at (0,0); \coordinate (C) at (3,0);
  \coordinate (E) at ($(B)!0.5!(C)$);
  \coordinate (A) at ($(E)+(0,3.2)$);
  \coordinate (D) at ($(B)+(0,3.2)$);
  \draw[aux] (D) -- (A); \draw[aux] (D) -- (B);
  \draw[fig] (B) -- (C) -- (A) -- cycle;
  \draw[fig] (A) -- (E);
  % the book marks AB and AC with DOUBLE ticks
  \foreach \s/\t in {A/B, A/C}
    {\foreach \f in {0.465,0.535}
       {\draw[tick] ($(\s)!\f!(\t)!2.2pt!90:(\t)$) --
                    ($(\s)!\f!(\t)!2.2pt!-90:(\t)$);}}
  \foreach \p/\n/\d in {A/A/above, B/B/below left, C/C/below right,
                        D/D/above left, E/E/below}
    {\node[vlab,\d] at (\p) {$\n$};}
  \node[slab,below=20pt] at (1.5,0) {(a)};
\end{tikzpicture}
\end{minipage}\hfill
\begin{minipage}[b]{0.31\linewidth}\centering
\begin{tikzpicture}[scale=0.62]
  % re-measured on scan13 p4: A'B' rises at 61 degrees (the old 49 made the
  % parallelogram far too flat), |A'B'| : |A'D'| = 0.54, and A'D' drops 6.7
  % degrees left to right.
  \coordinate (Ap) at (0,0.75); \coordinate (Bp) at (0.675,1.98);
  \coordinate (Dp) at (2.586,0.446);
  \coordinate (Cp) at ($(Dp)+(Bp)-(Ap)$);
  \coordinate (Ep) at ($(Dp)-(Bp)+(Ap)$);
  \draw[fig] (Ap) -- (Bp) -- (Cp) -- (Dp) -- cycle;
  \draw[aux] (Bp) -- (Dp); \draw[aux] (Ap) -- (Ep); \draw[aux] (Ep) -- (Dp);
  \foreach \p/\n/\d in {Ap/A'/left, Bp/B'/above left, Cp/C'/above right,
                        Dp/D'/right, Ep/E'/below}
    {\node[vlab,\d] at (\p) {$\n$};}
  \node[slab,below=20pt] at (1.3,-0.5) {(b)};
\end{tikzpicture}
\end{minipage}\hfill
\begin{minipage}[b]{0.31\linewidth}\centering
\begin{tikzpicture}[scale=0.62]
  % Re-measured on scan13 p4 (vertices, not label boxes): height = 1.46 x
  % base (not 1.94), B'' at 0.42 of the base, E'' at 0.91.  Line weights
  % corrected too: the book inks A''B''C''D'' and the top edge B''..F'',
  % and draws the SECOND parallelogram's slant sides A''E'' and D''F''
  % dashed, plus the cut E''D''.  There is no A''C'' diagonal in the book;
  % the old drawing had one, and had A''E''F''D'' solid.
  \coordinate (Aq) at (0,0);    \coordinate (Dq) at (1.6,0);
  \coordinate (Bq) at (0.67,2.33);
  \coordinate (Cq) at ($(Bq)+(Dq)-(Aq)$);
  \coordinate (Eq) at (1.46,2.33);
  \coordinate (Fq) at ($(Eq)+(Dq)-(Aq)$);
  \draw[fig] (Aq) -- (Bq);            % first parallelogram, left side
  \draw[fig] (Bq) -- (Fq);            % top edge, through E'' and C''
  \draw[fig] (Dq) -- (Cq);            % first parallelogram, right side
  \draw[fig] (Aq) -- (Dq);            % common base
  \draw[aux] (Aq) -- (Eq); \draw[aux] (Dq) -- (Fq);
  \draw[aux] (Eq) -- (Dq);
  \foreach \p/\n/\d in {Aq/A''/below left, Dq/D''/below right,
                        Bq/B''/above left, Cq/C''/above,
                        Eq/E''/above, Fq/F''/above right}
    {\node[vlab,\d] at (\p) {$\n$};}
  \node[slab,below=20pt] at (0.8,0) {(c)};
\end{tikzpicture}
\end{minipage}
\figcap{10--5}
\end{bkfigure}}

% ------------------------------------------------------------ Fig 10-6, 10-7
\newcommand{\FIGXSIXSEVEN}{\begin{bkfigure}[\linewidth]
\begin{minipage}{0.5\linewidth}\centering
\begin{tikzpicture}[scale=0.868]
  % the shared side DE is drawn dead vertical in the book (scan13 p5);
  % the other six vertices are re-measured off that page relative to D and
  % the length of DE.  A had been drawn too close to E, which crowded the
  % two labels; in the book A sits well to the left of E.
  \coordinate (D) at (1.6,0.3);  \coordinate (E) at (1.6,1.75);
  \coordinate (A) at (0.464,2.389); \coordinate (B) at (-0.779,1.378);
  \coordinate (C) at (0.365,-0.197);
  \coordinate (F) at (4.004,2.389); \coordinate (G) at (4.335,-0.347);
  \draw[fig] (A) -- (B) -- (C) -- (D) -- (G) -- (F) -- (E) -- cycle;
  \draw[aux] (D) -- (E);
  % E is lettered straight ABOVE its vertex in the book; anchoring the label
  % "left" put it on the side EA, which runs up and to the left from E.
  \foreach \p/\n/\d in {A/A/above, B/B/left, C/C/below left, D/D/below,
                        E/E/above, F/F/above right, G/G/below right}
    {\node[vlab,\d] at (\p) {$\n$};}
\end{tikzpicture}
\figcap{10--6}
\end{minipage}\hfill
\begin{minipage}{0.48\linewidth}\centering
\begin{tikzpicture}[scale=0.868]
  % a : b = 0.43 on scan13 p5 (was 0.53)
  \coordinate (C) at (0,0); \coordinate (A) at (3,0);
  \coordinate (B) at (0,1.3); \coordinate (Cp) at (3,1.3);
  \draw[aux] (B) -- (Cp) -- (A);
  \draw[fig] (C) -- (A) -- (B) -- cycle;
  \draw[rmark] (0,0.28) -- (0.28,0.28) -- (0.28,0);
  \node[slab,left] at (0,0.65) {$a$};
  \node[slab,below] at (1.5,0) {$b$};
  \foreach \p/\n/\d in {B/B/above left, C/C/below left, A/A/below right,
                        Cp/C'/above right}
    {\node[vlab,\d] at (\p) {$\n$};}
\end{tikzpicture}
\figcap{10--7}
\end{minipage}
\end{bkfigure}}

% ---------------------------------------------------------------- Fig 10-8
% parallelogram ABCD (base b = AD, height h) inside rectangle AECF; EB = x.
\newcommand{\FIGXEIGHT}{\begin{bkfigure}[0.66\linewidth]
% Proportions read off scan13 p6: x : b = 217 : 137 (1.58) and h : b = 143 : 137
% (1.04).  The old drawing had x : b = 0.4 and h : b = 0.57 -- far too little
% slant, which makes the x-strip of the proof almost invisible.
\begin{tikzpicture}[scale=1.105]
  \coordinate (A) at (0,0);   \coordinate (D) at (2,0);
  \coordinate (B) at (3.1,2.05);
  \coordinate (C) at ($(B)+(D)-(A)$);
  \coordinate (E) at (0,2.05); \coordinate (F) at (5.1,0);
  \draw[aux] (E) -- (C); \draw[aux] (A) -- (F); \draw[aux] (E) -- (A);
  \draw[aux] (C) -- (F);
  \draw[fig] (A) -- (B) -- (C) -- (D) -- cycle;
  \draw[rmark] (0,1.77) -- (0.28,1.77) -- (0.28,2.05);
  \draw[rmark] (5.1,0.28) -- (4.82,0.28) -- (4.82,0);
  \node[slab,above] at (1.55,2.05) {$x$};
  \node[slab,below] at (1,0) {$b$};
  \node[slab,right] at (5.1,1.02) {$h$};
  \foreach \p/\n/\d in {E/E/above left, B/B/above, C/C/above right,
                        A/A/below left, D/D/below, F/F/below right}
    {\node[vlab,\d] at (\p) {$\n$};}
\end{tikzpicture}
\figcap{10--8}
\end{bkfigure}}

% ---------------------------------------------------------------- Fig 10-9
% triangle of base b and height h, completed to a parallelogram.
\newcommand{\FIGXNINE}{\begin{bkfigure}[0.58\linewidth]
\begin{tikzpicture}[scale=1.235]
  % re-measured on scan13 p7: apex at 0.56 of the base, height 0.54 x base
  \coordinate (B) at (0,0); \coordinate (C) at (2.6,0);
  \coordinate (A) at (1.46,1.41);
  \coordinate (Bp) at ($(A)+(C)-(B)$);
  \coordinate (H) at ($(B)!(A)!(C)$);
  \draw[aux] (A) -- (Bp) -- (C);
  \draw[fig] (B) -- (C) -- (A) -- cycle;
  \draw[fig] (A) -- (H);
  \node[slab,left] at ($(A)!0.5!(H)$) {$h$};
  \node[slab,below] at (1.75,0) {$b$};
  \foreach \p/\n/\d in {A/A/above left, Bp/B'/above right, B/B/below left}
    {\node[vlab,\d] at (\p) {$\n$};}
\end{tikzpicture}
\figcap{10--9}
\end{bkfigure}}

% ---------------------------------------------------------- Fig 10-10, 10-11
\newcommand{\FIGXTENELEVEN}{\begin{bkfigure}[\linewidth]
\begin{minipage}{0.5\linewidth}\centering
% Ex. Group 10-5 no. 5 states l || m, 3 ft apart, and AB = CD = 5 ft; the book
% draws it to that scale, with C directly over B and the zigzag standing for
% the 1-mile gap from B to C.  The old drawing had AB = 1.4 and CD = 1.3 (not
% congruent) and a 1.07 : 1 gap-to-AB ratio instead of 3 : 5.  One tikz unit
% is now one foot.
\begin{tikzpicture}[scale=0.40]
  \draw[fig] (0,3) -- (11.6,3);  \node[vlab,right] at (11.6,3) {$l$};
  \draw[fig] (0,0) -- (11.6,0);  \node[vlab,right] at (11.6,0) {$m$};
  \coordinate (A) at (0.6,0);  \coordinate (B) at (5.6,0);
  \coordinate (C) at (5.6,3);  \coordinate (D) at (10.6,3);
  \draw[fig] (C) -- (6.05,2.45) -- (5.15,1.95) -- (6.05,1.45)
             -- (5.15,0.95) -- (5.9,0.45) -- (B);
  \foreach \p in {A,B,C,D} {\dt{\p}}
  \node[vlab,above] at (C) {$C$};  \node[vlab,above] at (D) {$D$};
  \node[vlab,below] at (A) {$A$};  \node[vlab,below] at (B) {$B$};
\end{tikzpicture}
\figcap{10--10}
\end{minipage}\hfill
\begin{minipage}{0.47\linewidth}\centering
\begin{tikzpicture}[scale=0.48]
  \coordinate (A) at (0,0);  \coordinate (D) at (4,0);
  \coordinate (B) at (4,3);  \coordinate (C) at ($(B)+(D)-(A)$);
  \draw[fig] (A) -- (B) -- (C) -- (D) -- cycle;
  \draw[fig] (B) -- (D);
  \draw[rmark] (4,0.4) -- (3.6,0.4) -- (3.6,0);
  \draw[rmark] (4,2.6) -- (4.4,2.6) -- (4.4,3);
  % "above left" grazed AB; step 0.7 off the midpoint along AB's normal
  \node[slab] at ($(A)!0.5!(B)!0.55!90:(B)$) {5};
  \node[slab,below] at ($(A)!0.5!(D)$) {4};
  \foreach \p/\n/\d in {A/A/below left, D/D/below, B/B/above left,
                        C/C/above right}
    {\node[vlab,\d] at (\p) {$\n$};}
\end{tikzpicture}
\figcap{10--11}
\end{minipage}
\end{bkfigure}}

% ---------------------------------------------------------- Fig 10-12, 10-13
\newcommand{\FIGXTWELVETHIRTEEN}{\begin{bkfigure}[\linewidth]
\begin{minipage}{0.5\linewidth}\centering
\begin{tikzpicture}[scale=0.734]
  \coordinate (B) at (0,0);    \coordinate (C) at (4.5,0);
  \coordinate (A) at (1,3);    \coordinate (D) at (3.5,3);
  \draw[fig] (B) -- (C) -- (D) -- (A) -- cycle;
  \coordinate (T) at (2.25,3);
  \draw[fig] (T) -- (2.25,0);
  \draw[aux] (A) -- (C);
  % plain "above" left the label's box touching the top side; lift it clear
  \node[slab,above] at ($($(A)!0.5!(D)$)+(0,0.28)$) {$b' = 10$};
  \node[slab,below] at ($(B)!0.5!(C)$) {$b = 18$};
  % the book prints "h = 12" INSIDE the trapezoid, in the panel between the
  % left side and the altitude, not out in the margin (scan13 p8)
  \node[slab] at (1.40,1.5) {$h = 12$};
  \foreach \p/\n/\d in {A/A/above left, D/D/above right, B/B/below left,
                        C/C/below right}
    {\node[vlab,\d] at (\p) {$\n$};}
\end{tikzpicture}
\figcap{10--12}
\end{minipage}\hfill
\begin{minipage}{0.47\linewidth}\centering
\begin{tikzpicture}[scale=0.568]
  % measured on scan13 p8: b' : b = 0.563, h : b = 0.72 (the old drawing had
  % 0.5 and 0.5).  Both b' and b - b' are printed INSIDE the trapezoid just
  % above the base, and the b of the dimension line sits in a break of it.
  \coordinate (B) at (0,0);   \coordinate (C) at (6,0);
  \coordinate (A) at (1.2,4.3);
  \coordinate (D) at ($(A)+(3.38,0)$);
  \coordinate (P) at (3.38,0);
  \draw[fig] (B) -- (C) -- (D) -- (A) -- cycle;
  \draw[aux] (D) -- (P);
  \node[slab,above] at ($($(A)!0.5!(D)$)+(0,0.30)$) {$b'$};
  \node[slab] at ($($(B)!0.5!(P)$)+(0,0.50)$) {$b'$};
  \node[slab] at ($($(P)!0.5!(C)$)+(0.15,0.50)$) {$b - b'$};
  \draw[fig,|<->|] (0,-1.3) -- node[slab,fill=white,inner sep=1.5pt]
        {$b$} (6,-1.3);
\end{tikzpicture}
\figcap{10--13}
\end{minipage}
\end{bkfigure}}

% --------------------------------------------------------------- Fig 10-14
\newcommand{\FIGXFOURTEEN}{\begin{bkfigure}[0.52\linewidth]
\begin{tikzpicture}[scale=1.118]
  % base : height = 1.75 and lean : height = 0.26 on scan13 p9
  \coordinate (B) at (0,0); \coordinate (C) at (3,0);
  \coordinate (A) at (0.45,1.72);
  \coordinate (D) at ($(A)+(C)-(B)$);
  \draw[fig] (A) -- (B) -- (C) -- (D) -- cycle;
  \draw[fig] (A) -- (C); \draw[fig] (B) -- (D);
  \coordinate (O) at (intersection of A--C and B--D);
  % "above right" put the label along diagonal BD.  The book sets O straight
  % above the crossing, in the wedge between OA and OD.
  \node[vlab,above] at ($(O)+(0,0.16)$) {$O$};
  \foreach \p/\n/\d in {A/A/above left, D/D/above right, B/B/below left,
                        C/C/below right}
    {\node[vlab,\d] at (\p) {$\n$};}
\end{tikzpicture}
\figcap{10--14}
\end{bkfigure}}

% --------------------------------------------------------------- Fig 10-15
\newcommand{\FIGXFIFTEEN}{\begin{bkfigure}[0.52\linewidth]
\begin{tikzpicture}[scale=1.170]
  % apex at 0.30 of the base, h : b = 0.50 (scan13 p9)
  \coordinate (A) at (0,0); \coordinate (C) at (3.4,0);
  \coordinate (B) at (1.0,1.7);
  \draw[fig] (A) -- (C) -- (B) -- cycle;
  \node[slab,above left]  at ($(A)!0.5!(B)$) {$c$};
  \node[slab,above right] at ($(B)!0.5!(C)$) {$a$};
  \node[slab,below]       at ($(A)!0.5!(C)$) {$b$};
  \foreach \p/\n/\d in {A/A/below left, C/C/below right, B/B/above}
    {\node[vlab,\d] at (\p) {$\n$};}
\end{tikzpicture}
\figcap{10--15}
\end{bkfigure}}

% ---------------------------------------------------------- Fig 10-16, 10-17
\newcommand{\FIGXSIXTEENSEVENTEEN}{\begin{bkfigure}[\linewidth]
\begin{minipage}{0.5\linewidth}\centering
\begin{tikzpicture}[scale=0.873]
  % Re-measured off scan13 p9 at 150 dpi (vertices A(232,65) D(700,82)
  % B(145,342) C(620,335)): h : BC = 0.58 and lean : h = 0.31.  An earlier
  % pass had 0.78 / 0.20, which stretched the parallelogram upright.
  \coordinate (B) at (0,0);   \coordinate (C) at (3.5,0);
  \coordinate (A) at (0.64,2.04);
  \coordinate (D) at ($(A)+(C)-(B)$);
  \coordinate (E) at ($(B)!0.25!(D)$);     % BE : ED = 2 : 6
  \draw[fig] (A) -- (B) -- (C) -- (D) -- cycle;
  \draw[fig] (B) -- (D);
  \draw[aux] (E) -- (C);
  \foreach \p/\n/\d in {A/A/above left, D/D/above right, B/B/below left,
                        C/C/below right, E/E/above left}
    {\node[vlab,\d] at (\p) {$\n$};}
\end{tikzpicture}
\figcap{10--16}
\end{minipage}\hfill
\begin{minipage}{0.48\linewidth}\centering
\begin{tikzpicture}[scale=0.786]
  \coordinate (A) at (0,0); \coordinate (B) at (4.333,0);
  \coordinate (C) at (0.641,1.538);        % AC = 5, CB = 12, AB = 13 (/3)
  \coordinate (D) at ($(A)!(C)!(B)$);
  \draw[fig] (A) -- (B) -- (C) -- cycle;
  \draw[fig] (C) -- (D);
  % No right-angle box: the book draws neither the one at C nor the one at D
  % here, though it marks right angles freely in 10-7, 10-8, 10-11 and 10-20.
  % Both facts are stated in the exercise text instead (scan13 p9).
  \foreach \p/\n/\d in {A/A/below left, B/B/right, C/C/above, D/D/below}
    {\node[vlab,\d] at (\p) {$\n$};}
\end{tikzpicture}
\figcap{10--17}
\end{minipage}
\end{bkfigure}}

% ---------------------------------------------------------- Fig 10-18, 10-19
% 10-18: squares on the three sides of a right triangle.
% 10-19: the (a+b) square dissected into four right triangles and a c-square.
\newcommand{\FIGXEIGHTEENNINETEEN}{\begin{bkfigure}[\linewidth]
\begin{minipage}{0.5\linewidth}\centering
\begin{tikzpicture}[scale=1.187]
  \coordinate (R) at (0,0); \coordinate (S) at (0,1.0); \coordinate (T) at (1.4,0);
  \draw[fig] (R) -- (S) -- (T) -- cycle;
  \draw[fig] (R) -- (-1.0,0) -- (-1.0,1.0) -- (S);          % square 1
  \draw[fig] (R) -- (0,-1.4) -- (1.4,-1.4) -- (T);          % square 2
  \coordinate (N) at (1.0,1.4);                              % normal (s1,s2)
  \draw[fig] (S) -- ($(S)+(N)$) -- ($(T)+(N)$) -- (T);       % square 3
  \node[slab] at (-0.5,0.5) {1};
  \node[slab] at (0.7,-0.7) {2};
  \node[slab] at (1.2,1.2)  {3};
\end{tikzpicture}
\figcap{10--18}
\end{minipage}\hfill
\begin{minipage}{0.48\linewidth}\centering
\begin{tikzpicture}[scale=0.736]
  % In the book (scan13 p11, close-up) the SHORT piece of every side is the
  % one lettered a and the long piece is b: going clockwise from the top-left
  % corner the top edge reads a then b, the right edge a then b, the bottom
  % edge b then a, the left edge b then a.  The letters were in the right
  % places here but the two lengths were interchanged, which mirrored the
  % inner square's tilt.  Swapping the values fixes both at once.
  \def\aa{1.6}\def\bb{2.4}\def\ss{4.0}
  \coordinate (P1) at (\bb,0);   \coordinate (P2) at (\ss,\bb);
  \coordinate (P3) at (\aa,\ss); \coordinate (P4) at (0,\aa);
  \draw[fig] (0,0) rectangle (\ss,\ss);
  \draw[fig] (P1) -- (P2) -- (P3) -- (P4) -- cycle;
  \node[slab,below] at (\bb/2,0)        {$b$};
  \node[slab,below] at ({(\bb+\ss)/2},0){$a$};
  \node[slab,right] at (\ss,\bb/2)      {$b$};
  \node[slab,right] at (\ss,{(\bb+\ss)/2}) {$a$};
  \node[slab,above] at (\aa/2,\ss)      {$a$};
  \node[slab,above] at ({(\aa+\ss)/2},\ss) {$b$};
  \node[slab,left]  at (0,\aa/2)        {$a$};
  \node[slab,left]  at (0,{(\aa+\ss)/2}){$b$};
  \node[slab] at (2.83,1.45) {$c$};
  \node[slab] at (2.55,2.83) {$c$};
  \node[slab] at (1.17,2.55) {$c$};
  \node[slab] at (1.45,1.17) {$c$};
\end{tikzpicture}
\figcap{10--19}
\end{minipage}
\end{bkfigure}}

% --------------------------------------------------------------- Fig 10-20
% Bhaskara's dissection: square BAED of side c, four right triangles
% (legs a, b) round a central square of side a - b.  DF || AC, EG || BC.
\newcommand{\FIGXTWENTY}{\begin{bkfigure}[0.56\linewidth]
\begin{tikzpicture}[scale=1.495]
  \coordinate (B) at (0,0);   \coordinate (A) at (3,0);
  \coordinate (E) at (3,3);   \coordinate (D) at (0,3);
  \coordinate (C)  at (1.92,1.44);     % right angle at C, BC = a, CA = b
  \coordinate (Cp) at (1.56,1.92);
  \coordinate (G)  at (1.08,1.56);
  \coordinate (F)  at (1.44,1.08);
  \draw[aux] (B) -- (D) -- (E) -- (A);
  \draw[fig] (B) -- (A);
  \draw[fig] (B) -- (C) -- (A);
  % The central square is NOT inked as a solid quadrilateral in the book
  % (scan13 p14, close-up): three of its sides fall along lines already drawn
  % -- GF lies on DF, C'G lies on EG, and FC lies on the solid leg BC -- and
  % only CC' is added, dashed.  Drawing the square solid put ink on the page
  % that the source does not have and broke the given/auxiliary convention.
  \draw[aux] (D) -- (F); \draw[aux] (E) -- (G);
  \draw[aux] (C) -- (Cp);
  % right-angle marks at G and at C', as the book has them
  \draw[rmark] (1.256,1.692) -- (1.388,1.516) -- (1.212,1.384);
  \draw[rmark] (1.384,1.788) -- (1.516,1.612) -- (1.692,1.744);
  % "a" stepped 0.25 units off BC along its normal (it was grazing the leg
  % at 1.2pt); the book likewise floats it clear, above the midpoint of BC.
  \node[slab] at (0.81,0.92) {$a$};
  \node[slab] at (2.62,0.95) {$b$};
  \node[slab,below] at (1.5,0) {$c$};
  % C is lettered up and to the RIGHT of the vertex in the book; a plain
  % "right" anchor slid the letter down onto the leg CA (0.5pt clearance).
  \foreach \p/\n/\d in {B/B/below left, A/A/below right, D/D/above left,
                        E/E/above right, C/C/above right, F/F/below right}
    {\node[vlab,\d] at (\p) {$\n$};}
  \node[vlab] at (1.38,2.10) {$G$};
\end{tikzpicture}
\figcap{10--20}
\end{bkfigure}}

% --------------------------------------------------------------- Fig 10-21
% quadrilateral ABCD; DE drawn parallel to the diagonal AC, E on line BC.
\newcommand{\FIGXTWENTYONE}{\begin{bkfigure}[0.6\linewidth]
\begin{tikzpicture}[scale=1.235]
  % A and D re-measured on scan13 p14 against BC: A at (0.38, 0.68) and
  % D at (0.97, 0.83) in units of BC.  The quadrilateral had been drawn
  % noticeably taller than the book's.
  \coordinate (B) at (0,0);   \coordinate (C) at (2.2,0);
  \coordinate (A) at (0.84,1.50); \coordinate (D) at (2.13,1.83);
  \coordinate (Dp) at ($(D)+(C)-(A)$);
  \coordinate (Xf) at (6,0);
  \coordinate (E) at (intersection of D--Dp and B--Xf);
  % the base is drawn as the RAY BC extended, arrowhead at E (scan13 p14)
  \draw[fig,->] (B) -- (E);
  \draw[fig] (B) -- (A) -- (D) -- (C);
  \draw[aux] (A) -- (C); \draw[aux] (A) -- (E); \draw[aux] (D) -- (E);
  \foreach \p/\n/\d in {A/A/above left, D/D/above, B/B/below left,
                        C/C/below, E/E/below right}
    {\node[vlab,\d] at (\p) {$\n$};}
\end{tikzpicture}
\figcap{10--21}
\end{bkfigure}}

% ---------------------------------------------------------- Fig 10-22, 10-23
\newcommand{\FIGXTWENTYTWOTHREE}{\begin{bkfigure}[\linewidth]
\begin{minipage}{0.5\linewidth}\centering
\begin{tikzpicture}[scale=1.048]
  % vertices measured off scan13 p15 and rescaled on AB: the book's D sits
  % almost straight above A and C just outside B, giving a much broader top
  % than the drawing had.
  \coordinate (A) at (0.9,0);   \coordinate (B) at (2.6,0);
  \coordinate (C) at (2.74,2.52); \coordinate (D) at (0.96,2.78);
  \coordinate (E) at (0.08,1.29);
  \draw[fig] (A) -- (B) -- (C) -- (D) -- (E) -- cycle;
  \foreach \p/\n/\d in {A/A/below left, B/B/below right, C/C/above right,
                        D/D/above left, E/E/left}
    {\node[vlab,\d] at (\p) {$\n$};}
\end{tikzpicture}
\figcap{10--22}
\end{minipage}\hfill
\begin{minipage}{0.48\linewidth}\centering
\begin{tikzpicture}[scale=0.839]
  \coordinate (B) at (0,0); \coordinate (C) at (4,0);
  \coordinate (A) at (2,3.464);            % equilateral
  \coordinate (P) at (1.9,1.0);
  \coordinate (L) at ($(B)!(P)!(C)$);
  \coordinate (M) at ($(A)!(P)!(C)$);
  \coordinate (N) at ($(A)!(P)!(B)$);
  \draw[fig] (A) -- (B) -- (C) -- cycle;
  \draw[aux] (P) -- (A); \draw[aux] (P) -- (B); \draw[aux] (P) -- (C);
  \draw[aux] (P) -- (L); \draw[aux] (P) -- (M); \draw[aux] (P) -- (N);
  \foreach \p/\n/\d in {A/A/above, B/B/below left, C/C/below right,
                        L/L/below, M/M/right, N/N/left}
    {\node[vlab,\d] at (\p) {$\n$};}
  % the book sets P in the wedge between PL and PC; at our trim the label has
  % to sit further out along that wedge's bisector to clear both dashed lines
  \node[vlab] at (2.19,0.54) {$P$};
\end{tikzpicture}
\figcap{10--23}
\end{minipage}
\end{bkfigure}}

% --------------------------------------------------------------- Fig 10-24
\newcommand{\FIGXTWENTYFOUR}{\begin{bkfigure}[0.46\linewidth]
\begin{tikzpicture}[scale=1.066]
  % scan13 p16, deskewed by the 4.1 degree page rotation: apex at 0.56 of
  % the base, height 0.72 x base (it had been a symmetric 0.5 / 0.5).
  \coordinate (A) at (2.13,2.74); \coordinate (B) at (0,0);
  \coordinate (C) at (3.8,0);
  \draw[fig] (A) -- (B) -- (C) -- cycle;
  \draw[aux] (A) -- ($(B)!0.5!(C)$);
  \draw[aux] (B) -- ($(A)!0.5!(C)$);
  \draw[aux] (C) -- ($(A)!0.5!(B)$);
\end{tikzpicture}
\figcap{10--24}
\end{bkfigure}}

% figures11.tex -- Chapter 11 figures, built 2026-08-17 from the iPad scans.
%
% Every figure below was measured off a scan crop (scan13 pp.19-23, scan14
% pp.3,4,5,7) rather than eyeballed from the photo PDF.  Constrained points
% are DERIVED, never placed by hand:
%   regular polygons -> polar coordinates (k*360/n), never by eye
%   perpendicular bisectors -> (intersection of ...) of two rotated midlines
%   midpoints / apothem feet -> ($(X)!0.5!(Y)$)
%   tangency  -> radius drawn to the point, line laid perpendicular to it
% so the hypothesis stated in the book holds exactly in the drawing.
% Hypotheses are re-asserted numerically in tools/constraints/ch11.py.
%
% Side counts were measured, not guessed.  Fig 11-5's polygon is a HEXAGON:
% fitting the drawn circumscribed and inscribed circles from one common centre
% gives a/R = 0.8646 against cos(pi/6) = 0.8660 (n = 5.97); a heptagon would
% need 0.9010.  See chapters/ch11-UNCERTAIN.md.
%
% Macro names are chapter-prefixed: \FIGXI<NUMBERWORD>.
%
% Line treatment follows the book's own convention: solid where the book
% prints solid (given), dashed where the book prints dashed (aux).

\usetikzlibrary{decorations.pathreplacing,patterns,arrows.meta}

% Primed labels.  The math prime sits ~4pt above the letter, which makes the
% collision audit read it as a separate word overlapping its own letter; a
% 1.3pt drop keeps the pair in one word and is visually indistinguishable.
\newcommand{\XIpr}{\raisebox{-1.3pt}{\ensuremath{{}^{\prime}}}}

% Chapter-local label clearance, as in ch09: a labelled point lying ON a line
% needs a little more outer sep than brumfiel.sty's 3.4pt or the rule clips
% the corner of the label box.
\tikzset{
  vlab/.append style={outer sep=4.8pt},
  slab/.append style={outer sep=4.8pt},
  klab/.append style={outer sep=4.8pt},
  bkarr/.style={fig, -{Stealth[length=3.6pt,width=3pt]}},
  bkarrb/.style={fig, {Stealth[length=3.6pt,width=3pt]}-{Stealth[length=3.6pt,width=3pt]}},
  % Fig 11-6's four callouts are drawn in the book as ARROWS from the word to
  % the part of the polygon it names, so the leader carries a head at the
  % geometry end (the end of the path, since every leader starts at its node).
  lead/.style={fig, line width=0.4pt, -{Stealth[length=3.2pt,width=2.6pt]}},
}

% ------------------------------------------------------- Figs 11-1 and 11-2
% 11-1: circumcentre of a triangle.  m and n are the perpendicular bisectors
% of AB and BC; O is their intersection, hence equidistant from A, B and C.
% AB is drawn horizontal, so m is vertical, exactly as the book prints it.
% 11-2: regular hexagon inscribed in a circle, vertices at k*60 deg with
% C at 0 deg.  Angles 1,2 sit at B and 3,4 at C, each on the bisector of its
% own sub-angle (the 0.33 construction reproduces the book's placement).
\newcommand{\FIGXIONETWO}{\begin{bkfigure}[\linewidth]
\begin{minipage}{0.52\linewidth}\centering
\begin{tikzpicture}[scale=0.86]
  \coordinate (A) at (0,0);
  \coordinate (B) at (2.20,0);
  \coordinate (C) at (4.45,2.00);
  \coordinate (Mab) at ($(A)!0.5!(B)$);
  \coordinate (Mbc) at ($(B)!0.5!(C)$);
  % a second point on each perpendicular bisector, by 90-degree rotation
  \coordinate (mP) at ($(Mab)!1!90:(B)$);
  \coordinate (nP) at ($(Mbc)!1!90:(C)$);
  \coordinate (O)  at (intersection of Mab--mP and Mbc--nP);
  \draw[given] (A) -- (B) -- (C) -- cycle;
  \draw[aux] (O) -- ($(Mab)!-0.10!(O)$);
  \draw[aux] (O) -- ($(Mbc)!-0.18!(O)$);
  \node[vlab, below left] at (A) {$A$};
  \node[vlab, below] at (B) {$B$};
  \node[vlab, above right] at (C) {$C$};
  \node[vlab, above] at (O) {$O$};
  \node[vlab, below] at ($(Mab)!-0.10!(O)$) {$m$};
  \node[vlab, below right] at ($(Mbc)!-0.18!(O)$) {$n$};
\end{tikzpicture}
\figcap{11--1}
\end{minipage}\hfill
\begin{minipage}{0.44\linewidth}\centering
\begin{tikzpicture}[scale=0.86]
  \def\R{1.55}
  \coordinate (O) at (0,0);
  \coordinate (C) at (0:\R);   \coordinate (D) at (60:\R);
  \coordinate (E) at (120:\R); \coordinate (F) at (180:\R);
  \coordinate (A) at (240:\R); \coordinate (B) at (300:\R);
  \draw[given] (O) circle (\R);
  \draw[given] (C) -- (D) -- (E) -- (F) -- (A) -- (B) -- cycle;
  \foreach \p in {A,B,C,D,E} {\draw[given] (O) -- (\p);}
  \fill (O) circle (1.3pt);
  \node[vlab, above left]  at (E) {$E$};
  \node[vlab, above right] at (D) {$D$};
  \node[vlab, right]       at (C) {$C$};
  \node[vlab, below left]  at (A) {$A$};
  \node[vlab, below right] at (B) {$B$};
  \node[vlab, left]        at (O) {$O$};
  \node[slab] at ($($(B)!0.33!(A)$)!0.5!($(B)!0.33!(O)$)$) {1};
  \node[slab] at ($($(B)!0.33!(O)$)!0.5!($(B)!0.33!(C)$)$) {2};
  \node[slab] at ($($(C)!0.33!(B)$)!0.5!($(C)!0.33!(O)$)$) {3};
  \node[slab] at ($($(C)!0.33!(O)$)!0.5!($(C)!0.33!(D)$)$) {4};
\end{tikzpicture}
\figcap{11--2}
\end{minipage}
\end{bkfigure}}

% ------------------------------------------------------- Figs 11-3 and 11-4
% 11-3: l is laid perpendicular to the radius OP at P (direction = OP angle
% + 90), so the tangency hypothesis holds exactly.  A is any other point of l.
% 11-4: l is NOT perpendicular to OP.  m is the perpendicular from O to l,
% A its foot, and B the reflection of P in A -- so AB = AP and OB = OP, which
% is what puts B on the circle.
\newcommand{\FIGXITHREEFOUR}{\begin{bkfigure}[\linewidth]
\begin{minipage}{0.46\linewidth}\centering
% Both halves at 1.30, not 0.86: in 11-4 the letter A has to sit in the wedge
% between OP and l, and at the original scale that wedge was narrower than the
% 9pt letter box.  Same scale on 11-3 so the pair still matches in size.
\begin{tikzpicture}[scale=1.30]
  \def\R{1.25}
  \coordinate (O) at (0,0);
  \coordinate (P) at (-19:\R);
  \coordinate (A) at ($(P)+(71:0.92)$);
  \draw[given] (O) circle (\R);
  \draw[given] (O) -- (P);
  \draw[given] ($(P)+(71:-0.50)$) -- ($(P)+(71:1.84)$);
  \fill (O) circle (1.2pt);
  \fill (A) circle (1.2pt);
  \fill (P) circle (1.2pt);
  \node[vlab, left]        at (O) {$O$};
  \node[vlab, below right] at (P) {$P$};
  \node[vlab, right]       at (A) {$A$};
  \node[vlab, above]       at ($(P)+(71:1.84)$) {$l$};
\end{tikzpicture}
\figcap{11--3}
\end{minipage}\hfill
\begin{minipage}{0.50\linewidth}\centering
\begin{tikzpicture}[scale=1.30]
  \def\R{1.25}
  \coordinate (O) at (0,0);
  \coordinate (A) at (0,-0.39);
  \coordinate (P) at (1.188,-0.39);
  \coordinate (B) at (-1.188,-0.39);
  \draw[given] (O) circle (\R);
  \draw[aux]   (0,1.58) -- (0,-1.58);
  \draw[given] (-1.72,-0.39) -- (1.68,-0.39);
  \draw[given] (O) -- (P);
  \fill (O) circle (1.2pt); \fill (A) circle (1.2pt);
  \fill (B) circle (1.2pt); \fill (P) circle (1.2pt);
  \node[vlab, right]       at (0,0.80) {$m$};
  \node[vlab, above right] at (O) {$O$};
  % A sits in the narrow wedge between OP and l, exactly as the book sets it
  % (scan13 p20).  Anchored placement puts the letter on one line or the
  % other, so the offset is explicit and measured: 0.14 R right, 0.14 R up.
  \node[vlab]              at ($(A)+(0.162,0.191)$) {$A$};
  \node[vlab, above left]  at (B) {$B$};
  \node[vlab, below right] at (P) {$P$};
  \node[vlab, right]       at (1.68,-0.39) {$l$};
\end{tikzpicture}
\figcap{11--4}
\end{minipage}
\end{bkfigure}}

% ------------------------------------------------------- Figs 11-5 and 11-6
% 11-5: regular hexagon (measured, see header) with both circles.  C and E are
% the midpoints of the adjacent sides AB and BD, so OC and OE are apothems and
% the inner circle -- radius R cos(30) -- is tangent at exactly those points.
% 11-6: the four named parts of a regular polygon, with leader lines.
\newcommand{\FIGXIFIVESIX}{\begin{bkfigure}[\linewidth]
\begin{minipage}{0.44\linewidth}\centering
% Both halves are drawn at 1.15 rather than 0.86.  The two figures carry a lot
% of lettering inside the outline (C, E here; four whole words in 11-6), and
% node text does not scale with the picture, so at 0.86 the letters were
% bigger relative to the drawing than the book sets them.
\begin{tikzpicture}[scale=1.15]
  \def\R{1.45}
  \coordinate (O) at (0,0);
  \coordinate (A)  at (84:\R);   \coordinate (B)  at (24:\R);
  \coordinate (D)  at (-36:\R);  \coordinate (V4) at (-96:\R);
  \coordinate (V5) at (-156:\R); \coordinate (V6) at (144:\R);
  \coordinate (C) at ($(A)!0.5!(B)$);
  \coordinate (E) at ($(B)!0.5!(D)$);
  \draw[given] (O) circle (\R);
  \draw[given] (A) -- (B) -- (D) -- (V4) -- (V5) -- (V6) -- cycle;
  \draw[given] (O) circle ({\R*cos(30)});
  \draw[aux]   (O) -- (A);
  \draw[given] (O) -- (C);
  \draw[aux]   (O) -- (B);
  \draw[given] (O) -- (E);
  \draw[aux]   (O) -- (D);
  \fill (O) circle (1.2pt);
  \node[vlab, above]       at (A) {$A$};
  \node[vlab, above right] at (B) {$B$};
  % C and E are tangency points: they sit ON the inscribed circle, so an
  % anchored label grazes that arc.  The book pulls both letters INSIDE the
  % inner circle instead -- C straight down from its side's midpoint, E down
  % and to the left of its own -- measured off scan13 p21 and reproduced here.
  \node[vlab]              at ($(C)+(0.04,-0.34)$) {$C$};
  \node[vlab, left]        at (O) {$O$};
  \node[vlab]              at ($(E)+(-0.29,-0.21)$) {$E$};
  \node[vlab, below right] at (D) {$D$};
\end{tikzpicture}
\figcap{11--5}
\end{minipage}\hfill
\begin{minipage}{0.52\linewidth}\centering
\begin{tikzpicture}[scale=1.15]
  \def\R{1.30}
  \coordinate (O) at (0,0);
  \coordinate (Vr) at (0:\R);
  \coordinate (Vb) at (60:\R);
  \coordinate (Va) at (120:\R);
  \coordinate (Ap) at (-30:{\R*cos(30)});
  \draw[given] (0:\R) -- (60:\R) -- (120:\R) -- (180:\R)
            -- (240:\R) -- (300:\R) -- cycle;
  \draw[aux] (O) -- (Va);
  \draw[aux] (O) -- (Vb);
  \draw[aux] (O) -- (Vr);
  \draw[aux] (O) -- (Ap);
  % The book marks the TOP central angle -- the one between the radii to the
  % 60- and 120-degree vertices -- not the one on the radius side.  Arc radius
  % 0.40 R, measured off scan13 p21.
  \draw[bkarrb] (60:{\R*0.40}) arc (60:120:{\R*0.40});
  \draw (O) circle (0.07);
  \node[slab, align=left, anchor=west] (cang) at (1.30,1.05) {Central\\angle};
  \node[slab, anchor=west] (crad) at (1.46,0.18) {Radius};
  % "Center" sits INSIDE the hexagon at its lower left in the book, with the
  % arrow running up and to the right to the centre mark.
  \node[slab, anchor=east] (ccen) at (0.09,-0.45) {Center};
  \node[slab, anchor=west] (capo) at (1.16,-1.16) {Apothem};
  \draw[lead] (cang.west) -- (66:{\R*0.44});
  \draw[lead] (crad.west) -- ($(O)!0.62!(Vr)$);
  \draw[lead] (-0.62,-0.26) -- (-0.10,-0.115);
  \draw[lead] (capo.west) -- ($(O)!0.72!(Ap)$);
\end{tikzpicture}
\figcap{11--6}
\end{minipage}
\end{bkfigure}}

% ---------------------------------------------------------------- Fig 11-7
% Exterior angle of a regular polygon.  The side DC is produced beyond C; the
% marked angle is between that extension and the next side CB.
\newcommand{\FIGXISEVEN}{\begin{bkfigure}[0.92\linewidth]
\begin{tikzpicture}[scale=0.80]
  \def\R{1.35}
  \coordinate (Cv) at (0:\R);
  \coordinate (Bv) at (60:\R);
  \coordinate (Dv) at (-60:\R);
  \draw[given] (0:\R) -- (60:\R) -- (120:\R) -- (180:\R)
            -- (240:\R) -- (300:\R) -- cycle;
  % Measured off scan13 p22: the extension runs a full side length beyond C,
  % and the marked angle's arc has radius ~0.53 R, not the 0.31 R first drawn.
  \coordinate (Ext) at ($(Cv)!-1.00!(Dv)$);
  \draw[aux] (Cv) -- (Ext);
  \draw[bkarrb] ($(Cv)+(120:0.72)$) arc (120:60:0.72);
\end{tikzpicture}
\figcap{11--7}
\end{bkfigure}}

% ------------------------------------------------------- Figs 11-8 and 11-9
% Two regular hexagons of different size, flat-topped (vertices at k*60 deg
% with C at 0 deg).  11-9 adds Q, the midpoint of the top side AB, so that
% OQ is the apothem a and OA the radius r.
\newcommand{\FIGXIEIGHTNINE}{\begin{bkfigure}[\linewidth]
\begin{minipage}{0.44\linewidth}\centering
% 0.70 here against 0.90 in 11-9: the two blocks share one text line, and 11-9
% needs the larger scale to keep r', a' off the radii it sets them between.
\begin{tikzpicture}[scale=0.70]
  \def\Ra{1.55}\def\Rb{1.10}
  \begin{scope}
    \coordinate (O) at (0,0);
    \coordinate (A) at (120:\Ra); \coordinate (B) at (60:\Ra);
    \coordinate (C) at (0:\Ra);   \coordinate (D) at (-60:\Ra);
    \coordinate (E) at (-120:\Ra);
    \draw[given] (0:\Ra) -- (60:\Ra) -- (120:\Ra) -- (180:\Ra)
              -- (240:\Ra) -- (300:\Ra) -- cycle;
    \draw[aux] (O) -- (A);
    \draw[aux] (O) -- (B);
    \fill (O) circle (1.2pt);
    \node[vlab, above left]  at (A) {$A$};
    \node[vlab, above right] at (B) {$B$};
    \node[vlab, right]       at (C) {$C$};
    \node[vlab, below right] at (D) {$D$};
    \node[vlab, below left]  at (E) {$E$};
    \node[vlab, below]       at (O) {$O$};
  \end{scope}
  \begin{scope}[xshift=3.35cm]
    \coordinate (Op) at (0,0);
    \coordinate (Ap) at (120:\Rb); \coordinate (Bp) at (60:\Rb);
    \coordinate (Cp) at (0:\Rb);   \coordinate (Dp) at (-60:\Rb);
    \coordinate (Ep) at (-120:\Rb);
    \draw[given] (0:\Rb) -- (60:\Rb) -- (120:\Rb) -- (180:\Rb)
              -- (240:\Rb) -- (300:\Rb) -- cycle;
    \draw[aux] (Op) -- (Ap);
    \draw[aux] (Op) -- (Bp);
    \fill (Op) circle (1.2pt);
    \node[vlab, above left]  at (Ap) {$A\XIpr$};
    \node[vlab, above right] at (Bp) {$B\XIpr$};
    \node[vlab, right]       at (Cp) {$C\XIpr$};
    \node[vlab, below right] at (Dp) {$D\XIpr$};
    \node[vlab, below left]  at (Ep) {$E\XIpr$};
    \node[vlab, below]       at (Op) {$O\XIpr$};
  \end{scope}
\end{tikzpicture}
\figcap{11--8}
\end{minipage}\hfill
\begin{minipage}{0.52\linewidth}\centering
% 0.90, not 0.80.  r, a, r' and a' are set inside the wedge between two
% radii, and on the smaller polygon that wedge is only a few points wider
% than the letters themselves; the extra 12% is what keeps them off the lines.
\begin{tikzpicture}[scale=0.90]
  \def\Ra{1.45}\def\Rb{1.02}
  \begin{scope}
    \coordinate (O) at (0,0);
    \coordinate (A) at (120:\Ra); \coordinate (B) at (60:\Ra);
    \coordinate (C) at (0:\Ra);
    \coordinate (Q) at ($(A)!0.5!(B)$);
    \draw[given] (0:\Ra) -- (60:\Ra) -- (120:\Ra) -- (180:\Ra)
              -- (240:\Ra) -- (300:\Ra) -- cycle;
    \draw[aux, bkarr] (A) -- (O);
    \draw[aux, bkarr] (Q) -- (O);
    \draw[aux, bkarr] (B) -- (O);
    \fill (O) circle (1.2pt);
    \node[vlab, above left]  at (A) {$A$};
    \node[vlab, above]       at (Q) {$Q$};
    \node[vlab, above right] at (B) {$B$};
    \node[vlab, right]       at (C) {$C$};
    \node[vlab, below]       at (O) {$O$};
    \node[slab, left]  at ($(A)!0.45!(O)$) {$r$};
    % The wedge between OQ and OB tapers toward O, so a sits high in it, as
    % the book has it (scan13 p23); coordinates solved for equal clearance.
    \node[slab] at (0.186,0.944) {$a$};
  \end{scope}
  \begin{scope}[xshift=3.05cm]
    \coordinate (Op) at (0,0);
    \coordinate (Ap) at (120:\Rb); \coordinate (Bp) at (60:\Rb);
    \coordinate (Cp) at (0:\Rb);
    \coordinate (Qp) at ($(Ap)!0.5!(Bp)$);
    \draw[given] (0:\Rb) -- (60:\Rb) -- (120:\Rb) -- (180:\Rb)
              -- (240:\Rb) -- (300:\Rb) -- cycle;
    \draw[aux, bkarr] (Ap) -- (Op);
    \draw[aux, bkarr] (Qp) -- (Op);
    \draw[aux, bkarr] (Bp) -- (Op);
    \fill (Op) circle (1.2pt);
    \node[vlab, above left]  at (Ap) {$A\XIpr$};
    \node[vlab, above]       at (Qp) {$Q\XIpr$};
    \node[vlab, above right] at (Bp) {$B\XIpr$};
    \node[vlab, right]       at (Cp) {$C\XIpr$};
    \node[vlab, below]       at (Op) {$O\XIpr$};
    % On the smaller hexagon the corridors these two letters live in -- r'
    % between side A'(180-deg vertex) and radius O'A', a' between O'Q' and
    % O'B' -- are only a couple of points wider than the primed letter boxes
    % themselves (5.4 x 10.0 pt, measured off the compiled page).  Anchored
    % placement cannot fit them, so both are set at coordinates solved for
    % equal clearance on every side; r' still hugs the radius it names, as
    % the book has it.
    \node[slab] at (-0.508,0.333) {$r\XIpr$};
    \node[slab] at ( 0.141,0.623) {$a\XIpr$};
  \end{scope}
\end{tikzpicture}
\figcap{11--9}
\end{minipage}
\end{bkfigure}}

% --------------------------------------------------------------- Fig 11-10
% Inscribed square; bisecting the central angle in the first quadrant gives
% the octagon vertex at 45 deg, and the two dashed chords are the octagon
% sides there.  The bisection is exact (45 = 90/2), not drawn by eye.
\newcommand{\FIGXITEN}{\begin{bkfigure}[0.56\linewidth]
\begin{tikzpicture}[scale=0.92]
  \def\R{1.50}
  \coordinate (O) at (0,0);
  \draw[given] (O) circle (\R);
  \draw[given] (90:\R) -- (0:\R) -- (-90:\R) -- (180:\R) -- cycle;
  % Checked against scan14 p3 at 3x: the radii to the square's own vertices
  % carry NO arrowheads; only the bisecting ray does, and that one is
  % double-headed (a head at the centre and a head at the circle).
  \draw[aux] (O) -- (90:\R);
  \draw[aux] (O) -- (0:\R);
  \draw[aux, bkarrb] (O) -- (45:\R);
  \draw[aux] (90:\R) -- (45:\R) -- (0:\R);
\end{tikzpicture}
\figcap{11--10}
\end{bkfigure}}

% --------------------------------------------------------------- Fig 11-11
% Circle of radius 1 with the inscribed hexagon (solid, vertices at 30+k*60)
% and the circumscribed hexagon (dashed, vertices at k*60, circumradius
% R/cos 30) -- so the dashed hexagon's sides are tangent to the circle at the
% solid hexagon's vertices, exactly as the book draws it.
\newcommand{\FIGXIELEVEN}{\begin{bkfigure}[0.52\linewidth]
\begin{tikzpicture}[scale=1.05]
  \def\R{1.30}
  \coordinate (O) at (0,0);
  \draw[given] (O) circle (\R);
  \draw[given] (30:\R) -- (90:\R) -- (150:\R) -- (210:\R)
            -- (270:\R) -- (330:\R) -- cycle;
  \draw[aux] (0:{\R/cos(30)})   -- (60:{\R/cos(30)})  -- (120:{\R/cos(30)})
          -- (180:{\R/cos(30)}) -- (240:{\R/cos(30)}) -- (300:{\R/cos(30)})
          -- cycle;
  \draw[bkarr] (O) -- (0:\R);
  \fill (O) circle (1.2pt);
  \node[slab, below] at (0:{\R*0.55}) {1};
\end{tikzpicture}
\figcap{11--11}
\end{bkfigure}}

% --------------------------------------------------------------- Fig 11-12
% A and C are ends of a side of the inscribed polygon; B is its midpoint, so
% OB is the apothem a_n and OB is perpendicular to the chord AC.
\newcommand{\FIGXITWELVE}{\begin{bkfigure}[0.50\linewidth]
\begin{tikzpicture}[scale=1.0]
  \def\R{1.40}
  % A and C at +-45 deg (was +-40): measured off scan14 p5, where the chord
  % subtends ~90 deg and the drawn apothem is 0.69 r, not 0.77 r.
  \coordinate (O) at (0,0);
  \coordinate (A) at (45:\R);
  \coordinate (C) at (-45:\R);
  \coordinate (B) at ($(A)!0.5!(C)$);
  \draw[given] (O) circle (\R);
  \draw[given] (O) -- (A);
  \draw[given] (O) -- (B);
  \draw[given] (O) -- (C);
  \draw[given] (A) -- (C);
  \fill (O) circle (1.2pt);
  \node[vlab, above right] at (A) {$A$};
  % B lies on the chord AC with the circle close behind it: the chapter-wide
  % 4.8pt outer sep pushed the letter onto the arc, so this one label uses the
  % style file's own 3.4pt and sits where the book sets it, between the two.
  \node[vlab, right, outer sep=2.6pt] at (B) {$B$};
  \node[vlab, below right] at (C) {$C$};
  \node[vlab, left]        at (O) {$O$};
  \node[slab, above left]  at ($(O)!0.55!(A)$) {$r$};
  \node[slab, above]       at ($(O)!0.66!(B)$) {$a_n$};
\end{tikzpicture}
\figcap{11--12}
\end{bkfigure}}

% --------------------------------------------------------------- Fig 11-13
% Annular ring.  Radii 2 and 4 as in Exercise 11-13.2, i.e. inner = outer/2.
\newcommand{\FIGXITHIRTEEN}{\begin{bkfigure}[\linewidth]
\begin{tikzpicture}[scale=1.0]
  \def\Ro{1.30}
  \path[pattern=vertical lines, even odd rule]
        (0,0) circle (\Ro) (0,0) circle ({\Ro/2});
  \draw[given] (0,0) circle (\Ro);
  \draw[given] (0,0) circle ({\Ro/2});
\end{tikzpicture}
\figcap{11--13}
\end{bkfigure}}

% --------------------------------------------------------------- Fig 11-14
% ORGANIC SHAPE -- built against scan14 p7 (the straight duplicate), per the
% source policy.  AC is a diameter of the big circle and B its centre, with
% AB = BC.  The S-curve is the semicircle on AB drawn ABOVE the diameter
% followed by the semicircle on BC drawn BELOW it; the shaded region is
% everything inside the big circle below that curve, which is why its area is
% exactly half the disc (Exercise 11-13.3).
\newcommand{\FIGXIFOURTEEN}{\begin{bkfigure}[\linewidth]
\begin{tikzpicture}[scale=0.95]
  \def\R{1.40}
  \coordinate (A) at (-\R,0);
  \coordinate (B) at (0,0);
  \coordinate (C) at (\R,0);
  \path[pattern=north east lines]
        (A) arc[start angle=180, end angle=0,   radius={\R/2}]
            arc[start angle=180, end angle=360, radius={\R/2}]
            arc[start angle=0,   end angle=-180, radius=\R] -- cycle;
  \draw[given] (B) circle (\R);
  \draw[given] (A) arc[start angle=180, end angle=0,   radius={\R/2}];
  \draw[given] (B) arc[start angle=180, end angle=360, radius={\R/2}];
  \fill (A) circle (1.3pt); \fill (B) circle (1.3pt); \fill (C) circle (1.3pt);
  % A stands inside the hatched region; the book clears a small white patch of
  % hatching around the letter rather than printing it over the ruling.
  \node[vlab, right, fill=white, inner sep=1.6pt] at (A) {$A$};
  \node[vlab, right] at (B) {$B$};
  \node[vlab, right] at (C) {$C$};
\end{tikzpicture}
\figcap{11--14}
\end{bkfigure}}

% --------------------------------------------------------------- Fig 11-15
% Square with both its circles: the circumscribed circle through the vertices
% (radius R) and the inscribed circle touching the sides (radius R/sqrt 2).
\newcommand{\FIGXIFIFTEEN}{\begin{bkfigure}[\linewidth]
\begin{tikzpicture}[scale=1.0]
  \def\R{1.30}
  \coordinate (O) at (0,0);
  \draw[given] (O) circle (\R);
  \draw[given] (45:\R) -- (135:\R) -- (225:\R) -- (315:\R) -- cycle;
  \draw[given] (O) circle ({\R/sqrt(2)});
\end{tikzpicture}
\figcap{11--15}
\end{bkfigure}}

% --------------------------------------------------------------- Fig 11-16
% Semicircles erected outward on the three sides of a right triangle; the
% right angle is at the lower-left vertex, so the two legs are the vertical
% and horizontal sides and the hypotenuse carries the largest semicircle.
\newcommand{\FIGXISIXTEEN}{\begin{bkfigure}[\linewidth]
\begin{tikzpicture}[scale=0.86]
  \coordinate (L) at (0,0);
  \coordinate (T) at (0,1.60);
  \coordinate (R) at (2.52,0);
  \coordinate (Mh) at ($(T)!0.5!(R)$);
  \draw[given] (L) -- (T) -- (R) -- cycle;
  \draw[aux] (T) arc[start angle=90,  end angle=270, radius=0.80];
  \draw[aux] (L) arc[start angle=180, end angle=360, radius=1.26];
  \draw[aux] (T) arc[start angle={atan2(0.80,-1.26)},
                     end angle={atan2(0.80,-1.26)-180}, radius=1.4925];
\end{tikzpicture}
\figcap{11--16}
\end{bkfigure}}

% figures12.tex -- Chapter 12 figures, built 2026-08-17 from the iPad scans
% (scan14 pp. 8-21 = book pp. 210-226; every figure crop-checked at 150-200 dpi).
%
% Constrained points are DERIVED, never eyeballed:
%   arc midpoints    -> half-angle arithmetic on the central angle
%   angle bisectors  -> half-angle arithmetic
%   crossings        -> (intersection of A--B and C--D)
%   tangency         -> radius-perpendicular construction
%   external points  -> derived as the crossing of the two secants
% Hypotheses are re-asserted numerically in tools/constraints/ch12.py.
%
% 2026-08-17 figure-repair pass (feedback/FIX-SPEC.md).  Figs 12-28, 12-29 and
% 12-30 now carry the BOOK's schematic proportions, NOT the lengths quoted in
% Exercises 15-18: drawing those lengths to scale handed the student the answer.
% Only the structural hypotheses (tangency, collinearity, betweenness, which
% points lie on the circle) are enforced there now.
%
% Macro names are chapter-prefixed: \FIGXII<NUMBERWORD>.
%
% Line treatment follows the book's own convention: solid where the book
% prints solid (given), dashed where the book prints dashed (aux).
%
% Labelling rule used throughout: a point ON the circle is labelled RADIALLY
% OUTWARD, which is always clear of the chords and radii that meet it; interior
% points are labelled along the bisector of the widest free wedge at that point.

\usetikzlibrary{decorations.pathreplacing,arrows.meta}

% Arc-measure marker used in Figs 12-14 to 12-17, and the angle arcs of
% 12-21/12-22.  The book strikes all of these MUCH thinner than the figure
% proper -- the sector's radii and arc read bold against them.
\tikzset{
  XIIhook/.style={fig, line width=0.45pt},
}

% A curly BRACE struck along a circular arc, which is how the book marks an
% arc's measure (Figs 12-14, 12-15, 12-16): a thin curve concentric with the
% arc, an inward hook at each end, and a nib at the middle pointing outward at
% the label.  Before this pass these were drawn as plain arcs with straight
% end ticks (or with nothing at all), which is what Caleb flagged.
%   #1 centre   #2 radius   #3 start angle   #4 end angle   #5 hook/nib depth
% The nib offset is signed off (#4-#3), so the argument order may run either way.
\newcommand{\XIIbrace}[5]{%
  \pgfmathsetmacro{\XIIbm}{0.5*((#3)+(#4))}%
  \pgfmathsetmacro{\XIIbd}{0.05*((#4)-(#3))}%
  \draw[XIIhook, line join=round]
      ($(#1)+({#3}:{(#2)-(#5)})$)
   -- ($(#1)+({#3}:{#2})$)
      arc[start angle={#3}, end angle={\XIIbm-\XIIbd}, radius={#2}]
   -- ($(#1)+({\XIIbm}:{(#2)+(#5)})$)
   -- ($(#1)+({\XIIbm+\XIIbd}:{#2})$)
      arc[start angle={\XIIbm+\XIIbd}, end angle={#4}, radius={#2}]
   -- ($(#1)+({#4}:{(#2)-(#5)})$);%
}

% Primes and degree marks carry a hair kern (1.3mu, about 0.65pt) so that the
% label audit tokenises "A'" and "n^\circ" as single symbols instead of two
% touching glyphs. The shift is well under a point and does not read on the page.
\newcommand{\XIIp}{\mkern1.3mu'}
\newcommand{\XIId}{\mkern1.3mu^\circ}

%----------------------------------------------------------------- Fig 12-1
% Diameter CB through O; A on the lower-left arc; radius OA.  The arc ABC is
% the arc exterior to the central angle AOC.
\newcommand{\FIGXIIONE}{\begin{bkfigure}[\linewidth]
\begin{tikzpicture}[scale=1.300]
  \coordinate (O) at (0,0);
  \def\R{1.25}
  \coordinate (C) at ($(O)+(180:\R)$);
  \coordinate (B) at ($(O)+(0:\R)$);
  \coordinate (A) at ($(O)+(227:\R)$);
  \draw[given] (O) circle[radius=\R];
  \draw[given] (C) -- (B);
  \draw[given] (O) -- (A);
  \node[vlab] at ($(O)+(180:\R+0.28)$) {$C$};
  \node[vlab] at ($(O)+(0:\R+0.28)$) {$B$};
  \node[vlab] at ($(O)+(227:\R+0.28)$) {$A$};
  \node[vlab] at ($(O)+(80:0.32)$) {$O$};
\end{tikzpicture}
\figcap{12--1}
\end{bkfigure}}

%---------------------------------------------------------------- Fig 12-2
% Rays OA, OC, OB with C interior to angle AOB, so angle AOB > angle COB.
% Split out of the old \FIGXIITWOTHREE for the 2026-08-17 pass: 12-2 belongs to
% Exercise 1 of Group 12-1 and 12-3 to Exercise 4, so they sit inline after
% different exercises and can no longer share one captioned block.
\newcommand{\FIGXIITWO}{\begin{bkfigure}[\linewidth]
\begin{tikzpicture}[scale=1.117]
  \coordinate (O) at (0,0);
  \def\R{1.45}
  \coordinate (A) at ($(O)+(78:\R)$);
  \coordinate (C) at ($(O)+(47:\R)$);
  \coordinate (B) at ($(O)+(0:\R)$);
  \draw[given] (O) circle[radius=\R];
  \draw[given] (O) -- (A);
  \draw[given] (O) -- (C);
  \draw[given] (O) -- (B);
  \node[vlab] at ($(O)+(78:\R+0.28)$) {$A$};
  \node[vlab] at ($(O)+(47:\R+0.28)$) {$C$};
  \node[vlab] at ($(O)+(0:\R+0.28)$) {$B$};
  \node[vlab] at ($(O)+(225:0.34)$) {$O$};
\end{tikzpicture}
\figcap{12--2}
\end{bkfigure}}

%---------------------------------------------------------------- Fig 12-3
% Chord EF; P and P' are where rays OQ, OQ' cut EF -- both DERIVED as
% intersections, never placed by hand.
\newcommand{\FIGXIITHREE}{\begin{bkfigure}[\linewidth]
\begin{tikzpicture}[scale=1.117]
  \coordinate (O) at (0,0);
  \def\R{1.90}
  \coordinate (E) at ($(O)+(75:\R)$);
  \coordinate (F) at ($(O)+(0:\R)$);
  \coordinate (Q) at ($(O)+(56:\R)$);
  % Q' at 27 rather than the plate's 22: the wedge Q'OF has to hold the label
  % P', and 22 degrees on this radius leaves it touching the ray.
  \coordinate (Qp) at ($(O)+(27:\R)$);
  \coordinate (P)  at (intersection of O--Q and E--F);
  \coordinate (Pp) at (intersection of O--Qp and E--F);
  \draw[given] (O) circle[radius=\R];
  \draw[given] (O) -- (E);
  \draw[given] (O) -- (F);
  \draw[given] (E) -- (F);
  \draw[given] (O) -- (Q);
  \draw[given] (O) -- (Qp);
  \node[vlab] at ($(O)+(75:\R+0.28)$) {$E$};
  \node[vlab] at ($(O)+(56:\R+0.28)$) {$Q$};
  \node[vlab] at ($(O)+(27:\R+0.30)$) {$Q\XIIp$};
  \node[vlab] at ($(O)+(0:\R+0.28)$) {$F$};
  \node[vlab] at ($(O)+(238:0.34)$) {$O$};
  \node[vlab] at ($(P)+(271.75:0.36)$) {$P$};
  % P' sits on the bisector of the wedge Q'OF, the one direction that clears
  % both the ray OQ' and the chord EF, which cross at P' itself.
  \node[vlab] at ($(Pp)+(250:0.46)$) {$P\XIIp$};
\end{tikzpicture}
\figcap{12--3}
\end{bkfigure}}

%----------------------------------------------------------------- Fig 12-4
% C is the midpoint of the arc EXTERIOR to central angle AOB: the exterior arc
% runs A(118) counterclockwise to B(363), so C sits at (118+363)/2 = 240.5.
\newcommand{\FIGXIIFOUR}{\begin{bkfigure}[\linewidth]
\begin{tikzpicture}[scale=1.300]
  \coordinate (O) at (0,0);
  \def\R{1.25}
  \coordinate (A) at ($(O)+(118:\R)$);
  \coordinate (B) at ($(O)+(3:\R)$);
  \coordinate (C) at ($(O)+(240.5:\R)$);
  \draw[given] (O) circle[radius=\R];
  \draw[given] (O) -- (A);
  \draw[given] (O) -- (B);
  \draw[aux]   (O) -- (C);
  \node[vlab] at ($(O)+(118:\R+0.28)$) {$A$};
  \node[vlab] at ($(O)+(3:\R+0.28)$) {$B$};
  \node[vlab] at ($(O)+(240.5:\R+0.28)$) {$C$};
  \node[vlab] at ($(O)+(60:0.34)$) {$O$};
  \node[slab] at ($(O)+(240.5:\R+0.78)$) {Midpoint};
\end{tikzpicture}
\figcap{12--4}
\end{bkfigure}}

%------------------------------------------------------- Figs 12-5 and 12-6
% The book prints these two side by side on p.211, so they stay one block.
%
% 12-5, rebuilt 2026-08-17 from the book plate (and Caleb's photo of it):
%   * the book draws the two RADII OA and OB, meeting at a shallow 170-degree
%     bend at O -- not the single chord AB the old version drew, which missed
%     the centre by 0.44 r and left O stranded;
%   * the arc therefore spans A(145) to B(-25), 170 degrees, which is what
%     gives the nine labels room.  They were crowded before because the arc
%     was only 146 degrees wide.
%   * C = midpoint of arc AB = 60.  Every later point is the midpoint of the
%     arc it bisects, so all nine positions are arithmetic:
%       D = mid(A,C) = 102.5    E = mid(C,B) = 17.5
%       F = mid(A,D) = 123.75   G = mid(D,C) = 81.25
%       H = mid(C,E) = 38.75    I = mid(E,B) = -3.75
%   * the eight-side path is struck heavier, as the book strikes it, and each
%     of its interior vertices carries a short radial nick.  Caleb's complaint
%     was that the letters read as a caption strip because nothing on the arc
%     marked where the points actually were; the nicks fix that.
% 12-6: same arc, showing an inscribed path against a circumscribed path whose
% corners are the intersections of the tangents at A, C and at C, B --
% distance R/cos(42.5) along the bisecting directions 102.5 and 17.5.
\newcommand{\FIGXIIFIVESIX}{\begin{bkfigure}[\linewidth]
\begin{minipage}[b]{0.54\linewidth}\centering
\begin{tikzpicture}[scale=1.130]
  \coordinate (O) at (0,0);
  \def\R{1.5}
  \coordinate (A) at ($(O)+(145:\R)$);
  \coordinate (B) at ($(O)+(-25:\R)$);
  \coordinate (C) at ($(O)+(60:\R)$);
  \coordinate (D) at ($(O)+(102.5:\R)$);
  \coordinate (E) at ($(O)+(17.5:\R)$);
  \coordinate (F) at ($(O)+(123.75:\R)$);
  \coordinate (G) at ($(O)+(81.25:\R)$);
  \coordinate (H) at ($(O)+(38.75:\R)$);
  \coordinate (I) at ($(O)+(-3.75:\R)$);
  \draw[given] (O) circle[radius=\R];
  \draw[given] (O) -- (A);
  \draw[given] (O) -- (B);
  % the eight-side path is struck HEAVIER in the book, not dashed, so that it
  % reads apart from the dashed two-side path it refines
  \draw[fig, line width=1.5pt]
        (A) -- (F) -- (D) -- (G) -- (C) -- (H) -- (E) -- (I) -- (B);
  % short radial nicks at the interior vertices, so each letter labels a place
  % on the arc and not a slot in a row
  \foreach \a in {123.75, 102.5, 81.25, 38.75, 17.5, -3.75}
    {\draw[fig, line width=0.5pt] ($(O)+(\a:\R)$) -- ($(O)+(\a:\R+0.09)$);}
  \draw[aux] (A) -- (C) -- (B);
  \foreach \p/\a in {A/145, B/-25, C/60, D/102.5, E/17.5,
                     F/123.75, G/81.25, H/38.75, I/-3.75}
    {\node[vlab] at ($(O)+(\a:\R+0.38)$) {$\p$};}
  \node[vlab] at ($(O)+(238:0.34)$) {$O$};
\end{tikzpicture}
\figcap{12--5}
\end{minipage}\hfill
\begin{minipage}[b]{0.40\linewidth}\centering
\begin{tikzpicture}[scale=1.060]
  \coordinate (O) at (0,0);
  \def\R{1.3}
  \coordinate (A) at ($(O)+(145:\R)$);
  \coordinate (B) at ($(O)+(-25:\R)$);
  \coordinate (C) at ($(O)+(60:\R)$);
  % tangent-intersection corners of the circumscribed path
  \pgfmathsetmacro{\Rt}{\R/cos(42.5)}
  \coordinate (T) at ($(O)+(102.5:\Rt)$);
  \coordinate (U) at ($(O)+(17.5:\Rt)$);
  \draw[given] (O) circle[radius=\R];
  \draw[given] (O) -- (A);
  \draw[given] (O) -- (B);
  \draw[aux] (A) -- (C) -- (B);
  \draw[aux] (A) -- (T) -- (C) -- (U) -- (B);
\end{tikzpicture}
\figcap{12--6}
\end{minipage}
\end{bkfigure}}

%---------------------------------------------------------------- Fig 12-7
% Regular hexagon, A(120), C(60), B(0); radii OA, OC, OB drawn.
% Split out of the old \FIGXIISEVENEIGHT: 12-7 belongs to Exercise 1 of Group
% 12-2 and 12-8 to Exercise 3, so each now stands inline after its own item.
\newcommand{\FIGXIISEVEN}{\begin{bkfigure}[\linewidth]
\begin{tikzpicture}[scale=1.196]
  \coordinate (O) at (0,0);
  \def\R{1.25}
  \coordinate (A) at ($(O)+(120:\R)$);
  \coordinate (C) at ($(O)+(60:\R)$);
  \coordinate (B) at ($(O)+(0:\R)$);
  \draw[given] (O) circle[radius=\R];
  \draw[given] ($(O)+(0:\R)$) -- ($(O)+(60:\R)$) -- ($(O)+(120:\R)$)
            -- ($(O)+(180:\R)$) -- ($(O)+(240:\R)$) -- ($(O)+(300:\R)$) -- cycle;
  \draw[given] (O) -- (A);
  \draw[given] (O) -- (C);
  \draw[given] (O) -- (B);
  \node[vlab] at ($(O)+(120:\R+0.30)$) {$A$};
  \node[vlab] at ($(O)+(60:\R+0.30)$) {$C$};
  \node[vlab] at ($(O)+(0:\R+0.30)$) {$B$};
  \node[vlab] at ($(O)+(255:0.34)$) {$O$};
\end{tikzpicture}
\figcap{12--7}
\end{bkfigure}}

%---------------------------------------------------------------- Fig 12-8
% Angle AOC is a right angle and angle AOB = angle BOC, so B bisects it.
% Rebuilt 2026-08-17.  Two things were wrong and Caleb caught both:
%   * the OUTER arc ran from OA all the way to OC, CROSSING ray OB -- it marked
%     angle AOC.  The book strikes TWO arcs that abut at OB, one on angle AOB
%     and one on angle BOC, at clearly different radii (the book's differ by
%     about 60%); nothing crosses a ray.  That adjacency is the whole point of
%     the exercise.
%   * the fan was axis-aligned (OA vertical, OC horizontal), so the render read
%     as though angle AOC were the axis frame rather than the right angle the
%     exercise names.  The book tilts the fan; A(78), B(33), C(-12) keeps
%     angle AOC = 90 and angle AOB = angle BOC = 45 exactly, with no ray level.
\newcommand{\FIGXIIEIGHT}{\begin{bkfigure}[\linewidth]
\begin{tikzpicture}[scale=1.196]
  \coordinate (O) at (0,0);
  \def\R{1.25}
  \coordinate (A) at ($(O)+(78:\R)$);
  \coordinate (B) at ($(O)+(33:\R)$);
  \coordinate (C) at ($(O)+(-12:\R)$);
  \draw[given] (O) circle[radius=\R];
  \draw[given] (O) -- (A);
  \draw[given] (O) -- (B);
  \draw[given] (O) -- (C);
  % angle AOB, inner radius; angle BOC, outer radius.  They meet on ray OB.
  \draw[XIIhook] ($(O)+(33:0.36)$) arc[start angle=33, end angle=78, radius=0.36];
  \draw[XIIhook] ($(O)+(-12:0.60)$) arc[start angle=-12, end angle=33, radius=0.60];
  \node[vlab] at ($(O)+(78:\R+0.28)$) {$A$};
  \node[vlab] at ($(O)+(33:\R+0.28)$) {$B$};
  \node[vlab] at ($(O)+(-12:\R+0.28)$) {$C$};
  \node[vlab] at ($(O)+(228:0.28)$) {$O$};
\end{tikzpicture}
\figcap{12--8}
\end{bkfigure}}

%----------------------------------------------------- Figs 12-9 and 12-10
% 12-9: an IRREGULAR inscribed polygonal path in arc AB (Theorem 12-2 drops the
%   regularity requirement), so the vertex spacing is deliberately uneven.
%   Rebuilt 2026-08-17: the path used to have SEVEN short sides, whose sag off
%   the circle was 0.017 r -- under a point on the page.  The dashes therefore
%   printed on top of the circle line and the whole thing read as a ragged
%   doubled edge, which is what Caleb saw.  The book draws four long sides,
%   whose sag is about 0.07 r, so the dashed path is plainly INSIDE the circle
%   with a visible gap.  Spacing 40/34/48/38 degrees keeps it irregular.
% 12-10: adjacent arcs AB and BC, sharing only B.  The emphasis is now a
%   heavier DASH-BROKEN arc struck just outside the circle, which is how the
%   book marks it; simply thickening the circle's own stroke (the old \key arc)
%   was invisible at 1.0pt against a 0.9pt circle.  The fan is tilted too --
%   the book sets no radius on an axis.
\newcommand{\FIGXIININETEN}{\begin{bkfigure}[\linewidth]
\begin{minipage}[b]{0.46\linewidth}\centering
\begin{tikzpicture}[scale=1.198]
  \coordinate (O) at (0,0);
  \def\R{1.25}
  \coordinate (A) at ($(O)+(152:\R)$);
  \coordinate (B) at ($(O)+(-8:\R)$);
  \draw[given] (O) circle[radius=\R];
  \draw[given] (O) -- (A);
  \draw[given] (O) -- (B);
  \draw[aux] (A)
    -- ($(O)+(112:\R)$) -- ($(O)+(78:\R)$) -- ($(O)+(30:\R)$) -- (B);
  \node[vlab] at ($(O)+(152:\R+0.28)$) {$A$};
  \node[vlab] at ($(O)+(-8:\R+0.28)$) {$B$};
  \node[vlab] at ($(O)+(248:0.30)$) {$O$};
\end{tikzpicture}
\figcap{12--9}
\end{minipage}\hfill
\begin{minipage}[b]{0.46\linewidth}\centering
\begin{tikzpicture}[scale=1.198]
  \coordinate (O) at (0,0);
  \def\R{1.25}
  \coordinate (A) at ($(O)+(82:\R)$);
  \coordinate (B) at ($(O)+(37:\R)$);
  \coordinate (C) at ($(O)+(-8:\R)$);
  \draw[given] (O) circle[radius=\R];
  \draw[given] (O) -- (A);
  \draw[given] (O) -- (B);
  \draw[given] (O) -- (C);
  \draw[fig, line width=1.5pt, dash pattern=on 3.6pt off 1.5pt]
        ($(O)+(-8:\R+0.05)$) arc[start angle=-8, end angle=82, radius=\R+0.05];
  \node[vlab] at ($(O)+(82:\R+0.30)$) {$A$};
  \node[vlab] at ($(O)+(37:\R+0.30)$) {$B$};
  \node[vlab] at ($(O)+(-8:\R+0.30)$) {$C$};
  \node[vlab] at ($(O)+(222:0.30)$) {$O$};
\end{tikzpicture}
\figcap{12--10}
\end{minipage}
\end{bkfigure}}

%----------------------------------------------------- Figs 12-11 and 12-12
% 12-11: C interior to angle AOB, so angle AOB (93) > angle COB (80).
%   Directions measured off the plate: A 104, C 90, B 10 (the book sets A and
%   C only about 14 degrees apart, which is why their labels sit side by side).
% 12-12: corresponding central angles congruent on both circles -- the same
% four directions 148, 90, 48, 0 are used at each centre, which is exactly the
% hypothesis of Theorem 12-7.
\newcommand{\FIGXIIELEVENTWELVE}{\begin{bkfigure}[\linewidth]
\begin{minipage}[b]{0.32\linewidth}\centering
\begin{tikzpicture}[scale=0.995]
  \coordinate (O) at (0,0);
  \def\R{1.25}
  \coordinate (A) at ($(O)+(104:\R)$);
  \coordinate (C) at ($(O)+(90:\R)$);
  \coordinate (B) at ($(O)+(10:\R)$);
  \draw[given] (O) circle[radius=\R];
  \draw[given] (O) -- (A);
  \draw[given] (O) -- (C);
  \draw[given] (O) -- (B);
  \node[vlab] at ($(O)+(112:\R+0.32)$) {$A$};
  \node[vlab] at ($(O)+(86:\R+0.32)$) {$C$};
  \node[vlab] at ($(O)+(10:\R+0.30)$) {$B$};
  \node[vlab] at ($(O)+(250:0.34)$) {$O$};
\end{tikzpicture}
\figcap{12--11}
\end{minipage}\hfill
\begin{minipage}[b]{0.64\linewidth}\centering
\begin{tikzpicture}[scale=0.922]
  \coordinate (Op) at (0,0);
  \def\Rp{1.35}
  \coordinate (Ap) at ($(Op)+(148:\Rp)$);
  \coordinate (Pp) at ($(Op)+(90:\Rp)$);
  \coordinate (Qp) at ($(Op)+(48:\Rp)$);
  \coordinate (Bp) at ($(Op)+(0:\Rp)$);
  \draw[given] (Op) circle[radius=\Rp];
  \draw[given] (Op) -- (Ap);
  \draw[given] (Op) -- (Bp);
  \draw[aux] (Op) -- (Pp);
  \draw[aux] (Op) -- (Qp);
  \draw[aux] (Ap) -- (Pp) -- (Qp) -- (Bp);
  \node[vlab] at ($(Op)+(152:\Rp+0.40)$) {$A\XIIp$};
  \node[vlab] at ($(Op)+(90:\Rp+0.34)$) {$P\XIIp$};
  \node[vlab] at ($(Op)+(44:\Rp+0.40)$) {$Q\XIIp$};
  \node[vlab] at ($(Op)+(0:\Rp+0.32)$) {$B\XIIp$};
  \node[vlab] at ($(Op)+(250:0.48)$) {$O\XIIp$};
  \node[slab] at ($(Op)+(0:0.68)+(0,-0.40)$) {$r\XIIp$};
  %
  \coordinate (Oo) at (3.75,0);
  \def\Rr{0.85}
  \coordinate (Aa) at ($(Oo)+(148:\Rr)$);
  \coordinate (Pq) at ($(Oo)+(90:\Rr)$);
  \coordinate (Qq) at ($(Oo)+(48:\Rr)$);
  \coordinate (Bb) at ($(Oo)+(0:\Rr)$);
  \draw[given] (Oo) circle[radius=\Rr];
  \draw[given] (Oo) -- (Aa);
  \draw[given] (Oo) -- (Bb);
  \draw[aux] (Oo) -- (Pq);
  \draw[aux] (Oo) -- (Qq);
  \draw[aux] (Aa) -- (Pq) -- (Qq) -- (Bb);
  \node[vlab] at ($(Oo)+(148:\Rr+0.30)$) {$A$};
  \node[vlab] at ($(Oo)+(90:\Rr+0.30)$) {$P$};
  \node[vlab] at ($(Oo)+(48:\Rr+0.30)$) {$Q$};
  \node[vlab] at ($(Oo)+(0:\Rr+0.30)$) {$B$};
  \node[vlab] at ($(Oo)+(250:0.36)$) {$O$};
  \node[slab] at ($(Oo)+(0:0.44)+(0,-0.38)$) {$r$};
\end{tikzpicture}
\figcap{12--12}
\end{minipage}
\end{bkfigure}}

%---------------------------------------------------- Figs 12-13 and 12-14
% 12-13: OC bisects the right angle AOB, OD bisects BOC, OE bisects DOC.
%   A = 0, B = 90, C = 45, D = 67.5, E = 56.25 -- pure halving, no eyeballing.
% 12-14: rebuilt 2026-08-17 from the book plate.  Four faults, all Caleb's:
%   * it was drawn as an axis-aligned 90-degree quarter.  The book draws a
%     roughly 68-degree sector, tilted so it is symmetric about the horizontal
%     and neither bounding radius is level: A(-34), B(+34).
%   * it carried five interior rays running clear THROUGH the outer marker, so
%     the crossings read as tick marks on it and the outer region looked
%     hatched.  The book has exactly four, overshooting the arc only slightly
%     and stopping well short of the marker.
%   * the interior rays are the same subdivisions 12-13 constructs, read as
%     fractions of the sector measured up from OA: 1/2, 5/8, 3/4, all dashed,
%     plus the SOLID ray at 0.662 -- "a little greater than 5/8", exactly what
%     the text claims of L.  That is why the book's third ray sits a hair below
%     the solid one.
%   * the l marker was a plain arc spanning the WHOLE sector with no hooks.  It
%     is a brace, it carries end hooks and a middle nib, and it spans only the
%     sub-arc AP whose length is l.
\newcommand{\FIGXIITHIRTEENFOURTEEN}{\begin{bkfigure}[\linewidth]
\begin{minipage}[b]{0.48\linewidth}\centering
\begin{tikzpicture}[scale=1.194]
  \coordinate (O) at (0,0);
  \def\R{1.6}
  \coordinate (A) at ($(O)+(0:\R)$);
  \coordinate (B) at ($(O)+(90:\R)$);
  \coordinate (C) at ($(O)+(45:\R)$);
  \coordinate (D) at ($(O)+(67.5:\R)$);
  \coordinate (E) at ($(O)+(56.25:\R)$);
  \draw[given] (O) circle[radius=\R];
  \draw[given] (O) -- (B);
  \draw[given] (O) -- (A);
  \draw[aux] (O) -- (D);
  \draw[aux] (O) -- (E);
  \draw[aux] (O) -- (C);
  \node[vlab] at ($(O)+(90:\R+0.28)$) {$B$};
  \node[vlab] at ($(O)+(67.5:\R+0.28)$) {$D$};
  \node[vlab] at ($(O)+(56.25:\R+0.28)$) {$E$};
  \node[vlab] at ($(O)+(45:\R+0.28)$) {$C$};
  \node[vlab] at ($(O)+(0:\R+0.28)$) {$A$};
  \node[vlab] at ($(O)+(250:0.34)$) {$O$};
  \node[slab] at ($(O)+(90:0.80)+(-0.28,0)$) {$r$};
\end{tikzpicture}
\figcap{12--13}
\end{minipage}\hfill
\begin{minipage}[b]{0.44\linewidth}\centering
\begin{tikzpicture}[scale=1.099]
  \coordinate (O) at (0,0);
  \def\R{1.7}
  % the sector opens 68 degrees, symmetric about the horizontal
  \def\Alo{-34}\def\Bhi{34}
  \coordinate (A) at ($(O)+(\Alo:\R)$);
  \coordinate (B) at ($(O)+(\Bhi:\R)$);
  % interior rays as fractions of the opening measured up from OA
  \pgfmathsetmacro{\rhalf}{\Alo+0.500*(\Bhi-\Alo)}
  \pgfmathsetmacro{\rfive}{\Alo+0.625*(\Bhi-\Alo)}
  \pgfmathsetmacro{\rP}{\Alo+0.662*(\Bhi-\Alo)}
  \pgfmathsetmacro{\rthree}{\Alo+0.750*(\Bhi-\Alo)}
  \coordinate (P) at ($(O)+(\rP:\R)$);
  \draw[given] (O) -- (B);
  \draw[given] (O) -- (A);
  \draw[given] ($(O)+(\Alo:\R)$)
      arc[start angle=\Alo, end angle=\Bhi, radius=\R];
  \foreach \a in {\rhalf, \rfive, \rthree}
    {\draw[aux] (O) -- ($(O)+(\a:\R+0.15)$);}
  \draw[given] (O) -- ($(O)+(\rP:\R+0.15)$);
  % brace marking the arc AP, of length l: hooks at both ends, nib at the
  % middle pointing out at the label, clear of the rays' overshoot
  \XIIbrace{O}{\R+0.36}{\rP}{\Alo}{0.14}
  \node[vlab] at ($(O)+(\Bhi:\R)+(0.06,0.28)$) {$B$};
  \node[vlab] at ($(O)+(\Alo:\R)+(0.06,-0.28)$) {$A$};
  \node[vlab] at ($(O)+(196:0.30)$) {$O$};
  % r sits INSIDE the sector, against OB, as the book sets it -- the render had
  % it outside, on the far side of the radius.  It is centred in the 17-degree
  % wedge between OB and the 3/4 ray, far enough out that the wedge is wide
  % enough to hold it clear of both.
  \node[slab] at ($(O)+(25.5:1.35)$) {$r$};
  \pgfmathsetmacro{\lmid}{0.5*(\rP+\Alo)}
  \node[slab] at ($(O)+(\lmid:\R+0.74)$) {$l$};
\end{tikzpicture}
\figcap{12--14}
\end{minipage}
\end{bkfigure}}

%---------------------------------------------------------------- Fig 12-15
% Successive 1-degree arcs counted from B toward A, exaggerated for legibility.
% A(48) lies strictly between the n-degree ray (32) and the (n+1)-degree ray
% (57), which is exactly the inequality the theorem asserts.
%
% Rebuilt 2026-08-17.  The two arc-measure markers were plain concentric arcs
% -- the n-degree one with straight end ticks, the (n+1)-degree one with
% nothing at all and its label floating free.  The book draws both as BRACES:
% inward hooks at each end and a nib at the middle pointing straight at the
% label.  The 1-degree marker was a bare line with a crossbar dropped far below
% the figure; the book tucks a short headed arrow into the 1-degree wedge with
% the label immediately under OB.
\newcommand{\FIGXIIFIFTEEN}{\begin{bkfigure}[\linewidth]
\begin{tikzpicture}[scale=1.300]
  \coordinate (O) at (0,0);
  \def\R{1.95}
  \coordinate (B) at ($(O)+(0:\R)$);
  \coordinate (A) at ($(O)+(48:\R)$);
  \draw[given] (O) -- (B);
  \draw[given] (O) -- (A);
  \draw[given] ($(O)+(0:\R)$) arc[start angle=0, end angle=48, radius=\R];
  % each dashed ray runs out only as far as the marker it serves: the 32-degree
  % ray closes the n-degree brace, the 57-degree ray closes the (n+1)-degree
  % brace, and the 10-degree ray merely clears the arc
  \foreach \a/\len in {10/0.10, 32/0.24, 57/0.92}
    {\draw[aux] (O) -- ($(O)+(\a:\R+\len)$);}
  \XIIbrace{O}{\R+0.20}{0}{32}{0.12}
  \XIIbrace{O}{\R+0.80}{0}{57}{0.14}
  % the 1-degree wedge: one short arrow rising from under OB, its head landing
  % inside the wedge between OB and the 10-degree ray
  \draw[XIIhook, -{Stealth[length=4pt,width=3pt]}] (1.28,-0.34) -- (1.28,0.15);
  \node[vlab] at ($(O)+(200:0.34)$) {$O$};
  \node[vlab] at ($(O)+(0:\R)+(0.16,-0.26)$) {$B$};
  \node[vlab] at ($(O)+(48:\R+0.34)$) {$A$};
  \node[slab] at ($(O)+(32:1.45)+(-58:0.32)$) {$r$};
  \node[slab, anchor=west] at ($(O)+(16:\R+0.38)$) {$n\XIId$};
  \node[slab, anchor=west] at ($(O)+(28.5:\R+1.00)$) {$(n+1)\XIId$};
  \node[slab] at (1.28,-0.60) {$1\XIId$};
\end{tikzpicture}
\figcap{12--15}
\end{bkfigure}}

%---------------------------------------------------------------- Fig 12-16
% The primed twin of Fig 12-15: angle A'O'B' congruent to angle AOB, counted
% the same way, so A' falls on the same (n+1)st arc.
% 2026-08-17: the radius was labelled a bare r here.  EVERY letter in this
% figure is primed in the book -- it is the primed twin -- so it is r', and
% Caleb flagged the bare r by name.  Markers rebuilt as braces, as in 12-15.
\newcommand{\FIGXIISIXTEEN}{\begin{bkfigure}[\linewidth]
\begin{tikzpicture}[scale=1.300]
  \coordinate (O) at (0,0);
  \def\R{1.95}
  \coordinate (B) at ($(O)+(0:\R)$);
  \coordinate (A) at ($(O)+(48:\R)$);
  \draw[given] (O) -- (B);
  \draw[given] (O) -- (A);
  \draw[given] ($(O)+(0:\R)$) arc[start angle=0, end angle=48, radius=\R];
  \foreach \a/\len in {10/0.10, 32/0.24, 57/0.92}
    {\draw[aux] (O) -- ($(O)+(\a:\R+\len)$);}
  \XIIbrace{O}{\R+0.20}{0}{32}{0.12}
  \XIIbrace{O}{\R+0.80}{0}{57}{0.14}
  \draw[XIIhook, -{Stealth[length=4pt,width=3pt]}] (1.28,-0.34) -- (1.28,0.15);
  \node[vlab] at ($(O)+(200:0.36)$) {$O\XIIp$};
  \node[vlab] at ($(O)+(0:\R)+(0.20,-0.26)$) {$B\XIIp$};
  \node[vlab] at ($(O)+(48:\R+0.34)$) {$A\XIIp$};
  \node[slab] at ($(O)+(32:1.45)+(-58:0.34)$) {$r\XIIp$};
  \node[slab, anchor=west] at ($(O)+(16:\R+0.38)$) {$n\XIId$};
  \node[slab, anchor=west] at ($(O)+(28.5:\R+1.00)$) {$(n+1)\XIId$};
  \node[slab] at (1.28,-0.60) {$1\XIId$};
\end{tikzpicture}
\figcap{12--16}
\end{bkfigure}}

%--------------------------------------------------- Figs 12-17 and 12-18
% 12-17: left, the arc subtended by a central angle of A degrees; right, the
% arc EXTERIOR to a congruent central angle, of 360 - A degrees.  The two
% central angles are drawn with the same opening (75 degrees).
% 12-18: one radian -- the central angle whose subtended arc equals the radius,
% so the opening is exactly 180/pi = 57.2958 degrees.
\newcommand{\FIGXIISEVENTEENEIGHTEEN}{\begin{bkfigure}[\linewidth]
\begin{minipage}[b]{0.62\linewidth}\centering
\begin{tikzpicture}[scale=0.76]
  % left circle: interior arc
  \coordinate (P) at (0,0);
  % 1.15, not 0.95: the interior A-degree numeral has to sit inside a 75-degree
  % wedge, and on the smaller radius its box overhangs the upper ray.
  \def\Ra{1.15}
  \draw[given] (P) circle[radius=\Ra];
  \draw[given] (P) -- ($(P)+(37.5:\Ra)$);
  \draw[given] (P) -- ($(P)+(-37.5:\Ra)$);
  \draw[given] ($(P)+(-37.5:0.26)$) arc[start angle=-37.5, end angle=37.5, radius=0.20];
  \node[slab] at ($(P)+(0:0.70)$) {$A\XIId$};
  \draw[XIIhook] ($(P)+(-37.5:\Ra+0.22)$)
      arc[start angle=-37.5, end angle=37.5, radius=\Ra+0.22];
  \draw[XIIhook] ($(P)+(37.5:\Ra+0.22)$) -- ($(P)+(37.5:\Ra+0.10)$);
  \draw[XIIhook] ($(P)+(-37.5:\Ra+0.22)$) -- ($(P)+(-37.5:\Ra+0.10)$);
  \node[slab] at ($(P)+(0:\Ra+0.58)$) {$A\XIId$};
  % right circle: exterior arc
  \coordinate (Q) at (3.95,0);
  \def\Rb{1.2}
  \draw[given] (Q) circle[radius=\Rb];
  \draw[given] (Q) -- ($(Q)+(142.5:\Rb)$);
  \draw[given] (Q) -- ($(Q)+(217.5:\Rb)$);
  \draw[given] ($(Q)+(142.5:0.28)$) arc[start angle=142.5, end angle=217.5, radius=0.24];
  \node[slab] at ($(Q)+(180:0.74)$) {$A\XIId$};
  \draw[XIIhook] ($(Q)+(142.5:\Rb+0.24)$)
      arc[start angle=142.5, end angle=-142.5, radius=\Rb+0.24];
  \draw[XIIhook] ($(Q)+(142.5:\Rb+0.24)$) -- ($(Q)+(142.5:\Rb+0.12)$);
  \draw[XIIhook] ($(Q)+(217.5:\Rb+0.24)$) -- ($(Q)+(217.5:\Rb+0.12)$);
  \node[slab, right] at ($(Q)+(0:\Rb+0.34)$) {$(360\XIId-A\XIId\mkern1.3mu)$};
\end{tikzpicture}
\figcap{12--17}
\end{minipage}\hfill
\begin{minipage}[b]{0.34\linewidth}\centering
\begin{tikzpicture}[scale=0.86]
  \coordinate (O) at (0,0);
  \def\R{1.3}
  \coordinate (U) at ($(O)+(50:\R)$);
  \coordinate (V) at ($(O)+(-7.2958:\R)$);
  \draw[given] (O) circle[radius=\R];
  \draw[given] (O) -- (U);
  \draw[given] (O) -- (V);
  \draw[given] ($(O)+(-7.2958:0.44)$)
      arc[start angle=-7.2958, end angle=50, radius=0.38];
  \node[slab] at ($(O)+(21.35:0.82)$) {$1$};
  \node[slab] at ($(O)!0.58!(U)+(-0.30,0.18)$) {$r$};
  \node[slab] at ($(O)+(21.35:\R+0.28)$) {$r$};
\end{tikzpicture}
\figcap{12--18}
\end{minipage}
\end{bkfigure}}

%--------------------------------------------------- Figs 12-19 and 12-20
% 12-19: two inscribed angles (vertex ON the circle) and one that is not
% (vertex OUTSIDE the circle) -- the distinction the definition turns on.
% 12-20: BC a diameter, so the inscribed angle B stands on the arc AC; OA is
% the radius drawn for the lemma's isosceles triangle OAB.
\newcommand{\FIGXIININETEENTWENTY}{\begin{bkfigure}[\linewidth]
\begin{minipage}[b]{0.66\linewidth}\centering
\begin{tikzpicture}[scale=0.74]
  \def\Rc{0.95}
  % (a) inscribed
  \coordinate (Oa) at (0,0);
  \coordinate (Ba) at ($(Oa)+(168:\Rc)$);
  \coordinate (Aa) at ($(Oa)+(12:\Rc)$);
  \coordinate (Ca) at ($(Oa)+(-55:\Rc)$);
  \draw[given] (Oa) circle[radius=\Rc];
  \draw[given] (Ba) -- (Aa);
  \draw[given] (Ba) -- (Ca);
  \node[vlab] at ($(Oa)+(168:\Rc+0.34)$) {$B$};
  \node[vlab] at ($(Aa)+(48:0.36)$) {$A$};
  \node[vlab] at ($(Oa)+(-55:\Rc+0.34)$) {$C$};
  \node[slab] at ($(Oa)+(0,-1.85)$) {Inscribed};
  % (b) inscribed
  \coordinate (Ob) at (3.15,0);
  \coordinate (Bb) at ($(Ob)+(232:\Rc)$);
  \coordinate (Ab) at ($(Ob)+(119:\Rc)$);
  \coordinate (Cb) at ($(Ob)+(-43:\Rc)$);
  \draw[given] (Ob) circle[radius=\Rc];
  \draw[given] (Bb) -- (Ab);
  \draw[given] (Bb) -- (Cb);
  \node[vlab] at ($(Ab)+(136:0.50)$) {$A$};
  \node[vlab] at ($(Ob)+(232:\Rc+0.34)$) {$B$};
  \node[vlab] at ($(Ob)+(-43:\Rc+0.34)$) {$C$};
  \node[slab] at ($(Ob)+(0,-1.85)$) {Inscribed};
  % (c) NOT inscribed.  Measured off scan14 p.17: the vertex B sits a whisker
  % OUTSIDE the circle (1.08 r) and the side BA runs AWAY from it, so the whole
  % ray BA misses the circle -- that is what fails the definition.  Side BC, by
  % contrast, cuts straight across.  (Directions 215.6 / 111.8 / -45.3 and the
  % 1.5 r length of BA are read straight off the plate.)
  \coordinate (Oc) at (6.30,0);
  \coordinate (Bc) at ($(Oc)+(215.6:1.08*\Rc)$);
  \coordinate (Ac) at ($(Bc)+(111.8:1.50*\Rc)$);
  \coordinate (Cc) at ($(Oc)+(-45.3:\Rc)$);
  \draw[given] (Oc) circle[radius=\Rc];
  \draw[given] (Ac) -- (Bc);
  \draw[given] (Bc) -- (Cc);
  \node[vlab] at ($(Ac)+(50:0.34)$) {$A$};
  \node[vlab] at ($(Bc)+(205:0.34)$) {$B$};
  \node[vlab] at ($(Cc)+(-40:0.34)$) {$C$};
  \node[slab] at ($(Oc)+(0,-1.85)$) {Not inscribed};
\end{tikzpicture}
\figcap{12--19}
\end{minipage}\hfill
\begin{minipage}[b]{0.30\linewidth}\centering
\begin{tikzpicture}[scale=0.86]
  \coordinate (O) at (0,0);
  \def\R{1.25}
  \coordinate (B) at ($(O)+(180:\R)$);
  \coordinate (C) at ($(O)+(0:\R)$);
  \coordinate (A) at ($(O)+(79:\R)$);
  \draw[given] (O) circle[radius=\R];
  \draw[given] (B) -- (C);
  \draw[given] (B) -- (A);
  \draw[aux]   (A) -- (O);   % the book prints this radius DASHED
  \node[vlab] at ($(O)+(79:\R+0.30)$) {$A$};
  \node[vlab] at ($(O)+(180:\R+0.30)$) {$B$};
  \node[vlab] at ($(O)+(0:\R+0.30)$) {$C$};
  \node[vlab] at ($(O)+(255:0.36)$) {$O$};
\end{tikzpicture}
\figcap{12--20}
\end{minipage}
\end{bkfigure}}

%--------------------------------------------------- Figs 12-21 and 12-22
% 12-21 (Case 1): the centre is INTERIOR to angle ABC.  The diameter BP splits
%   arc AC into AP(80) and PC(33), and angle ABC into x = 40 and y = 16.5.
%   Directions read off scan14 p.18: A 80, C -33.  The two angle arcs at B are
%   in the book and are what tie the letters x, y to the angles they name.
% 12-22 (Case 2): the centre is EXTERIOR to angle ABC -- both rays BA, BC lie
%   on the same side of the diameter BQ.  Measured: A 110, C 36.
%   Rebuilt 2026-08-17; three faults, all Caleb's, all confirmed against the
%   plate at high zoom:
%     * a chord AC was drawn that the book does NOT have, closing a triangle
%       ABC.  Between A and C the book has only the circle's arc.  Removed.
%     * BOTH angle arcs at B were missing.  The book strikes two thin
%       concentric arcs centred at B, the outer springing from A itself and the
%       inner from about 55% of the way along chord BA; each sweeps clockwise
%       ACROSS chord BC and lands on the dashed diameter BQ with a short
%       downward tick.  They are what carry the subtraction the proof needs,
%       angle ABC = angle ABQ - angle CBQ.
%     * the centre carried no dot.  The book marks it, and sets O above-left.
\newcommand{\FIGXIITWENTYONETWENTYTWO}{\begin{bkfigure}[\linewidth]
\begin{minipage}[b]{0.46\linewidth}\centering
\begin{tikzpicture}[scale=1.091]
  \coordinate (O) at (0,0);
  % 1.6: y names a 16.5-degree wedge, and its numeral needs about 1.5 units of
  % run from B before the wedge is tall enough to hold it clear of both sides.
  \def\R{1.6}
  \coordinate (B) at ($(O)+(180:\R)$);
  \coordinate (P) at ($(O)+(0:\R)$);
  \coordinate (A) at ($(O)+(80:\R)$);
  \coordinate (C) at ($(O)+(-33:\R)$);
  \draw[given] (O) circle[radius=\R];
  \draw[given] (B) -- (A);
  \draw[given] (A) -- (C);
  \draw[given] (B) -- (C);
  \draw[aux] (B) -- (P);
  % angle arcs: x = angle ABP (40 deg), y = angle PBC (16.5 deg)
  \draw[XIIhook] ($(B)+(0:0.60)$)     arc[start angle=0,     end angle=40, radius=0.60];
  \draw[XIIhook] ($(B)+(-16.5:0.90)$) arc[start angle=-16.5, end angle=0,  radius=0.90];
  \node[vlab] at ($(O)+(80:\R+0.30)$) {$A$};
  \node[vlab] at ($(O)+(180:\R+0.30)$) {$B$};
  \node[vlab] at ($(O)+(0:\R+0.30)$) {$P$};
  \node[vlab] at ($(O)+(-33:\R+0.30)$) {$C$};
  \node[vlab] at ($(O)+(90:0.30)$) {$O$};
  \node[slab] at ($(B)+(20:1.00)$) {$x$};
  \node[slab] at ($(B)+(-8.25:1.85)$) {$y$};
\end{tikzpicture}
\figcap{12--21}
\end{minipage}\hfill
\begin{minipage}[b]{0.46\linewidth}\centering
\begin{tikzpicture}[scale=1.091]
  \coordinate (O) at (0,0);
  \def\R{1.35}
  \coordinate (B) at ($(O)+(180:\R)$);
  \coordinate (Q) at ($(O)+(0:\R)$);
  \coordinate (A) at ($(O)+(110:\R)$);
  \coordinate (C) at ($(O)+(36:\R)$);
  \draw[given] (O) circle[radius=\R];
  \draw[given] (B) -- (A);
  \draw[given] (B) -- (C);
  \draw[aux] (B) -- (Q);
  % the two angle arcs at B.  Their common span is angle ABQ, so both are
  % DERIVED from the directions B->A and B->Q, never swept by eye; the outer
  % radius is |BA| exactly, so that arc springs from A itself.
  % chord direction and chord length on a circle, in closed form: the ray from
  % the point at angle p to the point at angle q runs at (p+q)/2 - 90, and the
  % chord measures 2 r sin((p-q)/2).  B is at 180, A at 110.
  \pgfmathsetmacro{\bA}{0.5*(180+110)-90}
  \pgfmathsetmacro{\bBA}{2*\R*sin(0.5*(180-110))}
  \foreach \f in {1.0, 0.55}{%
    \pgfmathsetmacro{\rr}{\f*\bBA}
    \draw[XIIhook] ($(B)+(\bA:\rr)$)
        arc[start angle=\bA, end angle=0, radius=\rr]
        -- ($(B)+(0:\rr)+(0,-0.11)$);}
  \dt{O}
  \node[vlab] at ($(O)+(110:\R+0.30)$) {$A$};
  \node[vlab] at ($(O)+(180:\R+0.30)$) {$B$};
  \node[vlab] at ($(O)+(36:\R+0.30)$) {$C$};
  \node[vlab] at ($(O)+(0:\R+0.30)$) {$Q$};
  % The book sets O above-left of the centre dot, but the wedge there -- between
  % the dashed diameter and chord BC -- is only 0.42 r deep at the centre and
  % will not hold a footnotesize letter clear of both.  Below-left is the one
  % free quadrant; the dot itself now fixes which point the letter names.
  \node[vlab] at ($(O)+(250:0.32)$) {$O$};
\end{tikzpicture}
\figcap{12--22}
\end{minipage}
\end{bkfigure}}

%---------------------------------------------------------------- Fig 12-23
% Angle inscribed in a semicircle: BC is a diameter, so angle BAC is right.
\newcommand{\FIGXIITWENTYTHREE}{\begin{bkfigure}[\linewidth]
\begin{tikzpicture}[scale=1.300]
  \coordinate (O) at (0,0);
  \def\R{1.25}
  \coordinate (B) at ($(O)+(180:\R)$);
  \coordinate (C) at ($(O)+(0:\R)$);
  \coordinate (A) at ($(O)+(105:\R)$);
  \draw[given] (O) circle[radius=\R];
  \draw[given] (B) -- (C);
  \draw[given] (B) -- (A);
  \draw[given] (A) -- (C);
  \node[vlab] at ($(O)+(105:\R+0.30)$) {$A$};
  \node[vlab] at ($(O)+(180:\R+0.30)$) {$B$};
  \node[vlab] at ($(O)+(0:\R+0.30)$) {$C$};
  \node[vlab] at ($(O)+(265:0.34)$) {$O$};
\end{tikzpicture}
\figcap{12--23}
\end{bkfigure}}

%---------------------------------------------------------------- Fig 12-24
% Split out of the old \FIGXIITWENTYFOURTWENTYFIVE: 12-24 belongs to Exercise 7
% of Group 12-11 and 12-25 to Exercise *9, so they sit inline after different
% exercises.
% 12-24: AB and CD are parallel secants.  Their offsets from the centre are
%   DELIBERATELY unequal (0.55 r above, 0.68 r below -- both measured off the
%   plate).  Equal offsets would make AB congruent to CD, a coincidence the
%   exercise never grants; arc AC = arc BD holds for ANY pair of parallel
%   chords, which is exactly what Exercise 7 asks the student to prove.
% 12-25: AC is tangent at B (perpendicular to OB by construction); D and E are
%   placed symmetrically about OB, which makes DE parallel to AC and the arcs
%   BD and BE equal -- both facts the exercise's hint needs.
\newcommand{\FIGXIITWENTYFOUR}{\begin{bkfigure}[\linewidth]
\begin{tikzpicture}[scale=1.162]
  \coordinate (O) at (0,0);
  \def\R{1.25}
  \coordinate (A) at (-1.0440, 0.6875);
  \coordinate (B) at ( 1.0440, 0.6875);
  \coordinate (C) at (-0.9165,-0.8500);
  \coordinate (D) at ( 0.9165,-0.8500);
  \draw[given] (O) circle[radius=\R];
  \draw[given] (-1.52, 0.6875) -- (1.52, 0.6875);
  \draw[given] (-1.44,-0.8500) -- (1.44,-0.8500);
  \node[vlab] at ($(A)+(112:0.36)$) {$A$};
  \node[vlab] at ($(B)+(68:0.36)$) {$B$};
  \node[vlab] at ($(C)+(248:0.36)$) {$C$};
  \node[vlab] at ($(D)+(292:0.36)$) {$D$};
\end{tikzpicture}
\figcap{12--24}
\end{bkfigure}}

%---------------------------------------------------------------- Fig 12-25
% AC is tangent at B (perpendicular to OB by construction); D and E are placed
% symmetrically about OB, which makes DE parallel to AC and the arcs BD and BE
% equal -- both facts the exercise's hint needs.
\newcommand{\FIGXIITWENTYFIVE}{\begin{bkfigure}[\linewidth]
\begin{tikzpicture}[scale=1.221]
  \coordinate (O) at (0,0);
  \def\R{1.05}
  \coordinate (B) at ($(O)+(133:\R)$);
  \coordinate (A) at ($(B)+(223:1.22)$);
  \coordinate (C) at ($(B)+(43:1.06)$);
  \coordinate (D) at ($(O)+(199:\R)$);
  \coordinate (E) at ($(O)+(67:\R)$);
  \draw[given] (O) circle[radius=\R];
  \draw[given] (A) -- (C);
  \draw[aux] (B) -- (D);
  \draw[aux] (B) -- (E);
  \draw[aux] (D) -- (E);
  \node[vlab] at ($(A)+(223:0.34)$) {$A$};
  \node[vlab] at ($(O)+(133:\R+0.36)$) {$B$};
  \node[vlab] at ($(C)+(43:0.34)$) {$C$};
  \node[vlab] at ($(O)+(199:\R+0.34)$) {$D$};
  \node[vlab] at ($(O)+(67:\R+0.34)$) {$E$};
\end{tikzpicture}
\figcap{12--25}
\end{bkfigure}}

%---------------------------------------------------------------- Fig 12-26
% Split out of the old \FIGXIITWENTYSIXTWENTYSEVEN: 12-26 is first cited by
% Exercise *10 of Group 12-11 and 12-27 by Exercise *12, so each stands inline
% after its own item.
% 12-26: two secants from the external point A.  A is DERIVED as the crossing
%   of lines CB and DE, so A, B, C and A, E, D are exactly collinear.  The
%   auxiliary chord is B--D, NOT C--D: the book draws BD, and both exercises
%   that lean on this figure need triangle ABD -- Ex *10 gets angle A as the
%   exterior angle of ABD at B, and Ex *14 proves ABD ~ AEC.  The upper secant
%   runs on past C, as in the book, to show it is a line and not a segment.
\newcommand{\FIGXIITWENTYSIX}{\begin{bkfigure}[\linewidth]
\begin{tikzpicture}[scale=0.999]
  \coordinate (O) at (0,0);
  \def\R{1.1}
  \coordinate (B) at ($(O)+(151.6:\R)$);
  \coordinate (C) at ($(O)+(79.5:\R)$);
  \coordinate (D) at ($(O)+(-31:\R)$);
  \coordinate (E) at ($(O)+(213:\R)$);
  \coordinate (A) at (intersection of C--B and D--E);
  \draw[given] (O) circle[radius=\R];
  \draw[given] (A) -- ($(A)!1.12!(C)$);
  \draw[given] (A) -- (D);
  \draw[given] (B) -- (D);
  \node[vlab] at ($(A)+(186:0.34)$) {$A$};
  \node[vlab] at ($(B)+(88:0.50)$) {$B$};
  \node[vlab] at ($(O)+(95:\R+0.36)$) {$C$};
  \node[vlab] at ($(E)+(265:0.44)$) {$E$};
  \node[vlab] at ($(O)+(-31:\R+0.34)$) {$D$};
\end{tikzpicture}
\figcap{12--26}
\end{bkfigure}}

%---------------------------------------------------------------- Fig 12-27
% Chords AC and BD crossing at E, again derived as an intersection; BC dashed
% is the hint line for Exercise *12.
\newcommand{\FIGXIITWENTYSEVEN}{\begin{bkfigure}[\linewidth]
\begin{tikzpicture}[scale=1.140]
  \coordinate (O) at (0,0);
  \def\R{1.1}
  \coordinate (A) at ($(O)+(190:\R)$);
  \coordinate (B) at ($(O)+(135:\R)$);
  \coordinate (C) at ($(O)+(10:\R)$);
  \coordinate (D) at ($(O)+(-68:\R)$);
  \coordinate (E) at (intersection of A--C and B--D);
  \draw[given] (O) circle[radius=\R];
  \draw[given] (A) -- (C);
  \draw[given] (B) -- (D);
  \draw[aux] (B) -- (C);
  \node[vlab] at ($(O)+(190:\R+0.34)$) {$A$};
  \node[vlab] at ($(O)+(155:\R+0.36)$) {$B$};
  \node[vlab] at ($(O)+(10:\R+0.34)$) {$C$};
  \node[vlab] at ($(O)+(-68:\R+0.34)$) {$D$};
  \node[vlab] at ($(E)+(247:0.36)$) {$E$};
\end{tikzpicture}
\figcap{12--27}
\end{bkfigure}}

%---------------------------------------------------------------- Fig 12-28
% Two chords crossing at O.  2026-08-17: this used to be SOLVED so that the
% drawing reproduced the numbers quoted in Exercise 15 (AO:OB = 12:4 = 3 and
% CO:OD = 8:6 = 4/3) -- which handed the student the answer off a ruler.
% Caleb asked for the book's own schematic cut instead.  Measured off the
% plate, book p.224: A 198, B 31, C 105, D -18, giving AO:OB = 2.76 and
% CO:OD = 1.11.  Neither matches the quoted ratio, which is the point; what
% the figure still asserts exactly is the STRUCTURE -- four points on the
% circle, two chords, one crossing strictly inside.
\newcommand{\FIGXIITWENTYEIGHT}{\begin{bkfigure}[\linewidth]
\begin{tikzpicture}[scale=1.495]
  \coordinate (Ctr) at (0,0);
  \def\R{1.1}
  \coordinate (A) at ($(Ctr)+(198:\R)$);
  \coordinate (B) at ($(Ctr)+(31:\R)$);
  \coordinate (C) at ($(Ctr)+(105:\R)$);
  \coordinate (D) at ($(Ctr)+(-18:\R)$);
  \coordinate (X) at (intersection of A--B and C--D);
  \draw[given] (Ctr) circle[radius=\R];
  \draw[given] (A) -- (B);
  \draw[given] (C) -- (D);
  \node[vlab] at ($(Ctr)+(198:\R+0.30)$) {$A$};
  \node[vlab] at ($(Ctr)+(31:\R+0.30)$) {$B$};
  \node[vlab] at ($(Ctr)+(105:\R+0.30)$) {$C$};
  \node[vlab] at ($(Ctr)+(-18:\R+0.30)$) {$D$};
  \node[vlab] at ($(X)+(250:0.36)$) {$O$};
\end{tikzpicture}
\figcap{12--28}
\end{bkfigure}}

%---------------------------------------------------------------- Fig 12-29
% Two secants from an external point A.  2026-08-17: this was SOLVED from the
% numbers in Exercise *16 (AD = 6, DE = 10, AB = 8, hence BC = 4), which is
% precisely what Caleb asked us to stop doing -- a scale drawing answers the
% exercise for the student.  It is now the book's own schematic cut, measured
% off the plate on p.225: D 186, E 47, B 203, C -21, and A DERIVED as the
% crossing of the two secants, so collinearity is exact rather than plotted.
% Two further faults from the same photo pair, both fixed by these angles:
%   * the lower secant used to graze the circle -- B and C landed close
%     together on the bottom arc.  The book runs BC as a long, near-horizontal
%     chord across almost the whole circle, with C on the RIGHT.
%   * A sat about 1.5 diameters out with a 25-30 degree wedge.  The book keeps
%     A close in (about 0.6 r clear) and the wedge compact.
% What the drawing still asserts exactly: four points on the circle, both
% secants straight through A, D between A and E, B between A and C.
\newcommand{\FIGXIITWENTYNINE}{\begin{bkfigure}[\linewidth]
\begin{tikzpicture}[scale=1.140]
  \coordinate (O) at (0,0);
  \def\R{1.4}
  \coordinate (D) at ($(O)+(186:\R)$);
  \coordinate (E) at ($(O)+(47:\R)$);
  \coordinate (B) at ($(O)+(203:\R)$);
  \coordinate (C) at ($(O)+(-21:\R)$);
  \coordinate (A) at (intersection of D--E and B--C);
  \draw[given] (O) circle[radius=\R];
  \draw[given] (A) -- (E);
  \draw[given] (A) -- (C);
  \node[vlab] at ($(A)+(196:0.34)$) {$A$};
  \node[vlab] at ($(D)+(120:0.32)$) {$D$};
  \node[vlab] at ($(B)+(255:0.32)$) {$B$};
  \node[vlab] at ($(O)+(47:\R+0.32)$) {$E$};
  \node[vlab] at ($(O)+(-21:\R+0.32)$) {$C$};
\end{tikzpicture}
\figcap{12--29}
\end{bkfigure}}

%---------------------------------------------------------------- Fig 12-30
% Tangent plus secant from an external point.  2026-08-17: likewise reverted
% from the solved cut (AC = 8, CD = 10, AB = 12) to the book's schematic.  The
% two structural facts the exercises actually lean on are now built in, and
% both were broken before:
%   * AB is TANGENT at B.  The old cut met the circle at a visibly oblique
%     angle to the radius, so it read as a clipped secant -- and Ex *17's whole
%     argument is the tangent case.  B is now derived from the tangent-length
%     construction, so OB is perpendicular to AB exactly.
%   * ACD is a true DIAMETER, level with A, as the book draws it; the old cut
%     sloped it down the right-hand side and it missed the centre entirely.
% With A on the line of the diameter, AB^2 = AC * AD then holds identically --
% but the drawn ratio AC:CD is 0.43, not the exercise's 0.8, so nothing is
% given away.
\newcommand{\FIGXIITHIRTY}{\begin{bkfigure}[\linewidth]
\begin{tikzpicture}[scale=1.140]
  \coordinate (O) at (0,0);
  \def\R{1.4}
  \def\OA{2.59}
  \coordinate (A) at (-\OA, 0);
  \coordinate (C) at ($(O)+(180:\R)$);
  \coordinate (D) at ($(O)+(0:\R)$);
  % tangency: the radius to B makes angle arccos(r/|OA|) with OA, which is the
  % construction, not a guess -- so OB is exactly perpendicular to AB.
  \pgfmathsetmacro{\Bang}{180-acos(\R/\OA)}
  \coordinate (B) at ($(O)+(\Bang:\R)$);
  \draw[given] (O) circle[radius=\R];
  \draw[given] (A) -- (B);
  \draw[given] (A) -- (D);
  \draw[aux] (B) -- (C);
  \draw[aux] (B) -- (D);
  \node[vlab] at ($(A)+(186:0.32)$) {$A$};
  \node[vlab] at ($(B)+(100:0.32)$) {$B$};
  \node[vlab] at ($(C)+(310:0.34)$) {$C$};
  \node[vlab] at ($(O)+(0:\R+0.32)$) {$D$};
\end{tikzpicture}
\figcap{12--30}
\end{bkfigure}}

%---------------------------------------------------------------- Fig 12-31
% A central right angle; the dashed chord is the segment of length sqrt(2) r
% against which the arc is compared.
\newcommand{\FIGXIITHIRTYONE}{\begin{bkfigure}[\linewidth]
\begin{tikzpicture}[scale=1.300]
  \coordinate (O) at (0,0);
  \def\R{1.25}
  \coordinate (T) at ($(O)+(90:\R)$);
  \coordinate (S) at ($(O)+(0:\R)$);
  \draw[given] (O) circle[radius=\R];
  \draw[given] (O) -- (T);
  \draw[given] (O) -- (S);
  \draw[aux] (T) -- (S);
  \node[slab] at ($(O)+(90:0.62)+(-0.26,0)$) {$r$};
\end{tikzpicture}
\figcap{12--31}
\end{bkfigure}}

% figures13.tex -- Chapter 13 (Loci and Sets) figures.
%
% Constrained points are DERIVED, never placed by hand:
%   parallels    -> equal ratios along the two sides   ($(X)!r!(Y)$)
%   perp. feet   -> ($(X)!(P)!(Y)$)
%   crossings    -> (intersection of A--B and C--D)
%   bisectors    -> ratio from the angle-bisector theorem
%   midpoints    -> ($(X)!0.5!(Y)$)
% so the hypothesis stated in the book holds exactly in the drawing.
% Hypotheses are re-asserted numerically in tools/constraints/ch13.py.
%
% Macro names are chapter-prefixed: \FIGXIII<NUMBERWORD>.
%
% Line treatment follows the book's own convention: solid where the book
% prints solid (given), dashed where the book prints dashed (aux).  In this
% chapter the dashed role is doing double duty, exactly as the book does it:
%   - a LOCUS that is being asserted rather than given (Figs 13-1, 13-2)
%   - auxiliary/proof lines (Figs 13-5 ... 13-7, 13-11, 13-12, 13-14)
%   - the second position of a rolling wheel (Fig 13-9)
%   - the constructed outer triangle A'B'C' (Fig 13-16)
% The book shades no region anywhere in this chapter, so no fill patterns.

\usetikzlibrary{arrows.meta}

% Primed labels (Fig 13-16).  Same 1.3pt drop as ch09: keeps the prime in the
% same "word" as its letter for the collision audit.
\newcommand{\XIIIpr}{\raisebox{-1.3pt}{\ensuremath{{}^{\prime}}}}

% Chapter-local label clearance, as ch09: 3.4pt clips when a labelled point
% lies ON a drawn line.
\tikzset{
  vlab/.append style={outer sep=2.9pt},
  slab/.append style={outer sep=2.9pt},
}

% Long, widely spaced dashes.  Caleb, 2026-08-17: in Figs 13-11/13-12/13-16 the
% standard `aux' pattern is so fine that at print size the dashed cevians read
% as solid and the figure turns into "a quadrilateral with diagonals".  Measured
% off the print, the book's dashes in those three figures run about 2.6 mm on
% and 1.4 mm off -- roughly double ours.
\tikzset{auxlong/.style={aux, dash pattern=on 6.6pt off 3.8pt}}

% Small filled arrowheads, at a FIXED printed size so they do not grow with the
% stroke weight or shrink with the picture's `scale='.  The book's arrowheads
% are small; ours were blotting where several met at one vertex (Fig 13-16).
\tikzset{
  bkarrow/.style  = {-{Latex[length=3.4pt,width=2.8pt]}},
  bkarrows/.style = {{Latex[length=3.4pt,width=2.8pt]}-{Latex[length=3.4pt,width=2.8pt]}},
}

% The book marks the concurrency point in Figs 13-11, 13-12 and the centre in
% Fig 13-14 with a small saltire/plus rather than a filled dot.  In 13-11 and
% 13-12 that mark is a bold six-rayed asterisk, clearly heavier than the lines
% that cross there -- \XIIIstar reproduces it.
\newcommand{\XIIIcross}[1]{%
  \draw[fig,line width=0.5pt] ($(#1)+(-0.10,-0.10)$) -- ($(#1)+(0.10,0.10)$);
  \draw[fig,line width=0.5pt] ($(#1)+(-0.10,0.10)$) -- ($(#1)+(0.10,-0.10)$);}
\newcommand{\XIIIplus}[1]{%
  \draw[fig,line width=0.5pt] ($(#1)+(-0.11,0)$) -- ($(#1)+(0.11,0)$);
  \draw[fig,line width=0.5pt] ($(#1)+(0,-0.11)$) -- ($(#1)+(0,0.11)$);}
\newcommand{\XIIIstar}[1]{%
  \foreach \a in {15,75,135}
    \draw[fig,line width=0.75pt] ($(#1)+(\a:0.15)$) -- ($(#1)+(\a+180:0.15)$);}

% ---------------------------------------------------------------- Fig 13-1
% Line l solid; the locus (a pair of parallels at distance d) dashed, one on
% each side.  The d-tick is drawn from l up to the upper parallel.
\newcommand{\FIGXIIIONE}{\begin{bkfigure}[\linewidth]
\begin{tikzpicture}[scale=1.300]
  \draw[given] (0,0) -- (6.10,0);
  \node[vlab, right] at (6.10,0) {$l$};
  \draw[aux] (0.28,0.82) -- (5.82,0.82);
  \draw[aux] (0.28,-0.82) -- (5.82,-0.82);
  % the distance tick, with the book's short cap at the top
  \draw[given] (1.62,0) -- (1.62,0.82);
  \draw[given] (1.50,0.82) -- (1.74,0.82);
  \node[slab, right] at (1.66,0.42) {$d$};
\end{tikzpicture}
\figcap{13--1}\end{bkfigure}}

% ---------------------------------------------------------------- Fig 13-2
% Two given parallels solid; the locus (the midline) dashed.  P sits on the
% midline, and the transversal tick carries d above and d below.
\newcommand{\FIGXIIITWO}{\begin{bkfigure}[\linewidth]
\begin{tikzpicture}[scale=1.300]
  \coordinate (T) at (0,0.86);
  \coordinate (Bm) at (0,-0.86);
  \draw[given] (0,0.86) -- (6.00,0.86);
  \draw[given] (0,-0.86) -- (6.00,-0.86);
  \draw[aux] (0.22,0) -- (5.78,0);
  % transversal: P is the midpoint of the cut, so the two d's are equal
  \coordinate (Pt) at (3.46,0.86);
  \coordinate (Pb) at (3.46,-0.86);
  \coordinate (P)  at ($(Pt)!0.5!(Pb)$);
  \draw[given] (Pt) -- (Pb);
  \node[vlab, below left] at (P) {$P$};
  \node[slab, right] at (3.50,0.43) {$d$};
  \node[slab, right] at (3.50,-0.43) {$d$};
\end{tikzpicture}
\figcap{13--2}\end{bkfigure}}

% ---------------------------------------------------------------- Fig 13-3
% Ex 13-1 #2: a 3-unit segment, ticked into units, with the locus only
% partly sketched -- five dots at height 1 over the segment, and one dot at
% distance 1 measured radially off the right endpoint.  Reproduced as drawn.
\newcommand{\FIGXIIITHREE}{\begin{bkfigure}[\linewidth]
\begin{tikzpicture}[scale=1.196]
  \draw[given] (0,0) -- (3,0);
  \foreach \x in {0,1,2,3} \draw[given] (\x,0) -- (\x,0.17);
  % five dots, spacing measured from the print (0.19 unit); the arrow lands on
  % the FOURTH dot, with three dots to its left and one to its right, as drawn
  \foreach \x in {1.48,1.67,1.86,2.05,2.24} \fill (\x,1.0) circle (0.052);
  \draw[given,<->] (2.05,0.06) -- (2.05,0.94);
  \node[slab, left] at (2.01,0.50) {$1$};
  % radial unit off the right end: 72 degrees, length exactly 1
  \coordinate (R) at (3,0);
  \coordinate (Rd) at ($(R)+(72:1)$);
  \fill (Rd) circle (0.052);
  \draw[given,<->] ($(R)!0.07!(Rd)$) -- ($(R)!0.93!(Rd)$);
  \node[slab, right] at ($(R)!0.5!(Rd)$) {$1$};
\end{tikzpicture}
\figcap{13--3}\end{bkfigure}}

% ---------------------------------------------------------------- Fig 13-4
% Ex 13-1 #5: fixed base of 4 units, free vertex V.
\newcommand{\FIGXIIIFOUR}{\begin{bkfigure}[\linewidth]
\begin{tikzpicture}[scale=1.196]
  \coordinate (L) at (0,0);
  \coordinate (R) at (3.00,0);
  % apex measured off the print: 0.741 of the base along, 0.463 of the base up
  \coordinate (V) at (2.223,1.389);
  \draw[given] (L) -- (R) -- (V) -- cycle;
  \node[vlab, above] at (V) {$V$};
  \node[slab, below] at ($(L)!0.5!(R)$) {$4$};
\end{tikzpicture}
\figcap{13--4}\end{bkfigure}}

% ---------------------------------------------------------------- Fig 13-5
% Thm 13-1.  C is the MIDPOINT of AB; the perpendicular bisector is solid and
% runs through C to P; PA and PB are the dashed proof lines.
\newcommand{\FIGXIIIFIVE}{\begin{bkfigure}[\linewidth]
\begin{tikzpicture}[scale=1.300]
  \coordinate (A) at (0,0);
  \coordinate (B) at (3.40,0);
  \coordinate (C) at ($(A)!0.5!(B)$);
  \coordinate (P) at ($(C)+(0,1.56)$);
  \draw[given] (A) -- (B);
  % the bisector overshoots P by about 0.15 of AB in the print, and drops
  % about 0.13 of AB below the segment
  \draw[given] ($(C)+(0,-0.44)$) -- ($(C)+(0,2.06)$);
  \draw[aux] (A) -- (P);
  \draw[aux] (P) -- (B);
  \fill (C) circle (0.055);
  \node[vlab, below left] at (A) {$A$};
  \node[vlab, below right] at (C) {$C$};
  \node[vlab, below right] at (B) {$B$};
  \node[vlab, above right] at (P) {$P$};
\end{tikzpicture}
\figcap{13--5}\end{bkfigure}}

% ---------------------------------------------------------------- Fig 13-6
% Thm 13-2.  The sides of the angle are solid; the bisector is dashed and
% carries P.  A and B are the FEET of the perpendiculars from P, so they are
% built with the projection operator -- never by eye.
% Re-drawn 2026-08-17 against the print (Caleb's photo of p.229).  Four things
% changed: (a) the arrowhead that pointed INTO the vertex is gone -- the book
% has no arrowhead anywhere in this figure; (b) the dashed bisector now stops a
% little past P instead of running to the edge of the picture; (c) PA and PB are
% two straight segments meeting in a visible corner at P, with PA dropping
% VERTICALLY to A (the book lays the lower side of the angle horizontal, which
% is what makes that read); (d) the book's short cross-strokes at the two feet
% and at P are drawn.  Opening at O is 32 degrees, as in the print.
\newcommand{\FIGXIIISIX}{\begin{bkfigure}[\linewidth]
\begin{tikzpicture}[scale=1.300]
  \coordinate (O) at (0,0);
  \coordinate (Uend) at (32:3.56);   % upper side of the angle
  \coordinate (Lend) at (0:3.56);    % lower side, horizontal as in the print
  % |OP| = 2.85/cos 16 puts each foot at 0.80 of its drawn side, as the book does
  \coordinate (P) at (16:2.9648);    % on the bisector of the 32-degree angle
  \coordinate (B) at ($(O)!(P)!(Uend)$);
  \coordinate (A) at ($(O)!(P)!(Lend)$);
  \draw[given] (O) -- (Uend);
  \draw[given] (O) -- (Lend);
  \draw[aux] (O) -- ($(O)!1.25!(P)$);
  \draw[aux] (B) -- (P);
  \draw[aux] (P) -- (A);
  % the print's cross-strokes: one at each perpendicular foot, one at P
  \draw[tick] ($(A)!-0.13!(P)$) -- ($(A)!0.13!(P)$);
  \draw[tick] ($(B)!-0.13!(P)$) -- ($(B)!0.13!(P)$);
  \draw[tick] ($(P)+(106:0.13)$) -- ($(P)+(-74:0.13)$);
  \node[vlab, below left] at (O) {$O$};
  \node[vlab, above left] at (B) {$B$};
  \node[vlab] at ($(A)+(-84:0.32)$) {$A$};
  \node[vlab] at ($(P)+(62:0.34)$) {$P$};
\end{tikzpicture}
\figcap{13--6}\end{bkfigure}}

% ---------------------------------------------------------------- Fig 13-7
% Thm 13-3.  AB is a DIAMETER (O the midpoint); P on the circle; PA, PB dashed
% so that angle APB is the right angle in a semicircle.
\newcommand{\FIGXIIISEVEN}{\begin{bkfigure}[\linewidth]
\begin{tikzpicture}[scale=1.300]
  \def\rr{1.42}
  \coordinate (O) at (0,0);
  \coordinate (A) at (-\rr,0);
  \coordinate (B) at (\rr,0);
  \coordinate (P) at (118:\rr);
  \draw[given] (O) circle (\rr);
  \draw[given] (A) -- (B);
  \draw[aux] (P) -- (A);
  \draw[aux] (P) -- (B);
  \fill (O) circle (0.055);
  \node[vlab, left] at (A) {$A$};
  \node[vlab, right] at (B) {$B$};
  \node[vlab, below] at (O) {$O$};
  \node[vlab, above left] at (P) {$P$};
\end{tikzpicture}
\figcap{13--7}\end{bkfigure}}

% ---------------------------------------------------------------- Fig 13-8
% Ex 13-2 #4: A and B fixed on a dashed line; P below, on one side, with the
% solid segments PA, PB subtending the 50-degree angle of the exercise.
\newcommand{\FIGXIIIEIGHT}{\begin{bkfigure}[\linewidth]
\begin{tikzpicture}[scale=1.118]
  % AB rises at 29 degrees, as measured off the print (PDF 244 at 200 dpi).
  \coordinate (A) at (0.62,0.50);
  \coordinate (B) at (2.3255,1.4454);
  % P is placed on the 50-degree arc over AB, so that angle APB is EXACTLY the
  % 50 degrees the exercise states.  Measured against the print, the drawn P
  % already subtends 50.1 degrees, and it sits at 0.062 of AB along and 0.903
  % of AB below the line -- the position reproduced here.
  \coordinate (P) at (1.5799,-0.9823);
  \draw[aux] ($(A)!-0.68!(B)$) -- ($(B)!-1.00!(A)$);
  \draw[given] (A) -- (P);
  \draw[given] (B) -- (P);
  \fill (B) circle (0.055);
  \node[vlab, above left] at (A) {$A$};
  \node[vlab, above left] at (B) {$B$};
  \node[vlab, below] at (P) {$P$};
\end{tikzpicture}
\figcap{13--8}\end{bkfigure}}

% ---------------------------------------------------------------- Fig 13-9
% Ex 13-2 #6: the wheel in two positions -- the second dashed -- with the
% tracing point P on the rim.  Equal radii; the centres sit at the same
% height, which is what rolling on a (horizontal) line forces.
% Re-drawn 2026-08-17.  The two positions of the wheel are the SAME size and are
% offset SIDEWAYS, so they cross cleanly near the top and near the bottom; the
% old drawing let them all but coincide, which read as one circle nested inside
% another.  The short solid connector that the book draws -- the first piece of
% the traced path, from P's earlier position on the dashed rim to its later
% position on the solid rim, with a cross-tick at the earlier end -- was missing
% entirely, leaving the P label floating with nothing to name.
% Offsets measured off the print: centre separation 0.43 r (= 0.215 diameter),
% P at 241 degrees on the earlier wheel.
\newcommand{\FIGXIIININE}{\begin{bkfigure}[\linewidth]
\begin{tikzpicture}[scale=1.300]
  \def\rr{1.05}
  \def\dd{0.4515}                 % = 0.43\rr, how far the wheel has rolled
  \coordinate (Kp) at (0,0);      % dashed wheel, EARLIER position
  \coordinate (K)  at (\dd,0);    % solid wheel, LATER position; centres level
  % rolling without slipping turns the wheel through d/r = 0.43 rad = 24.637 deg
  \coordinate (P0) at ($(Kp)+(241:\rr)$);
  \coordinate (P)  at ($(K)+(216.363:\rr)$);
  \draw[aux] (Kp) circle (\rr);
  \draw[given] (K) circle (\rr);
  \draw[given] (P0) -- (P);
  \draw[tick] ($(P0)!0.33!90:(P)$) -- ($(P0)!-0.33!90:(P)$);
  \node[vlab] at ($(P)+(38:0.36)$) {$P$};
\end{tikzpicture}
\figcap{13--9}\end{bkfigure}}

% --------------------------------------------------------------- Fig 13-10
% Ex 13-2 #8: a circle of radius R/4 rolling INSIDE a circle of radius R.
% Internal tangency is forced: centre distance = R - R/4.
\newcommand{\FIGXIIITEN}{\begin{bkfigure}[\linewidth]
\begin{tikzpicture}[scale=1.300]
  \def\RR{1.40}
  \coordinate (O) at (0,0);
  \coordinate (Q) at (-32:0.75*\RR);   % = R - R/4 = 0.75 R
  \draw[given] (O) circle (\RR);
  \draw[given] (Q) circle (0.25*\RR);
  \fill (O) circle (0.05);
  \fill (Q) circle (0.05);
\end{tikzpicture}
\figcap{13--10}\end{bkfigure}}

% --------------------------------------------------------------- Fig 13-11
% Thm 13-4, INCENTER.  Each bisector foot is placed by the angle-bisector
% theorem (ratio of the two adjacent sides), so the three dashed cevians are
% genuinely the angle bisectors and genuinely concurrent.
\newcommand{\FIGXIIIELEVEN}{\begin{bkfigure}[\linewidth]
\begin{tikzpicture}[scale=1.300]
  \coordinate (A) at (0,0);
  \coordinate (B) at (3.50,2.70);
  \coordinate (C) at (3.30,0.10);
  % |BC| = 2.6077, |CA| = 3.3015, |AB| = 4.4204
  % foot on BC from A divides BC as AB : AC = 4.4204 : 3.3015
  \coordinate (Da) at ($(B)!0.57243!(C)$);
  % foot on CA from B divides AC as AB : BC = 4.4204 : 2.6077
  \coordinate (Eb) at ($(A)!0.62896!(C)$);
  % foot on AB from C divides AB as AC : BC = 3.3015 : 2.6077
  \coordinate (Fc) at ($(A)!0.55871!(B)$);
  \coordinate (O) at (intersection of A--Da and B--Eb);
  \draw[given] (A) -- (B) -- (C) -- cycle;
  % long dashes (see auxlong above): at the standard pattern these three read
  % as solid and the figure turns into a quadrilateral with its diagonals
  \draw[auxlong] (A) -- (Da);
  \draw[auxlong] (B) -- (Eb);
  \draw[auxlong] (C) -- (Fc);
  % the print marks the concurrency with a bold asterisk, not a bare crossing
  \XIIIstar{O}
  \node[vlab, below left] at (A) {$A$};
  \node[vlab, above] at (B) {$B$};
  \node[vlab, below right] at (C) {$C$};
  \node[vlab] at ($(O)+(167:0.42)$) {$O$};
\end{tikzpicture}
\figcap{13--11}\end{bkfigure}}

% --------------------------------------------------------------- Fig 13-12
% Thm 13-5, CIRCUMCENTER.  NOTE: this figure is the circumcenter, not a second
% incenter -- the dashed lines leave the MIDPOINTS of the sides, not the
% vertices, and O falls outside the triangle because angle B is obtuse.
% O is built as an intersection of two perpendicular bisectors.
\newcommand{\FIGXIIITWELVE}{\begin{bkfigure}[\linewidth]
\begin{tikzpicture}[scale=1.300]
  \coordinate (A) at (0,0);
  \coordinate (B) at (3.20,0);
  \coordinate (C) at (3.75,2.75);
  \coordinate (Mab) at ($(A)!0.5!(B)$);
  \coordinate (Mbc) at ($(B)!0.5!(C)$);
  \coordinate (Mac) at ($(A)!0.5!(C)$);
  % perpendicular bisector of AB: through Mab, perpendicular to AB
  \coordinate (Pab) at ($(Mab)+(0,1)$);
  % perpendicular bisector of BC: through Mbc, perpendicular to BC
  \coordinate (Pbc) at ($(Mbc)!1!90:(C)$);
  \coordinate (O) at (intersection of Mab--Pab and Mbc--Pbc);
  \draw[given] (A) -- (B) -- (C) -- cycle;
  % Re-drawn 2026-08-17.  The book draws each perpendicular bisector SHORT: it
  % starts at the midpoint of its side and stops in a small stub just past O.
  % The old drawing ran all three clean through the triangle and out the far
  % side into empty space, which is what Caleb flagged.
  \draw[auxlong] (Mab) -- ($(Mab)!1.26!(O)$);
  \draw[auxlong] (Mbc) -- ($(Mbc)!1.22!(O)$);
  \draw[auxlong] (Mac) -- ($(Mac)!2.25!(O)$);
  % the cross-strokes at the midpoints -- these are what identify the three
  % lines as PERPENDICULAR bisectors rather than as arbitrary cevians
  \draw[tick] ($(Mab)!0.075!90:(B)$) -- ($(Mab)!-0.075!90:(B)$);
  \draw[tick] ($(Mbc)!0.0856!90:(C)$) -- ($(Mbc)!-0.0856!90:(C)$);
  \XIIIstar{O}
  \node[vlab, below left] at (A) {$A$};
  \node[vlab, below right] at (B) {$B$};
  \node[vlab, above right] at (C) {$C$};
  \node[vlab] at ($(O)+(40:0.40)$) {$O$};
\end{tikzpicture}
\figcap{13--12}\end{bkfigure}}

% --------------------------------------------------------------- Fig 13-13
% Ex 13-4 #1.  CD is PARALLEL to the base: C and D are taken at the SAME ratio
% down the two sides, which forces CD || AB exactly.  P is the midpoint of CD.
\newcommand{\FIGXIIITHIRTEEN}{\begin{bkfigure}[\linewidth]
\begin{tikzpicture}[scale=1.300]
  \coordinate (A) at (0,0);
  \coordinate (Bs) at (3.64,-0.18);
  \coordinate (T) at (0.73,2.09);
  \coordinate (C) at ($(T)!0.4785!(A)$);
  \coordinate (D) at ($(T)!0.4785!(Bs)$);
  \coordinate (P) at ($(C)!0.5!(D)$);
  \draw[given] (A) -- (Bs) -- (T) -- cycle;
  \draw[given] (C) -- (D);
  \fill (P) circle (0.055);
  \node[vlab, left] at (C) {$C$};
  \node[vlab, above right] at (D) {$D$};
  \node[vlab, below] at (P) {$P$};
\end{tikzpicture}
\figcap{13--13}\end{bkfigure}}

% --------------------------------------------------------------- Fig 13-14
% Ex 13-4 #4.  O is the intersection of the perpendicular bisector of AB with
% the perpendicular to BC at B -- so OB _|_ BC exactly, and |OA| = |OB|.
% Built the other way round for stability: put O at the centre, take A and B
% on the circle, then erect BC perpendicular to OB at B.
\newcommand{\FIGXIIIFOURTEEN}{\begin{bkfigure}[\linewidth]
\begin{tikzpicture}[scale=1.300]
  \def\rr{1.50}
  \coordinate (O) at (0,0);
  \coordinate (A) at (203:\rr);
  \coordinate (B) at (-19:\rr);
  \coordinate (Q) at (59:\rr);
  % C on the perpendicular to OB at B, on the far side from the centre
  \coordinate (C) at ($(B)!1.05!90:(O)$);
  \draw[given] (O) circle (\rr);
  \draw[given] (A) -- (B);
  \draw[given] (B) -- (C);
  % perpendicular bisector of AB (through O), and the radius OB
  \coordinate (Mab) at ($(A)!0.5!(B)$);
  % drawn as in the print: a short overshoot above O, then down past the
  % chord AB to the foot of the circle
  \draw[aux] ($(O)!-0.89!(Mab)$) -- ($(O)!2.85!(Mab)$);
  \draw[aux] (B) -- ($(O)!-0.62!(B)$);
  \XIIIplus{O}
  \fill (Q) circle (0.055);
  \node[vlab, left] at (A) {$A$};
  \node[vlab, right] at (B) {$B$};
  \node[vlab] at ($(O)+(216:0.34)$) {$O$};
  \node[vlab, above right] at (Q) {$Q$};
  \node[vlab, below] at (C) {$C$};
\end{tikzpicture}
\figcap{13--14}\end{bkfigure}}

% ------------------------------------------------------- Figs 13-15, 13-16
% Printed side by side on p.233.
%
% 13-15, Thm 13-6 CENTROID: D, E are midpoints of AB, BC; O is the crossing of
% the medians CD and AE; G, F are the midpoints of OA, OC (so that D, E, F, G
% are the parallelogram vertices of the proof outline).
%
% 13-16, Thm 13-7 ORTHOCENTER: through each vertex, the parallel to the
% opposite side.  That makes A, B, C exactly the midpoints of the sides of
% A'B'C', so A' = B+C-A, B' = C+A-B, C' = A+B-C -- computed, not eyeballed.
% O is the orthocentre of ABC (= circumcentre of A'B'C').
\newcommand{\FIGXIIIFIFTEENSIXTEEN}{\begin{bkfigure}[\linewidth]
\hfill
\begin{minipage}{0.43\linewidth}\centering
\begin{tikzpicture}[scale=0.92]
  \coordinate (A) at (0,0);
  \coordinate (B) at (2.00,2.14);
  \coordinate (C) at (4.10,-0.10);
  \coordinate (D) at ($(A)!0.5!(B)$);
  \coordinate (E) at ($(B)!0.5!(C)$);
  \coordinate (O) at (intersection of A--E and C--D);
  \coordinate (G) at ($(A)!0.5!(O)$);
  \coordinate (F) at ($(O)!0.5!(C)$);
  \draw[given] (A) -- (B) -- (C) -- cycle;
  \draw[given] (A) -- (E);
  \draw[given] (C) -- (D);
  \draw[given] (D) -- (E);
  \fill (G) circle (0.055);
  \fill (F) circle (0.055);
  \node[vlab, left] at (A) {$A$};
  \node[vlab, above] at (B) {$B$};
  \node[vlab, right] at (C) {$C$};
  \node[vlab, above left] at (D) {$D$};
  \node[vlab, above right] at (E) {$E$};
  \node[vlab] at (1.38446,0.20798) {$G$};
  \node[vlab] at (2.75246,0.16316) {$F$};
  \node[vlab] at ($(O)+(269:0.32)$) {$O$};
\end{tikzpicture}
\figcap{13--15}
\end{minipage}
\hfill
\begin{minipage}{0.55\linewidth}\centering
\begin{tikzpicture}[scale=0.82]
  \coordinate (A) at (0,0);
  \coordinate (B) at (1.10,0.95);
  \coordinate (C) at (2.86,-0.05);
  \coordinate (Ap) at ($(B)+(C)-(A)$);
  \coordinate (Bp) at ($(C)+(A)-(B)$);
  \coordinate (Cp) at ($(A)+(B)-(C)$);
  % orthocentre: intersection of two altitudes, each a genuine perpendicular
  \coordinate (Fa) at ($(B)!(A)!(C)$);
  \coordinate (Fb) at ($(A)!(B)!(C)$);
  \coordinate (O) at (intersection of A--Fa and B--Fb);
  % Re-drawn 2026-08-17 against the print.  The arrowheads had been put on the
  % WRONG lines: the book arrows the three SOLID sides (two tips converging on
  % C, one on A) and leaves the dashed altitudes AO and CO plain, with a single
  % arrowhead at O belonging to the altitude through B.  Ours arrowed AO and CO
  % into A and C, so three or four heads piled onto each vertex and blotted it
  % -- which is also what made the straight line C'-A-B' look kinked at A.
  \draw[auxlong,bkarrows] ($(Cp)!-0.04!(Ap)$) -- ($(Ap)!-0.04!(Cp)$);
  \draw[auxlong] (Cp) -- (Bp);
  \draw[auxlong] (Ap) -- (Bp);
  \draw[given,bkarrow] (B) -- (A);
  \draw[given,bkarrow] (B) -- (C);
  \draw[given,bkarrow] (A) -- (C);
  \draw[auxlong] (A) -- (O);
  \draw[auxlong] (C) -- (O);
  \draw[auxlong,bkarrow] (Fb) -- (O);
  % the print ticks AC where the altitude through B meets it
  \draw[tick] ($(Fb)!0.055!90:(C)$) -- ($(Fb)!-0.055!90:(C)$);
  \node[vlab, below left] at (A) {$A$};
  \node[vlab] at ($(B)+(44:0.46)$) {$B$};
  \node[vlab, below right] at (C) {$C$};
  \node[vlab, right] at ($(Ap)!-0.04!(Cp)$) {$A\XIIIpr$};
  \node[vlab, below] at (Bp) {$B\XIIIpr$};
  \node[vlab, left] at ($(Cp)!-0.04!(Ap)$) {$C\XIIIpr$};
  \node[vlab, above] at (O) {$O$};
\end{tikzpicture}
\figcap{13--16}
\end{minipage}
\hfill
\end{bkfigure}}

% --------------------------------------------------------------- Fig 13-17
% Ex 13-5 *3: the data for an SSA construction -- the angle at A, with the two
% given lengths tick-marked.  Reproduced as the book draws it: the second side
% of the angle at C runs free (the third vertex is not yet located), and the
% triple-ticked segment is drawn across from that free side to CA.
% Re-drawn 2026-08-17.  Caleb settled this figure: "YOU DID THIS CORRECT BESIDES
% THE 1 TICK NOT BEING PERPENDICULAR TO THE LINE CA."  The single tick had been
% struck at 115 degrees while CA runs at -43.6, i.e. only 25 degrees off lying
% ALONG the line -- which is what smeared it into the triple tick below it.
% Every tick group is now struck perpendicular to the segment it measures, the
% three ticks are shortened and spread so they stay countable, and the whole
% figure is re-proportioned to the print: angle at C = 26 degrees, the free ray
% 1.36|CA| long, |AB| = 1.745|CA|, the polyline C-A-B turning 39 degrees at A,
% the single tick at 0.468 of CA and the transversal foot at 0.70 (it used to be
% at 0.85, which crowded the two tick groups together).
\newcommand{\FIGXIIISEVENTEEN}{\begin{bkfigure}[\linewidth]
\begin{tikzpicture}[scale=1.450]
  \coordinate (C) at (0,1.30);
  \coordinate (A) at ($(C)+(-43.6:1.55)$);
  \coordinate (Cf) at ($(C)+(-17.6:2.108)$);   % free end of the second ray
  \coordinate (B) at ($(A)+(-4.9:2.705)$);
  \draw[given] (C) -- (A);
  \draw[given] (A) -- (B);
  \draw[given] (C) -- (Cf);
  % the triple-ticked transversal, from a point on CA across to the free ray
  \coordinate (Sf) at ($(C)!0.70!(A)$);
  \coordinate (S)  at ($(C)!0.51!(Cf)$);
  \draw[given] (Sf) -- (S);
  % ONE tick on CA, PERPENDICULAR to CA
  \coordinate (T1) at ($(C)!0.468!(A)$);
  \draw[tick] ($(T1)!0.1394!90:(A)$) -- ($(T1)!-0.1394!90:(A)$);
  % THREE ticks on the transversal, perpendicular to it, grouped toward the CA
  % end as the print groups them
  \coordinate (Qt) at ($(Sf)!0.144!90:(S)$);
  \foreach \t in {0.28,0.41,0.54}
    \draw[tick] ($($(Sf)!\t!(S)$)+(Qt)-(Sf)$) -- ($($(Sf)!\t!(S)$)-(Qt)+(Sf)$);
  % TWO ticks on AB, perpendicular to AB
  \coordinate (Qb) at ($(A)!0.0388!90:(B)$);
  \foreach \t in {0.47,0.53}
    \draw[tick] ($($(A)!\t!(B)$)+(Qb)-(A)$) -- ($($(A)!\t!(B)$)-(Qb)+(A)$);
  \node[vlab, above left] at (C) {$C$};
  \node[vlab, below left] at (A) {$A$};
  \node[vlab, below right] at (B) {$B$};
\end{tikzpicture}
\figcap{13--17}\end{bkfigure}}

% figures14.tex -- Chapter 14 (Space Geometry) figures redrawn in TikZ.
% Proportions follow the printed figures (iPad scan15 pp.3-19 = book pp.236-249);
% exact coordinates are reconstructions.
%
% Projection conventions copied from the book:
%   * a plane is a parallelogram, long axis horizontal, receding edge up-right;
%   * hidden edges are dashed (style "hid" below -- part of the GIVEN object,
%     so it stays black in the colour edition; it is not an "aux" line);
%   * solids are drawn in oblique projection with the depth axis up-right.

\tikzset{
  hid/.style  = {fig, dashed, dash pattern=on 2.6pt off 2.2pt},
  arw/.style  = {fig, ->, >=latex},
  % A dihedral/plane-angle arc, as the book strikes it.
  farc/.style = {fig, line width=0.6pt},
}
% Point dots are \dt{} (style/brumfiel.sty): a node of fixed 3.2pt diameter, so
% it prints the same size whatever `scale=' the picture carries.  The book's
% points are heavy round discs -- never arrowheads (Caleb, 2026-08-17).

% \XIVplane{name}{ox}{oy}{width}{dx}{dy} -- a plane as a parallelogram.
% Defines <name>sw, <name>se, <name>ne, <name>nw and strokes the outline.
\newcommand{\XIVplane}[6]{%
  \coordinate (#1sw) at (#2,#3);
  \coordinate (#1se) at (#2+#4,#3);
  \coordinate (#1ne) at (#2+#4+#5,#3+#6);
  \coordinate (#1nw) at (#2+#5,#3+#6);
  \draw[fig] (#1sw) -- (#1se) -- (#1ne) -- (#1nw) -- cycle;}

% \XIVplanec{name}{ox}{oy}{width}{dx}{dy} -- same, corners only, no stroke.
\newcommand{\XIVplanec}[6]{%
  \coordinate (#1sw) at (#2,#3);
  \coordinate (#1se) at (#2+#4,#3);
  \coordinate (#1ne) at (#2+#4+#5,#3+#6);
  \coordinate (#1nw) at (#2+#5,#3+#6);}

% ---------------------------------------------------------------- Fig 14-1
% Five familiar solids: two intersecting planes, cube, sphere, pentagonal
% prism, square pyramid.  Scan: the cube carries no hidden edges; the prism
% and the pyramid do.
\newcommand{\FIGXIVONE}{\begin{bkfigure}[\linewidth]
\begin{tikzpicture}[scale=0.760]
  % (1) two intersecting planes.  The book draws two long parallelogram
  %     "boards" that genuinely CROSS near their lower-right ends, leaving the
  %     two little spikes where each corner pokes out past the other.  (The old
  %     version built both on one common line, which put three near-coincident
  %     long edges down the middle -- the "smear" Caleb flagged.)
  \coordinate (oA) at (1.05,2.90);  \coordinate (dA) at (1.20,-2.60);
  \coordinate (wA) at (0.55,0.18);
  \coordinate (oB) at (0.05,1.75);  \coordinate (dB) at (2.55,-1.55);
  \coordinate (wB) at (0.42,0.62);
  \draw[fig] (oA) -- ($(oA)+(dA)$) -- ($(oA)+(dA)+(wA)$) -- ($(oA)+(wA)$)
             -- cycle;
  \draw[fig] (oB) -- ($(oB)+(dB)$) -- ($(oB)+(dB)+(wB)$) -- ($(oB)+(wB)$)
             -- cycle;
  % (2) cube
  \begin{scope}[xshift=3.55cm, yshift=0.45cm]
    \draw[fig] (0,0) -- (1.45,0) -- (1.45,1.45) -- (0,1.45) -- cycle;
    \draw[fig] (0,1.45) -- (0.62,1.95) -- (2.07,1.95) -- (1.45,1.45);
    \draw[fig] (1.45,0) -- (2.07,0.5) -- (2.07,1.95);
  \end{scope}
  % (3) sphere.  The book shades it with a DENSE fine dot screen filling a
  %     broad annulus just inside the rim, heaviest at the upper left and
  %     thinning round the bottom, leaving a clean highlight low and right.
  %     `pattern=dots' is far too sparse for that, so the screen is drawn.
  \begin{scope}[xshift=7.05cm, yshift=1.5cm]
    \begin{scope}[even odd rule]
      \clip (0,0) circle (0.86);
      \clip (0.86,-0.86) rectangle (-0.86,0.86) (0.26,-0.22) circle (0.72);
      \foreach \i in {0,...,31}{\foreach \j in {0,...,31}{
        \fill ({-0.87+0.056*\i+0.028*mod(\j,2)},{-0.87+0.056*\j})
              circle (0.30pt);}}
    \end{scope}
    \draw[fig] (0,0) circle (0.86);
  \end{scope}
  % (4) rectangular box, standing on end and tilted, drawn open at the top --
  %     the book's fourth solid.  e1 = width, e2 = depth, e3 = length.
  \begin{scope}[xshift=8.95cm, yshift=2.40cm]
    \coordinate (e1) at (1.02,0.09);
    \coordinate (e2) at (0.60,-0.83);
    \coordinate (e3) at (-1.02,-2.33);
    \coordinate (TL) at (0,0);
    \coordinate (TT) at ($(TL)+(e1)$);
    \coordinate (TR) at ($(TL)+(e1)+(e2)$);
    \coordinate (TB) at ($(TL)+(e2)$);
    \coordinate (BL) at ($(TL)+(e3)$);
    \coordinate (BT) at ($(TT)+(e3)$);
    \coordinate (BR) at ($(TR)+(e3)$);
    \coordinate (BB) at ($(TB)+(e3)$);
    \draw[fig] (TL) -- (TT) -- (TR) -- (TB) -- cycle;   % top face (rim)
    \draw[fig] (TL) -- (BL) -- (BB) -- (TB);            % near-left face
    \draw[fig] (BB) -- (BR) -- (TR);                    % near-right face
    \draw[hid] (BT) -- (TT);                            % hidden far edge
    \draw[hid] (BT) -- (BL);  \draw[hid] (BT) -- (BR);  % hidden bottom edges
    % the inner lip: the box reads as open at the top
    \coordinate (c) at ($(TL)!0.5!(TR)$);
    \draw[fig] ($(c)!0.80!(TL)$) -- ($(c)!0.80!(TB)$) -- ($(c)!0.80!(TR)$);
  \end{scope}
  % (5) square pyramid.  The book draws the apex-to-rear-vertex edge SOLID and
  %     fans the two hidden base edges from its foot as straight dashes.
  \begin{scope}[xshift=11.15cm, yshift=0.40cm]
    \coordinate (b1) at (0,0.34); \coordinate (b2) at (1.0,0);
    \coordinate (b3) at (2.0,0.34); \coordinate (b4) at (1.0,0.68);
    \coordinate (ap) at (1.0,1.92);
    \draw[fig] (b1) -- (b2) -- (b3);
    \draw[fig] (ap) -- (b1); \draw[fig] (ap) -- (b2); \draw[fig] (ap) -- (b3);
    \draw[fig] (ap) -- (b4);
    \draw[hid] (b3) -- (b4) -- (b1);
  \end{scope}
\end{tikzpicture}
\figcap{14--1}
\end{bkfigure}}

% ---------------------------------------------------------------- Fig 14-2
% Pictorial: a swept-wing jet above a sedan on a hatched ground line.  This is
% an illustration, not a geometric construction -- see ch14-UNCERTAIN.md.
\newcommand{\FIGXIVTWO}{\begin{bkfigure}[0.80\linewidth]
\begin{tikzpicture}[scale=0.806]
  % ---- jet, three-quarter view, nose to the right
  \begin{scope}[yshift=1.55cm]
    % fuselage
    \draw[fig] (0.30,0.42) -- (0.95,0.30) -- (4.55,0.28) -- (4.95,0.40)
               -- (4.95,0.56) -- (4.52,0.66) -- (1.05,0.70) -- (0.42,0.62)
               -- cycle;
    % canopy
    \draw[fig] (3.05,0.70) -- (3.35,0.94) -- (4.00,0.94) -- (4.22,0.70);
    \draw[fig] (3.42,0.70) -- (3.60,0.94);
    % swept main wing (far) and near wing
    \draw[fig] (1.85,0.70) -- (1.35,1.30) -- (2.55,1.24) -- (2.90,0.70);
    \draw[fig] (2.05,0.28) -- (1.50,-0.34) -- (2.75,-0.28) -- (3.05,0.28);
    % tailplane + fin
    \draw[fig] (0.62,0.70) -- (0.28,1.10) -- (1.05,1.06) -- (1.28,0.70);
    \draw[fig] (0.55,0.70) -- (0.45,1.42) -- (1.02,1.40) -- (1.30,0.70);
    % nose probe
    \draw[fig] (4.95,0.48) -- (5.35,0.48);
  \end{scope}
  % ---- sedan
  \begin{scope}[xshift=4.05cm]
    \draw[fig] (0.30,0.30) -- (0.36,0.62) -- (0.95,0.66) -- (1.55,1.02)
               -- (2.95,1.04) -- (3.30,0.68) -- (4.15,0.62) -- (4.22,0.30);
    \draw[fig] (0.30,0.30) -- (4.22,0.30);
    \draw[fig] (1.62,0.94) -- (2.22,0.96) -- (2.22,0.66);
    \draw[fig] (2.36,0.96) -- (2.90,0.96) -- (3.16,0.68);
    \draw[fig] (0.95,0.66) -- (3.30,0.68);
    \draw[fig] (1.28,0.30) circle (0.24);
    \draw[fig] (3.42,0.30) circle (0.24);
  \end{scope}
  % ---- ground line with hatching
  \draw[fig] (0.55,0.06) -- (9.20,0.06);
  \foreach \x in {0.75,1.05,...,9.05}
    {\draw[fig, line width=0.4pt] (\x,0.06) -- (\x-0.22,-0.16);}
\end{tikzpicture}
\figcap{14--2}
\end{bkfigure}}

% ---------------------------------------------------------------- Fig 14-3
\newcommand{\FIGXIVTHREE}{\begin{bkfigure}[0.52\linewidth]
\begin{tikzpicture}[scale=0.910]
  \XIVplane{P}{0}{0}{3.5}{0.95}{1.30}
  \dt{1.30,0.52}
  \dt{2.55,0.86}
  \dt{1.95,0.30}
\end{tikzpicture}
\figcap{14--3}
\end{bkfigure}}

% ------------------------------------------------------------ Fig 14-4, 14-5
\newcommand{\FIGXIVFOURFIVE}{\begin{bkfigure}[\linewidth]
\begin{minipage}{0.47\linewidth}\centering
\begin{tikzpicture}[scale=0.793]
  \XIVplane{P}{0}{0}{3.5}{0.95}{1.30}
  \dt{1.15,0.46}
  \dt{2.30,0.78}
  \dt{1.85,0.26}
  \dt{2.20,1.52}   % the fourth point, off the plane
\end{tikzpicture}
\figcap{14--4}
\end{minipage}\hfill
\begin{minipage}{0.47\linewidth}\centering
\begin{tikzpicture}[scale=0.793]
  \XIVplane{P}{0}{0}{3.5}{0.95}{1.30}
  \coordinate (Pp) at (2.05,0.55);
  \coordinate (Lu) at (3.10,2.00);           % upper end of l, above pi
  \coordinate (Lo) at ($(Pp)!-1.25!(Lu)$);   % lower end -- on the SAME line
  % The book hides the stretch of l that lies BEHIND the plane: solid until
  % l crosses the near edge, dashed from there to P, solid above P.
  \coordinate (Le) at (intersection of Lo--Lu and Psw--Pse);
  \draw[fig] (Lo) -- (Le);
  \draw[hid] (Le) -- (Pp);
  \draw[fig] (Pp) -- (Lu);
  \dt{Pp}
  \node[vlab, left=3pt] at (Pp) {$P$};
  \node[vlab, right=3pt] at (Lu) {$l$};
\end{tikzpicture}
\figcap{14--5}
\end{minipage}
\end{bkfigure}}

% ------------------------------------------------------------ Fig 14-6, 14-7
\newcommand{\FIGXIVSIXSEVEN}{\begin{bkfigure}[\linewidth]
\begin{minipage}{0.40\linewidth}\centering
\begin{tikzpicture}[scale=1.144]
  \XIVplane{P}{0}{0}{2.9}{0.85}{1.20}
  % PQ is lifted clear of the bottom edge so that "l" fits underneath it
  \coordinate (Pp) at (0.72,0.80); \coordinate (Qq) at (2.05,0.60);
  \draw[fig] (Pp) -- (Qq);
  \dt{Pp} \dt{Qq}
  \dt{2.26,0.96}
  \node[vlab, above left=3pt] at (Pp) {$P$};
  \node[vlab, right=3.5pt] at (Qq) {$Q$};
  \node[vlab, right=3.5pt] at (2.26,0.96) {$R$};
  \node[vlab, below=3pt] at ($(Pp)!0.5!(Qq)$) {$l$};
\end{tikzpicture}
\figcap{14--6}
\end{minipage}\hfill
\begin{minipage}{0.56\linewidth}\centering
\begin{tikzpicture}[scale=1.10]
  % Theorem 14-3: infinitely many planes contain one line.  The book draws the
  % axis l nearly end-on, so the four planes are parallelograms built on the
  % SAME short segment AB of l and fanned about it like a deck of cards; each
  % has real area, and its far corner is a heavy dot on the dashed transversal
  % m.  (The old drawing made the planes slivers and left l undrawn.)
  \coordinate (Av) at (-1.450,1.167);        % upper common vertex, on l
  \coordinate (Bv) at (0,0);                 % lower common vertex, on l
  \coordinate (d1) at (-2.917,1.450);        % first dot on m
  \coordinate (dm) at (2.867,0.650);         % direction of m
  \foreach \t in {0, 0.53, 0.695, 1.0}{
    \coordinate (D) at ($(d1)+\t*(dm)$);
    \coordinate (W) at ($(D)-(Av)$);
    \draw[fig] (Av) -- (D) -- ($(Bv)+(W)$) -- (Bv);}
  \draw[fig] (Av) -- (Bv) -- (0.233,-0.283);       % the axis l, overshooting B
  \draw[hid] ($(d1)-0.16*(dm)$) -- ($(d1)+1.30*(dm)$);   % m, dashed throughout
  \foreach \t in {0, 0.53, 0.695, 1.0}{\dt{$(d1)+\t*(dm)$}}
  \node[vlab, right=2pt] at ($(d1)+1.30*(dm)$) {$m$};
  \node[vlab, below right=1pt] at (0.233,-0.283) {$l$};
\end{tikzpicture}
\figcap{14--7}
\end{minipage}
\end{bkfigure}}

% ------------------------------------------------------------ Fig 14-8, 14-9
\newcommand{\FIGXIVEIGHTNINE}{\begin{bkfigure}[\linewidth]
\begin{minipage}{0.47\linewidth}\centering
\begin{tikzpicture}[scale=1.259]
  \XIVplane{P}{0}{0}{3.0}{0.75}{1.25}
  % l_1 runs Q -> E1, l_2 runs P -> E2; R is their true crossing.  The book
  % ends each segment on a HEAVY ROUND DOT at P and at Q -- no arrowheads --
  % and marks R with the largest dot of the three.
  \coordinate (Pp) at (1.10,1.00); \coordinate (Qq) at (0.95,0.45);
  \coordinate (E1) at (2.62,1.06); \coordinate (E2) at (2.68,0.36);
  \draw[fig] (Pp) -- (E2);
  \draw[fig] (Qq) -- (E1);
  \coordinate (R) at (intersection of Pp--E2 and Qq--E1);
  \dt{Pp}\dt{Qq}\dt{R}
  \node[vlab, left=4pt] at (Pp) {$P$};
  \node[vlab, left=4pt] at (Qq) {$Q$};
  \node[vlab, below=3pt] at (R) {$R$};
  \node[vlab, right=2.5pt] at (E1) {$l_1$};
  \node[vlab, right=2.5pt] at (E2) {$l_2$};
\end{tikzpicture}
\figcap{14--8}
\end{minipage}\hfill
\begin{minipage}{0.47\linewidth}\centering
\begin{tikzpicture}[scale=1.44]
  % Theorem 14-5.  The book draws an OPEN-BOOK hinge: the two planes share one
  % COMPLETE edge and swing away from each other, so nothing crosses anything.
  % (The old drawing had them straddle the line, so each plane ran out past the
  % other on both sides and the middle filled with a lens-shaped tangle.)
  \coordinate (I1) at (0,0);            % bottom-left end of the common edge
  \coordinate (I2) at (1.567,1.217);    % upper-right end of it
  \coordinate (ua) at (-0.333,1.017);   % pi_1 swings up-left
  \coordinate (ub) at (1.183,0.333);    % pi_2 swings out to the right
  \draw[fig] (I1) -- (I2) -- ($(I2)+(ua)$) -- ($(I1)+(ua)$) -- cycle;
  \draw[fig] (I1) -- (I2) -- ($(I2)+(ub)$) -- ($(I1)+(ub)$) -- cycle;
  \coordinate (Pp) at ($(I1)!0.31!(I2)$);
  \coordinate (Qq) at ($(I1)!0.81!(I2)$);
  \dt{Pp} \dt{Qq}
  \node[vlab] at ($(Pp)+(-0.28,0.24)$) {$P$};
  \node[vlab] at ($(Qq)+(-0.28,0.24)$) {$Q$};
  \node[vlab] at (0.10,0.95) {$\pi_1$};
  \node[vlab] at (1.25,0.68) {$\pi_2$};
\end{tikzpicture}
\figcap{14--9}
\end{minipage}
\end{bkfigure}}

% --------------------------------------------------------------- Fig 14-10
% Dihedral angle: the two half planes are solid, the rest of each plane dashed.
\newcommand{\FIGXIVTEN}{\begin{bkfigure}[0.60\linewidth]
\begin{tikzpicture}[scale=1.620]
  \coordinate (E1) at (0.55,0.35);  \coordinate (E2) at (2.95,1.35);  % edge l
  \coordinate (u) at (-0.55,0.85);   % offset of half plane 1 (up-left)
  \coordinate (v) at (0.75,-0.55);   % offset of half plane 2 (down-right)
  % half plane pi_1^+ (solid)
  \draw[fig] (E1) -- ($(E1)+(u)$) -- ($(E2)+(u)$) -- (E2);
  % half plane pi_2^+ (solid)
  \draw[fig] (E1) -- ($(E1)+(v)$) -- ($(E2)+(v)$) -- (E2);
  % the completing half planes, dashed
  \draw[hid] (E1) -- ($(E1)-(u)$) -- ($(E2)-(u)$) -- (E2);
  \draw[hid] (E1) -- ($(E1)-(v)$) -- ($(E2)-(v)$) -- (E2);
  \draw[fig] (E1) -- (E2);
  \draw[fig] (E2) -- (3.45,1.55);
  \draw[fig] (E1) -- (0.15,0.18);            % l runs past both vertices
  \node[vlab, above] at (3.45,1.55) {$l$};
  \node[vlab] at (1.25,1.13) {$\pi_1^+$};
  \node[vlab] at (1.90,0.48) {$\pi_2^+$};
\end{tikzpicture}
\figcap{14--10}
\end{bkfigure}}

% ---------------------------------------------------------------- Fig 14-11
% Split out of the old \FIGXIVELEVENTWELVE block: Ex.~14-3 \#1 cites 14-11 and
% \#2 cites 14-12, so each now stands on its own inside the exercise list.
\newcommand{\FIGXIVELEVEN}{\begin{bkfigure}[0.62\linewidth]
\begin{tikzpicture}[scale=1.55]
  % tetrahedron seen nearly edge-on: L apex at left, vertical edge T-B at right
  \coordinate (L) at (0,0.62); \coordinate (T) at (2.32,1.42);
  \coordinate (Bv) at (2.32,0.02); \coordinate (R) at (2.86,0.72);
  \draw[fig] (L) -- (T) -- (Bv) -- cycle;
  \draw[fig] (T) -- (R) -- (Bv);
  \draw[hid] (L) -- (R);
\end{tikzpicture}
\figcap{14--11}
\end{bkfigure}}

% ---------------------------------------------------------------- Fig 14-12
\newcommand{\FIGXIVTWELVE}{\begin{bkfigure}[0.46\linewidth]
\begin{tikzpicture}[scale=1.35]
  % irregular polyhedron (surface), no hidden edges in the book's drawing
  \draw[fig] (0.62,0.10) -- (0.06,1.02) -- (0.30,2.02) -- (1.28,2.52)
             -- (2.10,2.18) -- (2.30,1.08) -- (1.72,0.16) -- cycle;
  \draw[fig] (0.30,2.02) -- (1.42,1.42) -- (2.30,1.08);
  \draw[fig] (1.42,1.42) -- (1.28,2.52);
  \draw[fig] (1.42,1.42) -- (1.05,0.62) -- (1.72,0.16);
  \draw[fig] (1.05,0.62) -- (0.62,0.10);
\end{tikzpicture}
\figcap{14--12}
\end{bkfigure}}

% ----------------------------------------------------------- Fig 14-13, 14-14
\newcommand{\FIGXIVTHIRTEENFOURTEEN}{\begin{bkfigure}[\linewidth]
\begin{minipage}{0.52\linewidth}\centering
\begin{tikzpicture}[scale=1.05]
  % pi drawn as a rectangle (seen face on); angle AOB is a right angle
  \draw[fig] (0,0) rectangle (2.55,1.75);
  \coordinate (O) at (1.05,0.52);
  \coordinate (Aa) at (2.02,0.86);          % ray OA, direction (0.97,0.34)
  \coordinate (Bb) at (0.71,1.49);          % ray OB = OA turned exactly 90 deg
  \draw[arw] (O) -- (Aa); \draw[arw] (O) -- (Bb);
  \draw[hid] (O) -- ($(O)!-0.85!(Aa)$);      % the rest of the line of OA
  \draw[fig, line width=0.5pt]
        (1.3331,0.6192) -- (1.2339,0.9024) -- (0.9508,0.8031);
  \dt{O}
  \node[vlab] at (1.22,0.24) {$O$};
  \node[vlab] at (2.24,0.92) {$A$};
  \node[vlab] at (0.50,1.49) {$B$};
  \node[vlab] at (2.30,0.34) {$\pi$};
  % pi' drawn as a diamond
  \begin{scope}[xshift=2.62cm, yshift=-0.15cm]
    \draw[fig] (1.375,0) -- (2.75,1.375) -- (1.375,2.75) -- (0,1.375) -- cycle;
    \coordinate (Op) at (1.375,1.10);
    \coordinate (rp) at (0.50,1.10);        % r'
    \coordinate (rpp) at (1.375,2.21);      % r''  (perpendicular to r')
    % the book tips OA and OB with arrowheads but leaves r' and r'' plain
    \draw[fig] (Op) -- (rp); \draw[fig] (Op) -- (rpp);
    \draw[hid] (Op) -- ($(Op)!-1!(rp)$);     % the rest of the line of r'
    \draw[fig, line width=0.5pt] (1.03,1.10) -- (1.03,1.445) -- (1.375,1.445);
    \node[vlab] at (0.95,0.85) {$r'$};
    \node[vlab] at (1.68,1.72) {$r''$};
    \node[vlab] at (1.72,0.80) {$O'$};
    \node[vlab] at (1.375,0.42) {$\pi'$};
  \end{scope}
\end{tikzpicture}
\figcap{14--13}
\end{minipage}\hfill
\begin{minipage}{0.44\linewidth}\centering
\begin{tikzpicture}[scale=1.1]
  % Definition 14-3: l meets pi at P and is perpendicular to EVERY line of pi
  % through P, so the book draws two such lines, both passing through P.
  \XIVplane{P}{0}{0}{3.0}{1.15}{1.25}
  \coordinate (Pp) at (1.72,0.55);
  \coordinate (md) at (0.98,0.40);            % direction of m
  \coordinate (nd) at (-0.92,0.36);           % a second line of pi through P
  \draw[fig] (Pp) -- (1.72,2.05);
  \draw[fig] ($(Pp)-0.65*(md)$) -- ($(Pp)+1.15*(md)$);
  \draw[fig] ($(Pp)-0.85*(nd)$) -- ($(Pp)+0.80*(nd)$);
  % right-angle mark between l and m, sheared into the oblique projection
  \draw[fig, line width=0.5pt] ($(Pp)+(0,0.30)$)
        -- ($(Pp)+(0,0.30)+0.30*(md)$) -- ($(Pp)+0.30*(md)$);
  \dt{Pp}
  \node[vlab, above] at (1.72,2.05) {$l$};
  \node[vlab, below=1pt] at (Pp) {$P$};
  \node[vlab, right=2pt] at ($(Pp)+1.15*(md)$) {$m$};
  \node[vlab] at (0.72,0.26) {$\pi$};
\end{tikzpicture}
\figcap{14--14}
\end{minipage}
\end{bkfigure}}

% ----------------------------------------------------------- Fig 14-15, 14-16
\newcommand{\FIGXIVFIFTEENSIXTEEN}{\begin{bkfigure}[\linewidth]
\begin{minipage}{0.48\linewidth}\centering
\begin{tikzpicture}[scale=1.13]
  \XIVplane{P}{-0.35}{0}{3.35}{0.85}{1.20}
  % Theorem 14-6 / its hint: l_1 = PB and l_2 = PC are the two intersecting
  % lines of pi; B and C are chosen on them, Q is a point of the auxiliary
  % segment BC, and m is the line PQ.  Every one of l_1, l_2, m therefore runs
  % through P, and Q genuinely lies on BC.
  \coordinate (Pp) at (1.15,0.72);
  \coordinate (Aa) at (1.15,2.35);
  \coordinate (Ap) at (1.15,-0.91);          % A' -- P is the midpoint of AA'
  \coordinate (Bb) at (0.52,0.46);
  \coordinate (Cc) at (2.62,0.76);
  \coordinate (Qq) at ($(Bb)!0.52!(Cc)$);    % Q lies on BC
  \draw[fig] (1.15,2.72) -- (Aa);
  \draw[fig] (Aa) -- (Pp) -- (Ap);
  \draw[arw] (Pp) -- (Bb);                   % l_1
  \draw[arw] (Pp) -- (Cc);                   % l_2
  \draw[fig] (Pp) -- ($(Pp)!3.14!(Qq)$);     % m = PQ
  \draw[hid] (Bb) -- (Cc);
  \foreach \z in {Bb,Qq,Cc}{\draw[hid] (Aa) -- (\z); \draw[hid] (Ap) -- (\z);}
  \foreach \z in {Pp,Qq,Aa,Ap}{\dt{\z}}
  \node[vlab, above] at (1.15,2.72) {$l$};
  \node[vlab] at (1.42,2.42) {$A$};
  \node[vlab] at (1.46,-0.95) {$A'$};
  \node[vlab] at (0.90,0.89) {$P$};
  \node[vlab] at (1.74,0.39) {$Q$};
  \node[vlab] at (0.34,0.68) {$B$};
  \node[vlab] at (2.80,0.68) {$C$};
  \node[vlab] at (0.36,0.22) {$l_1$};
  \node[vlab] at (3.06,0.86) {$l_2$};
  \node[vlab] at (2.88,0.22) {$m$};
\end{tikzpicture}
\figcap{14--15}
\end{minipage}\hfill
\begin{minipage}{0.48\linewidth}\centering
\begin{tikzpicture}[scale=0.81]
  % (a) the plane perpendicular to l at P.  The book's plane is a parallelogram
  % with NEAR-VERTICAL sides; l runs solid to the heavy dot at P, is hidden
  % (dashed) for exactly the stretch that lies behind the plane -- from P to the
  % point where it leaves the parallelogram -- and solid again beyond.  The old
  % drawing put a wedge at P that read as an arrowhead and stranded the label.
  \coordinate (T1) at (0.30,2.05);  \coordinate (T2) at (1.35,1.50);
  \coordinate (T3) at (1.35,-0.10); \coordinate (T4) at (0.30,0.45);
  \draw[fig] (T1) -- (T2) -- (T3) -- (T4) -- cycle;
  \coordinate (L0) at (-0.55,1.32); \coordinate (L1) at (2.30,1.12);
  \coordinate (Pa) at ($(L0)!0.40!(L1)$);
  \coordinate (Xa) at (intersection of L0--L1 and T2--T3);
  \draw[fig] (L0) -- (Pa);
  \draw[hid] (Pa) -- (Xa);
  \draw[fig] (Xa) -- (L1);
  \dt{Pa}
  \node[vlab, left=2pt] at (L0) {$l$};
  \node[vlab, below=3pt] at (Pa) {$P$};
  \node[slab] at (0.90,-0.62) {(a)};
  % (b) l1, l2, l3 and l' all perpendicular to l at P.  The right-angle mark
  % belongs in the wedge between l and l' (it used to sit between l and l_3),
  % and every ray label sits just past the free end of its own ray.
  \begin{scope}[xshift=2.75cm]
    \coordinate (Pb) at (0.95,0.95);
    \coordinate (Lb) at (0.98,2.10);      % l
    \coordinate (Lp) at (1.88,1.88);      % l'
    \coordinate (L3) at (2.15,0.87);      % l_3
    \coordinate (L2) at (1.67,0.20);      % l_2
    \coordinate (L1b) at (0.10,0.26);     % l_1
    \draw[fig] (Pb) -- (Lb);
    \draw[fig] (Pb) -- (L3);
    \draw[fig] (Pb) -- (L1b);
    \draw[fig] (Pb) -- (L2);
    \draw[aux] (Pb) -- (Lp);
    \draw[farc] (0.9568,1.2099) -- (1.1406,1.3937) -- (1.1338,1.1338);
    \node[vlab, above] at (Lb) {$l$};
    \node[vlab, right=2pt] at (L3) {$l_3$};
    \node[vlab, below left=1pt] at (L1b) {$l_1$};
    \node[vlab, below right=1pt] at (L2) {$l_2$};
    \node[vlab, above right=0pt] at (Lp) {$l'$};
    \node[vlab] at (0.62,1.12) {$P$};
    \node[slab] at (1.05,-0.62) {(b)};
  \end{scope}
\end{tikzpicture}
\figcap{14--16}
\end{minipage}
\end{bkfigure}}

% ----------------------------------------------------------- Fig 14-17, 14-18
\newcommand{\FIGXIVSEVENTEENEIGHTEEN}{\begin{bkfigure}[\linewidth]
\begin{minipage}{0.47\linewidth}\centering
\begin{tikzpicture}[scale=1.55]
  % Theorem 14-8.  The book's plane stands close to upright (near-vertical left
  % and right edges, top and bottom edges sloping DOWN to the right) and the
  % hidden run of l goes from the piercing dot ALL THE WAY to the plane's right
  % edge -- ours used to go solid again while still inside the plane.
  \coordinate (T1) at (0.05,1.85);  \coordinate (T2) at (1.80,1.38);
  \coordinate (T3) at (1.74,-0.43); \coordinate (T4) at (-0.01,0.04);
  \draw[fig] (T1) -- (T2) -- (T3) -- (T4) -- cycle;
  \coordinate (L0) at (-0.64,0.98); \coordinate (L1) at (2.38,0.52);
  \coordinate (Pp) at ($(L0)!0.452!(L1)$);
  \coordinate (Xp) at (intersection of L0--L1 and T2--T3);
  \draw[fig] (L0) -- (Pp);
  \draw[hid] (Pp) -- (Xp);
  \draw[fig] (Xp) -- (L1);
  \dt{Pp}
  \dt{0.75,1.405}
  \node[vlab] at (0.95,1.25) {$Q$};
  \node[vlab, left=2pt] at (L0) {$l$};
\end{tikzpicture}
\figcap{14--17}
\end{minipage}\hfill
\begin{minipage}{0.47\linewidth}\centering
\begin{tikzpicture}[scale=1.180]
  \XIVplane{P}{0}{0}{2.6}{0.75}{1.05}
  \coordinate (F1) at (1.05,0.42); \coordinate (F2) at (1.95,0.66);
  \draw[fig] (F1) -- (1.05,1.72);
  \draw[fig] (F2) -- (1.95,1.96);
  \dt{F1} \dt{F2}
\end{tikzpicture}
\figcap{14--18}
\end{minipage}
\end{bkfigure}}

% --------------------------------------------------------------- Fig 14-19
% Left plane: two lines that are parallel, and two that intersect.
% Right pair: parallel planes, each carrying a line; the lines are not parallel.
\newcommand{\FIGXIVNINETEEN}{\begin{bkfigure}[0.90\linewidth]
\begin{tikzpicture}[scale=0.806]
  \XIVplane{A}{0}{0}{3.4}{0.95}{1.30}
  \draw[fig] (1.05,0.92) -- (2.55,1.10);
  \draw[fig] (1.15,0.62) -- (2.65,0.80);      % parallel to the line above
  \draw[fig] (1.72,0.10) -- (2.62,0.52);
  \draw[fig] (2.55,0.12) -- (1.85,0.52);
  % two parallel planes at the right
  \XIVplane{B}{4.55}{1.55}{3.4}{0.95}{1.30}
  \draw[fig] (6.30,2.05) -- (7.15,2.72);
  \XIVplane{C}{4.55}{-0.55}{3.4}{0.95}{1.30}
  \draw[fig] (5.65,0.12) -- (7.35,0.42);
\end{tikzpicture}
\figcap{14--19}
\end{bkfigure}}

% ---------------------------------------------------- Fig 14-20, 14-21, 14-22
% l1' and l2' (Fig 14-20) are drawn parallel to l1 and l2 (Fig 14-21).
\newcommand{\FIGXIVTWENTYTWENTYTWO}{\begin{bkfigure}[\linewidth]
\begin{minipage}{0.30\linewidth}\centering
\begin{tikzpicture}[scale=1.05]
  \coordinate (Pp) at (1.35,1.05);
  \draw[aux] ($(Pp)-(1.30,0.52)$) -- ($(Pp)+(1.30,0.52)$);   % l1' -- slope 0.40
  \draw[aux] ($(Pp)-(1.20,-0.72)$) -- ($(Pp)+(1.20,-0.72)$); % l2' -- slope -0.60
  \dt{Pp}
  \node[vlab] at (1.30,1.42) {$P$};
  \node[vlab] at (2.86,1.86) {$l_1'$};
  \node[vlab, right=4.5pt] at ($(Pp)+(1.20,-0.72)$) {$l_2'$};
\end{tikzpicture}
\figcap{14--20}
\end{minipage}\hfill
\begin{minipage}{0.34\linewidth}\centering
\begin{tikzpicture}[scale=0.88]
  \XIVplane{P}{0}{0}{3.2}{0.95}{1.45}
  \coordinate (Xx) at (1.85,0.75);
  % In the book both segments are SHORT and lie entirely inside the plane --
  % ours ran corner to corner and poked out through two edges.
  \draw[fig] ($(Xx)-(0.62,0.248)$) -- ($(Xx)+(0.62,0.248)$);   % l1, slope 0.40
  \draw[fig] ($(Xx)-(0.58,-0.348)$) -- ($(Xx)+(0.58,-0.348)$); % l2, slope -0.60
  \node[vlab] at (2.68,1.10) {$l_1$};
  \node[vlab] at (2.66,0.30) {$l_2$};
\end{tikzpicture}
\figcap{14--21}
\end{minipage}\hfill
\begin{minipage}{0.30\linewidth}\centering
\begin{tikzpicture}[scale=0.86]
  % Two parallel planes with their common perpendicular.  Each plane stands on
  % edge with VERTICAL front and back edges and top/bottom edges sloping down
  % to the right.  Dashing (the defect Caleb flagged, "fully inverted"): the
  % line is SOLID everywhere -- including the whole stretch BETWEEN the planes
  % -- and dashed only for the short run inside each plane, from its piercing
  % point back to that plane's far edge.  Each plane's near vertical edge is
  % broken with a small gap where the line passes in front of it; the far edge
  % is drawn straight through, since it stands in front of the hidden run.
  \coordinate (aL) at (0.10,2.60);  \coordinate (bL) at (0.93,2.125);
  \coordinate (cL) at (0.93,0.175); \coordinate (dL) at (0.10,0.65);
  \coordinate (aR) at (2.19,2.62);  \coordinate (bR) at (3.02,2.145);
  \coordinate (cR) at (3.02,0.195); \coordinate (dR) at (2.19,0.67);
  \draw[fig] (aL) -- (bL) -- (cL) -- (dL);
  \draw[fig] (aL) -- (0.10,1.205);  \draw[fig] (0.10,1.025) -- (dL);
  \draw[fig] (aR) -- (bR) -- (cR) -- (dR);
  \draw[fig] (aR) -- (2.19,1.268);  \draw[fig] (2.19,1.088) -- (dR);
  \draw[fig] (-0.395,1.1000) -- (0.465,1.1260);
  \draw[hid] (0.465,1.1260) -- (0.930,1.1400);
  \draw[fig] (0.930,1.1400) -- (2.576,1.1897);
  \draw[hid] (2.576,1.1897) -- (3.020,1.2031);
  \draw[fig] (3.020,1.2031) -- (3.330,1.2125);
\end{tikzpicture}
\figcap{14--22}
\end{minipage}
\end{bkfigure}}

% --------------------------------------------------------------- Fig 14-23
% Three parallel planes cut by l and l'; AB:BC = A'B':B'C' by construction.
\newcommand{\FIGXIVTWENTYTHREE}{\begin{bkfigure}[0.64\linewidth]
\begin{tikzpicture}[scale=1.60]
  \XIVplane{T}{0.55}{2.30}{2.85}{0.72}{0.62}
  \XIVplane{M}{0.28}{1.15}{2.85}{0.72}{0.62}
  \XIVplane{B}{0.00}{0.00}{2.85}{0.72}{0.62}
  % Two transversals.  In the book the LEFT-hand label at the top is l', but
  % its line runs down to the RIGHT through A', B', C'; the right-hand label is
  % l and its line runs down to the LEFT through A, B, C -- the two cross just
  % above the top plane.  Both are parameterised by height h (h = 1 in the top
  % plane, 0.5 in the middle, 0 in the bottom), so B and B' fall in the middle
  % plane and AB:BC = A'B':B'C' exactly, which is Theorem 14-12's claim.
  \coordinate (Aa) at (1.72,2.62);   \coordinate (Cc) at (0.60,0.32);
  \coordinate (Bb) at ($(Aa)!0.5!(Cc)$);
  \coordinate (Ap) at (2.05,2.62);   \coordinate (Cp) at (2.80,0.32);
  \coordinate (Bp) at ($(Ap)!0.5!(Cp)$);
  % Hidden-line runs, computed rather than guessed: the line is dashed exactly
  % where it passes BELOW a plane and its projection still falls inside that
  % plane's parallelogram.  That is a short break under each of A, B and C, and
  % solid everywhere else -- the reverse of the old drawing, which was dashed
  % almost end to end and so carried no depth information at all.
  \draw[fig] (2.101,3.402) -- (Aa);
  \draw[hid] (Aa) -- (1.564,2.300);
  \draw[fig] (1.564,2.300) -- (Bb);
  \draw[hid] (Bb) -- (1.004,1.150);
  \draw[fig] (1.004,1.150) -- (Cc);
  \draw[hid] (Cc) -- (0.444,0.000);
  \draw[fig] (0.444,0.000) -- (0.331,-0.232);
  \draw[fig] (1.795,3.402) -- (Ap);
  \draw[hid] (Ap) -- (2.154,2.300);
  \draw[fig] (2.154,2.300) -- (Bp);
  \draw[hid] (Bp) -- (2.529,1.150);
  \draw[fig] (2.529,1.150) -- (Cp);
  \draw[hid] (Cp) -- (2.837,0.205);
  \draw[fig] (2.837,0.205) -- (2.980,-0.232);
  \foreach \z in {Aa,Bb,Cc,Ap,Bp,Cp}{\dt{\z}}
  \node[vlab, left=3pt] at (Aa) {$A$};
  \node[vlab, left=3pt] at (Bb) {$B$};
  \node[vlab, below right=0pt] at (Cc) {$C$};
  \node[vlab, right=3pt] at (Ap) {$A'$};
  \node[vlab, right=3pt] at (Bp) {$B'$};
  \node[vlab] at (2.48,0.20) {$C'$};
  \node[vlab, above=1pt] at (1.795,3.402) {$l'$};
  \node[vlab, above=1pt] at (2.101,3.402) {$l$};
\end{tikzpicture}
\figcap{14--23}
\end{bkfigure}}

% ----------------------------------------------------------- Fig 14-24, 14-25
\newcommand{\FIGXIVTWENTYFOURTWENTYFIVE}{\begin{bkfigure}[\linewidth]
\begin{minipage}{0.50\linewidth}\centering
\begin{tikzpicture}[scale=1.34]
  \coordinate (V) at (1.72,3.05);
  \coordinate (Q1) at (0.42,0.72);   % P1
  \coordinate (Q2) at (0.95,1.28);   % P2
  \coordinate (Q3) at (2.62,1.22);   % P3
  \coordinate (Q4) at (1.42,0.42);   % P4
  \coordinate (Pt) at ($(Q4)!0.58!(Q3)$);   % P is a point OF the side P_4P_3
  % rays from V through each vertex, extended past the plane
  \foreach \z in {Q1,Q3,Q4,Pt}{\draw[fig] (V) -- ($(V)!1.42!(\z)$);}
  \draw[hid] (V) -- ($(V)!1.30!(Q2)$);
  % the polygon P1 P2 P3 P4
  \draw[fig] (Q1) -- (Q4) -- (Q3);
  \draw[hid] (Q3) -- (Q2) -- (Q1);
  \coordinate (Qq) at ($(V)!1.30!(Pt)$);
  \dt{Qq}
  \node[vlab, above=1.5pt] at (V) {$V$};
  \node[vlab, left=3.5pt] at (Q1) {$P_1$};
  \node[vlab] at (0.56,1.50) {$P_2$};
  \node[vlab, right=3.5pt] at (Q3) {$P_3$};
  \node[vlab, below left=3pt] at (Q4) {$P_4$};
  \node[vlab] at (2.44,0.84) {$P$};
  \node[vlab, right=3.5pt] at (Qq) {$Q$};
  \node[vlab] at (1.80,0.98) {$\pi$};
\end{tikzpicture}
\figcap{14--24}
\end{minipage}\hfill
\begin{minipage}{0.46\linewidth}\centering
\begin{tikzpicture}[scale=1.03]
  % Theorem 14-13's hint figure.  V sits at the upper LEFT and B far off to the
  % right, so VB and CB form the narrow wedge the book draws.  The rays V->A
  % and V->C overshoot A and C; the DASHED V->D stops dead at D.  (Ours used to
  % overshoot the dashed VD instead -- which read as a stray fragment -- and
  % carried no overshoot on VC at all.)
  \coordinate (Aa) at (0,0);
  \coordinate (V)  at (0.667,1.983);
  \coordinate (Bb) at (3.900,0.733);
  \coordinate (Cc) at (0.733,0.817);
  \coordinate (Dd) at ($(Aa)!0.402!(Bb)$);   % D lies on AB
  \draw[fig] (V) -- ($(V)!1.26!(Aa)$);
  \draw[fig] (V) -- ($(V)!2.17!(Cc)$);
  \draw[fig] (V) -- (Bb);
  \draw[hid] (V) -- (Dd);
  \draw[fig] (Aa) -- (Bb);
  \draw[fig] (Aa) -- (Cc); \draw[fig] (Cc) -- (Bb);
  \foreach \z in {V,Aa,Cc,Dd}{\dt{\z}}
  \node[vlab, above=2pt] at (V) {$V$};
  \node[vlab, left=3pt] at (Aa) {$A$};
  \node[vlab, right=3.5pt] at (Bb) {$B$};
  \node[vlab] at (0.54,0.99) {$C$};
  \node[vlab, below=3pt] at (Dd) {$D$};
\end{tikzpicture}
\figcap{14--25}
\end{minipage}
\end{bkfigure}}

% ----------------------------------------------------------- Fig 14-26, 14-27
\newcommand{\FIGXIVTWENTYSIXTWENTYSEVEN}{\begin{bkfigure}[\linewidth]
\begin{minipage}{0.44\linewidth}\centering
\begin{tikzpicture}[scale=1.407]
  % a tetrahedral angle flattened out: four triangles round C, which sits ON
  % the boundary -- the faces do not encircle the point.
  \coordinate (C)  at (1.52,1.30);
  \coordinate (T1) at (0.92,2.52);
  \coordinate (T2) at (2.38,2.52);
  \coordinate (Lf) at (0.15,1.30);
  \coordinate (B1) at (0.90,0.06);
  \coordinate (B2) at (2.32,0.05);
  \draw[fig] (T1) -- (T2) -- (C) -- (B2) -- (B1) -- (Lf) -- cycle;
  \draw[fig] (C) -- (T1); \draw[fig] (C) -- (Lf); \draw[fig] (C) -- (B1);
\end{tikzpicture}
\figcap{14--26}
\end{minipage}\hfill
\begin{minipage}{0.52\linewidth}\centering
\begin{tikzpicture}[scale=1.320]
  \XIVplane{P}{0}{0}{3.05}{1.15}{1.70}
  \coordinate (Q1) at (1.52,1.02);   % P1
  \coordinate (Q2) at (2.42,1.02);   % P2
  \coordinate (Q3) at (2.30,0.48);   % P3
  \coordinate (Q4) at (1.62,0.45);   % P4
  \coordinate (Q5) at (1.18,0.72);   % P5
  \fill[pattern=north east lines] (Q1) -- (Q2) -- (Q3) -- (Q4) -- (Q5) -- cycle;
  \draw[fig] (Q1) -- (Q2) -- (Q3) -- (Q4) -- (Q5) -- cycle;
  % labels sit outside the cell, on the outward normal from its centroid
  \node[vlab] at (1.285,1.265) {$P_1$};
  \node[vlab] at (2.735,1.170) {$P_2$};
  \node[vlab] at (2.610,0.315) {$P_3$};
  \node[vlab] at (1.435,0.255) {$P_4$};
  \node[vlab] at (0.830,0.710) {$P_5$};
\end{tikzpicture}
\figcap{14--27}
\end{minipage}
\end{bkfigure}}

% --------------------------------------------------------------- Fig 14-28
% The five regular polyhedra, laid out as in the book: tetrahedron,
% octahedron, icosahedron on the top row; cube and dodecahedron below.
\newcommand{\FIGXIVTWENTYEIGHT}{\begin{bkfigure}[\linewidth]
\begin{tikzpicture}[scale=1.076]
  % ---- tetrahedron
  \begin{scope}[xshift=0cm, yshift=3.30cm]
    \coordinate (t1) at (0.42,2.15); \coordinate (t2) at (0.00,0.42);
    \coordinate (t3) at (1.92,0.72); \coordinate (t4) at (1.05,0.00);
    \draw[fig] (t1) -- (t2) -- (t4) -- cycle;
    \draw[fig] (t1) -- (t3) -- (t4);
    \draw[hid] (t2) -- (t3);
    \node[slab, below] at (0.95,-0.34) {Tetrahedron};
  \end{scope}
  % ---- octahedron
  \begin{scope}[xshift=3.55cm, yshift=3.20cm]
    \coordinate (o1) at (1.15,2.42); \coordinate (o2) at (1.15,-0.10);
    \coordinate (o3) at (0.00,1.32); \coordinate (o4) at (2.32,1.28);
    \coordinate (o5) at (0.78,0.72); \coordinate (o6) at (1.62,1.72);
    \draw[fig] (o3) -- (o1) -- (o4) -- (o2) -- cycle;
    \draw[fig] (o3) -- (o5) -- (o4); \draw[fig] (o1) -- (o5) -- (o2);
    \draw[hid] (o3) -- (o6) -- (o4); \draw[hid] (o1) -- (o6) -- (o2);
    \node[slab, below] at (1.15,-0.52) {Octahedron};
  \end{scope}
  % ---- icosahedron (drawn as the book does: hexagonal silhouette)
  \begin{scope}[xshift=7.35cm, yshift=3.20cm]
    \foreach \i in {0,1,2,3,4,5}
      {\coordinate (h\i) at ({1.28+1.28*cos(30+60*\i)},{1.20+1.28*sin(30+60*\i)});}
    \draw[fig] (h0) -- (h1) -- (h2) -- (h3) -- (h4) -- (h5) -- cycle;
    % The book's projection: a horizontal chord across the top pair of
    % vertices carrying two interior nodes, a second chord across the bottom
    % pair carrying one, and no hidden edges drawn.
    \coordinate (u1) at ($(h2)!0.21!(h0)$);
    \coordinate (u2) at ($(h2)!0.81!(h0)$);
    \coordinate (ic) at ($(h3)!0.5!(h5)$);
    \draw[fig] (h2) -- (h0);
    \draw[fig] (h3) -- (h5);
    \draw[fig] (h1) -- (u1); \draw[fig] (h1) -- (u2);
    \draw[fig] (u1) -- (ic); \draw[fig] (u2) -- (ic);
    \draw[fig] (u1) -- (h3); \draw[fig] (u2) -- (h5);
    \draw[fig] (ic) -- (h4);
    \node[slab, below] at (1.28,-0.40) {Icosahedron};
  \end{scope}
  % ---- cube
  \begin{scope}[xshift=1.55cm, yshift=-0.20cm]
    \draw[fig] (0,0) -- (1.75,0) -- (1.75,1.75) -- (0,1.75) -- cycle;
    \draw[fig] (0,1.75) -- (0.68,2.35) -- (2.43,2.35) -- (1.75,1.75);
    \draw[fig] (1.75,0) -- (2.43,0.60) -- (2.43,2.35);
    \node[slab, below] at (1.20,-0.30) {Cube};
  \end{scope}
  % ---- dodecahedron
  \begin{scope}[xshift=5.65cm, yshift=-0.55cm]
    \foreach \i in {0,...,9}
      {\coordinate (d\i) at ({1.42+1.42*cos(18+36*\i)},{1.42+1.42*sin(18+36*\i)});}
    \draw[fig] (d0) -- (d1) -- (d2) -- (d3) -- (d4) -- (d5) -- (d6) -- (d7)
               -- (d8) -- (d9) -- cycle;
    \foreach \i in {0,...,4}
      {\coordinate (c\i) at ({1.42+0.62*cos(90+72*\i)},{1.42+0.62*sin(90+72*\i)});}
    \draw[fig] (c0) -- (c1) -- (c2) -- (c3) -- (c4) -- cycle;
    \draw[fig] (c0) -- (d2); \draw[fig] (c1) -- (d4); \draw[fig] (c2) -- (d6);
    \draw[fig] (c3) -- (d8); \draw[fig] (c4) -- (d0);
    \node[slab, below] at (1.42,-0.34) {Dodecahedron};
  \end{scope}
\end{tikzpicture}
\figcap{14--28}
\end{bkfigure}}

% ----------------------------------------------------------- Fig 14-29, 14-30
\newcommand{\FIGXIVTWENTYNINETHIRTY}{\begin{bkfigure}[\linewidth]
\begin{minipage}{0.50\linewidth}\centering
\begin{tikzpicture}[scale=1.10]
  % Theorem 14-17: a dihedral angle cut by two parallel planes.  The EDGE of
  % the dihedral is the left edge BL-TL of the "wall"; the near face F1 is the
  % solid quadrilateral, and the far face F2 lies behind it, so only its outer
  % edge shows -- dashed.  Each plane cuts the dihedral in a plane angle: at
  % the top it is the thin wedge TL-TR / TL-U, struck with an ARC and a
  % perpendicular tick; at the bottom the same wedge, its second side hidden
  % behind the wall and therefore dashed.  (The old drawing had two arrowheads
  % where those two arcs belong and no ticks at all.)
  \XIVplane{PB}{-1.30}{-0.85}{3.30}{0.95}{1.25}
  \XIVplane{PT}{-0.75}{1.55}{3.30}{0.95}{1.25}
  \coordinate (BLw) at (0,0);        \coordinate (TLw) at (0.55,2.40);
  \coordinate (BRw) at (1.25,0.12);  \coordinate (TRw) at (1.80,2.52);
  \coordinate (Ub)  at (1.20,0.34);  \coordinate (Ut)  at (1.75,2.74);
  \draw[fig] (TLw) -- (TRw) -- (BRw) -- (BLw) -- cycle;   % near face F1
  \draw[fig] (TLw) -- (Ut);                               % F2 trace, top plane
  \draw[hid] (Ut) -- (Ub);                                % F2 far edge, hidden
  \draw[hid] (BLw) -- (Ub);                       % F2 trace, bottom, hidden
  \draw[farc] ([shift=(5.48:0.78)]TLw) arc (5.48:15.82:0.78);
  \draw[farc] ([shift=(5.48:0.42)]BLw) arc (5.48:15.82:0.42);
  \coordinate (Qt) at ($(TLw)!0.66!(Ut)$);
  \draw[farc] (Qt) -- ($(TLw)!(Qt)!(TRw)$);
  \coordinate (Qb) at ($(BLw)!0.60!(Ub)$);
  \draw[farc] (Qb) -- ($(BLw)!(Qb)!(BRw)$);
\end{tikzpicture}
\figcap{14--29}
\end{minipage}\hfill
\begin{minipage}{0.46\linewidth}\centering
\begin{tikzpicture}[scale=1.05]
  % Definition 14-8.  A dihedral edge E0->E1 lies in the horizontal plane; one
  % face rises from it as the tilted flap, and the plane PERPENDICULAR to the
  % edge is the clean rectangle standing in front.  The edge passes BEHIND the
  % rectangle, so it is dashed, and it RISES to the right (ours used to run
  % flat along the rectangle's foot, which killed the depth reading); a
  % perpendicular tick marks the right angle at the rectangle's foot.
  \draw[fig] (-0.950,-0.523) -- (2.483,-0.110) -- (3.633,1.200)
             -- (0.183,0.800) -- cycle;                 % the horizontal plane
  \coordinate (E0) at (0,0);          \coordinate (E1) at (1.933,0.850);
  \coordinate (r)  at (1.517,0.150);  \coordinate (u)  at (0.067,1.140);
  \draw[fig] ($(E0)+(u)$) -- ($(E1)+(u)$) -- (E1);       % the flap
  \draw[hid] (E1) -- (E0);                              % the dihedral edge
  \draw[fig] (E0) -- ($(E0)+(r)$) -- ($(E0)+(r)+(u)$) -- ($(E0)+(u)$)
             -- cycle;                                  % the perpendicular
  \coordinate (Qe) at ($(E0)!0.25!(E1)$);
  \draw[farc] (Qe) -- ($(E0)!(Qe)!($(E0)+(r)$)$);
\end{tikzpicture}
\figcap{14--30}
\end{minipage}
\end{bkfigure}}

% ---------------------------------------------------------------- Fig 14-31
% Split out of the old \FIGXIVTHIRTYONETHIRTYTHREE block: Ex.~14-12 cites
% 14-31 in \#1, 14-32 in \#2 and 14-33 in \#3, so each stands on its own now.
\newcommand{\FIGXIVTHIRTYONE}{\begin{bkfigure}[0.74\linewidth]
\begin{tikzpicture}[scale=1.45]
  % A right hexagonal prism with its bases IN two parallel planes.  The old
  % hexagon put its vertices at 30+60i degrees, so the pairs (30,330) and
  % (150,210) shared an x -- their lateral edges landed on top of one another
  % and the solid read as a flat bar with a single dashed line down it.  These
  % vertex angles are the book's own irregular ones: no two share an x, so the
  % near side shows FOUR solid lateral edges and the far side TWO dashed.
  \XIVplane{PB}{0}{0}{4.00}{1.10}{1.05}
  \XIVplane{PT}{0}{2.60}{4.00}{1.10}{1.05}
  \foreach \i/\a in {0/-10, 1/36, 2/118, 3/181, 4/213, 5/301}
    {\coordinate (b\i) at ({2.20+1.00*cos(\a)},{0.52+0.44*sin(\a)});
     \coordinate (t\i) at ({2.20+1.00*cos(\a)},{3.12+0.44*sin(\a)});}
  % top base: seen from above, so the whole outline is visible
  \draw[fig] (t0) -- (t1) -- (t2) -- (t3) -- (t4) -- (t5) -- cycle;
  % bottom base: near half solid, far half hidden
  \draw[fig] (b3) -- (b4) -- (b5) -- (b0);
  \draw[hid] (b0) -- (b1) -- (b2) -- (b3);
  \foreach \i in {0,3,4,5}{\draw[fig] (b\i) -- (t\i);}
  \foreach \i in {1,2}{\draw[hid] (b\i) -- (t\i);}
\end{tikzpicture}
\figcap{14--31}
\end{bkfigure}}

% ---------------------------------------------------------------- Fig 14-32
\newcommand{\FIGXIVTHIRTYTWO}{\begin{bkfigure}[0.44\linewidth]
\begin{tikzpicture}[scale=1.10]
  % pentagonal pyramid.  q0 and q2 are the silhouette vertices and q1 the front
  % one, so those three lateral edges are solid; q3 and q4 lie behind, so their
  % lateral edges AND the three base edges joining them are dashed.  (The old
  % version drew the back edge q2--q3 and the lateral to q3 solid.)
  \coordinate (ap) at (1.42,2.72);
  \foreach \i in {0,...,4}
    {\coordinate (q\i) at ({1.42+1.32*cos(198+72*\i)},{0.62+0.46*sin(198+72*\i)});}
  \draw[fig] (q0) -- (q1) -- (q2);
  \draw[hid] (q2) -- (q3) -- (q4) -- (q0);
  \foreach \i in {0,1,2}{\draw[fig] (ap) -- (q\i);}
  \foreach \i in {3,4}{\draw[hid] (ap) -- (q\i);}
\end{tikzpicture}
\figcap{14--32}
\end{bkfigure}}

% ---------------------------------------------------------------- Fig 14-33
\newcommand{\FIGXIVTHIRTYTHREE}{\begin{bkfigure}[0.54\linewidth]
\begin{tikzpicture}[scale=1.15]
  % parallelepiped
  \coordinate (ea) at (2.65,0.32);   % length
  \coordinate (eb) at (1.15,0.95);   % depth
  \coordinate (ec) at (-0.62,1.05);  % height
  \coordinate (A0) at (0.62,0);
  \coordinate (A1) at ($(A0)+(ea)$);
  \coordinate (A2) at ($(A1)+(eb)$);
  \coordinate (A3) at ($(A0)+(eb)$);
  \coordinate (B0) at ($(A0)+(ec)$);
  \coordinate (B1) at ($(A1)+(ec)$);
  \coordinate (B2) at ($(A2)+(ec)$);
  \coordinate (B3) at ($(A3)+(ec)$);
  \draw[fig] (A0) -- (A1) -- (A2); \draw[fig] (A0) -- (B0);
  \draw[fig] (A1) -- (B1); \draw[fig] (A2) -- (B2);
  \draw[fig] (B0) -- (B1) -- (B2) -- (B3) -- cycle;
  \draw[hid] (A0) -- (A3) -- (A2);
  \draw[hid] (A3) -- (B3);
\end{tikzpicture}
\figcap{14--33}
\end{bkfigure}}

% ----------------------------------------------------------- Fig 14-34, 14-35
\newcommand{\FIGXIVTHIRTYFOURTHIRTYFIVE}{\begin{bkfigure}[\linewidth]
\begin{minipage}{0.48\linewidth}\centering
\begin{tikzpicture}[scale=1.046]
  % rectangular box, edges labelled l, w, h
  \coordinate (a0) at (0,0);       \coordinate (a1) at (2.85,0);
  \coordinate (a2) at (2.85,1.55); \coordinate (a3) at (0,1.55);
  \coordinate (d)  at (0.85,0.72);
  \draw[fig] (a0) -- (a1) -- (a2) -- (a3) -- cycle;
  \draw[fig] (a3) -- ($(a3)+(d)$) -- ($(a2)+(d)$) -- (a2);
  \draw[fig] (a1) -- ($(a1)+(d)$) -- ($(a2)+(d)$);
  \node[vlab] at (3.50,1.72) {$w$};
  \node[vlab] at (1.425,1.28) {$l$};
  \node[vlab] at (3.12,0.775) {$h$};
\end{tikzpicture}
\figcap{14--34}
\end{minipage}\hfill
\begin{minipage}{0.48\linewidth}\centering
\begin{tikzpicture}[scale=1.25]
  % Theorem 14-19: one polyhedron as the union of two, sharing the HATCHED
  % face.  The book's shared section is a big triangle spanning the whole near
  % face (ours was a small patch in one corner), its boundary heavy and solid,
  % its interior filled with dense steep hatching.  The one hidden vertex is F;
  % only its two hidden edges FA and FC are dashed.
  \coordinate (Av) at (-1.117,1.917);
  \coordinate (Bv) at ( 1.350,2.033);
  \coordinate (Cv) at ( 2.883,1.550);
  \coordinate (Dv) at ( 2.483,0.167);
  \coordinate (Ev) at ( 0,0);
  \coordinate (Fv) at (-0.017,1.167);
  \begin{scope}
    \clip (Bv) -- (Fv) -- (Ev) -- cycle;
    \foreach \k in {0,...,27}
      {\draw[fig, line width=0.35pt]
         ({-1.40+0.107*\k},-0.30) -- ({-1.40+0.107*\k+1.05},2.55);}
  \end{scope}
  \draw[fig] (Av) -- (Bv) -- (Cv) -- (Dv) -- (Ev) -- cycle;
  \draw[fig] (Bv) -- (Ev);  \draw[fig] (Bv) -- (Dv);
  \draw[fig] (Bv) -- (Fv) -- (Ev);
  \draw[hid] (Fv) -- (Av);  \draw[hid] (Fv) -- (Cv);
\end{tikzpicture}
\figcap{14--35}
\end{minipage}
\end{bkfigure}}

% ----------------------------------------------------------- Fig 14-36, 14-37
\newcommand{\FIGXIVTHIRTYSIXTHIRTYSEVEN}{\begin{bkfigure}[\linewidth]
\begin{minipage}{0.46\linewidth}\centering
\begin{tikzpicture}[scale=0.672]
  \coordinate (ea) at (2.85,0);      % length
  \coordinate (eb) at (0.95,0.62);   % depth
  \coordinate (ec) at (-0.55,0.95);  % height (leaning left)
  \coordinate (A0) at (0.62,0.10);
  \coordinate (A1) at ($(A0)+(ea)$);  \coordinate (A2) at ($(A1)+(eb)$);
  \coordinate (A3) at ($(A0)+(eb)$);
  \coordinate (B0) at ($(A0)+(ec)$);  \coordinate (B1) at ($(A1)+(ec)$);
  \coordinate (B2) at ($(A2)+(ec)$);  \coordinate (B3) at ($(A3)+(ec)$);
  \draw[fig] (A0) -- (A1) -- (A2); \draw[fig] (A0) -- (B0);
  \draw[fig] (A1) -- (B1); \draw[fig] (A2) -- (B2);
  \draw[fig] (B0) -- (B1) -- (B2) -- (B3) -- cycle;
  \draw[hid] (A0) -- (A3) -- (A2);  \draw[hid] (A3) -- (B3);
\end{tikzpicture}
\figcap{14--36}
\end{minipage}\hfill
\begin{minipage}{0.50\linewidth}\centering
\begin{tikzpicture}[scale=0.672]
  \coordinate (ea) at (3.15,0);
  \coordinate (eb) at (1.05,0.68);
  \coordinate (ec) at (-0.62,1.05);
  \coordinate (A0) at (0.68,0.10);
  \coordinate (A1) at ($(A0)+(ea)$);  \coordinate (A2) at ($(A1)+(eb)$);
  \coordinate (A3) at ($(A0)+(eb)$);
  \coordinate (B0) at ($(A0)+(ec)$);  \coordinate (B1) at ($(A1)+(ec)$);
  \coordinate (B2) at ($(A2)+(ec)$);  \coordinate (B3) at ($(A3)+(ec)$);
  \draw[fig] (A0) -- (A1) -- (A2); \draw[fig] (A0) -- (B0);
  \draw[fig] (A1) -- (B1); \draw[fig] (A2) -- (B2);
  \draw[fig] (B0) -- (B1) -- (B2) -- (B3) -- cycle;
  \draw[hid] (A0) -- (A3) -- (A2);  \draw[hid] (A3) -- (B3);
  % Altitude: a TRUE vertical between the two horizontal faces, as in the book
  % -- not a copy of the (leaning) lateral edge.  Its top sits on the edge
  % B1B2 of the top face; dropping by the height of ec lands it on the
  % hidden base face, so the segment is dashed.
  \coordinate (Ht) at ($(B1)!0.5!(B2)$);
  \coordinate (Hb) at ($(Ht)-(0,1.05)$);
  \draw[hid] (Ht) -- (Hb);
  \draw[fig, line width=0.4pt] ($(Ht)-(0.14,0)$) -- ($(Ht)+(0.14,0)$);
  \draw[fig, line width=0.4pt] ($(Hb)-(0.14,0)$) -- ($(Hb)+(0.14,0)$);
  \node[slab, right] at (4.95,1.72) {Altitude $h$};
  \node[slab, right] at (4.95,0.42) {Area $B$};
  \draw[arw] (4.90,1.72) -- ($(Ht)-(0.04,0.14)$);
  \draw[arw] (4.90,0.42) -- ($(A1)!0.55!(A2)+(0.12,-0.04)$);
\end{tikzpicture}
\figcap{14--37}
\end{minipage}
\end{bkfigure}}

% --------------------------------------------------------------- Fig 14-38
\newcommand{\FIGXIVTHIRTYEIGHT}{\begin{bkfigure}[0.94\linewidth]
\begin{tikzpicture}[scale=0.950]
  % (a) a parallelepiped of which a triangular prism is half
  \begin{scope}
    \coordinate (ea) at (2.55,0); \coordinate (eb) at (1.02,0.72);
    \coordinate (ec) at (0,1.42);
    \coordinate (A0) at (0,0);   \coordinate (A1) at ($(A0)+(ea)$);
    \coordinate (A2) at ($(A1)+(eb)$); \coordinate (A3) at ($(A0)+(eb)$);
    \coordinate (B0) at ($(A0)+(ec)$); \coordinate (B1) at ($(A1)+(ec)$);
    \coordinate (B2) at ($(A2)+(ec)$); \coordinate (B3) at ($(A3)+(ec)$);
    \draw[fig] (A0) -- (A1) -- (A2); \draw[fig] (A0) -- (B0);
    \draw[fig] (A1) -- (B1); \draw[fig] (A2) -- (B2);
    \draw[fig] (B0) -- (B1) -- (B2) -- (B3) -- cycle;
    \draw[hid] (A0) -- (A3) -- (A2); \draw[hid] (A3) -- (B3);
    % The cut runs BACK corner -> FRONT corner of the top face (B3 -> B1).  The
    % old B0--B2 diagonal ran almost along the top edge, so it read as a stray
    % duplicate edge; and the book draws no companion diagonal on the base.
    \draw[fig] (B3) -- (B1);
    \node[slab, below] at (1.72,-0.55) {(a)};
  \end{scope}
  % (b) an arbitrary prism, cut into triangular prisms
  \begin{scope}[xshift=5.35cm]
    \coordinate (ea) at (2.95,0); \coordinate (eb) at (1.02,0.72);
    \coordinate (ec) at (0,1.42);
    \coordinate (A0) at (0,0);   \coordinate (A1) at ($(A0)+(ea)$);
    \coordinate (A2) at ($(A1)+(eb)$); \coordinate (A3) at ($(A0)+(eb)$);
    \coordinate (B0) at ($(A0)+(ec)$); \coordinate (B1) at ($(A1)+(ec)$);
    \coordinate (B2) at ($(A2)+(ec)$); \coordinate (B3) at ($(A3)+(ec)$);
    \draw[fig] (A0) -- (A1) -- (A2); \draw[fig] (A0) -- (B0);
    \draw[fig] (A1) -- (B1); \draw[fig] (A2) -- (B2);
    \draw[fig] (B0) -- (B1) -- (B2) -- (B3) -- cycle;
    \draw[hid] (A0) -- (A3) -- (A2); \draw[hid] (A3) -- (B3);
    % a vertical cutting plane
    \coordinate (S0) at ($(A0)!0.62!(A1)$);
    \coordinate (S1) at ($(B0)!0.62!(B1)$);
    \coordinate (S2) at ($(B3)!0.62!(B2)$);
    \draw[fig] (S0) -- (S1) -- (S2);
    % the cut's hidden edge is the VERTICAL down the back face, not a line
    % across the base: it drops from the top-back edge to the hidden bottom one
    \draw[hid] (S2) -- ($(A3)!0.62!(A2)$);
    \node[slab, below] at (2.02,-0.55) {(b)};
  \end{scope}
\end{tikzpicture}
\figcap{14--38}
\end{bkfigure}}

% --------------------------------------------------------------- Fig 14-39
\newcommand{\FIGXIVTHIRTYNINE}{\begin{bkfigure}[0.50\linewidth]
\begin{tikzpicture}[scale=1.560]
  % base is a parallelogram; ct is where its diagonals meet, so the altitude
  % really does fall on the centre of the base.
  % The base is a parallelogram whose front vertex is NOT under the apex, so
  % the altitude reads as its own line instead of hiding inside the front
  % edge -- which is how the book draws it.
  \coordinate (b1) at (0,0.55); \coordinate (b2) at (1.30,0);
  \coordinate (b3) at (3.05,0.55); \coordinate (b4) at (1.75,1.10);
  \coordinate (ct) at (1.525,0.55);          % centre = crossing of the diagonals
  \coordinate (ap) at (1.525,2.85);
  \draw[fig] (b1) -- (b2) -- (b3);
  \draw[hid] (b3) -- (b4) -- (b1);
  \draw[fig] (ap) -- (b1); \draw[fig] (ap) -- (b2); \draw[fig] (ap) -- (b3);
  \draw[hid] (ap) -- (b4);
  \draw[aux, ->, >=latex] (ap) -- (ct);
  \node[vlab] at (1.74,0.92) {$h$};
\end{tikzpicture}
\figcap{14--39}
\end{bkfigure}}

% --------------------------------------------------------------- Fig 14-40
% Two similar pyramids: the right one is the left one scaled by 1.55 about V'.
\newcommand{\FIGXIVFORTY}{\begin{bkfigure}[0.66\linewidth]
\begin{tikzpicture}[scale=0.780]
  % ---- small pyramid, vertex V at the lower left of the base
  \begin{scope}
    \coordinate (V)  at (0,0);
    \coordinate (v2) at (1.35,0);
    \coordinate (v3) at (1.92,0.42);
    \coordinate (ap) at (0.78,1.62);
    \draw[fig] (V) -- (v2) -- (v3);
    \draw[fig] (ap) -- (V); \draw[fig] (ap) -- (v2); \draw[fig] (ap) -- (v3);
    \node[vlab, below=3.5pt] at (V) {$V$};
    \node[vlab, above right=3pt] at ($(ap)!0.45!(v3)$) {$l$};
  \end{scope}
  % ---- large pyramid, the same shape scaled by 1.55
  \begin{scope}[xshift=2.75cm]
    \coordinate (Vp)  at (0,0);
    \coordinate (w2) at (2.0925,0);
    \coordinate (w3) at (2.976,0.651);
    \coordinate (bp) at (1.209,2.511);
    \draw[fig] (Vp) -- (w2) -- (w3);
    \draw[fig] (bp) -- (Vp); \draw[fig] (bp) -- (w2); \draw[fig] (bp) -- (w3);
    \node[vlab, below=3.5pt] at (Vp) {$V'$};
    \node[vlab, above right=3pt] at ($(bp)!0.45!(w3)$) {$l'$};
  \end{scope}
\end{tikzpicture}
\figcap{14--40}
\end{bkfigure}}

% --------------------------------------------------------------- Fig 14-41
\newcommand{\FIGXIVFORTYONE}{\begin{bkfigure}[0.82\linewidth]
\begin{tikzpicture}[scale=0.858]
  % ---- cylinder, tilted
  \begin{scope}[rotate=-18]
    \draw[fig] (0,0) ellipse (0.62 and 0.24);
    \draw[fig] (-0.62,0) -- (-0.62,-2.35);
    \draw[fig] (0.62,0) -- (0.62,-2.35);
    \draw[fig] (-0.62,-2.35) arc (180:360:0.62 and 0.24);
    \draw[hid] (0.62,-2.35) arc (0:180:0.62 and 0.24);
  \end{scope}
  % ---- cone
  \begin{scope}[xshift=2.85cm, yshift=-2.05cm]
    \draw[fig] (-0.78,0) arc (180:360:0.78 and 0.28);
    \draw[hid] (0.78,0) arc (0:180:0.78 and 0.28);
    \draw[fig] (-0.78,0) -- (0,2.05) -- (0.78,0);
  \end{scope}
  % ---- sphere, with the equator
  \begin{scope}[xshift=5.75cm, yshift=-1.02cm]
    \draw[fig] (0,0) circle (1.08);
    \draw[fig] (-1.08,0) arc (180:360:1.08 and 0.38);
    \draw[hid] (1.08,0) arc (0:180:1.08 and 0.38);
  \end{scope}
\end{tikzpicture}
\figcap{14--41}
\end{bkfigure}}

% ---------------------------------------------------------------- Fig 14-42
% Split out of the old \FIGXIVFORTYTWOFORTYTHREE block: Ex.~14-16 \#2 cites
% 14-42, while 14-43 belongs to the body text about Theorem 14-24.
\newcommand{\FIGXIVFORTYTWO}{\begin{bkfigure}[0.50\linewidth]
\begin{tikzpicture}[scale=1.35]
  % Right circular cone with an inscribed pyramid.  The book's base polygon is
  % a HEXAGON with vertices at the two extremes of the ellipse plus two in
  % front and two behind: the three near base chords are solid, the three far
  % ones dashed, and the two hidden lateral edges stand clear of the solid
  % ones instead of smearing into them at the apex.  (Ours drew a single
  % dashed chord and no near chords at all, so nothing read as a pyramid, and
  % it carried a spurious dashed diameter across the base.)
  \coordinate (ap) at (0,2.07);
  \draw[fig] (0,0) ellipse (1.42 and 0.61);
  \foreach \i/\a in {0/180, 1/230, 2/310, 3/0, 4/50, 5/130}
    {\coordinate (e\i) at ({1.42*cos(\a)},{0.61*sin(\a)});}
  \draw[fig] (e0) -- (e1) -- (e2) -- (e3);          % near base chords
  \draw[hid] (e3) -- (e4) -- (e5) -- (e0);          % far base chords
  \foreach \i in {0,1,2,3}{\draw[fig] (ap) -- (e\i);}
  \foreach \i in {4,5}{\draw[hid] (ap) -- (e\i);}
\end{tikzpicture}
\figcap{14--42}
\end{bkfigure}}

% ---------------------------------------------------------------- Fig 14-43
\newcommand{\FIGXIVFORTYTHREE}{\begin{bkfigure}[0.48\linewidth]
\begin{tikzpicture}[scale=0.980]
  % a sphere carved into small "pyramids" with a common apex at the centre
  \begin{scope}[even odd rule]
    \clip (0,0) circle (1.42);
    \clip (0.28,-0.24) circle (1.28) (-1.6,-1.6) rectangle (1.6,1.6);
    \fill[pattern=north east lines] (-1.6,-1.6) rectangle (1.6,1.6);
  \end{scope}
  \draw[fig] (0,0) circle (1.42);
  % one small pyramid, apex at the centre
  \coordinate (O) at (0,0);
  \coordinate (k1) at (0.72,0.42); \coordinate (k2) at (1.02,0.62);
  \coordinate (k3) at (1.02,0.05); \coordinate (k4) at (0.72,-0.10);
  \draw[fig] (k1) -- (k2) -- (k3) -- (k4) -- cycle;
  \draw[fig] (O) -- (k1); \draw[fig] (O) -- (k4);
  \draw[hid] (O) -- (k2); \draw[hid] (O) -- (k3);
  \dt{O}
\end{tikzpicture}
\figcap{14--43}
\end{bkfigure}}

% figures15.tex -- Chapter 15 (Analytic Geometry) figures redrawn in TikZ.
% Coordinate-geometry chapter: every figure is exactly constructible from the
% mathematics in the text (axes, plotted points, lines, circles), so vertices
% are computed, never eyeballed. Proportions were measured off the printed
% figures at 400 dpi (PDF pp. 267-279) during the figure review; the readings
% are recorded in chapters/ch15-UNCERTAIN.md.
%
% Helper macros used only inside this file (check_labels greps FIG[A-Z]+ for
% the figure macros, so these helpers are not picked up as figures).
% `positioning' is needed for the "above left=<a> and <b>" node syntax; it is
% loaded here rather than in brumfiel.sty because this file is also compiled
% standalone by tools/check_labels.py.
\usetikzlibrary{positioning}

% Standard axes WITH the letters x and y: {xmin}{xmax}{ymin}{ymax}.
% The book letters the axes in Figs 15-6 to 15-19 and leaves them unlettered
% in Fig 15-4, so the two variants are kept apart.
\newcommand{\XVaxes}[4]{%
  \draw[fig,-latex] (#1,0) -- (#2,0);
  \node[vlab,right] at (#2,0) {$x$};
  \draw[fig,-latex] (0,#3) -- (0,#4);
  \node[vlab,right] at (0,#4) {$y$};}
% Same axes, unlettered (Fig 15-4).
\newcommand{\XVaxesnl}[4]{%
  \draw[fig,-latex] (#1,0) -- (#2,0);
  \draw[fig,-latex] (0,#3) -- (0,#4);}
% Origin label
\newcommand{\XVzero}{\node[vlab,below left=-1pt and 0pt] at (0,0) {$0$};}
% Unit ticks on the x-axis, from #1 to #2 step 1 (skipping 0)
\newcommand{\XVxticks}[2]{%
  \foreach \t in {#1,...,#2}{\ifnum\t=0\else
    \draw[tick] (\t,-0.11) -- (\t,0.11);\fi}}
\newcommand{\XVyticks}[2]{%
  \foreach \t in {#1,...,#2}{\ifnum\t=0\else
    \draw[tick] (-0.11,\t) -- (0.11,\t);\fi}}

% ---------------------------------------------------------------- Fig 15-1
% Coordinatized line: O, the unit point P_1, a positive point P_x and a
% negative point P'_{x'}.  Collinear by construction.  Spacings follow the
% printed figure (O -> P_1 : P_1 -> P_x : P'_{x'} -> O = 1 : 1.55 : 1.22).
\newcommand{\FIGXVONE}{\begin{bkfigure}[\linewidth]
\begin{tikzpicture}[scale=1.0]
  \draw[fig] (-0.3,0) -- (10.3,0);
  \foreach \x/\below/\above in {0.6/{P'_{x'}}/{x'}, 3.5/{O}/{},
                                5.9/{P_1}/{1}, 9.6/{P_x}/{x}}
    {\dt{\x,0}
     \node[vlab,below=2pt] at (\x,0) {$\below$};}
  \foreach \x/\above in {0.6/{x'}, 5.9/{1}, 9.6/{x}}
    {\node[vlab,above=2pt] at (\x,0) {$\above$};}
\end{tikzpicture}
\figcap{15--1}
\end{bkfigure}}

% ---------------------------------------------------------------- Fig 15-2
% Two PERPENDICULAR lines through O, coordinatized, drawn obliquely to make
% the point that the choice of axes is arbitrary.
% abscissa positive direction 210 deg; ordinate positive direction 300 deg.
% Measured off the plate: abscissa runs 2.6 units back / 5.0 forward, ordinate
% 1.5 back / 3.4 forward.  Every numeral sits on the side the book puts it on:
% the positive abscissa numerals on the 120 deg side, the positive ordinate
% numerals on the 30 deg side, and each -1 thrown to the far side of its own
% axis so the two do not crowd the 0.
\newcommand{\FIGXVTWO}{\begin{bkfigure}[\linewidth]
\begin{tikzpicture}[scale=1.040]
  % unit vectors
  \def\ax{-0.86603} \def\ay{-0.5}      % abscissa positive (210 deg)
  \def\ox{0.5}      \def\oy{-0.86603}  % ordinate positive (300 deg)
  % the two axes
  \draw[fig] (2.6*-\ax,2.6*-\ay) -- (5.0*\ax,5.0*\ay);
  \draw[fig] (1.5*-\ox,1.5*-\oy) -- (3.4*\ox,3.4*\oy);
  % ticks + numerals on the abscissa axis (tick direction = ordinate dir)
  \foreach \t/\n in {1/1, 2/2, 3/3}
    {\draw[tick] ({\t*\ax-0.20*\ox},{\t*\ay-0.20*\oy})
              -- ({\t*\ax+0.20*\ox},{\t*\ay+0.20*\oy});
     \node[slab] at ({\t*\ax-0.46*\ox},{\t*\ay-0.46*\oy}) {$\n$};}
  \draw[tick] ({-\ax-0.20*\ox},{-\ay-0.20*\oy})
           -- ({-\ax+0.20*\ox},{-\ay+0.20*\oy});
  \node[slab] at ({-\ax-0.50*\ox},{-\ay-0.50*\oy}) {$-1$};
  % ticks + numerals on the ordinate axis (tick direction = abscissa dir)
  \foreach \t/\n in {1/1, 2/2, 3/3}
    {\draw[tick] ({\t*\ox-0.20*\ax},{\t*\oy-0.20*\ay})
              -- ({\t*\ox+0.20*\ax},{\t*\oy+0.20*\ay});
     \node[slab] at ({\t*\ox-0.50*\ax},{\t*\oy-0.50*\ay}) {$\n$};}
  \draw[tick] ({-\ox-0.20*\ax},{-\oy-0.20*\ay})
           -- ({-\ox+0.20*\ax},{-\oy+0.20*\ay});
  \node[slab] at ({-\ox+0.66*\ax},{-\oy+0.66*\ay}) {$-1$};
  % 0 in the open sector between the two negative rays, as in the book
  \node[slab] at (0.16,0.60) {$0$};
  \node[slab,left]  at (5.15*\ax,5.15*\ay) {Abscissas};
  \node[slab,right] at (3.55*\ox,3.55*\oy) {Ordinates};
\end{tikzpicture}
\figcap{15--2}
\end{bkfigure}}

% ---------------------------------------------------------------- Fig 15-3
% Theorem 15-1: P,Q in the SAME quadrant => PQ meets neither axis;
% P,Q in DIFFERENT quadrants => PQ meets at least one axis.
\newcommand{\FIGXVTHREE}{\begin{bkfigure}[\linewidth]
\begin{tikzpicture}[scale=1.014]
  \draw[fig] (-4.3,0) -- (3.9,0);
  \draw[fig] (0,2.9) -- (0,-1.5);
  % same quadrant (I): segment meets neither axis
  \coordinate (P1) at (1.6,2.15);  \coordinate (Q1) at (2.65,1.35);
  \draw[key] (P1) -- (Q1);
  \node[vlab,above=1pt] at (P1) {$P$};
  \node[vlab,right=1pt] at (Q1) {$Q$};
  % quadrant II to quadrant III: crosses the axis of abscissas
  \coordinate (Q2) at (-2.45,1.5); \coordinate (P2) at (-3.75,-0.62);
  \draw[key] (Q2) -- (P2);
  \node[vlab,above=1pt] at (Q2) {$Q$};
  \node[vlab,left=1pt]  at (P2) {$P$};
  % quadrant II to quadrant IV: crosses both axes
  \coordinate (P3) at (-1.25,1.05); \coordinate (Q3) at (2.35,-0.72);
  \draw[key] (P3) -- (Q3);
  \node[vlab,above=1pt] at (P3) {$P$};
  \node[vlab,right=1pt] at (Q3) {$Q$};
  \foreach \p in {P1,Q1,Q2,P2,P3,Q3} {\dt{\p}}
\end{tikzpicture}
\figcap{15--3}
\end{bkfigure}}

% ---------------------------------------------------------------- Fig 15-4
% The four quadrants, numbered.  The book leaves this pair of axes unlettered
% (no x, no y) and runs the left branch longer than the right.
\newcommand{\FIGXVFOUR}{\begin{bkfigure}[\linewidth]
\begin{tikzpicture}[scale=0.936]
  \XVaxesnl{-3.6}{2.9}{-1.2}{2.7}
  \XVxticks{1}{2}
  \XVyticks{1}{2}
  \node[slab,above] at (1,0.12) {$1$};
  \node[slab,above] at (2,0.12) {$2$};
  \node[slab,left]  at (-0.12,1) {$1$};
  \node[slab,left]  at (-0.12,2) {$2$};
  \XVzero
  \node[vlab] at (1.40,1.50)  {\textup{I}};
  \node[vlab] at (-1.70,1.50) {\textup{II}};
  \node[vlab] at (-1.75,-0.87){\textup{III}};
  \node[vlab] at (1.36,-0.87) {\textup{IV}};
\end{tikzpicture}
\figcap{15--4}
\end{bkfigure}}

% ---------------------------------------------------------------- Fig 15-5
% Rotated (still perpendicular) axes; the reader is asked which quadrant is
% which.  Ordinate positive direction 150 deg, abscissa positive 240 deg --
% the printed pair measures about 92 deg apart, so they are squared up.
% Branch lengths follow the plate: the ordinate axis runs 1.8 units forward
% and 3.9 back, the abscissa 2.05 forward and 1.95 back.
\newcommand{\FIGXVFIVE}{\begin{bkfigure}[\linewidth]
\begin{tikzpicture}[scale=1.040]
  \def\ox{-0.86603} \def\oy{0.5}       % ordinate positive (150 deg)
  \def\ax{-0.5}     \def\ay{-0.86603}  % abscissa positive (240 deg)
  \draw[fig,-latex] (3.9*-\ox,3.9*-\oy) -- (1.8*\ox,1.8*\oy);
  \draw[fig,-latex] (1.95*-\ax,1.95*-\ay) -- (2.05*\ax,2.05*\ay);
  \dt{\ox,\oy}
  \dt{\ax,\ay}
  \node[slab,above] at (\ox,\oy) {$1$};
  \node[slab,right] at (\ax,\ay) {$1$};
  \node[slab,above] at (2.0*\ox,2.0*\oy) {Ordinates};
  \node[slab,below] at (2.25*\ax,2.25*\ay) {Abscissas};
\end{tikzpicture}
\figcap{15--5}
\end{bkfigure}}

% ---------------------------------------------------------------- Fig 15-6
% Definition 15-4: l_1 through P parallel to the axis of ordinates meets the
% first axis in A; l_2 through P parallel to the axis of abscissas meets the
% second axis in B.  A sits on the positive abscissa axis and B on the
% negative ordinate axis, so P falls in quadrant IV, as in the book.
\newcommand{\FIGXVSIX}{\begin{bkfigure}[\linewidth]
\begin{tikzpicture}[scale=1.040]
  \XVaxes{-1.6}{2.8}{-3.2}{1.65}
  \coordinate (A) at (1.55,0);
  \coordinate (B) at (0,-1.70);
  \coordinate (P) at (1.55,-1.70);
  \draw[aux] (1.55,1.55) -- (1.55,-3.5);
  \draw[aux] (-1.65,-1.70) -- (3.2,-1.70);
  \foreach \p in {A,B,P} {\dt{\p}}
  \node[vlab,above left=-1pt and 0pt] at (0,0) {$O$};
  \node[vlab,above right=1pt and 1pt] at (A) {$A$};
  \node[vlab,above left=1pt and 1pt]  at (B) {$B$};
  \node[vlab,above right=1pt and 1pt] at (P) {$P$};
  \node[vlab,right] at (3.2,-1.70) {$l_2$};
  \node[vlab,below right=-1pt and 0pt] at (1.55,-3.5) {$l_1$};
\end{tikzpicture}
\figcap{15--6}
\end{bkfigure}}

% ---------------------------------------------------------------- Fig 15-7
% The line x = 3: parallel to the y-axis, 3 units to its right.  The book
% ticks only 1 and 2 -- the line itself stands in for the third mark -- and
% sets the label to the RIGHT of the line.
\newcommand{\FIGXVSEVEN}{\begin{bkfigure}[\linewidth]
\begin{tikzpicture}[scale=0.910]
  \XVaxes{-0.9}{4.3}{-1.7}{3.55}
  \XVxticks{1}{2}
  \draw[key] (3,-1.65) -- (3,3.48);
  \XVzero
  \node[vlab,right] at (3.08,1.2) {$x=3$};
\end{tikzpicture}
\figcap{15--7}
\end{bkfigure}}

% ------------------------------------------------------------ Fig 15-8, 15-9
% 15-8: the points of 2x - 3y + 6 = 0 found in Example 1, namely
%       (0,2), (3,4), (-3,0) -- collinear by construction.  Both axes are
%       ticked in the book.
% 15-9: the line y = -2 of Example 2, parallel to the x-axis, 2 units below.
%       The axis of ordinates stops on that line, as it does in the book.
\newcommand{\FIGXVEIGHTNINE}{\begin{bkfigure}[\linewidth]
\begin{minipage}{0.48\linewidth}\centering
\begin{tikzpicture}[scale=0.44]
  \XVaxes{-4.6}{6.4}{-2.6}{5.4}
  \XVxticks{-4}{5}
  \XVyticks{-2}{4}
  % the line 2x - 3y + 6 = 0, i.e. y = (2/3)x + 2 -- endpoints computed
  \draw[aux] (-4.2,-0.8) -- (5.4,5.6);
  \foreach \p in {(-3,0),(0,2),(3,4)} {\node[bkdot] at \p {};}
\end{tikzpicture}
\figcap{15--8}
\end{minipage}\hfill
\begin{minipage}{0.48\linewidth}\centering
\begin{tikzpicture}[scale=0.44]
  \XVaxes{-3.4}{5.3}{-2.05}{3.1}
  \draw[aux] (-3.6,-2) -- (5.4,-2);
  \draw[tick] (-0.22,-1) -- (0.22,-1);
  \XVzero
  \node[vlab,above=1pt] at (3.15,-2) {$(x,-2)$};
\end{tikzpicture}
\figcap{15--9}
\end{minipage}
\end{bkfigure}}

% ---------------------------------------------------------- Fig 15-10, 15-11
% 15-10: proof of Theorem 15-2, Case 3.  (0,-C/B), (x1,y1), (x2,y2) are
%        collinear; the two vertical legs are perpendicular to the base.
%        The base is taken at height 0 and the x-axis sits below it, at the
%        same fraction of x2 as in the book.
% 15-11: proof of Theorem 15-3, Case 3.  (x1,y1), (x2,y2), (x,y) collinear;
%        right angles at (x2,y1) and (x,y1).  In the book (x1,y1) lies just
%        to the LEFT of the axis of ordinates and the whole construction sits
%        above the axis of abscissas.
\newcommand{\FIGXVTENELEVEN}{\begin{bkfigure}[\linewidth]
\begin{minipage}{0.49\linewidth}\centering
\begin{tikzpicture}[scale=0.55]
  % base line y = -C/B taken at height 0; the x-axis sits below it.
  \def\bY{0}
  \coordinate (O)  at (0,\bY);          % (0, -C/B)
  \coordinate (P1) at (2.81,1.47525);   % (x1, y1)   slope 0.525
  \coordinate (P2) at (4.2,2.205);      % (x2, y2)   -- collinear with O,P1
  \coordinate (F1) at (2.81,\bY);
  \coordinate (F2) at (4.2,\bY);
  \draw[fig,-latex] (-0.52,-1.65) -- (4.54,-1.65);
  \node[vlab,right] at (4.54,-1.65) {$x$};
  \draw[fig,-latex] (0,-2.20) -- (0,2.49);
  \node[vlab,right] at (0,2.49) {$y$};
  \draw[fig] (O) -- (F2);
  \draw[fig] (P1) -- (F1);
  \draw[fig] (P2) -- (F2);
  \draw[aux] (O) -- (P2);
  \foreach \p in {O,P1,P2} {\dt{\p}}
  \node[slab,left]  at (O)  {$(0,-C/B)$};
  \node[slab,anchor=south east,xshift=-1pt,yshift=3pt] at (P1) {$(x_1,y_1)$};
  \node[slab,anchor=south east,xshift=9pt,yshift=3pt]  at (P2) {$(x_2,y_2)$};
  % one node holds both foot labels, as the book sets them: side by side on a
  % single line below the base.  Two separate nodes overlap at this trim.
  \node[slab,below] at (3.3,\bY) {$(x_1,-C/B)\quad(x_2,-C/B)$};
\end{tikzpicture}
\figcap{15--10}
\end{minipage}\hfill
\begin{minipage}{0.49\linewidth}\centering
\begin{tikzpicture}[scale=0.60]
  \coordinate (P1) at (-0.61,1.36);     % (x1, y1)
  \coordinate (P2) at (1.78,1.77825);   % (x2, y2)   slope 0.175
  \coordinate (P3) at (4.67,2.284);     % (x, y)   -- collinear with P1,P2
  \coordinate (F2) at (1.78,1.36);
  \coordinate (F3) at (4.67,1.36);
  \draw[fig,-latex] (-1.65,0) -- (4.93,0);
  \node[vlab,right] at (4.93,0) {$x$};
  \draw[fig,-latex] (0,-1.02) -- (0,3.44);
  \node[vlab,right] at (0,3.44) {$y$};
  \draw[fig] (P1) -- (F3);
  \draw[fig] (P2) -- (F2);
  \draw[fig] (P3) -- (F3);
  \draw[fig] (P1) -- (P3);
  \foreach \p in {P1,P2,P3} {\dt{\p}}
  \node[slab,left] at (P1) {$(x_1,y_1)$};
  \node[slab] at ($(P2)+(-0.34,0.48)$) {$(x_2,y_2)$};
  \node[slab] at ($(P3)+(0.10,0.50)$)  {$(x,y)$};
  \node[slab,below] at (F2) {$(x_2,y_1)$};
  \node[slab,below] at (F3) {$(x,y_1)$};
\end{tikzpicture}
\figcap{15--11}
\end{minipage}
\end{bkfigure}}

% --------------------------------------------------------------- Fig 15-12
% The line through (-2,3) and (1,5) of the worked Example.  Both axes ticked;
% the book sets (-2,3) above-left of its point and (1,5) below-right of its
% own, which also keeps both clear of the line.
\newcommand{\FIGXVTWELVE}{\begin{bkfigure}[\linewidth]
\begin{tikzpicture}[scale=0.676]
  \XVaxes{-5.0}{5.6}{-1.2}{6.9}
  \XVxticks{-4}{5}
  \XVyticks{1}{5}
  % slope 2/3 through (-2,3) and (1,5) -- endpoints computed, not eyeballed
  \draw[key] (-4.1,1.6) -- (3.45,6.63333);
  \dt{-2,3}
  \dt{1,5}
  \XVzero
  \node[vlab,above left=0pt and 0pt]  at (-2,3) {$(-2,3)$};
  \node[vlab,below right=0pt and 0pt] at (1,5)  {$(1,5)$};
\end{tikzpicture}
\figcap{15--12}
\end{bkfigure}}

% ---------------------------------------------------------- Fig 15-13, 15-14
% 15-13: distances along lines parallel to an axis.  AB = 1 - (-3) = 4;
%        CD = |x2 - x1|.  All four dashed pairs are axis-parallel.  C and D
%        straddle the axis of ordinates, at (-1,-1) and (2,-1), as in the
%        book, and each of the four lettered points carries its letter on one
%        side of the dot and its coordinate pair on the other.
% 15-14: the right triangle P1 P2 Q of the distance formula, right angle at Q.
\newcommand{\FIGXVTHIRTEENFOURTEEN}{\begin{bkfigure}[\linewidth]
\begin{minipage}{0.58\linewidth}\centering
\begin{tikzpicture}[scale=0.52]
  \draw[fig,-latex] (-2.8,0) -- (7.0,0);
  \node[vlab,right] at (7.0,0) {$x$};
  \draw[fig,-latex] (0,-3.6) -- (0,4.6);
  \node[vlab,right] at (0,4.6) {$y$};
  \node[slab,above left=-1pt and 0pt] at (0,0) {$0$};
  % horizontal pair (2,3)-(6,3)
  \draw[aux] (2,3) -- (6,3);
  \dt{2,3} \dt{6,3}
  \node[slab,above] at (2,3.1) {$(2,3)$};
  \node[slab,above] at (6,3.1) {$(6,3)$};
  % horizontal pair (-2,2)-(3,2)
  \draw[aux] (-2,2) -- (3,2);
  \dt{-2,2} \dt{3,2}
  \node[slab,left]  at (-2.1,2) {$(-2,2)$};
  \node[slab,right] at (3.1,2)  {$(3,2)$};
  % vertical pair A(1,1)-B(1,-3)
  \draw[aux] (1,1) -- (1,-3);
  \dt{1,1} \dt{1,-3}
  \node[slab,left]  at (0.88,1)  {$A$};
  \node[slab,right] at (1.12,1)  {$(1,1)$};
  \node[slab,left]  at (0.88,-3) {$B$};
  \node[slab,right] at (1.12,-3) {$(1,-3)$};
  % horizontal pair C(x1,y1)-D(x2,y1), straddling the axis of ordinates.
  % Dropped to y = -1.5: at this trim a label set above a point at y = -1
  % would run back over the axis of abscissas.
  \draw[aux] (-1,-1.5) -- (2,-1.5);
  \dt{-1,-1.5} \dt{2,-1.5}
  \node[slab,above] at (-1,-1.4) {$C$};
  \node[slab,above] at (2,-1.4)  {$D$};
  \node[slab,left]  at (-1.12,-1.5) {$(x_1,y_1)$};
  \node[slab,right] at (2.12,-1.5)  {$(x_2,y_1)$};
\end{tikzpicture}
\figcap{15--13}
\end{minipage}\hfill
\begin{minipage}{0.38\linewidth}\centering
\begin{tikzpicture}[scale=0.62]
  \draw[fig,-latex] (-0.95,0) -- (4.57,0);
  \node[vlab,right] at (4.57,0) {$x$};
  \draw[fig,-latex] (0,-0.26) -- (0,3.55);
  \node[vlab,right] at (0,3.55) {$y$};
  \node[slab,below right=-1pt and 1pt] at (0,0) {$0$};
  \coordinate (P1) at (0.62,1.15);
  \coordinate (P2) at (3.72,3.05);
  \coordinate (Q)  at (3.72,1.15);
  \draw[key] (P1) -- (P2);
  \draw[aux] (P1) -- (Q);
  \draw[aux] (Q)  -- (P2);
  \foreach \p in {P1,P2,Q} {\dt{\p}}
  \node[vlab,above,xshift=-2pt] at (P1) {$P_1$};
  \node[vlab,anchor=south east,xshift=-1pt] at (P2) {$P_2$};
  \node[vlab,right,yshift=5pt] at (Q) {$Q$};
  % nudged right so it clears the axis of ordinates, as the book's does
  \node[slab,below,xshift=8pt] at (P1) {$(x_1,y_1)$};
  % the book prints (x_2,y_1) at P_2 as well as at Q -- see ch15-UNCERTAIN.md
  \node[slab,anchor=south west,xshift=1pt] at (P2) {$(x_2,y_1)$};
  \node[slab,below=3pt] at (Q) {$(x_2,y_1)$};
\end{tikzpicture}
\figcap{15--14}
\end{minipage}
\end{bkfigure}}

% ---------------------------------------------------------- Fig 15-15, 15-16
% 15-15: M is the midpoint of P1P2; parallels to the axes through M bisect
%        the sides of the right triangle.
% 15-16: the worked example -- M(0,1) is the midpoint of (-1,3) and (1,-1).
\newcommand{\FIGXVFIFTEENSIXTEEN}{\begin{bkfigure}[\linewidth]
\begin{minipage}{0.52\linewidth}\centering
\begin{tikzpicture}[scale=0.756]
  \draw[fig,-latex] (-0.40,0) -- (3.60,0);
  \node[vlab,right] at (3.60,0) {$x$};
  \draw[fig,-latex] (0,-0.35) -- (0,3.90);
  \node[vlab,right] at (0,3.90) {$y$};
  \node[slab,below right=-1pt and 1pt] at (0,0) {$0$};
  \coordinate (P1) at (0.46,1.05);
  \coordinate (P2) at (3.36,2.52);
  \coordinate (M)  at ($(P1)!0.5!(P2)$);
  \coordinate (C)  at (3.36,1.05);
  \coordinate (Mh) at (3.36,1.785);    % foot of the parallel to the x-axis
  \coordinate (Mv) at (1.91,1.05);     % foot of the parallel to the y-axis
  \draw[key] (P1) -- (P2);
  \draw[aux] (P1) -- (C) -- (P2);
  \draw[aux] (M) -- (Mh);
  \draw[aux] (M) -- (Mv);
  \foreach \p in {P1,P2,M} {\dt{\p}}
  \node[slab,anchor=north west,xshift=-9pt,yshift=-3pt] at (P1) {$P_1\,(x_1,y_1)$};
  \node[slab,right] at (P2) {$P_2\,(x_2,y_2)$};
  \node[vlab,above left=-3pt and -1pt] at (M) {$M$};
  % Solidus form: a built-up fraction draws a rule inside the label box, which
  % the collision audit reads as geometry.  The solidus runs about two and a
  % half times as wide as the book's stacked fraction, so the label is lifted
  % clear of P_2 rather than tucked beside it.  See chapters/ch15-UNCERTAIN.md.
  \node[slab,anchor=south west] at (0.05,2.95)
    {$\left((x_1+x_2)/2,\ (y_1+y_2)/2\right)$};
\end{tikzpicture}
\figcap{15--15}
\end{minipage}\hfill
\begin{minipage}{0.44\linewidth}\centering
\begin{tikzpicture}[scale=0.651]
  \XVaxes{-2.0}{3.2}{-1.95}{3.5}
  \node[slab,below left=-1pt and 0pt] at (0,0) {$0$};
  \dt{-1,3}
  \dt{0,1}
  \dt{1,-1}
  \node[slab,left]  at (-1.15,3) {$P_1\,(-1,3)$};
  \node[slab,right] at (0.18,1)  {$M\,(0,1)$};
  \node[slab,right] at (1.18,-1) {$P_2\,(1,-1)$};
\end{tikzpicture}
\figcap{15--16}
\end{minipage}
\end{bkfigure}}

% --------------------------------------------------------------- Fig 15-17
% Exercise *7: the segment joining the midpoints of two sides is parallel to
% the third side and half as long.  Vertices (0,0), (a,0), (b,c); the book
% draws b barely larger than a, so the right side is all but vertical, and
% the midline is a plain dashed segment with no arrowhead.
\newcommand{\FIGXVSEVENTEEN}{\begin{bkfigure}[\linewidth]
\begin{tikzpicture}[scale=0.936]
  \draw[fig,-latex] (-0.70,0) -- (3.70,0);
  \node[vlab,right] at (3.70,0) {$x$};
  \draw[fig,-latex] (0,-0.60) -- (0,2.75);
  \node[vlab,right] at (0,2.75) {$y$};
  \coordinate (A) at (0,0);
  \coordinate (B) at (3.10,0);
  \coordinate (C) at (3.15,2.48);
  \draw[fig] (A) -- (B) -- (C) -- cycle;
  \coordinate (M1) at ($(A)!0.5!(C)$);
  \coordinate (M2) at ($(B)!0.5!(C)$);
  \draw[aux] (M1) -- (M2);
  \foreach \p in {M1,M2} {\dt{\p}}
  \node[slab,below right=-1pt and 0pt] at (A) {$(0,0)$};
  \node[slab,below] at (B) {$(a,0)$};
  \node[slab,right] at (C) {$(b,c)$};
\end{tikzpicture}
\figcap{15--17}
\end{bkfigure}}

% ---------------------------------------------------------- Fig 15-18, 15-19
% 15-18: the circle of Eq. (15-6): centre (h,k), radius r, a point (x,y) on it.
%        The book's centre sits just above the axis of abscissas and about
%        half a radius right of the axis of ordinates, so the circle crosses
%        both axes; the axis of ordinates stops well short of the bottom of
%        the circle.
% 15-19: the circle x^2+y^2+6x-8y+19 = 0, centre (-3,4), radius sqrt(6).
\newcommand{\FIGXVEIGHTEENNINETEEN}{\begin{bkfigure}[\linewidth]
\begin{minipage}{0.50\linewidth}\centering
\begin{tikzpicture}[scale=0.66]
  \draw[fig,-latex] (-1.82,0) -- (4.90,0);
  \node[vlab,right] at (4.90,0) {$x$};
  \draw[fig,-latex] (0,-0.77) -- (0,2.65);
  \node[vlab,right] at (0,2.65) {$y$};
  % the book sets 0 below left of the origin; the origin here lies inside the
  % circle, and the audit approximates each quarter-arc by its chord, so a
  % label dropped below left reads as touching that chord.  Set above left.
  \node[slab,above left=-1pt and 0pt] at (0,0) {$0$};
  \coordinate (K) at (1.04,0.38);
  \def\rr{1.85}
  \draw[key] (K) circle (\rr);
  % (x,y) placed on the circle at 148.5 deg from the centre -- exact by
  % construction, so the drawn radius really has length r.
  \coordinate (X) at ($(K)+(148.5:\rr)$);
  \draw[fig] (K) -- (X);
  \dt{K}
  \node[slab,right] at (K) {$(h,k)$};
  \node[slab,left=1pt]  at (X) {$(x,y)$};
  % r set off the radius along the perpendicular, so it does not sit on the line
  \node[slab] at ($(K)!0.5!(X)+(58.5:0.34)$) {$r$};
\end{tikzpicture}
\figcap{15--18}
\end{minipage}\hfill
\begin{minipage}{0.46\linewidth}\centering
\begin{tikzpicture}[scale=0.48]
  \draw[fig,-latex] (-6.0,0) -- (3.2,0);
  \node[vlab,right] at (3.2,0) {$x$};
  \draw[fig,-latex] (0,-1.0) -- (0,7.0);
  \node[vlab,right] at (0,7.0) {$y$};
  \node[slab,below right=-1pt and 1pt] at (0,0) {$0$};
  \XVxticks{-5}{-1}
  \XVyticks{1}{5}
  \draw[key] (-3,4) circle (2.4494897);
  \dt{-3,4}
  \node[slab,below=3pt] at (-3,4) {$(-3,4)$};
\end{tikzpicture}
\figcap{15--19}
\end{minipage}
\end{bkfigure}}

% figures17.tex -- Appendix figures redrawn in TikZ.
%
% The Appendix (book pp.275-282) contains exactly ONE figure: the unnumbered
% and uncaptioned diagram on book p.277 that accompanies Postulate 16. It
% shows two three-ray configurations, unprimed and primed, drawn congruent.
% (Pages 275, 276, 278-281 are text only; 282 is blank. Verified page by page
% against sources/Geometry.pdf pp.289-296.)
%
% GEOMETRY MEASURED off sources/Geometry.pdf p.291 rasterised at 600 dpi.
% Method: threshold the ink, separate the heavy rays from the light arcs by
% binary erosion (rays survive a 5x5 erosion, arcs do not), take the vertex as
% the left end of the horizontal ray, then fit ray directions from the eroded
% ray cores only. Reading the whole ink mask instead biases every direction
% toward the arcs that run tangent to it -- that is what produced the earlier
% 45.0/117.0 figures. Raw and de-skewed (photo tilt removed by subtracting the
% measured tilt of the horizontal ray) results:
%
%     ray          raw unprimed  raw primed   de-skewed mean
%     horizontal       +0.48        +0.67          0 (datum)
%     diagonal         44.38        43.87        43.55
%     upper-left      116.33       116.13       115.66
%
% Drawn at 0 / 43.6 / 115.7 deg. Ray lengths and arc radii as fractions of the
% horizontal ray (mean of the two sub-figures):
%
%     horizontal ray 1.000   arc A 0.446   (drawn 3.30 / 1.47)
%     diagonal ray   0.994   arc C 0.699   (drawn 3.28 / 2.31)
%     upper-left ray 0.939   arc B 0.774   (drawn 3.10 / 2.55)
%
% The three arc radii are DISTINCT and ordered A < C < B, and each label sits
% on its own arc in a gap in that arc. Arc spans read off the source by
% scanning each arc's circle for ink (rays and label glyphs masked out):
%
%     arc A   4..60  and  78..114   gap 60..78, label at  69 deg
%     arc B  47..77  and  87..114   gap 77..87, label at  81.5 deg
%     arc C   4..16  and  30..42    gap 16..30, label at  23 deg
%
% The arrowhead ends below are the source's exactly; the GAP ends are opened a
% few degrees wider (A 52..86, B 71.5..91.5, C 12..34) because our labels are
% set with more padding than the book's tight hand lettering -- at the source's
% own gap widths check_labels.py reports the glyph boxes touching the arcs.
%
% Arrowheads therefore land at 114 (upper-left ray), 47/42 (diagonal ray) and
% 4 (horizontal ray) -- always just short of the ray they point at, and the two
% hits on the horizontal ray (arc A at 1.47, arc C at 2.31) and the two on the
% diagonal (arc B at 2.55, arc C at 2.31) are at different radii, as in the book.
%
% LINE WEIGHT: the book draws the three rays roughly 1.75x the weight of the
% annotation arcs (measured 13.0 px vs 8.5 px of ink at 600 dpi on the same
% photo). The rays keep the house `given' weight so the Appendix matches the
% rest of the book; the arcs are dropped to 0.42pt to restore the contrast.
%
% What the postulate rests on: the diagonal ray is INTERIOR to angle A, so
% angle A = angle B + angle C. See tools/constraints/ch17.py.

% ------------------------------------------------ Appendix figure (p.277)
\newcommand{\FIGAPPONE}{\begin{bkfigure}[0.52\linewidth]
% --- unprimed configuration
\begin{tikzpicture}[scale=0.72,
    mk/.style={fig, line width=0.42pt},          % light annotation arcs
    alab/.style={vlab, outer sep=1.5pt}]         % label sits in its arc's gap
  \coordinate (O) at (0,0);
  \draw[given] (O) -- (0:3.30);        % horizontal ray
  \draw[given] (O) -- (43.6:3.28);     % diagonal ray, interior to angle A
  \draw[given] (O) -- (115.7:3.10);    % upper-left ray
  % angle A : upper-left ray <-> horizontal ray  (the whole angle)
  \draw[mk,->] (86:1.47) arc[start angle=86, end angle=114, radius=1.47];
  \draw[mk,->] (52:1.47) arc[start angle=52, end angle=4,   radius=1.47];
  \node[alab] at (69:1.47) {$A$};
  % angle B : upper-left ray <-> diagonal ray
  \draw[mk,->] (91.5:2.55) arc[start angle=91.5, end angle=114, radius=2.55];
  \draw[mk,->] (71.5:2.55) arc[start angle=71.5, end angle=47,  radius=2.55];
  \node[alab] at (81.5:2.55) {$B$};
  % angle C : diagonal ray <-> horizontal ray
  \draw[mk,->] (34:2.31) arc[start angle=34, end angle=42, radius=2.31];
  \draw[mk,->] (12:2.31) arc[start angle=12, end angle=4,  radius=2.31];
  \node[alab] at (23:2.31) {$C$};
\end{tikzpicture}

\vspace{4mm}

% --- primed configuration (drawn congruent to the unprimed one, as the book
%     draws it; the postulate itself only supposes congruence of two pairs)
\begin{tikzpicture}[scale=0.72,
    mk/.style={fig, line width=0.42pt},
    alab/.style={vlab, outer sep=1.5pt}]
  \coordinate (O) at (0,0);
  \draw[given] (O) -- (0:3.30);
  \draw[given] (O) -- (43.6:3.28);
  \draw[given] (O) -- (115.7:3.10);
  \draw[mk,->] (86:1.47) arc[start angle=86, end angle=114, radius=1.47];
  \draw[mk,->] (52:1.47) arc[start angle=52, end angle=4,   radius=1.47];
  \node[alab] at (69:1.47) {$A'$};
  \draw[mk,->] (91.5:2.55) arc[start angle=91.5, end angle=114, radius=2.55];
  \draw[mk,->] (71.5:2.55) arc[start angle=71.5, end angle=47,  radius=2.55];
  \node[alab] at (81.5:2.55) {$B'$};
  \draw[mk,->] (34:2.31) arc[start angle=34, end angle=42, radius=2.31];
  \draw[mk,->] (12:2.31) arc[start angle=12, end angle=4,  radius=2.31];
  \node[alab] at (23:2.31) {$C'$};
\end{tikzpicture}
\end{bkfigure}}

% ================= local macros from frontmatter =================
% Halftone photographs in the original are reproduced as PLATES: a framed
% placeholder carrying the original caption. Never redrawn.
%   #1 = placeholder height, #2 = caption block (already formatted)
\newcommand{\FMplate}[2]{%
  \par\medskip
  \begin{center}
  \fbox{\begin{minipage}[c][#1][c]{0.68\linewidth}\centering\footnotesize\itshape
    halftone photograph\\ not reproduced in this edition
  \end{minipage}}%
  \par\smallskip
  {\footnotesize #2}
  \end{center}
  \par\medskip}

% Title-page author line: name in letterspaced caps, affiliation in italic.
\newcommand{\FMauthor}[2]{%
  \par\medskip
  {\large\scshape #1}\par
  {\footnotesize\itshape #2}\par}
% Narrow measure (145 mm trim) + Palatino: give TeX a little slack so the
% justified prose paragraphs set without overfull lines.
\emergencystretch=1.5em

% Blank versos forced by \chapter* (openright) must carry no folio or head.
\makeatletter
\def\cleardoublepage{\clearpage
  \if@twoside\ifodd\c@page\else
    \null\thispagestyle{empty}\newpage
    \if@twocolumn\hbox{}\newpage\fi
  \fi\fi}
\makeatother
% ================= local macros from ch01 =================
% Halftone plates: the book prints full-page photographs facing several pages of
% this chapter. They are reproduced as framed placeholders carrying the original
% caption; see ch01-UNCERTAIN.md.
\newcommand{\Iplate}[2]{%
  \par\medskip
  \begin{center}
  \fbox{\begin{minipage}[c][0.42\textheight][c]{0.78\linewidth}\centering
    \footnotesize\itshape [halftone photograph --- plate not reproduced]
    \par\vspace{1.2em}
    \normalfont\textsc{#1}\par
    \smallskip
    {\footnotesize #2}
  \end{minipage}}
  \end{center}
  \par\medskip}
% ================= local macros from ch03 =================
% Unnumbered star footnote, same reasoning as ch02's \IIstarnote: \starnote in
% brumfiel.sty emits a bare \footnotetext, which prints the footnote counter;
% the book marks all three of this chapter's notes with "*" and no number.
\newcommand{\IIIstarfn}[1]{\begingroup\renewcommand{\thefootnote}{}%
  \footnotetext{*#1}\endgroup}
% Starred theorems. The book prints "*Theorem 3-4." for the optional theorems
% of Section 3-6; this shares the `theorem' counter so numbering stays exact.
\theoremstyle{bkplain}
\newtheorem{IIIsthm}[theorem]{*Theorem}
% ================= local macros from ch04 =================
% The book's run-in "Problem." paragraphs.
\newenvironment{IVproblem}{\par\medskip\noindent\emph{Problem.}\ }{\par\medskip}
% Centred display headings ("REVIEW OF CHAPTER 4", "ALGEBRA REVIEW").
\newcommand{\IVhead}[1]{%
  \par\bigskip\begin{center}\bfseries\large\MakeUppercase{#1}\end{center}%
  \par\nopagebreak\medskip}
% ================= local macros from ch05 =================
% Centred display headings ("REVIEW OF CHAPTER 5", "ALGEBRA REVIEW").
\newcommand{\Vhead}[1]{%
  \par\bigskip\begin{center}\bfseries\large\MakeUppercase{#1}\end{center}%
  \par\nopagebreak\medskip}
% Unnumbered star footnote, same reasoning as ch02's \IIstarnote: \footnote
% would set a numbered mark, but the book prints a bare "*". Used for the
% "irrational" footnote on book p.74.  NOTE: this chapter does NOT carry the
% "*Starred ... are optional" note -- book p.75, which holds this chapter's
% first starred exercises (Ex. Group 5-3, nos. 4 and 7), has no footnote at
% all; ch02 owns the book-wide first use.  Verified against PDF p.89.
\newcommand{\Vstarfoot}[1]{\begingroup\renewcommand{\thefootnote}{}%
  \footnotetext{*\,#1}\endgroup}
% Fig. 5-1 is a pictorial vignette, not a geometric diagram -> placeholder,
% per the plate rule (same treatment as Fig. 4-1). See ch05-UNCERTAIN.md.
\newcommand{\FIGVONE}{\begin{bkfigure}[0.46\linewidth]
\fbox{\begin{minipage}{0.92\linewidth}\centering\footnotesize
\vspace{1.6em}
\textit{[Illustration: a man striding along the\\
hatched ground beside a building, pacing\\
off its length.]}
\vspace{1.6em}
\end{minipage}}
\figcap{5--1}
\end{bkfigure}}
% ================= local macros from ch06 =================
% Unnumbered star footnote, same reasoning as ch02/ch09: \starnote in
% brumfiel.sty emits a bare \footnotetext, which prints the footnote counter;
% the book marks these notes with the symbol itself and no number.
\newcommand{\VIstarnote}[1]{\begingroup\renewcommand{\thefootnote}{}%
  \footnotetext{#1}\endgroup}
% The book's stacked "Given: / Prove:" exercise layout (Exercise Groups 6--5,
% 6--11, 6--14 and the review exercises).  Each argument may hold its own
% \\-separated lines.
% The leading break must NOT be taken when \VIgp is the first thing in an
% \item.  LaTeX's \@endpe machinery (armed by any list/center environment that
% just ended -- e.g. the centred display closing Ex. 6-6 #3) hands the NEXT
% paragraph an \everypar that removes its last box, and that box is the item
% label.  The label is then eaten and re-emitted on the following item, which
% printed "4" and "5" on top of each other in Exercise Group 6-6.  Entering
% horizontal mode with \leavevmode when we are still in vertical mode lays the
% label down first and puts "Given:" on the same line, which is how every
% correctly-set item in this chapter already reads.
\newcommand{\VIgp}[2]{%
  \ifvmode\leavevmode\else\par\smallskip\noindent\fi
  \begin{tabular}[t]{@{}l@{\hspace{0.6em}}l@{}}
    Given: & \begin{tabular}[t]{@{}l@{}}#1\end{tabular}\\[2pt]
    Prove: & \begin{tabular}[t]{@{}l@{}}#2\end{tabular}
  \end{tabular}\par\smallskip}
% ================= local macros from ch07 =================
% Numbered remark head. Sections 7-4 (pp.122) and 7-5 (pp.123-124) print
% "Remark 1.", "Remark 2." ... as running heads; brumfiel.sty's {remark}
% emits a bare "Remark." and would drop the book's numbering.
% Same shape as {remark}, with the number restored.
\newenvironment{VIIremarkn}[1]%
  {\par\medskip\noindent\emph{Remark #1.}\ }{\par\medskip}
% No \VIIstarnote: ch07 has starred exercises (7-2 #3, 7-4 #4-7, 7-7 #3,
% 7-8 #8-9, Algebra Review #7-9) but the book prints NO star footnote on any
% of pp.116-132 -- verified against the source pages, incl. the foot of
% p.118 where the chapter's first "*3." falls. ch02 owns the book-wide first
% use of the convention note (\IIstarnote); ch05 and ch10 omit it for the
% same reason.
% ================= local macros from ch08 =================
% Plate placeholder: a halftone photograph in the original, never redrawn.
\newcommand{\VIIIplate}[2]{%
  \begin{bkfigure}[0.78\linewidth]
  \setlength{\fboxsep}{10pt}%
  \fbox{\begin{minipage}{0.80\linewidth}\centering\footnotesize
    \textsc{plate}\par\medskip
    #1\par\medskip
    \emph{Halftone photograph in the original; not redrawn.}
  \end{minipage}}
  \figcap{#2}
  \end{bkfigure}}
% Two unnumbered "*" footnotes belonging to this chapter (book pp.135 and 139).
% \starnote in brumfiel.sty emits a bare \footnotetext, which prints the
% footnote counter; the book marks both of these with "*" and no number, so
% they get the same counter-suppressing treatment ch10 gives \Xlimitnote.
\newcommand{\VIIIproofnote}{\begingroup\renewcommand{\thefootnote}{}%
  \footnotetext{*\,If the teacher wishes, the second and third proofs may be
  omitted. Note that the second proof employs a lemma, or helping
  theorem.}\endgroup}
\newcommand{\VIIIgaussnote}{\begingroup\renewcommand{\thefootnote}{}%
  \footnotetext{*\,See, for example, the discussion of Gauss' experiment in
  \emph{The World of Mathematics}, by J.~R. Newman, Vol.~3, p.~1644.}\endgroup}
% The three-utilities puzzle on book p.145 is set as bare letters over dots in
% the running text, not as a numbered figure.
\newcommand{\VIIIutilities}{%
  \par\medskip\begin{center}
  \begin{tabular}{@{}c@{\hspace{3.2em}}c@{\hspace{3.2em}}c@{}}
    $A$ & $B$ & $C$ \\[-1pt]
    $\cdot$ & $\cdot$ & $\cdot$ \\[4pt]
    $G$ & $W$ & $E$ \\[-1pt]
    $\cdot$ & $\cdot$ & $\cdot$ \\
  \end{tabular}
  \end{center}\par\medskip}
% Narrow multicol columns plus this chapter's primed/subscripted math atoms
% ($\ang{A'O'B'}$, $l_1 \pll l_2$) leave TeX few legal break points.
\setlength{\emergencystretch}{2.5em}
% ================= local macros from ch09 =================
% Unnumbered star footnote, same reasoning as ch06's \VIstarnote: a bare
% \footnotetext prints the footnote counter, but the book marks its footnotes
% with "*" and no number.  Used for the "mean proportional" note on book p.171
% (its in-text mark is \textsuperscript{*}, per ch14's convention).
%
% NOTE: no "Starred exercises ... are optional" note here.  Ch 9 has starred
% exercises (Groups 9--3, 9--4, 9--5, 9--10, 9--13 and Review Exs. 3, 4, 15),
% but the book prints that convention note only once, at its first book-wide
% occurrence, which ch02 owns; pp.158--159 carry no footnote at all.  Same
% call as ch05 and ch10.
\newcommand{\IXstarnote}[1]{\begingroup\renewcommand{\thefootnote}{}%
  \footnotetext{#1}%
  \endgroup}
% ================= local macros from ch10 =================
% (No \Xstarnote here: ch10 has starred exercises but the book prints no star
% footnote on any of its pages; ch02:311 owns the book-wide first use.)
% Unnumbered "limit" footnote on book p.182 — same reasoning as ch02's
% \IIstarnote: \starnote in brumfiel.sty emits a bare \footnotetext, which
% prints the footnote counter, and the book marks this note with "*" and no
% number.
\newcommand{\Xlimitnote}{\begingroup\renewcommand{\thefootnote}{}%
  \footnotetext{*The word ``limit'' has a precise mathematical sense, which
  need not concern us now. Roughly, the idea is that the areas of the
  approximating rectangles get close to a number that we call the \emph{area}
  of the given rectangle.}\endgroup}
% Figure-name glyphs. The book writes "area (<glyph>ABCD)". In the running
% proofs it uses a plain square for both rectangles and parallelograms; the
% Review of Chapter 10 summary distinguishes four glyphs, so the parallelogram
% and trapezoid are drawn here rather than faked with \square.
\newcommand{\Xrect}{\ensuremath{\square}}
\newcommand{\Xpgram}{\mathord{\tikz[baseline=-0.05ex,line width=0.5pt]{%
  \draw (0,0) -- (1.35ex,0) -- (1.75ex,1.05ex) -- (0.40ex,1.05ex) -- cycle;}}}
\newcommand{\Xtrap}{\mathord{\tikz[baseline=-0.05ex,line width=0.5pt]{%
  \draw (0,0) -- (1.85ex,0) -- (1.45ex,1.05ex) -- (0.42ex,1.05ex) -- cycle;}}}
% Narrow multicol columns plus this chapter's long primed/subscripted math
% atoms ($A''B''C''D''$, $\sg{A_1B_1}/\sg{A_2B_2}$) leave TeX no legal break
% point. Give it a little slack rather than re-wording the book.
\setlength{\emergencystretch}{2.5em}
% ================= local macros from ch11 =================
% Theorem head carrying an unparenthesised note.  The book prints
% "Theorem 11-4 Converse to Theorem 11-3." with no parentheses, but
% brumfiel.sty's bkplain style wraps \thmnote in parentheses.  Shares the
% `theorem' counter so the numbering stays in one sequence.
\newtheoremstyle{XIplainnote}{1.1\topsep}{1.1\topsep}{\itshape}{}%
  {\bfseries}{.}{ }{\thmname{#1}\thmnumber{ #2}\thmnote{ \normalfont#3}}
\theoremstyle{XIplainnote}
\newtheorem{XIthmnote}[theorem]{Theorem}
\theoremstyle{bkplain}
% (No star-convention footnote here: ch11 has starred exercises -- Ex. Gp.
% 11-3 #2, 11-8 ##8-10, 11-12 ##2-3, 11-13 #13 -- but the book prints no such
% footnote on any page of the chapter, book pp. 196, 200, 205, 207.
% ch02:311 owns the book-wide first use; same call as ch10.)
% Numbered remarks. The book sets italic "Remark" + roman number with NO
% period after the word ("Remark 1."), unlike the shared \begin{remark} in
% brumfiel.sty, which prints "Remark." Verified against book pp. 196, 205.
\newenvironment{XIremarkn}[1]{\par\medskip\noindent\emph{Remark} #1.\ }%
  {\par\medskip}
% Plural run-in label, as the book sets the note under Theorem 11-12
% (book p. 201) -- same case as ch14's \XIVremarks.
\newenvironment{XIremarks}{\par\medskip\noindent\emph{Remarks.}\ }%
  {\par\medskip}
% Footnotes marked with the book's symbols (* then \dag) rather than digits;
% both fall on book p. 204.  Same device as ch12's \XIIsymfoot.
\newcommand{\XIsymfoot}[1]{%
  \begingroup\renewcommand{\thefootnote}{\fnsymbol{footnote}}%
  \footnote{#1}\endgroup}
% The book sets "n >= 3" with the double-bar relation, not \geq.
\newcommand{\XIgeqq}{\geqq}
% ================= local macros from ch12 =================
% Arc accent, as the book sets it: a frown over the endpoint letters.
\newcommand{\XIIarc}[1]{\overset{\raisebox{-0.25ex}{\scalebox{1.15}[0.85]{$\frown$}}}{#1}}
% Equation numbers as (12--1), matching the book's en-dash style.
\renewcommand{\theequation}{\thechapter--\arabic{equation}}
% Numbered remarks. The book sets these as italic "Remark" + roman numeral,
% with NO period after the word (e.g. "Remark 1."), unlike the unnumbered
% \begin{remark} in brumfiel.sty, which prints "Remark." Verified against
% book pp. 218, 221.
\newenvironment{XIIremarkn}[1]{\par\medskip\noindent\emph{Remark} #1.\ }%
  {\par\medskip}
% Footnotes marked with * rather than a digit (the book's convention).
\newcommand{\XIIsymfoot}[1]{%
  \begingroup
  \renewcommand{\thefootnote}{\fnsymbol{footnote}}%
  \setcounter{footnote}{0}%
  \footnote{#1}%
  \endgroup}
% ================= local macros from ch13 =================
% The "*" on _locus_ on p.227 is a GLOSSARY footnote, not the starred-exercise
% convention, so \starnote is wrong here.  As with ch02's \IIstarnote and
% ch09's \IXstarnote, \footnotetext alone would print the footnote counter;
% the book marks this note with the "*" itself and no number.
\newcommand{\XIIIlocusnote}{\begingroup\renewcommand{\thefootnote}{}%
  \footnotetext{*\emph{Locus} is a Latin word meaning, literally, \emph{place}
  in the sense that we use the word \emph{location}.}%
  \endgroup}
% ================= local macros from ch14 =================
% Centred, letter-spaced part heading ("A. LINES, PLANES, AND ...").
\newcommand{\XIVpart}[1]{%
  \par\addvspace{1.5\baselineskip}\nopagebreak
  \begin{center}\bfseries\spaceskip=1.05em\relax #1\end{center}%
  \nopagebreak\addvspace{0.8\baselineskip}}
% Unnumbered footnote (the chapter-title and "sum" notes carry a bare *).
\newcommand{\XIVnote}[1]{{\let\thefootnote\relax\footnotetext{#1}}}
% The book's run-in label on book p. 242 is PLURAL ("Remarks."), whereas the
% shared \begin{remark} in brumfiel.sty prints the singular "Remark."
% Verified against book p. 242 (PDF 256, scan15 p. 10).
\newenvironment{XIVremarks}{\par\medskip\noindent\emph{Remarks.}\ }%
  {\par\medskip}
% Postulate Group VII. The book sets a bold run-in label, roman body -- the
% same shape as the bkdefn theorem style but hand-numbered (VII-1 ... VII-4).
\newenvironment{postulate}[1]%
  {\par\addvspace{1.1\topsep}\noindent\textbf{Postulate #1.}\ \ignorespaces}%
  {\par\addvspace{1.1\topsep}}
% ================= local macros from ch15 =================
% Starred (optional) theorems, as the book marks Theorems 15-2 and 15-3.
% Shares the theorem counter so the numbering runs 15-1, 15-2, 15-3.
\theoremstyle{bkplain}
\newtheorem{XVstheorem}[theorem]{*Theorem}
% Worked examples, set as the book sets them: italic lead-in, no number when
% a section carries only one.
\newenvironment{XVexample}[1][]{\par\medskip\noindent\emph{Example#1.}\ }{\par\medskip}
% NUMBERED remarks. brumfiel.sty's `remark' hard-codes an unnumbered
% "Remark."; book p.261 prints "Remark 1." and "Remark 2." in sequence, so the
% numbered form needs its own environment. The unnumbered `remark' is still
% used for the two lone remarks on pp.257 and 262.
\newenvironment{XVremark}[1]{\par\medskip\noindent\emph{Remark #1.}\ }{\par\medskip}
% Symbol-marked footnotes. The book marks this chapter's footnotes with * and
% dagger, never with numbers; \footnotetext would print the footnote counter,
% so the counter is blanked and the symbol is typed into the body text.
% (Same device as ch02's \IIstarnote and ch09's \IXstarnote.)
\newcommand{\XVnote}[1]{\begingroup\renewcommand{\thefootnote}{}%
  \footnotetext{#1}\endgroup}
% Halftone photographic plate: framed placeholder, never redrawn (B&W pass).
% Same shape as ch01's \Iplate and ch02's \IIplate.
\newcommand{\XVplate}[2]{%
  \par\medskip
  \begin{center}
  \fbox{\begin{minipage}{0.72\linewidth}\centering\vspace{2.4cm}
    {\footnotesize\itshape[halftone photographic plate --- not redrawn]}
    \vspace{2.4cm}\end{minipage}}%
  \par\smallskip
  {\footnotesize\textsc{#1}\\ #2}
  \end{center}
  \par\medskip}
% (No star-convention footnote here: ch15 has starred exercises and two starred
% theorems, but none of book pp.251-265 carries the "*Starred exercises ... are
% optional" note; ch02:311 owns the book-wide first use.)
% ================= local macros from ch16 =================
% Unnumbered star footnote, as the book sets its five notes in this chapter.
% brumfiel.sty's \starnote emits a bare \footnotetext, which prints the
% footnote counter ("0"); this chapter's notes carry a bare "*" and no number.
\newcommand{\XVIstarnote}[1]{\begingroup\renewcommand{\thefootnote}{}%
  \footnotetext{*#1}\endgroup}
% In-text marker for those notes (a raised asterisk, as printed).
\newcommand{\XVIstar}{\textsuperscript{*}}
% ================= local macros from appendix =================
% The Appendix has no numbered "N--M" sections. Its three major heads are
% set centred, bold and small, not flush left like a chapter's \section*.
% \APPhead also files the TOC entry the assembled book needs.
\newcommand{\APPhead}[2]{%
  \par\bigskip\begin{center}\bfseries #1\end{center}%
  \addcontentsline{toc}{section}{#2}\par\nopagebreak\medskip}
% Bold centred subhead grouping a run of postulates.
\newcommand{\APPsub}[1]{%
  \par\medskip\begin{center}\bfseries\small #1\end{center}\par\nopagebreak}
% Continuously-numbered reference list, with (a)/(b)/(c) subparts.
\newlist{applist}{enumerate}{2}
\setlist[applist,1]{label=\arabic*., leftmargin=1.7em, labelsep=0.4em,
                    itemsep=2.5pt, topsep=4pt, parsep=0pt}
\setlist[applist,2]{label=(\alph*), leftmargin=1.9em, labelsep=0.4em,
                    itemsep=1.5pt, topsep=2.5pt, parsep=0pt}
% short pages end short instead of stretching inter-para glue
\raggedbottom
% several units relied on standalone-body emergencystretch; the
% narrow 145mm measure needs it book-wide (appendix flag, 2026-08-17)
\emergencystretch=2em
\begin{document}
\pagenumbering{roman}

%========================================================= half title
\thispagestyle{empty}
\begin{center}
  \vspace*{0.28\textheight}
  {\Huge GEOMETRY}
\end{center}
\clearpage

%========================================================= p. i  series page
\thispagestyle{empty}
\null\vfill
\begin{center}
  {\footnotesize This book is in the}\par\smallskip
  {\normalsize ADDISON-WESLEY SCIENCE EDUCATION SERIES}\par
  \rule{0.28\linewidth}{0.4pt}\par\smallskip
  {\footnotesize\itshape Consulting Editors}\par\smallskip
  {\footnotesize
   \begin{tabular}{@{}ll@{}}
   \textsc{Richard S. Pieters} \quad & \textsc{Paul Rosenbloom}\\
   \textsc{George B. Thomas, Jr.} \quad & \textsc{John Wagner}
   \end{tabular}}
\end{center}
\vfill
\clearpage

%========================================================= p. ii  frontispiece
% Euclid frontispiece: engraved/halftone portrait plate facing the title page.
\thispagestyle{empty}
\null\vspace*{0.06\textheight}
\FMplate{52mm}{%
  \textsc{Euclid}\\
  Courtesy of\\
  \textit{Scripta Mathematica}}
\vfill
\clearpage

%========================================================= p. iii  title page
\thispagestyle{empty}
\null\vspace*{0.10\textheight}
\begin{center}
  {\Huge GEOMETRY}\par
  \vspace{2.2em}
  \FMauthor{Charles F. Brumfiel}{Ball State Teachers College}
  \FMauthor{Robert E. Eicholz}{Ball State Teachers College, Burris Laboratory School}
  \FMauthor{Merrill E. Shanks}{Purdue University}
  \vspace{3.2em}

  {\normalsize ADDISON-WESLEY PUBLISHING COMPANY, INC.}\par
  {\footnotesize READING, MASSACHUSETTS, U.S.A.}\par
  {\footnotesize LONDON, ENGLAND}
\end{center}
\vfill
\clearpage

%========================================================= p. iv  copyright
% Original copyright statement, transcribed as printed in the 1961 second
% printing. Reproduced as a historical record of the source edition.
\thispagestyle{empty}
\null\vspace*{0.22\textheight}
\begin{center}
  {\itshape Copyright \textcopyright\ 1960}\par
  {\normalsize ADDISON-WESLEY PUBLISHING COMPANY, INC.}\par
  \rule{0.30\linewidth}{0.4pt}\par
  {\itshape Printed in the United States of America}\par
  \vspace{1.1em}
  {\footnotesize\textsc{All rights reserved. This book, or parts thereof,
   may not be reproduced in any form without written permission of the
   publisher.}}\par
  \vspace{0.9em}
  {\itshape Library of Congress Catalog Card No. 60-8336}\par
  \vspace{1.4em}
  {\normalsize Second printing, February, 1961}
\end{center}
\vfill
\clearpage

%========================================================= contents
% The printed book's CONTENTS list (book pp. v--viii) is deliberately NOT
% transcribed: the assembled edition generates its own table of contents from
% the \addcontentsline entries in each unit. The integration agent places
\tableofcontents

%========================================================= PREFACE  (p. ix)
\chapter*{\centering\Large PREFACE}
\addcontentsline{toc}{chapter}{Preface}
\markboth{PREFACE}{PREFACE}

It is a common belief that plane geometry was completed by Euclid 2000 years ago
and that nothing has been added to it or taken from it since. This is simply not
true. That there are logical gaps in Euclid's presentation has been known for a
long time. Means to remedy these deficiencies have been known for about sixty
years, but strangely enough a mathematically adequate and yet elementary
treatment of plane geometry in the spirit of Euclid has not appeared in print.
This text represents an earnest effort to do just this. The current interest in
improvement of the secondary school curriculum makes this an appropriate time
for such an attempt.

This book, in earlier unpublished versions, has been taught to \emph{normal}
classes in \emph{average} high schools for five years. Many of the classes were
taught by one of us in the Burris School, the laboratory school of Ball State
Teachers College, but the material has been taught in other schools by teachers
who had had no previous special training in geometry. Our experience is that
both student and teacher find this treatment of geometry stimulating. The good
students become excited about mathematics and the students of average ability
learn the basic geometric facts.

Besides its mathematical completeness, an attribute that makes for clarity and
ease of understanding, there are several features which we think important.

There is a full chapter on logic early in the text. It pays the teacher to
return to this logic again and again in order to clarify the principles
underlying deductive proof. For the nonscience student, this chapter on logic
may be the item of greatest permanent value.

The postulates are presented in groups, as needed, with a resulting gain in
clarity and simplicity. Although it is true that at the start most students are
hardly ``provers,'' by the end of the year a considerable part of the class will
exhibit real mathematical perception. In the early stages the small number of
postulates makes proofs easier to construct.

Definitions are given explicitly; that is, they are formally labeled as such. It
is one of our prime objectives to get the student to see that precise
definitions are necessary for proofs.

There is plenty of material for a full year's work for all kinds of classes.
More difficult sections, theorems, or exercises are marked with an asterisk, and
some of the more difficult theorems may be treated as postulates. The chapters
on analytic geometry, locus, solid geometry, and philosophy of mathematics
provide supplementary material beyond the standard course. The chapter on solid
geometry alone could occupy a good portion of a semester if desired.

There is a review of important concepts at the end of each chapter. In addition,
there are algebra review sections, designed not only to refresh skills but also
to bring out the logical structure of algebra.

The Appendix lists all the postulates, plus the principal definitions and
theorems. Besides being a convenient reference, the Appendix is worth studying
during the year, for it gives, in brief, the step-by-step construction of the
full geometry.

We believe that the text raises students' reading ability. There is little
verbiage. Each section has a definite purpose, and the text is directed to this
purpose. Good classroom discussion, with careful attention to the use of words,
will add much to each student's ability to express himself clearly.

A separate Teachers' Manual gives greater detail on selection of topics,
emphasis, and time schedule. There is also an Answer Book for the exercises.

It is a great pleasure to thank all those who have helped us with encouragement
or constructive criticism. Special thanks are due to Dr.\ Curtis Howd, Principal
of the Burris Laboratory School, who encouraged us and made possible the
experimental sections; to Professor Jack Forbes, who observed many of the
classes in various schools and who insisted that the mathematics be correct; to
teachers at the Burris School, and many others, who by teaching the material
have improved the exposition; to The National Science Foundation, which, through
a grant for in-service institutes at Ball State, enabled us to reach a wider
audience and see the text taught by many teachers in a variety of school
situations; and finally, to our wives, who have waited impatiently for the
publication day.

\vspace{1.2em}
\noindent
\begin{minipage}[t]{0.45\linewidth}
\itshape January 1960
\end{minipage}\hfill
\begin{minipage}[t]{0.45\linewidth}
\raggedleft C.\,B., R.\,E., M.\,S.
\end{minipage}

%========================================================= FOREWORD  (p. xi)
\chapter*{\centering\Large FOREWORD TO THE STUDENT}
\addcontentsline{toc}{chapter}{Foreword to the Student}
\markboth{FOREWORD TO THE STUDENT}{FOREWORD TO THE STUDENT}

Mathematics is much more than a collection of interesting facts and useful tools
for solving special problems. Mathematics is primarily concerned with the
\emph{development} of important ideas. Frequently mathematicians develop an
interesting mathematical theory that is apparently useless so far as the
solution of practical problems is concerned. Later someone finds that this
mathematics fits very nicely into the pattern of some complicated practical
problem. The mathematical theory then becomes a useful tool.

As a branch of mathematics, geometry is first of all concerned with the
development of geometrical theory. As you study geometry you will see again and
again how the geometrical concepts fit into patterns in the real world and
become helpful in the solution of practical problems.

Mathematics requires thought, and thinking deserves your best effort. Do not
hesitate to ask questions.

\vspace{1.0em}
\noindent
\begin{flushright}
Good luck,\par
\textit{The Authors}
\end{flushright}

%========================================================= p. xii  Hilbert plate
% Halftone photograph on book p. xii, facing the Chapter 1 opener (book p. 1).
% Kept here because p. xii falls inside the front-matter page range; if the
% ch01 unit also carries it, drop one copy at integration.
\clearpage
\thispagestyle{empty}
\null\vspace*{0.05\textheight}
\FMplate{58mm}{%
  \textsc{David Hilbert}\\
  \textit{(1862--1943)}\\
  Professor of Mathematics,\\
  University of G\"ottingen, Germany\\[2pt]
  Perhaps the world's foremost mathematician during the\\
  first half of the twentieth century.\\[2pt]
  \textit{Courtesy of the New York Public Library.}}
\vfill
\cleardoublepage
\pagenumbering{arabic}
% ==================== ch01 ====================
\setcounter{theorem}{0}\setcounter{definition}{0}\setcounter{construction}{0}\setcounter{exgroup}{0}\setcounter{footnote}{0}
\setcounter{chapter}{1}
\renewcommand{\thefootnote}{\fnsymbol{footnote}}

\chapter*{\raggedright\Huge INTRODUCTION}
\addcontentsline{toc}{chapter}{1.\ Introduction}
\markboth{INTRODUCTION}{}
\vspace{-1em}\noindent\rule{\linewidth}{0.6pt}\vspace{1em}
\figurekey

% The Hilbert plate (leaf facing book p.1) is typeset at the END of the front
% matter (frontmatter.tex, p. xii slot), matching its physical position in the
% book. It was duplicated here until the 2026-08-17 integration pass.

%========================================================= 1-1
\section*{1--1\quad WHY STUDY GEOMETRY?}
\addcontentsline{toc}{section}{1--1\quad Why study geometry?}

\setcounter{footnote}{0}
There are at least four reasons why geometry is studied by so many students in
the high schools of the United States and, in fact, throughout the world. The
first of these reasons is that the logical beauty and precision of geometry has
fascinated men since the time of the ancient Greeks. Geometry is one of the
great achievements of the human mind, and for 2000 years men have regarded its
study as a necessary part of the making of a truly educated person. The feeling
has been that here is at least one thing whose truth is clear, because every
statement can be demonstrated beyond a doubt. People also felt that because the
reasoning in geometry had to be careful and accurate, the study of it would help
one to be careful and accurate in other
activities.\footnote{The Dutch philosopher Benedictus de Spinoza
(1632--1677) was so impressed by the clarity of geometrical reasoning that he
arranged his reasoning on religion in a ``geometric style.''}
Because this great invention of the mind has value as an introduction to correct
logical reasoning, geometry is a required subject for entrance to many colleges
regardless of what one plans to study.

A second important reason for studying geometry is its practical importance.
People in all kinds of occupations have occasional need for a bit of geometry,
and in some fields the study of geometry is a most important step in professional
training. Fields in which it is required are physics, chemistry, engineering,
pure mathematics, statistics, some of the biological sciences, and, more
recently, certain branches of economics and psychology.

Perhaps most students start the study of geometry because of its practical value
or because they are told it is required. But a third reason is that after you
have studied geometry you will have the knowledge that is necessary to give you a
deeper insight into the wonderful details and complexities of our world, both the
natural world and the constructions of mankind. It is this deeper understanding
of the world which can be furthered by the study of geometry that makes its study
important to all people, whether scientists or not. Nature abounds in geometric
designs, and geometric knowledge gives the observer keener interest.

A fourth reason for studying geometry is that we live in what is often called a
scientific age, and even if you yourself are not planning on a career in science
it will be necessary for you to have some understanding and appreciation for the
ways in which a scientist thinks. The study of geometry is a big step toward
getting such understanding.

%========================================================= 1-2
\section*{1--2\quad BRIEF HISTORY OF GEOMETRY}
\addcontentsline{toc}{section}{1--2\quad Brief history of geometry}

\setcounter{footnote}{0}
Knowledge of isolated geometric facts goes back before the beginning of recorded
history. Certainly the early Egyptians and Babylonians (4000--3000 B.C.) knew
many practically useful geometric relations. The construction of the pyramids
alone required considerable practical geometry.

But it was the ancient Greeks who collected the known geometric facts, discovered
new ones, and arranged them into a logically consistent system. (The word
\emph{geometry} comes from two Greek words, \emph{ge} meaning ``earth,'' and
\emph{metrei} meaning ``to measure,'' showing that originally it was thought of
as ``earth measurement.'') This process of organization and discovery took
centuries and occupied the minds of a long list of able men. Remember that the
centuries just before, during, and after the period of greatest Greek political
influence (4th and 5th centuries B.C.) were times of intense intellectual
activity. The leading minds were interested in ideas of all kinds. Mathematics
was but one of their interests.

Some of these early geometers deserve special mention either because of their
mathematical eminence or because of their influence on later men of
philosophy\footnote{The word \emph{philosophy} comes from Greek words meaning
``love of wisdom.'' Look this up in your dictionary. The early Greek
philosophers tried to understand the entire world. More on philosophy will be
found in Chapter 16 of this text.}
or science. One of these early men was Thales of Miletos (6th century B.C.) who
was considered one of the seven wise men of his day. He was the first Greek
mathematician and the first Greek astronomer. Besides knowing some geometric
facts, he probably knew proofs for many of them. The next important Greek for us
was Pythagoras (late 6th century B.C.), who established a school, or brotherhood,
in Croton, Sicily (a part of Greater Greece then). It is often hard to
distinguish whether a discovery originated with Pythagoras or one of his
followers. All ideas were communal property of the brotherhood. Besides the
Pythagorean theorem concerning right triangles, he (or his school) made notable
discoveries in music, astronomy, and arithmetic.

\Iplate{Pythagoras}{From a fresco by Raphael.\\ Courtesy of \emph{Scripta Mathematica}.}

\Iplate{The Death of Archimedes}{From a mosaic found in the ancient city of
Pompeii.\\ Courtesy of \emph{Scripta Mathematica}.}

\setcounter{footnote}{0}
One Greek geometer in particular deserves mention here because his fame is so
slight compared with his worth. This man is Eudoxus of Cnidus (4th century
B.C.), one of the foremost mathematicians of all time. His fame rests upon his
general theory of proportion, certain geometric constructions, and his ``method
of exhaustion.'' He lived during the time of
Plato.\footnote{Plato and Aristotle are the two most renowned of the Greek
philosophers. Although Eudoxus is less well known, he was a much better
mathematician than either of these men. However, both Plato and Aristotle have
had tremendous influence on later thought. It is said that above the entrance to
Plato's academy were inscribed the words, ``Let no one ignorant of geometry enter
here.''}

We find that the Greeks organized and extended geometry for their own mental
enjoyment. Deductive geometry (that is, with proofs from simple assumptions) is
truly a Greek invention. These men were not seeking practical applications at
all; in fact, an educated Greek regarded applications of geometry as beneath his
dignity. By the time of Plato (4th century B.C.), Greek scholars had a fairly
well-organized body of geometry and perhaps even a geometric text written for
members of Plato's academy.

By far the best known of the ancient Greek geometers is Euclid, who wrote the
\emph{Elements} about 300 B.C. This work is the most successful textbook of all
time and was used throughout the world until well into the present century. It
consists of 13 ``books,'' the first six of which are concerned with plane
geometry. The others deal with arithmetic and solid geometry. The
\emph{Elements} has not only been widely used as written (or as translated); it
has also been the model for innumerable other books. Isaac Newton, the renowned
English mathematician and physicist, wrote his great book, the \emph{Principia},
in the ``geometric style'' even though this often tended to hide the path he took
in making his discoveries. Most of the geometry books used in American high
schools today are somewhat altered and simplified versions of Euclid's
\emph{Elements}. It is astonishing, and a great tribute to Euclid, that this
book has retained its value for over 2000 years.

Little is known of Euclid's life except that he probably learned his mathematics
from pupils of Plato's academy and that he later founded his own school at
Alexandria.\footnote{The city of Alexandria was founded by Alexander the Great in
332 B.C. Alexander was a supporter of learning, as was his successor, Ptolemy I.
It was Ptolemy who started the great library at Alexandria, and it was during his
rule that Euclid opened a school there.}
Undoubtedly he made full use of earlier writings, but the organization of the
\emph{Elements} is his own, as are many significant improvements.

One of the many amusing stories concerning Euclid is the following. Upon being
asked by a student of what profit was the study of geometry, Euclid directed a
slave to ``give the man a coin that he may profit from what he learns.''

\setcounter{footnote}{0}
Many able Greek mathematicians followed Euclid, among whom were Apollonius (3rd
century B.C.) and Archimedes (3rd century B.C.). Archimedes, who had perhaps the
greatest mind of ancient times, did fine work in mathematics, physics, and
engineering. He was employed as a kind of technical consultant in warfare by the
king of Syracuse and was killed by a soldier in the siege of that city in 212
B.C.

With the rise of the Roman Empire, Greece declined as a political power but
retained its eminence as a center of learning and culture well into the Christian
era. Rome added little (compared with Greece) to the arts, science, or
literature. It is in the realm of practical government, warfare, and architecture
that the Romans made their contributions. But when Rome fell to the barbarians
from the north in the 5th century A.D., the ``dark ages'' began. During the dark
ages, mathematical research was largely in the hands of the Arabs. (Our word
\emph{algebra} is of Arabic origin.) Serious scientific work in Europe was almost
at a standstill until the Renaissance, or revival of learning, in the 14th
century.

Beginning with the 16th century, mathematical activity was primarily in fields
other than classical geometry, particularly after the invention of
calculus\footnote{Calculus is a powerful mathematical tool useful in solving
certain kinds of problems. It is a subject required of most science majors in
college.}
by Newton and Leibniz in the 17th century. Nevertheless, significant studies in
geometry were made which culminated in the discovery, in the 19th century by
K.~F. Gauss, N.~Lobachevski, and J.~Bolyai, independently, that geometries other
than Euclid's are possible and just as true (that is, logically correct). It was
during the 19th century that men became aware that there were certain assumptions
made by Euclid which he did not actually state. Therefore Euclidean geometry was
re-examined and made complete and precise according to modern standards. The
first really complete treatment of Euclidean geometry was published by the German
mathematician David Hilbert in 1899. Our book is based upon Hilbert's treatment
(with some simplifications). Thus you will be studying geometry that is almost
2300 years old but dressed up in modern form.

Hilbert (who died in 1943 at the age of 81) had wide-ranging mathematical
interests and made outstanding contributions to many fields. His influence on
American mathematics has been considerable, because at the time when American
mathematicians were beginning to achieve eminence, during the latter part of the
19th century, many of them studied under Hilbert at the University of
G\"ottingen in Germany.

\begin{center}\bfseries Suggested Discussion Topics\end{center}
\nopagebreak
\begin{multicols}{2}
\begin{exlist}
\item Discuss the four reasons listed for studying geometry. Can you think of
      any others?
\item What do you know about Greek civilization other than what has been
      mentioned in this chapter?
\item What were some of the major scientific developments of the 19th century?
      In what countries did they occur?
\item Look up and report on contributions of the Arabs to mathematics.
\item Look up Euclid in an encyclopedia.
\item Look up Pythagoras in an encyclopedia.
\item Look up Archimedes in \emph{The World of Mathematics}, by J.~R. Newman.
\end{exlist}
\end{multicols}

%========================================================= 1-3
\section*{1--3\quad FAMILIAR GEOMETRIC OBJECTS}
\addcontentsline{toc}{section}{1--3\quad Familiar geometric objects}

In your previous study of arithmetic and algebra you have met many ``geometric
objects,'' such as the ones shown in Fig.~1--1, where the examples are all
\emph{plane} figures. Some \emph{solid} figures that you may have encountered are
pictured in Fig.~1--2.

\FIGIONE

\FIGITWO

All these shapes are easily recognized as familiar; you would have no difficulty
in recognizing one of them on sight. However, you may have noticed that we have
not \emph{defined} the shapes, but only drawn them. Part of our job later in this
book will be to give accurate definitions. As you study geometry, it will become
more and more clear that we cannot prove anything without a clear definition. But
in this chapter, we shall most often use drawings to introduce the things you
will study later.

We have used the words ``geometric object'' or ``geometric figure,'' but have
only given examples by drawings on a page of the book. A precise definition is
the following:

\begin{definition}
A geometric figure is any \emph{set of points}.
\end{definition}

Some explanation of the words ``set of points'' is necessary. A set of points is
just a ``bunch'' of points (or ``collection'' of points, or ``aggregate'' of
points). This definition allows some rather strange sets of points to be called
geometric figures, for example those in Fig.~1--3. However, we usually deal with
simpler sets of points than these, namely, with lines, line segments, triangles,
circles, etc. A word of caution is necessary regarding our use of the word
``line.'' By ``line'' we mean a \emph{straight line} indefinite in extent in both
directions. We shall see in Chapter 3 exactly how to state what we need to know
about lines in order to prove other things about them.

\FIGITHREE

\exercisegroup{1--1}
\begin{multicols}{2}
\begin{exlist}
\item Describe in your own words what is meant by a \emph{set} of points.
\item Invent some geometric figures of your own and then describe these sets of
      points.
\item Mark a point on a sheet of paper and consider the set of all points on your
      paper such that each point in the set is two inches from this first point.
      Describe this set of points in another way.
\item Some of the shapes drawn in Fig.~1--3 are intended to represent infinitely
      many line segments or circles. Which are they? Can you invent other figures
      with infinitely many segments or circles?
\end{exlist}
\end{multicols}

%========================================================= 1-4
\section*{1--4\quad WHAT ARE POINTS?}
\addcontentsline{toc}{section}{1--4\quad What are points?}

Up to now we have been talking about points just as if we knew what they are. But
no one has ever seen a point. The smallest dot you can make on a sheet of paper
is certainly not a point. The dot would look huge to a microbe (if it could see),
and the microbe, which can be seen under a microscope, is very large compared
with the molecules of chemical substances of which it is composed. What, then, is
a point? The answer is that points are an \emph{invention} of the human mind. We
\emph{imagine} that there \emph{ought to be} points, which themselves cannot be
split up into smaller pieces. A given ``hunk'' or portion of space can be split
into two hunks. And each of these can be split in two, and so on indefinitely. A
point is what \emph{should} be left when our hunk has no extension in space or,
to use other words, when it is ``infinitely small.'' Euclid, in fact, gave
something like that as a definition of a point. But such a ``definition'' is not
really a help in geometry because from it we cannot derive anything more. All it
does is make us feel more or less happy about the existence of the things we
\emph{think} we are talking about. The modern point of view (and really Euclid's
too) is that points are undefined. The things that we do define are certain
relations between points or sets of points.

It is essentially the same regarding lines and planes. No one has ever seen a
line or a plane. A thin ray of light or a thin line on the page is only an
approximation to what we imagine that a line should be like. And a very flat
surface, like a mirror, is but an approximation to what we imagine a plane to be.
Therefore right at the beginning of geometry we have to make a strong effort of
the imagination. We must imagine a large set of things, called points, which we
call a plane. Certain parts of this plane, that is, certain special sets of
points in the plane, are called straight lines, or just lines. The study of
geometry begins when we tell \emph{how these points and lines are related to each
other}.

\exercisegroup{1--2}
\begin{multicols}{2}
\begin{exlist}
\item Draw two intersecting lines on the blackboard. You will then have some
      chalk hanging on the board. How many reasons can you give to explain that
      what you have drawn are not lines? Does the chalk on both your drawn lines
      correspond to what you think of as a point?
\item Mark two dots on the blackboard. Are they points? Is a ``blob'' of chalk a
      point? How many lines do you feel there should be which lie on two given
      points? If points were blobs of chalk, could you draw more than one line
      through them?
\item When you draw a figure to illustrate a geometric idea, what is the relation
      of the drawn figure to the geometric idea you have in mind?
\item Is it reasonable to expect a drawing to give a proof of something in
      geometry?
\item If an object had length or width, would you call it a point? Can you see
      points and lines?
\item Can you describe what you mean by straightness of a line?
\end{exlist}
\end{multicols}

%========================================================= 1-5
\section*{1--5\quad TOOLS YOU WILL NEED}
\addcontentsline{toc}{section}{1--5\quad Tools you will need}

\setcounter{footnote}{0}
Even though points and lines are creations of the mind it is necessary for us to
draw pictures of them in order to see what is ``going on.'' Thus you will need a
\emph{ruler} to draw pictures of straight lines. A ruler without a marked scale
on it is called a \emph{straightedge}, and for the constructions we shall make
this is all that is absolutely necessary. But it is very convenient to be able to
\emph{measure} the length of a line segment, and for this reason an ordinary foot
ruler, or a six-inch one, will do. However, we shall not actually deal with
measurement\footnote{It is interesting and surprising that Euclid in his
\emph{Elements} did not discuss measurement, that is, the process of getting a
\emph{number} for the distance between two points.}
of lengths until Chapter 4.

Angles will be discussed more fully later, and we shall consider measurement of
angles in Chapter 12. But you are already familiar with a device, called the
\emph{protractor}, for measuring pictures of angles drawn on paper or on the
blackboard.

You will also need a device, called a \emph{compass}, for drawing pictures of
circles. Undoubtedly you are familiar with its use.

With these simple tools you can make a number of ``constructions.'' These
constructions will seem to be obviously correct. However, we will prove later in
the book that they actually are correct.

You have learned in your previous mathematics how to use these tools. The
exercises below are for review of definitions and practice in the use of ruler,
compass, and protractor. However, throughout this course your best tool will be
your ability to think straight.

\exercisegroup{1--3}
\begin{multicols}{2}
\begin{exlist}
\item Draw a picture of a line, using your straightedge. Now if you sharpened
      your pencil \emph{very} sharp, could you draw a picture of a line inside the
      first one? Was your drawing really what you intended a line to be?
\item Draw a picture of a circle with your compass. How many differences can you
      find between your drawing and an imaginary perfect circle?
\item Draw a picture of an angle. Measure it with your protractor. Have a friend
      measure the angle. Do you get \emph{exactly} the same answers? How accurate
      do you think your measurement is?
\item How long are the sides of an angle? When you draw a picture of an angle,
      you cannot draw the sides infinitely long because your straightedge and
      paper are limited in size.
\item Draw a picture of a ray from a point $A$. Draw another ray from $A$. What
      we call the \emph{angle} consists of the points on the two rays. The points
      between the rays form the \emph{interior} of the angle. The other points of
      the plane form the \emph{exterior} of the angle. (Fig.~1--4.)
\item A \emph{right angle} has a measure of $90^\circ$. Construct a right angle
      with your protractor.
\end{exlist}
\end{multicols}

\FIGIFOUR

\begin{multicols}{2}
\begin{exlist}\setcounter{exlisti}{6}
\item Draw an angle of $37^\circ$, of $120^\circ$, of $85^\circ$.
\item An angle is \emph{bisected} if a ray from the vertex splits the angle into
      two angles of the same size. In Fig.~1--5, the ray $AC$ bisects
      $\ang{BAD}$ because $\ang{CAB}$ and $\ang{DAC}$ are the same size. The ray
      $AC$ is the bisector of the angle. Draw an angle of $45^\circ$ and bisect
      it, using your protractor.
\item Draw a picture of a triangle. Measure each of the angles of the triangle.
      What is their sum? Repeat for several different triangles. Can you guess
      something that is true for all triangles?
\item As in Fig.~1--6, an \emph{acute angle} is an angle of less than
      $90^\circ$; an \emph{obtuse angle} is an angle of more than $90^\circ$.
      Draw two acute angles and two obtuse angles with your protractor.
\item Draw a triangle and measure the sides and the angles. Which side is
      opposite the largest angle? opposite the smallest angle? Draw another
      triangle and repeat the measurements. Can you guess something which now
      seems to be true for all triangles?
\item A triangle that has one right angle is called a \emph{right triangle}. Draw
      a picture of a right triangle, as in Fig.~1--7. Measure the two acute
      angles of the triangle. What is their sum? Draw another right triangle and
      make the same measurements. Is the sum of the two angles the same as
      before?
\end{exlist}
\end{multicols}

\FIGIFIVESIX

\FIGISEVENEIGHT

\begin{multicols}{2}
\begin{exlist}\setcounter{exlisti}{12}
\item One angle is said to be \emph{complementary} to another if their sum is
      $90^\circ$. Construct two such angles with your protractor. In Fig.~1--8,
      if $\ang{AOB}$ is a right angle, what can you say about $\ang{COA}$? about
      $\ang{COB}$?
\item One angle is said to be \emph{supplementary} to another if their sum is
      $180^\circ$. Construct two such angles with your protractor. What can you
      say about $\ang{AOB}$ and about $\ang{BOC}$ in Fig.~1--9?
\item Two distinct rays $OA$ and $OB$ which lie on one straight line form a
      \emph{straight angle} (Fig.~1--10). What is the measure of a straight
      angle?
\item A \emph{quadrilateral} (Fig.~1--11) is a figure with four sides, just as a
      triangle is a figure with three sides. Draw a picture of a quadrilateral and
      measure its angles. What is their sum?
\end{exlist}
\end{multicols}

\FIGININETEN

\FIGIELEVEN

\begin{multicols}{2}
\begin{exlist}\setcounter{exlisti}{16}
\item Add the following:
  \begin{enumerate}[label=(\alph*), leftmargin=1.9em, itemsep=1pt, topsep=2pt]
  \item $12^\circ\ 10'\ 25''$, $31^\circ\ 48'\ 30''$, and $43^\circ\ 50'\ 40''$.
  \item $17^\circ\ 18'\ 19''$, $90^\circ\ 9'\ 9''$, and $41^\circ\ 48'\ 48''$.
  \item $121^\circ\ 18'\ 17''$, $40^\circ\ 40'\ 40''$, and $18^\circ\ 1'\ 3''$.
  \item $66^\circ\ 55'\ 44''$, $55^\circ\ 44'\ 33''$, and $44^\circ\ 33'\ 22''$.
  \end{enumerate}
\item Subtract:
  \begin{enumerate}[label=(\alph*), leftmargin=1.9em, itemsep=1pt, topsep=2pt]
  \item $41^\circ\ 15'\ 13''$ from $63^\circ\ 28'\ 44''$.
  \item $41^\circ\ 15'\ 13''$ from $50^\circ\ 18'\ 7''$.
  \item $61^\circ\ 28'\ 37''$ from $80^\circ\ 18'\ 23''$.
  \end{enumerate}
\item What is the complement of an angle of
  \begin{enumerate}[label=(\alph*), leftmargin=1.9em, itemsep=1pt, topsep=2pt]
  \item $51^\circ\ 17'\ 18''$?
  \item $42^\circ\ 38'\ 29''$?
  \item $A^\circ$?
  \end{enumerate}
\item What is the supplement of an angle of
  \begin{enumerate}[label=(\alph*), leftmargin=1.9em, itemsep=1pt, topsep=2pt]
  \item $51^\circ\ 17'\ 18''$?
  \item $138^\circ\ 4'\ 5''$?
  \item $A^\circ$?
  \end{enumerate}
\item If an angle is twice as large as its complement, what is its degree
      measure? [\emph{Hint}: Set up an equation to solve.]
\item Two angles have a sum of $56^\circ\ 28'$ and a difference of
      $8^\circ\ 46'$. Find the angles.
\item Find the complement and supplement of an angle of $56^\circ\ 28'\ 31''$.
      What is the difference between the supplement and complement? Can you guess
      something that is true for all acute angles?
\item What kind of angle is the supplement of an acute angle? of an obtuse angle?
      Does an obtuse angle have a complement?
\end{exlist}
\end{multicols}

%========================================================= 1-6
\section*{1--6\quad RULER AND COMPASS}
\addcontentsline{toc}{section}{1--6\quad Ruler and compass}

During their study of geometry, the ancient Greeks found that many interesting
figures could be drawn with a straightedge and compass alone. Drawings needing
only these tools are called \emph{ruler and compass constructions}. The Greeks
tried to solve all construction problems using only these tools, but they
encountered some that could not be done this way.

The exercises below involve some simple ruler and compass constructions that the
Greeks knew. Exercise 2, in fact, is the first proposition in Euclid's
\emph{Elements}. We shall return to these constructions again in Chapter 7, where
we shall \emph{prove} that some of them are correct.

\exercisegroup{1--4}
\begin{multicols}{2}
\begin{exlist}
\item A triangle having two sides of equal length is called an \emph{isosceles
      triangle}. Draw a picture of an isosceles triangle. Measure the angles
      opposite the equal sides. Draw several isosceles triangles, each with the
      equal sides three inches long.
\item A triangle with all sides the same length is called an \emph{equilateral
      triangle} (Fig.~1--12). Draw a picture of an equilateral triangle. This is
      easy to do with your compass and straightedge. What is the degree measure of
      each angle in an equilateral triangle? A triangle with all angles equal is
      called an \emph{equiangular triangle}. Are equilateral triangles
      equiangular? Are equiangular triangles equilateral?
\item Draw an angle like that in Fig.~1--13. Use your compass to copy $\ang{AOB}$
      as $\ang{A'O'B'}$ in the figure. Measure the two angles with your protractor
      to see if they are the same size.
\end{exlist}
\end{multicols}

\FIGITWELVETHIRTEEN

\begin{multicols}{2}
\begin{exlist}\setcounter{exlisti}{3}
\item Draw the figures shown in Fig.~1--14.
\item Use your straightedge and compass to draw and bisect an angle
      (Fig.~1--15). Check your work by measurement with a protractor. State in
      your own words why you think the construction is correct.
\item Angles which have a common vertex and a common side and have no interior
      points in common are said to be \emph{adjacent}. Thus, in Fig.~1--16,
      $\ang{AOB}$ is adjacent to $\ang{BOC}$. Is $\ang{AOC}$ adjacent to
      $\ang{BOC}$? Are the angles in Fig.~1--17 adjacent to one another?
\item Draw a picture of a segment. From some point $A$ on a line, construct with
      your compass a segment $AB$ as long as the segment you first drew. (See
      Fig.~1--18.)
\end{exlist}
\end{multicols}

\FIGIFOURTEEN

\FIGIFIFTEENSIXTEEN

\FIGISEVENTEENEIGHTEEN

\begin{multicols}{2}
\begin{exlist}\setcounter{exlisti}{7}
\item Two lines are \emph{perpendicular} if they intersect at right angles.
      Another way of defining perpendicular would be to say that all the adjacent
      angles formed by the intersecting lines are the same size. To indicate that
      $l$ and $m$ are perpendicular, we write $l \perp m$.

      A line $l$ is called the \emph{perpendicular bisector} of the segment $AB$
      if the line $l$ is perpendicular to $AB$ and intersects the segment $AB$ in
      a point $C$ such that $AC$ and $CB$ have equal lengths. The point $C$ is
      called the \emph{midpoint} of $AB$. (See Fig.~1--19.)

      Draw a line segment and construct its perpendicular bisector with
      straightedge and compass. Check by measurement. State why you think the
      construction is correct.
\item Draw a line $l$ and mark a point $O$ on it. With straightedge and compass,
      construct a line $m \perp l$ and passing through $O$. (See Fig.~1--20.)
\item Draw a line $l$ and mark a point $P$ not on the line. With your straightedge
      and compass, construct a line $m$ through $P$ perpendicular to $l$. Check by
      measurement with your protractor. State why you think the construction is
      correct. (See Fig.~1--21.)
\item A line segment drawn from a vertex of a triangle to the midpoint of the
      opposite side is called a \emph{median} of the triangle. A triangle has
      three medians that meet in one point $O$ (see Fig.~1--22). We shall prove
      this later.

      Draw a triangle and its medians. Can you describe where the medians
      intersect? Try to find how to describe $O$ by measurement. Make this
      construction for several triangles.
\item Cut a triangle out of cardboard. Find where the medians intersect. If your
      work is accurate you will find that the triangle will balance on your finger
      at the place where the medians intersect (Fig.~1--23).
\end{exlist}
\end{multicols}

\FIGININETEENTWENTY

\FIGITWENTYONETWENTYTWO

\FIGITWENTYTHREE

\begin{multicols}{2}
\begin{exlist}\setcounter{exlisti}{12}
\item An \emph{altitude} of a triangle is a line segment drawn from a vertex
      perpendicular to the line \emph{containing} the opposite side [see
      Fig.~1--24, (a) and (b)]. A triangle has three altitudes. The lines
      containing the altitudes of a triangle always intersect in one point [see
      Fig.~1--24(c)]. Draw several triangles and their altitudes.
\item Draw a triangle and bisect each of the angles of the triangle. The three
      bisectors will meet in one point (Fig.~1--25).
\end{exlist}
\end{multicols}

\FIGITWENTYFOUR

\FIGITWENTYFIVE

\begin{multicols}{2}
\begin{exlist}\setcounter{exlisti}{14}
\item Two lines, $l$ and $m$, are said to be \emph{parallel} if they do not
      intersect. We indicate ``$l$ is parallel to $m$'' by writing $l \pll m$.
      (Remember that lines are infinitely long.) If another line $n$ intersects
      the parallel lines, then the angles 1 and 2 in Fig.~1--26 are the same size.
      Conversely, if angles 1 and 2 are the same size, then the lines are
      parallel.

      As in Fig.~1--27, draw a line $l$ and a line $n$ intersecting it. Choose a
      point $P$ on $n$ but not on $l$. At $P$, copy angle 2 in order to get a line
      parallel to $l$. Make this construction with your protractor and then repeat
      it with your straightedge and compass.

      We shall prove later that this construction with straightedge and compass is
      correct.
\item Draw a line $l$ and mark a point $A$ not on $l$. Construct a line through
      the point $A$ parallel to $l$.
\item Draw a right triangle (Fig.~1--28). The side opposite the right angle is
      called the \emph{hypotenuse}. The other two sides are called the \emph{legs}
      of the triangle. Measure the lengths of the sides in your right triangle. If
      $a$, $b$, and $c$ are the lengths of the sides, where $c$ is the hypotenuse,
      you will find that $a^2 + b^2 = c^2$. This formula is called the
      \emph{Pythagorean theorem}. It was discovered (probably) by the Greek
      Pythagoras in the 6th century B.C. We will prove this important theorem
      later in the book. Test the truth of the Pythagorean theorem by measuring
      several different right triangles.
\end{exlist}
\end{multicols}

\FIGITWENTYSIXTWENTYSEVEN

\FIGITWENTYEIGHT

%========================================================= Review
\section*{REVIEW OF CHAPTER 1}
\addcontentsline{toc}{section}{Review of Chapter 1}
\markboth{REVIEW OF CHAPTER 1}{}

\begin{exlist}
\item Make sure that you are familiar with all the following terms:
\end{exlist}

\begin{center}
\begin{tabular}{@{}p{0.42\linewidth}p{0.42\linewidth}@{}}
line segment            & perpendicular    \\
angle                   & right angle      \\
angle bisector          & equilateral      \\
triangle                & equiangular      \\
degree, minute, second  & isosceles        \\
midpoint                & median           \\
complementary           & altitude         \\
supplementary           & hypotenuse       \\
straight angle          & diameter         \\
obtuse angle            & radius           \\
acute angle             & adjacent angles  \\
ray                     & intersect        \\
set                     & contain          \\
\end{tabular}
\end{center}

\begin{exlist}\setcounter{exlisti}{1}
\item Make sure that you can do the following straightedge and compass
      constructions:
  \begin{enumerate}[label=(\alph*), leftmargin=2.2em, itemsep=1pt, topsep=2pt]
  \item Copy a segment.
  \item Draw the bisector of a segment.
  \item Draw a line $\perp$ to a given line from a point not on the given line.
  \item Draw a line $\perp$ to a given line from a point on a given line.
  \item Draw a line $\pll$ to a given line, through a point not on the line.
  \item Bisect an angle.
  \item Copy an angle.
  \end{enumerate}
\end{exlist}

\begin{center}\bfseries Brief Review Tests\end{center}
\nopagebreak

Decide whether the following statements are true or false. If in doubt, look up
the answer. Do not guess.

\medskip
\noindent\textbf{I.} \emph{History}
\begin{exlist}
\item Euclid lived before Pythagoras.
\item Euclid collected, organized, and added to the geometry of his day.
\item Euclid's geometry was incomplete as he wrote it. It was completed in the
      19th century.
\item Archimedes lived at the same time as Euclid.
\item Thales lived before Euclid.
\item Greece was superior to Rome in all ways.
\item Rome was a great center for the study of geometry.
\end{exlist}

\medskip
\noindent\textbf{II.} \emph{Geometry}
\begin{exlist}
\item A straightedge is used for measurement.
\item A line may have finite length.
\item A straight angle contains $90^\circ$.
\item If angles have a common vertex and a common side, they are adjacent.
\item Complementary angles have a sum of $180^\circ$.
\item In a right triangle the side opposite the right angle is called the
      hippopotamus.
\item The altitude of a triangle is its distance from the center of the earth.
\item If two lines intersect they are parallel.
\end{exlist}

\begin{center}\bfseries ARITHMETIC AND ALGEBRA REFRESHER\end{center}
\nopagebreak

\begin{exlist}
\item Carry out the arithmetic operations indicated:
  \begin{multicols}{2}
  \begin{enumerate}[label=(\alph*), leftmargin=2.2em, itemsep=1pt, topsep=2pt]
  \item $[(12)(8)] \div 6$
  \item $2992 \div 68$
  \item $326 + (-257)$
  \item $\frac{1}{2} + \frac{1}{4}$
  \item $\frac{1}{3} \div \frac{11}{3}$
  \item $56789 \div 357$
  \item $\frac{2}{3} - \frac{2}{5}$
  \item $3.762 \div 5.7$
  \item $\frac{1}{2} \cdot \frac{1}{3}$
  \item $\frac{1}{2} \div \frac{1}{3}$
  \item $(-24) \div (-6/5)$
  \item $[(-6.1) + (-8.5)] - 2.7$
  \item $3/7 \cdot 35/24$
  \item $2.5679/0.1372$
  \item $[3.4 - 8.5] + (-2.7)$
  \item $(5 + \frac{1}{2})/33$
  \end{enumerate}
  \end{multicols}
\item Write the answers to the following. If you take more than a minute and miss
      any part of this problem, you need more review.
  \begin{multicols}{3}
  \begin{enumerate}[label=(\alph*), leftmargin=2.2em, itemsep=1pt, topsep=2pt]
  \item $2 \times 2 =$
  \item $2 \div 2 =$
  \item $2 \times 0.2 =$
  \item $2 \div 0.2 =$
  \item $2 \div \frac{1}{2} =$
  \item $2 \times \frac{1}{2} =$
  \item $\frac{1}{2} + \frac{1}{3} =$
  \end{enumerate}
  \end{multicols}
\item Write formulas for the following. Be sure that you tell what each letter
      stands for. For example: If $A$ is the area of a rectangle $l$ its length
      and $w$ its width, then $A = lw$.
  \begin{enumerate}[label=(\alph*), leftmargin=2.2em, itemsep=1pt, topsep=2pt]
  \item The volume of a box.
  \item The area of a circle.
  \item The circumference of a circle.
  \item The area of a triangle.
  \item The product of a number $z$, by one more than the number.
  \item Twice a number plus 5 is equal to 13.
  \item One number is 4 times another.
  \end{enumerate}
\item Find $x$ if
  \begin{multicols}{2}
  \begin{enumerate}[label=(\alph*), leftmargin=2.2em, itemsep=1pt, topsep=2pt]
  \item $2x - 5 = 3$
  \item $3x - 5 = x + 7$
  \item $2x = 3$
  \item $ax = b$, $a \neq 0$
  \item $x^2 = 4$
  \item $\frac{1}{3} \cdot (12x) = 44$
  \item $2x/3 + 1 = 5$
  \item $4x + 5x - x = 96$
  \item $x \div 6 = 5 \div 24$
  \item $6.4 = 2x - 0.4$
  \end{enumerate}
  \end{multicols}
\item If $x$ and $y$ are numbers, what do we call $xy$? $x + y$?
\item Write equations for the following statements:
  \begin{enumerate}[label=(\alph*), leftmargin=2.2em, itemsep=1pt, topsep=2pt]
  \item The sum of two numbers is 17 and their difference is 7.
  \item A number $z$ is equal to the sum of the product of $x$ and $y$, and the
        product of $a$ and $b$.
  \item The number $\pi$ is equal to the quotient of the circumference of a
        circle divided by its diameter.
  \item The volume of a sphere is equal to $\frac{4}{3}$ of the product of $\pi$
        and the cube of the radius.
  \item Einstein's famous formula for the equivalence of mass and energy is as
        follows: If $E$ is the energy of a mass $m$, and $c$ is the velocity of
        light, in suitable units, then $E$ is equal to the product of the mass by
        the square of the velocity of light. Write this formula.
  \end{enumerate}
\end{exlist}

% The Aristotle plate (book p.20) is NOT emitted here. It is the facing-page
% plate for the Chapter 2 opener and ch02.tex already emits it immediately
% before \chapter*{LOGIC}. Keeping it in both units printed it twice in the
% assembled book. Ownership: ch02. See ch01-UNCERTAIN.md "Text review".
\clearpage
% ==================== ch02 ====================
\setcounter{theorem}{0}\setcounter{definition}{0}\setcounter{construction}{0}\setcounter{exgroup}{0}\setcounter{footnote}{0}
\setcounter{chapter}{2}

% ---- local macros (hoist into book preamble) ----
% Halftone photographic plate: framed placeholder, never redrawn (B&W pass).
\newcommand{\IIplate}[2]{%
  \par\medskip
  \begin{center}
  \fbox{\begin{minipage}{0.72\linewidth}\centering\vspace{2.4cm}
    {\footnotesize\itshape[halftone photographic plate --- not redrawn]}
    \vspace{2.4cm}\end{minipage}}%
  \par\smallskip
  {\footnotesize\textsc{#1}\\ #2}
  \end{center}
  \par\medskip}
% Unnumbered star footnote. \starnote in brumfiel.sty emits a bare
% \footnotetext, which prints the footnote counter ("0"); the book marks this
% note with the same "*" the starred exercises carry, and no number.
\newcommand{\IIstarnote}{\begingroup\renewcommand{\thefootnote}{}%
  \footnotetext{*Starred exercises, sections, theorems, etc., are optional.}%
  \endgroup}
% Ch 2's eight "Problem." blocks. These are NOT the book's "Remark." blocks:
% p.22 etc. print an italic "Problem." head and nothing else, so wrapping them
% in {remark} (which emits "Remark.") printed a word the book never sets.
% Same shape as brumfiel.sty's {remark}, with the correct head.
\newenvironment{IIproblem}{\par\medskip\noindent\emph{Problem.}\ }{\par\medskip}
% ---- end local macros ----

% Facing-page plate for the chapter opener (book p.20, facing p.21).
\IIplate{Aristotle}{A view of a statue in a Vienna museum.\\
Courtesy of \emph{Scripta Mathematica}.}

\chapter*{\raggedright\Huge LOGIC}
\addcontentsline{toc}{chapter}{2.\ Logic}
\markboth{LOGIC}{}
\vspace{-1em}\noindent\rule{\linewidth}{0.6pt}\vspace{1em}

%========================================================= 2-1
\section*{2--1\quad STATEMENTS}
\addcontentsline{toc}{section}{2--1\quad Statements}

Everyone thinks he knows what a \emph{statement} is. A politician issues a
\emph{statement} when he concedes defeat during an election. The plumber sends
your father a \emph{statement} after installing a new sink. But in mathematics we
use the term in a narrow technical sense. We mentioned in the first chapter that
clear, precise definitions are necessary. Accordingly, we make the following
definition, which turns out to be very useful.

\begin{definition}
A statement is a sentence which is either true or false, but not both.
\end{definition}

Consider the following:

\begin{enumerate}[label=(\alph*), leftmargin=2.2em, itemsep=1pt, topsep=3pt]
\item Go back home.
\item $2 \times 2 = 5$.
\item $2 \times 2 = 4$.
\item New York is a suburb of Los Angeles.
\item Thursday up tar paper.
\item The silliness of speedily.
\end{enumerate}

Some of the above are statements and some are not. Before reading on, decide for
yourself which ones satisfy the definition.

Let us look at some slightly different examples.

\begin{enumerate}[label=(\alph*), leftmargin=2.2em, itemsep=1pt, topsep=3pt]
\item I am the most beautiful girl in town.
\item If I don't get a good night's sleep, I can't pass the examination.
\item $x + 2 = 5$.
\end{enumerate}

All of these are sensible sentences. The only question is whether or not they
satisfy our condition of being either true or false but not both. Before we can
decide whether (a) is a true statement or a false statement we must know who
``I'' is and what town is meant. Then we can decide upon the truth or falsity.
(We could hold a beauty contest!) In (c) if we replace the letter $x$ by ``3,''
we have a true statement. If we replace $x$ by ``2'' the resulting statement is
false. In statement (b), we must again know who ``I'' is and what examination is
meant in order to decide whether it is a true or false statement.

Besides rather simple statements, we can have more complicated ones, such as:

\begin{enumerate}[label=(\alph*), leftmargin=2.2em, itemsep=1pt, topsep=3pt]
\item I cannot go to the dance because I have nothing to wear, my feet hurt, and
      I can't dance.
\item A triangle has three equal sides if and only if it has three equal angles.
\item Jones is a mathematician and Smith is an engineer, while Brown doesn't work
      at all.
\item There is one and only one line through two points.
\item Either I will marry you or get a new job.
\item Two lines either intersect or are parallel.
\end{enumerate}

In some of these examples the full statement is made up of simpler statements. We
see that statements may be combined to give more complicated statements. For
example:

\begin{center}
``Mr.\ Baker is a barber,''\qquad ``Mr.\ Barber is a baker,''
\end{center}

\noindent are two statements. We can combine them into a single statement by
joining them with ``and.''

\begin{center}
``Mr.\ Baker is a barber and Mr.\ Barber is a baker.''
\end{center}

Here are some examples of combining two statements. Make up some others of your
own.

\begin{enumerate}[label=(\alph*), leftmargin=2.2em, itemsep=1pt, topsep=3pt]
\item We are having a picnic and it is not raining.
\item It is night and I can hear the wind.
\item Mary is going with Fred and Suzie is going with John.
\item I like spinach and I dislike candy.
\item You are rich and I am poor.
\item Thursday is April second and the moon is made of green cheese.
\end{enumerate}

To be perfectly clear about combining statements, we indicate statements by
letters, using Greek letters like $\al$ (alpha), $\be$ (beta),
$\ga$ (gamma). Now if we let $\al$ indicate the statement ``Mr.\ Barber is a
baker'' and let $\be$ indicate ``Mr.\ Baker is a barber,'' then the statement
``Mr.\ Barber is a baker and Mr.\ Baker is a barber'' may be written
``$\al$ and $\be$.''

\begin{IIproblem}
In each of the combined statements (a) through (f) above, pick
out the two statements which form it. Note that each statement has the form
``$\al$ and $\be$.''
\end{IIproblem}

The combination of two statements with ``and'' to get a new statement has a name.
We emphasize this by making the definition:

\begin{definition}
If $\al$ and $\be$ are statements, the statement ``$\al$ and $\be$'' is called the
\emph{conjunction} of $\al$ and $\be$.
\end{definition}

You should not accept definitions blindly. It is possible to sneak a false
statement into a definition! Such a definition would be worthless. From the way
Definition 2--2 is worded it is clear that we have \emph{taken for granted} that
``$\al$ and $\be$'' will always be a statement. But a sentence is not a statement
(according to Definition 2--1) unless it is either true or false. Therefore, we
must decide upon some \emph{rule} which will enable us to say when ``$\al$ and
$\be$'' is true and when it is false. If we cannot agree upon such a rule,
Definition 2--2 will be meaningless.

The following statements are all of the form ``$\al$ and $\be$.'' Decide which
ones are true and which ones are false, and then try to state a general rule for
deciding upon the truth of such statements.

\begin{enumerate}[label=(\alph*), leftmargin=2.2em, itemsep=1pt, topsep=3pt]
\item $2 \times 2 = 4$ and $2 + 2 = 5$.
\item $3 + 2 = 5$ and there are at least 52 Sundays in every year.
\item A triangle has three sides and a pentagon has six sides.
\item Every rectangle is a square and all men have purple hair.
\item The positive square root of 16 is 4 and the negative square root of 16 is
      $-4$.
\end{enumerate}

You will probably have no trouble coming to the same conclusion that
mathematicians have agreed upon. That is, ``$\al$ and $\be$'' is considered true
when both $\al$ and $\be$ are true. If one of the statements is false (or both
are), then ``$\al$ and $\be$'' is false. We can describe this rule most clearly by
the \emph{truth table} below.

\begin{truthtable}{ccc}
\tth{\al} & \tth{\be} & \multicolumn{1}{c}{$\al$ and $\be$} \\[2pt]
\TT & \TT & \TT \\
\TT & \FF & \FF \\
\FF & \TT & \FF \\
\FF & \FF & \FF \\
\end{truthtable}

The table is self-explanatory. It points out that ``$\al$ and $\be$'' is true (T)
in only one case and false (F) in all others.

The next step in our development of logic will be harder. We can combine
statements in another way by putting the word ``or'' between them. For example:

\begin{quote}
I will go to the movies \emph{or} I will stay home.\\
I will receive an $A$ in geometry \emph{or} Latin.\\
State University will not admit you unless you have had a year of Spanish
\emph{or} a year of Latin.
\end{quote}

Study the three statements above carefully. Everyone would agree that in the
first statement the speaker means that he will go to the movie \emph{or} he will
stay home \emph{but he will not do both}. In the second statement does the
speaker mean that he will get \emph{either} an $A$ in geometry \emph{or} an $A$
in Latin \emph{but not both}? Or does he mean that he is sure of getting one $A$
and perhaps he will get two? Of course there is no question about the last
statement. State University will not refuse to admit you because you had one year
of Spanish \emph{and} one year of Latin.

We are pointing out that in everyday speech we use the word ``or'' ambiguously.
Sometimes when we say ``$\al$ or $\be$'' we mean that \emph{one and only one} of
the statements is true. Sometimes we mean that \emph{at least} one is true and
possibly both are. We cannot use the word ``or'' with this double meaning in
mathematics. We must explain precisely what we mean when we say ``$\al$ or
$\be$.'' Mathematicians have agreed to call statements of the form ``$\al$ or
$\be$'' true in all cases except when $\al$ and $\be$ are both false. Of course
this differs from the ordinary usage of the word ``or.'' But we use everyday
language too carelessly for our purpose in mathematics. We must be able to
express ourselves precisely and know that our meaning will not be misunderstood.
You must remember that when we say ``$\al$ or $\be$'' we mean that either $\al$
is true and $\be$ is false, or $\al$ is false and $\be$ is true, or $\al$ is true
and $\be$ is true. We make the following formal definition.

\begin{definition}
The statement ``$\al$ or $\be$'' is called the \emph{disjunction} of $\al$ and
$\be$. It is false when both $\al$ and $\be$ are false and true in all other
cases.
\end{definition}

\begin{IIproblem}
Construct the truth table for the statement ``$\al$ or $\be$.''
\end{IIproblem}

\exercisegroup{2--1}

In each exercise below there are two statements. Form their conjunction and
disjunction. Discuss the resulting statements with respect to truth and falsity.

\begin{multicols}{2}
\begin{exlist}
\item Some roses are red. Some violets are blue.
\item September has 30 days. There are eight days in a week.
\item 3 is less than 4. 4 is less than 3.
\item Napoleon was an emperor of Japan. Marco Polo discovered America.
\item $2 + 2 = 5$; \quad $6 + 3 = 9$.
\item $7 + 4 = 11$; \quad $8 + 0 = 9$.
\item $1 + 1 = 2$; \quad $1 \times 1 = 2$.
\item $5 + 1 = 5$; \quad $5 \times 1 = 6$.
\end{exlist}
\end{multicols}

%========================================================= 2-2
\section*{2--2\quad NEGATION}
\addcontentsline{toc}{section}{2--2\quad Negation}

Sometimes we make statements about other statements. One of the simplest and most
useful statements of this type has the form ``$\al$ is false.'' At times you may
make a statement that you believe true, and someone else may show his
disagreement by saying, ``That is not true.'' Just as we use the terms
\emph{conjunction} and \emph{disjunction} to describe the statements ``$\al$ and
$\be$'' and ``$\al$ or $\be$,'' so also we have a particular name for the
statement ``$\al$ is false.''

\begin{definition}
The statement ``$\al$ is false'' is called the \emph{negation} of $\al$ (and will
be written, and read, ``not $\al$'').
\end{definition}

\begin{IIproblem}
State a rule for determining when ``not $\al$'' is true and when
it is false. Make a truth table for ``not $\al$.''
\end{IIproblem}

When we form the negation of a statement in words we do not usually place ``not''
in front of it. Logically it would be correct to do so, but the resulting
sentence might sound awkward. For example, let $\al$ be ``a square is a
rectangle.'' Anyone untrained in logic would have a right to complain that he did
not understand our meaning if we said, ``not a square is a rectangle.'' The
examples below indicate how negations may be formed. After studying these, make
up several examples of your own.

\begin{enumerate}[label=(\alph*), leftmargin=2.6em, itemsep=3pt, topsep=3pt]
\item $\al$: A square is a rectangle.\\
      not $\al$: A square is not a rectangle.
\item $\al$: $2 + 2 = 4$.\\
      not $\al$: $2 + 2 \neq 4$.
\item $\al$: A triangle has five sides.\\
      not $\al$: A triangle does not have five sides.
\item $\al$: It is false that my donkey has three legs.\\
      not $\al$: My donkey has three legs.
\end{enumerate}

\exercisegroup{2--2}

Form negations of the following statements.

\begin{multicols}{2}
\begin{exlist}
\item The grass in my yard is green.
\item This rose is black.
\item I am tired.
\item The temperature of the sun is more than 1000 degrees.
\item Geometry is interesting.
\item $5 + 4 = 9$.
\item $6 + 3 = 7$.
\item[\stex 8.] $x = 2$ and $y = 3$.
\item[\stex 9.] $x = 2$ or $y = 3$.
\item[\stex 10.] All pigs are fat.
\item[\stex 11.] Some Americans are liars.
\end{exlist}
\end{multicols}
% p.25 carries the "*" footnote at the foot of the page holding the first
% starred exercises (Ex. Group 2-2, nos. 8-11); anchor it here, not mid-prose.
\IIstarnote{}

Sometimes we wish to state the negation of a statement like ``All cats are
black.'' If we can find \emph{one} cat which is not black we will have proved
this statement false. Hence we form the negation of ``All cats are black'' by
saying ``There is at least one cat which is not black'' or ``Some cat is not
black.'' The most common error made in forming the negation of a
statement like ``All students in this class like geometry'' is to say ``No
student in this class likes geometry.'' These statements certainly do not agree
but they are definitely not negations of each other. Why? Remember that if $\al$
is false then ``not $\al$'' must be true, for ``not $\al$'' states that $\al$ is
false. Show that the statements, ``All students in this class like geometry'' and
``No student in this class likes geometry,'' could both be false.

Another common error is to negate ``All cats are black'' by saying ``All cats are
not black.'' Although in everyday speech this sentence might be understood to
mean ``Not all cats are black,'' we do not interpret it so. If we say ``All cats
are red,'' we mean \emph{every cat has the property of being red}. So we interpret
``All cats are not black'' to mean \emph{every cat has the property of being not
black}. This is a subtle distinction, and again it points out that the precise
language of mathematics differs from careless, everyday speech.

\exercisegroup{2--3}

Give the negation of each statement.

\begin{multicols}{2}
\begin{exlist}
\item All men in Georgia are at least six feet tall.
\item All students in France like geometry.
\item All students in Russia hate geometry.
\item Every rectangle is a square.
\item Every triangle is equilateral.
\item All gadgets are widgets.
\item All miggles are maggles.
\item[\stex 8.] For every pair of numbers $a$ and $b$, $a + b = b + a$.
\item[9.] Every circle has a center.
\end{exlist}
\end{multicols}

Consider the statement ``\emph{Some} cats are black.'' We usually understand the
word ``some'' to mean ``more than one.'' But in logic it is more convenient to
agree that it means ``one or more.'' (We must have precise definitions!) The
negation of ``Some cats are black'' then becomes ``No cat is black.'' We could
also say, ``Every cat is a color different from black'' or ``Every cat is not
black.'' Notice again that this last statement, ``Every cat is not black,'' means
something altogether different from ``Not every cat is black.''

\exercisegroup{2--4}

Give negations of the following statements. After doing this, make up some
examples of your own.

\begin{multicols}{2}
\begin{exlist}
\item Some men play golf.
\item Some boys play football.
\item There is at least one equilateral triangle.
\item For some numbers $x$, $x^2 + 2 = 5$.
\item There is at least one number $x$ such that $x + 3 = 7$.
\item There is at least one number $x$ such that $x^2 + 2 = 0$.
\item Some dogs have fewer than four legs.
\item[\stex 8.] Some numbers are negative.
\item[9.] Some fishermen always tell the truth.
\item[10.] Some people work hard.
\end{exlist}
\end{multicols}

When we have two statements, $\al$ and $\be$, and wish to form the \emph{negation}
of their \emph{conjunction}, we must think a little harder than before. Thus if we
wish to form the negation of ``Smith is a dentist and Jones is a teacher,'' we
must say that the statement is false. By our agreement in Section 2--1, we must
say that at least one of the statements ``Smith is a dentist'' and ``Jones is a
teacher'' is false. We do this by saying, ``Smith is \emph{not} a dentist or Jones
is \emph{not} a teacher.''

As another illustration, suppose one student says, ``We will win our game with
Lincoln \emph{and} we will win our game with Bay City.'' To call his statement
false is to say, ``We will \emph{not} win our game with Lincoln \emph{or} we will
\emph{not} win our game with Bay City.'' It should be clear now that the negation
of ``$\al$ and $\be$'' is the statement ``not $\al$ or not $\be$.''

\begin{IIproblem}
Construct a truth table for each of the statements ``$\al$ and
$\be$'' and ``not $\al$ or not $\be$.'' Show that when one statement is true the
other is false, and vice versa.
\end{IIproblem}

\exercisegroup{2--5}

Give the negations of the statements below. Make up some additional exercises of
your own.

\begin{multicols}{2}
\begin{exlist}
\item Today is Monday and the weather is cold.
\item $x = 2$ and $y = 5$.
\item $2 \times 2 = 4$ and $3 \times 5 = 15$.
\item $\tri{ABC}$ is isosceles and $\tri{ABC}$ is a right triangle.
\item The point $A$ is on the line $l$ and the point $B$ is on the line $m$.
\item[\stex 6.] All cats are black and all men are mortal.
\item[\stex 7.] All cats are black and some birds are blue.
\item[\stex 8.] Some rectangles are squares and some triangles are equilateral.
\item[\stex 9.] Some widgets are gadgets and no wiggles are waggles.
\end{exlist}
\end{multicols}

We may wish to form the \emph{negation} of the \emph{disjunction} of two
statements. For example, we may wish to deny the statement ``Johnson is a tailor
or Brown is a doctor.'' You should have no trouble expressing the negation of this
statement.

\begin{IIproblem}
State the general rule for forming the negation of the
disjunction of two statements. Do so by completing the sentence, ``The negation of
`$\al$ or $\be$' is \underline{\hspace{5em}}.'' Construct a truth table for each of
the statements ``$\al$ or $\be$'' and ``not $\al$ \emph{and} not $\be$.''
\end{IIproblem}

\exercisegroup{2--6}

Give the negation of each statement below. Invent some more examples.

\begin{multicols}{2}
\begin{exlist}
\item I will go to the game or I will go to the movie.
\item I will study hard or I will fail.
\item We will start at 6:30 or we will be late.
\item A donkey has three legs or a rabbit has long ears.
\item $2 + 2 = 5$ or $2 + 2 = 6$.
\item $3 \times 7 = 21$ or $4 \times 5 = 20$.
\item Billings is right or Jones is wrong.
\item[\stex 8.] Some will work or all will starve.
\item[\stex 9.] Every triangle is isosceles or every triangle is right.
\item[\stex 10.] No roses are red or some violets are blue.
\end{exlist}
\end{multicols}

%========================================================= 2-3
\section*{2--3\quad WHAT ARE PROOFS?}
\addcontentsline{toc}{section}{2--3\quad What are proofs?}

In everyday life when trying to convince a person that something is true, we may
say, ``\emph{This} is true because \emph{that} is true.'' For such a statement to
be convincing, both sides in the argument must agree that the \emph{that} is true
and the \emph{this} must then \emph{necessarily} follow. In other words, there
must be agreement as to what information we start with and how we draw
conclusions from this information. \emph{Logic is the study of the rules for
making correct statements.} Proof always consists of making statements of the
following type: If so and so is true, then such and such must also be true. For
example, consider the statement, ``If the sky is blue (no clouds at all), then it
is not raining.'' Everyone would agree to this statement because from our
experience we know that rain invariably comes from clouds. However, the following
statement, which you might hear in a political speech, is not logically correct.
``The United States is the best country because it is a democracy.'' That we are a
democracy is true; that we are the best country (whatever that means) may be true
although it would not be accepted by all. The logical error comes from the fact
that there are other democratic countries, and hence they too by our careless
reasoning should be best. Therefore, that the United States is best does not
follow from the fact that we are a democracy.

Here are simple examples from algebra:

\begin{center}
``If $3x = 4$, then $x = \tfrac{4}{3}$.''\\[2pt]
``If $a = 2$ and $b = 3$, then $a^2 + ab = 10$.''
\end{center}

Both of these statements are correct and both are of the form, ``If $\al$ then
$\be$.''
\emph{All mathematical proofs employ statements of this type.} We express such a
statement another way by saying

\begin{center}
$\al$ \emph{implies} $\be$.
\end{center}

For example, let $\al$ be the statement ``Today is Monday'' and let $\be$ be the
statement ``Tomorrow is Tuesday.'' Then the statement ``if $\al$ then $\be$'' is
certainly true. However, the statement ``If the sun comes up tomorrow, then it
will not rise the next day'' is very probably false. At least we shall hope that
it turns out to be false. It is fortunate that statements like ``If he takes her
to the dance, then I'll never speak to him again'' and ``If our party loses the
election, then the country will be ruined'' are usually false. It is convenient to
have a symbolic way to indicate that $\al$ implies $\be$. We shall write this as

\begin{center}
$\al \imp \be$,
\end{center}

\noindent where you should read the arrow as the word ``implies.'' We call the
statement ``$\al$ implies $\be$'' an \emph{implication}.

\exercisegroup{2--7}

Put the following statements in an ``if $\al$ then $\be$'' form so that $\al$
implies $\be$.

\begin{multicols}{2}
\begin{exlist}
\item If I eat strawberries I get hives.
\item The picnic will be canceled in case of rain.
\item We will win the game if Jonesevitch plays.
\item Numbers larger than 5 are positive.
\item $x^2 = y^2$ because $x = y$.
\item $3x^2 - 2 = 10$ when $x = 2$.
\item A geometric figure which is a square is not a circle.
\item Squares are rectangles.
\item The area of a rectangle with sides of length $a$ and $b$ is $ab$.
\item With a million dollars, I'd eat steak for breakfast.
\item People who don't study mathematics don't have good sense.
\item When it rains my rheumatism bothers me.
\item When I study I do well in school.
\item People who do not vote are poor citizens.
\item Two lines which intersect are not parallel.
\item I know that $x = 2$ because $3x = 6$.
\item Satisfaction guaranteed or your money back.
\item I know $3x = 6$ because $x = 2$.
\item I will not go to the party if Jimmy is coming.
\item This shirt will not shrink or your money will be refunded.
\item This shirt will shrink or your money will be refunded.
\item To remain free we must vote in elections.
\item The crops will be saved if we have rain this week.
\item The crops will be saved only if we have rain this week.
\item I will go to the dance if Joe asks me.
\item I will go to the dance only if Joe asks me.
\item It will rain tomorrow if it turns cold.
\item It will rain tomorrow only if it turns cold.
\end{exlist}
\end{multicols}

Working with statements of the form ``$\al \imp \be$'' is called \emph{deductive
reasoning} or, in everyday language, \emph{drawing conclusions}. The statement
$\al$ is called the \emph{hypothesis} and the statement $\be$ is called the
\emph{conclusion}. We use this method of reasoning in daily life when we try to
prove to someone that some statement $\be$ is true. We try to find another
statement $\al$ which implies $\be$, and which is true. Then we say, ``$\al$ is
true, therefore $\be$ is true.''

\exercisegroup{2--8}

In Exercises 1 through 7, a statement $\al \imp \be$ is given which is to be
accepted as true. What conclusion can you draw if (a) is true? if (b) is true?

\begin{exlist}
\item If Fred is on the team he must go to bed early.\\
      (a) Fred goes to bed early.\qquad (b) Fred does not go to bed early.
\item If a man is a postman he has foot trouble.\\
      (a) This man has perfect feet.\qquad (b) That man has foot trouble.
\item You will be given a ticket if you are a sophomore.\\
      (a) Jim was not given a ticket.\qquad (b) Joe was given a ticket.
\item You will be given a ticket only if you are a sophomore.\\
      (a) Jean was given a ticket.\qquad (b) George is a sophomore.
\item If two triangles have equal bases and heights then they have equal areas.\\
      (a) These triangles have equal areas.\qquad (b) These triangles do not have
      equal areas and they do have equal heights.
\item The house will burn down if we don't get some water.\\
      (a) We did not get any water.\qquad (b) We did get some water.
\item If a man is beheaded in June then he will die before July 4.\\
      (a) This man died before July 4.\qquad (b) Smith was not beheaded in June.
\end{exlist}

Is the reasoning in Exercises 8 through 18 correct? Invent some other examples
illustrating both correct and incorrect reasoning.

\begin{exlist}\setcounter{exlisti}{7}
\item All men are mortal. Socrates is a man.\\
      \emph{Conclusion:} Socrates is mortal.
\item All widgets are gadgets. This is a gadget.\\
      \emph{Conclusion:} This is a widget.
\item All widgets are gadgets. This is not a gadget.\\
      \emph{Conclusion:} This is not a widget.
\item Only widgets are gadgets. This is a gadget.\\
      \emph{Conclusion:} This is a widget.
\item Only widgets are gadgets. This is a widget.\\
      \emph{Conclusion:} This is a gadget.
\item Every equilateral triangle is isosceles. $\tri{ABC}$ is isosceles.\\
      \emph{Conclusion:} $\tri{ABC}$ is equilateral.
\item Parallel lines do not intersect. Lines $r$ and $s$ intersect.\\
      \emph{Conclusion:} Lines $r$ and $s$ are not parallel.
\item All of these dresses will shrink if washed in hot water. This dress did
      shrink.\\
      \emph{Conclusion:} It was washed in hot water.
\item All sweet food has honey in it. This food is sweet.\\
      \emph{Conclusion:} It has honey in it.
\item All sweet food has honey in it. This food doesn't have honey in it.\\
      \emph{Conclusion:} It is not sweet.
\item All sweet food has honey in it. This food is not sweet.\\
      \emph{Conclusion:} It does not have honey in it.
\end{exlist}

%========================================================= 2-4
\section*{2--4\quad THE TRUTH TABLE FOR IMPLICATION}
\addcontentsline{toc}{section}{2--4\quad The truth table for implication}

If $\al$ and $\be$ are statements, then ``$\al \imp \be$'' should be another
statement. In other words, ``$\al \imp \be$'' must be either true or false. In
that case we must agree on a way of determining whether ``$\al \imp \be$'' is true
or false.

The definition we give may seem a bit strange at first, but keep in mind two
things that may help you. In the first place the definition below is precise. In
the second place we want to use the implication $\al \imp \be$ only when the
hypothesis $\al$ is presumed true.

\begin{definition}
$\al \imp \be$ is true whenever the statement ``(not $\al$) or $\be$'' is true. It
is false when the statement ``(not $\al$) or $\be$'' is false.
\end{definition}

The truth table, then, is as follows.

\begin{truthtable}{cccc}
\tth{\al} & \tth{\be} & \multicolumn{1}{c}{$\al \imp \be$}
          & \multicolumn{1}{c}{(not $\al$) or $\be$} \\[2pt]
\TT & \TT & \TT & \TT \\
\TT & \FF & \FF & \FF \\
\FF & \FF & \TT & \TT \\
\FF & \TT & \TT & \TT \\
\end{truthtable}

The last line may surprise you. Just remember that $\al \imp \be$ means that
$\be$ is true or $\al$ is false and that $\be$ is true \emph{if} $\al$ is true.

\exercisegroup{2--9}

The exercises referred to below are those in Exercise Group 2--8.

\begin{multicols}{2}
\begin{exlist}
\item Write each statement $\al \imp \be$ in Exercises 1 through 7 in the
      ``(not $\al$) or $\be$'' form.
\item Re-examine the validity of the conclusions in Exercises 8 through 18.
\item A statement $\al \imp \be$ is also sometimes expressed as ``$\al$ is
      sufficient for $\be$,'' or ``$\al$ suffices for $\be$,'' or again, ``for
      $\be$ to be true it is sufficient that $\al$ be true.'' Thus for $x^2$ to
      equal 4 it is sufficient that $x = -2$. Express the implications in
      Exercises 1 through 7 in this language.
\item A statement $\al \imp \be$ is also sometimes expressed as ``$\be$ is
      necessary for $\al$'' or ``in order that $\al$ be true it is necessary that
      $\be$ be true.'' Thus for $x$ to equal $-2$ it is necessary that $x^2 = 4$.
      Express the implications in the Exercises 1 through 7 in this language.
\end{exlist}
\end{multicols}


%========================================================= 2-5
\section*{2--5\quad THE CONVERSE OF AN IMPLICATION}
\addcontentsline{toc}{section}{2--5\quad The converse of an implication}

From one implication ``$\al \imp \be$'' we can form others which may or may not be
true. One way of doing this gives what we call the \emph{converse}.

\begin{definition}
The \emph{converse} of ``$\al \imp \be$'' is ``$\be \imp \al$.''
\end{definition}

\emph{Example 1.} $\al \imp \be$: If an animal is a man then it is a human being.\\
\emph{Converse:} $\be \imp \al$: If an animal is a human being then it is a man.\\
In this example ``$\al \imp \be$'' is true but ``$\be \imp \al$'' is false.

\medskip
\emph{Example 2.} $\al \imp \be$: If $\tri{ABC}$ is equilateral then $\tri{ABC}$
is equiangular.\\
\emph{Converse:} $\be \imp \al$: If $\tri{ABC}$ is equiangular, then $\tri{ABC}$
is equilateral.\\
In this example both ``$\al \imp \be$'' and ``$\be \imp \al$'' are true.

\exercisegroup{2--10}

Write the converses of the following statements. In each exercise decide whether
the statement is true or false. Do the same thing for the converse.

\begin{multicols}{2}
\begin{exlist}
\item Every man born in Indiana is an American citizen.
\item If $x - y = 3$, then $x$ is greater than $y$.
\item If two rectangles have equal bases and equal altitudes, they have equal
      areas.
\item If $x^2 = 9$, then $x = 3$ or $x = -3$.
\item If a man is a Polynesian, he is at least six feet tall.
\item If plants are green, they must have had sunshine.
\item If a point is on the perpendicular bisector of a line segment, it is
      equidistant from the endpoints of the segment.
\item If I am sleeping, then I am breathing.
\item If $2 + 4 = 5$, then $4 + 7 = 10$.
\item If $2 + 4 = 5$, then $4 + 7 = 11$.
\item If $2 + 4 = 6$, then $4 + 7 = 10$.
\item If $2 + 4 = 6$, then $4 + 7 = 11$.
\end{exlist}
\end{multicols}

%========================================================= 2-6
\section*{2--6\quad EQUIVALENT STATEMENTS}
\addcontentsline{toc}{section}{2--6\quad Equivalent statements}

When two people give the same mathematical proof they probably will not make
exactly the same statements. But it is very likely that at certain places in their
proofs they will make statements which we say are \emph{equivalent} to one
another. In order to understand precisely what this means we make the following
definition.

\begin{definition}
Two statements $\al$, $\be$ are said to be \emph{equivalent} to one another if
\emph{both} the statement $\al \imp \be$ and the converse $\be \imp \al$ are true.
If $\al$ and $\be$ are equivalent statements, we indicate this by writing
$\al \iffx \be$, which we read as ``$\al$ implies $\be$ and $\be$ implies $\al$.''
\end{definition}

This definition may be difficult to understand. An explanation is that
``equivalent statements present the same information---they may simply be stated
differently.''

An example of a pair of equivalent statements is

\begin{quote}
$\al$: Line $l$ is perpendicular to line $m$.\\
$\be$: Line $m$ is perpendicular to line $l$.
\end{quote}

\exercisegroup{2--11}

In 1 through 7 below decide whether the statements are equivalent to one another.
Invent a few examples of pairs of statements which are equivalent to one another
and some which are not.

\begin{multicols}{2}
\begin{exlist}
\item $\al$: Point $P$ lies on the circle with center at $O$ and radius five
      inches.\\
      $\be$: The length of the line segment $OP$ is five inches.
\item $\al$: Line $l$ is parallel to line $m$.\\
      $\be$: Line $n$ is parallel to line $m$.
\item $\al$: $x$ is greater than $y$.\\
      $\be$: $y$ is less than $x$.
\item $\al$: $2x + 3 = 17$.\\
      $\be$: $x = 7$.
\item $\al$: This geometric figure is a square.\\
      $\be$: This geometric figure is a special kind of parallelogram.
\item $\al$: $\tri{ABC}$ is equilateral.\\
      $\be$: $\tri{ABC}$ is equiangular.
\item $\al$: $x = 3$ or $x = -3$.\\
      $\be$: $x^2 = 9$.
\item[\stex 8.] A statement $\al \iffx \be$ is also sometimes expressed by
      saying, ``In order that $\al$ be true it is necessary and sufficient that
      $\be$ be true.'' In each of the above problems state the equivalence of
      $\al$ and $\be$ in this language.
\end{exlist}
\end{multicols}

Exercises 4 and 7 above illustrate that equivalent statements have an important
place in algebra. See if you can prove that the following statement is true:

\begin{quote}
\emph{If two statements $\al$, $\be$ are equivalent to each other, then either
both are true or both are false.}
\end{quote}

An easy way to prove this is to complete the following truth table. Can you
understand how this table really presents a proof?

\begin{truthtable}{ccccc}
\tth{\al} & \tth{\be} & \multicolumn{1}{c}{$\al \imp \be$}
          & \multicolumn{1}{c}{$\be \imp \al$}
          & \multicolumn{1}{c}{$\al \iffx \be$} \\[2pt]
\TT & \TT & \TT & \TT & ? \\
\TT & \FF & \FF & \TT & ? \\
\FF & \TT & \TT & \FF & ? \\
\FF & \FF & \TT & \TT & ? \\
\end{truthtable}

This truth table must be used together with Definition 2--7. By this definition we
know whether to mark $\al \iffx \be$ as true or false. Studying the table, we see
that according to Definition 2--7, the statement $\al \iffx \be$ must be labeled
true on lines 1 and 4 and labeled false on lines 2 and 3. It should be easy for
you to complete the proof.

\exercisegroup{2--12}

Form statements equivalent to those listed below. If you are not sure of the
equivalence, make a truth table and check your work.

\begin{multicols}{2}
\begin{exlist}
\item $x + y = 7$ and $x - y = 3$.
\item $x$ is greater than $y$.
\item Line segments $AB$ and $CD$ do not intersect.
\item There is a point $P$ which lies on both of the lines $l$ and $m$.
\item Whether you are a smoker or whether you are not, whether you are married or
      single, whether you are a church member or not, you will like Dum-Dums.
\item This geometric figure has four sides, and pairs of opposite sides are
      parallel to one another.
\item That figure is a triangle.
\item This geometric figure consists of two points $A$ and $B$ and all the points
      between $A$ and $B$.
\item The points $A$ and $B$ lie on the same side of the line $l$.
% p.33 prints "opposite lines of the line l" -- kept verbatim; see UNCERTAIN.
\item The points $A$ and $B$ lie on opposite lines of the line $l$.
\end{exlist}
\end{multicols}

%========================================================= 2-7
\section*{2--7\quad THE CONTRAPOSITIVE OF AN IMPLICATION}
\addcontentsline{toc}{section}{2--7\quad The contrapositive of an implication}

A second important implication related to $\al \imp \be$ is its
\emph{contrapositive}.

\begin{definition}
The statement ``not $\be \imp$ not $\al$'' is called the \emph{contrapositive} of
the statement ``$\al \imp \be$.''
\end{definition}

\emph{Examples:}

\begin{enumerate}[label=(\alph*), leftmargin=2.6em, itemsep=3pt, topsep=3pt]
\item $\al \imp \be$: $x = 3 \imp x^2 = 9$.\\
      not $\be \imp$ not $\al$: $x^2 \neq 9 \imp x \neq 3$.
\item $\al \imp \be$: If two lines intersect they are not parallel.\\
      not $\be \imp$ not $\al$: If two lines are parallel they do not intersect.
\item $\al \imp \be$: If I make an $A$ on the examination I will receive an $A$ in
      the course.\\
      not $\be \imp$ not $\al$: If I do not receive an $A$ in the course then I
      will not have made an $A$ on the examination.
\item $\al \imp \be$: If I go to the game I will not go to the movie.\\
      not $\be \imp$ not $\al$: If I go to the movie I will not go to the game.
\end{enumerate}

The most interesting property of the statement ``not $\be \imp$ not $\al$'' is
that it is equivalent to $\al \imp \be$. We can prove this by showing that these
two implications have identical truth tables.

\begin{IIproblem}
Complete the following truth table and show that ``$\al \imp
\be$'' and ``not $\be \imp$ not $\al$'' are equivalent to each other.
\end{IIproblem}

\begin{truthtable}{cccccc}
\tth{\al} & \tth{\be} & \multicolumn{1}{c}{not $\al$}
          & \multicolumn{1}{c}{not $\be$}
          & \multicolumn{1}{c}{$\al \imp \be$}
          & \multicolumn{1}{c}{not $\be \imp$ not $\al$} \\[2pt]
\TT & \TT & \FF & \FF & --- & --- \\
\TT & \FF & \FF & \TT & --- & --- \\
\FF & \TT & \TT & \FF & --- & --- \\
\FF & \FF & \TT & \TT & --- & --- \\
\end{truthtable}

\exercisegroup{2--13}

Form the contrapositive of each statement in 1 through 8 below.

\begin{multicols}{2}
\begin{exlist}
\item If $2 \times 2 = 4$, then oranges are composed of uranium.
\item If you can run the mile in four minutes, then I'll eat my hat.
\item If the ground hog sees his shadow, then we will have a cold spring.
\item If there is no atmosphere on Mars, then there is no life there.
\item If point $P$ is on line $r$, then it is not on line $s$.
\item If you are not satisfied, then your money will be returned.
\item If food is sweet, then it has honey in it.
\item If you always obey traffic laws, then you will not be given a ticket.
\item[\stex 9.] Using the word \emph{necessary}, state the contrapositive of the
      implication of Exercise 4.
\item[\stex 10.] Using the word \emph{sufficient}, state the contrapositive of the
      implication of Exercise 4.
\end{exlist}
\end{multicols}

%========================================================= 2-8
\section*{2--8\quad ``IF AND ONLY IF''}
\addcontentsline{toc}{section}{2--8\quad ``If and only if''}

Consider the two statements ``$\al$ if $\be$'' and ``$\al$ only if $\be$.'' Our
experience in working with statements like these tells us that ``$\al$ if $\be$''
means ``$\be \imp \al$'' and that ``$\al$ only if $\be$'' means ``$\al \imp \be$.''
If both of these statements are true, we have ``$\al \iffx \be$.'' If both of
these statements are true, their conjunction is true.

\begin{definition}
``$\al$ if and only if $\be$'' means ``$\al \imp \be$ and $\be \imp \al$,'' that
is, ``$\al \iffx \be$.''
\end{definition}

Consider the following two statements. ``Two lines are parallel if they do not
intersect.'' ``Two lines do not intersect if they are parallel.'' According to
Definition 2--9 we may write the conjunction of these two statements as ``Two
lines are parallel if and only if they do not intersect.''

\exercisegroup{2--14}

In each exercise below there are two statements. Formulate a single statement,
using ``if and only if,'' that has the same meaning as the conjunction of these
statements.

\begin{multicols}{2}
\begin{exlist}
\item If the weather is bad we will stay home. If we stay home the weather is bad.
\item If you are a sophomore then you will get a ticket. If you get a ticket then
      you are a sophomore.
\item If the weather is bad we will stay home. If the weather is not bad we will
      not stay home.
\item Everyone who receives more than 50\% on this test will pass the course. All
      others will fail.
\item Your dog will obey you if you treat him well. If you do not treat him well
      he will not obey you.
\item I will marry you if you buy me a mink coat. If you don't buy me a mink coat
      I won't marry you.
\item If food is sweet then it has honey in it. Food with honey in it is sweet.
\end{exlist}
\end{multicols}

%========================================================= 2-9
\section*{2--9\quad THEOREMS AND PROOFS}
\addcontentsline{toc}{section}{2--9\quad Theorems and proofs}

In Section 2--3 we said logic is a study of the rules for making correct
statements. The Greek philosopher Aristotle was the first great student of logic.
His work was not improved upon until modern times. In fact, all the logical
concepts we have studied in this chapter belong to what we call
\emph{Aristotelian logic}, which is over 2000 years old and is just as useful
today as it ever was. The logical principles of this chapter are all that you will
need for many years. If you have found these ideas interesting, you may wish to
take courses in logic later on in college. Whether or not you continue your study
of logic, the work of this chapter should help you become careful in both speech
and writing. In particular, during the work in geometry this year, do not tolerate
false statements. Remember that if one false statement is accepted as true then
strange things can be proved. Let us strive then to express ourselves precisely.

Most of our geometrical work will consist of \emph{proving theorems}. Theorems are
statements. To \emph{prove a theorem} is to show that it follows logically from
our assumptions. To do so we employ our principles of logic. For example, \emph{if}
we know that the statement $\al \imp \be$ is true \emph{and if} we know that $\al$
is true, \emph{then} we can infer that $\be$ is true. In symbols,

\begin{center}
$\al$ and $(\al \imp \be) \imp \be$.
\end{center}

Another technique of proof is the following: \emph{If} we know that $\al \imp \be$
is true \emph{and if} we know that $\be$ is false, \emph{then} we can infer that
$\al$ is false. Symbolically,

\begin{center}
not $\be$ and $(\al \imp \be) \imp$ not $\al$.
\end{center}

You need only look at the truth table for $\al \imp \be$ to realize that the two
rules above are correct.

\begin{truthtable}{ccc}
\tth{\al} & \tth{\be} & \multicolumn{1}{c}{$\al \imp \be$} \\[2pt]
\TT & \TT & \TT \\
\TT & \FF & \FF \\
\FF & \TT & \TT \\
\FF & \FF & \TT \\
\end{truthtable}

\begin{IIproblem}
Explain how this truth table justifies the two methods of proof
described above.
\end{IIproblem}

Let us illustrate how these proof techniques may be employed (although you have
already reasoned in this way in many exercises in this chapter).

\medskip
\emph{Example 1.} The notice said, ``If it rains on commencement day then the
graduation ceremony will be held in the gymnasium.'' ($\al \imp \be$.) Today is
commencement day and it is raining. ($\al$ is true.) Therefore, since both of the
statements $\al$ and $\al \imp \be$ are true, we know that $\be$ is true. That is,
the graduation ceremony will be held in the gymnasium.

\medskip
\emph{Example 2.} Jim says that $-3$ is a root of the equation $x^2 - x - 6 = 0$.
If $-3$ is a root of the equation $x^2 - x - 6 = 0$, then $(-3)^2 - (-3) - 6 = 0$.
($\al \imp \be$.) But, $(-3)^2 - (-3) - 6 = 6$. ($\be$ is false.) Now, since
$\al \imp \be$ is true and $\be$ is false, I know that $\al$ is false and Jim was
wrong. The number $-3$ is not a root of the equation $x^2 - x - 6 = 0$.

\medskip
In these examples we would say that we have proved the theorems,

\begin{exlist}
\item The graduation ceremony will be held in the gym.
\item The number $-3$ is not a root of $x^2 - x - 6 = 0$.
\end{exlist}

We might have written out our proofs as follows, in a ``statement-reason'' form.

\begin{steps}
\step{Rain on commencement day $\imp$ ceremony in gym.}{Given.}
\step{It is raining on commencement day.}{Given.}
\step{The ceremony will be in the gym.}{Statements 1 and 2 and laws of logic.}
\end{steps}

\begin{steps}
\step{$-3$ is a root of $x^2 - x - 6 = 0 \imp (-3)^2 - (-3) - 6 = 0$.}{Definition
      of root.}
\step{$(-3)^2 - (-3) - 6 \neq 0$.}{Arithmetic.}
\step{$-3$ is not a root of $x^2 - x - 6 = 0$.}{Statements 1 and 2 and laws of
      logic.}
\end{steps}

\exercisegroup{2--15}

In each problem below you are given certain statements which we shall call true.
Then a theorem is stated. In some cases the theorem may be false; in others there
may not be enough information given to decide whether the theorem is true or
false. Write out proofs. We refer to our given true statements in each problem as
our \emph{hypothesis} (hyp).

\begin{exlist}
\item \begin{hyp}Uncle George will take us skating today if Mother says we may go.
      Mother has said we may go.\end{hyp}
      \begin{thmclaim}Uncle George will take us skating today.\end{thmclaim}
\item \begin{hyp}If there is no oxygen upon the moon then there can be no animal
      life there. Tests have shown conclusively that there is no oxygen upon the
      moon.\end{hyp}
      \begin{thmclaim}There is no animal life on the moon.\end{thmclaim}
\item \begin{hyp}People with green eyes cannot be trusted. Joan has green
      eyes.\end{hyp}
      \begin{thmclaim}Joan cannot be trusted.\end{thmclaim}
\item \begin{hyp}People with blue eyes are honest. Henry is dishonest.\end{hyp}
      \begin{thmclaim}Henry does not have blue eyes.\end{thmclaim}
\item \begin{hyp}If Joe parts his hair on the left then he does not like spinach.
      Joe does not like spinach.\end{hyp}
      \begin{thmclaim}Joe parts his hair on the left.\end{thmclaim}
\item \begin{hyp}Cats with green whiskers have long tails. Cats who do not like
      milk do not have long tails. Simba is a cat who does not like milk.\end{hyp}
      \begin{thmclaim}Simba does not have green whiskers.\end{thmclaim}
\item \begin{hyp}A student who attends M.I.T. will spend more than \$1200 during
      the school year. Kate spent \$980 last year in school.\end{hyp}
      \begin{thmclaim}Kate did not attend M.I.T. last year.\end{thmclaim}
\item \begin{hyp}$x + y = 20$. $x - y = 4$.\end{hyp}
      \begin{thmclaim}$x \neq 11$.\end{thmclaim}
\item \begin{hyp}$2x - 3y = 7$. $x + 2y = 3$.\end{hyp}
      \begin{thmclaim}$3x - y = 10$.\end{thmclaim}
\item \begin{hyp}Joe has 11 marbles: four blue ones and seven yellow ones. Joe
      gave his sister six marbles. Joe did not give all his blue marbles
      away.\end{hyp}
      \begin{thmclaim}Joe does not have more than four yellow marbles
      left.\end{thmclaim}
\item \begin{hyp}Only those students whose test grades average more than 90\% will
      receive $A$'s for their semester's work in geometry. Jim's test grades have
      been 89, 84, 92, 85, and 100. Sally's scores are 91, 90, 89, 94, and
      87.\end{hyp}
      \begin{thmclaim}Sally will receive an $A$ and Jim will not.\end{thmclaim}
\item \begin{hyp}The lovely actress Amanda Thistle has promised to marry her
      honest admirer Silas Moneybags only if he buys her a mink stole, a yacht,
      and the great Calcutta Emerald. Silas has purchased for his Amanda these
      items and many others. Amanda always keeps her promises.\end{hyp}
      \begin{thmclaim}Amanda will wed Silas.\end{thmclaim}
\item \begin{hyp}$x + y = 10$. $x$ and $y$ are positive whole numbers. $x$ is
      greater than $y$. $y + 4$ is greater than $x$.\end{hyp}
      \begin{thmclaim}$y = 4$.\end{thmclaim}
\item \begin{hyp}The square of a positive number is positive. The square of a
      negative number is positive. The square of zero is zero. Zero is neither
      positive nor negative. $k$ is a number. $k^2 = -1$.\end{hyp}
      \begin{thmclaim}There are numbers which are not positive or negative or
      equal to zero.\end{thmclaim}
\item \begin{hyp}The product of two positive numbers is positive. The number $a$
      is positive. The product of numbers $a$ and $b$ is not positive.\end{hyp}
      \begin{thmclaim}The number $b$ is negative.\end{thmclaim}
\item[\stex 16.] \begin{hyp}If $a$ is any number, then $1 \cdot a = a$. If $a$,
      $b$, $c$ are any numbers, then $ba + ca = (b + c)a$. $1 + 1 = 2$.\end{hyp}
      \begin{thmclaim}For every number $a$, it is true that $a + a = 2 \cdot
      a$.\end{thmclaim}
\item[\stex 17.] \begin{hyp}$4 + (-4) = 0$. $6 = 2 + 4$. For all numbers $a$, $b$,
      $c$, it is true that $(a + b) + c = a + (b + c)$. We always have $x + 0 =
      x$.\end{hyp}
      \begin{thmclaim}$6 + (-4) = 2$.\end{thmclaim}
\item[\stex 18.] \begin{hyp}If $8 + x = 0$ then $x = -8$. $8 = 3 + 5$. Numbers can
      be rearranged when we add them. $3 + (-3) = 0$. $5 + (-5) = 0$. $0 + 0 =
      0$.\end{hyp}
      \begin{thmclaim}$(-3) + (-5) = -8$.\end{thmclaim}
\item[\stex 19.] \begin{hyp}If neither of two numbers is equal to zero then their
      product is not zero. The product of the two numbers $a$ and $b$ is zero.
      $b \neq 0$.\end{hyp}
      \begin{thmclaim}$a = 0$.\end{thmclaim}
\item[\stex 20.] \begin{hyp}Dogs with long ears have short tails. Dogs who chase
      rabbits never get measles. Dogs who do not chase rabbits have long tails. My
      dog Rover died of the measles.\end{hyp}
      \begin{thmclaim}Rover did not have long ears.\end{thmclaim}
\end{exlist}

%========================================================= review
\section*{REVIEW OF CHAPTER 2}
\addcontentsline{toc}{section}{Review of Chapter 2}

The ideas in this chapter are not to be memorized as you may have memorized the
addition and multiplication tables when you studied arithmetic. It is very
important that you realize that \emph{all} our complicated mathematical statements
may be built up from simple statements using only the connectives ``and,'' ``or,''
``not,'' and ``if \dots\ then.'' You should relate the precise meaning of other
connectives like ``only if'' and ``if and only if'' to these four fundamental
terms. \emph{Do not} make any effort to memorize the meanings of these latter
terms. You should remember the meaning of conjunction, disjunction, negation,
implication, and equivalence. You should be able to form the converse or
contrapositive of an implication. You should be able to give examples showing that
two statements, ``$\al \imp \be$,'' and the converse, ``$\be \imp \al$,'' are not
necessarily equivalent to each other, that is, that one might be true and the
other false. You should see clearly that the implication ``$\al \imp \be$'' and its
contrapositive, ``not $\be \imp$ not $\al$,'' really assert the same thing; hence
they are equivalent statements. You should understand and be able to use the
important rules of logic employed so often in proofs, namely,

\begin{exlist}
\item $\al$ and $(\al \imp \be) \imp \be$.
\item not $\be$ and $(\al \imp \be) \imp$ not $\al$.
\end{exlist}

You should understand how to set up and use truth tables. For example, to set up a
truth table for the statement ``$\al$ and $\al \imp \be$,'' we could list the
following:

\begin{truthtable}{cccc}
\tth{\al} & \tth{\be} & \multicolumn{1}{c}{$\al \imp \be$}
          & \multicolumn{1}{c}{$\al$ and $(\al \imp \be)$} \\[2pt]
\TT & \TT & \TT & ? \\
\TT & \FF & \FF & ? \\
\FF & \TT & \TT & ? \\
\FF & \FF & \TT & ? \\
\end{truthtable}

\begin{IIproblem}
Fill in the last column in the truth table above and show that
when the statement ``$\al$ and $\al \imp \be$'' is true, then $\be$ is true also.
Make a similar truth table to illustrate the second rule of logic listed above.
Can you invent a rule of logic like these two and show that it is correct by
setting up a truth table?
\end{IIproblem}

\subsection*{Review Exercises}

\begin{exlist}
\item Under what condition is ``$\al$ and $\be$'' a true statement? Give examples.
\item Give an example to show why we call ``$\al$ or $\be$'' true when both $\al$
      and $\be$ are true.
\item What is the negation of not $\al$? Give examples.
\item Make up a statement of the form ``$\al$ and $\be$.'' State its negation.
\item Make up a statement of the form ``$\al$ or $\be$.'' State its negation.
\item What is meant by ``$\al \imp \be$''? Give two statements having this form,
      one of which is true and the other false.
\item Why is the statement ``$\al$ or not $\al$'' true no matter whether $\al$ is
      true or false? Give examples to illustrate this.
\item[\stex 8.]\stepcounter{exlisti} Construct a truth table to show under what condition
      ``($\al$ and $\be$) or $\ga$'' is a true statement. Invent an example of
      this type.
\item[9.]\stepcounter{exlisti} Give an example of an implication ``$\al \imp \be$'' whose converse is
      true and one whose converse is false.
\item Form the converse of:\\
      (a) If a man is from Chicago then he is rich.\\
      (b) If she sings in the choir then I will not go to church.\\
      (c) If Smith is elected president he will get ulcers.\\
      (d) You must be a sophomore if you are taking geometry.\\
      (e) Only the brave are free.
\item Form the contrapositive of each implication in Exercise 10.
\item Write a statement equivalent to:\\
      (a) Sweet food always has honey in it.\\
      (b) My rheumatism never bothers me in warm weather.\\
      (c) In order to be healthy you must drink milk.\\
      (d) A grade of 80\% on the final examination will be sufficient to pass the
      course.\\
      (e) Sophomores may come to the dance, but no one else will be permitted to
      do so.
\item Write the negation of:\\
      (a) Honesty is the best policy.\\
      (b) All math teachers are brutes.\\
      (c) The potato crop will fail and everyone will starve.\\
      (d) Roses are red and violets are blue, sugar is sweet and I love you.\\
      (e) Some cats can swim but no chickens can.
\end{exlist}

There are two statements in each exercise below. Calling the first statement $\al$
and the second one $\be$, rewrite them in the symbolic form ``$\al$ and $\be$,''
``$\al$ or $\be$,'' ``$\al \imp \be$,'' ``$\be \imp \al$,'' ``$\al \imp$ not
$\be$,'' ``$\al \iffx \be$,'' etc., as the case may be. For example, the statement,
``A man is honest if he has big feet'' would be written ``$\be \imp \al$,'' not
``$\al \imp \be$.''

\begin{exlist}\setcounter{exlisti}{13}
\item Vacation begins May 12 but we are not going any place.
\item A man on the faculty must park his car in the main parking lot.
\item You will not receive an $A$ in geometry if you do not work this example
      correctly.
\item Mother will let me go to the dance if and only if I wash the dishes for a
      week.
\item Students who laugh at his jokes are the only ones who get good grades.
\item It is necessary that we have rain this week if the melon crop is to be good.
\item You will write a term paper or you will fail.
\item This is the last exercise and I am glad that it is.
\end{exlist}

\subsection*{ALGEBRA REVIEW}

In the exercises below we indicate some of the ways that logical reasoning is used
in algebra. First we shall list some statements that we agree to call \emph{true}
in algebra. Then we shall set up algebraic problems which provide opportunities to
use these true statements and reason logically to correct conclusions. In what
follows, all letters refer to integers, that is, whole numbers in the set
$(\dots, -3, -2, -1, 0, 1, 2, \dots)$. Let us agree that the following statements
are true:

\begin{exlist}
\item For all integers $a$, $b$, $c$ it is true that if $a + b = a + c$, then
      $b = c$.
\item For all integers $a$, $b$, $c$ it is true that if $ab = ac$ and $a \neq 0$,
      then $b = c$.
\end{exlist}

Let us prove the theorem: \emph{If $x = 5$ then $3x + 4 = 19$.}

\medskip
\noindent\emph{Proof.} If $x = 5$ then $3x + 4 = 3 \cdot 5 + 4 = 15 + 4 = 19$.

\medskip
Let us prove the converse of the above theorem, namely: \emph{If $3x + 4 = 19$
then $x = 5$.}

\medskip
\noindent\emph{Proof.} $3x + 4 = 19 \imp 3x + 4 = 15 + 4$. Now, by statement (1)
above we know that $3x + 4 = 15 + 4 \imp 3x = 15$ and that $3x = 15 \imp 3x =
3 \cdot 5$. Using statement (2) above, we have $3x = 3 \cdot 5 \imp x = 5$. We
could present this proof briefly in the form:

\begin{center}
$3x + 4 = 19 \imp 3x + 4 = 3 \cdot 5 + 4 \imp 3x = 3 \cdot 5 \imp x = 5$.
\end{center}

\subsection*{Exercises}

\begin{multicols}{2}
\begin{exlist}
\item Prove that if $x = 7$ then $4x + 5 = 33$.
\item Prove the converse of the statement in Exercise 1.
\item Prove that if $7x + 9 = 65$ then $x = 8$, and conversely.
\item Prove that $5x + (-3) = 27$ only if $x = 6$.
\item Prove that $x = 11$ is a necessary condition for $5x + 4 = 59$.
\item Prove that if $x = 11$, $5x + 4 = 59$.
\item Prove that $2x + 7 = 19$ if and only if $4x - 3 = 21$.
\item Prove that there is no number $x$ such that $3x + 5 = 11$ and $4x - 1 = 3$.
\item What can you infer if it is known that $a \neq 0$ and $a(3b - 8) = a$?
\item What can you prove if it is known that $r(7b - 5) = r$?
\item It is given that $(3a + 12)(2b - 8) = 0$ and $b \neq 4$. What can be proved?
\item[\stex 12.] It is known that $(2a - 3)(b - 3) = 2a - 2$. Prove that if
      $a = 1$ then $b = 3$. Prove that if $a \neq 1$ then $b \neq 3$. What are the
      values of $a$ and $b$ if $a \neq 1$? (Remember that $a$ and $b$ are whole
      numbers.)
\end{exlist}
\end{multicols}
\clearpage
% ==================== ch03 ====================
\setcounter{theorem}{0}\setcounter{definition}{0}\setcounter{construction}{0}\setcounter{exgroup}{0}\setcounter{footnote}{0}
\setcounter{chapter}{3}

\chapter*{\raggedright\Huge CONCERNING LINES}
\addcontentsline{toc}{chapter}{3.\ Concerning Lines}
\markboth{CONCERNING LINES}{}
\vspace{-1em}\noindent\rule{\linewidth}{0.6pt}\vspace{1em}
\figurekey

%========================================================= 3-1
\section*{3--1\quad INTRODUCTION}
\addcontentsline{toc}{section}{3--1\quad Introduction}

We have seen in Chapter 1 how from our everyday experience it is natural to
imagine the existence of things called points and lines. Even though one cannot
``see'' a point or a line, the study of such imaginary objects turns out to be
practically useful as well as interesting.

Because the objects of study in geometry are inventions of the mind, we must
state very carefully the properties we wish them to possess. Then if we start
with enough information about points and lines, other properties of points and
lines can be \emph{proved}. Many of these provable properties will be surprising,
and not at all obvious at the beginning. It is, of course, impossible to prove
anything about points and lines unless we agree in advance about some properties
that they are to have. These agreed-upon properties, or assumed properties, are
called \emph{postulates}, or \emph{axioms}.* We begin the proper study of plane
geometry in this chapter when we list our beginning postulates in Section 3--3.
The postulates should be simple enough to seem almost obvious and yet must be
sufficient for us to prove geometric facts that seem to be true in our everyday
life.%
\IIIstarfn{~The word ``postulate'' is derived from a Latin word meaning
``demand.'' ``Axiom'' comes from a Greek word meaning ``worthy.''}

A different method of procedure would be to list all the known geometric facts
and assume that they are true. Then when we had memorized all these facts, our
study of geometry would be completed. There are at least three objections to
such a plan. In the first place, we might not have listed all the facts needed
and we would have no experience in finding new ones. In the second place, it
would be a terrible strain upon the memory. To avoid memorizing everything, we
need to study how different facts are related and how the more complicated ones
depend upon other simpler ones. When we prove theorems from postulates or other
theorems, we are studying the way the geometric facts are related to each other;
this helps us to remember them. In the third place, we would lose the great
amount of pleasure that we can get out of proving from very simple and obvious
assumptions that certain other things are true. It is because of this pleasure
that geometry has been so attractive to men since the time of Euclid. The fact
that all geometry can be established from a small list of postulates gives to
geometry a kind of beauty that has appealed to men as different as Abraham
Lincoln and the Greek philosopher Plato.

%========================================================= 3-2
\section*{3--2\quad STRAIGHTNESS}
\addcontentsline{toc}{section}{3--2\quad Straightness}

Soon, when we list our postulates, we shall be talking about \emph{lines}, and of
course what we have in mind are things which are like the familiar \emph{straight}
lines of our experience. But we will call them simply \emph{``lines.''} There is
really nothing to be gained by putting in the extra word ``straight'' because it
would add nothing to our use of the postulates. It will turn out much later that
our lines \emph{are} straight in the sense that any path between two points,
other than a line, would be longer than the line.

When we draw pictures of lines we will, of course, draw them to \emph{appear}
straight, for this will help us see how the points and lines are related. But we
will not be concerned with the straightness of the lines because the figures are
meant to be only a guide to the proofs and to help us see what is going on.
\emph{The proofs are all completely independent of the figures.} By this we mean
that we must not say something is true because of the appearance of a figure;
rather, we must give reasons for its truth. The things that are true in our
geometry are only those things which can be proved by our use of our assumptions
and laws of logic. But as we progress in our study of geometry we will gradually
use figures more and more freely when it is quite clear how to fill in a complete
argument.

To understand why it is that we are not concerned with straightness at all, at
least to begin with, we shall look at a plane in two different ways. First,
imagine a sheet of paper with lines drawn on it so that some of the lines might
look like Fig.~3--1. Now suppose you have a pane of uneven glass between your
eyes and the paper, so that when you look at the drawing it looks like
Fig.~3--2. Note that in looking at the drawing, with or without the glass, you
are seeing the same lines. In one case they appear straight; in the other case
they do not. In the first figure $A$ and $B$ may appear to be the same distance
apart as $A'$ and $B'$, but this may not be so in the second figure. Therefore,
\emph{if it were possible to see the plane only through the glass, you would have
to be told which sets of points are ``straight'' lines, and whether or not $A$
and $B$ are the same distance apart as $A'$ and $B'$.}

\FIGIIIONETWO

In our study of geometry we shall be talking about imaginary things---points and
lines. We cannot really \emph{see} our lines. Hence we shall not call them
``straight'' but refer to them only as \emph{lines}. All the properties that we
need to know about these lines will be stated clearly in the postulates.

%========================================================= 3-3
\section*{3--3\quad THE POSTULATES}
\addcontentsline{toc}{section}{3--3\quad The postulates}

We imagine that there is a \emph{set} (or ``bunch,'' or ``collection'') of things
which we call \emph{points}. This set of points will be called \emph{the plane}.
In the plane certain other sets of points are called \emph{lines}. We represent
points by capital letters and lines by small letters. When a point $P$ is one of
the points of the line $l$ (Fig.~3--3), we say that ``$P$ is on $l$'' and also
that the line ``goes through $P$,'' or ``passes through $P$,'' or ``$l$ contains
$P$.'' If two lines $l$ and $m$ have a point $P$ in common (Fig.~3--4), that is,
if $P$ is on both $l$ and $m$, then we say that the lines \emph{intersect}.
Sometimes we will label two points with the same letter but different
\emph{subscripts}, such as $P_1$ and $P_2$. We read $P_1$ as ``$P$ sub-one'' and
$P_2$ as ``$P$ sub-two.''

\FIGIIITHREEFOURFIVE

The postulates will be divided into six groups, of which only the first group
will be given in this section. We shall give the other groups as we need them.
It is extremely important at the beginning of your study of geometry that you
understand the relationship between our mathematical postulates and the real
world. We decide \emph{what} we are going to assume about the imaginary points
and lines of our geometry by \emph{looking} at the real world. As a matter of
fact, in order to choose the postulates for our geometry, we need only to look at
some simple drawings. For example, the postulates in our first group are
suggested by Fig.~3--5. Study this figure to see if you can decide what these
postulates will be.

Of course, after we have decided what postulates to use, we must prove our
theorems by \emph{using these postulates and not by looking at other figures}.

\medskip
\noindent\textbf{Postulate Group I.}\quad This group contains three postulates,
indicated by I--1, I--2, and I--3.

\medskip
\noindent I--1 \emph{The three-point postulate}. There are at least three points
which do not all lie on one line.

\medskip
\noindent I--2 \emph{The point-line postulate}. For any two different points
there is exactly one line containing these points.

\medskip
\noindent I--3 \emph{The line-point postulate}. Every line contains at least two
points.

\medskip
The point-line postulate tells us one way we want lines to behave. The
line-point postulate \emph{starts} us on the way to getting points on a line. The
three-point postulate tells us that not all points are on one line. It gets us
off the line.

With only these postulates we cannot as yet prove very much. The following two
theorems are about all that can be formulated.

\begin{theorem}
Two different lines do not intersect in more than one point.
\end{theorem}

\noindent\emph{Proof.} The theorem is either true or false. We shall prove that
it cannot be false. Hence it must be true. We therefore suppose that the lines
intersect in at least two points and then show that this is impossible. Suppose
the lines are $l$ and $m$ and that they intersect in the two points $A$ and $B$
(Fig.~3--6). We would then have two different lines, $l$ and $m$, through $A$ and
$B$. This contradicts the point-line postulate. Therefore our assumption that
$l$ and $m$ intersect in more than one point is false and our theorem is true.

\begin{remark}
Several ideas from logic are involved in this proof. Comment on them.
\end{remark}

\begin{theorem}
There are at least three lines in the plane.
\end{theorem}

\FIGIIISIXSEVEN

\noindent\emph{Proof.}
\begin{steps}
\step{There are at least three points in the plane not all on one line. Call
      them $A$, $B$, and $C$.}{Three-point postulate.}
\step{Let $l$ be the line through $A$ and $B$, $m$ the line through $B$ and $C$,
      and $n$ the line through $C$ and $A$.}{Point-line postulate.}
\step{The three lines $l$, $m$, and $n$ are all different from one
      another.}{If any two were the same line, all three points $A$, $B$, and
      $C$ would be on one line, contrary to the way they were chosen.}
\end{steps}

The three points and lines are shown in Fig.~3--7. Of course, it looks just like
the drawing (Fig.~3--5) that suggested our first postulates.

\begin{remark}
At this point we cannot prove from our postulates that there are any more points
or lines than shown in Fig.~3--7.
\end{remark}

\exercisegroup{3--1}
\begin{multicols}{2}
\begin{exlist}
\item Restate the point-line and the three-point postulates in other words.
\item Prove that if $l$ is any line, there is at least one point not on $l$.
\item Explain why both $a$ and $b$ in Fig.~3--8 cannot represent lines if $P$ and
      $Q$ represent points.
\item Prove that if $l$ is a line and $P$ is on $l$ then there is another point
      on $l$.
\item Why can we not prove there are more than three lines in our plane?
\item State what ``intersect'' means.
\item Suppose that $l$ is a line. List all the things you can say about this line
      and points of the plane. Give your reasons.
\item Given Fig.~3--9, what can you conclude about $l$ and $m$? Why?
\item A line $l$ contains the points $P$ and $Q$. The line $m$ also contains $P$
      and $Q$. What can you conclude?
\item Two lines have a point $A$ in common. What can you say about the lines?
\item Prove that there are at least two lines on each point.
\item Why, when looking at Fig.~3--10, can we not prove that there are at least
      five points in the plane?
\end{exlist}
\end{multicols}

\FIGIIIEIGHT

\FIGIIININE

\FIGIIITEN

%========================================================= 3-4
\section*{3--4\quad POSTULATES OF BETWEENNESS}
\addcontentsline{toc}{section}{3--4\quad Postulates of betweenness}

We come now to our assumptions about the order, or arrangement, of points on a
line. It is interesting that Euclid did not state these assumptions, but used
them, without comment, wherever necessary. It was the German mathematician
Moritz Pasch (pronounced ``posh'') who first, in the latter part of the 19th
century, made a study of the properties of betweenness actually needed in
geometry. The postulates given below are a modification of those of Pasch. You
will have observed that the postulates of Group I seemed very obvious, and
perhaps those of Group II will also. But remember that we are trying to state
clearly what we imagine to be true about a plane, and hence we must be careful to
state \emph{everything} we are going to use.

\medskip
\noindent\textbf{Postulate Group II. Betweenness.}\quad There is a relation,
called \emph{between}, such that one of three points is sometimes between the
other two. Of course, if we consider \emph{any} three points, it does not
\emph{have} to happen that any one is between the others. But if for three points
$A$, $B$, and $C$ it is true that $B$ is between $A$ and $C$, we shall indicate
this by writing ``$ABC$'' or ``$CBA$,'' that is, with $B$ in the middle.

\medskip
\noindent II--1 \emph{The ABC betweenness postulate}. If $ABC$, that is, if $B$
is between $A$ and $C$, then $A$, $B$, and $C$ are three different (distinct)
points on a line. (A drawing of the line and points might look like Fig.~3--11.)

\FIGIIIELEVENTWELVE

\medskip
\noindent II--2 \emph{The three-point betweenness postulate}. For every three
points on a line, exactly one of them is between the other two. That is, if $R$,
$S$, and $T$ are on one line, then exactly one of the statements $RST$, $STR$,
and $TRS$ is true.

\medskip
\noindent II--3 \emph{The four-point betweenness postulate}. Any four points on a
line may be labeled $A_1$, $A_2$, $A_3$, $A_4$ so that the only betweenness
relations are the same as the order of the subscript numbers. That is, the only
betweenness relations are $A_1A_2A_3$, $A_1A_2A_4$, $A_1A_3A_4$, $A_2A_3A_4$. Of
course these may all be reversed as $A_3A_2A_1$, $A_4A_2A_1$, $A_4A_3A_1$,
$A_4A_3A_2$. A similar statement is to hold for any five points on a line, and in
fact for any finite number of points on a line.

\medskip
\noindent II--4 \emph{The line-building postulate}. If $A$ and $B$ are two
points, then there is at least one point $C$ such that $ABC$, and at least one
point $D$ such that $ADB$. (See Fig.~3--12.)

\medskip
Before stating the last of the betweenness postulates, we need some definitions.

\begin{definition}
Two points not on a line $l$ are said to be \emph{on the same side of $l$} if no
point of $l$ is between the two points. (See Fig.~3--13.) Two points not on a
line $l$ are said to be \emph{on opposite sides of $l$} if there is a point of
$l$ between the two points. (See Fig.~3--14.)
\end{definition}

\FIGIIITHIRTEENFOURTEEN

\medskip
\noindent II--5 \emph{The two-side postulate}. Each line separates the plane. By
this we mean that all points of the plane not on the line are divided into two
sets, called the two sides of the line, such that if $P$ and $Q$ belong to one of
these sets, then $P$ and $Q$ are on the same side of the line, whereas if $P$
belongs to one set and $Q$ belongs to the other, then $P$ and $Q$ are on opposite
sides of the line. (Briefly, this says that a line has two sides.)

\medskip
Before using these postulates, let us first discuss their meaning in everyday
life. If three men are standing in a line (Fig.~3--15) and $A$ cannot see $C$
because $B$ is in the way, then we say that $B$ is between $A$ and $C$. On the
other hand, $B$ and $C$ can see each other perfectly. Hence $A$ is not between
$B$ and $C$. That is, we do not have $BAC$.

\FIGIIIFIFTEEN

The $ABC$ betweenness postulate says that to talk about betweenness, we must have
the men ``lined up.'' The three-point betweenness postulate says that when three
men are lined up, exactly one of them is between the other two.

The four-point betweenness postulate gives a definite meaning to what we expect
to see when we say, ``Line up these four people.'' It tells exactly who is
between whom. An easy way to remember the content of this postulate is to observe
that of four points on a line, two of them are never between any others. These
are called the \emph{endpoints}. The other two points are between points exactly
twice.

The line-building postulate tells us first that there is always a point $C$
``beyond'' $B$ from $A$. In fact, Euclid stated (vaguely) something very much
like this as a postulate. But our postulate also tells us that there are points
between any two. Euclid never stated this. It is this postulate that gives the
existence of more than two points on a line, and therefore it is called the
line-building postulate. You might think about the relation between this
postulate and the real world. When you draw pictures of points, can you always
mark a picture of a third point between any two?

The two-side postulate says that a line has two sides. It merely expresses the
way we tell whether or not two of us are on the same or different sides of a
fence. If, when we look at each other, we see the fence between us, we know we
are on opposite sides of it. If we do not see the fence between us, we are on the
same side.

Now, using these postulates, we can define some familiar geometric objects.

\begin{definition}
If $A$ and $B$ are any two points, then the set of points consisting of $A$ and
$B$ and all the points between $A$ and $B$ will be called the \emph{line segment}
(or \emph{segment}) determined by $A$ and $B$. The segment will be indicated by
``$AB$.'' The points $A$ and $B$ are called the \emph{ends} of the segment. In
Fig.~3--16, $P$ is a point of the segment $AB$.
\end{definition}

\FIGIIISIXTEENSEVENTEEN

\begin{definition}
If $A$, $B$, and $C$ are three points not all on one line,* the set of points of
the segments $AB$, $BC$, $AC$ is called a \emph{triangle} (Fig.~3--17). We call
the points $A$, $B$, and $C$ \emph{vertices} of the triangle and the segments
$AB$, $AC$, $BC$ \emph{sides} of the triangle. To indicate ``triangle $ABC$,''
we write ``$\tri{ABC}$.''
\end{definition}
\IIIstarfn{~Points all lying on a line are said to be \emph{collinear}. Hence in
the definition we could have said, ``If $A$, $B$, and $C$ are not collinear
\ldots{}''.}

Before working the exercises, read all the postulates again to be sure that you
understand them.

\exercisegroup{3--2}
\begin{multicols}{2}
\begin{exlist}
\item Draw a line and mark five points on it. Write out all the betweenness
      relations.
\item If $ABC$ and $ACD$, prove that the four points $A$, $B$, $C$, and $D$ lie
      on one line.
\item Give in your own words the definitions of ``segment,'' ``triangle,'' and
      ``two points are on the same side of a line.''
\item If $A$, $B$, $C$, and $D$ are four points of a line and we know that $CAD$
      and $ADB$, what are the other betweenness relations? Draw a picture. Is the
      picture proof?
\item[\stex 5.] If $A$, $B$, $C$, and $D$ are four points of a line and we know
      that $CAD$ and $CAB$, what are the other betweenness relations?
\item[6.] If $E$, $F$, $G$, and $H$ are four points on a line and we know that
      $EFG$, how many possible positions are there for $H$ that can be described
      in terms of our betweenness relations?
\end{exlist}
\end{multicols}

\begin{exlist}[start=7]
\item $A$ and $B$ are points on a line $l$. Write three true statements about
      points either on $l$ or not on $l$. Give reasons, as below.
\end{exlist}

\begin{steps}
\step{There is a point $C$ such that $ACB$.}{The line-building postulate.}
\step{\rule{0.92\linewidth}{0.4pt}}{\rule{0.92\linewidth}{0.4pt}}
\step{\rule{0.92\linewidth}{0.4pt}}{\rule{0.92\linewidth}{0.4pt}}
\step{\rule{0.92\linewidth}{0.4pt}}{\rule{0.92\linewidth}{0.4pt}}
\end{steps}

\begin{multicols}{2}
\begin{exlist}[start=8]
\item If $RST$ and $VRS$, relabel the points $R$, $S$, $T$, and $V$ as $A_1$,
      $A_2$, $A_3$, and $A_4$ (in some order) in accordance with II--3. If
      $CAD$, $BDE$, and $BCA$, relabel the points $A_1$, $A_2$, $A_3$, $A_4$, and
      $A_5$ in accordance with II--3.
\item Prove that there are at least four lines on each point.
\item If $DEF$ and $EGF$, then the four-point betweenness postulate implies that
      \rule{2.2em}{0.4pt} and \rule{2.6em}{0.4pt}.
\item If $RST$, then what postulate tells us that $R$, $S$, and $T$ are distinct
      points on a line?
\item[\stex 12.] Draw a line and choose two points on it, $A$ and $B$ (Postulate
      I--3). Mark a point $C$ such that $ABC$ (Postulate II--4). Mark a point $D$
      such that $BCD$ (Postulate II--4). Is it true that $ACD$? Why? Use II--3 to
      prove that $ACD$.
\item[\stex 13.] Draw a line and choose two points on it, $A$ and $B$. Choose a
      point $C$ such that $ACB$ (Postulate II--4). Now choose a point $D$ such
      that $CDB$. Is it true that $ACD$? that $ADB$? Use II--3 to prove this.
      Now choose a point $E$ such that $DEB$. What are all the betweenness
      relations between these five points? Can this process be indefinitely
      repeated? Describe, in your own words, the meaning of ``this process.'' How
      many points are there on a line segment?
\item[\stex 14.] Draw a line and choose a point $A$ on one side and a point $B$ on
      the other side. Now mark a point $C$ on the opposite side of the line from
      $A$. Points $B$ and $C$ will be on the same side of the line. How does this
      follow from II--5?
\item[\stex 15.] A line $l$ (Fig.~3--18) does not contain any of the vertices of
      a triangle. The line $l$ does intersect the side $AB$, but does not
      intersect the side $BC$. Prove that $l$ must intersect $AC$.
\item[\stex 16.] Draw a line and choose two points $A$ and $B$ on one side of the
      line. Choose $C$ so that $ACB$. Can you prove that $C$ and $A$ are on the
      same side of the line?
\item[\stex 17.] If $A$ is on line $l$ and $B$ is not on $l$ and $ACB$, prove
      that $C$ and $B$ are on the same side of $l$.
\item[\stex 18.] Prove that for every line $l$ there is a point on each side of
      $l$.
\end{exlist}
\end{multicols}

\FIGIIIEIGHTEEN

%========================================================= 3-5
\section*{3--5\quad RAYS AND ANGLES}
\addcontentsline{toc}{section}{3--5\quad Rays and angles}

It is a familiar fact from experience that a point $O$ on a line $l$ divides the
other points of the line into two sets---those on one side and those on the other
side of the point $O$ on the line $l$. (See Fig.~3--19.) From our postulates on
betweenness we can now make precise definitions of this observation. First we
need to know that a point on a line separates the line.

\FIGIIININETEEN

\begin{theorem}
{\normalfont The two-side theorem.} Let $O$ be any point of a line $l$. The
points of $l$, different from $O$, are separated into two sets such that:
(1) If $P$ and $Q$ are in the same set, then either $OPQ$ or $OQP$.
(2) If $P$ is in one set and $Q$ is in the other set, then $POQ$.
\end{theorem}

\noindent*\emph{Proof.} Choose a point $A$ not on $l$ and consider the line $m$
through $O$ and $A$ (see Fig.~3--20). We are using I--1 and I--2. By the two-side
postulate, II--5, the line $m$ separates the plane into the two sides of $m$. Let
one of the sets, which we seek, consist of the points of $l$ which are on one
side of $m$, and let the other set consist of the points of $l$ on the other side
of $m$.%
\IIIstarfn{~Starred proofs are difficult and may be omitted. However, the meaning
of the theorems must be understood.}

If $P$ and $Q$ are in one of these sets, they are by definition on the same side
of $m$, and therefore $O$ is not between them. Hence by II--2 we must have either
$OPQ$ or $OQP$. This proves (1).

Now suppose that $P$ is in one set and $Q$ in the other set (see Fig.~3--21).
Then by definition of the sets, $P$ and $Q$ are on different sides of $m$. By the
two-side postulate, II--5, there is a point of $m$ between $P$ and $Q$. But $l$
and $m$ have only point $O$ in common. Hence the point between must be $O$, and
we have $POQ$. This proves (2).

\FIGIIITWENTYTWENTYONE

\begin{remark}
The two sets of the theorem are not dependent on how we choose the line $m$,
because we end up with the same betweenness relations.
\end{remark}

\begin{definition}
The two sets of Theorem 3--3 are called the two \emph{sides} of $O$ on the line
$l$.
\end{definition}

\begin{definition}
A \emph{ray} from $O$, or a ray with end $O$, is a set of points consisting of
$O$ and all the points on one side of $O$ of a line $l$ through $O$ (see
Fig.~3--22).
\end{definition}

If $P$ is any point of the ray, we may speak of ``ray $OP$.'' Note that we may
describe the ray as the set of points consisting of $O$ and $P$ and all points
$Q$ such that $OQP$ and all points $R$ such that $OPR$. Refer back to
Theorem 3--3.

\FIGIIITWENTYTWOTWENTYTHREE

\exercise{3--3}
Consider a line $l$ and a ray $r$, not on $l$, from a point $P$ on $l$, as in
Fig.~3--23. What do we mean by ``$r$'s side of $l$''?

\begin{definition}
The set of all points on two rays from a point is called an \emph{angle}. The
point is called the \emph{vertex} of the angle, and the two rays are called the
\emph{sides} of the angle. If the two rays are on the same line, the angle is
called a \emph{straight angle}. If $O$ is the vertex and $A$ and $B$ are points
other than $O$ on the sides of the angle, we shall speak of ``angle $AOB$'' and
write this as ``$\ang{AOB}$.'' (See Fig.~3--24.)
\end{definition}

\FIGIIITWENTYFOURTWENTYFIVE

\begin{remark}
Note that what we call the angle is the set of points in the \emph{sides} of the
angle. Points ``between'' the rays are \emph{not} part of the angle. Later we
discuss the inside and the outside of an angle. It will be convenient to have a
way to clearly indicate the sides of the angles. This can be done as in
Fig.~3--25(a), and the angle can be numbered. Since the arrowheads indicating the
sides are not vital, we can as well use the simpler notation shown in
Fig.~3--25(b) and (c). \emph{Remember that this notation is just a way of
indicating the sides of the angle.}
\end{remark}

\begin{definition}
If two lines $l$ and $m$ intersect at $O$ and the four rays from $O$ are as
indicated in Fig.~3--26, then angles 1 and 2 are said to be \emph{vertical} with
respect to each other. For brevity, we say that 1 and 2 are a pair of vertical
angles. Angles 3 and 4 are also vertical to each other.
\end{definition}

\FIGIIITWENTYSIX

\begin{definition}
If two angles have a common vertex and a common side and if the other sides are
on opposite sides of the line on the common side, then the angles are said to be
\emph{adjacent} to each other (Fig.~3--27).
\end{definition}

\FIGIIITWENTYSEVEN

\begin{definition}
Two adjacent angles whose noncommon sides form a straight angle are said to be
\emph{supplementary and adjacent} to each other (Fig.~3--28).
\end{definition}

\FIGIIITWENTYEIGHT

\exercisegroup{3--4}
\begin{multicols}{2}
\begin{exlist}
\item State in your own words the definitions of ``ray,'' of ``angle,'' and of
      ``straight angle.''
\item How ``long'' are the sides of an angle?
\item If $\ang{AOB}$ and $\ang{DOC}$ are vertical to each other with $B$, $O$,
      $D$ collinear, what kind of angle is $\ang{AOC}$?
\item If $\ang{AOC}$ is a straight angle and $OR$ is a ray from $O$, how are
      angles $AOR$ and $ROC$ related?
\item Two lines intersect at $O$. With the four rays from $O$, how many different
      angles with vertex $O$ can be found? Pick out the pairs of vertical angles.
\item Three lines intersect at $O$. With the six rays from $O$, how many
      different angles with vertex $O$ can be formed? How many pairs of vertical
      angles are there?
\item Draw two angles with a common vertex but which are not adjacent to each
      other.
\item Draw two angles with a common vertex and side but which are not adjacent.
\item State Definitions 3--7, 3--8, and 3--9 in your own words.
\item If $\ang{AOB}$ is not a straight angle, how many angles are there that are
      vertical with respect to $\ang{AOB}$?
\item Why is it not sensible to say, ``$\ang{AOB}$ is a vertical angle''?
\end{exlist}
\end{multicols}

%========================================================= 3-6
\section*{3--6\quad INTERIOR AND EXTERIOR}
\addcontentsline{toc}{section}{3--6\quad Interior and exterior}

When we draw a picture of an angle (not a straight angle) on a sheet of paper, we
can tell by looking at it which points of the plane are inside and which are
outside the \emph{picture} of the angle. But in our development of geometry
starting from but a few postulates, we must \emph{define} what these sets are.
(Why?) Definition 3--10 is one that can be understood by a blind person provided
he knows what an angle is and that a line separates the plane. To appreciate the
definition, consider Fig.~3--29. The point $D$ lies on one side of $l$, which
side has been shaded with vertical lines. The point $C$ lies on one side of $m$,
which side has been shaded with horizontal lines. The points in the plane which
are shaded twice (i.e., where the vertical and horizontal lines cross) fill up
the \emph{interior} of $\ang{COD}$.

\FIGIIITWENTYNINETHIRTY

\exercisegroup{3--5}
\begin{multicols}{2}
\begin{exlist}
\item Copy the lines $l$ and $m$ and the points $A$, $B$, $C$, $D$ on a sheet of
      paper and shade all the points of the plane which are (a) on the same side
      of $l$ as $B$ and the same side of $m$ as $A$, (b) on the same side of $l$
      as $B$ and on the same side of $m$ as $C$, (c) on the same side of $l$ as
      $D$ and on the same side of $m$ as $A$.
\item In each case name the angle whose interior has been described.
\end{exlist}
\end{multicols}

\begin{definition}
Consider an angle, not a straight angle, with vertex $O$ and whose sides are the
rays $r_1$ and $r_2$ on the lines $l_1$ and $l_2$, respectively. The
\emph{interior} of the angle is the set of all points on the same side of $l_1$
as $r_2$ and on the same side of $l_2$ as $r_1$. The \emph{exterior} of the angle
is the set of all points in the plane which are not points of the angle and are
not in the interior of the angle. (Fig.~3--30.)
\end{definition}

There are several properties of the interior and exterior of angles that we can
now prove. These properties are all quite obvious when we look at a figure. But
even when something seems obvious from a picture, we must be sure that we have
the right postulates so that our geometry agrees with our feelings about what
should be true---that is, with the picture. It is really not at all obvious that
our \emph{geometry} agrees with our \emph{picture}. Let us see that this is so by
proving some ``obvious'' theorems.

In all these theorems we shall be concerned with an angle whose vertex is $O$,
which is not a straight angle, and whose sides are the rays $r_1$ and $r_2$.

\begin{IIIsthm}
If $P$ and $Q$ are both in the interior of $\ang{O}$, then the segment $PQ$ is in
the interior of $\ang{O}$.
\end{IIIsthm}

\noindent\emph{Proof.} We must show that if $A$ is any point between $P$ and $Q$,
then $A$ is in the interior of the angle (Fig.~3--31). This is easy if we use the
result of Exercise 16 in Exercise Group 3--2. By that exercise, since $P$ and $Q$
are on one side of $l_1$, then $A$ is also on that same side of $l_1$. Similarly,
$A$ is on the same side of $l_2$, as are both $P$ and $Q$. Hence, by the
definition of \emph{interior} of an angle, $A$ is in the interior of the angle
formed by rays $r_1$ and $r_2$.

\FIGIIITHIRTYONE

\begin{IIIsthm}
If $r$ is a ray from $O$ other than $r_1$ and $r_2$, then either all the points
of $r$ other than $O$ lie in the interior of $\ang{O}$ or they lie in the
exterior of $\ang{O}$.
\end{IIIsthm}

\noindent\emph{Proof.} Let $r_1'$ and $r_2'$ be the rays from $O$ on $l_1$ and
$l_2$ other than $r_1$ and $r_2$. (See Fig.~3--32.) If $r = r_1'$ or
$r = r_2'$, then $r$ is in the exterior. If $r$ is neither $r_1'$ nor $r_2'$,
then $r$ lies on a definite side of $l_1$ and a definite side of $l_2$ because of
the way a ray is defined. Hence $r$ is in either the interior or exterior of
$\ang{O}$.

\FIGIIITHIRTYTWO

\begin{IIIsthm}
If $P$ is in the interior of $\ang{O}$ and $Q$ is in the exterior of $\ang{O}$,
then the segment $PQ$ contains a point of the angle.
\end{IIIsthm}

\noindent\emph{Proof.} There are five possible locations of the point $Q$,
represented by the points $Q_1, \ldots, Q_5$ in Fig.~3--33. (Describe these
positions in terms of sides of lines.) For positions $Q_1$, $Q_2$, $Q_4$, $Q_5$
the proofs of the theorem are quite easy and are assigned as exercises below.
Hence we restrict our attention here to the position $Q_3$.

\FIGIIITHIRTYTHREE

Since $Q_3$ is on opposite sides of both $l_1$ and $l_2$ from $P$ (why?), the
segment $Q_3P$ must contain a point $R$ of $l_1$ and a point $S$ of $l_2$. We
assert that either $R$ or $S$ is a point of the angle, that is, a point of either
$r_1$ or $r_2$. If $R$ is on $r_1$, our segment contains a point of the angle.
Suppose that $R$ is not on $r_1$. Then $R$ is on $r_1'$ and hence has the same
relative position as $Q_4$. (Why?) But then the segment $RP$ contains a point of
the angle which must be $S$. (Why?) Since $Q_3RP$ and $RSP$, we have $Q_3SP$.
(Why?)

\begin{definition}
The \emph{interior} (or \emph{inside}) of a triangle is the set of all points
which are interior to each of the angles of the triangle. The \emph{exterior} of
a triangle is the set of points which are neither on the triangle nor in the
interior.
\end{definition}

If the triangle is $ABC$ (see Fig.~3--34), an alternative definition of interior
would be ``the set of all points on the same side of $AB$ as $C$, on the same
side of $BC$ as $A$, and on the same side of $CA$ as $B$.'' Show that these two
definitions say the same thing.

We can prove the following remark by using Theorem 3--6. The remark is an
``obvious'' theorem that is hard to prove. We will not need this result later.

\begin{remark}
If $P$ is in the interior of $\tri{ABC}$ and $Q$ is in the exterior, then the
segment $PQ$ contains a point of the triangle (Fig.~3--35).
\end{remark}

\FIGIIITHIRTYFOURTHIRTYFIVE

\exercisegroup{3--6}
\begin{multicols}{2}
\begin{exlist}
\item Referring to Fig.~3--36, draw an angle with vertex $O$. Pick a point $A$ on
      one side of $\ang{O}$. Draw an angle with vertex $A$ with one side the ray
      $AO$ and the other side a ray intersecting the other side of $\ang{O}$ in
      $B$. Shade the interior of $\ang{O}$ and the interior of $\ang{A}$. Are
      points in both interiors also in the interior of $\tri{OBA}$?
\item Draw a $\tri{ABC}$. If a point $P$ is on the same side of $AB$ as $C$ and
      on the same side of $BC$ as $A$, show that $P$ is not necessarily on the
      same side of $AC$ as $B$.
\item Draw two angles whose interiors have no points in common.
\item A straight angle does not have an interior. Tell why we do not try to
      define interior for a straight angle.
\item State in your own words the definitions of ``interior of an angle'' and
      ``interior of a triangle.''
\item For a $\tri{ABC}$, the point $P$ is in the interior of $\ang{A}$ and not in
      the interior of $\ang{B}$. Can you be sure that $P$ is in the exterior of
      $\tri{ABC}$?
\item Consider the angle of Fig.~3--37, with sides $l_1$ and $m_1$ indicated on
      lines $l$ and $m$. Tell where the point $P$ is if:
      \begin{itemize}[label=,leftmargin=1.1em,itemsep=0pt,topsep=1pt]
      \item (a) $P$ is on the opposite side of $m$ from $l_1$.
      \item (b) $P$ is on the opposite side of $l$ from $m_1$.
      \item (c) $P$ is on $l_1$.
      \item (d) $P$ is on $m_1$.
      \item (e) $P$ is on the same side of $l$ as $m_1$.
      \item (f) $P$ is on the same side of $l$ as $m_1$ and on the same side of
            $m$ as $l_1$.
      \end{itemize}
\item Prove that the following definition is equivalent to Definition 3--11: The
      interior of $\tri{ABC}$ is the set of all points in the interior of
      $\ang{A}$ and also in the interior of $\ang{B}$. (With this definition it
      is not too difficult to prove the truth of the remark preceding this group
      of exercises.)
\item In $\ang{AOB}$, not a straight angle, point $P$ is between $A$ and $B$.
      Prove that $P$ is in the interior of the angle.
\item In triangle $ABC$ we have the relations $AXB$, $AYC$, and $XZY$. Prove
      that $Z$ is in the interior of the triangle.
\item If $\ang{AOB}$ is not a straight angle and we have the relations $XOB$ and
      $XYA$, prove that $Y$ is in the exterior of $\ang{AOB}$.
\item[\stex 12.] In Theorem 3--6 suppose $Q$ is in position $Q_1$, that is, on
      the same side of $l_2$ as $r_1$ and on the opposite side of $l_1$ as $r_2$
      (see Fig.~3--38). Show that there is a point $R$ of $r_1$ between $Q$ and
      $P$. [\emph{Hint:} Certainly there is a point $R$ of $l_1$ between $Q$ and
      $P$. Show that it is impossible that $R$ be on $r_1'$.]
\item[\stex 13.] In Theorem 3--6 suppose $Q$ is in position $Q_2$, that is, on
      $r_2'$ (see Fig.~3--39). Show that there is a point $R$ of $r_1$ between
      $Q$ and $P$.
\item[\stex 14.] Tell why the positions $Q_4$ and $Q_5$ of Fig.~3--33 are much
      the same as positions $Q_2$ and $Q_1$, respectively. If you have done
      positions $Q_1$ and $Q_2$, why are positions $Q_5$ and $Q_4$ easy?
\end{exlist}
\end{multicols}

\FIGIIITHIRTYSIX

\FIGIIITHIRTYSEVEN

\FIGIIITHIRTYEIGHTTHIRTYNINE

%========================================================= end matter
\section*{REVIEW OF CHAPTER 3}
\addcontentsline{toc}{section}{Review of Chapter 3}

\begin{exlist}
\item The fundamental relations of betweenness are expressed in the postulates.
      Study them again to be sure you know them.
\item Definitions of the following concepts were given in this chapter. Be sure
      that you know all of them.
\end{exlist}

\begin{multicols}{2}
\begin{itemize}[label=,leftmargin=1.2em,itemsep=0pt,topsep=2pt]
\item two sides of a line (see II--5)
\item segment
\item triangle
\item collinear (see footnote to Def.~3--3)
\item two sides of a point on a line
\item ray
\item angle
\item straight angle
\item vertical angles
\item adjacent angles
\item supplementary angles
\item interior of an angle
\item exterior of an angle
\item interior of a triangle
\item exterior of a triangle
\end{itemize}
\end{multicols}

\begin{exlist}[start=3]
\item The following statements have been proved in this chapter. First re-read
      them, and then study their proofs again.
      \begin{itemize}[label=,leftmargin=1.4em,itemsep=1pt,topsep=2pt]
      \item (a) Two different lines intersect in at most one point.
      \item (b) There are at least three lines.
      \item (c) A point $O$ on a line separates the line into two sides.
      \item (d) If $P$ and $Q$ are in the interior of $\ang{O}$, then the segment
            $PQ$ is in the interior of $\ang{O}$.
      \item (e) A ray from the vertex of an angle lies, except for the vertex,
            either in the interior or in the exterior of the angle.
      \item (f) If $P$ is in the interior of $\ang{O}$ and $Q$ is in the
            exterior, then the segment $PQ$ intersects the angle.
      \end{itemize}
\end{exlist}

\begin{center}\bfseries Important Remark\end{center}

By assuming a few properties of points and lines, we have been able to
\emph{define} many familiar types of figures and to \emph{prove} that these
figures have properties that may seem obvious. \emph{Therefore we have proved
that some of our figures agree with our geometry.} From now on in the book, we
will often determine properties of betweenness by just looking at figures. The
definitions and theorems of this chapter will enable us to use our figures in the
confident belief that we could always go back and supply proofs. Therefore, we
have completed the first stage in our development of geometry. The theorems
proved in this chapter need not be memorized because the facts are all familiar
from drawings. However, the postulates and definitions of this chapter are very
important, and therefore should be memorized.

\begin{center}\bfseries Review Questions\end{center}

\begin{exlist}
\item If $A$ and $B$ are distinct points, what can you say about a point $P$
      which is on the segment $AB$? Why?
\item What postulate did we use to define the sides of a point on a line?
\item A line $l$ contains the distinct points $R$ and $S$. A line $m$ also
      contains these points. What do you conclude?
\item If $ABC$, and $B$ is on line $l$ but $A$ is not on $l$, is the point $C$ on
      $l$? Why?
\item Tell whether the following statements are true or false, and give
      explanations for your answers.
      \begin{itemize}[label=,leftmargin=1.4em,itemsep=1pt,topsep=2pt]
      \item (a) For $\tri{ABC}$, if the point $S$ is in the interior of $\ang{A}$,
            $\ang{B}$, and $\ang{C}$, then it is a theorem that $S$ is in the
            interior of the triangle.
      \item (b) If $APB$, then $P$ belongs to the segment $AB$, by definition.
      \item (c) If $ARQ$ and $AGR$, then $AGQ$.
      \item (d) A point $P$ not on an angle must be in either the interior or
            exterior.
      \item (e) We postulate that a point on a line separates the line.
      \item (f) If $AQB$, then $A$ and $B$ are on a ray from $Q$.
      \item (g) A line segment contains a definite number of points.
      \end{itemize}
\item What does it mean to say that three or more points are noncollinear? How
      can we be sure that there are such sets of points?
\item If $ADB$ and $BAC$, what are the other betweenness relations?
\item Lines $s$ and $t$ intersect at $Q$ forming rays $s_1$, $s_2$ and $t_1$,
      $t_2$. Draw a figure and indicate the pairs of vertical angles. Which pairs
      of angles are adjacent? supplementary?
\item If $A$ and $B$ are on opposite sides of a line $l$, and $B$ and $C$ are on
      opposite sides of $l$, then $A$ and $C$ are on the same side of $l$. Why is
      this true?
\item Explain in your own words what it means for a point to be in the interior
      of an angle.
\end{exlist}

\begin{center}\bfseries Review Test\end{center}

\begin{exlist}
\item Define the following: (a) angle, (b) segment, (c) when two points are on
      the same side of a line, (d) triangle.
\item Tell whether true or false:
      \begin{itemize}[label=,leftmargin=1.4em,itemsep=1pt,topsep=2pt]
      \item (a) $A$, $B$, $C$ are three points on a line and $ABC$ and $CAB$.
      \item (b) If $PRU$, then $P$, $R$, and $U$ are distinct points on a line.
      \item (c) It is a theorem that a line separates the plane.
      \item (d) It is a theorem that a point separates a line.
      \item (e) It is a theorem that there is a point between any two points.
      \item (f) Betweenness is an undefined relation.
      \item (g) If $ABC$, then $CBA$.
      \item (h) If $E$, $F$, $G$ are not on $l$ and such that the segment $EF$
            does not intersect $l$, while the segment $FG$ does intersect $l$,
            then the segment $EG$ intersects $l$.
      \end{itemize}
\item State in your own words:
      \begin{itemize}[label=,leftmargin=1.4em,itemsep=1pt,topsep=2pt]
      \item (a) The point-line postulate.
      \item (b) The $ABC$ betweenness postulate.
      \item (c) The line-building postulate.
      \end{itemize}
\end{exlist}

\begin{center}\bfseries ALGEBRA REVIEW\end{center}

In this review you are supposed to use the properties of addition and
multiplication of numbers (both positive and negative), along with the
definitions of subtraction and division, to make simple proofs. All letters here
refer to rational numbers (fractions). These exercises show how definitions are
used in algebra.

\medskip
\noindent\textbf{Definition.}\quad If $a = b + x$, then $x$ is called ``$a$
\emph{minus} $b$'' and is written ``$x = a - b$.'' It is also called ``$b$
\emph{subtracted from} $a$.''

\begin{exlist}
\item If $c = p - q$, what is $p$? Why?
\item Why does $7 - 3 = 4$? Write $7 - 3 = x$, and ask what property $x$ must
      have.
\item Using the definition, what can you say about these numbers:
      $\frac{9}{5} - \frac{3}{5}$? $\frac{1}{2} - \frac{1}{3}$?
      $\frac{7}{8} - \frac{3}{4}$? $4.678 - 2.589$?
\end{exlist}

\medskip
\noindent\textbf{Definition.}\quad $-b$ (called ``\emph{negative} $b$'') is
defined to be $0 - b$. That is, $-b = 0 - b$ and is a number such that
$b + (-b) = 0$.

\begin{exlist}[start=4]
\item Why does $-7 = 0 - 7$? $-3 = 0 - 3$?
      $-\frac{1}{2} = 0 - \frac{1}{2}$?
\item Show that $9 + (-3) = 9 - 3$, $11 + (-8) = 11 - 8$,
      $\frac{5}{6} + (-\frac{1}{3}) = \frac{5}{6} - \frac{1}{3}$.
\item[\stex 6.] Prove that $a + (-b) = a - b$. [\emph{Hint:} Let $a - b = x$ so
      that $a = b + x$. Then $a + (-b) = b + x + (-b)$.]
\item[7.] Show that $-(7 + 3) = (-7) + (-3)$, $-(2 + 7) = (-2) + (-7)$.
\item[8.] Show that $-(a + b) = (-a) + (-b)$.
\end{exlist}

\medskip
\noindent\textbf{Definition.}\quad If $a = bx$ and $b \neq 0$, then $x$ is called
``$a$ \emph{divided by} $b$'' and is written ``$x = a \div b$'' or
``$x = a/b$.''

\begin{exlist}[start=9]
\item Explain why in this last definition we required that $b \neq 0$.
      [\emph{Hint:} Would $x$ be definite otherwise?]
\item According to the definition, $\frac{1}{2}$ is a number such that
      $1 = 2(\frac{1}{2})$. What then are the numbers $\frac{3}{4}$,
      $-\frac{2}{3}$, $4/{-5}$, $\frac{7}{16}$, $\frac{5}{10}$, $\frac{1}{1}$,
      $\frac{2}{2}$?
\item Why does $b \cdot a/b = a$?
\item Prove that $a/b + c/b = (a + c)/b$.
\item Show that $\frac{1}{2} = \frac{2}{4}$. [\emph{Hint:} Let
      $\frac{1}{2} = x$. Then $1 = 2 \cdot x$. Why? Hence $2 = 4x$. Why? Hence
      $\frac{2}{4} = x$. Why?]
\item Prove that if $b \neq 0$ and $c \neq 0$, then $a/b = ac/bc$.
\end{exlist}
\clearpage
% ==================== ch04 ====================
\setcounter{theorem}{0}\setcounter{definition}{0}\setcounter{construction}{0}\setcounter{exgroup}{0}\setcounter{footnote}{0}
\setcounter{chapter}{4}

\chapter*{\raggedright\Huge CONGRUENCE OF SEGMENTS}
\addcontentsline{toc}{chapter}{4.\ Congruence of Segments}
\markboth{CONGRUENCE OF SEGMENTS}{}
\vspace{-1em}\noindent\rule{\linewidth}{0.6pt}\vspace{1em}
\figurekey

%========================================================= 4-1
\section*{4--1\quad INTRODUCTION}
\addcontentsline{toc}{section}{4--1\quad Introduction}

In this chapter we shall be concerned with comparing the ``sizes'' of segments.
All of us have learned to understand the terms ``distance'' and ``length'' as
used in everyday life. For example, we may say that it is 12 miles to the
nearest town, or that a fence is 153 feet long. In each case we give a number to
indicate a distance or length. To give a number is a most convenient means of
describing the distance between two points or, in other words, of telling how far
apart they are.

But even before men learned to count, they must have wished at times to tell how
far apart points were, or to decide if two objects had the same length. For
example, a hunter visiting a friend in a neighboring village might have said to
his friend, ``My son is as tall as this mark on my spear.'' Then the men could
see, by comparing with the mark, whether the host's son was as tall as the
visitor's son. In this example the marked spear served as a kind of
``yardstick'' for comparing the height of the hunter's son with the height of
others. Even in very primitive societies men must have had a means for comparing
the sizes of objects.

\FIGIVONE

\begin{IVproblem}
Describe other examples where objects are compared in size without the use of
measurement, that is, a number. Some of your examples may use a marked stick (or
``yardstick'') and some not.
\end{IVproblem}

We are so used to speaking of distance or of length in terms of numbers that it
is hard for us to find other words to describe how far apart two points are.
Although we usually associate numbers with the words ``distance'' and
``length,'' what we study in this chapter is so simple that it would be
understandable to a person who knew nothing of numbers at all. We will be talking
about the ``size'' of a segment and whether two segments are ``the same size.''
As we have seen with the example of the hunter, a practical method for comparing
segments is to use a marked stick and to carry it from one place to another. We
shall see in the Section 4--2 how in geometry we can avoid carrying such sticks.

\exercisegroup{4--1}

\begin{multicols}{2}
\begin{exlist}
\item Your foot is a natural ``yardstick'' you carry with you. Describe several
      examples of using your foot for comparing lengths.
\item Think of several ways to use your hands as ``yardsticks.''
\end{exlist}
\end{multicols}

%========================================================= 4-2
\section*{4--2\quad CONGRUENCE}
\addcontentsline{toc}{section}{4--2\quad Congruence}

In daily life we are often interested in size, and we have ways of testing
whether or not two things are the same size. But our geometry so far has in it
\emph{no notion of size at all}. Because our geometry deals with imaginary
concepts (points and lines), to talk about size it is necessary to assume
(postulate) properties that we desire. We do not wish to talk about ``carrying a
segment around in the plane,'' because this involves the idea of movement, which
is not a part of geometry proper. Consequently \emph{we will just suppose that by
some means (it does not matter how) we can tell whether or not two segments are
the same size.} In other words, we will assume that we can pick out all the
segments of the same size.

We have been using the words ``the same size.'' We now use another word which is
to have exactly the same intuitive significance. The word is \emph{congruent},
and whenever you see it you can, if you like, say to yourself, ``the same size.''
Congruence refers to the \emph{relation} of being the same size. The word comes
from a Latin word meaning ``agree'' or ``conform.''

Euclid used the word ``equal'' instead of congruent. However, \emph{in this book
we will use the word ``equal'' (and the sign ``$=$'') to mean ``identically the
same.''} Thus, if we write $A = B$ we mean that $A$ and $B$ are one and the same
point. Or, if $x$ and $y$ are numbers, $x = y$ means that they are one and the
same number. And in dealing with segments, if we write $AB = CD$, we mean that
$A = C$ and $B = D$ or $A = D$ and $B = C$, that is, $AB$ and $CD$ are one and
the same segment.

%========================================================= 4-3
\section*{4--3\quad POSTULATE GROUP III. LINEAR CONGRUENCE}
\addcontentsline{toc}{section}{4--3\quad Postulate Group III. Linear Congruence}

We assume that there is a relation between segments called \emph{congruence}.
This relation is to have the following properties.

III--1 \emph{The existence postulate for segments.} Given points $A$ and $B$ on a
line $l$ and a point $A'$ on a line $l'$ (Fig.~4--2), then on each side of $A'$
on $l'$ there is exactly one point $B'$ such that $AB$ is congruent to $A'B'$. To
indicate that $AB$ is congruent to $A'B'$, we write ``$AB \cong A'B'$.''

\FIGIVTWOTHREE

In congruence the order of the points does not matter, that is, if
$AB \cong A'B'$, then also $AB \cong B'A'$, $BA \cong A'B'$, and
$BA \cong B'A'$.

III--2 \emph{The three-part segment postulate.} Congruence also satisfies the
following:

\begin{enumerate}[label=(\alph*), leftmargin=2.2em, itemsep=2pt, topsep=3pt]
\item $AB \cong AB$. (A segment is congruent to itself; that is, congruence of
      segments is reflexive.)
\item If $AB \cong A'B'$, then $A'B' \cong AB$. (Congruence of segments is
      symmetric.)
\item If $AB \cong A'B'$ and $A'B' \cong A''B''$, then $AB \cong A''B''$.
      (Congruence of segments is transitive.)
\end{enumerate}

\begin{IVproblem}
Prove that if $AB \cong CD$ and $EF \cong CD$, then $AB \cong EF$.
[\emph{Hint:} Use parts (b) and (c) of III--2.]
\end{IVproblem}

III--3 \emph{The addition and subtraction of segments postulate.} Suppose that
$B$ is between $A$ and $C$ on a line $l$, and that $B'$ is between $A'$ and $C'$
on a line $l'$ (Fig.~4--3).

\begin{enumerate}[label=(\alph*), leftmargin=2.2em, itemsep=1pt, topsep=3pt]
\item If $AB \cong A'B'$ and $BC \cong B'C'$, then $AC \cong A'C'$.
\item If $AC \cong A'C'$ and $BC \cong B'C'$, then $AB \cong A'B'$.
\end{enumerate}

These postulates are somewhat like our earlier ones. They seem so obvious that
you might feel they do not need to be stated. As before, however, you must
remember that points and lines are inventions of the mind, and so we must state
clearly what properties they are to have. Euclid did not state them as clearly as
we have. The postulates given in this book are easier to use than those of Euclid
because they tell you exactly what you can do. They are similar to ones given by
David Hilbert (see Section 1--2).

In less exact language, the postulates say the following things. The first says
that on the line $l'$ there are two points $B'$ and $B''$, one on each side of
$A'$, such that $AB$ is the same size as both $A'B'$ and $A'B''$. The second one
says three things: that a segment is the same size as itself; that if $AB$ is the
same size as $A'B'$, then $A'B'$ is the same size as $AB$; and that if two
segments are the same size as a third segment, then the two segments are the same
size. The third postulate tells how to combine segments by what we call
\emph{addition and subtraction of segments}.

\exercisegroup{4--2}

\begin{multicols}{2}
\begin{exlist}
\item What does it mean to say $RS = TU$?
\item If $RS = TU$, is $RS \cong TU$? Justify your answer.
\item Is it true for every pair of segments $RS$ and $TU$ that if
      $RS \cong TU$ then $RS = TU$?
\item Two segments $AB$ and $CD$ are not congruent to each other. (We write this
      as $AB \ncong CD$.) Prove that $AB$ and $CD$ are not equal to each other
      ($AB \neq CD$).
\item If $ABC$, what postulates must you use to prove that $AB \ncong AC$?
\item Is the following statement true or false? If $AB$ is any segment and $A'$
      is any point on a line $l'$, then there is one and only one point $B'$ on
      $l'$ such that $AB \cong A'B'$.
\item In Fig.~4--4 the following congruences hold: $AB \cong BC$,
      $CD \cong CE$, $CE \cong ED$, $EF \cong EB$. List all other congruences
      that you recognize.
\item The sign $=$ was defined in this section. State the definition in your own
      words.
\item State the three postulates of linear congruence in your own words, using
      the phrase ``the same size'' instead of ``congruent.'' Then go back and
      read again the paragraph preceding this group of exercises.
\item Show (using postulates) that if $A$, $B$, $C$, $D$ are points, in this
      order, on a line $l$, and if $A'$, $B'$, $C'$, $D'$ are points, in this
      order, on a line $l'$, and if $AB \cong A'B'$, $BC \cong B'C'$,
      $CD \cong C'D'$, then $AD \cong A'D'$. [\emph{Hint:} Draw a figure. You
      will have to use part (a) of Postulate III--3 twice.]
\item Assuming that the points have the same order as in Exercise 10, prove that
      $(AD \cong A'D'$ and $CD \cong C'D'$ and
      $BC \cong B'C') \rightarrow AB \cong A'B'$.
\item Suppose $E$, $F$, $G$ and $E'$, $F'$, $G'$ are \emph{any} points on lines
      $l$ and $l'$, respectively, such that $EF \cong E'F'$ and
      $FG \cong F'G'$. Show, by drawing a figure, that it does not need to be
      true that $EG \cong E'G'$. What do we need to know about the position of
      the points? Does III--3 take care of this?
\item[\stex 13.] Part (b) of Postulate III--3 can be \emph{proved} from the
      other postulates. We have included this as a postulate to make things
      easier. Try to prove (b). You will have to use (a) and Postulate III--1.
\item[\stex 14.] Try to reword III--3 in terms of length.
\item[15.] Each part of Fig.~4--5 indicates certain congruence relations. The
      marks show which segments are congruent. What other relations can be
      deduced? Point out how the postulates are used.
\item[\stex 16.] Show how to relabel the points in parts (b) and (c) of
      Fig.~4--5 as $A$, $B$, $C$ and $A'$, $B'$, $C'$ so that the application of
      Postulate III--3 is clearly shown.
\end{exlist}
\end{multicols}

\FIGIVFOUR

\FIGIVFIVE

%========================================================= 4-4
\section*{4--4\quad COMPARING SEGMENTS}
\addcontentsline{toc}{section}{4--4\quad Comparing Segments}

The postulates in Group III tell us about segments that are the same size or
congruent to each other. We can now \emph{define}, in terms of congruence and
betweenness, the ideas of larger and smaller segments. The definition will simply
describe what the hunter did when he compared the boy with the mark on the spear.

We need a symbol to stand for ``is greater than'' and another symbol to stand for
``is less than.'' These symbols are: $>$ and $<$. Whenever you see $>$, read
``is greater than,'' and whenever you see $<$, read ``is less than.''

\begin{definition}
Let $AB$ be any segment and $A'B'$ any other segment on line $l'$ (Fig.~4--6). On
the same side of $A'$ as $B'$, let $B''$ be the point such that
$AB \cong A'B''$. If $A'B'B''$ (that is, if $B'$ is between $A'$ and $B''$), we
say that $AB > A'B'$. If $A'B''B'$ (that is, $B''$ is between $A'$ and $B'$), we
say that $AB < A'B'$.
\end{definition}

\FIGIVSIX

\begin{remark}
Because of III--1 there is only one point $B''$ on the same side of $A'$ as the
point $B'$ and such that $AB \cong A'B''$. The three points $A'$, $B'$, and
$B''$ are therefore in some definite order (unless $B' = B''$, in which case
$AB \cong A'B'$). We cannot have $A'$ between $B'$ and $B''$ because $B'$ and
$B''$ are on the same side of $A'$. Therefore if $AB \ncong A'B'$, we must have
either $A'B'B''$ or $A'B''B'$. In other words, either $AB > A'B'$ or
$AB < A'B'$. This remark proves the following theorem.
\end{remark}

\begin{theorem}
If $AB$ and $A'B'$ are any two segments, then either $AB > A'B'$, or
$AB \cong A'B'$, or $AB < A'B'$, and exactly one of these is true.
\end{theorem}

\exercisegroup{4--3}

\begin{multicols}{2}
\begin{exlist}
\item Does Fig.~4--7 indicate that $AB < CD$ or that $CD > AB$? Read
      Definition 4--1 again.
\item What relationships are indicated by Fig.~4--8? Use Definition 4--1.
\item What relationships are indicated by Fig.~4--9? Use Definition 4--1.
\item If we know that $AB \cong CD$ and $ABX$, what can we conclude? Use
      Definition 4--1.
\item It is known that $RS \cong UT$ and $RVS$. What conclusion can be drawn?
      Use Definition 4--1.
\item Does Definition 4--1 say that if $AB > CD$ then $CD < AB$? Does looking at
      a figure prove that this is so?
\end{exlist}
\end{multicols}

\FIGIVSEVENEIGHTNINE

The exercises above are meant to prepare you for the main theorem of this
section, namely: if one segment is greater than a second, then the second is less
than the first. Of course, we know that if we draw two pencil marks so that the
first is longer than the second, then the second is shorter than the first. But
we have not yet defined \emph{length of a segment}. As a matter of fact we shall
need to introduce a new postulate into our geometry before length can be defined.
Since our definitions of ``greater than'' and ``less than'' have been formulated
in terms of betweenness, the truth of any statement about these relations
(``greater than'' and ``less than'') must be established by using properties of
betweenness.

% Theorem 4-2 carries a dagger and a marker-less footnote in the book, so it is
% set by hand rather than through the theorem environment; the counter is still
% stepped so that Theorem 4-3 numbers correctly.
\refstepcounter{theorem}%
\par\medskip\noindent
\textbf{$\dagger$Theorem \thetheorem.}\ \emph{If $AB > A'B'$ then
$A'B' < AB$, and conversely, if $A'B' < AB$ then $AB > A'B'$.}%
\begingroup\let\thefootnote\relax
\footnotetext{$\dagger$ Although this theorem is intuitively obvious, its proof
is quite difficult. If the teacher wishes, it may be treated as a postulate.}%
\endgroup
\par\medskip

\noindent\emph{Proof.} Figure 4--10 illustrates what must be proved in order to
establish the first part of the theorem.

\begin{enumerate}[label=(\arabic*), leftmargin=2.2em, itemsep=1pt, topsep=3pt]
\item $AB > A'B'$ means $AB \cong A'C'$ and $A'B'C'$.
\item $A'B' < AB$ means $A'B' \cong AC$ and $ACB$.
\end{enumerate}

\FIGIVTENELEVEN

To show that $AB > A'B' \rightarrow A'B' < AB$, we must prove that if $C'$ falls
``beyond'' $B'$, then $C$ falls between $A$ and $B$.

From (1) above we have $AB \cong A'C'$ and $A'B'C'$. We know also that on a
given side of $A$ there is a point $C$ such that $AC \cong A'B'$. To establish
that we do indeed have $ACB$, consider the point $X$ in Fig.~4--11 such that
$ACX$ and $CX \cong B'C'$. Now, by addition of segments, $AX \cong A'C'$. But
$AB \cong A'C'$, so $X = B$. (Why?) And since $ACX$, we have $ACB$; and from (2)
above we have $A'B' < AB$.

The proof of the converse is similar to the above proof and is left as an
exercise.

From now on we shall freely interchange the ``greater than'' and ``less than''
terminology since the statements $AB > CD$ and $CD < AB$ are equivalent to each
other.

\begin{theorem}
If $AC \cong A'C'$ and $A'B'C'$ and $B$ is the point on $C$'s side of $A$ such
that $AB \cong A'B'$, then $ABC$ (Fig.~4--12).
\end{theorem}

\noindent\emph{Proof.} Since $AC > A'B'$, $A'B' < AC$, and so by Definition 4--1
we have $B$ between $A$ and $C$.

\medskip
This theorem is proved at this time because we shall need it in the next chapter
when we discuss the length of a segment. We can then show that congruent segments
have the same length. It is also a very convenient tool to use in many subsequent
proofs.

\FIGIVTWELVETHIRTEEN

\begin{definition}
If we wish to say ``$AB \cong A'B'$ or $AB > A'B'$,'' we write
``$AB \geqq A'B'$'' and read the symbol $\geqq$ as ``is greater than or
congruent to.'' Similarly, if we wish to say ``$AB \cong A'B'$ or
$AB < A'B'$,'' we write ``$AB \leqq A'B'$'' and read the symbol $\leqq$ as ``is
less than or congruent to.''
\end{definition}

\begin{theorem}
If $AB > CD$ and $CD \cong EF$, then $AB > EF$.
\end{theorem}

\noindent\emph{Proof.} Since $CD < AB$ (why?) there is a point $D'$
(Fig.~4--13) on $AB$ such that $CD \cong AD'$. But then, by III--2,
$EF \cong AD'$. Hence, by Definition 4--1, $EF < AB$ and $AB > EF$.

\exercisegroup{4--4}

\begin{multicols}{2}
\begin{exlist}
\item Using the definition of $>$, explain what must be done in order to prove
      that $AB > CD$. Draw a figure to illustrate this.
\item What must one show in order to prove $RS < TU$? Draw a figure to
      illustrate this.
\item What basic undefined relations are used to define the phrases ``is greater
      than'' and ``is less than''?
\item Is it correct to say that $AB \geqq CD \rightarrow AB \ncong CD$? Is the
      converse of this implication true?
\item Is it correct to say that $AB > CD \rightarrow AB \ncong CD$? State both
      the converse and contrapositive of this statement. Which one of these
      statements is true?
\item What conclusion can you draw if you know that $A$, $B$, $C$, and $D$ are
      points on line $l$ and that $AB > CD$ and $A = C$?
\item If $AB$ is neither equal to nor greater than $CD$, can you conclude
      $AB < CD$?
\item If $AB$ is neither congruent to nor less than $CD$, can you conclude
      $AB > CD$?
\item[\stex 9.] Prove that if $AB < EF$ and $EF \cong CD$, then $AB < CD$.
      State the converse and contrapositive of this theorem.
\item[\stex 10.] Prove that if $AB > CD$ and $CD > EF$, then $AB > EF$.
\item[\stex 11.] Prove that if $AB < CD$ and $CD < EF$, then $AB < EF$.
\item[\stex 12.] Prove that if $B$ is between $A$ and $C$, and $B'$ is between
      $A'$ and $C'$, and $AB \cong A'B'$, and $BC > B'C'$, then $AC > A'C'$.
\item[13.] If $AB > CD$ and $CD < EF$, can you prove that $AB > EF$?
\item[\stex 14.] Prove that if $ABC$ and $A'C'B'$ and $AC \cong A'C'$, then
      $AB < A'B'$.
\end{exlist}
\end{multicols}

%========================================================= Review
\IVhead{Review of Chapter 4}
\addcontentsline{toc}{section}{Review of Chapter 4}
\markboth{REVIEW OF CHAPTER 4}{}

Several important ideas have been developed in this chapter. You should realize
that \emph{congruence} is really an undefined relation. Of course, you are so
accustomed to comparing the sizes of things by comparing numbers that it is
probably quite difficult to see clearly that numbers are not actually needed. We
know that a $6'8''$ basketball center is taller than a
$6'2''$ center. We know that a 240-lb football tackle is
larger than a 180-lb tackle. But such comparisons as these can be made
\emph{without using numbers at all}. Surely men were able to recognize
differences in sizes long before they had a number system to measure them.
Historically, the Greeks developed their geometry without the use of number
concepts.

In the next chapter, when we define what is meant by the length of a segment, it
will naturally turn out that congruent segments have the same length. That is, if
$AB \cong CD$, and $x$ and $y$ are numbers which are the lengths of these
segments, then $x = y$.

Again, the postulates of this chapter are so obvious to our intuition that you
might wonder why we bother to make them. All we can say is that postulates should
be obvious. The amazing thing about mathematics is that out of such simple
assumptions as we are making in our study of geometry one can construct theorems
of marvelous complexity which provide the foundation for all work in science.

The proofs in this chapter are quite difficult. There is no need to memorize the
results of the theorems since they merely show that our geometry agrees with the
figures that we draw.

In working with Definition 4--1, you should have learned an important
mathematical lesson, namely, that definitions must be clearly understood. It
would be very easy to read Definition 4--1 carelessly and to assume that if
$AB > CD$ then $CD < AB$. This is not stated in the definition, but of course we
were able to use the definition, our postulates, and the rules of logic to prove
that it is indeed true.

The exercises below will help you review the concepts of this chapter.

\medskip
Label the following statements true or false.

\begin{exlist}
\item If $AB \cong CD$, then $A = C$ and $B = D$.
\item If $A = C$ and $B = D$, then $AB \cong CD$.
\item If $AB \cong CD$ or $AB < CD$, then $AB \leqq CD$.
\item If $A$, $B$, and $C$ are points on a line $l$ and $A'$, $B'$, and $C'$ are
      points on a line $l'$, and $AB \cong A'B'$, and $BC \cong B'C'$, then
      $AC \cong A'C'$.
\item If $AB$ is any segment and $A'$ is any point on a line $l'$, then there are
      exactly two points $B'$ and $B''$ on $l'$ such that $A'B'$ and $A'B''$ are
      congruent to $AB$.
\item If $AB$ is not less than $CD$, then it must be greater than $CD$.
\item If $RS$ is neither greater than nor less than $TU$, then $RS \cong TU$.
\item If $A$, $B$, $C$, and $D$ are points in this order on a line $l$ and
      $AC \cong BD$, then $AB \cong CD$.
\item If $A$, $B$, $C$, and $D$ are points in this order on a line $l$ and
      $AC > BD$, then $AB < CD$.
\item If $A$, $B$, $C$, and $D$ are points in this order on a line $l$ and
      $AC < BD$, then $AB \leqq CD$.
\end{exlist}

\medskip
This chapter has included three postulates, two definitions, and four theorems.
The following statements are all true. Designate a postulate or definition or
theorem which could be used in justifying each statement.

\begin{exlist}[start=11]
\item For every line segment $AB$, it is true that $AB \cong AB$.
\item If $AB \cong CD$ and $ABX$, then $AX > CD$.
\item If $AB$ is not less than $CD$, then $AB \geqq CD$.
\item If $AB \cong CD$ and $CD \cong EF$ and $XY > AB$, then $XY > EF$.
\item If $RS \cong TU$ and $XY \cong TU$, then $RS \cong XY$.
\item If $ABC$ and $BCD$, then $AC < AD$.
\item If $AB \leqq CD$, then there is no point $E$ such that $CDE$ and
      $AB \cong CE$.
\item If in triangle $ABC$ we have side $AC$ greater than side $AB$, then there
      is a point $E$ on side $AC$ such that $AE \cong AB$.
\item If $ABC$ and $AB \cong BC$ and $XYZ$ and $XY \cong YZ \cong AB$, then
      $AC \cong XZ$.
\item If $A = X$ and $B = Y$, then $AY \cong BX$.
\end{exlist}

%========================================================= Algebra review
\IVhead{Algebra Review}
\addcontentsline{toc}{section}{Algebra Review}

In algebra the relations of ``greater than'' and ``less than'' are defined for
pairs of numbers. We say that $5 > 3$ because $5 = 3 + 2$, and 2 is a positive
number. More generally we make the following definitions:

\begin{enumerate}[label=(\arabic*), leftmargin=2.2em, itemsep=2pt, topsep=3pt]
\item If $a$, $b$, and $c$ are numbers and $a = b + c$ and $c$ is positive, then
      we say that $a$ is greater than $b$ and write $a > b$.
\item If $a$, $b$, and $c$ are numbers and $a = b - c$ and $c$ is a positive
      number, then we say that $a$ is less than $b$ and write $a < b$.
\end{enumerate}

For example, since $-\frac{3}{2} = -\frac{5}{2} + \frac{2}{2}$, we say
$-\frac{3}{2} > -\frac{5}{2}$; and since $0 = \frac{3}{4} - \frac{3}{4}$ we say
that $0 < \frac{3}{4}$.

\begin{exlist}
\item Use the first definition above to prove that $0 > -\frac{4}{3}$.
\item Use the second definition to prove that $-\frac{4}{3} < 0$.
\item Use the two definitions to prove that if $a > b$ then $b < a$.
\item If $a = b + 1$ and $b = c + 2$, prove that $a > c$.
\item Prove that if $a > b$ and $b > c$, then $a > c$.
\item If $a = b - 2$ and $b = c - 3$, then prove that $a < c$.
\item Prove that if $a < b$ and $b < c$, then $a < c$.
\item Show that if $a = b + 3$, then $3a > 3b$ and $-2a < -2b$.
\item Prove that if $a > b$, then $ca > cb$ if $c$ is positive and $ca < cb$ if
      $c$ is negative.
\item Prove that if $a < b$, then $ca < cb$ if $c$ is positive and $ca > cb$ if
      $c$ is negative.
\end{exlist}
\clearpage
% ==================== ch05 ====================
\setcounter{theorem}{0}\setcounter{definition}{0}\setcounter{construction}{0}\setcounter{exgroup}{0}\setcounter{footnote}{0}
\setcounter{chapter}{5}

\chapter*{\raggedright\Huge MEASUREMENT\\OF SEGMENTS}
\addcontentsline{toc}{chapter}{5.\ Measurement of Segments}
\markboth{MEASUREMENT OF SEGMENTS}{}
\vspace{-1em}\noindent\rule{\linewidth}{0.6pt}\vspace{1em}
\figurekey

%========================================================= 5-1
\section*{5--1\quad INTRODUCTION}
\addcontentsline{toc}{section}{5--1\quad Introduction}

In this chapter we are going to get a number for the length of a segment, that
is for the distance between two points. The way we get the length will be
exactly the same as the method we use in practical work. We will see that
finding a length is based upon simple \emph{counting}. For example, when you
pace off the length of a building (Fig.~5--1), you count how many steps you
take. Or when you say that the room is 15 feet wide, you mean that a foot ruler
can be laid down 15 times, one after the other, on a line. If the width of the
room is between 14 and 15 feet, the ruler can be laid down 14 times, but 15
times would put the end of the ruler past the wall. To measure the room more
accurately, you must choose a smaller unit of measure, say an inch, so that if
the width is 14 feet and three inches, the foot ruler can be laid down 14 times
and beyond that the inch ruler can be laid down three times. If even then you
have not come exactly to the wall of the room, you must choose an even smaller
unit of measure, some fraction of an inch, and so on.

\FIGVONE

\exercisegroup{5--1}

\begin{multicols}{2}
\begin{exlist}
\item Measure the length of your desk. Describe the process of measurement in
      terms of counting.
\item Measure the length of a page. Describe the counting you do to get the
      length.
\item Think of other examples of measurement and describe them in terms of
      counting.
\end{exlist}
\end{multicols}

%========================================================= 5-2
\section*{5--2\quad NUMBERS}
\addcontentsline{toc}{section}{5--2\quad Numbers}

Because we express length in terms of numbers, we need certain facts about
numbers that you learned in algebra and arithmetic. You have studied different
kinds of numbers. In the beginning grades you first used whole numbers
$1$, $2,\ldots$ and learned how to add, subtract, and multiply them. These
numbers we now call the \emph{positive integers} instead of ``whole numbers.''
When you tried to divide integers you found that it was necessary to have other
numbers, which you called \emph{fractions}, for example,
$\frac{1}{2}$, $\frac{2}{3}$, $\frac{5}{8}$, $\frac{7}{37}$. These are numbers
obtained by dividing one integer by another integer. \emph{The numbers which
are either positive integers or positive fractions are called positive rational
numbers.}

You also wrote numbers as \emph{decimals}, for example, 1.25, 2.537, 173.612.
Sometimes a rational number when written as a decimal ``comes out even,'' as
for example, $\frac{1}{2} = 0.5$, $\frac{3}{8} = 0.375$. Sometimes it does not
come out even, so that the decimal goes on forever, for example,
$\frac{1}{3} = 0.333\ldots$, $\frac{2}{3} = 0.6666\ldots$,
$\frac{1}{11} = 0.090909\ldots$. You also met with other numbers which did not
come out even, for example, $\sqrt{2} = 1.4142\ldots$,
$\sqrt{3} = 1.7321\ldots$. These numbers are different from $\frac{1}{3}$,
$\frac{2}{3}$, and $\frac{1}{11}$ because they not only do not come out even,
but \emph{the decimal never repeats}. In other words, no matter to how many
decimal places you were to calculate $\sqrt{2}$, you would never reach a place
where, from then on, one group of digits would occur over and over again in the
same order.

Decimals that, after some place, repeat one group of digits over and over are
called \emph{repeating decimals}. Decimals which come out even are called
\emph{terminating decimals}, and may be considered as repeating with zeros.

\exercise{5--2}

Write the following rational numbers as decimals to see that they are repeating
decimals. If they come out even (terminate), the repeating part is all zeros.
1/6, 5/6, 122/99, 1/13, 1/7, 6/5, 3/2.

\medskip
It is proved in advanced algebra that rational numbers are always repeating
decimals, and that, conversely, repeating decimals are rational numbers. (These
proofs are hard, so don't be discouraged if you cannot construct them; your
teacher may show you.) Hence, rational numbers and repeating decimals are the
same. What, then, about the decimals which never repeat? These numbers are not
rational numbers, and so they have been called \emph{irrational
numbers}.*%
\Vstarfoot{The word ``irrational'' is not a very pleasant choice
because it sounds like ``not sensible.'' What it really means is this. Rational
numbers are obtained by dividing one integer by another; that is, a rational
number is the ratio of two integers. Hence the word ``rational'' in ``rational
number'' means that the number is a \emph{ratio} of integers. Therefore, an
\emph{irrational} number is one which is \emph{not} the ratio of two integers.}
Irrational numbers are defined to be \emph{those numbers which are not
rational}. Examples of irrational numbers are $\sqrt{2}$, $\sqrt{3}$,
$\sqrt{5}$, $\sqrt[3]{2}$, $\sqrt[3]{17}$, and, in fact, any root of an integer
which is not itself an integer.

All the decimals together are defined to be the \emph{real numbers}. Thus a
real number is a decimal. If the decimal repeats, the number is a rational
number. If the decimal does not repeat, it is an irrational number.

\exercisegroup{5--3}
% Book p.75 holds this chapter's first starred exercises (nos. 4 and 7) but
% prints NO "*Starred ... are optional" footnote -- verified against PDF p.89.
% ch02 owns the book-wide first use; deliberately no \starnote here.

\begin{multicols}{2}
\begin{exlist}
\item Write the symbols for several rational numbers and for several irrational
      numbers.
\item Is $3.131313\ldots$, where the 1's and 3's keep alternating, a rational
      number?
\item Is $0.234234234\ldots$, where 2, 3, and 4 keep following each other, a
      rational number?
\item[\stex 4.] The number $1.010010001\ldots$ is formed from just 0's and 1's.
      After each 1 there is one more 0 than the previous time. Is this number
      rational?
\item[5.] Give the definitions of rational number, irrational number, real
      number.
\item[6.] A real number is either rational or irrational. Why?
\item[\stex 7.] There is a rule for deciding when a fraction written as a
      decimal will have its repeating part all zeros and when it will not. For
      example, each of the fractions 3/20, 17/250, 5/16, and 4/25 has its
      repeating decimal part all zeros. The fractions 7/15, 5/12, and 11/23 do
      not. Can you discover the rule?
\end{exlist}
\end{multicols}

When we write a decimal, for example 1.23, we mean $1 + 0.2 + 0.03$ or
$1 + 2/10 + 3/100$. Similarly, $4.383 = 4 + 3/10 + 8/100 + 3/1000$. Thus, any
terminating decimal (one which comes out even) may be written as a sum of
ordinary fractions with denominators 10, $100,\ldots$. If the decimal is
repeating, this sum is an \emph{infinite} one, for example,
$1/3 = 0.333\ldots = 3/10 + 3/100 + 3/1000 + 3/10000 + \cdots$.

\exercisegroup{5--4}

\begin{multicols}{2}
\begin{exlist}
\item Write the following decimals as sums of rational numbers: 4.82, 1.29,
      2.75, 7.122, 0.8875, 2.40. In order to express each number as a single
      fraction, add the rational numbers.
\item Write the following numbers as decimals: $3 + 2/10 + 4/100$,
      $2 + 3/10 + 5/100 + 6/1000$, $4 + 3/100$, $1 + 1/10 + 1/10000$.
\item Write the following repeating decimals as sums of rational numbers:
      $0.6666\ldots$, $0.7777\ldots$, $0.121212\ldots$.
\item Write the following numbers as decimals:
      $1 + 2/10 + 2/100 + 2/1000 + \cdots$,
      $3 + 1/10 + 4/100 + 1/1000 + 4/10000 + \cdots$,
      $8/10 + 8/100 + 8/1000 + 8/10000 + \cdots$.
\end{exlist}
\end{multicols}

%========================================================= 5-3
\section*{5--3\quad INEQUALITIES}
\addcontentsline{toc}{section}{5--3\quad Inequalities}

In this book you are expected to know how to do arithmetic, that is, how to
add, subtract, multiply, and divide real numbers. But there is one part of
arithmetic you have not yet studied systematically. That part is concerned with
comparing the sizes of real numbers. Since you already know what it means to
say that the number $x$ is larger than the number $y$, all we need is a symbol
to indicate this. We therefore make the definition: $x > y$ means ``$x$ is
greater than $y$.'' Another way to say this is ``$x - y$ is a positive
number.'' We also define $x < y$ to mean ``$x$ is less than $y$,'' or, in other
words ``$y$ is greater than $x$.'' Listed below are several properties of real
numbers. These properties can be established as indicated in the review of
Chapter 4.

\begin{enumerate}[label=(\alph*), leftmargin=2.2em, itemsep=1pt, topsep=3pt]
\item If $x > y$ and $y = z$, then $x > z$.
\item If $x > y$ and $y > z$, then $x > z$.
\item $x$ is a positive number if and only if $x > 0$.
\item $x$ is a negative number if and only if $x < 0$.
\item If $x > y$, then $x - y > 0$, and conversely.
\end{enumerate}

\medskip
We define $x \geqq y$ to mean ``either $x > y$ or $x = y$,'' and we read
``$\geqq$'' as ``is greater than or equal to.'' We make a similar definition
for $x \leqq y$ (``$x$ is less than or equal to $y$'').

\begin{enumerate}[label=(\alph*), leftmargin=2.2em, itemsep=1pt, topsep=3pt,
                  start=6]
\item If $x > y$ and $u \geqq v$, then $x + u > y + v$.
\end{enumerate}

\exercisegroup{5--5}

\begin{multicols}{2}
\begin{exlist}
\item Substitute numbers for the variables in properties (a), (b), (e), and (f)
      above, and thus verify their truth for certain cases.
\item Write in symbols: $a$ is greater than $b$; $a$ is greater than or equal
      to $b$; $u$ is less than $w$; $a + c$ is less than or equal to $r + s$;
      $(a + b)^2$ is positive; $z$ is a negative number.
\item If $x$ is any real number, $x^2 \geqq 0$. Why?
\item Why is it true that if $x > 0$ and $y > 0$, then $xy > 0$?
\item Tell why $xy < 0$ if $x > 0$ and $y < 0$.
\end{exlist}
\end{multicols}

%========================================================= 5-4
\section*{5--4\quad MEASUREMENT}
\addcontentsline{toc}{section}{5--4\quad Measurement}

To measure a line segment we must first tell what our unit of length is. Some
units of length often used in practical work are inch, foot, yard, mile, meter,
centimeter, etc. Having picked a unit of length, we will describe the geometric
process whereby we get a real number, that is, a decimal for the length of a
line segment.

Choose any segment $CD$ as a unit (Fig.~5--2). This segment and all segments
congruent to it are defined to have length 1. Suppose $AB$ is any segment on a
line $l$. On the same side of $A$ as $B$ there is a point $A_1$ such that
$AA_1 \cong CD$. If $A_1 = B$, then the length of $AB$ is 1.

\FIGVTWO

\FIGVTHREE

\noindent If $B$ lies between $A$ and $A_1$ (Fig.~5--3), then our unit was too
large to measure the length of $AB$ in integers (whole numbers). We will return
to this possibility later. If $A_1$ lies between $A$ and $B$, then there is a
point $A_2$, on the same side of $A_1$ as $B$, such that either $A_2 = B$, or
$B$ is between $A_1$ and $A_2$, or $A_2$ is between $A_1$ and $B$. If
$A_2 = B$, the length of $AB$ is 2. If $B$ is between $A_1$ and $A_2$, then our
unit will not measure the length in integers. If $A_2$ is between $A_1$ and
$B$, there is a point $A_3$, on the same side of $A_2$ as $B$, such that
$A_2A_3 \cong CD$.

The process described above is what we actually follow when we measure. If we
lay out segments congruent to $CD$ from $A$ toward $B$, we know from experience
that we eventually either reach $B$ exactly, so that the length of $AB$ is an
integer, or reach a point \emph{beyond} $B$. Thus in Fig.~5--2, $A_4$ is beyond
$B$ but $A_3$ is still between $A$ and $B$. The length in this case will be a
number between 3 and 4. If the above process can always be carried out in our
geometry, then we have described the beginning of the process of measurement.

The only part of the above process that does not follow from our postulates is
the proof that we can always get to or beyond $B$. Because we need to do this,
we now add a new postulate to our list. It was the Greek mathematician
Archimedes (287--212 B.C.) who first stated this postulate. He apparently did
not see how to get it from Euclid's postulates and, needing it, assumed it was
true. It is one of the things that Euclid forgot to state as a postulate. It is
called \emph{Archimedes' axiom} or \emph{Archimedes' postulate} in spite of its
use by Eudoxus much earlier. This is the first of the two new postulates in
Postulate Group IV which we need for measurement. The other one will be given
later.

\medskip
IV--1 \emph{Archimedes' postulate.} If $AB$ is any segment on a line $l$ and
$CD$ any other segment, then there is a definite number $n$ of points
$A_1$, $A_2,\ldots$, $A_n$ on $l$ such that the points $A = A_0$, $A_1$,
$A_2,\ldots$, $A_n$ are arranged in this order, and all the segments $AA_1$,
$A_1A_2,\ldots$, $A_{n-1}A_n$ are congruent to $CD$, and either $B = A_n$ or
$B$ lies between $A_{n-1}$ and $A_n$.

\medskip
The integer $n$ is called \emph{Archimedes' integer}. When $n$ is 4, an
illustration would look like Fig.~5--4 or Fig.~5--5.

\FIGVFOURFIVE

\exercisegroup{5--6}

\begin{multicols}{2}
\begin{exlist}
\item If $AB$ is this segment \raisebox{0.18ex}{$A$\rule[0.35ex]{2.2em}{0.4pt}$B$}
      and $CD$ is this segment
      \raisebox{0.18ex}{$C$\rule[0.35ex]{1.3em}{0.4pt}$D$}, find how many
      points are needed to get beyond $B$.
\item Draw a segment $AB$ and a segment $CD$. Find the integer $n$ of
      Archimedes' axiom for your segments.
\end{exlist}
\end{multicols}

Using Archimedes' postulate, we find that the length of $AB$ is either an
integer or is between two integers. If the length lies between two integers,
the next step in the process is to divide the unit interval into tenths
(Fig.~5--6). (We would not have to keep to tenths if we were not trying to come
out with a decimal.) Then we count ``tenths'' until either we reach $B$ exactly
or $B$ lies between two of the tenth marks. If we reach $B$ exactly, then the
length of $AB$ comes out exactly in tenths. If $B$ is not reached exactly, then
we must divide the unit interval into one-hundredths and repeat the process. In
this way we obtain a real number or decimal, for the length of $AB$. In
Fig.~5--7 the length of $AB$ is about 2.3.

\FIGVSIXSEVEN

We have indicated a segment by writing the two endpoints of the segment as
$AB$. At this point we also need a notation for the \emph{length} of the
segment. For this we will use $\sg{AB}$, that is $\sg{AB} = $ \emph{the length
of the segment} $AB$. Whenever you see the symbol $\sg{AB}$ you must remember
that it represents a \emph{number}, while the symbol $AB$ stands for a \emph{set
of points}.

Except for dividing a segment into tenths, the entire procedure of this section
can be carried out (that is, described) in our geometry. At the moment we do
not have the means for dividing a segment into tenths. As soon as we have
studied parallel lines (Chapter 8), we will be able to make this construction.
Since we do not need the results of this chapter until after we have studied
parallel lines, our procedure is logically correct.

\exercisegroup{5--7}

\begin{multicols}{2}
\begin{exlist}
\item Draw a line segment. Then draw a segment that will be your unit. Now
      measure the first segment, describing how you use Archimedes' postulate
      and relations of betweenness in the process.
\item Draw a line segment. Measure it with your ruler to the nearest 1/16 of an
      inch. Describe your measurement process, using Archimedes' postulate and
      betweenness.
\end{exlist}
\end{multicols}

%========================================================= 5-5
\section*{5--5\quad POINTS ON A LINE}
\addcontentsline{toc}{section}{5--5\quad Points on a Line}

We have just described the measurement process for getting a number for the
length of a segment. We now turn the question around and ask whether, given a
number, there is a segment with that number for its length. Of course, we would
like this to be true, but our postulates so far do not tell us that it is true.
The process of measurement describes a method for ``sneaking up'' on the point
$B$, but if we start out with a number we are not assured that there is a point
to sneak up on! To have this property, we must postulate it. This postulate
tells us that there are enough points on the line.

\medskip
IV--2. \emph{The completeness postulate.} For every line $l$ and any point $A$
on $l$, and for any positive real number $x$, there is, on a given side of $A$,
a point $B$ such that $\sg{AB} = x$.

\medskip
The name of the last postulate is meant to signify that our lines are complete,
that is, that they have all the points on them we expect them to have.

\exercisegroup{5--8}

\begin{multicols}{2}
\begin{exlist}
\item Draw a segment that is to be the unit segment. Draw another segment that
      is 1.3 units in length.
\item With the unit segment of Exercise 1, draw a segment of length
      $2 + \frac{1}{2} + \frac{1}{4} + \frac{1}{8}$ units.
\end{exlist}
\end{multicols}

%========================================================= 5-6
\section*{5--6\quad LENGTH AND CONGRUENCE}
\addcontentsline{toc}{section}{5--6\quad Length and Congruence}

When we started talking about congruent segments we had in mind a notion of
size. Now that we have lengths for segments it should turn out that congruent
segments have the same length, and conversely that segments of the same length
are congruent. We now prove these theorems. To do so we need the following
theorem (4--3) that we proved in Chapter 4. \emph{If $AB \cong A'B'$ and $ACB$
and if $C'$ is on the same side of $A'$ as $B'$ and $AC \cong A'C'$, then
$A'C'B'$.}

\medskip\noindent
\textbf{*Theorem 5--1.}\ \emph{If $AB \cong A'B'$, then
$\sg{AB} = \sg{A'B'}$.}
\refstepcounter{theorem}
\par\medskip

\noindent\emph{Proof.} To get the length of segment $AB$, we establish a
sequence of points \emph{between} $A$ and $B$. Let us designate this sequence
by $C_1$, $C_2$, $C_3,\ldots$. Now, consider the corresponding sequence,
$C_1'$, $C_2'$, $C_3',\ldots$, where $AC_1 \cong A'C_1'$,
$AC_2 \cong A'C_2'$, $AC_3 \cong A'C_3',\ldots$, and all the points $C_1'$,
$C_2'$, $C_3',\ldots$ are on the same side of $A'$ as $B'$. Using Theorem
4--3, we see that all the points $C_1'$, $C_2'$, $C_3',\ldots$, are between
$A'$ and $B'$; therefore, this sequence determines the same real number
(decimal) as does the sequence $C_1$, $C_2$, $C_3,\ldots$.

\medskip\noindent
\textbf{*Theorem 5--2.}\ \emph{If $\sg{AB} = \sg{A'B'}$, then
$AB \cong A'B'$.}
\refstepcounter{theorem}
\par\medskip

\noindent\emph{Proof.} This theorem is the converse of Theorem 5--1. We will
prove it by choosing a point $B''$, on the same side of $A$ as $B$, such that
$AB'' \cong A'B'$, and then proving that $B'' = B$ (Fig.~5--8). This will show
that $AB \cong A'B'$. Either $B'' = B$ or $B'' \neq B$. Since
$AB'' \cong A'B'$, then from Theorem 5--1 $\sg{AB''} = \sg{A'B'}$. We are given
that $\sg{AB} = \sg{A'B'}$. Hence $\sg{AB''} = \sg{AB}$. Thus we must show that
if $B'' \neq B$, we cannot have $\sg{AB''} = \sg{A'B'}$. But if $B'' \neq B$,
then, if the unit segment is smaller than $BB''$, the corresponding Archimedes
integers for $AB$ and $AB''$ would be different. This would lead to different
lengths for $AB$ and $AB''$. Hence $B''$ must be the same as $B$ and
$AB \cong A'B'$.

\FIGVEIGHT

These two theorems give us the right to use the words ``the same length'' in
place of ``congruent.'' Therefore we can use everyday language in talking about
congruence of segments.

\exercisegroup{5--9}

\begin{multicols}{2}
\begin{exlist}
\item $AB \cong CD$ and $CD \cong EF$, and $\sg{AB} = 7.26$ inches. What is
      $\sg{EF}$?
\item If $\sg{PQ}$ is 1.23 ft, and $VU \cong RS$, and $\sg{RS} = 1.23$ ft, is
      $QP \cong VU$?
\item If $AB \cong CD$, $CD \cong EF$, $EF \cong GH$, is $\sg{AB} = \sg{GH}$?
\item An isosceles triangle can be defined now in two ways. State the two
      definitions.
\end{exlist}
\end{multicols}

%========================================================= 5-7
\section*{5--7\quad MORE ABOUT INEQUALITIES}
\addcontentsline{toc}{section}{5--7\quad More about Inequalities}

In Chapter 4 we defined $>$ and $<$ for segments. In this chapter (Section
5--3) we defined $>$ and $<$ for numbers. Since ``same length'' means
``congruent,'' it should also be true that if $AB > CD$, then
$\sg{AB} > \sg{CD}$, and conversely.

\medskip\noindent
\textbf{*Theorem 5--3.}\ \emph{If $AB > CD$ then $\sg{AB} > \sg{CD}$, and
conversely, if $\sg{AB} > \sg{CD}$ then $AB > CD$.}
\refstepcounter{theorem}
\par\medskip

\noindent\emph{Proof.} The theorem could also be stated as ``$AB > CD$ if and
only if $\sg{AB} > \sg{CD}$.'' Or, again, as ``If $AB > CD$ then
$\sg{AB} > \sg{CD}$, and conversely.'' We first prove that $AB > CD$ implies
$\sg{AB} > \sg{CD}$. The hypothesis, $AB > CD$, means that if $E$ is on the
same side of $C$ as $D$, and $AB \cong CE$, then $D$ is between $C$ and $E$.
(See Fig.~5--9.) By Theorem 5--1, $\sg{AB} = \sg{CE}$. In the measurement
process, if the unit is chosen smaller than $DE$, the corresponding Archimedes'
integers for $CD$ and $CE$ will be different, and in fact the integer for $CD$
will be smaller than that for $CE$. Hence $\sg{CD} < \sg{CE}$ and, since
$\sg{CE} = \sg{AB}$, $\sg{CD} < \sg{AB}$.

\FIGVNINE

The converse is proved in a similar manner. We have given that
$\sg{AB} > \sg{CD}$. Choose $E$, on the same side of $C$ as $D$, such that
$AB \cong CE$. By Theorem 5--1, $\sg{AB} = \sg{CE}$. Now either $CDE$,
$D = E$, or $CED$. (Why?) We cannot have $D = E$. (Why?) An argument like the
first part of the proof will show we cannot have $CED$. Hence we must have
$CDE$. It follows from the definition of ``less than'' for segments that
$CD < CE \cong AB$.

\exercisegroup{5--10}

\begin{multicols}{2}
\begin{exlist}
\item If $ABC$ and $\sg{AB} = 5$, what can you say about $\sg{AC}$?
\item If $AB > EF$ and $\sg{EF} = 4$, what can you say about $\sg{AB}$?
\item If $AB > CD > EF$, what can you say about $\sg{AB}$ and $\sg{EF}$?
\item If $ABCD$, what can you say about $\sg{AD}$ and $\sg{BC}$?
\item If $\sg{AB} > \sg{AC}$, and $C$ is a point on the line containing $A$ and
      $B$, and $C$ lies on the same side of $A$ as $B$, where is the point $C$?
\end{exlist}
\end{multicols}

%========================================================= 5-8
\section*{5--8\quad CHANGE OF SCALE}
\addcontentsline{toc}{section}{5--8\quad Change of Scale}

In measurement we choose a fixed segment for a unit of length. What happens if
a different segment is chosen for the unit? This is a familiar situation
because we are all used to dealing with inches, feet, yards, or miles to
measure length. In Fig.~5--10, if $CD$ is the unit, the length of $AB$ is 3. If
$ED$ is the unit, the length of $AB$ is 6. This illustrates the following
theorem.

\FIGVTEN

\medskip\noindent
\textbf{Theorem 5--4.}\ \emph{If $\sg{AB} = x$ when $CD$ is the unit, and if
$\sg{CD} = y$ when $EF$ is the unit, then $AB = xy$ when $EF$ is the unit.}
\refstepcounter{theorem}
\par\medskip

We will not give the proof of this theorem for arbitrary $x$ and $y$ because it
is long to write out and quite hard if $y$ is an irrational number. But if $y$
is an integer, the proof is very easy and is really clear from Fig.~5--10,
where $y = 2$.

\exercisegroup{5--11}

\begin{multicols}{2}
\begin{exlist}
\item $\sg{AB} = 5$ when $CD$ is the unit, and $\sg{CD} = 2.3$ when $EF$ is the
      unit. What is $\sg{AB}$ when $EF$ is the unit?
\item The width of the room is 14 ft. What is the width in inches? How does
      this follow from Theorem 5--4?
\item One kilometer is 1000 meters. One meter is 100 centimeters. How many
      centimeters in a kilometer? Show how this follows from Theorem 5--4.
\item $\sg{PQ} = 6.5$ when $\sg{RS} = 1$. $\sg{RS} = \frac{1}{2}$ when
      $\sg{UV} = 1$. What is $\sg{PQ}$ when $\sg{UV} = 1$?
\item If $CD < EF$, how does $\sg{AB}$ with unit $CD$ compare with $\sg{AB}$
      with unit $EF$? In other words, if we use two different units and compute
      two lengths for one segment, does the smaller unit give the smaller or
      the larger length?
\end{exlist}
\end{multicols}

%========================================================= Review
\Vhead{Review of Chapter 5}
\addcontentsline{toc}{section}{Review of Chapter 5}
\markboth{REVIEW OF CHAPTER 5}{}

\begin{exlist}
\item In this chapter we have given definitions of the following: real number,
      rational number, irrational number, positive integer, greater than, less
      than, the process of measurement. Be sure you know these definitions.
\item This chapter has two new postulates, Archimedes' postulate and the
      completeness postulate. Go over them again and be sure you see what they
      mean.
\item In this chapter we have proved the following theorems. Study them again
      to be sure you know what they say. You should be able to state clearly
      the facts upon which each proof depends.
\end{exlist}

\begin{align*}
(AB \cong CD) &\iffx \sg{AB} = \sg{CD},\\
(AB > CD) &\iffx \sg{AB} > \sg{CD},\\
(AB < CD) &\iffx \sg{AB} < \sg{CD}.
\end{align*}

\begin{note}
In this chapter we have shown that congruence is equivalent to equality of
length. From this point on in our development of geometry you may use the
ordinary language of length in speaking of congruence.
\end{note}

%========================================================= Algebra review
\Vhead{Algebra Review}
\addcontentsline{toc}{section}{Algebra Review}

\begin{exlist}
\item Segment $XY$ has a length of 28 units. Point $T$ is between $X$ and $Y$
      and $\sg{XT} = \sg{YT} + 9$. Determine the lengths of $XT$ and $YT$.
\item Points $R$, $T$, and $S$ are in this order on a line. Segment $RT$ is
      twice as long as segment $TS$. The difference in their lengths is 11.
      Find the length of $RS$.
\item Points $A$, $B$, and $C$ are on one line. The length of $AC$ is 12 more
      than that of $AB$. The length of $BC$ is four times the length of $AB$.
      Find the lengths of $AB$ and $AC$ if (a) $A$ is between $B$ and $C$, and
      (b) $B$ is between $A$ and $C$.
\item Points $A$, $B$, $C$, and $D$ are in this order on one line.
      $\sg{BC} = \sg{AB} + 4$; $\sg{CD} = \sg{AB} + \sg{BC}$; $\sg{AD} = 32$.
      Find $\sg{AB}$.
\item The length of $XY$ is $\frac{2}{3}$ of the length of $XZ$. The length of
      $XZ$ is $\frac{3}{5}$ of the length of $YZ$. If these three points are
      collinear: (a) What is the order of the points on $l$? (b) Express
      $\sg{XY}$ as a fraction of $\sg{YZ}$. (c) Determine $\sg{XY}$ if
      $\sg{XZ} - \sg{XY} = 6$.
\item The point $M$ is the midpoint of segment $AB$, $N$ is the midpoint of
      $MB$, and $O$ is the midpoint of $NB$. If $\sg{AN} - \sg{NO} = 20$, find
      the length of $AB$.
\item $\sg{AC} = 2\sg{AB}$, $\sg{AB} + \sg{BC} = \sg{AC} + 2$. If
      $\sg{AB} > 1$, prove that the three points $A$, $B$, and $C$ are not on
      one line.
\item A square and equilateral triangle have equal perimeters. The length of a
      side of the triangle exceeds the length of a side of the square by six
      inches. Find the perimeter of each.
\item The sum of the perimeters of a square and a rectangle is 36 ft. The
      rectangle is twice as long as it is wide, and the width of the rectangle
      is 1 ft more than the length of a side of the square. Find the dimensions
      of each.
\item The point $B$ is the midpoint of segment $AC$. Point $X$ is between $B$
      and $C$, and $\sg{BX} = 3$ inches. Show that a square having a side
      congruent to $AB$ has an area of 9 in$^2$ more than a rectangle with
      sides congruent to $AX$ and $XC$.
\end{exlist}
\clearpage
% ==================== ch06 ====================
\setcounter{theorem}{0}\setcounter{definition}{0}\setcounter{construction}{0}\setcounter{exgroup}{0}\setcounter{footnote}{0}
\setcounter{chapter}{6}

\chapter*{\raggedright\Huge CONGRUENCE OF\\ANGLES AND TRIANGLES}
\markboth{CONGRUENCE OF ANGLES AND TRIANGLES}{}
\addcontentsline{toc}{chapter}{6.\ Congruence of Angles and Triangles}
\vspace{-1em}\noindent\rule{\linewidth}{0.6pt}\vspace{1em}
\figurekey

%========================================================= 6-1
\section*{6--1\quad INTRODUCTION}
\addcontentsline{toc}{section}{6--1\quad Introduction}

In the last two chapters we have been concerned with congruence of segments
only. But we also have a feeling for size of an angle---and, in fact, in
Chapter 1 we used a protractor to measure angles. If our geometry is to be
like the world of our experience, we must include in it postulates that will
enable us to compare the sizes of angles. These postulates form Postulate
Group V. There are four postulates in the group, and we take them up one at a
time. As in Chapter 4, we shall use the word ``congruent'' to give the idea of
\emph{the same size}.

%========================================================= 6-2
\section*{6--2\quad POSTULATE V--1, THE EXISTENCE POSTULATE FOR ANGLES}
\addcontentsline{toc}{section}{6--2\quad Postulate V--1, the existence postulate for angles}

This first postulate of Group V tells of the existence of angles congruent to
a given angle. Remember to think of the word ``congruent'' as meaning
\emph{the same size}.

\medskip
V--1 If $\ang{A}$ is not a straight angle, and if $r\VIpr$ is a ray from
$A\VIpr$ on a line $l\VIpr$, then on a given side of $l\VIpr$, there is exactly
one ray $s\VIpr$ from $A\VIpr$ such that the angle $A\VIpr$, with sides
$r\VIpr$ and $s\VIpr$, is congruent to angle $A$ (Fig.~6--1). We write this as
$\ang{A\VIpr} \cong \ang{A}$. If $\ang{A}$ is a straight angle, then the
straight angle at $A\VIpr$, with one side $r\VIpr$, is the only angle with one
side $r\VIpr$ congruent to angle $A$.
\medskip

\FIGVIONE

\exercisegroup{6--1}
\begin{multicols}{2}
\begin{exlist}
\item If $\ang{A}$ is not a straight angle, how many angles with side $r\VIpr$
  are congruent to $\ang{A}$? How many if $\ang{A}$ is a straight angle?
\item In a figure like 6--2 draw the possible positions of the side $s\VIpr$
  such that $\ang{A\VIpr} \cong \ang{A}$.
\item In a figure like 6--3, draw all possible positions of $s\VIpr$ such that
  $\ang{A\VIpr} \cong \ang{A}$.
\item If in Fig.~6--4, $A$ is a straight angle and $r\VIpr$ is a ray from
  $A\VIpr$, draw an angle, with one side $r\VIpr$, that is congruent to $A$.
  Is it the only such angle?
\item If $\ang{O}$ is a straight angle and $\ang{O\VIpr} \cong \ang{O}$, is
  $\ang{O\VIpr}$ a straight angle? Why?
\end{exlist}
\end{multicols}

\FIGVITWOTHREEFOUR

%========================================================= 6-3
\section*{6--3\quad POSTULATE V--2, THE THREE-PART ANGLE POSTULATE}
\addcontentsline{toc}{section}{6--3\quad Postulate V--2, the three-part angle postulate}

The next postulate is for angles what Postulate III--2 is for segments. It may
seem odd to have to state such seemingly obvious things, but remember we are
dealing with those imaginary things called points and lines, and so we must
state \emph{all} the properties we need.

\medskip
\noindent
\begin{tabular}{@{}p{2.2em}@{}p{0.86\linewidth}@{}}
V--2 & (a)\ \ $\ang{A} \cong \ang{A}$.\ \ (Congruence of angles is reflexive.)\\[2pt]
 & (b)\ \ If $\ang{A} \cong \ang{B}$, then $\ang{B} \cong \ang{A}$.
   (Congruence of angles is symmetric.)\\[2pt]
 & (c)\ \ If $\ang{A} \cong \ang{B}$, and $\ang{B} \cong \ang{C}$, then
   $\ang{A} \cong \ang{C}$. (Congruence of angles is transitive.)
\end{tabular}
\medskip

\exercisegroup{6--2}
\begin{multicols}{2}
\begin{exlist}
\item State each part of Postulate V--2 in words.
\item If $\ang{E} \cong \ang{F}$, and $\ang{F} \cong \ang{G}$, and
  $\ang{G} \cong \ang{H}$, show, by using part (c) of Postulate V--2, that
  $\ang{E} \cong \ang{H}$.
\item If $\ang{R} \cong \ang{S}$ and $\ang{T} \cong \ang{S}$, prove that
  $\ang{R} \cong \ang{T}$.
\item[\stex 4.] Use Postulate V--1 and (b) and (c) of Postulate V--2 to prove
  that for every $\ang{A}$, $\ang{A} \cong \ang{A}$.
\end{exlist}
\end{multicols}

%========================================================= 6-4
\section*{6--4\quad POSTULATE V--3, THE ADDITION AND SUBTRACTION OF ANGLES POSTULATE}
\addcontentsline{toc}{section}{6--4\quad Postulate V--3, the addition and subtraction of angles postulate}

The next postulate does for angles what Postulate III--3 does for segments. It
enables us to ``add'' and ``subtract'' angles. To avoid a complicated statement
in words, we explain the postulate by using Fig.~6--5, where the side common to
$\ang{B}$ and $\ang{C}$, except the vertex, is in the interior of $\ang{A}$,
and the same for $\ang{A\VIpr}$, $\ang{B\VIpr}$, and $\ang{C\VIpr}$.

\FIGVIFIVE

\medskip
\noindent
V--3 If the angles in any two of the pairs $\ang{A}$ and $\ang{A\VIpr}$,
$\ang{B}$ and $\ang{B\VIpr}$, $\ang{C}$ and $\ang{C\VIpr}$ are congruent to
each other, then the angles of the third pair are also congruent to each other.
In other words:
\begin{enumerate}[label=(\alph*), leftmargin=2.6em, itemsep=1pt, topsep=3pt]
\item If $\ang{A} \cong \ang{A\VIpr}$ and $\ang{B} \cong \ang{B\VIpr}$, then
  $\ang{C} \cong \ang{C\VIpr}$.
\item If $\ang{B} \cong \ang{B\VIpr}$ and $\ang{C} \cong \ang{C\VIpr}$, then
  $\ang{A} \cong \ang{A\VIpr}$.
\item If $\ang{C} \cong \ang{C\VIpr}$ and $\ang{A} \cong \ang{A\VIpr}$, then
  $\ang{B} \cong \ang{B\VIpr}$.
\end{enumerate}
\noindent The angle $A$ may be a straight angle.
\medskip

\exercisegroup{6--3}
\begin{multicols}{2}
\begin{exlist}
\item We have not yet in our geometry studied the measurement of angles in
  degrees. However, in this exercise, use the fact (which will be proved later)
  that angles congruent to each other have the same measure in degrees, and add
  or subtract degree measures as you are accustomed to do. Refer to Fig.~6--6.
  \begin{enumerate}[label=(\alph*), leftmargin=2.2em, itemsep=1pt, topsep=3pt]
  \item If $\ang{A}$ is $130^\circ$ and $\ang{B}$ is $45^\circ$, how many
    degrees in $\ang{C}$?
  \item If $\ang{B\VIpr}$ is $41^\circ$ and $\ang{C\VIpr}$ is $75^\circ$, how
    many degrees in $\ang{A\VIpr}$?
  \item If $\ang{A} \cong \ang{A\VIpr}$ and $\ang{A}$ is $150^\circ$, how many
    degrees in $\ang{A\VIpr}$?
  \item If $\ang{A} \cong \ang{A\VIpr}$ and $\ang{C} \cong \ang{C\VIpr}$, are
    $\ang{B}$ and $\ang{B\VIpr}$ congruent to each other?
  \item If $\ang{A}$ is $87^\circ$ and congruent to $\ang{A\VIpr}$, and if
    $\ang{C} \cong \ang{C\VIpr}$ and $\ang{C\VIpr}$ is $33^\circ$, how many
    degrees in $\ang{B}$?
  \end{enumerate}
\item In Fig.~6--7, if $\ang{POR} \cong \ang{SVU}$, and $\ang{QOR} \cong
  \ang{TVU}$, what can you say about $\ang{POQ}$?
\item In Fig.~6--8, if $\ang{BAC} \cong \ang{DAE}$, what can you say about
  $\ang{BAD}$?
\item In Fig.~6--9, if $\ang{SPU} \cong \ang{RPT}$, what can you say about
  $\ang{RPS}$?
\item Is Exercise 2 an example of ``adding'' or ``subtracting'' angles? What
  about Exercises 3 and 4?
\end{exlist}
\end{multicols}

\FIGVISIXSEVEN

\FIGVIEIGHTNINE

Supplementary and adjacent angles were defined in Definition 3--9. The
statement that $\ang{A}$ is \emph{supplementary} to $\ang{B}$ shall mean that
$\ang{A}$ is congruent to an angle that is \emph{supplementary and adjacent} to
$\ang{B}$.

\begin{theorem}
If two angles are congruent to each other, then their supplements are also
congruent to each other.
\end{theorem}

\emph{Proof.} Let the given angles be $A$ and $A\VIpr$, and let their
supplements be $B$ and $B\VIpr$, respectively, as in Fig.~6--10. We are given
that $\ang{A} \cong \ang{A\VIpr}$. Since any two straight angles are congruent
to each other, it follows at once from Postulate V--3 that $\ang{B} \cong
\ang{B\VIpr}$. (Why?)

\FIGVITENELEVENTWELVE

\begin{theorem}
Angles that are vertical to each other are also congruent to each other.
\end{theorem}

\emph{Proof.} Let the angles $A$ and $A\VIpr$ be vertical to each other, as in
Fig.~6--11, and consider $\ang{B}$, one of the other angles formed by the
lines. Then $\ang{B}$ is the supplement of $\ang{A}$ and also the supplement of
$\ang{A\VIpr}$. Since $\ang{B} \cong \ang{B}$, it follows from Theorem 6--1
that $\ang{A} \cong \ang{A\VIpr}$.

\begin{definition}
An angle (Fig.~6--12) that is congruent to its supplement is called a
\emph{right angle}.
\end{definition}

\begin{remark}
Note that the words \emph{congruent}, \emph{vertical}, \emph{equal}, and
\emph{supplementary} refer to relations between two objects, not to their
properties. That is, it makes no sense to say that $\ang{A}$ is a vertical
angle, or a congruent angle, or a supplementary angle. At times, we refer to
two angles $A$ and $B$ as equal angles, or vertical angles, or congruent
angles, etc. What we really mean is that $\ang{A}$ is equal to $\ang{B}$,
$\ang{A}$ is vertical to $\ang{B}$, $\ang{A}$ is congruent to $\ang{B}$, etc.
It is important to distinguish between properties and relations.
\end{remark}

\exercisegroup{6--4}
\begin{multicols}{2}
\begin{exlist}
\item If an angle is a right angle, what is its supplement?
\item If $\ang{1} \cong \ang{2}$ in Fig.~6--13, prove that $\ang{EAC} \cong
  \ang{DAB}$.
\item In Fig.~6--14 it is given that $\ang{B}$ is congruent to the supplement
  of $\ang{C}$ and $\ang{C} \cong \ang{D}$. Prove that $\ang{A} \cong \ang{E}$.
\item[\stex 4.] An angle is congruent to a right angle. Show that it, too, is a
  right angle.
\item[5.] In Fig.~6--15, the line $CD$ bisects $\ang{AOB}$. Show that it also
  bisects $\ang{EOF}$. (By ``bisect'' we mean that $\ang{AOC} \cong \ang{COB}$.)
\item[6.] Decide whether the following words refer to properties or relations:
  perpendicular, parallel, vertex, altitude, greater than, angle, segment,
  median, right (angle), straight (angle), bisector.
\end{exlist}
\end{multicols}

\FIGVITHIRTEENFOURTEENFIFTEEN

%========================================================= 6-5
\section*{6--5\quad POSTULATE V--4, THE SIDE-ANGLE-SIDE POSTULATE}
\addcontentsline{toc}{section}{6--5\quad Postulate V--4, the side-angle-side postulate}

This postulate (abbreviated S.A.S.) is one that combines congruence of segments
with congruence of angles. Up to this point, we have the two kinds of
congruence but no relationship between them. This is the connecting link that
will enable us to prove many theorems.

\medskip
\noindent
V--4 If two triangles, $\tri{ABC}$ and $\tri{A\VIpr B\VIpr C\VIpr}$, are such
that $AB \cong A\VIpr B\VIpr$, $AC \cong A\VIpr C\VIpr$, and $\ang{A} \cong
\ang{A\VIpr}$, then $\ang{B} \cong \ang{B\VIpr}$ and $\ang{C} \cong
\ang{C\VIpr}$ (Fig.~6--16).
\medskip

\FIGVISIXTEEN

This postulate tells us that if two triangles have two sides of one congruent
to two sides of the other, and if the angles formed by these two sides (called
the \emph{included} angles) are congruent to each other, then the other pairs
of corresponding angles are congruent.

\exercisegroup{6--5}
\begin{multicols}{2}
\begin{exlist}
\item In Fig.~6--17, if $AB \cong DF$, $BC \cong DE$, and $\ang{B} \cong
  \ang{D}$, what angle is congruent to $\ang{A}$? to $\ang{C}$?
\item In Fig.~6--18, if $PQ = US$, $QR = UT$, and $\ang{Q} \cong \ang{U}$, what
  congruences hold?
\item Draw figures to show that if $AB \cong A\VIpr B\VIpr$, and $\ang{A} \cong
  \ang{A\VIpr}$, but $AC \not\cong A\VIpr C\VIpr$ (the symbol $\not\cong$ is
  read ``is not congruent to''), then it need not be true that $\ang{C} \cong
  \ang{C\VIpr}$.
\item Draw figures to show that if $AB \cong A\VIpr B\VIpr$, and $AC \cong
  A\VIpr C\VIpr$, but $\ang{A} \not\cong \ang{A\VIpr}$, then it need not be
  true that $\ang{B} \cong \ang{B\VIpr}$.
\item If, in the figure $AC \cong AD$, and $\ang{CAB} \cong \ang{BAD}$, prove
  that $\ang{C} \cong \ang{D}$. (Fig.~6--19)
\item If, in the figure, $DO \cong OB$, and $AO \cong OC$, prove that $\ang{A}
  \cong \ang{C}$. (Fig.~6--20)
\item If, in the figure, $AB \cong CD$, and $\ang{CBA} \cong \ang{BCD}$, prove
  that $\ang{CAB} \cong \ang{BDC}$. (Fig.~6--21)
\item If, in the figure, $AD \cong BD$, and $\ang{ADC}$ is a right angle, prove
  that $\ang{A} \cong \ang{B}$. (Fig.~6--22)
\item \VIgp{$\ang{EBC} \cong \ang{ABD}$,\\ $AB \cong BC$,\\ $BE \cong BD$.}%
  {$\ang{A} \cong \ang{C}$.\\ \hbox{}\hfill(Fig.~6--23)}
\item \VIgp{$AD \cong DC$,\\ $\ang{1} \cong \ang{C}$,\\ $\ang{2} \cong
  \ang{C}$.}{$\ang{3} \cong \ang{4}$.\\ \hbox{}\hfill(Fig.~6--24)}
\item \VIgp{$\ang{1} \cong \ang{2}$,\\ $\ang{BCF} \cong \ang{ADF}$,\\ $CG \cong
  HD$,\\ $CF \cong DF$.}{$\ang{3} \cong \ang{4}$.\\ \hbox{}\hfill(Fig.~6--25)}
\item \VIgp{$\ang{1} \cong \ang{2}$,\\ $AE \cong CD$,\\ $AB \cong BC$.}%
  {$\ang{E} \cong \ang{D}$.\\ \hbox{}\hfill(Fig.~6--26)}
\item \VIgp{$\ang{1} \cong \ang{2}$,\\ $AC \cong BC$,\\ $AM \cong BM$.}%
  {$\ang{3} \cong \ang{4}$.\\ \hbox{}\hfill(Fig.~6--27)}
\item \VIgp{$\ang{ADC} \cong \ang{ABC}$,\\ $\ang{1} \cong \ang{2}$,\\ $AD \cong
  BC$.}{$\ang{A} \cong \ang{C}$.\\ \hbox{}\hfill(Fig.~6--28)}
\item \VIgp{$AD \cong DC$,\\ $\ang{1} \cong \ang{2}$.}{$\ang{A} \cong
  \ang{C}$.\\ \hbox{}\hfill(Fig.~6--29)}
\item[\stex 16.] \VIgp{$BF \cong ED$,\\ $AE \cong CF$,\\ $\ang{1} \cong
  \ang{2}$.}{$\ang{3} \cong \ang{4}$,\\ $\ang{5} \cong \ang{6}$.\\
  \hbox{}\hfill(Fig.~6--30)}
\item[\stex 17.] \VIgp{$AB \cong AD$,\\ $BO \cong OD$,\\ $\ang{OAB} \cong
  \ang{OAD}$.}{$\ang{1} \cong \ang{2}$.\\ \hbox{}\hfill(Fig.~6--31)}
\end{exlist}
\end{multicols}

\FIGVISEVENTEENEIGHTEEN

\FIGVININETEENTWENTY

\FIGVITWENTYONETWENTYTWO

\FIGVITWENTYTHREETWENTYFOUR

\FIGVITWENTYFIVETWENTYSIXTWENTYSEVEN

\FIGVITWENTYEIGHTTWENTYNINE

\FIGVITHIRTYTHIRTYONE

%========================================================= 6-6
\section*{6--6\quad ISOSCELES TRIANGLES}
\addcontentsline{toc}{section}{6--6\quad Isosceles triangles}

We can now prove a theorem about isosceles triangles, for which we need a
definition. The definition of isosceles triangle is the same as that given in
Chapter 1, Exercise Group 1--4:

\begin{definition}
A triangle is called \emph{isosceles} if two of its sides are congruent (of
equal length). The \emph{base angles} of the triangle are those opposite the
congruent sides.
\end{definition}

\begin{theorem}
In every isosceles triangle, the base angles are congruent to each other.
\end{theorem}

\emph{Proof.} Consider any isosceles triangle, and label the vertices in two
different ways, as shown in Fig.~6--32. Let $AB \cong AC$. Let us look at
the triangle both as $\tri{ABC}$ and as $\tri{A\VIpr B\VIpr C\VIpr}$. We have
$AB \cong AC \cong A\VIpr B\VIpr$, and $AC \cong AB = A\VIpr C\VIpr$, and
$\ang{A} \cong \ang{A\VIpr}$. Therefore, by the S.A.S. postulate, $\ang{B}
\cong \ang{B\VIpr}$, that is, $\ang{B} \cong \ang{C}$.

This proof is quite interesting because it illustrates very well how good
notation can help. In the proof we look at the same triangle in two different
ways, and once we have done this, the proof is very short.

\FIGVITHIRTYTWO

\begin{definition}
A \emph{midpoint} of a segment $AB$ is a point $M$ of the segment, such that
$AM \cong MB$.
\end{definition}

\exercisegroup{6--6}
\begin{multicols}{2}
\begin{exlist}
\item Using Theorem 6--3, prove that an equilateral triangle is equiangular.
\item If $PQ = PR$, and $\ang{Q}$ is $44^\circ$, how many degrees in $\ang{R}$?
  (We have not yet studied degrees in our geometry, but for this problem you
  may use what you already know from Chapter 1.) (Fig.~6--33)
\item If, in the figure, $AB \cong AC$ and $BD \cong DC$, prove that
  \begin{center}
  $\ang{ABC} \cong \ang{ACB}$,\\ $\ang{CBD} \cong \ang{BCD}$,\\
  $\ang{ABD} \cong \ang{ACD}$.\\ \hbox{}\hfill(Fig.~6--34)
  \end{center}
\item \VIgp{$AC \cong BC$,\\ $AD \cong DB$.}{$\ang{1} \cong \ang{2}$.\\
  \hbox{}\hfill(Fig.~6--35)}
\item \VIgp{$\ang{ABC} \cong \ang{ABE}$,\\ $CB \cong BE$.}{$\ang{C} \cong
  \ang{E}$.\\ \hbox{}\hfill(Fig.~6--36)}
\item \VIgp{$AC \cong DF$,\\ $BD \cong CE$,\\ $\ang{1} \cong \ang{2}$.}%
  {$\ang{A} \cong \ang{F}$.\\ \hbox{}\hfill(Fig.~6--37)}
\item In the figure $AD \cong EC$, and $BE \cong BD$. Show that $\ang{A} \cong
  \ang{C}$. (Fig.~6--38)
\item \VIgp{$RS \cong RT$,\\ $US \cong UT$.}{$\ang{RSU} \cong \ang{RTU}$.\\
  \hbox{}\hfill(Fig.~6--39)}
\item \VIgp{$AC \cong BC$,\\ $CD \cong CE$.}{$\ang{1} \cong \ang{2}$.\\
  \hbox{}\hfill(Fig.~6--40)}
\item \VIgp{$AO \cong OD$,\\ $AF \cong FD$,\\ $BO \cong OC$,\\ $BF \cong FC$.}%
  {$\ang{EAB} \cong \ang{EDC}$.\\ \hbox{}\hfill(Fig.~6--41)}
\item In the figure, $AC \cong BC$, $CG \cong CF$, and $\ang{1} \cong \ang{2}$.
  Prove that $\ang{3} \cong \ang{4}$. (Fig.~6--42)
\end{exlist}
\end{multicols}

\FIGVITHIRTYTHREETHIRTYFOURTHIRTYFIVE

\FIGVITHIRTYSIXTHIRTYSEVEN

\FIGVITHIRTYEIGHTTHIRTYNINE

\FIGVIFORTYFORTYONEFORTYTWO

%========================================================= 6-7
\section*{6--7\quad CONGRUENCE OF TRIANGLES}
\addcontentsline{toc}{section}{6--7\quad Congruence of triangles}

We now define congruence for triangles and prove some theorems.

\begin{definition}
Two triangles are said to be \emph{congruent} to each other if they can be
labeled $\tri{ABC}$ and $\tri{A\VIpr B\VIpr C\VIpr}$, such that
$AB \cong A\VIpr B\VIpr$, $AC \cong A\VIpr C\VIpr$, $BC \cong B\VIpr C\VIpr$,
$\ang{A} \cong \ang{A\VIpr}$, $\ang{B} \cong \ang{B\VIpr}$, and $\ang{C} \cong
\ang{C\VIpr}$. We then write $\tri{ABC} \cong \tri{A\VIpr B\VIpr C\VIpr}$ and
read this as ``triangle $ABC$ is congruent to triangle $A\VIpr B\VIpr
C\VIpr$.'' $\ang{A}$ corresponds to $\ang{A\VIpr}$, $\ang{B}$ to
$\ang{B\VIpr}$, side $AC$ to side $A\VIpr C\VIpr$, etc.
\end{definition}

\exercise{6--7}
Prove that congruence of triangles is transitive.

\begin{theorem}
If the vertices of two triangles may be labeled $A$, $B$, $C$ and $A\VIpr$,
$B\VIpr$, $C\VIpr$, such that $AB \cong A\VIpr B\VIpr$, $AC \cong A\VIpr
C\VIpr$, $\ang{A} \cong \ang{A\VIpr}$, then $\tri{ABC} \cong \tri{A\VIpr
B\VIpr C\VIpr}$.
\end{theorem}

\begin{remark}
This is the first of the three congruence theorems for triangles. It may also
be stated as follows: \emph{if two triangles are such that two sides and the
included angle of one are congruent, respectively, to two sides and the
included angle of the other, then the triangles are congruent.} This theorem
is often called the ``side-angle-side theorem'' and is abbreviated S.A.S.
\end{remark}

\emph{Proof.} From the S.A.S. postulate, we know that $\ang{B} \cong
\ang{B\VIpr}$ and $\ang{C} \cong \ang{C\VIpr}$. It only remains to be shown
that the third sides are congruent, that is, that $BC \cong B\VIpr C\VIpr$. We
prove this by considering the point $D\VIpr$ on ray $C\VIpr B\VIpr$, such that
$CB \cong C\VIpr D\VIpr$, as in Fig.~6--43, and then showing that $D\VIpr =
B\VIpr$.

\FIGVIFORTYTHREE

\begin{steps}
\step{There is a point $D\VIpr$ on the ray $C\VIpr B\VIpr$ such that $CB \cong
      C\VIpr D\VIpr$.}{Existence postulate for segments.}
\step{$AC \cong A\VIpr C\VIpr$, and $AB \cong A\VIpr B\VIpr$.}{Given.}
\step{$\ang{B\VIpr A\VIpr C\VIpr} \cong \ang{BAC}$.}{Given.}
\step{$\ang{C} \cong \ang{C\VIpr}$, and $\ang{B} \cong \ang{B\VIpr}$.}{S.A.S. postulate.}
\step{$\ang{BAC} \cong \ang{D\VIpr A\VIpr C\VIpr}$.}{By S.A.S. postulate, considering
      $\tri{ACB}$ and $\tri{A\VIpr C\VIpr D\VIpr}$.}
\step{Ray $A\VIpr B\VIpr$ is the same as the ray $A\VIpr D\VIpr$.}{Existence postulate for angles.}
\step{$B\VIpr = D\VIpr$.}{The lines $A\VIpr B\VIpr$ and
      $B\VIpr C\VIpr$ intersect in exactly one point by Theorem 3--1. This
      point is $B\VIpr$ and is also $D\VIpr$.}
\step{$BC \cong B\VIpr C\VIpr$.}{Statements 1 and 7.}
\step{$\tri{ABC} \cong \tri{A\VIpr B\VIpr C\VIpr}$.}{Definition 6--4.}
\end{steps}

In the exercises that follow we will make use of the congruence theorems. In
some exercises it will be convenient to know that \emph{``halves'' of congruent
segments are congruent to each other}. That is, if $AB \cong CD$ and $M$ and $N$
are midpoints of $AB$ and $CD$ respectively, then $AM \cong CN$. To see this we
note that $\sg{AB} = \sg{CD}$ and by the properties of numbers $\sg{AM} =
\sg{CN}$ so that $AM \cong CN$.

\exercisegroup{6--8}
\begin{multicols}{2}
\begin{exlist}
\item Given the figure, let $\ang{1} \cong \ang{2}$ and $AB \cong BC$. Prove
  that $AD \cong DC$. (Fig.~6--44)
\item Given the figure, let $\ang{1} \cong \ang{2}$ and $EF \cong GH$. Prove
  that $\tri{EGH} = \tri{HFE}$. (Fig.~6--45)
\item Given the figure, let $AC \cong AD$ and $AB$ bisect $\ang{A}$. Prove that
  $\ang{CBA} \cong \ang{DBA}$. (Fig.~6--46)
\item Given the figure, let $AB \cong AD$ and $BC \cong CD$. Prove that
  $\tri{ABC} \cong \tri{ADC}$. (Fig.~6--47)
\item If, in the figure, $AB \cong AC$, $D$ is the midpoint of $AB$, and $E$ is
  the midpoint of $AC$, prove that $\tri{ABE} \cong \tri{ACD}$. (Fig.~6--48)
\item If, in the figure, $AB \cong AD$, and $BC \cong CD$, prove that
  $\ang{ABC} \cong \ang{ADC}$ and $\ang{BAC} \cong \ang{DAC}$. (Fig.~6--49)
\item In the figure, $AB \cong BC$, and $D$, $E$ are the midpoints of $AB$ and
  $BC$ respectively. Prove that $AE \cong CD$. (Fig.~6--50)
\item Given: $EB \cong BD$, $B$ is the midpoint of $AC$, and $\ang{EBC} \cong
  \ang{DBA}$. Prove that $AE \cong CD$. (Fig.~6--51)
\item Prove that $BD \cong AC$ if $\ang{DAB} \cong \ang{CBA}$ and $AD \cong
  BC$. (Fig.~6--52)
\item In the figure, $AD \cong EC$, and $BE \cong BD$. Prove that $\tri{ABC}$
  is isosceles. (Fig.~6--53)
\item A farmer, wishing to find the distance between posts at $A$ and $B$ on
  opposite sides of a barn, does the following. He places a stake at $C$, and
  on the line through $B$ and $C$ places a stake at $D$, such that $\sg{BC} =
  \sg{DC}$. Then on the line through $A$ and $C$, he places a stake at $E$,
  such that $\sg{AC} = \sg{CE}$. He measures $DE$ and finds the length is
  78~ft. Prove that $\sg{AB} = 78$~ft. (Fig.~6--54)
\end{exlist}
\end{multicols}

\FIGVIFORTYFOURFORTYFIVE

\FIGVIFORTYSIXFORTYSEVEN

\FIGVIFORTYEIGHTFORTYNINE

\FIGVIFIFTYFIFTYONE

\FIGVIFIFTYTWOFIFTYTHREEFIFTYFOUR

%========================================================= 6-8
\section*{6--8\quad THE SECOND CONGRUENCE THEOREM}
\addcontentsline{toc}{section}{6--8\quad The second congruence theorem}

The first congruence theorem was the side-angle-side theorem (S.A.S.). The
second one says that two triangles are congruent if two angles and the included
side of one are, respectively, congruent to two angles and the included side of
the other. This theorem is called the ``angle-side-angle theorem'' and is
abbreviated A.S.A.

\begin{theorem}
If two triangles $\tri{ABC}$ and $\tri{A\VIpr B\VIpr C\VIpr}$ are such that
$\ang{A} \cong \ang{A\VIpr}$, $\ang{C} \cong \ang{C\VIpr}$, and $AC \cong
A\VIpr C\VIpr$, then $\tri{ABC} \cong \tri{A\VIpr B\VIpr C\VIpr}$.
\end{theorem}

\FIGVIFIFTYFIVE

\emph{Proof.} We consider the point $D\VIpr$ on the ray $A\VIpr B\VIpr$
(Fig.~6--55), such that $AB \cong A\VIpr D\VIpr$. Then it is true that
$\tri{ABC} \cong \tri{A\VIpr D\VIpr C\VIpr}$. (Why?) Then, if $D\VIpr \neq
B\VIpr$, we have both $\ang{D\VIpr C\VIpr A\VIpr}$ and $\ang{B\VIpr C\VIpr
A\VIpr}$ congruent to $\ang{C}$. This is impossible, so $B\VIpr = D\VIpr$.

\begin{steps}
\step{There is a point $D\VIpr$ on the ray $A\VIpr B\VIpr$, such that $AB \cong
      A\VIpr D\VIpr$.}{Existence postulate for segments.}
\step{$\ang{A} \cong \ang{A\VIpr}$, $AC \cong A\VIpr C\VIpr$.}{Given.}
\step{$\ang{C} \cong \ang{D\VIpr C\VIpr A\VIpr}$.}{S.A.S. postulate, applied to
      $\tri{ABC}$ and $\tri{A\VIpr D\VIpr C\VIpr}$.}
\step{$\ang{C} \cong \ang{B\VIpr C\VIpr A\VIpr}$.}{Given.}
\step{Ray $C\VIpr D\VIpr$ $=$ ray $C\VIpr B\VIpr$.}{Existence postulate for angles.}
\step{$D\VIpr = B\VIpr$.}{Two lines intersect in at most
      one point.}
\step{$\tri{ABC} \cong \tri{A\VIpr B\VIpr C\VIpr}$.}{S.A.S. theorem.}
\end{steps}

\exercisegroup{6--9}
\begin{multicols}{2}
\begin{exlist}
\item If, in the figure, $\ang{1} \cong \ang{2}$, and $\ang{3} \cong \ang{4}$,
  prove that $\tri{ABD} \cong \tri{CBD}$. (Fig.~6--56)
\item If, in the figure, $\ang{1} \cong \ang{2}$, and $\ang{3} \cong \ang{4}$,
  prove that $AB = CD$. (Fig.~6--57)
\item If, in the figure, $\ang{1} \cong \ang{2}$, and $AB \cong BC$, prove that
  $AE \cong DC$. (Fig.~6--58)
\item Draw a triangle. Let one angle be $31^\circ$, a second angle be
  $53^\circ$, and their included side be 7 inches.
\item If, in the figure, $\ang{ABC} \cong \ang{ABE}$, and $\ang{CDB} \cong
  \ang{EDB}$, prove that $CD \cong ED$. (Fig.~6--59)
\item If, in the figure, $AB \cong BC$, and $\ang{A} \cong \ang{C}$, prove that
  $\tri{ABD} \cong \tri{CBE}$. (Fig.~6--60)
\item If, in the figure, $\ang{A} \cong \ang{D}$ and $AB \cong BD$, prove that
  $AE \cong CD$. (Fig.~6--61)
\item If, in the figure, $AF \cong CD$, $\ang{A} \cong \ang{D}$, and
  $\ang{ACB} \cong \ang{DFE}$, prove that $\tri{ABC} \cong \tri{DEF}$.
  (Fig.~6--62)
\item Draw triangles $ABC$ whose parts are:
  \begin{enumerate}[label=(\alph*), leftmargin=2.2em, itemsep=1pt, topsep=3pt]
  \item $\ang{A} = 36^\circ$; $\ang{B} = 45^\circ$; $AB = 3$ in.
  \item $\ang{A} = 60^\circ$; $\ang{C} = 30^\circ$; $AC = 4$ in.
  \item $\ang{C} = 33^\circ$; $\ang{B} = 33^\circ$; $BC = 5/2$ in.
  \end{enumerate}
\item \VIgp{$\ang{3} \cong \ang{4}$,\\ $AB \cong AD$,\\ $AC$ bisects
  $\ang{BAD}$.}{$BC \cong CD$.\\ \hbox{}\hfill(Fig.~6--63)}
\item In the figure, $AF \cong FB$, and $AE \cong DB$. Prove that $\ang{3}
  \cong \ang{4}$. (Fig.~6--64)
\item Prove that $\ang{1} \cong \ang{2}$ if $AB \cong BC$ and $AD \cong DC$.
  (Fig.~6--65)
\item If $AH \cong HB$, and $AG \cong GB$, prove that $\ang{3} \cong \ang{4}$.
  (Fig.~6--66)
\item \VIgp{$\ang{1} \cong \ang{2}$,\\ $AC \cong BC$,\\ $DF \cong EG$.}%
  {$\tri{HIJ}$ is isosceles.\\ \hbox{}\hfill(Fig.~6--67)}
\end{exlist}
\end{multicols}

\FIGVIFIFTYSIX

\FIGVIFIFTYSEVENFIFTYEIGHT

\FIGVIFIFTYNINESIXTY

\FIGVISIXTYONESIXTYTWOSIXTYTHREE

\FIGVISIXTYFOURSIXTYFIVE

\FIGVISIXTYSIXSIXTYSEVEN

In Section 6--6 we proved that the base angles of an isosceles triangle are
congruent. We can now easily prove the converse.

\begin{theorem}
If a triangle has two angles congruent to each other, then the sides opposite
these angles are congruent.
\end{theorem}

\emph{Proof.} Suppose that the triangle is $\tri{ABC}$, and that $\ang{A} \cong
\ang{C}$. Let the vertices be labeled $A\VIpr$, $B\VIpr$, $C\VIpr$, as in
Fig.~6--68. Then we have $\ang{A} \cong \ang{C} \cong \ang{A\VIpr}$, and
$\ang{C} \cong \ang{A} \cong \ang{C\VIpr}$, and $AC \cong CA \cong A\VIpr
C\VIpr$. Hence by the A.S.A. congruence theorem, $\tri{ABC} \cong \tri{A\VIpr
B\VIpr C\VIpr}$. Therefore $AB \cong A\VIpr B\VIpr \cong CB$.

\FIGVISIXTYEIGHT

\exercise{6--10}
Prove that an equiangular triangle is equilateral.

%========================================================= 6-9
\section*{6--9\quad THE THIRD CONGRUENCE THEOREM}
\addcontentsline{toc}{section}{6--9\quad The third congruence theorem}

The two previous congruence theorems were the S.A.S. and A.S.A. theorems. The
third of the three fundamental congruence theorems states that two triangles
are congruent if the sides of one are congruent to the sides of the other. It
is called the ``side-side-side theorem'' and is abbreviated S.S.S.

\begin{theorem}
If the sides of one triangle are congruent, respectively, to the sides of
another triangle, then the triangles are congruent.
\end{theorem}

\FIGVISIXTYNINE

\emph{Proof.} Suppose that the triangles are labeled $\tri{ABC}$ and
$\tri{A\VIpr B\VIpr C\VIpr}$, as in Fig.~6--69(a) and 6--69(b), and that $AB
\cong A\VIpr B\VIpr$, $AC \cong A\VIpr C\VIpr$, and $BC \cong B\VIpr C\VIpr$.
We have only to prove that the angles are congruent. There is a ray $AD$, as in
Fig.~6--69(c), from $A$ on the opposite side of $AC$ from $B$, such that
$\ang{DAC} \cong \ang{A\VIpr}$. On $AD$, there is a point $E$, such that $AE
\cong A\VIpr B\VIpr$.

\begin{steps}
\step{$\ang{EAC} \cong \ang{A\VIpr}$.}{The ray $AE$ was chosen this way.}
\step{$AE \cong A\VIpr B\VIpr$.}{$E$ was so chosen.}
\step{$AC \cong A\VIpr C\VIpr$.}{Given.}
\step{$\tri{AEC} \cong \tri{A\VIpr B\VIpr C\VIpr}$.}{S.A.S. theorem.}
\end{steps}

Now consider the segment $BE$ and the triangles $ABE$ and $BEC$, as in
Fig.~6--70.

\FIGVISEVENTY

\begin{steps}
\stepx{5}{$\tri{ABE}$ is isosceles.}{$AB \cong A\VIpr B\VIpr
      \cong AE$.}
\stepx{6}{$\ang{1} \cong \ang{3}$.}{Statement 5 and Theorem 6--3.}
\stepx{7}{$\ang{2} \cong \ang{4}$.}{Same argument as used for
      $\ang{1}$ and $\ang{3}$.}
\stepx{8}{$\ang{B} \cong \ang{E}$.}{Addition of angles postulate.}
\stepx{9}{$\tri{ABC} \cong \tri{AEC}$.}{S.A.S. theorem.}
\stepx{10}{$\tri{ABC} \cong \tri{A\VIpr B\VIpr C\VIpr}$.}{Statements 4 and 9.}
\end{steps}

\begin{remark}
In our proof, we have tacitly assumed that $BE$ intersects the segment $AC$ at
a point between $A$ and $C$. This need not be the case. We could have either
Fig.~6--71(a) or 6--71(b).

\par\medskip\noindent\emph{Problem.} Prove both cases.
\end{remark}

\FIGVISEVENTYONE

\exercisegroup{6--11}
\begin{multicols}{2}
\begin{exlist}
\item You can conclude from Theorem 6--7 that a triangle is \emph{rigid}. Why?
  Why is a brace put on a gate? (Fig.~6--72)
\item If, in the figure, $\sg{AB} = \sg{BC}$, and $\sg{AD} = \sg{DC}$, prove
  that $\ang{ABD} \cong \ang{DBC}$. (Fig.~6--73)
\item If, in the figure, $\sg{AB} = \sg{DC}$, and $AD \cong BC$, prove that
  $\ang{CAB} \cong \ang{ACD}$ and $\ang{D} = \ang{B}$. (Fig.~6--74)
\item If, in the figure, $AB \cong BC$, and $AD \cong DC$, prove that $BD$
  bisects $\ang{ABC}$. (Fig.~6--75)
\item If, in the figure, $AB \cong AD$, and $BC \cong CD$, prove that $\ang{1}
  \cong \ang{2}$, $\tri{ABE} \cong \tri{ADE}$, and $\ang{AEB}$ is a right
  angle. (Fig.~6--76)
\item If, in the figure, $\sg{AB} = \sg{CD}$, and $\sg{BC} = \sg{AD}$, and
  $\sg{AF} = \sg{EC}$, prove that $\sg{BF} = \sg{ED}$. (Fig.~6--77)
\item If, in the figure, $\tri{ABC}$ is isosceles with vertex $B$, and $AD
  \cong EC$, prove that $BE \cong BD$. (Fig.~6--78)
\item If, in the figure, $AB \cong AD$, and $BC \cong CD$, prove that $\ang{B}
  \cong \ang{D}$. (Fig.~6--79)
\item \VIgp{$CE \cong CF$,\\ $AE \cong BF$,\\ $AD \cong BD$.}{$\ang{1} \cong
  \ang{2}$.\\ \hbox{}\hfill(Fig.~6--80)}
\item \VIgp{$\ang{1} \cong \ang{2}$,\\ $\ang{DCE} \cong \ang{DCF}$,\\ $AC \cong
  BC$.}{$AD \cong BD$.\\ \hbox{}\hfill(Fig.~6--81)}
\item If $AC \cong BD$, and $AD \cong BC$, show that $\ang{C} \cong \ang{D}$.
  (Fig.~6--82)
\item If $AB \cong BC$, and $BD \cong BE$, show that $\ang{A} \cong \ang{C}$.
  (Fig.~6--83)
\item \VIgp{$AC \cong BC$,\\ $AD \cong DB$.}{$AO \cong BO$.\\
  \hbox{}\hfill(Fig.~6--84)}
\item Prove that $AD \cong BE$ if $AC \cong BC$ and $\ang{1} \cong \ang{2}$.
  (Fig.~6--85)
\item \VIgp{$AD \cong BC$,\\ $AB \cong DC$,\\ $AM \cong MC$.}{$EM \cong MF$.\\
  \hbox{}\hfill(Fig.~6--86)}
\item \VIgp{$\ang{1} \cong \ang{2}$,\\ $AE \cong DB$.}{$\ang{3} \cong
  \ang{4}$.\\ \hbox{}\hfill(Fig.~6--87)}
\item \VIgp{$BE \cong CE$,\\ $\ang{1} \cong \ang{2}$.}{$\ang{3} \cong
  \ang{4}$.\\ \hbox{}\hfill(Fig.~6--88)}
\item Prove that the medians from the base angles of an isosceles triangle are
  congruent.
\item \VIgp{$AD \cong BC$,\\ $AB \cong CD$,\\ $\ang{1} \cong \ang{2}$.}%
  {$DE \cong BF$.\\ \hbox{}\hfill(Fig.~6--89)}
\item \VIgp{$AD \cong CB$,\\ $BD \cong AC$.}{$\tri{DEC}$ is isosceles.\\
  \hbox{}\hfill(Fig.~6--90)}
\item \VIgp{$AB \cong CD$,\\ $AE \cong CF$,\\ $\ang{3} \cong \ang{4}$.}%
  {$\ang{1} \cong \ang{2}$.\\ \hbox{}\hfill(Fig.~6--91)}
\item If triangles $ABC$ and $ADE$ are isosceles with vertex $A$, and $\ang{1}
  \cong \ang{2}$, prove that $\ang{3} \cong \ang{4}$, $FG \cong FH$, and $GI
  \cong HJ$. (Fig.~6--92)
\item Prove that the line joining the vertex of an isosceles triangle to the
  midpoint of the base bisects the vertex angle, and conversely.
\end{exlist}
\end{multicols}

\FIGVISEVENTYTWOSEVENTYTHREESEVENTYFOUR

\FIGVISEVENTYFIVESEVENTYSIX

\FIGVISEVENTYSEVENSEVENTYEIGHT

\FIGVISEVENTYNINEEIGHTY

\FIGVIEIGHTYONEEIGHTYTWO

\FIGVIEIGHTYTHREEEIGHTYFOUR

\FIGVIEIGHTYFIVEEIGHTYSIX

\FIGVIEIGHTYSEVENEIGHTYEIGHT

\FIGVIEIGHTYNINE

\FIGVININETYNINETYONE

\FIGVININETYTWO

%========================================================= 6-10
\section*{6--10\quad RIGHT ANGLES}
\addcontentsline{toc}{section}{6--10\quad Right angles}

In Definition 6--1 we defined right angles, but so far \emph{we do not know
that there are such angles}. We now show how to use our congruence theorem
(S.A.S.) to show that there is a right angle in our geometry.

\FIGVININETYTHREE

\emph{Proof.} On a line $l$, choose a point $O$ and let $\ang{AOB}$ be an angle
(not straight) with one side on $l$ (Fig.~6--93). Now consider the angle
congruent to $\ang{AOB}$, with vertex $O$, one side $OB$ on the line $l$ and
the other on the opposite side of $l$ from $A$. Choose the point $A\VIpr$ on
this last side, such that $OA \cong OA\VIpr$. Because $A$ and $A\VIpr$ are on
opposite sides of $l$, the segment $AA\VIpr$ intersects the line in some point
$C$. There are three cases to consider, as shown in Fig.~6--94.

\FIGVININETYFOUR

If $C$ is on the same side of $O$ as $B$, then by the S.A.S. theorem,
$\tri{AOC} \cong \tri{A\VIpr OC}$. Hence $\ang{OCA} \cong \ang{OCA\VIpr}$.
Since these angles are congruent and supplementary to each other, they are
right angles. In case $C = O$, we have $\ang{AOB} \cong \ang{A\VIpr OB}$, and
hence they are right angles. If $C$ is on the opposite side of $O$ from $B$,
then the angles $AOC$ and $A\VIpr OC$ are congruent. (Why?) We may apply the
S.A.S. congruence theorem. $\ang{ACO} \cong \ang{A\VIpr CO}$, and these are
right angles.

\exercisegroup{6--12}
\begin{multicols}{2}
\begin{exlist}
\item Perform a construction to draw a right angle.
\item If, in the figure, $PQ \cong PS$, and $\ang{QPR} \cong \ang{SPR}$, prove
  that $\ang{PRQ}$ is a right angle. (Fig.~6--95)
\item If, in the figure, $AB \cong AD$, and $BC \cong CD$, prove that
  $\ang{AEB}$ is a right angle. (Fig.~6--96)
\end{exlist}
\end{multicols}

\FIGVININETYFIVENINETYSIX

%========================================================= 6-11
\section*{6--11\quad INEQUALITIES FOR ANGLES}
\addcontentsline{toc}{section}{6--11\quad Inequalities for angles}

Even before studying geometry you used angles of the same size and compared
angles of different sizes (by saying, for example, that one angle is larger
than another). In this section, we see how to compare the sizes of angles. To
do this, we have to define $>$ (is greater than) and $<$ (is less than) for
angles. The idea we use in forming this definition is illustrated in
Fig.~6--97. If $\ang{B\VIpr}$ is congruent to $\ang{B}$, and $\ang{B\VIpr}$ and
$\ang{A}$ are as drawn (that is, the other side of $\ang{B\VIpr}$ is \emph{in
the interior} of $\ang{A}$), then we would say that $\ang{B} < \ang{A}$. In
Fig.~6--97, we have, so to speak, ``moved'' $\ang{B}$ to the position of
$\ang{B\VIpr}$ for comparison. Notice how the definition below parallels the
definition of ``greater than'' and ``less than'' for segments.

\FIGVININETYSEVEN

\begin{definition}
Let $A$ and $B$ be any angles not straight angles. Let $\ang{B\VIpr}$ be an
angle congruent to $\ang{B}$, such that one side of $\ang{B\VIpr}$ is the same
as the first side of $\ang{A}$, and the other side of $\ang{B\VIpr}$ and the
second side of $\ang{A}$ both lie on the same side of the first side of
$\ang{A}$.
\begin{enumerate}[label=(\alph*), leftmargin=2.2em, itemsep=2pt, topsep=3pt]
\item If the second side of $\ang{B\VIpr}$ is interior to $\ang{A}$, we say
  that $\ang{B} < \ang{A}$. (See Fig.~6--98.)
\item If the second side of $\ang{B\VIpr}$ is exterior to $\ang{A}$, we say
  that $\ang{B} > \ang{A}$. (See Fig.~6--99.)
\item If $\ang{A}$ is a straight angle, and $\ang{B}$ is not, we say that
  $\ang{A} > \ang{B}$ and that $\ang{B} < \ang{A}$.
\end{enumerate}
\end{definition}

\FIGVININETYEIGHTNINETYNINE

Our intuition tells us that if $\ang{B} < \ang{A}$, then $\ang{A} > \ang{B}$,
and if $\ang{B} > \ang{A}$, then $\ang{A} < \ang{B}$. These statements require
proof. Since the proofs of these statements are essentially the same, we prove
only the first.

\begin{theorem}
If $\ang{B}$ and $\ang{A}$ are two angles, such that $\ang{B} < \ang{A}$, then
$\ang{A} > \ang{B}$.
\end{theorem}

\VIstarnote{*This proof is difficult. If the teacher wishes, Theorem 6--8 may
be treated as a postulate.}

\stex\emph{Proof.} Figure 6--100 shows that $\ang{B} < \ang{A}$. To show that
$\ang{A} > \ang{B}$, we must show that the second side of $\ang{A\VIpr}$ is in
the exterior of $\ang{B}$, as in Fig.~6--101.

\FIGVIONEHUNDREDONEHUNDREDONE

\FIGVIONEHUNDREDTWO

Choose points $D$ and $E$ on the sides of $\ang{A}$ and points $D\VIpr$ and
$E\VIpr$ on the sides of $\ang{A\VIpr}$ such that $AD \cong BD\VIpr$ and $AE
\cong BE\VIpr$. (See Fig.~6--102.) $DE$ intersects the second side of
$\ang{B\VIpr}$ in some point $F$ between $D$ and $E$ because the second side of
$\ang{B\VIpr}$ is in the interior of $\ang{A}$. $\tri{ADE} \cong \tri{BD\VIpr
E\VIpr}$ by the S.A.S. theorem, and hence $DE \cong D\VIpr E\VIpr$. Choose
$F\VIpr$ between $D\VIpr$ and $E\VIpr$, such that $DF \cong D\VIpr F\VIpr$.
This is possible because $D\VIpr E\VIpr \cong DE > DF$. The ray $BF\VIpr$ is in
the interior of $\ang{A\VIpr}$. (Why?) $\ang{D} \cong \ang{D\VIpr}$, and hence
$\tri{FDA} \cong \tri{F\VIpr D\VIpr B}$. $\ang{DAF} \cong \ang{D\VIpr
BF\VIpr}$, and $\ang{DAF}$ was given congruent to $\ang{B}$. By the three-part
angle postulate, $\ang{B} \cong \ang{D\VIpr BF\VIpr}$; it follows that the ray
$BF\VIpr$ is the same as the second side of $\ang{B}$. Therefore, the second
side of $\ang{B}$ is interior to $\ang{A\VIpr}$ or, what is the same, the
second side of $\ang{A\VIpr}$ is in the exterior of $\ang{B}$. This completes
the proof that $\ang{A} > \ang{B}$.

\begin{remark}
Theorem 6--8 tells us that we get the same result regardless of whether we
``move'' $B$ to compare with $A$, or ``move'' $A$ to compare with $B$.
\end{remark}

\begin{theorem}
For any two angles $\ang{A}$ and $\ang{B}$, either $\ang{A} < \ang{B}$, or
$\ang{A} \cong \ang{B}$, or $\ang{A} > \ang{B}$; and exactly one of these is
true.
\end{theorem}

\emph{Proof.} If $\ang{A}$ is a straight angle, either $\ang{A} \cong \ang{B}$
and $\ang{B}$ is also a straight angle or, by Definition 6--5, $\ang{A} >
\ang{B}$. Similarly, if $\ang{B}$ is a straight angle, either $\ang{B} \cong
\ang{A}$ or $\ang{B} > \ang{A}$. If neither $\ang{A}$ nor $\ang{B}$ is
straight, then the second side of $\ang{B\VIpr}$ coincides with the second side
of $\ang{A}$, and $\ang{B} \cong \ang{A}$; or the second side of $\ang{B\VIpr}$
is interior to $\ang{A}$, and $\ang{A} > \ang{B}$; or the second side of
$\ang{B\VIpr}$ is exterior to $\ang{A}$, and $\ang{A} < \ang{B}$. Since only
one of these possibilities can hold, the proof is complete.

Some simple properties of inequalities are collected in the following theorem.

\begin{theorem}
\begin{enumerate}[label=(\alph*), leftmargin=2.2em, itemsep=1pt, topsep=3pt]
\item If $\ang{A} \cong \ang{B}$, and $\ang{B} < \ang{C}$, then $\ang{A} < \ang{C}$.
\item If $\ang{A} < \ang{B}$, and $\ang{B} \cong \ang{C}$, then $\ang{A} < \ang{C}$.
\item If $\ang{A} \cong \ang{B}$, and $\ang{B} > \ang{C}$, then $\ang{A} > \ang{C}$.
\item If $\ang{A} < \ang{B}$, and $\ang{B} < \ang{C}$, then $\ang{A} < \ang{C}$.
\item If $\ang{A} \geqq \ang{B}$, and $\ang{B} \geqq \ang{C}$, then
  $\ang{A} \geqq \ang{C}$.*
\end{enumerate}
\end{theorem}

\VIstarnote{*The symbol $\geqq$ means ``greater than or congruent to,'' as in
Chapter 4.}

\emph{Proof of} (a). Since $\ang{B} < \ang{C}$, on ``moving'' $\ang{B}$ to
compare with $\ang{C}$, the second side of $\ang{B\VIpr}$ is interior to
$\ang{C}$. Since $\ang{A} \cong \ang{B}$, the second sides of $\ang{A\VIpr}$
and $\ang{B\VIpr}$ would be the same. Hence $\ang{A} < \ang{C}$. The proofs in
all the other cases are essentially the same.

\exercisegroup{6--13}
\begin{multicols}{2}
\begin{exlist}
\item Can you think of any other properties not included in Theorem 6--10?
  There are several. Write them down. Draw figures to illustrate them.
\item Draw figures to indicate that parts (b) through (e) of Theorem 6--10 are
  true. Be sure you understand how to use Definition 6--5.
\end{exlist}
\end{multicols}

\begin{definition}
An angle less than a right angle is called an \emph{acute} angle. An angle
greater than a right angle is called an \emph{obtuse} angle.
\end{definition}

%========================================================= 6-12
\section*{6--12\quad RIGHT ANGLES AGAIN}
\addcontentsline{toc}{section}{6--12\quad Right angles again}

We can now use Theorem 6--9 to prove that all right angles are congruent to
each other. It is interesting, from a historical point of view, to know that
Euclid \emph{postulated} this fact. The proof we give below could have been
formulated by Euclid just as well as he formulated other proofs. Not only did
Euclid ignore postulates that were needed (like the betweenness postulates),
but he also introduced postulates that were not needed.

\begin{theorem}
All right angles are congruent to each other.
\end{theorem}

\emph{Proof.} Suppose that two right angles are $\ang{A}$ and $\ang{B}$, and
that their supplements are $\ang{A\VIpr}$ and $\ang{B\VIpr}$, as in
Fig.~6--103(a) and 6--103(b). There are only three possibilities, according to
Theorem 6--9. Either $\ang{A} < \ang{B}$, or $\ang{A} \cong \ang{B}$, or
$\ang{A} > \ang{B}$, and one of these must be true. We will show that $\ang{A}
< \ang{B}$ and $\ang{A} > \ang{B}$ are false. Suppose that $\ang{A} < \ang{B}$,
and that $\ang{B\VIppr}$ is congruent to $\ang{B}$ and located as in
Fig.~6--103(c). Let $B\VIpppr$ be the supplement of $\ang{B\VIppr}$.

\FIGVIONEHUNDREDTHREE

\begin{steps}
\step{$\ang{B\VIppr} \cong \ang{B}$.}{By definition of $\ang{B\VIppr}$.}
\step{$\ang{B\VIpppr}$ is the supplement of $\ang{B\VIppr}$, and
      $\ang{B\VIpr}$ is the supplement of $\ang{B}$.}{By definition of supplement.}
\step{$\ang{B\VIpr} \cong \ang{B\VIpppr}$.}{If angles are congruent, their
      supplements are congruent.}
\step{$\ang{B} \cong \ang{B\VIpr}$.}{$\ang{B}$ is a right angle.}
\step{$\ang{B} \cong \ang{B\VIpppr}$.}{Angle congruence is transitive.}
\step{$\ang{A} < \ang{B}$.}{Hypothesis.}
\step{$\ang{A} < \ang{B\VIpppr}$.}{$\ang{A} < \ang{B} \cong
      \ang{B\VIpppr}$.}
\step{$\ang{B\VIpppr} < \ang{A\VIpr}$.}{The ``second'' side of
      $\ang{B\VIpppr}$ is in the interior of $\ang{A\VIpr}$.}
\step{$\ang{B\VIpppr} < \ang{A}$.}{$\ang{B\VIpppr} < \ang{A\VIpr}
      \cong \ang{A}$.}
\step{$\ang{A} < \ang{A}$.}{Statements 7 and 9.}
\end{steps}

Since $\ang{A} \cong \ang{A}$, we know by Theorem 6--9 that it is impossible
for $\ang{A}$ to be less than itself. Therefore, it is impossible that $\ang{A}
< \ang{B}$. In the same way, it can be shown that it is impossible that
$\ang{B} < \ang{A}$. Therefore, we must have $\ang{A} \cong \ang{B}$. This
completes the proof.

\begin{definition}
Two lines, $l$ and $m$, are said to be \emph{perpendicular} to each other if
they intersect and the angles formed are right angles, as in Fig.~6--104(a).
To indicate that lines are perpendicular we write ``$l \perp m$.''
\end{definition}

\FIGVIONEHUNDREDFOUR

The line $l$ is called the \emph{perpendicular bisector} of the segment $AB$ if
$l \perp AB$ and $l$ passes through the midpoint of $AB$, as in
Fig.~6--104(b). This is written ``$l \perp$ bis $AB$,'' and read ``$l$ is the
perpendicular bisector of $AB$.''

\exercisegroup{6--14}
\begin{multicols}{2}
\begin{exlist}
\item Prove the theorem that if the two legs of one right triangle are
  congruent, respectively, to the two legs of another right triangle, then the
  triangles are congruent.
\item In the figure, angles $ADC$ and $BCD$ are right, $AD \cong BC$, and $AB
  \cong CD$. Prove that $\tri{ADC} \cong \tri{BCD}$ and $\tri{ABD} \cong
  \tri{BAC}$. (Fig.~6--105)
\item If, in the figure, $AD \cong DC$, and $BD \perp AC$, prove that
  $\ang{BAD} \cong \ang{BCD}$. (Fig.~6--106)
\item If, in the figure, $AD \cong BC$, angles $ADC$ and $BCD$ are right, and
  $DE \cong EC$, prove that $AE \cong BE$. (Fig.~6--107)
\item \VIgp{$BD \perp AC$,\\ $\ang{ABD} \cong \ang{DBC}$.}{$\ang{A} \cong
  \ang{C}$.\\ \hbox{}\hfill(Fig.~6--108)}
\item Using Fig.~6--108, prove that $\ang{A} \cong \ang{C}$ if $BD \perp$ bis
  $AC$.
\item In the figure, $AB \perp BD$, $DE \perp BD$, and $BC \cong CD$. Prove
  that $AB \cong DE$. (Fig.~6--109)
\item \VIgp{$AB \perp DC$,\\ $BE \cong BC$,\\ $\ang{DEB} \cong \ang{C}$.}%
  {$DE \cong AC$.\\ \hbox{}\hfill(Fig.~6--110)}
\item Prove that the bisector of the vertex angle of an isosceles triangle is
  the perpendicular bisector of the base.
\item Prove that the median from the vertex of an isosceles triangle is
  perpendicular to the base.
\item \VIgp{$AC \cong AD$,\\ $AC \perp CB$,\\ $AD \perp DB$.}{$\tri{ACB} \cong
  \tri{ADB}$.\\ \hbox{}\hfill(Fig.~6--111)}
\end{exlist}
\end{multicols}

\FIGVIONEHUNDREDFIVEONEHUNDREDSIX

\FIGVIONEHUNDREDSEVENONEHUNDREDEIGHT

\FIGVIONEHUNDREDNINEONEHUNDREDTEN

\FIGVIONEHUNDREDELEVEN

%========================================================= Review
\section*{REVIEW OF CHAPTER 6}
\addcontentsline{toc}{section}{Review of Chapter 6}

In this chapter there are many important ideas and geometric facts. It is most
important that you review the topics listed below. Each topic contains many
ideas, so study each one separately. Then study the whole list to be sure that
you understand how the topics are related.

1. First in importance are the four postulates of Group V. It is not necessary
to memorize the postulates word for word, but you must understand the meaning
of each one so that you will know when you can use it.

2. We also gave several important definitions. Be sure that you know
definitions for the following: right angle, congruent triangles, $>$ and $<$
for angles, acute angle, obtuse angle.

3. Below we list the most important theorems proved in this chapter. You should
be able to prove many of them. In reviewing the proofs of these theorems, first
see whether you can formulate a proof without looking at the text.

\begin{quote}
If two angles are congruent, their supplements are congruent.\\
Vertical angles are congruent.\\
In an isosceles triangle, the base angles are congruent.\\
The first congruence theorem (S.A.S.).\\
For any two angles $\ang{A}$ and $\ang{B}$, exactly one of the following is
true: $\ang{A} > \ang{B}$, $\ang{A} \cong \ang{B}$, $\ang{A} < \ang{B}$.\\
If $\ang{A} \cong \ang{B}$, and $\ang{B} < \ang{C}$, then $\ang{A} < \ang{C}$,
etc.\\
All right angles are congruent.\\
The second congruence theorem (A.S.A.).\\
The third congruence theorem (S.S.S.).
\end{quote}

\subsection*{Review Exercises}
\begin{multicols}{2}
\begin{exlist}
\item Draw an $\ang{A}$. Draw a ray $r$ from a point $O$. Draw the two angles
  congruent to $A$ that have the ray $r$ as one side. See Postulate V--1.
\item If $\ang{A} \cong \ang{D}$, and $\ang{D} \cong \ang{E}$, and $\ang{E}
  \cong \ang{F}$, prove that $\ang{A} \cong \ang{F}$. See Postulate V--2.
\item If in Fig.~6--112, $\ang{A_1} \cong \ang{R}$, and $\ang{A_2} \cong
  \ang{S}$, prove that $\ang{B} \cong \ang{T}$. See Postulate V--3.
\item If in Fig.~6--113, $\ang{A} \cong \ang{B}$, prove that $\ang{C} \cong
  \ang{D}$. What is the converse of this? Is the converse true? What is the
  contrapositive of this? Is it true?
\item If in Fig.~6--114, $\ang{AOB} \cong \ang{BOC}$, what is $\ang{AOB}$
  called?
\item If (Fig.~6--115) $FE \cong ST$, $EG \cong RT$, and $\ang{E} \cong
  \ang{T}$, to what angle is $\ang{F}$ congruent?
\item If (Fig.~6--116) $\sg{RS} = \sg{TS}$, prove that $\ang{R} \cong \ang{T}$.
\item Using Fig.~6--115, prove that the triangles are congruent.
\item Figure 6--117 shows a way to measure the width $\sg{PQ}$ of a river. From
  any point $R$ not on the line $PQ$, $RU$ is sighted so that $\ang{URP} \cong
  \ang{PRQ}$. From $P$ a line $PT$ is sighted so that $\ang{TPR} \cong
  \ang{QPR}$. The point $S$, where $PT$ and $RU$ intersect, is determined, and
  $PS$ is measured. Why is this the distance from $P$ to $Q$?
\item The first congruence theorem says that if two sides and the included
  angle of one triangle are congruent to two sides and the included angle of
  another triangle, then the triangles are congruent. What is the converse of
  this theorem? Is it true?
\item If a triangle is equilateral, then it is isosceles. What is the converse
  of this statement? Is it true? What is the contrapositive? Is it true?
\item If in Fig.~6--118, $\tri{ABC}$ and $\tri{ADC}$ are isosceles, with
  vertices $B$ and $D$, prove that $BD$ bisects $\ang{ABC}$.
\item If (Fig.~6--119) $\tri{ABC}$ is isosceles, and $\ang{ABD} \cong
  \ang{EBC}$, prove that $\sg{AD} = \sg{EC}$.
\item If (Fig.~6--120) $\sg{AB} = \sg{BC}$, $\sg{AD} = \sg{DC}$, and
  $\ang{ADE} \cong \ang{CDF}$, prove that $BE \cong BF$.
\item When we build a house, the sub-flooring and the sheathing boards are
  usually laid diagonally. Why?
\item If (Fig.~6--121) $QS \cong SR$, and $PS \perp QR$, prove that $PQ \cong
  PR$.
\item If the point $O$ bisects the segments $AB$ and $CD$, prove that $\ang{A}
  \cong \ang{B}$. See Fig.~6--122.
\item Show that if two isosceles triangles have congruent vertex angles, and
  their congruent sides have the same length, then the triangles are congruent.
\item If in Fig.~6--123, $AB \cong AD$, $BC \cong DE$, and $AC \cong AE$, to
  what angle is $\ang{B}$ congruent? Why?
\item If in Fig.~6--124 $\tri{ABC}$ is equilateral, and $D$, $E$, $F$ bisect
  the sides, prove that $\tri{EFD}$ is equilateral.
\item If (Fig.~6--125) $\ang{CAD} \cong \ang{BDA}$, and $\ang{CDA} \cong
  \ang{BAD}$, prove that $\sg{AB} = \sg{CD}$.
\item Prove that if two points are equidistant from the ends of a segment, they
  determine the perpendicular bisector of the segment.
\item Prove that if the bisector of an angle of a triangle is perpendicular to
  the opposite side, then it bisects that side, and the triangle is isosceles.
\item Prove that any point on the perpendicular bisector of a segment is
  equidistant from the ends of the segment.
\item \VIgp{$\ang{1} \cong \ang{2}$,\\ $AC \cong BC$,\\ $IJ \cong IK$.}%
  {$GJ \cong HK$,\\ $GF \cong FH$.\\ \hbox{}\hfill(Fig.~6--126)}
\end{exlist}
\end{multicols}

\FIGVIONEHUNDREDTWELVEONEHUNDREDTHIRTEEN

\FIGVIONEHUNDREDFOURTEENONEHUNDREDFIFTEEN

\FIGVIONEHUNDREDSIXTEENONEHUNDREDSEVENTEEN

\FIGVIONEHUNDREDEIGHTEENONEHUNDREDNINETEEN

\FIGVIONEHUNDREDTWENTYONEHUNDREDTWENTYONE

\FIGVIONEHUNDREDTWENTYTWOONEHUNDREDTWENTYTHREE

\FIGVIONEHUNDREDTWENTYFOUR

\FIGVIONEHUNDREDTWENTYFIVEONEHUNDREDTWENTYSIX

\subsection*{Test 1}
\begin{center}\bfseries True-False\end{center}
\begin{exlist}
\item If a theorem is true, its contrapositive is true.
\item The supplement of a $40^\circ$ angle is a $50^\circ$ angle.
\item If $\ang{A} < \ang{B}$, and $\ang{B} < \ang{C}$, and $C$ is a $50^\circ$
  angle, then $\ang{A}$ is less than a $49^\circ$ angle.
\item If $ABC$ ($B$ is between $A$ and $C$), $DEF$, $AC \cong DF$, and $AB >
  DE$, then $BC > EF$.
\item If $\tri{ABC}$ is isosceles, with vertex $A$, and $D$, $E$ are the
  midpoints of $AB$, $AC$ respectively, then $BE \cong CD$.
\item If in triangles $ABC$ and $A\VIpr B\VIpr C\VIpr$, $\ang{A} \cong
  \ang{A\VIpr}$, $\ang{B} \cong \ang{B\VIpr}$, $\ang{C} \cong \ang{C\VIpr}$,
  and $AB \cong A\VIpr C\VIpr$, then $\tri{ABC} \cong \tri{A\VIpr B\VIpr
  C\VIpr}$.
\item If $\ang{ABC}$ is $40^\circ$, and $\ang{DBC}$ is $50^\circ$, then
  $\ang{DBA}$ is a right angle.
\item If $\ang{ABC}$ is $60^\circ$, and $\ang{DBC}$ is $120^\circ$, then either
  $\ang{DBA} \cong \ang{ABC}$, or $\ang{DBA}$ is a straight angle.
\item If segments $AB$ and $CD$ intersect at $O$, then $\ang{AOB} \cong
  \ang{BOD}$.
\item For all statements $\al$, $\be$, $\ga$, such that $\al \to \be$ is true
  and $\be \to \ga$ is false, the implication $\al \to \ga$ is false.
\end{exlist}

\subsection*{Test 2}
\begin{exlist}
\item Lines $l$ and $m$ intersect so that a pair of adjacent angles differ by
  $18^\circ$. Determine the number of degrees in each angle.
\item If $\sg{AB} = 24$, and $ACB$, and $\sg{AC} = 3 \cdot \sg{CB}$, compute
  $\sg{CB}$.
\item If $\sg{AB} = 24$, $ABC$, and $\sg{AC} = 3 \cdot \sg{CB}$, compute
  $\sg{CB}$.
\item Line segments $AB$ and $CD$ are congruent. The length of $AB$ as measured
  by $EF$ is 12. The length of $CD$ as measured by $GH$ is 36. What is the
  length of $GH$ as measured by $EF$?
\item The length of $AB$ as measured by $CD$ is $\tfrac{3}{4}$. What is the
  length of $CD$ as measured by $AB$?
\item Let $AB$ be a unit segment. Designate the midpoint of $AB$ by $A_1$, the
  midpoint of $AA_1$ by $A_2$, the midpoint of $AA_2$ by $A_3$, etc. What is
  the length of $A_5A_3$?
\end{exlist}

\subsection*{Algebra Review}
\begin{exlist}
\item The perimeter of an isosceles triangle is 75 ft. One side is 15 ft longer
  than the second side. Prove that the shortest side of the triangle is either
  15 ft long or 20 ft long.
\item A triangle has two equal angles. The length of the included side is 7 in.
  The perimeter of the triangle is 43 in. Determine the lengths of the other
  two sides of the triangle.
\item $\tri{ABC}$ is equilateral. Points $X$, $Y$, and $Z$ are the midpoints of
  sides $AB$, $BC$, and $CA$, respectively. The perimeter of $\tri{XYZ}$ is
  30 in. Determine the length of each side of $\tri{ABC}$.
\item In Fig.~6--127, $\sg{AE} = \sg{AD}$, and $\sg{AB} = \sg{AC}$. $\sg{BE} =
  18$ in., and $\sg{DO} = \tfrac{1}{2}\sg{OC}$. Determine the length of segment
  $OC$.
\item In $\tri{ABC}$, the bisector of vertex angle $A$ is perpendicular to the
  base $BC$ at point $D$. $\sg{BD} = \sg{CD}$. $\sg{AB} - \sg{BD} = 8$ in., and
  the perimeter of $\tri{ABC}$ is 40 in. Determine the length of each side of
  $\tri{ABC}$.
\item Prove that if $(x + 2)(x - 3) = 0$, and $x$ is positive, then $x = 3$.
\item Prove that if $(x + 2)(x - 3) = 0$, and $x$ is negative, then $x = -2$.
\item A businessman bought several suits of clothes. He paid \$40 each for some
  of the suits and \$50 each for the rest. If he paid a total of \$1500 for
  these suits, prove that he could not possibly have bought 22 of the \$40
  suits.
\item[\stex 9.] Prove that the following problem has no solution: Find two
  numbers whose sum is 32, such that if the larger of the numbers is divided by
  6, and the smaller is divided by 5, then the sum of the quotients is 6.
\item[\stex 10.] Describe all possible pairs of numbers $a$ and $b$ such that
  $(a - 2)/(2b - 9) = a - 2$.
\end{exlist}

\FIGVIONEHUNDREDTWENTYSEVEN
\clearpage
% ==================== ch07 ====================
\setcounter{theorem}{0}\setcounter{definition}{0}\setcounter{construction}{0}\setcounter{exgroup}{0}\setcounter{footnote}{0}
\setcounter{chapter}{7}

\chapter*{\raggedright\Huge USE OF THE\\CONGRUENCE THEOREMS}
\addcontentsline{toc}{chapter}{7.\ Use of the Congruence Theorems}
\markboth{USE OF THE CONGRUENCE THEOREMS}{}
\vspace{-1em}\noindent\rule{\linewidth}{0.6pt}\vspace{1em}
\figurekey

%========================================================= 7-1
\section*{7--1\quad INTRODUCTION}
\addcontentsline{toc}{section}{7--1\quad Introduction}

In this chapter we do two things. In the first part we prove several
inequalities involving sides and angles of a triangle, ending up with the
\emph{triangle inequality}. This enables us to prove that a line segment is
the shortest path between two points. In the second part we prove that some of
the constructions we made with ruler and compass in Chapter 1 are correct.

%========================================================= 7-2
\section*{7--2\quad EXTERIOR ANGLES}
\addcontentsline{toc}{section}{7--2\quad Exterior angles}

In this section we prove a theorem about triangles which Euclid found to be
very useful. We need a definition.

\begin{definition}
If an angle is adjacent and supplementary to an angle of a triangle, it is
called an \emph{exterior angle} of the triangle. The other two angles of the
triangle are called the \emph{opposite interior angles}. For example,
$\ang{DBC}$ in Fig.~7--1 is an exterior angle and $\ang{A}$ and $\ang{C}$ are
interior angles and opposite to $\ang{DBC}$.
\end{definition}

\FIGVIIONE

\exercisegroup{7--1}
\begin{multicols}{2}
\begin{exlist}
\item In $\tri{ABC}$, how many exterior angles are there having vertex $B$?
  How are these exterior angles related to $\ang{B}$?
\item How many exterior angles does a triangle have?
\item Do the words \emph{opposite}, \emph{interior}, and \emph{exterior}, used
  in Definition 7--1, refer to properties or to relations?
\end{exlist}
\end{multicols}

\begin{theorem}
An exterior angle of a triangle is greater than either opposite interior angle.
\end{theorem}

We give two separate proofs of this theorem. The first proof may be found in
Hilbert's work, the second in Euclid's \emph{Elements}.

\noindent\emph{First Proof.}\quad
We know that, for any two angles, the first angle is either $>$, or $<$, or
$\cong$ the second. Therefore, we need only show that the exterior angle
cannot be $\cong$ or $<$ an opposite interior angle, and it must follow that
the exterior angle is $>$ an opposite interior angle. Suppose that the triangle
is labeled as in Fig.~7--2.

\FIGVIITWOTHREE

We first show that the exterior angle cannot be congruent to $\ang{C}$. Choose
$D$, as in the figure, such that $BD \cong AC$, and suppose that
$\ang{DBC} \cong \ang{C}$. Then $\tri{ACB} \cong \tri{DBC}$ because of the
S.A.S. theorem. Hence $\ang{CBA} \cong \ang{BCD}$. Therefore, because of the
addition of angles postulate, $\ang{ACD} \cong \ang{ABD}$. But $\ang{ABD}$ is
a straight angle, and so $\ang{ACD}$ is a straight angle. This is impossible.
(Why?) Hence $\ang{DBC}$ is \emph{not} congruent to $\ang{C}$.

We now show that it is impossible that $\ang{DBC} < \ang{ACB}$. Assume that
$\ang{DBC} < \ang{ACB}$. With $C$ as a vertex, let $\ang{ECB}$ be the angle, as
drawn in Fig.~7--3, that is congruent to $\ang{DBC}$. (Why is there such an
angle?) Then the ray $CE$ is in the interior of $\ang{ACB}$. (Why?) Hence the
segment $AB$ contains a point $F$ in the ray $CE$. (This is obvious from the
figure, and it may be proved with what we know about betweenness and interior
of angles.) Now consider $\tri{FCB}$. $\ang{DBC}$ is an exterior angle of this
triangle which is congruent to the opposite interior angle $FCB$. This has
already been shown to be impossible. Hence we cannot have
$\ang{DBC} < \ang{ACB}$. The only possibility then is
$\ang{DBC} > \ang{ACB}$. This completes the proof.

\noindent\emph{Second Proof.}\quad
Using the same labeling for the triangle, let $G$ be the point on $BC$, such
that $BG \cong GC$. Let $H$ be the point on the line $AG$ that is on the other
side of $G$ from $A$, such that $AG \cong GH$. Then
$\tri{GCA} \cong \tri{GBH}$. (Why?) Hence $\ang{C} \cong \ang{GBH}$. The ray
$BH$ is interior to $\ang{DBC}$. (This is obvious from Fig.~7--4, but it needs
proof. It is interesting to note that Euclid neglects to prove this. See
Exercise 3 below.) Therefore, $\ang{DBC} > \ang{HBC}$, and so
$\ang{DBC} > \ang{C}$. (Why?)

\FIGVIIFOUR

\begin{remark}
In Chapter 8, when we study parallel lines, we shall be able to prove that an
exterior angle is congruent to the ``sum'' of the opposite interior angles, and
from this fact, Theorem 7--1 follows easily. What is interesting is that
Theorem 7--1 may be proved without any knowledge of properties of parallel
lines.
\end{remark}

\exercisegroup{7--2}
\begin{multicols}{2}
\begin{exlist}
\item What kind of angles are the exterior angles of a right triangle?
\item Use Theorem 7--1 to prove that the angles opposite the legs of a right
  triangle are acute angles.
\item[\stex 3.]\stepcounter{exlisti} Prove that the ray $BH$ is interior to $\ang{DBC}$.
  [\emph{Hint:} first show that $H$ is on the same side of line $BC$ as $D$.
  Since $G$ is on the same side of $BD$ as $C$, if $H$ were on the opposite
  side of $BD$ from $C$, then the segment $GH$ would contain a point of the
  line $BD$. This is absurd.]
\item Prove the following corollary to Theorem 7--1. There is at most one line
  perpendicular to a given line $l$ from a given point $P$. [\emph{Hint:}
  Suppose that there were two perpendiculars, as in Fig.~7--5.]
\end{exlist}
\end{multicols}

\FIGVIIFIVE

%========================================================= 7-3
\section*{7--3\quad INEQUALITIES IN TRIANGLES}
\addcontentsline{toc}{section}{7--3\quad Inequalities in triangles}

In this section we prove some theorems about the relative sizes of sides and
angles in a triangle. We arrive at a proof of a theorem that, in effect, says
that a line segment is the shortest path between two points.

\begin{theorem}
If, in $\tri{ABC}$, $AB > AC$, then $\ang{C} > \ang{B}$, and conversely, if
$\ang{C} > \ang{B}$, then $AB > AC$.
\end{theorem}

\begin{remark}
The meaning of this theorem can be stated briefly as ``the greater side is
opposite the greater angle.''
\end{remark}

\FIGVIISIX

\noindent\emph{Proof.}\quad
To prove the first part of the theorem, suppose that $AB > AC$, as shown in
Fig.~7--6. Choose $C'$ on $AB$, such that $AC' \cong AC$. Then by
Theorem 7--1, $\ang{AC'C} > \ang{B}$. But $\ang{AC'C} < \ang{C}$. Hence
$\ang{C} > \ang{B}$.

The converse follows readily from the fact that any two segments, are either
congruent to each other or one is greater than the other. The argument is given
below.

Given that $\ang{C} > \ang{B}$, we wish to show that $AB > AC$. Exactly one of
the relations $AB < AC$, $AB \cong AC$, or $AB > AC$ is true. If
$AB \cong AC$, then we would have $\ang{C} \cong \ang{B}$, which contradicts
the given fact that $\ang{C} > \ang{B}$. If $AB < AC$, then by the first part
of the theorem, we would have $\ang{C} < \ang{B}$, which also contradicts the
given fact that $\ang{C} > \ang{B}$. Hence the only possibility is that
$AB > AC$.

\exercisegroup{7--3}
\begin{multicols}{2}
\begin{exlist}
\item Which angle in Fig.~7--7 is the smallest? Which is the largest? The
  numbers are lengths of the sides.
\item If in $\tri{EFG}$, $\ang{E}$ has a measure of $61^\circ$, and $\ang{F}$ a
  measure of $54^\circ$, which side is the longer, $FG$ or $EG$? Can you say
  anything yet about the length of $EF$?
\item In $\tri{PQR}$, $S$ is a point of $PR$, such that $QS$ bisects
  $\ang{PQR}$ (see Fig.~7--8). Prove that $PQ > PS$. [\emph{Hint:} Prove first
  that $\ang{PSQ} > \ang{PQS}$.]
\end{exlist}
\end{multicols}

\FIGVIISEVENEIGHT

We now come to a very important and familiar inequality for triangles, called
the \emph{triangle inequality}.

\begin{theorem}
In any $\tri{ABC}$, $\sg{AB} + \sg{BC} > \sg{AC}$.
\end{theorem}

\FIGVIININETEN

\noindent\emph{Proof.}\quad
If $AC$ is not the longest of the sides, the inequality is obvious. Therefore
we may confine our attention to the case in which $\sg{AC} > \sg{AB}$ and
$\sg{AC} > \sg{BC}$, as in Fig.~7--9.

\begin{steps}
\step{There is a point $B'$ on $AC$, such that $AB \cong AB'$.}{Why?}
\step{$\ang{1} \cong \ang{2}$.}{Why?}
\step{$\ang{3} > \ang{1}$.}{Why?}
\step{$\ang{2} > \ang{4}$.}{Why?}
\step{$\ang{3} > \ang{4}$.}{This follows from Statements 2, 3, and 4. Why?}
\step{$\sg{B'C} < \sg{BC}$.}{Why?}
\step{$\sg{AB'} + \sg{B'C} < \sg{AB} + \sg{BC}$.}{Statements 1 and 6 and
      properties of numbers.}
\step{$\sg{AC} < \sg{AB} + \sg{BC}$.}{$\sg{AC} = \sg{AB'} + \sg{B'C}$.}
\end{steps}

\begin{remark}
This theorem is sometimes loosely stated as ``a straight line is the shortest
distance between two points.'' In fact, this last statement is sometimes used
as a postulate, although incorrectly, since one must first have a notion of
distance or of length of a segment.
\end{remark}

Having proved the triangle inequality, we now see that the lines we have been
talking about all along are, after all, ``straight'' in the precise sense of
the triangle inequality. Perhaps this can be made clearer by defining
\emph{polygonal path} and \emph{length of a polygonal path}.

\begin{definition}
Let $P_1$, $P_2,\ldots, P_n$ be any set of $n$ points (Fig.~7--10). The set of
points on all the segments $P_1P_2$, $P_2P_3,\ldots, P_{n-1}P_n$ is called a
\emph{polygonal path}. The path is said to join, or connect, the
\emph{endpoints} $P_1$ and $P_n$. The points $P_1,\ldots, P_n$ are called the
\emph{vertices}. In particular, if $P_1 = P_n$, the polygonal path is called a
\emph{polygon} and designated as polygon $P_1P_2\ldots P_{n-1}$.
\end{definition}

\begin{definition}
The \emph{length of the polygonal path} $P_1\ldots P_n$ is defined to be
$\sg{P_1P_2} + \sg{P_2P_3} + \cdots + \sg{P_{n-1}P_n}$.
\end{definition}

\begin{corollary}
The length of a polygonal path is greater than, or equal to, the length of the
segment joining its ends.
\end{corollary}

\exercisegroup{7--4}
\begin{multicols}{2}
\begin{exlist}
\item Two sides of a triangle have lengths of 5 and 7 in. What can you say
  about the third side?
\item Can the numbers 6, 8, and 15 be the lengths of the three sides of a
  triangle?
\item Prove that the difference between the lengths of two sides of a triangle
  is less than the length of the third side.
\item[\stex 4.]\stepcounter{exlisti} Use the triangle inequality and Fig.~7--11 to show that if
  $\ang{C} > \ang{B}$, then $AB > AC$.
\item[\stex 5.]\stepcounter{exlisti} In Fig.~7--12, $AC \cong A'C'$, $BC \cong B'C'$, and
  $\ang{ACB} > \ang{A'C'B'}$. Prove that $\sg{AB} > \sg{A'B'}$.
  [\emph{Hint:} Show how to find a point $Q$ on $AB$, such that
  $QB \cong QB'$, and use $\sg{AB} = \sg{AQ} + \sg{QB} = \sg{AQ} + \sg{QB'}$.]
\item[\stex 6.]\stepcounter{exlisti} Use the result of Exercise 5 to prove that if, in two
  triangles, $\tri{ABC}$ and $\tri{A'B'C'}$, $AC \cong A'C'$,
  $BC \cong B'C'$, and $AB > A'B'$, then $\ang{C} > \ang{C'}$. Show that the
  statements $\ang{C} \cong \ang{C'}$ and $\ang{C} < \ang{C'}$ are false.
\item[\stex 7.]\stepcounter{exlisti} If $D$ in Fig.~7--13 is a point inside $\tri{ABC}$, show that
  $\sg{AD} + \sg{DC} < \sg{AB} + \sg{BC}$. [\emph{Hint:} Consider first
  $\tri{ABE}$ and $\tri{DEC}$.] Show also that $\ang{ADC} > \ang{ABC}$.
\item By using Theorem 7--2, prove that an equiangular triangle is equilateral.
\item If, in Fig.~7--14, $l \perp m$, and $l' \perp m$, prove that $l$ and $l'$
  do not intersect. [\emph{Hint:} Suppose they do.]
\item Show that a polygonal path connecting points $A$ and $B$ could have four
  vertices and still have the same length as $AB$.
\end{exlist}
\end{multicols}

\FIGVIIELEVEN
\FIGVIITWELVE
\FIGVIITHIRTEENFOURTEEN

%========================================================= 7-4
\section*{7--4\quad CIRCLES}
\addcontentsline{toc}{section}{7--4\quad Circles}

There are two ways that we could define a circle. One is to use the idea of
congruent segments, and the other is to use the idea of distance. Because two
segments are congruent if and only if they have the same length, these two ways
are equivalent.

\begin{definition}
A circle is determined by a point $O$, called the \emph{center}, and a positive
number $r$. The \emph{circle} is the set of all points $A$ of the plane, such
that $\sg{OA} = r$. If $A$ is a point of the circle, then the segment $OA$ is
called a \emph{radius} (plural, \emph{radii}). If $A$ and $B$ are on the
circle, and the segment $AB$ contains $O$, then $AB$ is called a
\emph{diameter}.
\end{definition}

\begin{VIIremarkn}{1}
You will note that we have not said that a circle is a ``curve,'' because
nowhere in our geometry have we defined what ``curve'' might mean. The circle
is a set of points, and in the definition, we have told only how to decide
whether any point $A$ is, or is not, a point of the circle.
\end{VIIremarkn}

\begin{VIIremarkn}{2}
According to the definition, a radius is a segment, that is, a set of points.
It is common usage to also call the real number $r$ the radius. We will also
use this language, so that the circle can be described as the ``circle with
center $O$ and radius $r$.'' Similarly, if $AB$ is a diameter, the length
$\sg{AB} = 2r$ will be called the diameter. No confusion will arise, because it
will always be clear whether we are talking about a set of points or a number.
\end{VIIremarkn}

\begin{VIIremarkn}{3}
We have not mentioned the inside or outside of the circle as yet, because these
sets have not yet been defined.
\end{VIIremarkn}

\begin{definition}
If a circle (Fig.~7--15) has center $O$ and radius $r$, the \emph{interior} of
the circle is the set of all points $B$, such that $\sg{OB} < r$. The
\emph{exterior} of the circle is the set of all points $C$, such that
$\sg{OC} > r$.

A \emph{chord} is a segment determined by two points $E$ and $F$ on the circle;
the line through $E$ and $F$ is a \emph{secant}.

A \emph{central angle} of a circle is an angle whose vertex is at the center of
the circle. An \emph{inscribed angle} (such as $\ang{EFG}$) is an angle whose
vertex is on the circle, and having points of the circle on each of the two
sides of the angle.
\end{definition}

\FIGVIIFIFTEEN

\exercisegroup{7--5}
\begin{multicols}{2}
\begin{exlist}
\item A circle has center $A$ and radius 5. The points $B$, $C$, $D$, $E$, are
  such that $\sg{AB} = 4$, $\sg{AC} = 5$, $\sg{AD} = 17$, and $\sg{AE} = 5$.
  Where are these points? That is, in what sets are they?
\item If $AB$ is a diameter of a circle with center $C$, what is $\ang{ACB}$?
\item In Fig.~7--16, the circle has center $O$. What kind of angles are
  $\ang{A}$, $\ang{B}$, $\ang{C}$, $\ang{D}$, and $\ang{E}$?
\item Give a definition of circle in terms of congruence instead of distance.
\item If $A$ is a point of a circle and $O$ is the center, what can be said
  about a point $B$, such that $\sg{OB} \leqq \sg{OA}$?
\end{exlist}
\end{multicols}

\FIGVIISIXTEEN

%========================================================= 7-5
\section*{7--5\quad RULER AND COMPASS}
\addcontentsline{toc}{section}{7--5\quad Ruler and compass}

In the preceding section we defined a circle in terms of distance. When drawing
a circle on a piece of paper, we use the fact, familiar from experience, that
the distance between the two points of the compass stays fixed. In this section
we define those drawings that are called ``ruler and compass constructions.''
Do not think that these are the only constructions possible. They are only the
simplest kind. The ruler used in these constructions is just an unmarked
straightedge.

\FIGVIISEVENTEEN

\begin{definition}
To make a \emph{ruler and compass} (\RC) \emph{construction} is to draw a
picture of a set of points, using only a ruler and compass. The ruler can only
be used to draw a line through two points. The compass can only be used to
draw a circle with a given center and radius.
\end{definition}

\begin{VIIremarkn}{1}
\RC{} constructions are sets that can be illustrated with certain mechanical
devices. Circles and lines exist in our geometry independently of these
mechanical devices. Interest in \RC{} constructions originates in the
simplicity of the tools needed to make them.
\end{VIIremarkn}

\begin{VIIremarkn}{2}
To draw a line or circle in an \RC{} construction, two points are needed. We
may not slide the ruler, mark it, or guess the opening of the compass. We must
have two points already constructed in order to draw a line or circle.
\end{VIIremarkn}

\begin{VIIremarkn}{3}
The study of \RC{} constructions at this point is not really necessary for the
logical development of our geometry, but it does help us to visualize geometric
figures, and it demonstrates the usefulness of what has already been studied.
\end{VIIremarkn}

\begin{VIIremarkn}{4}
It will be clear that the knowledge of \RC{} constructions is quite useful to a
draftsman.
\end{VIIremarkn}

%========================================================= 7-6
\section*{7--6\quad HISTORICAL NOTE}
\addcontentsline{toc}{section}{7--6\quad Historical note}

The ancient Greeks were much attracted by \RC{} constructions and attempted to
solve all sorts of problems with these two tools. However, they encountered
several problems that they could not handle using only ruler and compass. This
does not mean that these problems are not solvable---only that a ruler and
compass are not enough. Three of the problems that puzzled the Greeks have had
a long and interesting history, which we cannot pursue here. They are:

\begin{enumerate}[label=(\alph*), leftmargin=2.2em, itemsep=3pt]
\item \emph{The trisection of an angle.} Given any $\ang{A}$, the problem is to
  construct angles $B$, $C$, and $D$, as in Fig.~7--18, such that
  $\ang{B} \cong \ang{C} \cong \ang{D}$.
\item \emph{The duplication of a cube.} Suppose that a cube has side of length
  $x$. The problem is to construct the side of a cube having twice the volume
  of the given cube.
\item \emph{Squaring the circle.} Given a circle, the problem is to construct a
  square with the same area as the circle (Fig.~7--19).
\end{enumerate}

\FIGVIIEIGHTEENNINETEEN

All these problems have been completely resolved in modern times, and
\emph{all are impossible of solution with ruler and compass alone}. This is a
fact beyond dispute. It does not mean that in the trisection problem, for
example, angles $B$, $C$, and $D$ do not exist. It means that one cannot draw
them using ruler and compass alone. When you think about this a moment, it may
not be so surprising. No one is surprised that it is impossible to build a
large building with only a tack hammer. A ruler and compass alone are like a
tack hammer in so far as these problems are concerned.

There are, however, many constructions possible with ruler and compass. We will
study some of them.

%========================================================= 7-7
\section*{7--7\quad COPYING ANGLES AND TRIANGLES}
\addcontentsline{toc}{section}{7--7\quad Copying angles and triangles}

$\ang{A}$ and a ray from point $B$ are given, as in Fig.~7--20. We wish to
construct an angle at $B$ congruent to $\ang{A}$. Of course, from the existence
postulate for angles, there are two such angles, as shown by the dotted lines
in the figure. When we draw the angle at $B$ congruent to $\ang{A}$, we say
that we have ``copied'' $\ang{A}$.

\FIGVIITWENTYTWENTYONE

\begin{construction}
Draw a circle of radius $r$, with center $A$. This circle will intersect the
sides of $\ang{A}$ in two points $D$ and $E$ (Fig.~7--21). (Why?) Draw a circle
of the same radius $r$ with center $B$. This circle will intersect the ray from
$B$ in a point $F$. Draw a circle, with center $F$ and radius $\sg{DE}$. On a
given side of the line $BF$, the two circles will intersect in a point $G$.
$\ang{GBF}$ is congruent to $\ang{A}$.
\end{construction}

You should recognize that there are two things to prove. Because the entire
construction is centered around finding a point in which our circles intersect,
we must first prove that such a point exists in our geometry. Of course, when
we draw the picture, it is clear that there should be such a point. The real
question, however, is whether we can take the assumptions we have made and
prove that there is such a point. After we establish the existence of this
point, we must prove further that the angles are congruent to each other. These
two proofs should be considered separately.

\noindent\emph{Proof.} (a)\quad
(The existence of a point of intersection of the circles.) By the existence
postulate for angles, there is exactly one ray $BH$, on a given side of $BF$,
such that $\ang{FBH} \cong \ang{A}$. The existence of the \emph{ray} is thus
assured. The existence of the points $D$, $E$, and $F$ follows from theorems of
measurement and congruence. There is a point $G$ on the ray $BH$, such that
$BG \cong BF \cong AE \cong AD$. Since $\ang{A} \cong \ang{B}$,
$\tri{ADE} \cong \tri{BGF}$ by the S.A.S. theorem. Therefore,
$GF \cong DE$. Hence, the point $G$ is on the circle with radius $\sg{DE}$ and
center $F$. Since $G$ is also on the circle with center $B$ and radius
$\sg{BF}$, it is therefore on both circles and is a point of intersection of
the two circles.

\noindent\emph{Proof.} (b)\quad
(The congruence of the angles.) We now consider any point $G$ on both circles.
$AD \cong AE \cong BG \cong BF$. Since both circles have the same radius,
$DE \cong GF$, and $\tri{AED} \cong \tri{BFG}$ by the S.S.S. theorem. Hence
$\ang{A} \cong \ang{B}$.

In a similar manner we can copy a triangle.

\begin{construction}
Given $\tri{ABC}$ (Fig.~7--22) and a ray from a point $A'$, use the previous
construction to copy $\ang{A}$. Then draw a circle, with center $A'$ and radius
$\sg{AB}$, which intersects the other side of $\ang{A'}$ at $B'$. Then draw a
circle, with center $A'$ and radius $\sg{AC}$, which intersects the given ray
in $C'$. Then $\tri{ABC} \cong \tri{A'B'C'}$.
\end{construction}

\noindent\emph{Proof.}\quad
The proof is left to the student.

\FIGVIITWENTYTWOTWENTYTHREE

\begin{construction}
Given $\tri{ABC}$ (Fig.~7--23) and a ray from a point $A'$, draw a circle, with
center $A'$ and radius $\sg{AC}$, which intersects the ray in $C'$. Then draw
circles, with centers at $A'$ and $C'$ and radii $\sg{AB}$ and $\sg{BC}$,
respectively. These two circles will intersect, on a given side of the ray, in
exactly one point $B'$. We will then have $\tri{ABC} \cong A'B'C'$.
\end{construction}

\noindent\emph{Proof.}\quad
The proof is left to the student. Remember, you must first prove that $B'$
exists and then prove $\tri{ABC} \cong \tri{A'B'C'}$.

%========================================================= 7-8
\section*{7--8\quad OTHER CONSTRUCTIONS}
\addcontentsline{toc}{section}{7--8\quad Other constructions}

It should now be clear that proving the validity of certain constructions
involves first the proof of the existence of certain points and lines, and
secondly the proof that the construction produces the desired result. In the
remainder of our work on constructions, we will concentrate our efforts on the
second part. The proof of existence of certain points can actually be quite
hard at times. For example, when we try to prove that the construction for the
perpendicular bisector of a segment is valid, we need a theorem, Theorem 7--4,
that is very difficult to prove at this time.

\begin{theorem}
If $AB$ is a segment, and circles are drawn, with centers at $A$ and $B$ and
radius $r > \frac{1}{2}\sg{AB}$, then these circles intersect in exactly two
points $P$ and $Q$, on opposite sides of the line through $A$ and $B$
(Fig.~7--24).
\end{theorem}

\begin{remark}
From what we know about measurement of length, we know that every line segment
$AB$ has exactly one midpoint $C$, that is, a point $C$ such that
$\sg{AC} = \sg{BC} = \frac{1}{2}\sg{AB}$. Likewise, it is easy to show that
every angle has a bisector, that is, from the vertex of the angle, there is a
ray that forms, with the sides of the given angle, two congruent angles.
\end{remark}

\FIGVIITWENTYFOURTWENTYFIVE

The constructions of this section show how to make \RC{} constructions for
midpoints of segments and bisectors of angles; but the \emph{existence} of
these midpoints and bisectors has been obtained independently of these
constructions.

\begin{construction}
Draw circles, with centers $A$ and $B$ and equal radii greater than
$\frac{1}{2}\sg{AB}$. These circles intersect at points $P$ and $Q$. The line
$PQ$ is the desired perpendicular bisector (Fig.~7--25).
\end{construction}

\noindent\emph{Proof.}
\begin{steps}
\step{$AP \cong PB \cong AQ \cong QB$.}{Construction.}
\step{$\tri{APQ} \cong \tri{BPQ}$.}{S.S.S. theorem.}
\step{$PQ$ intersects line $AB$ in $C$.}{$P$ and $Q$ are on opposite sides of
      $AB$.}
\step{$\tri{APC} \cong \tri{BPC}$.}{S.A.S. theorem.}
\step{$AC \cong CB$, and $\ang{ACP} \cong \ang{BCP}$.}{Statement 4.}
\step{$PQ$ is the $\perp$ bis $AB$.}{Statement 5.}
\end{steps}

\exercisegroup{7--6}
\begin{multicols}{2}
\begin{exlist}
\item Prove that each segment has exactly one perpendicular bisector.
\item Prove that every point on the perpendicular bisector of $AB$ is
  equidistant from $A$ and $B$.
\item Prove that if a point $R$ is equidistant from $A$ and $B$, then it lies
  on the perpendicular bisector of $AB$.
\item Construct the perpendicular bisector of a segment by constructing two
  points on the same side of the line and equidistant from the ends of the
  segment.
\end{exlist}
\end{multicols}

\begin{construction}
To construct the perpendicular to a line at a point $P$ on the line
(Fig.~7--26), draw a circle, with center $P$ and arbitrary radius, intersecting
the line in points $A$ and $B$. Then construct the perpendicular bisector of
$AB$.
\end{construction}

\begin{construction}
To construct a perpendicular to a line from a point not on the line, choose a
point $A$ on the line, and draw the circle with center $P$ and radius $AP$
(Fig.~7--27). The circle will intersect the line in a second point $B$ (if $A$
is properly chosen). Then construct the perpendicular bisector of $AB$, which
passes through $P$.
\end{construction}

\FIGVIITWENTYSIXSEVENEIGHT

\noindent\emph{Proof.}\quad
We first show that there actually is a line through $P$ perpendicular to the
line. Then we show that the \RC{} construction produces this line.

First choose any point $A$ on $l$ and any other point $C$ on $l$ (Fig.~7--28).
Let $AQ$ be the ray, on the opposite side of $l$ from $P$, such that
$\ang{PAC} \cong \ang{QAC}$. (There is only one such ray, by the existence
postulate for angles.) Choose $P'$ on $AQ$, such that $AP' \cong AP$. Then
$PP'$ intersects $l$ in a point $D$. (Why?) If $A \neq D$, then triangles are
formed, and by S.A.S., $\tri{APD} \cong \tri{AP'D}$. Therefore, $\ang{ADP}$ and
$\ang{ADP'}$ are right angles, and $PP' \perp l$. If $A = D$, there are no
triangles, but then, because of the way $AQ$ was selected, $PP' \perp l$.
(Why?)

Having established that a perpendicular from $P$ exists, we now show that the
\RC{} construction gives us this perpendicular. If $D \neq A$, there is a point
$B$ on the opposite side of $D$ from $A$, so that $AD \cong BD$. It follows
that the circle with center $P$ and radius $AP$ does intersect the line in a
second point $B$, and this was the main part to be proved in the construction.
Now the perpendicular bisector of $AB$ can be drawn, and $P$ lies on it.
(Why?)

\exercisegroup{7--7}
\begin{multicols}{2}
\begin{exlist}
\item Construct a right triangle, having been given the lengths of the two
  legs.
\item Construct points $D$, $E$, and $F$ on the segment $AB$, such that
  $\sg{AD} = \sg{DE} = \sg{EF} = \sg{FB} = \frac{1}{4}\sg{AB}$.
\item[\stex 3.]\stepcounter{exlisti} Suppose that $P$ is not on a line $l$, and that a given circle
  with center at $P$ intersects $l$ in exactly one point $Q$. Prove that $PQ$
  is perpendicular to $l$. [\emph{Hint:} If $PQ$ were not perpendicular, then
  the construction in the proof of Construction 7--6 would lead to a
  contradiction.]
\item Construct the perpendicular from $P$ by choosing $A$, as in Fig.~7--29,
  constructing $\ang{2} \cong \ang{1}$, and point $P'$, such that
  $AP \cong AP'$.
\end{exlist}
\end{multicols}

\FIGVIITWENTYNINETHIRTY

\begin{construction}
To bisect $\ang{A}$, draw a circle, with center $A$, intersecting the sides of
$\ang{A}$ in points $B$ and $C$ (Fig.~7--30). Then construct a point $Q$ on the
perpendicular bisector of $BC$ and in the interior of $\ang{A}$. The ray $AQ$
bisects $\ang{A}$.
\end{construction}

\noindent\emph{Proof.}\quad
The point $Q$ can be found by a previous construction.
$\tri{ABQ} \cong \tri{ACQ}$, by S.S.S. Hence $\ang{BAQ} \cong \ang{CAQ}$.

\exercisegroup{7--8}
\begin{multicols}{2}
\begin{exlist}
\item Draw an acute angle, and construct its bisector.
\item Construct and bisect a right angle.
\item Bisect an obtuse angle.
\item Draw a triangle, and construct the bisectors of its angles. (They should
  meet in one point.)
\item If $P$ is on the bisector of $\ang{A}$, and $B$ and $C$ are points on the
  sides of $\ang{A}$, such that $AB = AC$, prove that $PB \cong PC$.
\item Prove that an angle has only one bisector.
\item Suppose that $B$ and $C$ are points on the two sides of $\ang{A}$, such
  that $AB \cong AC$. Show that if $M$ is the midpoint of $BC$, then $AM$
  bisects $\ang{A}$.
\item[\stex 8.]\stepcounter{exlisti} Figure 7--31 shows another construction for the angle bisector.
  Consider two circles, with center $A$ and radii $r$ and $r'$, such that
  $r' > r$. These circles intersect the sides of $A$ in four points $C$, $C'$,
  $D$, and $D'$. The lines through these points intersect in a point $P$ on the
  bisector. (Why?)
\item[\stex 9.]\stepcounter{exlisti} If, in Fig.~7--32, $\ang{CAD} \cong \ang{EBF}$,
  $\ang{1} \cong \ang{2}$, and $\ang{3} \cong \ang{4}$, then
  $\ang{2} \cong \ang{4}$. \emph{Note:} This theorem is often stated as
  ``halves of congruent angles are congruent to each other.'' Prove this
  theorem.
\end{exlist}
\end{multicols}

\FIGVIITHIRTYONETHIRTYTWO

%========================================================= Review
\section*{REVIEW OF CHAPTER 7}
\addcontentsline{toc}{section}{Review of Chapter 7}

In this chapter we have used the congruence theorems to study inequalities in
triangles and verify constructions.

\begin{enumerate}[label=(\alph*), leftmargin=2.2em, itemsep=3pt]
\item Make a list of all the theorems on inequalities in triangles, then study
  their proofs.
\item Make a list of the different \RC{} constructions. Prove them correct
  without looking at the book.
\end{enumerate}

Review the following definitions, which you met for the first time in this
chapter: exterior angle, circle, interior of circle, exterior of circle,
radius, diameter, chord, secant, central angle, inscribed angle, \RC{}
construction.

\begin{center}\bfseries Review Exercises\end{center}\nopagebreak
\begin{multicols}{2}
\begin{exlist}
\item Construct an angle of $45^\circ$.
\item Construct an angle of $60^\circ$, of $30^\circ$.
\item Construct an angle of $15^\circ$, of $75^\circ$.
\item Construct an equilateral triangle.
\item Construct a right triangle having acute angles of $60^\circ$ and
  $30^\circ$.
\item Draw an obtuse triangle, and construct the bisectors of its angles.
\item Draw an angle. Copy it, using ruler and compass.
\item Draw three separate segments. Construct a triangle having sides congruent
  to the given segments.
\item Draw a line, and mark a point $P$ not on the line. Construct the
  perpendicular to the line, passing through $P$. How do you know that there is
  only this one perpendicular?
\item Draw a triangle. Construct the altitudes of the triangle.
\item Draw a triangle. Construct the perpendicular bisectors of its sides.
\item Construct an isosceles triangle if the vertex angle and the congruent
  sides are given.
\item Draw a segment. Construct a circle having this segment as a diameter.
\item In a right triangle, which side is longest? Why?
\item Prove that the sum of the lengths of three sides of a quadrilateral is
  greater than the length of the fourth side.
\item $\ang{ABC}$ is a right angle, and $C$ is between $B$ and $D$. Prove that
  $AD > AC$. (Fig.~7--33)
\item Prove that $\sg{AD} + \sg{CB} > \sg{AB} + \sg{CD}$. (Fig.~7--34)
\item $AE$ bisects $\ang{DAF}$, and $AC$ bisects $\ang{DAB}$. (Fig.~7--35)
  Prove that $\ang{CAE}$ is right. [\emph{Hint:} Extend $AE$.]
\item $AC$ bisects $\ang{DAB}$, and $\ang{CAE}$ is right. Prove that $AE$
  bisects $\ang{DAF}$. (Fig.~7--36)
\end{exlist}
\end{multicols}

\FIGVIITHIRTYTHREETHIRTYFOUR
\FIGVIITHIRTYFIVETHIRTYSIX

\begin{center}\bfseries ALGEBRA REVIEW\end{center}\nopagebreak
\begin{multicols}{2}
\begin{exlist}
\item In $\tri{ABC}$, $AB < BC < CA$. The greatest angle in this triangle has a
  measure of $32^\circ$ more than the measure of the smallest angle. The sum of
  the measures of these two angles is $120^\circ$. Determine the number of
  degrees in the smallest angle in this triangle.
\item An exterior angle at $A$ of $\tri{ABC}$ has a measure of $65^\circ$. The
  sum of the measures of angles $A$ and $B$ is $145^\circ$, and the sum of the
  measures of angles $A$ and $C$ is $150^\circ$. Determine the number of
  degrees in each angle of the triangle.
\item Points $O$, $A$, and $B$ are collinear. $\sg{OA}$ is 7 in. more than
  $\sg{BA}$. $A$ is not between $O$ and $B$. Determine the length of $OB$.
\item $\sg{AB} = 17$, and $\sg{AC} = 12$. The number $\sg{BC}$ is a root of the
  equation $(x-6)(x-4) = 0$. Determine $\sg{BC}$.
\item If $B$ is between $A$ and $C$, $\sg{AC}/\sg{BC} = \frac{5}{3}$, and
  $\sg{AB} = 10$ in., determine the lengths of $AC$ and $BC$.
\item $B$ is between $A$ and $C$, and $\sg{AB}/\sg{AC} = \frac{3}{5}$. The
  midpoint of $AC$ is $M$, and $\sg{MB} = 3$. Determine $\sg{AC}$.
\item[\stex 7.]\stepcounter{exlisti} $A$, $B$, $C$ and $D$ are points in this order on a line. The
  length of $AC$ is twice the length of $BD$. The length of $AB$ is four times
  the length of $CD$. The length of $BC$ is 4. Determine the lengths of $AB$
  and $AD$.
\item[\stex 8.]\stepcounter{exlisti} The difference of the squares of two numbers is 33. The sum of
  these numbers is 11. Determine the numbers.
\item[\stex 9.]\stepcounter{exlisti} Analyze the following mathematical trick. Choose any number;
  multiply this number by 2; add 4 to this product; multiply by 5; divide by
  10; subtract from this result the original number. The result will be 2, no
  matter what number was originally chosen.
\end{exlist}
\end{multicols}
\clearpage
% ==================== ch08 ====================
\setcounter{theorem}{0}\setcounter{definition}{0}\setcounter{construction}{0}\setcounter{exgroup}{0}\setcounter{footnote}{0}
\setcounter{chapter}{8}

\chapter*{\raggedright\Huge PARALLEL LINES}
\addcontentsline{toc}{chapter}{8.\ Parallel Lines}
\markboth{PARALLEL LINES}{}
\vspace{-1em}\noindent\rule{\linewidth}{0.6pt}\vspace{1em}
\figurekey

%========================================================= 8-1
\section*{8--1\quad DEFINITION OF PARALLEL LINES}
\addcontentsline{toc}{section}{8--1\quad Definition of parallel lines}

In earlier mathematics courses, you learned some facts about parallel lines.
Before reading this chapter, try to recall some properties of parallel lines
that you have accepted as true. The following exercises will help your memory
and prepare you for proofs.

\exercisegroup{8--1}
\begin{multicols}{2}
\begin{exlist}
\item What do you mean when you say, ``Line $l$ is \emph{parallel} to line
  $m$''?
\item Name some geometric figures that have pairs of \emph{parallel} sides.
\item Explain why two sides of a triangle are not \emph{parallel}.
\item Two rocks are dropped, one from the Brooklyn Bridge, and the other from
  the Golden Gate Bridge. Would you say that these rocks fall along
  \emph{parallel} paths?
\item If you drop two rocks, one from each hand, do they traverse
  \emph{parallel} paths?
\item Roads $r$ and $s$ run ``due north and south.'' Are they \emph{parallel}?
\item Roads $t$ and $u$ run ``due east and west.'' Are they \emph{parallel}?
\item How many different pairs of \emph{parallel} edges does a cube have?
\item Both lines $l$ and $m \perp$ line $n$. Does this suggest a theorem about
  \emph{parallels}?
\item Each of the lines $l$ and $m$ is \emph{parallel} to line $n$. What
  theorem does this suggest?
\item Lines $p$ and $q$ are \emph{parallel}. Line $r \perp p$. What theorem
  does this suggest?
\item Draw two angles, $\ang{AOB}$ and $\ang{A'O'B'}$, such that the lines
  $AO$ and $A'O'$ are \emph{parallel}, and also lines $BO$ and $B'O'$ are
  \emph{parallel}. Does this figure suggest a theorem? Be careful! Draw more
  than one figure.
\item Draw a line $l$, and mark a point $A$ not on $l$. What might it mean to
  say that a ray from $A$ is \emph{parallel} to line $l$?
\item Which of the above exercises do not strictly belong to plane geometry?
\end{exlist}
\end{multicols}

If you discussed Exercises 6 and 7, someone probably pointed out that roads on
the surface of the earth are not quite the same as lines. Also, in connection
with Exercises 4 and 5, the paths along which objects fall are hardly the
imaginary lines of our geometry. Since the earth rotates about its axis and at
the same time revolves around the sun, an object falling toward the center of
the earth follows a path extremely hard to visualize. Our task is to choose a
definition of parallel lines, such that we can use the definition in our plane
geometry along with our postulates to prove the truth of statements that
concern parallel lines and that we feel ought to be true. The definition below
was perhaps suggested by someone in the class during the discussion of
Exercise Group 8--1.

\begin{definition}
Two lines $l$ and $m$ are parallel if and only if they do not intersect. We
indicate that two lines are parallel by writing ``$l \pll m$.''
\end{definition}

\exercisegroup{8--2}
\begin{multicols}{2}
\begin{exlist}
\item Give definitions of the following. In later sections, it will be
  convenient to have these definitions. Good definitions are not hard to make,
  and so are given as exercises here.
  \begin{enumerate}[label=(\alph*), leftmargin=1.8em, topsep=2pt, itemsep=1pt]
    \item A ray parallel to a line.
    \item Parallel rays.
    \item A line segment parallel to a line.
    \item Parallel line segments.
  \end{enumerate}
\item Draw a figure of a line not intersecting a segment and not parallel to
  the segment.
\item Can a ray and a line not intersect and yet not be parallel?
\end{exlist}
\end{multicols}

Now that we have parallel lines, it is possible for us to talk about rays that
have the same direction. This requires another definition.

\begin{definition}
Two rays have the same \emph{direction} if and only if they are parallel and
both rays are on the same side of the line through their vertices
(Fig.~8--1).
\end{definition}

\FIGVIIIONE

\exercise{8--3}
\noindent What would have been wrong had we defined lines to be parallel if
they had the same direction?

%========================================================= 8-2
\section*{8--2\quad EXISTENCE OF PARALLELS}
\addcontentsline{toc}{section}{8--2\quad Existence of parallels}

Perhaps it will seem strange to you that we are going to prove that parallel
lines actually exist. We could add an extra postulate and assume that there
are parallels, but the proof is not difficult. We give three different proofs
of the following theorem.

\begin{theorem}
If $l$ is any line, and $A$ any point not on $l$, then there exists at least
one line through $A$ parallel to $l$.
\end{theorem}

\emph{First Proof.} Referring to Fig.~8--2, let $B$ be any point on $l$, and
let $n$ be the line through $A$ and $B$. Let $E$ be another point on $l$, and
let $D$ be a point on the same side of $n$ as $E$, such that
$\ang{DAC} \cong \ang{EBA}$. Lines $l$ and $m$ cannot intersect, for if they
did, say at $P$, we should have $\tri{PBA}$ with exterior
$\ang{PAC} \cong \ang{PBA}$. But we know this is impossible, for an exterior
angle of a triangle is always greater than either opposite interior angle.

\FIGVIIITWOTHREE

\emph{*Second Proof. Lemma.} Let line $l$, point $A$, and line $n$ be as in
Theorem 8--1.\VIIIproofnote{} Now, referring to Fig.~8--2, if $m$ is not
parallel to $l$, then $\ang 1$ is not congruent to $\ang 2$. (For ``not
parallel to,'' we use the symbol $\nparallel$.) \emph{Proof.} If $m$ is not
parallel to $l$, these lines intersect, a triangle is formed, and
$\ang 1 \ncong \ang 2$. Now Theorem 8--1 follows quickly from the lemma. The
lemma states that $l \nparallel m \imp \ang 1 \ncong \ang 2$. The
contrapositive of the lemma is $\ang 1 \cong \ang 2 \imp l \pll m$, and the
theorem obviously follows.

\emph{*Third Proof.} Referring to Fig.~8--4, let $l$, $m$, and $n$ refer to
the lines of the preceding proofs with $n$ cutting $m$ at $A$, and $l$ at $B$,
and the indicated angles at $A$ and $B$ congruent to each other. If $l$ and
$m$ intersected at $P$, we could choose $Q$ on $m$, as above, so that
$AQ \cong BP$. It follows at once, by the S.A.S. theorem, that
$\tri{BAQ} \cong \tri{ABP}$. Hence, $\ang{ABQ} \cong \ang{BAP}$.

\FIGVIIIFOUR

Since also, $\ang{ABP} \cong \ang{BAQ}$, we have
$\ang{QBP} \cong \ang{QAP}$. $\ang{QAP}$ is a straight angle, so $\ang{QBP}$
is a straight angle. It follows that $Q$ lies on line $l$, and so if lines $l$
and $m$ intersect in one point, they have a second point in common. This is
impossible, so lines $l$ and $m$ do not intersect.

\exercisegroup{8--4}
\begin{multicols}{2}
\begin{exlist}
\item Draw a line, and mark a point not on it. Use the first proof of Theorem
  8--1 to draw a line, parallel to the first line, through your point.
\item Prove that if $l \perp m$, and $n \perp m$, then, if $l \neq n$,
  $l \pll n$.
\item Is a line parallel to itself?
\item In the first proof of Theorem 8--1, several details were omitted. First
  a point $B$ was chosen and then a point $E$. How was $C$ selected? How would
  we have argued if $l$ and $m$ were to intersect on the other side of $AB$?
\item The second proof of Theorem 8--1 is really just a reworded version of
  the first proof. Explain why this is so.
\item Without pointing at Fig.~8--4, describe how the point $Q$ is chosen in
  the third proof. What postulate guarantees the existence of $Q$?
\item Prove that if $l_1 \pll l_2$, then all the points of $l_2$ are on the
  same side of $l_1$.
\item Use the result of Exercise 7 to give a definition of ``the set of points
  between $l_1$ and $l_2$, where $l_1 \pll l_2$.''
\end{exlist}
\end{multicols}

%========================================================= 8-3
\section*{8--3\quad THE PARALLEL POSTULATE}
\addcontentsline{toc}{section}{8--3\quad The parallel postulate}

We proved in Section 8--2 that through a point $A$ not on line $l$, there is
\emph{at least one line parallel to $l$}. Now we raise the question, \emph{is
there more than one such line?} Our geometric intuition tells us that there is
only the one, and it is natural to attempt to prove this from our other
postulates. For 2000 years, many fine mathematicians tried and failed to find
such a proof. At last, in the 19th century, it was shown that \emph{it is
impossible to prove from the postulates we have already listed that there is
no more than one line through a given point parallel to a given line}. Euclid
himself assumed that there was no more than one such parallel, and we shall
make the same assumption.

\par\medskip\noindent
\textbf{Postulate VI--1.} \emph{The parallel postulate.} Through a point not
on a line $l$, there is no more than one line parallel to $l$.
\par\medskip

\emph{Historical note.} It may seem strange that it is possible to prove that
the parallel postulate is not a consequence of the other postulates. The
understanding of the fact that the parallel postulate is independent of the
others came about in the following way.

After vainly trying to prove the truth of the parallel postulate by using the
other postulates of Euclid, Johann Bolyai (Yo'-hahn Bull'-y\={\i}), a young
Hungarian mathematician, decided to assume that \emph{a second parallel could
be drawn through A}. Using this postulate, he proved many strange and
beautiful theorems. Later it was shown that Bolyai's geometry was just as
acceptable as ordinary Euclidean geometry. We call this geometry of Bolyai's a
\emph{non-Euclidean} geometry to contrast it with our Euclidean geometry which
employs the parallel postulate of Euclid. Whenever you encounter the term
non-Euclidean geometry, you can be sure that some postulate different from the
parallel postulate is used. All the theorems we have proved up to now are true
in the non-Euclidean geometry of Bolyai.

When these things are first explained, many students ask which postulate is
right. Can only one parallel be drawn, or is it possible to draw several?
While these questions are quite natural ones to ask, they do not have the
simple answers that one might expect. Our postulates are statements that we
invent in order to prove interesting and useful theorems. The theorems of
Euclidean geometry are useful in carpentry, drafting, surveying, navigation,
astronomy, and in many other fields. The theorems of some non-Euclidean
geometries have been found to be extremely useful in astronomy and physics.
Today mathematicians study many different geometries. Each of the geometries
has a postulate system different from that of any other geometry. All are
interesting. Many have proved to be of practical use. It does not make sense
to ask which geometry is true. Mathematicians design geometries as architects
design houses. You might as well ask which house is true. It \emph{is}
appropriate to ask which house is better adapted to one's living requirements.
And so it is also appropriate to ask which geometry is most convenient to
describe the universe.*\VIIIgaussnote

\VIIIplate{\textsc{Karl Friedrich Gauss}\\(1777--1855)\\
  Probably the first mathematician to recognize the\\
  existence of non-Euclidean geometries.\\
  Courtesy of \emph{Scripta Mathematica}.}{--- (book p.~137)}

\VIIIplate{\textsc{Nicolai Ivanovich Lobachevsky}\\(1793--1856)\\
  Co-inventor, with Johann Bolyai, of\\
  non-Euclidean geometry.\\
  Courtesy of \emph{Scripta Mathematica}.}{--- (book p.~138)}

%========================================================= 8-4
\section*{8--4\quad CONSEQUENCES OF THE\\PARALLEL POSTULATE}
\addcontentsline{toc}{section}{8--4\quad Consequences of the parallel postulate}

We can now prove several interesting theorems quite easily.

\begin{theorem}
If $l$, $m$, and $n$ are distinct lines, with $l \pll m$ and $m \pll n$, then
$l \pll n$ \textup{(Fig.~8--5)}.
\end{theorem}

\emph{Proof.} If $l$ were not parallel to $n$, then $l$ and $n$ would
intersect, say at $P$ in Fig.~8--5. Two lines, namely $l$ and $n$, would then
pass through $P$, parallel to $m$, thus contradicting the parallel postulate.

\begin{theorem}
If $l \pll m$, and line $n \neq l$ intersects $l$ at $P$, then $n$ intersects
$m$ \textup{(Fig.~8--6)}.
\end{theorem}

\FIGVIIIFIVESIX

\noindent\emph{Proof.}
\begin{steps}
\step{$l \pll m$.}{Given.}
\step{$P$ is on $n$ and $l$, and $n \neq l$.}{Given.}
\step{$l$ is the only line on $P$ parallel to $m$.}{Statements 1, 2, and the
  parallel postulate.}
\step{$n$ intersects $m$.}{Statements 2 and 3.}
\end{steps}

\exercisegroup{8--5}

\FIGVIIISEVENEIGHT

\begin{multicols}{2}
\begin{exlist}
\item Prove that the construction of a parallel to $l$ through a given point
  $A$ is an \RC{} construction.
\item In Fig.~8--7, $\ang{ABC} \cong \ang{BCD}$. Prove that $AB \pll CD$.
\item Prove that if $l \perp n$, $n \perp m$, and $l \neq m$, then
  $l \pll m$. Have you used the parallel postulate?
\item If $l \pll m$, and $l \perp n$, prove that $n \perp m$. Did you use the
  parallel postulate?
\item If $l \pll m$, $k \pll m$, and $l$ and $k$ intersect at $P$, what can
  you say?
\item[\stex 6.] In Fig.~8--8, $\sg{AB} = \sg{CD}$, and $\ang{ABC}$ and
  $\ang{BCD}$ are right. Prove that $AD \pll BC$. [\emph{Hint:} Prove that
  $\ang A \cong \ang D$.]
\item[\stex 7.] If a rectangle is defined as a quadrilateral with four right
  angles, prove that rectangles actually exist.
\end{exlist}
\end{multicols}

%========================================================= 8-5
\section*{8--5\quad CORRESPONDING AND\\ALTERNATE ANGLES}
\addcontentsline{toc}{section}{8--5\quad Corresponding and alternate angles}

When a line intersects two other lines, we often wish to refer to certain
pairs of the angles formed. In order to be precise, we need a definition.

\begin{definition}
A line $l$ (Fig.~8--9) that intersects lines $l_1$ and $l_2$ in distinct
points is called a \emph{transversal} of $l_1$ and $l_2$. We also say that
$l_1$ and $l_2$ are cut by the transversal $l$. Angles 3, 4, 5, and 6 are
called \emph{interior angles}. Angles 1, 2, 7, and 8 are called \emph{exterior
angles}. These eight angles are grouped into four pairs referred to as pairs
of \emph{corresponding angles}. Angles 1 and 5 form such a pair of
corresponding angles, as do also 2 and 6, 3 and 7, and 4 and 8. Angles such as
3 and 5 are called \emph{alternate angles}. Other pairs of alternate angles
are 2 and 8, 1 and 7, and 4 and 6. We speak of the pair 3 and 5 as
\emph{alternate interior angles}. We call the pair 2 and 8 \emph{alternate
exterior angles}. We observe that \emph{alternate} angles have the points of
their interiors on \emph{opposite} sides of $l$, while \emph{corresponding}
angles have their interiors on the \emph{same} side of $l$. In a pair of
corresponding angles, one angle is always exterior and one interior.
\end{definition}

\FIGVIIININETEN

\begin{theorem}
Let $l_1$ and $l_2$ be lines cut by line $l$. If $l_1 \pll l_2$, then two
alternate interior angles are congruent to each other, and conversely, if two
alternate interior angles are congruent, then $l_1 \pll l_2$
\textup{(Fig.~8--10)}.
\end{theorem}

\emph{Proof.} The last part of the theorem follows at once from the first
proof of Theorem 8--1, for if $\ang{CAB} \cong \ang{ABD}$, then
$\ang{ABD} \cong \ang{EAF}$, and $l_1 \pll l_2$.

The first part of the theorem requires the parallel postulate, and may be
proved as follows. Choose $G$, located in the interior of $\ang{DBE}$, such
that $\ang{EAG} \cong \ang{ABD}$. Then line $AG \pll l_2$. (Why?) But $A$ is
also on $l_1$, and $l_1 \pll l_2$. Hence the line through $A$ and $G$ is the
same as $l_1$, so $G$ is on $l_1$. Now,
$\ang{CAB} \cong \ang{EAG} \cong \ang{ABD}$. Thus the alternate interior
angles are congruent.

As important consequences of Theorem 8--4, we state the following corollaries,
which you can easily prove.

\begin{corollary}
Let $l$ cut $l_1$ and $l_2$ at distinct points. If $l_1 \pll l_2$, then
alternate exterior angles are congruent, and conversely.
\end{corollary}

\begin{corollary}
Let $l$ cut $l_1$ and $l_2$ at distinct points. If $l_1 \pll l_2$, then
corresponding angles are congruent, and conversely.
\end{corollary}

\exercisegroup{8--6}
\begin{multicols}{2}
\begin{exlist}
\item Prove that if a line is perpendicular to one of two parallel lines, it
  is perpendicular to the other.
\item Prove that if $l \pll m$, $n \perp l$, $p \perp m$, and $n \neq p$, then
  $n \pll p$.
\item In Fig.~8--11, $\sg{AO} = \sg{CO}$, and $\sg{BO} = \sg{OD}$. Prove that
  $AB \pll CD$.
\item Show that $\sg{BC} = \sg{AD}$, and $BC \pll AD$ in Fig.~8--11.
\item Show that if $AB \cong CD$, and $AB \pll CD$, then $BO \cong OD$, and
  $BC \pll AD$ in Fig.~8--11.
\item If $l \pll m$, and $A$ and $B$ are any two points on $l$, show that the
  perpendicular distances from $A$ and $B$ to $m$ are equal. We may state this
  result in the form: if $l \pll m$, then $l$ and $m$ are ``everywhere
  equidistant.'' State and prove the converse of this theorem. State also the
  contrapositive.
\end{exlist}
\end{multicols}

\FIGVIIIELEVEN

%========================================================= 8-6
\section*{8--6\quad ANGLE SUMS}
\addcontentsline{toc}{section}{8--6\quad Angle sums}

One of the most interesting consequences of the parallel postulate is a
theorem relating to the angles of a triangle. We need a definition.

\begin{definition}
If $\ang{AOB}$ and $\ang{BOC}$ are adjacent angles (Fig.~8--12), and either
$\ang{AOC}$ is a straight angle, or $OB$ is in the interior of $\ang{AOC}$,
then we say that $\ang{AOC}$ is a \emph{sum} of $\ang{AOB}$ and $\ang{BOC}$.
Likewise, if $\ang{A'O'B'} \cong \ang{AOB}$, and
$\ang{B''O''C''} \cong \ang{BOC}$, then we say that $\ang{AOC}$ is a
\emph{sum} of $\ang{A'O'B'}$ and $\ang{B''O''C''}$.
\end{definition}

\FIGVIIITWELVE

\begin{remark}
At this stage in our geometry, we cannot yet ``add'' angles by adding their
measures in degrees, because we have not as yet described how angles are
measured. The measurement of angles is more complicated than measurement of
length, and will be described in a later chapter.
\end{remark}

Draw a figure to illustrate the last part of Definition 8--4.

\begin{definition}
Two angles are said to be \emph{complementary} if their sum is a right angle.
\end{definition}

\begin{definition}
Two angles are said to be \emph{supplementary} if their sum is a straight
angle.
\end{definition}

\exercisegroup{8--7}
\begin{multicols}{2}
\begin{exlist}
\item Draw an obtuse angle. Draw an angle that is supplementary to this angle
  and is not adjacent to it. Draw one that is adjacent and supplementary.
\item Draw an acute angle. Draw an angle that is complementary to this angle
  and is adjacent to it. Draw one that is not adjacent, but is complementary,
  to the acute angle.
\item Prove that if two acute angles are congruent, then their complements are
  congruent.
\item Prove that if the sides of one angle are parallel, respectively, to the
  sides of a second angle, then either the angles are congruent or they are
  supplementary.
\end{exlist}
\end{multicols}

\begin{theorem}
The sum of the angles of a triangle is a straight angle.
\end{theorem}

\FIGVIIITHIRTEEN

\emph{Proof.} In Fig.~8--13, let $l$ be the line through $A$ and parallel to
the side $BC$ of $\tri{ABC}$. Let $D$ be a point on $l$ that is on the
opposite side of $AB$ from $C$. Let $E$ be a point of $l$ on the same side of
$AB$ as $C$. Then $\ang{DAB} \cong \ang{ABC}$, and
$\ang{EAC} \cong \ang{ACB}$. The straight angle $DAE$ is the sum of angles
$A$, $B$, and $C$.

\exercisegroup{8--8}

\FIGVIIIFOURTEEN

\begin{multicols}{2}
\begin{exlist}
\item Show that the acute angles of a right triangle are complementary.
\item Two angles of a triangle are $40^\circ\ 17'\ 30''$ and
  $78^\circ\ 31'\ 38''$, respectively. Determine the number of degrees in the
  third angle.
\item How many degrees are there in each angle of an equilateral triangle?
\item Show that a median of an equilateral triangle divides the triangle into
  two right triangles, each of which has acute angles of $30^\circ$ and
  $60^\circ$.
\item Prove that in a right triangle, with acute angles of $30^\circ$ and
  $60^\circ$, the side opposite the $30^\circ$ angle is half as long as the
  hypotenuse.
\item In $\tri{ABC}$, an exterior angle at $B$ is $124^\circ$. $\ang A$ is
  $58^\circ$. How many degrees in $\ang C$?
\item Prove that if two angles of one triangle are congruent, respectively, to
  two angles of another triangle, then the third angles are congruent.
\item In Fig.~8--14, $AB \cong A'B'$, $\ang A \cong \ang{A'}$, and
  $\ang C \cong \ang{C'}$. Prove that $\tri{ABC} \cong \tri{A'B'C'}$.
\item Is the following true? Two triangles are congruent to each other if one
  side and two angles of the one triangle are congruent to one side and two
  angles of the other.
\item Prove that two right triangles are congruent if the hypotenuse and an
  acute angle of one triangle are congruent, respectively, to the hypotenuse
  and acute angle of the other.
\item In $\tri{ABC}$ (Fig.~8--15), $\ang A$ is right, and $AD \perp BC$. Prove
  that $\ang{BAD} \cong \ang{ACD}$, and $\ang{ABD} \cong \ang{DAC}$.
\item[\stex 12.] In Fig.~8--16, $D$ is the midpoint of $BC$, and
  $BD \cong AD$. Prove that $\ang{BAC}$ is right.
\item[\stex 13.] Prove that two right triangles are congruent if the
  hypotenuse and one side of one triangle are congruent to the hypotenuse and
  one side of the other.
\end{exlist}
\end{multicols}

\FIGVIIIFIFTEENSIXTEEN

%========================================================= 8-7
\section*{8--7\quad POLYGONS}
\addcontentsline{toc}{section}{8--7\quad Polygons}

In this and the remaining sections, we will be studying certain polygons.
First, we need to recall the definition of a polygonal path.

A \emph{polygonal path} (Fig.~8--17) is determined by a number of points,
$P_1, P_2, \ldots, P_n$, called the \emph{vertices}, given in a definite
order. The path is the set of points on the segments
$P_1P_2, P_2P_3, \ldots, P_{n-1}P_n$. The point $P_1$ is called the
\emph{beginning} of the path, and $P_n$ is called the \emph{end}. The segments
$P_1P_2$, etc., are called the \emph{sides} of the path.

\FIGVIIISEVENTEENEIGHTEEN

\begin{remark}
There is no requirement that the sides of a polygonal path not intersect.
Hence the path can be quite complicated. See Fig.~8--18 for a six-sided path.
\end{remark}

\exercisegroup{8--9}
\begin{multicols}{2}
\begin{exlist}
\item Draw a polygonal path with five sides that does not cross itself. Draw
  one that does.
\item If the length of a polygonal path is the sum of the lengths of its
  sides, how does the length of a polygonal path $P_1, \ldots, P_n$ compare
  with $\sg{P_1P_n}$?
\end{exlist}
\end{multicols}

\begin{definition}
A \emph{polygon} is a polygonal path whose beginning and end coincide. If the
vertices are all different and no two sides have a point in common (other than
a vertex of two adjacent sides), then the polygon is called \emph{simple}
(Fig.~8--19).
\end{definition}

\FIGVIIININETEEN

Each simple polygon separates the points of the plane, other than the polygon
itself, into two sets called the \emph{interior} and \emph{exterior} of the
polygon. If $A$ is a point of the interior, and $B$ is a point of the
exterior, every polygonal path connecting $A$ to $B$ contains at least one
point of the polygon. If $A$ and $B$ both belong to either the interior or
exterior, then there exist polygonal paths joining $A$ to $B$ that contain no
point of the polygon. The proofs of these statements are quite difficult, and
will not be given here. It would be instructive to carry out the proofs in the
case in which the polygon is a triangle. The student interested in doing so
should restudy Section 3--6.

\emph{*Problem.} The theorem that a simple polygon separates the plane is
powerful and has surprising applications. It can be used to show that the
well-known puzzle below has no solution.

It is desired to connect each of three houses $A$, $B$, and $C$ with three
outlets, one for gas $G$, one for water $W$, and one for electricity $E$.

\VIIIutilities

Is it possible to run the lines, one from each house to each outlet, in such a
way that no lines cross? In other words, can polygonal paths from $A$ to $G$,
$A$ to $W$, etc., be drawn so that no two paths cross?

\begin{definition}
Vertices of a polygon that are endpoints of the same side are called
\emph{adjacent} vertices (Fig.~8--20). Angles whose vertices are adjacent to
each other are called \emph{consecutive} angles. A \emph{diagonal} of a
polygon is a segment joining any two nonadjacent vertices. Sides with a common
vertex are called \emph{adjacent} sides.
\end{definition}

\FIGVIIITWENTYTWENTYONE

\begin{definition}
A simple polygon is called \emph{convex} if, whenever $A$ and $B$ are points
of the polygon, every point of the segment $AB$ lies either on the polygon or
in its interior (see Fig.~8--21).
\end{definition}

\begin{remark}
Unless otherwise stated, we will confine our attention to simple convex
polygons.
\end{remark}

\begin{definition}
If all sides of a polygon have the same length, the polygon is
\emph{equilateral}. If all angles are congruent, it is \emph{equiangular}. A
convex polygon that is equilateral and equiangular is called \emph{regular}.
Polygons are also classified according to the number of sides. Polygons of 3,
4, 5, 6, 7, 8, 9, and 10 sides are called \emph{triangles},
\emph{quadrilaterals}, \emph{pentagons}, \emph{hexagons}, \emph{heptagons},
\emph{octagons}, \emph{nonagons}, and \emph{decagons}, respectively. A polygon
of $n$ sides will be called an \emph{$n$-gon}.
\end{definition}

\exercisegroup{8--10}
\begin{multicols}{2}
\begin{exlist}
\item What do we call a regular polygon of three sides? of four sides?
\item Draw a convex polygon of five sides. Draw a quadrilateral that is not
  convex. Draw a polygon that is not simple.
\item[\stex 3.] Is it true that every interior point of a convex polygon is
  also interior to every angle of the polygon?
\item[4.] We can divide a convex polygon into triangles by drawing the
  diagonals from any vertex. Into how many triangles does this procedure
  divide a convex polygon of four sides? of ten sides? of $n$ sides?
\item[\stex 5.] How many diagonals has a convex polygon of four sides? five
  sides? ten sides? $n$ sides?
\item[6.] Show that an equiangular polygon need not be equilateral.
\item[\stex 7.] Find an example of an equilateral, equiangular polygon that is
  not regular.
\end{exlist}
\end{multicols}

A diagonal of a convex quadrilateral divides the quadrilateral into two
triangles. Since the sum of the angles in each triangle is a straight angle,
we shall say that \emph{the sum of the (four) angles of the quadrilateral is
two straight angles}. The two diagonals from one vertex of a convex pentagon
divide the pentagon into three triangles, and we say that \emph{the sum of the
(five) angles of the pentagon is three straight angles}.

When we have defined degree measure for angles, we will see that the sum of
the degree measures of the angles of a convex quadrilateral is
$2 \cdot 180^\circ = 360^\circ$, of a pentagon,
$3 \cdot 180^\circ = 540^\circ$, and in general for a polygon of $n$ sides,
the sum of degree measures is $(n - 2) \cdot 180^\circ$.

Do not misunderstand the statement that the sum of the angles of a pentagon is
three straight angles. This is using the word ``sum'' in a different sense
than when we speak of the sum of \emph{two angles}. Recall that when we add
two angles, the sum is less than or congruent to a straight angle.

\FIGVIIITWENTYTWOTHREE

We observe that the sum of the five interior angles and the five indicated
exterior angles in Fig.~8--22 is precisely five straight angles. Since the
interior angles account for three straight angles, the sum of the exterior
angles must be two straight angles. Similar reasoning shows that for any
convex polygon of $n$ sides, the sum of the $n$ exterior angles (one at each
vertex) is exactly two straight angles (Fig.~8--23).

\exercisegroup{8--1}

\FIGVIIITWENTYFOUR

\begin{multicols}{2}
\begin{exlist}
\item A quadrilateral has three $95^\circ$ angles. How many degrees are there
  in the fourth angle?
\item How many degrees are there in each angle of a regular hexagon?
\item How many degrees are there in each exterior angle of a regular pentagon?
\item A convex quadrilateral has three angles of the same size. The fourth
  angle is a $120^\circ$ angle. How many degrees are there in each of the
  other angles?
\item Explain why it is impossible for a convex polygon to have as many as six
  angles smaller than $120^\circ$.
\item[\stex 6.] What is the sum of the angles at the vertices of a regular
  five-pointed star? (The points of the star are vertices of a regular
  pentagon.)
\item[7.] An exterior angle of a regular polygon is $20^\circ$. How many sides
  has the polygon?
\item[8.] Pentagon $ABCDE$ is regular. Prove that a line $l$, bisecting the
  angle at $A$, is perpendicular to side $CD$.
\item[9.] How many sides does a regular polygon have if each angle is
  $120^\circ$? $140^\circ$? $160^\circ$? $179^\circ$?
\item[10.] Make up some theorems of your own relating to regular pentagons or
  regular hexagons.
\end{exlist}
\end{multicols}

%========================================================= 8-8
\section*{8--8\quad PARALLELOGRAMS}
\addcontentsline{toc}{section}{8--8\quad Parallelograms}

The following definitions make certain familiar terms precise.

\begin{definition}
A \emph{parallelogram} is a quadrilateral having two pairs of parallel sides.
A \emph{rectangle} is a parallelogram with a right angle. A \emph{square} is a
rectangle with two congruent perpendicular sides. A \emph{rhombus} is an
equilateral parallelogram.
\end{definition}

\FIGVIIITWENTYFIVE

Our knowledge of parallels and these definitions make it possible to prove
many results, as the following exercises will demonstrate.

\exercisegroup{8--12}
\noindent The exercises below present many interesting theorems concerning
polygons. Give proofs for these theorems.

\begin{multicols}{2}
\begin{exlist}
\item A diagonal of a parallelogram divides it into two congruent triangles.
\item Opposite sides of a parallelogram are congruent to each other.
\item The diagonals of a parallelogram bisect each other.
\item Opposite angles of a parallelogram are congruent.
\item Consecutive angles of a parallelogram are supplementary to each other.
\item Every angle of a rectangle is right.
\item All sides of a square are congruent.
\item The diagonals of a rectangle are congruent.
\item The diagonals of a square are perpendicular to each other.
\item The diagonals of a rhombus are perpendicular.
\item If the diagonals of a parallelogram are congruent, the parallelogram is
  a rectangle.
\item If the diagonals of a rectangle are perpendicular, the rectangle is a
  square.
\item The quadrilateral formed by connecting midpoints of the sides of a
  parallelogram is also a parallelogram.
\item The parallelogram formed by connecting midpoints of the sides of a
  square is also a square.
\item If the diagonals of a quadrilateral bisect each other, the quadrilateral
  is a parallelogram.
\item If a quadrilateral has a pair of sides congruent and parallel to each
  other, it is a parallelogram.
\item A parallelogram has one angle of $70^\circ$. Give the number of degrees
  in each angle.
\item One angle of a parallelogram is twice as large as a second angle.
  Determine the number of degrees in each angle of the parallelogram.
\item A diagonal of a rhombus bisects two angles of the rhombus.
\item If a rectangle is not a square, a diagonal does not bisect an angle of
  the rectangle.
\item If each angle of a quadrilateral is congruent to the angle opposite it,
  then the quadrilateral is a parallelogram.
\item Is a rhombus a regular polygon?
\item For every quadrilateral, it is true that either the bisectors of two
  consecutive angles of the quadrilateral are perpendicular, or the
  quadrilateral is not a parallelogram.
\item If a diagonal of a parallelogram bisects one angle of the parallelogram,
  it bisects a second angle, and the parallelogram is a rhombus.
\item If a parallelogram is not a rhombus, then the four bisectors of the
  angles of the parallelogram are not concurrent.
\item If a parallelogram is not a rhombus, then the four bisectors of the
  angles intersect to form a rectangle.
\item The four bisectors of the angles of a rectangle that is not a square
  intersect to form a square.
\item How many degrees are there in the angle formed by the two diagonals
  drawn from one vertex of a regular pentagon?
\item In Fig.~8--26, $DE \pll AC$, and $DF \pll AB$. Use the figure to prove
  that the sum of the angles of $\tri{ABC}$ is a straight angle.
\item In the regular hexagon $ABCDEF$, prove that diagonal $FD$ is
  perpendicular to side $DC$ (Fig.~8--27).
\item Prove that $\sg{FC} = 2 \cdot \sg{DC}$ in Fig.~8--27. Show that the
  diagonals $FC$ and $AD$ bisect each other.
\end{exlist}
\end{multicols}

\FIGVIIITWENTYSIXSEVEN

%========================================================= Review of Chapter 8
\section*{REVIEW OF CHAPTER 8}
\addcontentsline{toc}{section}{Review of Chapter 8}

This chapter has dealt with consequences of the parallel postulate, by far the
most famous of all the postulates of Euclid. (Euclid did not state it exactly
as we have given it.) It was through the careful investigation of this
postulate that non-Euclidean geometries were invented early in the 19th
century. This investigation also led indirectly to the first set of correct
postulates from which the theorems of Euclidean geometry could be rigorously
proved. These postulates, presented about 1900 by the German mathematician
David Hilbert, are equivalent to the ones that we have listed in this text. We
list below some of the most important consequences of the parallel postulate
shown in this chapter.

\begin{quote}
If two lines are parallel to one line, they are parallel to each other.

When parallel lines are cut by a line, alternate interior angles are
congruent.

The sum of the angles of any triangle is a straight angle.

The sum of the interior angles of any convex polygon is $(n - 2)$ straight
angles.
\end{quote}

Several new terms were defined in this chapter. In addition, precise
definitions of some familiar concepts were formulated. (For example, though we
have referred freely to squares and rectangles in earlier chapters, they had
not yet been exactly defined. Indeed, it is not even possible to prove that
squares and rectangles exist until the parallel postulate is available.) You
should know the definitions given in this chapter for the following:

\begin{center}
\begin{tabular}{@{}l@{\hspace{3.2em}}l@{}}
alternate interior angles & octagon \\
corresponding angles      & $n$-gon \\
polygon                   & regular polygon \\
simple polygon            & parallelogram \\
convex polygon            & rectangle \\
parallel lines            & square \\
pentagon                  & rhombus \\
hexagon                   & diagonal \\
equilateral polygon       & equiangular polygon \\
\end{tabular}
\end{center}

%--------------------------------------------------------- Algebra Review
\section*{ALGEBRA REVIEW}
\addcontentsline{toc}{section}{Algebra Review (Chapter 8)}

\begin{exlist}
\item Two angles of a triangle are congruent, and the measure of each exceeds
  that of the third angle by $10^\circ$. Determine the number of degrees in
  each angle of the triangle.
\item In $\tri{ABC}$, the sum of $\ang B$ and an exterior angle at $A$ has a
  measure of $170^\circ$. The sum of this exterior angle and $\ang C$ is
  $160^\circ$. Determine the measure of each angle of the triangle.
\item An exterior angle of a triangle has a measure of $117^\circ$. One angle
  of this triangle has a measure of $42^\circ$. Determine the degree measure
  of each angle of the triangle.
\item Prove that there is no $\tri{ABC}$ such that the measure of an exterior
  angle at $A$ is $103^\circ$ and an exterior angle at $B$ is $74^\circ$.
\item Diagonal $AC$ of parallelogram $ABCD$ forms with sides $AB$ and $AD$ two
  angles whose measures differ by $10^\circ$. The measure of $\ang B$ is
  $110^\circ$. Determine the measure of $\ang A$.
\item Two angles of a triangle are complementary. Prove that the measure of
  the third angle is $90^\circ$.
\item The measure of the vertex angle of an isosceles triangle is $84^\circ$.
  Find the measure of each angle of the triangle.
\item The first angle of a triangle has a measure three-fourths of the second.
  The first angle is also one-third of the sum of the other two. Compute each
  angle.
\item In Fig.~8--28 the letters represent the degree measures of the indicated
  angles. Prove that $x + y + t = r + s + z$.
\item In Fig.~8--29, $AB \cong BC$, $BE \cong BD$, and the measure of
  $\ang{CBD}$ is $40^\circ$. Determine the measure of $\ang{ADE}$.
\end{exlist}

\FIGVIIITWENTYEIGHTNINE

%========================================================= Review, Chs 1-8
\section*{REVIEW EXERCISES,\\CHAPTERS 1 THROUGH 8}
\addcontentsline{toc}{section}{Review Exercises, Chapters 1 through 8}

\begin{exlist}
\item Give truth values of $\al$, $\be$, and $\ga$ for which a statement of the
  form ($\al$ or $\be$) and $\ga$ is true; is false.
\item Can you find truth values of $\al$, $\be$, and $\ga$ for which ($\al$ or
  $\be$) and $\ga$ is false, and $\al$ or ($\be$ and $\ga$) is true? Can you
  find truth values for which ($\al$ or $\be$) and $\ga$ is true, and $\al$ or
  ($\be$ and $\ga$) is false? Is ($\al$ or $\be$) and
  $\ga \imp \al$ or ($\be$ and $\ga$) true? Is $\al$ or ($\be$ and $\ga$)
  $\imp$ ($\al$ or $\be$) and $\ga$ true?
\item Formulate an ``if\ldots, then\ldots'' statement equivalent to (a) no
  fishermen lie, (b) only fishermen lie, (c) fishermen necessarily lie.
\item Give the contrapositive of each of your implications in Exercise 3.
\item One set of points contains, among others, the points $A$, $B$, $C$, and
  $D$. A second set of points contains $A$, $B$, and $D$, but does not contain
  the point $C$. Prove that these two sets of points cannot both be lines.
\item Point $C$ is in the interior of $\ang{AOB}$. Is every point of ray $OC$
  in the interior of this angle?
\item Define \emph{interior of a triangle} in at least two ways.
\item Explain what is meant by saying that \emph{congruence of segments} has
  the \emph{reflexive} property; the \emph{symmetric} property; the
  \emph{transitive} property.
\item Explain why one of these statements is true, and the other false:\\
  (a) $AB \cong CD \imp AB = CD$, (b) $AB = CD \imp AB \cong CD$.
\end{exlist}

\FIGVIIITHIRTY

\begin{exlist}
\item[10.] Figure 8--30 indicates that $AB \cong A'B'$, and $ABC$. Using only
  the definitions of \emph{greater than} and \emph{less than}, does the figure
  indicate that $AC > A'B'$, or that $A'B' < AC$?
\item[11.] The length of $AB$, measured with $CD$ as unit, is 17.23. What is
  Archimedes' integer?
\item[12.] The length of $CD$ in Exercise 11, as measured by $EF$, is 10. If
  $AB$ is measured by $EF$, what is its length? What is Archimedes' integer in
  this case?
\item[13.] Express the rational number $\tfrac{3}{19}$ as a repeating decimal.
  Find the rational number whose decimal representation is the repeating
  decimal $0.272727\ldots$; $0.027027027\ldots$.
\item[14.] State the S.A.S. postulate for triangles.
\item[15.] Two lines that intersect to form congruent adjacent angles are said
  to be perpendicular to each other. Does this definition in itself prove that
  there is a pair of perpendicular lines in the plane? How do we know that
  right angles exist?
\item[16.] Are there points $A$, $B$, and $C$, such that $\sg{AB} = 5$,
  $\sg{BC} = 7$, and $\sg{AC} = 13$?
\item[17.] What can be said about the position of the points $R$, $S$, and
  $T$, if $\sg{RS} = 8$, $\sg{ST} = 19$, and $\sg{RT} = 11$?
\item[18.] Consider a circle whose center is $A$ and which contains the point
  $B$.
  \begin{enumerate}[label=(\alph*), leftmargin=1.8em, topsep=2pt, itemsep=1pt]
    \item If $\sg{AX} > \sg{AB}$, what can be said about point $X$?
    \item If $P$ and $Q$ are on the circle, how do we refer to $\ang{PAQ}$? to
      segment $PQ$?
    \item If $O$ and $C$ are on the circle, how do we refer to $\ang{AOC}$?
    \item If $R$ and $S$ are points of the circle, and $A$ is on the segment
      $RS$, what do we call segment $RS$?
  \end{enumerate}
\item[19.] What does it mean to say that two rays have the same direction?
\item[\stex 20.] Without using the parallel postulate or any of its
  consequences, prove the following theorem.
  \begin{quote}
  \emph{If in triangles $ABC$ and $A'B'C'$, $AB \cong A'B'$,
  $\ang B \cong \ang{B'}$, and $\ang C \cong \ang{C'}$, then
  $\tri{ABC} \cong \tri{A'B'C'}$.}
  \end{quote}
\item[\stex 21.] Without using the parallel postulate, prove the following
  theorem, which relates to Fig.~8--31.
  \begin{quote}
  \emph{If $\ang{CAB}$ and $\ang{ABD}$ are right, $CA \cong DB$, and $M$ and
  $N$ are midpoints of $CD$ and $AB$, respectively, then $MN$ is perpendicular
  both to $AB$ and to $CD$.}
  \end{quote}
\end{exlist}

\FIGVIIITHIRTYONETWO

\begin{exlist}
\item[\stex 22.] Carry through the proofs of Exercises 20 and 21 using the
  parallel postulate.
\item[\stex 23.] Euclid's formulation of the parallel postulate is different
  from ours. The following statement is essentially the postulate that Euclid
  used:
  \begin{quote}
  \emph{If the points $X$ and $Y$ are on one side of line $AB$, and if the sum
  of $\ang{XAB}$ and $\ang{YBA}$ is less than a straight angle, then the rays
  $AX$ and $BY$ intersect.}
  \end{quote}
  Figure 8--32 illustrates Euclid's parallel postulate. (a) Prove Euclid's
  parallel postulate as a theorem, using our parallel postulate. (b) Assuming
  the truth of Euclid's parallel postulate, prove ours as a theorem.
\end{exlist}
\clearpage
% ==================== ch09 ====================
\setcounter{theorem}{0}\setcounter{definition}{0}\setcounter{construction}{0}\setcounter{exgroup}{0}\setcounter{footnote}{0}
\setcounter{chapter}{9}

\chapter*{\raggedright\Huge SIMILARITY OF\\TRIANGLES AND POLYGONS}
\markboth{SIMILARITY OF TRIANGLES AND POLYGONS}{}
\addcontentsline{toc}{chapter}{9.\ Similarity of Triangles and Polygons}
\vspace{-1em}\noindent\rule{\linewidth}{0.6pt}\vspace{1em}
\figurekey

%========================================================= 9-1
\section*{9--1\quad INTRODUCTION}
\addcontentsline{toc}{section}{9--1\quad Introduction}

In this chapter we study relations between two polygons that have ``the same
shape.'' (Later on we explain precisely what we mean by ``the same shape,'' and
say that such figures are \emph{similar to each other}.) Then we prove some
theorems about similar triangles, and using these theorems we prove the
Pythagorean theorem, which states that if $a$ and $b$ are lengths of the sides of
a right triangle, and $c$ is the length of the hypotenuse, then $c^2 = a^2 + b^2$.
We see as we study this chapter, that the theorems on similar triangles depend
very much on the properties of parallel lines. For this reason, we first continue
our study of parallel lines.

%========================================================= 9-2
\section*{9--2\quad TRANSVERSALS}
\addcontentsline{toc}{section}{9--2\quad Transversals}

The results of this chapter are based on the important theorem about parallel
lines given below. In the theorem, one of the lines is \emph{between} the other
two. You should be able to tell what this means in terms of the ideas of
betweenness that you have studied.

\exercise{9--1}
If $l$, $l'$, and $l''$ are three parallel lines, what should it mean to say that
$l'$ is between $l$ and $l''$?

\begin{theorem}
If $l$, $l'$, and $l''$ are parallel lines, and $m$ is a transversal intersecting
the lines in points $A$, $B$, and $C$, respectively, such that
$\sg{AB} \cong \sg{BC}$, and if $m'$ is any other transversal intersecting $l$,
$l'$, and $l''$ in $A'$, $B'$, and $C'$, respectively, then
$\sg{A'B'} \cong \sg{B'C'}$.
\end{theorem}

\begin{remark}
A brief way of stating the theorem is to say that if parallel lines cut off
congruent segments on one transversal, then they cut off congruent segments on
every transversal.
\end{remark}

\noindent\emph{Proof.}
\begin{steps}
\step{Consider lines $m''$ and $m'''$ through $A$ and $B$, respectively, and
      parallel to $m'$ in Fig.~9--1.}{Which postulate?}
\step{$m''$ intersects $l'$ in $D$, and $m'''$ intersects $l''$ in $E$.}{Why?}
\step{$A$, $A'$, $B'$, and $D$ are vertices of a parallelogram, as also are $B$,
      $B'$, $C'$, and $E$.}{Since $l$, $l'$, and $l''$ are parallel, as also are
      $m'$, $m''$, and $m'''$.}
\step{$\sg{AD} \cong \sg{A'B'}$, $\sg{BE} \cong \sg{B'C'}$.}{Why?}
\step{$\ang{ABD} \cong \ang{BCE}$,\newline $\ang{BAD} \cong \ang{CBE}$.}{Why?}
\step{$\sg{AB} \cong \sg{BC}$.}{Given.}
\step{$\tri{ABD} \cong \tri{BCE}$.}{A.S.A.}
\step{$\sg{AD} \cong \sg{BE}$.}{Statement 7.}
\step{$\sg{A'B'} \cong \sg{B'C'}$.}{Statements 4 and 8 and properties of
      congruence.}
\end{steps}

\FIGIXONE

\FIGIXTWOTHREE

\begin{corollary}
If a line is parallel to one side of a triangle and bisects another side, then it
bisects the third side.
\end{corollary}

\begin{corollary}
If a line bisects two sides of a triangle, it is parallel to the third side, and
the segment formed is half as long as the third side. (In Fig.~9--2,
$DE \pll AC$, and $\sg{DE} = \frac{1}{2}\sg{AC}$.)
\end{corollary}

\noindent\emph{Proof.} It is given that $\sg{AD} = \sg{DB}$, and
$\sg{BE} = \sg{EC}$. Through $D$, let $m$ be a line parallel to $AC$, as in
Fig.~9--3. By Corollary 9--1--1, $m$ meets $BC$ in a point $F$, which bisects
$BC$. Since a segment has only one bisector, $F = E$. To see that
$\sg{DE} = \frac{1}{2}\sg{AC}$, consider a line $m'$ through $E$, parallel to
$AB$, as in Fig.~9--4. Then $m'$ meets $AC$ in $G$, which bisects $AC$. It is
left for the student to show that $\sg{AC} = 2\sg{DE}$.

\FIGIXFOUR

\exercisegroup{9--2}
\begin{multicols}{2}
\begin{exlist}
\item The three sides of a triangle have lengths 5, 12, and 13. What are the
      lengths of the segments joining the midpoints of the sides?
\item In Fig.~9--5, $FC \cong CG$, $CD \pll EG$, and $DG \pll EF$. Prove that
      $\sg{EG} = 2\sg{CD}$.
\item Prove Corollary 9--1--1.
\item In Fig.~9--6, if $\sg{AB} = 2\sg{BC}$, how are $\sg{A'B'}$ and
      $\sg{B'C'}$ related? Suppose that $\sg{AB} = 3\sg{BC}$; $3\sg{AB} =
      5\sg{BC}$. Prove your statements.
\item In Fig.~9--6, suppose that $\sg{AC} = 10$ in., $\sg{BC} = 2$ in., and
      $\sg{B'C'} = 3$ in. How long is $A'B'$?
\item In Fig.~9--6, if $\sg{AB} = 6$ in., $\sg{A'B'} = 8$ in., and $\sg{AC} =
      12$ in., what is $\sg{B'C'}$?
\item In Fig.~9--7, if $\sg{EF} = 2$, $\sg{FG} = 3$, and $\sg{EP} = 3$, what is
      $\sg{EQ}$?
\item In the proof of Theorem 9--1, we assumed that $m$ was not parallel to
      $m'$. Prove Theorem 9--1 if $m \pll m'$.
\end{exlist}
\end{multicols}

\FIGIXFIVESEVEN

%========================================================= 9-3
\section*{9--3\quad A CONSTRUCTION}
\addcontentsline{toc}{section}{9--3\quad A Construction}

In this section we show how to divide a segment into any number of congruent
segments. The construction can be made with ruler and compass.

\FIGIXEIGHT

\begin{construction}
Given the segment $AB$, the problem is to divide it into a number of congruent
segments. (In Fig.~9--8, the number is four.)
\end{construction}

\begin{steps}
\step{Draw a ray $AC$, forming an angle with $AB$ that is not a straight
      angle.}{How do you know that such a ray exists?}
\step{With a compass, mark off points $D$, $E$, $F$, and $G$, in order $ADEFG$,
      such that $\sg{AD} \cong \sg{DE} \cong \sg{EF} \cong \sg{FG}$.}{This can be
      done because of the congruence postulate and betweenness postulates.}
\step{Draw $BG$.}{Postulate I.}
\step{Draw lines through $D$, $E$, and $F$, $\pll BG$.}{These lines exist by
      Theorem 8--1. The construction is \RC. Why?}
\step{The parallel lines through $D$, $E$, and $F$ intersect $AB$ in points $H$,
      $I$, and $J$, respectively.}{Why?}
\step{$\sg{AH} \cong \sg{HI} \cong \sg{IJ} \cong \sg{JB}$.}{Theorem 9--1 on
      transversals.}
\end{steps}

\begin{remark}
In discussing measurement in Chapter 5, we had to divide the unit segment into
tenths, and then hundredths, etc. The above construction shows that, with the
parallel postulate, such a division of the unit segment is possible. At no time
between Chapter 5 and now have we needed measurement for the development of our
geometry.
\end{remark}

\exercisegroup{9--3}
\begin{multicols}{2}
\begin{exlist}
\item Draw a segment, and divide it into three segments of equal length. This is
      called \emph{trisecting} a segment.
\item Draw a segment, and divide it into five congruent segments.
\item In Fig.~9--8, suppose that points $I$ and $J$ have been found. State two
      different ways to find $H$.
\item Can you divide a segment into four equal parts without using the above
      construction? eight equal parts? Is this an \RC\ construction?
\item If, in Fig.~9--9, $\sg{AB} = 10$, $\sg{AD} = 2$, and $\sg{EC} = 3$, what is
      $\sg{BC}$?
\item If, in Fig.~9--9, $\sg{AD} = 6$, $\sg{DB} = 8$, and $\sg{EC} = 2$, what is
      $\sg{BC}$?
\item In Fig.~9--10, $\sg{AB} = 18$, $\sg{BC} = 3$, and $\sg{AD} = 24$. What is
      $\sg{EF}$?
\item Using Fig.~9--10, find $\sg{EF}$ if $\sg{BC} = 1$, $\sg{AB} = 2$, and
      $\sg{AD} = 3$.
\item In Fig.~9--11, if $BE \pll CF \pll DG$, and the lengths are as shown, how
      long is $AG$?
\item Solve Exercise 9 if everything is the same except that $\sg{AE} = 6$.
\item In Fig.~9--12, how long is $AB$?
\item[\stex 12.] In Fig.~9--13, $ABCD$ is a parallelogram, and $E$ and $F$ are
      midpoints of $AB$ and $CD$, respectively. Prove that $AC$ intersects $ED$
      and $BF$ in $G$ and $H$, and that $G$ and $H$ trisect $AC$. Prove also that
      $\sg{EG} = \frac{1}{3}\sg{ED}$.
\item[\stex 13.] In Fig.~9--14, $x$, $y$, $z$, and 1 refer to the lengths of the
      segments. If $x$ is a whole number, show that $z = xy$.
\end{exlist}
\end{multicols}

\FIGIXNINEFOURTEEN

%========================================================= 9-4
\section*{9--4\quad PROPORTION}
\addcontentsline{toc}{section}{9--4\quad Proportion}

Before discussing similar triangles, we recall some facts about proportions that
you learned in arithmetic and algebra.

\begin{definition}
If $x$ and $y$ are numbers, with $y \neq 0$, the ratio of $x$ to $y$ is the
number $x/y$, that is, $x$ divided by $y$. A \emph{proportion} is a statement
that two (or more) ratios are equal. For example, $a/b = c/d$ is called a
proportion. In this proportion, we say that $a$ and $c$ are proportional to $b$
and $d$, respectively, and also that ``$a$ is to $b$ as $c$ is to $d$.'' More
generally, the statement that $x$, $y$, $z$, $u$, \dots\ are proportional to
$x'$, $y'$, $z'$, $u'$, \dots\ means that $x/x' = y/y' = z/z' = u/u' = \cdots$.
In all proportions, we must require that the denominators are not zero, because
division by zero is impossible.
\end{definition}

You may recall that in algebra you learned a number of theorems about
proportions, some of which are troublesome to remember. All statements about
proportions are simple consequences of the basic laws of algebra. The theorems on
proportions in this chapter are not to be memorized, but rather their proofs are
to be understood.

\exercisegroup{9--4}
\begin{multicols}{2}
\begin{exlist}
\item Give other ratios equal to the following:\\
      (a) $9/12$, (b) $2x^2/3x$, (c) $12abc/27c^2$, (d) $(x^2-x-6)/(4x^2+8x)$,
      (e) $(x^2-5x+6)/(x^2-4)$.
\item Are the following proportions true?\\
      (a) $\dfrac{a}{b} = \dfrac{a^2}{ab}$,\quad (b) $\dfrac{x}{y} =
      \dfrac{x^2}{y^2}$,\quad (c) $\dfrac{3}{4} = \dfrac{9}{16}$,\quad
      (d) $x = \dfrac{6x}{6}$.
\item Any number can be regarded as a ratio because $N = N/1$. Hence, can $x + 2$
      be regarded as a ratio for every number $x$?
\item Determine $x$, so that the following proportions are true:\\
      (a) $\dfrac{2x}{8} = \dfrac{5}{3}$,\quad (b) $\dfrac{(3-x)}{x} =
      \dfrac{5}{3}$,\quad (c) $\dfrac{2x+3}{5} = \dfrac{x-3}{2}$,\quad
      (d) $\dfrac{ax}{5} = \dfrac{a}{10}$,\quad
      \stex(e) $\dfrac{2x-7}{2} = \dfrac{3x-5}{3}$,\quad
      (f) $\dfrac{9}{x} = \dfrac{24}{8}$,\quad (g) $\dfrac{3}{x} =
      \dfrac{x}{12}$,\quad (h) $\dfrac{a}{x} = \dfrac{x}{b}$.
\item Determine the ratio of $u$ to $v$ if (a) $4u = 5v$, (b) $v = 3u$,
      (c) $au = bv$, (d) $u = 3v$, (e) $4u - 18v = 0$,
      (f) $(2u - v)/(u + v) = \frac{1}{2}$.
\item If $x$ is to 5 as 5 is to $x$, find $x$.
\end{exlist}
\end{multicols}

We list below some theorems about proportions that are simple consequences of the
laws of algebra. Since a ratio is a number obtained by division, we first recall
the meaning of division.

\begin{definition}
If $b \neq 0$, then $a/b$ (or $a \div b$) is the number $q$, such that $a = qb$.
In other words, for $b \neq 0$, $a/b = q$ if and only if $a = qb$.
\end{definition}

\begin{theorem}
If $b \neq 0$, and $c \neq 0$, then $a/b = (ac/bc)$.
\end{theorem}

\noindent\emph{Proof.} If $q = a/b$, then $a = qb$. Therefore, $ac = q \cdot bc$.
Whence $q = (ac/bc)$.

\begin{theorem}
If neither $b$ nor $d = 0$, $a/b = c/d$ if and only if $ad = bc$.
\end{theorem}

\noindent\emph{Proof.} The statement $a/b = c/d$ says that two numbers are equal,
that is, they are the same number. Multiplying this number by $bd$, we have
$bd \cdot a/b = bd \cdot c/d$, or $da = bc$. Conversely, if $ad = bc$ and
$bd \neq 0$, then dividing this number by $bd$, we have $(ad/bd) = (bc/bd)$ or
$a/b = c/d$.

\begin{theorem}
If $c \neq 0$, then $a/b = c/d$ if and only if $a/c = b/d$.
\end{theorem}

\noindent\emph{Proof.} Both proportions are equivalent to $ad = bc$, by
Theorem 9--3.

\exercisegroup{9--5}
\begin{multicols}{2}
\begin{exlist}
\item If $dx = ey$, what is the ratio of $x$ to $y$?
\item If $3u = 7v$, what is the ratio of $u$ to $v$? of $v$ to $u$?
\item If $hx - ky = ry + ty$, what is the ratio of $x$ to $y$?
\item[\stex 4.] Prove that if $x/y = u/v$, $x = u$, and $x \neq 0$, then $v = y$.
\item[\stex 5.] Prove that if $a/b = c/d$, and $b + d \neq 0$, then
      $a/b = (a+c)/(b+d)$.
\item[\stex 6.] Prove that if $a/b = c/d$, then $(a+b)/b = (c+d)/d$.
\item[7.] Use Exercise 6 to form a new proportion from $x/3 = z/7$.
\item[8.] Use Exercise 5 to form a new proportion from $6/15 = 10/25$.
\item[\stex 9.] State and prove the converse of Exercise 5.
\item[\stex 10.] State and prove the converse of Exercise 6.
\end{exlist}
\end{multicols}

%========================================================= 9-5
\section*{9--5\quad RATIO AND PROPORTION IN GEOMETRY}
\addcontentsline{toc}{section}{9--5\quad Ratio and Proportion in Geometry}

The objects of study in geometry are sets of points. Some of these objects have
numbers attached to them; for example, a segment has a \emph{length}. Later, we
also discuss area of polygons and measurement of angles. In geometry, when we
speak of the ratio of two segments, we mean the \emph{ratio of the numbers} that
are the lengths of the segments. It will be much the same when we speak of the
ratio of any other geometric objects.

\begin{definition}
If $AB$ and $CD$ are segments, the ratio of $AB$ to $CD$ is the number
$\sg{AB}/\sg{CD}$. The statement that the segments $AB$ and $CD$ are
\emph{proportional} to the segments $A'B'$ and $C'D'$ means that
$\sg{AB}/\sg{A'B'} = \sg{CD}/\sg{C'D'}$. We say that the sides of $\tri{ABC}$ are
proportional to the sides of $\tri{A'B'C'}$ if $\sg{AB}/\sg{A'B'} =
\sg{BC}/\sg{B'C'} = \sg{CA}/\sg{C'A'}$.
\end{definition}

\exercisegroup{9--6}
\begin{multicols}{2}
\begin{exlist}
\item Draw two segments. By measurement, determine the ratio of one segment to
      the other.
\item In Fig.~9--15, the sides of $\tri{EFG}$ are proportional to the sides of
      $\tri{E'F'G'}$. What is the ratio of $EG$ to $E'G'$, if $\sg{E'F'} = 2$,
      and the lengths of the sides of $\tri{EFG}$ are as shown? What is the ratio
      of $F'G'$ to $E'F'$?
\item If $\sg{AB} = 5$, $\sg{CD} = 7$, and $\sg{EF} = 3$, what is the ratio of
      $AB$ to $CD$? of $CD$ to $EF$? of $AB$ to $FE$?
\item The sides of $\tri{ABC}$ are proportional to the sides of $\tri{A'B'C'}$
      and $\sg{AB} = 5$, $\sg{BC} = 6$, $\sg{CA} = 8$. If $\sg{C'A'} = 9$, find
      $\sg{A'B'}$.
\item If the sides of $\tri{EFG}$ are proportional to the sides of $\tri{PQR}$
      and $\sg{EF} = 7$, $\sg{FG} = 8$, $\sg{QR} = 6$, $\sg{RP} = 8$, find
      $\sg{GE}$ and $\sg{PQ}$.
\end{exlist}
\end{multicols}

\FIGIXFIFTEEN

Before proceeding, we should observe that a question might be raised concerning
Definition 9--3. The number that we get for the length of a segment depends on
what we choose for a segment of unit length. Thus a foot ruler has length 1 if
the unit length is one foot, but has length 12 if the unit is one inch. We now
justify Definition 9--3 by proving that the ratio does not depend on the unit of
length.

\begin{theorem}
No matter what unit of length is chosen for measuring two segments, the ratio of
these segments is the same number.
\end{theorem}

\noindent\emph{Proof.} Suppose that the two segments are $AB$ and $A'B'$, and
suppose that we have two different unit segments $CD$ and $EF$, as in Fig.~9--16.

\FIGIXSIXTEEN

\noindent Suppose the lengths of $AB$ and $A'B'$, when $CD$ is the unit, are $l$
and $l'$ respectively, and their lengths, when $EF$ is the unit, are $L$ and $L'$
respectively. We have seen, in Theorem 5--4, that changing the unit of measure
just multiplies the length by a number. That is, if the length of $CD$, when $EF$
is the unit, is $x$, then $L = l \cdot x$, and $L' = l' \cdot x$. Hence, if we
form the ratio of $AB$ to $A'B'$ when $CD$ is the unit, we get $l/l'$; and if we
form the ratio when $EF$ is the unit, we get
\[
  \frac{L}{L'} \;=\; \frac{l \cdot x}{l' \cdot x} \;=\; \frac{l}{l'}.
\]
Hence, the two ratios are the same, and the theorem is proved.

\exercisegroup{9--7}
\begin{multicols}{2}
\begin{exlist}
\item Write the ratios above when $CD$ is twice as long as $EF$; three times as
      long.
\item One inch is approximately 2.54 centimeters. What is the length of a foot
      rule in centimeters?
\item Two sides of a triangle have lengths 12 in. and 16 in. What would their
      ratio be if the sides were measured in miles?
\item If the ratio of segment $AB$ to segment $CD$ is 0.75, what is the ratio of
      $CD$ to $AB$?
\item Is the ratio of one segment to another ever the number zero?
\end{exlist}
\end{multicols}

We have defined a ratio to be a number, and numbers can be either \emph{rational}
or \emph{irrational}. When we encounter these two kinds of ratios in geometry, we
use the words \emph{commensurable} and \emph{incommensurable}. We inherit this
language from the Greeks, who discovered irrational numbers and were much
disturbed by their discovery. We make the language clear in the following
definition.

\begin{definition}
One segment is said to be \emph{commensurable} with another if the ratio of the
first to the second is a rational number. They are said to be
\emph{incommensurable} if this ratio is an irrational number.
\end{definition}

\begin{theorem}
Two segments, $AB$ and $CD$, are commensurable with each other if and only if
there is a segment $EF$ which, if used as a unit of measure, makes $\sg{AB}$ and
$\sg{CD}$ whole numbers (integers).
\end{theorem}

\begin{remark}
Such a segment $EF$ is a common unit of measure for the two segments, hence the
word \emph{commensurable}.
\end{remark}

\noindent\emph{Proof.} We have to show that if such a segment $EF$ exists, then
$\sg{AB}/\sg{CD}$ is a rational number, and conversely, that if
$\sg{AB}/\sg{CD}$ is a rational number, then such a segment $EF$ exists.

Suppose first that such a segment $EF$ exists. Then $\sg{AB} = m$, and
$\sg{CD} = n$, where $m$ and $n$ are integers. Therefore $\sg{AB}/\sg{CD} = m/n$,
and is a rational number.

To prove the converse, suppose that $\sg{AB}/\sg{CD}$ is a rational number, that
is, $\sg{AB}/\sg{CD} = m/n$, where $m$ and $n$ are integers. Divide the segment
$CD$ into $n$ equal subintervals, and call one of these subintervals $EF$
(Fig.~9--17). With $EF$ as a unit of measure, $\sg{CD} = n$, and $\sg{AB} = m$.

\FIGIXSEVENTEEN

\begin{note}
In this book, our theorems deal only with commensurable pairs of segments, and we
assume their truth for incommensurable cases. Proofs for incommensurable cases
require some knowledge of \emph{limits}. The ideas involved are difficult and we
avoid limits whenever possible.
\end{note}

\begin{histremark}
Pythagoras (6th century B.C.) believed that all quantitative relations could be
described by pure numbers. By this, he meant whole numbers, or integers. He was
(so the story goes) profoundly shocked when he discovered that the side and
diagonal of a square are incommensurable. Because of this discovery, later Greeks
(notably Eudoxus and Euclid) invented an elaborate theory of ratio and
proportion that is a model of ingenuity and precision. The interesting fact
about this theory, from our point of view, is
that \emph{there were no numbers in Greek geometry, that is, no measurement}, and
hence the meaning of the words ``two ratios are equal'' had to be defined in a
special way that we cannot consider here.
\end{histremark}


%========================================================= 9-6
\section*{9--6\quad SIMILAR TRIANGLES AND POLYGONS}
\addcontentsline{toc}{section}{9--6\quad Similar Triangles and Polygons}

In this section we are primarily concerned with similar triangles, but we also
define similarity for polygons, and solve some problems concerning pairs of
similar polygons. Remember that we are trying to make precise what we mean when
we say that two polygons \emph{have the same shape}.

\begin{definition}
Two \emph{triangles} are said to be \emph{similar} to each other if corresponding
angles are congruent and corresponding sides are proportional. For example, in
Fig.~9--18, $\tri{ABC}$ is similar to $\tri{A'B'C'}$ if $\ang{A} \cong \ang{A'}$,
$\ang{B} \cong \ang{B'}$, $\ang{C} \cong \ang{C'}$, and $\sg{AB}/\sg{A'B'} =
\sg{BC}/\sg{B'C'} = \sg{CA}/\sg{C'A'}$. To indicate that triangles are similar,
we write $\tri{ABC} \simto \tri{A'B'C'}$.
\end{definition}

Two \emph{convex} polygons are said to be \emph{similar} to each other if the
corresponding angles are congruent and the corresponding sides are proportional.
For example, the polygons in Fig.~9--19 are similar if $\ang{A} \cong \ang{A'}$,
$\ang{B} \cong \ang{B'}$, $\ang{C} \cong \ang{C'}$, $\ang{D} \cong \ang{D'}$,
$\ang{E} \cong \ang{E'}$, and $\sg{AB}/\sg{A'B'} = \sg{BC}/\sg{B'C'} =
\sg{CD}/\sg{C'D'} = \sg{DE}/\sg{D'E'} = \sg{EA}/\sg{E'A'}$.

\FIGIXEIGHTEENNINETEEN

\exercisegroup{9--8}
\begin{multicols}{2}
\begin{exlist}
\item Two triangles are similar, and two sides of one triangle are 7 and 11 in.
      In the other triangle, the side corresponding to the 7-in.\ side has length
      12 in. Find the length of the side corresponding to the 11-in.\ side.
\item Two triangles are similar, and the sides of one triangle have lengths 4, 6,
      and 8. In the other triangle, the shortest side has length 6. Find the
      lengths of the other sides.
\item If $\tri{ABC} \simto \tri{DEF}$ and $\sg{AC} = 21$, $\sg{DF} = 15$,
      $\sg{FE} = 18$, and $\sg{DE} = 21$, find the lengths of the other sides of
      $\tri{ABC}$.
\item Two quadrilaterals are similar, and one quadrilateral has sides of 5, 7,
      12, and 16 in. The other quadrilateral has its longest side 12 ft. How long
      are the other sides?
\item The sides of a triangle are 4, 5, and $x$ in. The corresponding sides of a
      triangle similar to the first triangle are 6, $y$, and 8 in. What are the
      sides of each triangle?
\item A pentagon has sides 4, 6, 10, 15, and $x$ ft. In a pentagon similar to the
      first pentagon, the corresponding sides are $u$, $v$, 8, $w$, and 16. Find
      $u$, $v$, $w$, and $x$.
\item Two right triangles are similar, and in the first, the ratio of the shorter
      to the longer leg is $\frac{2}{3}$. In the other triangle, the shorter leg
      is 13 in.\ long. How long is the longer leg of this triangle?
\item Find as many examples of similar polygons from the everyday world as you
      can.
\item Prove that if $\tri{ABC} \simto \tri{DEF}$ and $\tri{DEF} \simto
      \tri{PQR}$, then $\tri{ABC} \simto \tri{PQR}$.
\item Are all squares similar? Why?
\item Are all rectangles similar? Why?
\item State conditions under which two parallelograms are similar.
\item Prove that congruent polygons are similar polygons. Is the converse true?
\item Prove that all equilateral triangles are similar.
\item In Fig.~9--20, $E$ and $F$ are midpoints of $AB$ and $BC$. Prove that
      $\tri{BEF} \simto \tri{BAC}$.
\item If, in Fig.~9--21, $A'$, $B'$, $C'$, and $D'$ bisect the corresponding
      segments from $O$, prove that the quadrilaterals $ABCD$ and $A'B'C'D'$ are
      similar.
\end{exlist}
\end{multicols}

\FIGIXTWENTYTWENTYONE

%========================================================= 9-7
\section*{9--7\quad THEOREMS ON SIMILAR TRIANGLES}
\addcontentsline{toc}{section}{9--7\quad Theorems on Similar Triangles}

In this section we give the fundamental theorems on similar triangles. They are
extremely important because they have wide application in mathematics, other
sciences, and engineering.

\begin{theorem}
If the angles of one triangle are congruent, respectively, to the angles of
another triangle, then the triangles are similar.
\end{theorem}

\noindent\emph{Proof.} Suppose that the triangles are $\tri{ABC}$ and
$\tri{A'B'C'}$, as in Fig.~9--22, and that $\ang{A} \cong \ang{A'}$,
$\ang{B} \cong \ang{B'}$, and $\ang{C} \cong \ang{C'}$. We must prove that the
sides are proportional, that is, $\sg{AB}/\sg{A'B'} = \sg{BC}/\sg{B'C'} =
\sg{CA}/\sg{C'A'}$.

\FIGIXTWENTYTWO

In forming the proof, we will suppose that the segments $AB$ and $A'B'$ are
commensurable. In other words, $\sg{AB}/\sg{A'B'}$ is a rational number. The
proof for the incommensurable case is difficult, since it involves knowing about
\emph{limits}. We therefore suppose that, with a suitable unit, the lengths of
$AB$ and $A'B'$ are positive integers. The figure is drawn with $\sg{AB} = 5$ and
$\sg{A'B'} = 3$. We may suppose that $AB > A'B'$, for if $AB < A'B'$, the
argument can be reversed, and if $AB \cong A'B'$, the triangles are congruent
(why?), and hence similar.

\begin{steps}
\step{There is a point $A''$ on $AB$, such that $\sg{A''B} \cong \sg{A'B'}$, and
      $\sg{A''B} = \sg{A'B'} = 3$.}{Congruence postulate.}
\step{$A''$ is between $A$ and $B$.}{We have assumed that $AB > A'B'$.}
\step{There is a line through $A''$ parallel to $AC$ which intersects $BC$ in
      $C''$.}{Why?}
\step{$\tri{BA''C''} \cong \tri{B'A'C'}$.}{A.S.A.}
\step{$\sg{BC''} \cong \sg{B'C'}$.}{Statement 4.}
\step{There are points $D$, $E$, and $F$, on $AB$, such that the points $A$, $F$,
      $A''$, $E$, $D$, and $B$, in that order, are at unit distance from their
      neighbors.}{It is given that $\sg{AB} = 5$. By Statement 1,
      $\sg{A''B} = \sg{A'B'} = 3$.}
\step{There are lines through $D$, $E$, and $F$, parallel to $AC$. These lines
      intersect $BC$ in points $I$, $H$, and $G$, respectively.}{Why?}
\stepx{8}{$\sg{CG} \cong \sg{GC''} \cong \sg{C''H} \cong \sg{HI} \cong
      \sg{IB}$.}{Theorem 9--1 on transversals.}
\stepx{9}{If $CG$ is chosen as unit of length, $\sg{BC} = 5$, and $\sg{BC''} =
      \sg{B'C'} = 3$.}{Definition of length and Statement 5.}
\stepx{10}{$\sg{BC}/\sg{B'C'} = 5/3$.}{Statement 9.}
\stepx{11}{$\sg{AB}/\sg{A'B'} = 5/3$.}{Given.}
\stepx{12}{$\sg{AB}/\sg{A'B'} = \sg{BC}/\sg{B'C'}$.}{Statements 10 and 11.}
\stepx{13}{In an analogous way, by considering a line parallel to $BC$, it may be
      shown that $\sg{AB}/\sg{A'B'} = \sg{AC}/\sg{A'C'}$.}{Same reasons as in
      Statements 1 through 12.}
\stepx{14}{$\tri{ABC} \simto \tri{A'B'C'}$.}{Statements 12 and 13.}
\end{steps}

\begin{remark}
It should be clear that the proof is essentially the same if $AB$ and $A'B'$ are
any other positive integers besides 5 and 3.
\end{remark}

\begin{corollary}
If two angles of one triangle are congruent to two angles of another triangle,
then the triangles are similar.
\end{corollary}

\begin{corollary}
If an acute angle of one right triangle is congruent to an acute angle of another
right triangle, then the triangles are similar.
\end{corollary}

The proofs are left to the student.

\exercisegroup{9--9}
\begin{multicols}{2}
\begin{exlist}
\item If, in Fig.~9--23, $BE \pll CD$, $\sg{AE} = 4$, and $\sg{ED} = 6$, what is
      $\sg{AB}/\sg{AC}$?
\item A church steeple casts a shadow 100 ft long, and a 6-ft post casts, at the
      same time, a shadow 7 ft long. How high is the steeple?
\item A line from the top of a cliff to the ground just passes over the top of a
      pole 8 ft high and meets the ground at a point 5 ft from the base of the
      pole. If it is 60 ft from this point to the base of the cliff, how high is
      the cliff?
\item As in Fig.~9--24, a lake lies between $A$ and $C$. If lines $AB$ and $DE$
      run north-south, $\sg{AB} = 8$ mi, $\sg{DE} = 1.2$ mi, and $\sg{DC} = 1.5$
      mi, how far is it from $A$ to $C$?
\item In Fig.~9--25, $DE \pll AC$. Prove that $\tri{ABC} \simto \tri{DBE}$.
\item If in $\tri{DEF}$, $G$ and $H$ are points on sides $EF$ and $DF$,
      respectively, and $\ang{GHF} \cong \ang{EDF}$, prove that $\tri{DEF} \simto
      \tri{HGF}$.
\item In Fig.~9--26, $AB \perp AE$, and $DE \perp AE$. Prove that (a)
      $\sg{AB}/\sg{DE} = \sg{AC}/\sg{CE}$; (b) $\sg{AB}/\sg{BC} =
      \sg{DE}/\sg{DC}$.
\item In $\tri{ABC}$, let $E$, $F$, and $G$ be the midpoints of the sides $AB$,
      $BC$, and $CA$, respectively. Prove that $\tri{ABC} \simto \tri{FGE}$.
\item If $\tri{ABC} \simto \tri{A'B'C'}$, and $BD$ and $B'D'$ are the altitudes
      from $B$ and $B'$ respectively, prove that $\sg{AC}/\sg{A'C'} =
      \sg{BD}/\sg{B'D'}$.
\item In Fig.~9--27, $l \pll m$. Prove that $\tri{EFG} \simto \tri{IHG}$.
\item If, in Fig.~9--28, $\tri{PQR} \simto \tri{P'Q'R'}$, and $PS$ and $P'S'$
      bisect angles $P$ and $P'$, prove that $\sg{PQ}/\sg{PS} =
      \sg{P'Q'}/\sg{P'S'}$.
\item In $\tri{ABC}$, $CD$ and $AE$ are altitudes, as in Fig.~9--29. Prove that
      (a) $\tri{BEA} \simto \tri{BDC}$; (b) $\tri{ADF} \simto \tri{CEF}$.
\end{exlist}
\end{multicols}

\FIGIXTWENTYTHREETWENTYNINE

Before proving the next of the fundamental theorems on similar triangles, we
prove a theorem that we need later and which is interesting in itself.

\begin{theorem}
A line that divides two sides of a triangle proportionally is parallel to the
third side. This means that if $\sg{AD}/\sg{AB} = \sg{AE}/\sg{AC}$ in
Fig.~9--30, then $DE \pll BC$.
\end{theorem}

\FIGIXTHIRTY

\noindent\emph{Proof.}
\begin{steps}
\step{Let $DF$ be a line parallel to $BC$ and meeting $AC$ in a point $F$.}{Why
      is there such a line?}
\step{The point $F$ is on the segment $AC$.}{Why?}
\step{$\tri{ABC} \simto \tri{ADF}$.}{Explain how Theorem 9--7 is used here.}
\step{$\sg{AF}/\sg{AC} = \sg{AD}/\sg{AB}$.}{From the definition of similar
      triangles.}
\step{$\sg{AD}/\sg{AB} = \sg{AE}/\sg{AC}$.}{Given.}
\step{$\sg{AF} = \sg{AE}$.}{Statements 4 and 5.}
\step{$F = E$.}{$AF \cong AE$ and congruence postulates.}
\step{$DE \pll BC$.}{Statements 7 and 1.}
\end{steps}

\begin{theorem}
Given $\tri{ABC}$ and $\tri{A'B'C'}$ with $\ang{A} \cong \ang{A'}$ and
$\sg{AB}/\sg{A'B'} = \sg{AC}/\sg{A'C'}$, then $\tri{ABC} \simto \tri{A'B'C'}$.
\end{theorem}

\noindent\emph{Proof.} We show that $\ang{B} \cong \ang{B'}$, and use
Corollary 9--7--1. We may suppose that $AB > A'B'$. (Why?)

\FIGIXTHIRTYONE

\begin{steps}
\step{In Fig.~9--31, there are points $B''$ on $AB$ and $C''$ on $AC$ such that
      $\sg{A'B'} \cong \sg{AB''}$ and $\sg{A'C'} \cong \sg{AC''}$.}{Why?}
\step{$\tri{AB''C''} \cong \tri{A'B'C'}$.}{Why?}
\step{$\sg{AB}/\sg{AB''} = \sg{AC}/\sg{AC''}$.}{Why?}
\step{$B''C'' \pll BC$.}{Theorem 9--8.}
\step{$\ang{AB''C''} \cong \ang{B}$.}{Why?}
\step{$\ang{B} \cong \ang{B'}$.}{Statements 2 and 5.}
\step{$\tri{ABC} \simto \tri{A'B'C'}$.}{Why?}
\end{steps}

\exercisegroup{9--10}
\begin{multicols}{2}
\begin{exlist}
\item Prove that if the legs of one right triangle are proportional to the legs
      of another right triangle, then the triangles are similar.
\item In Fig.~9--32, $\sg{AO} = 10$, $\sg{BO} = 12$, $\sg{OD} = 5$, and
      $\sg{OC} = 6$. Show that $\tri{AOB} \simto \tri{DOC}$.
\item Prove that two isosceles triangles are similar if their vertex angles are
      congruent; likewise if their base angles are all congruent to one another.
\item In Fig.~9--33, if $\sg{AD} = x \cdot \sg{AB}$, and $\sg{AE} = x \cdot
      \sg{AC}$, prove that $\tri{ABC} \simto \tri{ADE}$.
\item If $\tri{ABC}$ is a right triangle, with the right angle at $C$, and if, in
      Fig.~9--34, $\sg{AD}/\sg{AC} = \sg{AC}/\sg{AB}$, prove that $CD \perp AB$.
\item[\stex 6.] In Fig.~9--35, if $\tri{ABC}$ is a right triangle, with right
      angle at $C$, and if $CD \perp AB$, prove that $\tri{ABC} \simto
      \tri{ACD}$.
\end{exlist}
\end{multicols}

\FIGIXTHIRTYTWOTHIRTYFIVE

The next theorem is the last of the three basic theorems on similar triangles.

\begin{theorem}
If the sides of one triangle are proportional to the sides of another triangle,
then the triangles are similar. In other words, referring to Fig.~9--36, if
$\sg{AB}/\sg{A'B'} = \sg{AC}/\sg{A'C'} = \sg{BC}/\sg{B'C'}$, then
$\tri{ABC} \simto \tri{A'B'C'}$.
\end{theorem}

\noindent\emph{Proof.} It is sufficient to show that $\ang{B} \cong \ang{B'}$.
Why? We may assume that $AB > A'B'$, for if $AB < A'B'$, the argument can be
reversed, and if $\sg{AB} = \sg{A'B'}$, then the triangles are congruent (why?),
and hence similar.

\FIGIXTHIRTYSIX

\begin{steps}
\step{Choose $A''$ on $AB$ and $C''$ on $BC$, such that $\sg{A''B} = \sg{A'B'}$
      and $\sg{C''B} = \sg{C'B'}$.}{Why possible?}
\step{$\sg{AB}/\sg{A''B} = \sg{BC}/\sg{BC''}$.}{The hypothesis of Theorem 9--10
      and Statement 1.}
\step{$\tri{ABC} \simto \tri{A''BC''}$.}{Theorem 9--9.}
\step{$\sg{AC}/\sg{A''C''} = \sg{AB}/\sg{A''B} = \sg{AB}/\sg{A'B'}$.}{Statements
      3 and 1.}
\step{$\sg{AC}/\sg{A''C''} = \sg{AC}/\sg{A'C'}$.}{Statement 4 and the given fact
      $\sg{AB}/\sg{A'B'} = \sg{AC}/\sg{A'C'}$.}
\step{$\sg{A''C''} = \sg{A'C'}$.}{Statement 5.}
\step{$\tri{A''BC''} \cong \tri{A'B'C'}$.}{Statements 1 and 6 and S.S.S.}
\step{$\ang{B} \cong \ang{B'}$.}{Statement 7.}
\step{$\tri{ABC} \simto \tri{A'B'C'}$.}{Theorem 9--9.}
\end{steps}

\exercisegroup{9--11}
\begin{multicols}{2}
\begin{exlist}
\item State, in your own words, the three basic theorems on similar triangles,
      Theorems 9--7, 9--9, and 9--10.
\item In Fig.~9--37, the numbers are the lengths of the segments. How are the
      lines $AB$ and $CD$ related? Why?
\item If the hexagons of Fig.~9--38 are similar, prove that (a) $\tri{BCD} \simto
      \tri{B'C'D'}$; (b) quadrilateral $ABCD \simto$ quadrilateral $A'B'C'D'$.
\item In Fig.~9--39, $BE \perp AC$, and $CD \perp AB$. Find two similar
      triangles. How many proportions can you write?
\item If $\tri{ABC} \simto \tri{A'B'C'}$, and $AD$ and $A'D'$ are medians, prove
      that the ratio of the medians is equal to the ratio of any two
      corresponding sides. Restate this using the word \emph{proportion} (or
      \emph{proportional}).
\item Prove that if two triangles are similar, then the ratio of their perimeters
      is the same as the ratio of two corresponding sides.
\item Is a statement like that of Exercise 6 true for similar polygons? How would
      you prove this?
\end{exlist}
\end{multicols}

\FIGIXTHIRTYSEVENTHIRTYNINE

%========================================================= 9-8
\section*{9--8\quad THE PYTHAGOREAN THEOREM}
\addcontentsline{toc}{section}{9--8\quad The Pythagorean Theorem}

With our knowledge of similar triangles, we are able to prove a theorem
associated with the Greek mathematician Pythagoras, who lived in the 6th century
B.C. It is believed that he first gave a proof of the theorem that bears his
name. First, we state a theorem that is useful for the proof of the Pythagorean
theorem.

\begin{theorem}
The altitude to the hypotenuse of a right triangle forms two right triangles,
which are similar to the given triangle. For example, in Fig.~9--40,
$\tri{ABC} \simto \tri{ACD} \simto \tri{CBD}$.
\end{theorem}

\FIGIXFORTYFORTYONE

\noindent\emph{Proof.} In the first place, a line through $C$ perpendicular to
$AB$ must intersect the segment $AB$. This fact may be seen as follows: Consider
lines $l$, $l'$, and $l''$ in Fig.~9--41, through $C$, $A$, and $B$, and
perpendicular to $AB$. Then $l' \pll l \pll l''$. The point $C$ is between $l'$
and $l''$, because angles $A$ and $B$ are acute angles. Therefore, $l$ is between
$l'$ and $l''$, and hence $l$ intersects $AB$ in a point $D$. $\tri{ACD} \simto
\tri{ABC}$, since each triangle has a right angle, and the triangles have the
acute angle $A$ in common. By the same reasoning, $\tri{BCD} \simto \tri{BAC}$.

\begin{corollary}
The altitude to the hypotenuse is a mean proportional\textsuperscript{*}
between the two segments formed on the hypotenuse.
\end{corollary}
\IXstarnote{*$x$ is a ``mean proportional'' between $a$ and $b$ if
$a/x = x/b$, or, which is the same, $x^2 = ab$.}

\noindent\emph{Proof.} We have to show that, using Fig.~9--40,
$\sg{AD}/\sg{CD} = \sg{CD}/\sg{DB}$. Since $\tri{ADC} \simto \tri{CDB}$, this is
an immediate consequence of the definition of similar triangles.

\begin{corollary}
A leg of a right triangle is the mean proportional between the hypotenuse and the
adjacent segment on the hypotenuse formed by the altitude from the right angle.
That is, in Fig.~9--40, $\sg{AC}^2 = \sg{AD} \cdot \sg{AB}$.
\end{corollary}

\exercisegroup{9--12}
\begin{multicols}{2}
\begin{exlist}
\item In Fig.~9--42, $EF \perp FG$, and $FH \perp EG$. If $\sg{EF} = 8$, and
      $\sg{FG} = 5$, what is $\sg{FH}/\sg{EH}$? $\sg{FH}/\sg{HG}$?
\item In Fig.~9--43, the sides of a right triangle are 5, 12, and 13. How long is
      the altitude on the hypotenuse?
\item In Fig.~9--44, $Q$ is a right angle, and $QS \perp PR$. Find $\sg{PQ}$, if
      $\sg{PS} = 4$, and $\sg{SR} = 5$.
\item Using Fig.~9--44, if $\sg{QS} = 7$, what is the product
      $\sg{PS} \cdot \sg{SR}$?
\item In Fig.~9--45, $\ang{C}$ is a right angle, as is $\ang{ADC}$. Express $h$
      in terms of lengths $a$, $b$, and $c$ of the sides. Obtain $\sg{AD}$ in
      terms of $a$, $b$, and $c$.
\end{exlist}
\end{multicols}

\FIGIXFORTYTWOFORTYFIVE

\begin{theorem}[Pythagorean theorem]
The square of the length of the hypotenuse of a right triangle is equal to the
sum of the squares of the lengths of the legs.
\end{theorem}

If the triangle is $\tri{ABC}$, with the right angle at $C$, and if the lengths
of the sides opposite angles $A$, $B$, and $C$ are $a$, $b$, and $c$,
respectively, as in Fig.~9--46, then we are able to show that $c^2 = a^2 + b^2$.

\FIGIXFORTYSIX

\noindent\emph{Proof.}
\begin{steps}
\step{Let $CD$ be the altitude and $x = \sg{AD}$ and $y = \sg{DB}$. Then
      $x/b = b/c$ and $y/a = a/c$.}{Corollary 9--11--2.}
\step{$x + y = c$.}{Properties of length.}
\step{$b^2/c + a^2/c = c$.}{Statements 1 and 2.}
\step{$a^2 + b^2 = c^2$.}{Statement 3 and algebra.}
\end{steps}

\begin{theorem}[Converse to Pythagorean theorem]
If a triangle has sides of length $a$, $b$, and $c$, such that $c^2 = a^2 + b^2$,
then the triangle is a right triangle, and the side of length $c$ is the
hypotenuse.
\end{theorem}

\noindent\emph{Proof.} The method of proof is to show that there is a right
triangle that is congruent to the given triangle. Suppose that the given triangle
is $\tri{ABC}$.

\FIGIXFORTYSEVEN

\begin{steps}
\step{There is a segment $PQ \cong BC$, as in Fig.~9--47.}{Congruence
      postulates.}
\step{There is a ray perpendicular to $PQ$ at $P$.}{From the congruence
      postulates, there is an angle at $P$ congruent to a right angle.}
\step{There is a point $R$ on this ray such that $PR \cong AC$.}{Congruence
      postulates.}
\step{$\sg{PQ} = a$, and $\sg{PR} = b$.}{Statements 1 and 3.}
\step{$\sg{RQ}^2 = a^2 + b^2$.}{Pythagorean theorem.}
\step{$c^2 = a^2 + b^2$.}{Given.}
\step{$\sg{RQ}^2 = c^2$, so $\sg{RQ} = c$.}{Statements 5 and 6 and algebra.}
\step{$RP \cong AC$, $PQ \cong CB$, and $RQ \cong AB$.}{By construction and
      Statement 7.}
\step{$\tri{ABC} \cong \tri{RQP}$.}{S.S.S.}
\step{$\ang{C}$ is a right angle.}{$\ang{P}$ is a right angle, by construction,
      and $\ang{C} \cong \ang{P}$, by Statement 9.}
\end{steps}

\exercisegroup{9--13}
\begin{multicols}{2}
\begin{exlist}
\item It is said that the ancient Egyptians constructed right angles for
      surveying purposes by forming a triangle using ropes of length 3, 4, and 5
      units. Why would this be correct? Can you think of other integers for the
      sides of a right triangle?
\item In a right triangle, the hypotenuse is 7 and one leg is 3. How long is the
      other leg?
\item Write a formula for the length of one leg of a right triangle in terms of
      the lengths of the hypotenuse and other leg.
\item[\stex 4.]\stepcounter{exlisti} In Fig.~9--48, $AB \pll CD$, $BC \pll AD$, and $AC \perp BD$. If
      $\sg{BD} = 24$, and $\sg{AD} = 13$, find $\sg{AC}$.
\item[5.]\stepcounter{exlisti} In a right triangle, the two legs are 6 ft and 9 ft. How long is the
      hypotenuse?
\item The two legs of a right triangle are equal, and the hypotenuse is 4 in.\
      long. How long are the legs?
\item The base of an isosceles triangle is 12 in.\ long. The two congruent sides
      have a combined length of 13 in. How long is the altitude to the base?
\item Point $A$ in Fig.~9--49 is 21 mi due north of $O$, while point $B$ is 20 mi
      due east of $O$. What is the distance from $A$ to $B$, across the lake
      shown?
\item Find the altitude of an equilateral triangle, the length of whose side is
      $a$.
\item A rope 100 ft long reaches from the top of a pole to a point 40 ft from the
      foot of the pole. How tall is the pole?
\item Prove that the median to the hypotenuse of a right triangle is half as long
      as the hypotenuse.
\item Prove that in a $30^\circ$--$60^\circ$ right triangle, the shorter leg is
      half as long as the hypotenuse.
\item Prove that in an isosceles right triangle, the ratio of leg to hypotenuse
      is $1/\sqrt{2}$.
\item Construct a segment of length $\sqrt{5}$ in.
\item What is the length of a diagonal of a square of side 1? of side $a$?
\item Construct a segment of length $\sqrt{6}$.
\item[\stex 17.] If $n$ is any positive integer, can you construct (\RC) a
      segment of length $\sqrt{n}$?
\end{exlist}
\end{multicols}

\FIGIXFORTYEIGHTFORTYNINE

%========================================================= 9-9
\section*{9--9\quad ON CONSTRUCTIONS}
\addcontentsline{toc}{section}{9--9\quad On Constructions}

In our previous study of \RC\ constructions, we assumed that two circles with
equal radii would intersect if the distance between their centers were less than
the sum of their radii. We can easily prove this by using the Pythagorean
theorem. To do this, we first calculate where they \emph{should} intersect, and
then prove that this point \emph{is} on both circles.

\FIGIXFIFTYFIFTYONE

Consider Fig.~9--50, where $r$ is the radius of the circles, and $d = \sg{AB}$ is
the distance between their centers. If the circles intersect at $P$, and $C$ is
the midpoint of $AB$, then
$\tri{PAC} \cong \tri{PBC}$ (why?), and hence $\ang{C}$ is a right angle.
Therefore, $r^2 = \sg{PC}^2 + \sg{AC}^2$, so that
$\sg{PC} = \sqrt{r^2 - (\frac{1}{2}\sg{AB})^2}$. We now know ``where'' the point
of intersection $P$ must be, and we can prove that such a point of intersection
exists.

\begin{theorem}
Suppose that circles, with centers $A$ and $B$, have equal radii,
$r > \frac{1}{2}AB$. These circles intersect.
\end{theorem}

\noindent\emph{Proof.} Let $C$ be the midpoint of $AB$. Through $C$ there is a
line perpendicular to $AB$. On this line, let $P$ be a point at the distance
$\sg{PC} = \sqrt{r^2 - (\frac{1}{2}\sg{AB})^2}$ from $C$. (How do we know there
is such a point $P$?) Because $\tri{PAC}$ and $\tri{PBC}$ are right triangles, it
follows from the Pythagorean theorem that $\sg{PA} = \sg{PB} = r$. Hence $P$ lies
on both circles.

\begin{remark}
It also follows from the above proof that the two circles can intersect in two,
and only two, points, for the following reason. On the line through $C$,
perpendicular to $AB$, there are only two points at the correct distance from
$C$, namely one on each side of the line $AB$.
\end{remark}

\begin{theorem}
An angle inscribed in a semicircle is a right angle. That is, if in Fig.~9--51,
$AB$ is a diameter of the circle, with center at $O$, and $C$ is any point on the
circle other than $A$ and $B$, then $\ang{ACB}$ is a right angle.
\end{theorem}

\noindent\emph{Proof outline.} We drop a perpendicular from $C$ to $AB$, and
prove that the right triangles $ADC$ and $CDB$ are similar. It then follows that
$\ang{C}$ is a right angle.

\begin{steps}
\step{Let $CD$ be the perpendicular from $C$ to $AB$, and suppose that $D$ is
      between $A$ and $O$.}{The point $D$ must be between $A$ and $B$. (Why?)
      $D$, then, either is $O$, or is between $A$ and $O$, or $O$ and $B$. If
      $D = O$, the proof is easy. We discuss only the case in which $D$ is
      between $A$ and $O$. If $D$ is between $O$ and $B$, a similar proof can be
      made.}
\step{Let $x$, $y$, $z$, and $r$ be the lengths of the segments in Fig.~9--51.
      Then $x = r - y$, and $z = \sqrt{r^2 - y^2}$.}{Why?}
\step{$\dfrac{x}{z} = \dfrac{r-y}{\sqrt{r^2-y^2}} =
      \sqrt{\dfrac{r-y}{r+y}}$.}{Statement 2 and algebra.}
\step{$\dfrac{z}{r+y} = \dfrac{\sqrt{r^2-y^2}}{r+y} =
      \sqrt{\dfrac{r-y}{r+y}}$.}{Statement 2 and algebra.}
\step{$\tri{ACD}$ and $\tri{CDB}$ are right triangles.}{By construction.}
\step{$\tri{ACD} \simto \tri{CDB}$.}{Statements 3, 4, and 5 and Theorem 9--9.}
\step{$\ang{ACD} \cong \ang{CBD}$.}{Statement 6.}
\step{$\ang{ACB}$ is a right angle.}{Why?}
\end{steps}

\exercisegroup{9--14}
\begin{multicols}{2}
\begin{exlist}
\item A circle passes through the vertices of a right triangle whose sides are 3,
      4, and 5 in. What is the radius of the circle?
\item If $PQ$ is a diameter of a circle of radius 13, and $R$ is a point on the
      circle, such that $\sg{PR} = 10$, how long is $\sg{RQ}$?
\item In Fig.~9--52, $AB$ is a diameter of the circle, and $CD \perp AB$. If
      $\sg{AB} = 30$ and $\sg{EO} = 12$, find $\sg{CD}$.
\end{exlist}
\end{multicols}

\FIGIXFIFTYTWO

%========================================================= review
\section*{REVIEW OF CHAPTER 9}
\addcontentsline{toc}{section}{Review of Chapter 9}

Be sure you know the definitions given in this chapter of the following concepts:
transversal, ratio, proportion, ratio of segments, proportional sides,
commensurable, similar triangles.

Some of the important theorems of the chapter were concerned with these
definitions. Try to recall the theorems.

The main ideas of the chapter are listed briefly below. Be sure that you know
them, and review the proofs.

\begin{exlist}
\item If parallel lines cut off congruent segments on one transversal, they do so
      on every transversal.\\
      This theorem depended upon what facts about parallel lines and triangles?\\
      This theorem enables one to divide a segment into $m$ congruent parts with
      ruler and compass. Show how to do this.
\item What is a proportion? What is the ratio of $x$ to $y$ if $x$ and $y$ are
      numbers? if $x$ and $y$ are segments?\\
      If neither $b$ nor $d$ is zero, $a/b = c/d$ if and only if $ad = bc$.
      Illustrate this theorem by numerical examples.
\item The ratio of two segments does not depend on the unit of length.\\
      Illustrate this fact by getting the ratio of depth to width of your desk,
      first measuring in feet and then in inches.\\
      On what theorem about measurement of length does this fact depend?
\item[\stex 4.] Two segments are defined to be commensurable with each other if
      their ratio is a rational number. What is a rational number? Give an
      example of one. Give an example of an irrational number. What are
      incommensurable segments?\\
      State a theorem concerning commensurable segments.
\item[5.] What is the definition of similar triangles and polygons? Draw two
      similar triangles; pentagons.\\
      What do an architect's drawings have to do with similar figures?\\
      What does a map have to do with similarity?
\item[6.] Two triangles with corresponding angles congruent are similar.\\
      Illustrate the proof of this theorem, using two similar triangles $ABC$
      and $A'B'C'$, with $AB = 12$, and $A'B' = 8$.
\item[7.] If two triangles have their corresponding sides proportional, they
      are similar.\\
      What must be established to prove this theorem?
\item[8.] The Pythagorean theorem is proved, using a theorem about similar
      triangles in a right triangle. State this theorem.\\
      Prove the Pythagorean theorem.\\
      State its converse.
\end{exlist}

%--------------------------------------------------------- Review Exercises
\begin{center}\bfseries Review Exercises\end{center}
\addcontentsline{toc}{section}{Review Exercises (Chapter 9)}

\begin{exlist}
\item $\sg{AB} = 24$, and $C$ is a point on the line $AB$. If the ratio of $AC$
      to $CB = 3/5$, determine $\sg{AC}$ (a) if $C$ is between $A$ and $B$;
      (b) if $A$ is between $C$ and $B$.
\item Explain how to locate the point $C$ in Exercise 1(a) by an RC
      construction.
\item[\stex 3.] Can you give an RC construction for the point $C$ in Exercise
      1(b)?
\item[4.] A flagpole throws a 20-ft shadow when a yardstick throws a 10-in.
      shadow. Determine the height of the flagpole.
\item[5.] The length of the sides of a triangle are 4, 5, and 6. A triangle
      similar to this one has a perimeter of 30 in. Give the length of each side
      of this triangle.
\item[6.] In $\tri{ABC}$, the altitude from vertex $A$ divides the triangle into
      two similar triangles. $\tri{ABC}$ is not isosceles. What conclusion may
      be drawn?
\item[7.] In $\tri{ABC}$, $\ang{A} = 40^\circ$, and $\ang{B} = 80^\circ$. In
      $\tri{A'B'C'}$, $\ang{B'} = 40^\circ$, and $\ang{C'} = 60^\circ$. Are the
      two triangles similar?
\item[8.] What is the mean proportional of 4 and 9? of 2 and 4? of 6 and 6?
\item[9.] A triangle whose sides have lengths 3, 4, and 5 is a right triangle.
      So also is the triangle with sides of 5, 12, and 13; 7, 24, and 25; 9, 40,
      and 41; 11, 60, and 61; etc. Explain the pattern underlying the formation
      of these triples of numbers.
\item[10.] On a certain map, a distance of $\frac{5}{8}$ in. represents 80 mi.
      On this map, the points representing two cities are $5\frac{7}{16}$ in.
      apart. What is the airline distance between these cities?
\item[11.] The parallel sides of a trapezoid have lengths 18 and 12, as in
      Fig.~9--53. The altitude is 8, and one of the nonparallel sides has length
      10. If the nonparallel sides are extended, a triangle is formed by these
      extended sides and the greater base. (a) Show that this triangle is right.
      (b) Compute the lengths of the sides of this triangle.
\item[12.] If, in Exercise 11, the altitude is 6, rather than 8, and the other
      dimensions are as given, compute the length of the altitude of the
      triangle from the vertex opposite the side of length 18.
\end{exlist}

\FIGIXFIFTYTHREEFIFTYFOUR

\begin{exlist}
\item[13.] A parallelogram has adjacent sides of lengths 12 and 8, and these
      sides form a $60^\circ$ angle. Find the lengths of the diagonals of the
      parallelogram. Find the lengths of the diagonals, if the included angle is
      $30^\circ$.
\item[14.] Draw three segments, and label them $AB$, $CD$, and $EF$. Locate by a
      ruler and compass construction a point $X$ on $AB$, such that the ratio of
      $AX$ to $XB$ is equal to the ratio of $CD$ to $EF$.
\item[\stex 15.] In Fig.~9--54, line $l$ cuts sides $AB$ and $BC$ of $\tri{ABC}$
      in points $D$ and $E$, and cuts $AC$ extended in point $F$. Prove that the
      product of the three ratios, $AD$ to $DB$, $BE$ to $EC$, and $CF$ to $FA$,
      is 1. [\emph{Hint:} Consider perpendiculars $AX$, $BY$, and $CZ$ from $A$,
      $B$, and $C$ to $l$. Then $\sg{AD}/\sg{DB} = \sg{AX}/\sg{BY}$, etc.]
\item[16.] Construct a triangle similar to a given triangle and having a given
      perimeter.
\item[17.] Upon a given segment as side, construct a triangle similar to a given
      triangle.
\end{exlist}

\FIGIXFIFTYFIVE

\begin{exlist}
\item[18.] In $\tri{ABC}$, $CD$ and $BE$ are altitudes, as in Fig.~9--55. Prove
      that $\sg{CD}/\sg{EB} = \sg{AC}/\sg{AB}$.
\item[19.] Prove that if, in triangles $ABC$ and $A'B'C'$, $AB \pll A'B'$,
      $BC \pll B'C'$, and $CA \pll C'A'$, then $\tri{ABC} \simto \tri{A'B'C'}$.
\item[20.] Prove that the line joining the midpoints of two opposite sides of a
      parallelogram is parallel to the two other sides and passes through the
      point of intersection of the diagonals.
\end{exlist}
\clearpage
% ==================== ch10 ====================
\setcounter{theorem}{0}\setcounter{definition}{0}\setcounter{construction}{0}\setcounter{exgroup}{0}\setcounter{footnote}{0}
\setcounter{chapter}{10}

\chapter*{\raggedright\Huge AREA}
\addcontentsline{toc}{chapter}{10.\ Area}
\markboth{AREA}{}
\vspace{-1em}\noindent\rule{\linewidth}{0.6pt}\vspace{1em}
\figurekey

%========================================================= 10-1
\section*{10--1\quad INTRODUCTION}
\addcontentsline{toc}{section}{10--1\quad Introduction}

You already have an intuitive understanding of the concept of \emph{area}. In
this chapter we will describe \emph{area} in much the same way as we described
length in Chapter 5. There we determined a number that we called the length of
a line segment by choosing an arbitrary \emph{unit segment} and
\emph{counting}. Here, for a given geometric figure, we determine a number
that we call the \emph{area} of the figure by first choosing an arbitrary
\emph{unit square} and then \emph{counting}.

It is much more difficult to describe the \emph{counting process} for area
than it is for length. We content ourselves with a crude description of the
counting process as it relates to finding areas of rectangles. Once the area
formula for rectangles is determined, we employ it to derive area formulas for
other simple polygons.

%========================================================= 10-2
\section*{10--2\quad AREA BY COUNTING}
\addcontentsline{toc}{section}{10--2\quad Area by counting}

We understand by \emph{polygon}, of course, the set of points on the sides.
When we refer to the \emph{area} of a \emph{polygon}, we really are considering
the set of interior points. It would be more descriptive to speak of the
\emph{area} of this \emph{interior set}. But for brevity, we simply speak of
the area of a polygon. First, we study rectangular polygons.

We need to explain what we mean by the area of a rectangle. It will help to
consider a few examples.

Having chosen a unit length, a square 1 unit of length on a side is said to
have 1 unit of area. Clearly, if we had two squares side by side, as in
Fig.~10--1, we would want to count the area as 2. Suppose we had a rectangle
with sides of lengths 3 and 2. Figure 10--1 shows that we would want the area
to be 6 units.

If we had a square of side $\tfrac{1}{2}$ unit, then placing it in a unit
square as in Fig.~9--2,\footnote{\emph{Transcriber's note.} The 1961 printing
reads ``Fig.~9--2'' here; the figure meant is plainly Fig.~10--2. Reproduced
as printed.} we see that the area should be only $\tfrac{1}{4}$ unit.

If we had a rectangle with sides $2\tfrac{1}{2}$ and 3 units, what would the
area be?

\FIGXONETWO

\exercisegroup{10--1}
% No star footnote here. Text review 2026-08-17 checked the foot of book p.181
% at 200 dpi in BOTH sources (photo PDF 195 and scan13 p2): the page ends with
% Fig. 10-3 and carries no footnote. ch02:311 owns the book-wide first use of
% the "*Starred exercises ... are optional" note. Do not re-add it here.
\begin{multicols}{2}
\begin{exlist}
\item Draw a rectangle whose sides are 3 and 5 units. Determine its area by
  decomposing it into unit squares and counting.
\item Sketch a figure to illustrate why the area of an $m$ by $n$ rectangle is
  $mn$ units. Here, $m$ and $n$ are positive integers.
\item A rectangle is $3\tfrac{1}{3}$ units long and 3 units wide. Find its
  area. What is the largest size square, such that the rectangle can be
  covered by an integral number of squares? Suppose that the rectangle were
  $3\tfrac{1}{3}$ by 2. What is the answer then?
\item[\stex 4.] Suppose that the length of a rectangle is $a$, and its width
  is $b/c$, where $a$, $b$, and $c$ are positive integers. Obtain the area by
  dividing the rectangle into squares $1/c$ units on a side and then counting.
\item[5.] A rectangle has sides 2.2 and 3.3 in. Obtain its area by dividing it
  into small squares and counting.
\item[6.] A rectangle has sides 3.24 and 4.83 ft. Describe how to obtain the
  area by counting squares.
\item[\stex 7.] The sides of a rectangle are $a/b$ and $c/b$ units, where $a$,
  $b$, and $c$ are positive integers. Show how its area is obtained by counting
  squares.
\end{exlist}
\end{multicols}

%========================================================= 10-3
\section*{10--3\quad DEFINITION OF RECTANGLE AREA}
\addcontentsline{toc}{section}{10--3\quad Definition of rectangle area}

In the previous section we saw how to obtain the area of a rectangle by
counting squares if the rectangle has sides whose lengths are integers.
Remember that, so far, we have not explicitly defined what area is, but have
used only our previous experience.

\FIGXTHREE

Suppose the lengths of the sides of a rectangle are $l$ and $w$, where $l$ and
$w$ are terminating decimals. If $l$ and $w$ end in 10ths, then we chop up the
rectangle into squares 0.1 on a side. If $l$ and $w$ end in 100ths, we chop the
rectangle into squares 0.01 on a side, and so on. In Fig.~10--3, we have
started to chop up the rectangle, which is 1.4 by 2.3. Is it necessary to draw
more of the small squares in order to count them?

It is rather obvious that the method of counting could be applied to a
rectangle whose sides are any rational numbers.

\exercise{10--2}
\noindent Find, by counting, the area of the rectangle whose sides are
(a)~2.5 and 3.3, (b)~2.23 and 1.77, (c)~5/3 and 7/3, (d)~5/3 and 1/2,
(e)~5/3 and 3/4.

\medskip
In all our problems so far, we have found, by counting squares, that the area
of the rectangle is $lw$.

Now suppose that at least one of the sides of the rectangle is not a rational
number, that is, this side and the unit are not commensurable. To be definite,
let us consider a rectangle whose length is $\sqrt{2}$ and whose width is 1.
(Remember that $\sqrt{2}$ is not rational, and that \emph{approximately}
$\sqrt{2} = 1.4142$.) Examination of Fig.~10--4 shows that the area of the
rectangle is greater than 1.4 and less than 1.5. We may see this by counting
squares. If we were to subdivide the length into 100ths, we would also see
that the area is greater than 1.41 and less than 1.42. As we approximate to
$\sqrt{2}$ more and more closely by rational numbers, for example, by 1.4,
1.41, 1.414, 1.4142,\dots, we will get rectangles whose areas are ever closer
to $\sqrt{2}$.

% Fig. 10-4 sits at the foot of book p.182, i.e. AFTER Exercise 10-2 and after
% the incommensurable-case prose that references it. It was ahead of the
% exercise here; restored to source order.
\FIGXFOUR

Therefore, the situation in which we find ourselves is as follows. If the
sides of the rectangle are rational numbers, we can get the area by counting.
If the sides are not rational, then simple counting does not give us the area.
In this latter case, we must approximate to the sides by rational numbers and
consider the \emph{limit}*\Xlimitnote{} of the areas of the rectangles whose
sides have rational numbers for lengths. It is possible to prove that this
limit exists and is equal to the product $lw$.

You see that to define the area of a rectangle by starting with counting may be
quite troublesome. Therefore, in this book we do not give a detailed definition
of area of a rectangle, but only state the following theorem, which tells us
how to calculate the area.

\begin{theorem}
The area of a rectangle whose sides have lengths $l$ and $w$ is $lw$.
\end{theorem}

\exercisegroup{10--3}
\begin{multicols}{2}
\begin{exlist}
\item A square has side $x$ yards. What is its area in square feet?
\item What is the area of a rectangle $3\sqrt{2}$ units long and $2\sqrt{2}$
  units wide?
\end{exlist}
\end{multicols}

%========================================================= 10-4
\section*{10--4\quad AREA OF POLYGONS}
\addcontentsline{toc}{section}{10--4\quad Area of polygons}

So far, the only polygons whose areas we have studied are rectangles. In
discussing areas of other polygons, and in particular triangles, perhaps we
should define area in terms of some kind of counting process. If we did, we
would be forced to concern ourselves with limits, so we will determine areas of
polygons by a method used by the Greek geometers. Our basic formula for the
area of a rectangle will enable us to find areas of triangles, and other
figures, without any counting of squares.

Remember that although the Greeks used numbers in commerce, they had no numbers
other than integers in their geometry. They did not compute numerical measures
of angles, lines, or areas. Just as in a primitive society a hunter might
compare his son's height with that of the son of his friend by marking his
spear, so the Greeks devised a method for comparing the sizes of polygons. Of
course, the Greeks had no word in their geometry corresponding to our term
``area,'' since area is a \emph{number}. Instead of saying that two geometric
figures had \emph{equal areas}, they merely remarked that the figures were
\emph{equal}. The Greeks showed that two polygons were \emph{equal} by cutting
them up into congruent triangles. Note that the word \emph{equal} meant
\emph{the same size} to the Greeks. This is different from the word's meaning
in modern mathematics, where \emph{equal} means \emph{identical}. Figure
10--5(a) illustrates what is meant. The Greek geometers would call rectangle
$ADBE$ and isosceles triangle $ABC$ \emph{equal} because
$\tri{ABE} \cong \tri{ABE}$ and $\tri{ABD} \cong \tri{ACE}$.

\FIGXFIVE

\exercise{10--4}
\noindent In Fig.~10--5(b), $A'B'C'D'$ is a parallelogram, and
$\sg{A'B'} \cong \sg{D'E'}$. Show that $A'B'D'E'$ and $A'B'C'D'$ are
\emph{equal} in the sense of Greek geometry. Do the same for parallelograms
$A''B''C''D''$ and $A''E''F''D''$ in Fig.~10--5(c).

\medskip
We use the methods of the Greek geometers to develop area formulas. But before
doing so, we need a few theorems about areas of geometric figures. These
theorems can be proved, using properties of numbers and our geometric
postulates, but their proofs are long and difficult. Accordingly, we treat them
as postulates even though proofs can be given.

\begin{theorem}
Congruent triangles have equal areas.
\end{theorem}

\begin{theorem}
If two simple polygons have one side and only the points of this side in
common, and if, moreover, they have no interior points in common, then the
simple polygon formed by the sides that remain after the common side is
deleted has an area equal to the sum of the areas of the original polygons.
\end{theorem}

Theorem 10--2 is intuitively obvious, since any method of counting squares to
get the area of one triangle can be duplicated for the other triangle. Draw a
figure to illustrate this. Figure 10--6 will make the meaning of Theorem 10--3
clear. The theorem tells us that the area of polygon $ABCDGFE$ is the sum of
the areas of polygons $ABCDE$ and $DEFG$.

\FIGXSIXSEVEN

With these two theorems and our rectangle formula, we can establish theorems on
the areas of triangles and parallelograms.

\begin{theorem}
If $a$ and $b$ are the lengths of the legs of a right triangle, then the
area of the triangle is $\frac{1}{2}ab$.
\end{theorem}

\noindent\emph{Proof.} Let the right triangle be $\tri{ABC}$, as in
Fig.~10--7, and let $C'$ be chosen as indicated in the figure. (How is $C'$
chosen?)

\begin{steps}
\step{$\tri{ABC} \cong \tri{BAC'}$.}{Why?}
\step{Area $(\tri{ABC})$ = area $(\tri{ABC'})$.}{Theorem 10--2.}
\step{Area $(\tri{ABC})$ $+$\newline area $(\tri{ABC'})$ $=$\newline
      area $(\Xrect AC'BC) = ab$.}{Theorem 10--3 and formula for area of
      rectangle.}
\step{2 area $(\tri{ABC}) = ab$.}{Statements 2 and 3.}
\step{Area $(\tri{ABC}) = \frac{1}{2}ab$.}{Statement 4.}
\end{steps}

In a parallelogram, if we call the length of one side the \emph{base}, then the
corresponding \emph{altitude} (or \emph{height}) is the perpendicular distance
from this side to the line containing the opposite side. From Theorem 10--4, we
can deduce the area of a parallelogram.

\begin{theorem}
The area of a parallelogram is equal to the product of its base and its
height.
\end{theorem}

\FIGXEIGHT

\noindent\emph{Proof.} Let the parallelogram be $ABCD$, as in Fig.~10--8, its
base $b = \sg{AD}$, and its height $h$. Consider, as in the figure, the
rectangle $AECF$. (How can the vertices $E$ and $F$ be found?) Let
$\sg{EB} = x$.

\begin{steps}
\step{$\tri{AEB} \cong \tri{CFD}$.}{Why?}
\step{Area $(\tri{AEB})$ = area $(\tri{CFD})$ = $\frac{1}{2}hx$.}{Why?}
\step{Area $(\Xrect AECF)$ $=$\newline area $(\tri{AEB})$ $+$\newline
      area $(\Xrect ABCD)$ $+$\newline area $(\tri{CFD})$.}{Why?}
\step{$\sg{AF} = \sg{EC} = b + x$.}{Why?}
\step{$h(b + x) = \frac{1}{2}hx$ $+$\newline area $(\Xrect ABCD)$ $+$
      $\frac{1}{2}hx$.}{Statements 3 and 4.}
\step{$hb + hx$ = area $(\Xrect ABCD)$ $+$ $hx$.}{Statement 5 and algebra.}
\step{Area $(\Xrect ABCD) = bh$.}{Statement 6 and algebra.}
\end{steps}

We can now quite easily derive the familiar formula,
$A = \frac{1}{2}b \cdot h$, for computing the area of a triangle. Here, $b$
(base) is the length of a side of a triangle, and $h$ (height) is the
perpendicular distance from the opposite vertex to the line containing the
base.

\FIGXNINE

\begin{theorem}
The area of a triangle is given by the formula
$A = \frac{1}{2}b \cdot h$ (Fig.~10--9).
\end{theorem}

The proof is left to the student.

From these last two theorems, we may conclude that parallelograms (or
triangles) with equal bases and heights have equal areas. The area formula for
triangles provides us with a tool for determining areas of polygons in general,
for the area of any polygon can be computed as the sum of the areas of certain
triangles.

\exercisegroup{10--5}
\begin{multicols}{2}
\begin{exlist}
\item A triangle has an area of 30 and a base of 6. Determine its height.
\item Two adjacent sides of a parallelogram have lengths of 12 and 8 in. These
  sides determine a $30^\circ$ angle. Compute the area of the parallelogram.
\item A triangle has area 50 and altitude 8. What is the length of its base?
\item A rhombus has one angle of $45^\circ$ and is 2 ft on a side. Find its
  area.
\item In Fig.~10--10, $l$ and $m$ are parallel and 3 ft apart.
  $\sg{AB} = \sg{CD} = 5$ ft. If the distance from $B$ to $C$ is 1 mi, what is
  the area of the quadrilateral $ABDC$?
\item What is the area of $\Xrect ABCD$ if the sides are as shown in
  Fig.~10--11 and $BD \perp AD$?
\item The area of a triangle is $A$, and its base is $x$. Find the altitude in
  terms of $A$ and $x$.
\item The base of a parallelogram can be expressed in terms of its area and
  height. What is this formula?
\item Two triangles have the same side for base. The height of one triangle
  exceeds that of the other by 4 in. The areas differ by 40 in$^2$. What is the
  length of the common base?
\item The midpoints of the sides of a parallelogram determine a polygon.
  Compare the area of this polygon to the area of the original parallelogram.
\item[\stex 11.] The midpoints $A_1$, $B_1$, and $C_1$ of the sides of
  $\tri{ABC}$ are joined to form a second triangle. Consider also the midpoints
  $A_2$, $B_2$, and $C_2$ of the sides of triangle $A_1B_1C_1$. This process is
  continued. Compare the area of $\tri{A_5B_5C_5}$ with that of $\tri{ABC}$.
\item[\stex 12.] Triangles $ABC$ and $A'B'C'$ are similar. The lengths of two
  corresponding sides are $s$ and $s'$, where $s = 5s'$. If $A$ and $A'$
  represent the respective areas of the triangles, prove that $A = 25A'$.
  Generalize this problem, and establish the theorem that \emph{if two
  triangles are similar, the ratio of their areas is equal to the ratio of the
  squares of the lengths of two corresponding sides.}
\item[13.] $\tri{ABC}$ is similar to $\tri{A'B'C'}$, and $AB = 2A'B'$. How do
  their areas compare? Draw a figure.
\item[\stex 14.] Plans for a house are drawn, so that $\frac{1}{4}$ in. on the
  plan corresponds to 1 ft in the house. If the plans cover an area of
  105 in$^2$, how many square feet in the house?
\item[15.] Two equilateral triangles have sides of 8 and 10 in.,
  respectively. How do their areas compare?
\item[16.] A rhombus has sides 10 in.\ long. If one angle is $30^\circ$, what
  is its area? if $60^\circ$? if $120^\circ$?
\item[17.] Find the area of an isosceles right triangle if one leg is 20 ft
  long.
\item[18.] In Fig.~10--12, $AD \pll BC$. Compute the area of the figure.
\item[\stex 19.] Generalize Exercise 18, designating the lengths of the
  parallel sides, or bases, by $b$ and $b'$ and the height by $h$. Of course,
  you are developing a formula for the area of a trapezoid.
\item[20.] Besides the proof suggested by the dotted line in Fig.~10--12,
  Fig.~10--13 also suggests a proof. Here, it is assumed that $b > b'$. Make a
  proof.
\end{exlist}
\end{multicols}

\FIGXTENELEVEN

\FIGXTWELVETHIRTEEN

\begin{multicols}{2}
\begin{exlist}\setcounter{exlisti}{20}
\item The bases of a trapezoid are 3 and 12, and the altitude is 6. What is
  its area?
\item The area of a trapezoid of height 6 in.\ is 39 in$^2$. If one base is
  6 in.\ what is the other?
\item One base of a trapezoid is twice as long as the other. If the area is
  30 ft$^2$ and the altitude is 5 ft, what are the two bases?
\item An equilateral triangle has sides 18 in.\ long. Find its area in square
  feet.
\item An equilateral triangle has an area of $(9\sqrt{3})/4$. How long is its
  side?
\item In Fig.~10--14, $ABCD$ is a parallelogram. Prove that the four small
  triangles, $\tri{AOB}$, $\tri{BOC}$, etc., all have equal areas.
\item The areas of two similar triangles are 360 and 250 square units. If a
  side of the first triangle is 8 units, how long is the corresponding side of
  the second triangle?
\item[\stex 28.] Prove that the ratio of the areas of similar polygons is equal
  to the square of the ratio of corresponding sides.
\item[29.] Find the area of a regular hexagon 6 in.\ on a side. What will be
  the area of a regular hexagon whose side is 8 in.?
\item[\stex 30.] If $a$, $b$, and $c$ are the lengths of the sides of any
  triangle (Fig.~10--15), and $s = \frac{1}{2}(a + b + c)$, then the area is
  given by
  \[ A = \sqrt{s(s\!-\!a)(s\!-\!b)(s\!-\!c)}. \]
  This is ``Heron's formula,'' named after an Alexandrian Greek by that name of
  about 100 B.C., possibly later. Show that this formula is correct.
\item[31.] The sides of a triangle are 8, 11, and 13 in. Find its area.
\item[32.] The sides of a triangle are 17, 25, and 26 ft. Find its area. Find
  the height to the 25-ft side.
\item[33.] Find the area of an isosceles triangle if the base is 12 in.\ and
  the sides 10 in.
\item[34.] Find the area of an equilateral triangle if one side is 6 in.
\item[35.] If, in Fig.~10--16, $ABCD$ is a parallelogram, with $\sg{BE} = 2$
  and $\sg{ED} = 6$, what is the number
  \[ \frac{\text{Area } (BEC)}{\text{Area } (ABCD)}\,? \]
\item[36.] A square has a diagonal of length 12 in. Compute the area of the
  square. Determine a formula for the area of a square in terms of the length
  $d$ of its diagonal.
\item[37.] The diagonals of a rhombus are 10 in.\ and 6 in., respectively. What
  is its area? Generalize this problem by formulating and proving a theorem.
\item[38.] In Fig.~10--17, $\ang{C}$ is right, $CD \perp AB$, $\sg{AC} = 5$,
  and $\sg{CB} = 12$. Compute the areas of triangles $ABC$, $ACD$, and $BCD$.
  Check your work by verifying that the sum of the area of the latter two
  triangles is equal to the area of $\tri{ABC}$.
\end{exlist}
\end{multicols}

\FIGXFOURTEEN

\FIGXFIFTEEN

\FIGXSIXTEENSEVENTEEN

%========================================================= 10-5
\section*{10--5\quad AREA PROOF OF THE\\PYTHAGOREAN THEOREM}
\addcontentsline{toc}{section}{10--5\quad Area proof of the Pythagorean theorem}

The Greek geometers thought of the Pythagorean theorem, not as a numerical
relationship between lengths of the sides of a right triangle, but rather as a
relationship between areas. Viewed in this light, the theorem may be stated as:

\begin{theorem}
For every right triangle, it is true that the square on the hypotenuse has
an area equal to the sum of the areas of the squares on the other two sides
(see Fig.~10--18).
\end{theorem}

\FIGXEIGHTEENNINETEEN

\begin{center}
Area 1 $+$ area 2 $=$ area 3.
\end{center}

Probably more different proofs have been given for the Pythagorean theorem than
for any other theorem in mathematics. Our earlier proof did not refer to area
at all. It was essentially an ``algebraic style'' proof. We give a proof below
that is a kind of algebraic and area proof. If the sides of the right triangle
are $a$, $b$, and $c$, consider Fig.~10--19.

\noindent\emph{Proof.} Evidently, the area of the large square, of side
$a + b$, is equal to the area of the square of side $c$ plus the areas of the
four triangles. That is,
\[ (a + b)^2 = c^2 + 4(\tfrac{1}{2}ab). \]
But $(a + b)^2 = a^2 + 2ab + b^2$, so that we get
\[ a^2 + 2ab + b^2 = c^2 + 2ab. \]
From algebra then,
\[ a^2 + b^2 = c^2. \]

There are many more proofs of this theorem. Your teacher can show you some if
you are interested.

\exercisegroup{10--6}
\begin{multicols}{2}
\begin{exlist}
\item Triangles $ABC$ and $A'B'C'$ are similar. $\sg{AB} = 3$,
  $\sg{A'B'} = 5$, and area $(ABC) = 18$. Compute area $(A'B'C')$.
\item Determine the area of an equilateral triangle if the length of one side
  is 10.
\item Show that the area of an equilateral triangle is given by the formula
  $A = (s^2/4) \cdot \sqrt{3}$, where $s$ is the length of a side.
\item Develop a formula for the area of a regular hexagon if each side has
  length $c$.
\item The lengths of two adjacent sides of a parallelogram are 12 in.\ and
  8 in. The included angle is $45^\circ$. Determine the area of the
  parallelogram.
\item The altitude of a rectangle is 12 ft. A diagonal has length 20 ft.
  Compute the area of the rectangle.
\item[\stex 7.] A square and a regular hexagon have equal perimeters. Find the
  ratio of their areas.
\item[\stex 8.] It is said that the Hindu mathematician Bhaskara, born about
  1100 A.D., discovered a proof of the Pythagorean theorem and presented the
  figure like that in Fig.~10--20 with the single word, ``Behold.'' Give a
  proof. In Fig.~10--20 $DF \pll AC$, and $EG \pll BC$.
  [\emph{Hint:} $\sg{FC} = a - b$.]
\end{exlist}
\end{multicols}

\FIGXTWENTY

%========================================================= 10-6
\section*{*10--6\quad CONSTRUCTION PROBLEMS}
\addcontentsline{toc}{section}{*10--6\quad Construction problems}

There are many RC constructions related to area concepts. It is interesting to
note that it is possible to construct a triangle whose area is equal to the
area of any given polygon. We indicate in Fig.~10--21 a procedure for
constructing a triangle equal in area to a quadrilateral. $DE$ has been drawn
parallel to diagonal $AC$. Show that area $(ABE)$ = area $(ABCD)$.
[\emph{Hint:} Show that area $(ADC)$ = area $(AEC)$.]

\FIGXTWENTYONE

\exercisegroup{10--7}
\begin{multicols}{2}
\begin{exlist}
\item Construct a triangle whose area is equal to that of pentagon $ABCDE$
  (Fig.~10--22). Use the method employed above to construct a triangle equal in
  area to a given quadrilateral.
\item Show how to construct a square whose area is equal to the sum of the
  areas of two given squares.
\item Solve Exercise 2 with ``sum'' replaced by ``difference.''
\item Construct a square whose area is $\frac{3}{4}$ of the area of a given
  square.
\item Construct a square whose area is equal to the area of a given rectangle.
\item Construct a rectangle whose area is equal to the area of a given
  triangle.
\item Construct a square whose area is the same as the area of a given
  triangle.
\item Construct a triangle similar to a given triangle and having exactly twice
  the area of the given triangle.
\item Show how to construct a polygon of $n - 1$ sides whose area is the same
  as that of an $n$-gon. Assume $n > 3$.
\item[\stex 10.] Prove that if $P$ is any point in the interior of an
  equilateral triangle (Fig.~10--23), the sum of the lengths of the
  perpendiculars from $P$ to the sides of the triangle is equal to the length
  of an altitude of the triangle. [\emph{Hint:} Consider the areas of triangles
  $APB$, $APC$, and $BPC$.]
\end{exlist}
\end{multicols}

\FIGXTWENTYTWOTHREE

%========================================================= review
\section*{REVIEW OF CHAPTER 10}
\addcontentsline{toc}{section}{Review of Chapter 10}

\begin{exlist}
\item In this chapter we have studied methods for finding the areas of
  rectangles, triangles, and polygons. All area formulas were derived from the
  formula for the area of a rectangle.
  \begin{flushleft}
  Area $\Xrect$ $=$ length $\cdot$ width $=$ $lw$.\\
  Area $\Xpgram$ $=$ base $\cdot$ altitude $=$ $bh$.\\
  Area $\triangle$ $=$ $\frac{1}{2}$ base $\cdot$ altitude $=$
    $\frac{1}{2}bh$.\\
  Area $\Xtrap$ $=$ $\frac{1}{2}$ altitude $\cdot$ sum of bases $=$
    $h(b + b')/2$.
  \end{flushleft}
\item Review how the above area formulas were obtained.
\item Areas of similar figures are proportional to the squares of corresponding
  sides.
\end{exlist}

\par\medskip\begin{center}\bfseries Review Exercises\end{center}\par\nopagebreak
\begin{exlist}
\item $\tri{A_1B_1C_1} \sim \tri{A_2B_2C_2}$. If area $(A_2B_2C_2) = 72$, and
  $\sg{A_1B_1}/\sg{A_2B_2} = \tfrac{3}{4}$, find area $(A_1B_1C_1)$.
\item Prove that a median of a triangle divides the triangle into two equal
  areas.
\item[\stex 3.] Prove that the medians divide a triangle into six triangles of
  equal areas (Fig.~10--24).
\item[4.] Prove that the area of a square is half the area of a square whose
  side is the diagonal of the given square.
\item[5.] Two similar triangles have areas of 108 and 48 ft$^2$. The first has
  altitude $7\frac{1}{2}$ ft; find the altitude of the second.
\item[6.] Prove that the area of a trapezoid is equal to the altitude times the
  length of the segment joining the midpoints of the nonparallel sides.
\item[7.] A triangle has one angle of $30^\circ$. Prove that the area is
  $\frac{1}{4}$ of the product of the lengths of the sides adjacent to the
  $30^\circ$ angle.
\item[8.] In $\tri{ABC}$, if $D$ is the foot of the perpendicular from $C$ to
  $AB$, prove that $\sg{AC}^2 - \sg{BC}^2 = \sg{AD}^2 - \sg{BD}^2$.
\item[9.] The hypotenuse of a right triangle is 41 ft. One leg is 40 ft. Find
  the length of the other leg.
\item[10.] The median of a trapezoid is 8 in., and one of the bases is 12 in.
  How long is the other base? If the area is $\frac{1}{2}$ft$^2$, what is the
  altitude?
\item[11.] In a right triangle, the altitude from the right angle is 14 in.
  From the foot of the altitude to a vertex is 8 in. How long is the
  hypotenuse?
\end{exlist}

\FIGXTWENTYFOUR

\par\medskip\begin{center}\bfseries Review Test\end{center}\par\nopagebreak
\begin{exlist}
\item What is the area of a parallelogram with sides 10 and 8 and an included
  angle of $30^\circ$?
\item What is the area of a rhombus with side 10 if an angle is $45^\circ$?
\item Find the area of an isosceles right triangle whose legs are 20.
\item One of the bases of a trapezoid is 15, and its height is 6. Find the
  other base if the area is 75.
\item Two similar triangles, $T$ and $T'$, have two corresponding sides 6 and
  15, respectively. What is the ratio of the area of $T$ to the area of $T'$?
\item The length of a diagonal of a square is 15. Find the area.
\item Two similar polygons have areas of 250 and 160. If the length of a side
  of the first is 6, what is the length of the corresponding side of the
  second?
\item Two triangles are similar, and the lengths of the sides of the first
  triangle are 3, 5, and 6. If the length of the greatest side of the other
  triangle is 15, compute the lengths of the other sides.
\item If the legs of a right triangle are in the ratio of 5 to 4, find the
  ratio of the segments formed on the hypotenuse by the altitude to the
  hypotenuse.
\item If one leg of a right triangle has length 25, and the altitude to the
  hypotenuse is 7, what is the area of the triangle?
\end{exlist}
\clearpage
% ==================== ch11 ====================
\setcounter{theorem}{0}\setcounter{definition}{0}\setcounter{construction}{0}\setcounter{exgroup}{0}\setcounter{footnote}{0}
\setcounter{chapter}{11}

\chapter*{\raggedright\Huge CIRCLES AND\\REGULAR POLYGONS}
\markboth{CIRCLES AND REGULAR POLYGONS}{}
\addcontentsline{toc}{chapter}{11.\ Circles and Regular Polygons}
\vspace{-1em}\noindent\rule{\linewidth}{0.6pt}\vspace{1em}
\figurekey

%========================================================= 11-1
\section*{11--1\quad INTRODUCTION}
\addcontentsline{toc}{section}{11--1\quad Introduction}

In this chapter we define what we mean by the \emph{circumference} (or
\emph{length}) of a \emph{circle} and then prove that the circumference $C$ of
a circle of radius $r$ is given by the familiar formula, $C = 2\pi r$. This
involves giving a definition of the number $\pi$. To do these things, we must
learn a few facts about limits, because the study of circumference is made by
approximating to the circle by polygons and considering polygons with a very
large number of sides. It is therefore necessary, first, to prove some theorems
about regular polygons.

%========================================================= 11-2
\section*{11--2\quad CIRCLES AND REGULAR POLYGONS}
\addcontentsline{toc}{section}{11--2\quad Circles and regular polygons}

Our study of the circle depends upon knowing about regular polygons inscribed
in circles. We therefore study these matters first.

\begin{definition}
A regular polygon is \emph{inscribed} in a circle if all the vertices of the
polygon are on the circle. We also say that the circle \emph{circumscribes the
polygon}.
\end{definition}

\exercisegroup{11--1}
\begin{multicols}{2}
\begin{exlist}
\item With ruler and compass, inscribe a square in a given circle.
\item Using Exercise 1, show how to construct regular inscribed polygons of
      8 sides; 16 sides; 32 sides.
\item With ruler and compass, inscribe a regular hexagon (6-sided polygon) in
      a given circle.
\item Using Exercise 3, show how to construct regular polygons of 3 sides;
      12 sides; 24 sides; 48 sides.
\end{exlist}
\end{multicols}

\begin{theorem}
There is exactly one circle circumscribing a given triangle.
\end{theorem}

\noindent\emph{Proof.} In Fig.~11--1, suppose that the triangle is
$\tri{ABC}$, and choose any two of the sides, say $AB$ and $BC$. Let $m$ and
$n$ be the perpendicular bisectors of these sides. The lines $m$ and $n$ must
intersect at a point $O$, for if $m$ and $n$ did not intersect, they would be
parallel, $CB$ would be perpendicular to $m$, and the points $A$, $B$, and $C$
would be collinear. The point $O$ is equidistant from $A$ and $B$, and also
equidistant from $B$ and $C$. (Why?) Hence $\sg{OA} = \sg{OB} = \sg{OC}$, and
the circle with center $O$ and radius $\sg{OA}$ passes through $A$, $B$, and
$C$. This proves that there is at least one circle through $A$, $B$, and $C$.
To see that there can be only one, we observe that the center of such a circle
must be equidistant from $A$, $B$, and $C$, and hence must lie on the
perpendicular bisectors of $AB$ and $BC$. These lines, $m$ and $n$, intersect
in only one point. This completes the proof.

\FIGXIONETWO

\begin{theorem}
There is exactly one circle that circumscribes a regular polygon.
\end{theorem}

\noindent\emph{Proof.} In Fig.~11--2, suppose the polygon has successive
vertices $A$, $B$, $C$, $D,\ldots$. We show that the circle circumscribing
$\tri{ABC}$ actually circumscribes the polygon. Let $O$ be the center of the
circle circumscribing $\tri{ABC}$, and consider the segments $OA$, $OB$, $OC$,
$OD$, $OE,\ldots$. By definition of circumscribing circle,
$\sg{OA} = \sg{OB} = \sg{OC} = r$. We show that $\sg{OD} = r$. The same method
can be used to prove that $\sg{OE} = \cdots = r$. Hence we need only show that
$\sg{OD} = r$.

\begin{steps}
\step{$\ang{2} \cong \ang{3}$.}{Why?}
\step{$\ang{ABC} \cong \ang{BCD}$, and $AB \cong BC$.}{The polygon is regular.}
\step{$\ang{1} \cong \ang{4}$.}{Addition and subtraction of angles postulate.}
\step{$\tri{OAB} \cong \tri{ODC}$.}{S.A.S.}
\step{$OD \cong OA$.}{Statement 4.}
\step{$\sg{OD} = r$.}{Statement 5.}
\end{steps}

\exercise{11--2}
Suppose that in Theorem 11--1 we had considered the perpendicular bisectors of
$AB$ and $AC$ instead of $AB$ and $BC$. Would the same point $O$ have been
found? What does the theorem say about the perpendicular bisectors of the
sides of a triangle?

%========================================================= 11-3
\section*{11--3\quad TANGENTS TO CIRCLES}
\addcontentsline{toc}{section}{11--3\quad Tangents to circles}

To talk about circles \emph{inscribed} in regular polygons, we need to know
about tangents to circles.

\begin{theorem}
If $OP$ is a radius of a circle with center $O$, and a line $l$ passes through
$P$ and is perpendicular to $OP$, then $l$ has exactly one point in common with
the circle.
\end{theorem}

\noindent\emph{Proof.} In Fig.~11--3, let $A$ be any point on $l$ different
from $P$. Then $\tri{OPA}$ is a right triangle, and $\sg{OA} > \sg{OP}$.
(Why?) Hence $A$ is outside the circle, and $l$ intersects the circle only at
the point $P$.

\begin{definition}
A line $l$ is \emph{tangent} to a circle if it has exactly one point in common
with the circle; the point is called \emph{the point of tangency}.
\end{definition}

\FIGXITHREEFOUR

\begin{XIthmnote}[Converse to Theorem 11--3]
A line tangent to a circle at a point $P$ is perpendicular to the radius
containing $P$.
\end{XIthmnote}

\noindent\emph{Proof.} We will prove the contrapositive, that is, if the line
is not perpendicular to the radius at $P$, then it is not tangent to the
circle. Consider a circle (Fig.~11--4) with center $O$ and a point $P$ on the
circle. Let $l$ be a line through $P$, such that $l$ is not perpendicular to
$OP$.

\begin{steps}
\step{Let $m$ be the line $\perp l$ passing through $O$. Let $A$ be the
      intersection of $m$ and $l$.}{Why is there such a line?}
\step{Let $B$ be the point on $l$ on the opposite side of $A$ from $P$, and
      such that $\sg{AB} = \sg{AP}$.}{Why does such a point exist?}
\step{$\sg{OB} = \sg{OP}$.}{Why? Several different arguments may be given.}
\step{$B$ is on the circle and also on $l$.}{Statements 2 and 3 and definition
      of a circle.}
\step{$l$ intersects the circle twice, and hence $l$ is not tangent to the
      circle. This completes the proof.}{Statement 4 and definition of
      tangent.}
\end{steps}

\begin{XIremarkn}{1}
In the above proof, the point $A$ had to be inside the circle because
$\sg{OA} < \sg{OP}$, by the Pythagorean theorem. The Pythagorean theorem shows
that the perpendicular distance from a point to a line is less than the
distance to any other point on the line.
\end{XIremarkn}

\begin{XIremarkn}{2}
An argument similar to the proof of Theorem 11--4, and using the Pythagorean
theorem, can be made to show that a line can intersect a circle at most twice.
Hence, a line will either be tangent to a circle, or intersect the circle in
exactly two points, or fail to intersect the circle.
\end{XIremarkn}

\exercisegroup{11--3}
\begin{multicols}{2}
\begin{exlist}
\item Draw a circle, and construct (\RC) the line tangent to the circle at a
      point of the circle.
\item[\stex 2.] From a point outside a circle, construct a line tangent to the
      circle.
\item[3.] Prove that if $A$ is a point not on a line $l$, and $B$ and $C$ are
      distinct points on $l$, with $AB \perp l$, then $\sg{AB} < \sg{AC}$.
\end{exlist}
\end{multicols}

\begin{definition}
A circle is said to be \emph{inscribed} in a polygon, if it is tangent to each
of the sides of the polygon. We also say that the polygon \emph{circumscribes
the circle}.
\end{definition}

\begin{theorem}
There is a circle inscribed in every regular polygon.
\end{theorem}

\begin{steps}
\step{Let $O$ be the center of the circle circumscribing the polygon
      (Fig.~11--5).}{Theorem 11--2.}
\step{Let $AB$ be a side of the polygon, and $C$ the midpoint of $AB$.}{$AB$
      has one, and only one, midpoint.}
\step{$OC \perp AB$.}{Why?}
\step{The circle with center $O$ and radius $OC$ is tangent to $AB$ at
      $C$.}{Theorem 11--4.}
\step{Now let $BD$ be a side of the polygon adjacent to $AB$, and $E$ the
      midpoint of $BD$. Then $OE \cong OC$.}{Why? This takes several steps.}
\step{Hence the circle with radius $OC$ is also tangent to $BD$ at
      $E$.}{Statement 5 and the fact that $OE \perp BD$.}
\end{steps}

\noindent 7. The above argument can be continued around the polygon to complete
the proof.

\FIGXIFIVESIX

\exercise{11--4}
Show that congruent chords in a circle are equidistant from the center of the
circle.

\begin{definition}
The \emph{center} of a regular polygon is the center of the inscribed or
circumscribed circle. The \emph{apothem} of a regular polygon is the length of
the radius of the inscribed circle (that is, a number). The \emph{radius} of a
regular polygon is the radius of the circumscribed circle. A \emph{central
angle} of a regular polygon is an angle, with vertex at the center, whose sides
contain adjacent vertices of the polygon. See Fig.~11--6.
\end{definition}

\exercisegroup{11--5}

\FIGXISEVEN

\begin{multicols}{2}
\begin{exlist}
\item Inscribe a regular hexagon in a circle. Draw the circle inscribing the
      hexagon.
\item Using your experience in measuring angles in degrees (we will study this
      topic in the next chapter), how many degrees in a central angle of a
      regular pentagon? hexagon? $n$-gon?
\item What is the sum of the degree measures of the angles of a regular
      pentagon? hexagon? $n$-gon?
\item If an exterior angle of a regular polygon is as shown in Fig.~11--7, how
      many degrees in it if the polygon has 5 sides? 6 sides? $n$ sides?
\item Prove that a central angle of a regular polygon is the supplement of any
      angle of the polygon.
\item Show that an exterior angle of a regular polygon is congruent to a
      central angle.
\item In Definition 11--4, how do we know that the center of the inscribed
      polygon and the center of the circumscribed polygon are identical?
\end{exlist}
\end{multicols}

%========================================================= 11-4
\section*{11--4\quad REGULAR POLYGONS AND SIMILARITY}
\addcontentsline{toc}{section}{11--4\quad Regular polygons and similarity}

All our proofs about area and circumference of circles depend upon the
following theorem.

\begin{theorem}
Two regular polygons are similar if, and only if, they have the same number of
sides.
\end{theorem}

\noindent\emph{Preliminary remark.} Recall that two polygons are similar if
there is a correspondence between their angles and sides, such that
corresponding angles are congruent and corresponding sides are proportional.
Thus, if the polygons are similar, they necessarily have the same number of
sides, and the ``only if'' part of the theorem is obvious. Therefore, we have
only to prove that if they have the same number of sides, they must be similar.

\noindent\emph{Proof.} Suppose that two regular polygons have consecutive
vertices $A$, $B$, $C$, $D$, $E,\ldots$ and $A'$, $B'$, $C'$, $D'$,
$E',\ldots$, respectively, as in Fig.~11--8. Suppose also that the polygons are
not similar. Since they are regular,
\[
\frac{\sg{AB}}{\sg{A'B'}} = \frac{\sg{BC}}{\sg{B'C'}}
= \frac{\sg{CD}}{\sg{C'D'}} = \frac{\sg{DE}}{\sg{D'E'}} = \cdots.
\]
It must be, then, that $\ang{A} \ncong \ang{A'}$, $\ang{B} \ncong \ang{B'}$,
etc. If $\ang{A} \ncong \ang{A'}$, then $\ang{A} > \ang{A'}$ or
$\ang{A} < \ang{A'}$. Suppose that $\ang{A} < \ang{A'}$. Since the polygons are
regular, $\ang{B} < \ang{B'}$, $\ang{C} < \ang{C'}$, etc. The exterior angle at
$A$ is greater than the exterior angle at $A'$. (Why?) This is true also for
the exterior angles at $B$ and $B'$, $C$ and $C'$, etc. The sum of $n$ exterior
angles in polygon $ABCDE\ldots$ is two straight angles. Clearly then, the sum
of $n$ exterior angles in polygon $A'B'C'D'E'\ldots$ is not two straight
angles. This contradiction completes the proof.

\begin{remark}
If we had studied measurement of angles, we could more easily have proved that
$\ang{AOB} \cong \ang{A'O'B'}$, because the measure of these angles depends
only upon the number of sides $n$. In degrees, the measure of $\ang{AOB}$ would
be equal to $360/n$.
\end{remark}

\FIGXIEIGHTNINE

\begin{theorem}
The perimeters of regular polygons, with the same number of sides, are
proportional to their radii (or apothems).
\end{theorem}

\noindent\emph{Proof.} Suppose that two polygons (of $n$ sides) are labeled as
in Fig.~11--9, and have perimeters $p$ and $p'$. We must show that
$p/p' = r/r' = a/a'$.

\begin{steps}
\step{$\tri{AOQ} \simto \tri{A'O'Q'}$.}{Why?}
\step{$r/r' = a/a' = \sg{AQ}/\sg{A'Q'}$.}{Statement 1.}
\step{$p' = n \cdot \sg{A'B'} = 2n \cdot \sg{A'Q'}$, and
      $p = n \cdot \sg{AB} = 2n \cdot \sg{AQ}$.}{Definition of perimeter and
      the fact that $Q$ and $Q'$ are midpoints.}
\step{$p/p' = \sg{AQ}/\sg{A'Q'}$.}{Statement 3 and algebra.}
\step{$p/p' = r/r' = a/a'$.}{Statements 2 and 4.}
\end{steps}

\begin{theorem}
The area of a regular polygon is equal to one-half the product of apothem and
perimeter; area $= \frac{1}{2}a \cdot p$.
\end{theorem}

\noindent The proof is left as an exercise.

\begin{theorem}
The areas of two regular polygons, with the same number of sides, are
proportional to the squares of their radii (or the squares of their apothems,
or the squares of their sides).
\end{theorem}

\noindent The proof is left as an exercise.

\exercisegroup{11--6}
\begin{multicols}{2}
\begin{exlist}
\item Two regular heptagons (7-gons) have sides 3 in. and 4 in., respectively.
      What is the ratio of their perimeters? of their areas?
\item A regular pentagon has twice the area of another regular pentagon. What
      is the ratio of their perimeters?
\item Find the radius and apothem of an equilateral triangle of side 2 in.
\item Two regular polygons of the same number of sides have perimeters of 28
      and 42 in. The radius of the first polygon is $x$ in. What is the radius
      of the second?
\item The apothems of two regular hexagons are 3 in. and 5 in. What is the
      ratio of their areas?
\item The area of a regular 17-gon is 16/9 times the area of a second 17-gon.
      What is the ratio of their perimeters? radii?
\item If a regular hexagon and an equilateral triangle are inscribed in a
      circle, the side of the hexagon is twice the apothem of the triangle.
      Show this.
\item Find the areas of the squares inscribed in and circumscribed about a
      circle of radius 1. Also find the perimeters of the squares.
\item Repeat Exercise 8, using hexagons instead of squares.
\end{exlist}
\end{multicols}

%========================================================= 11-5
\section*{11--5\quad LIMITS}
\addcontentsline{toc}{section}{11--5\quad Limits}

To investigate circumference and area of circles, it is necessary for us to
discuss \emph{limits}. Precise definitions relating to limits are difficult to
grasp, and a careful study is first made in college calculus courses. For our
purposes, only a little information about limits is required. All our
statements about limits are statements about \emph{numbers}. Therefore you may
regard these statements as giving information about numbers that we will
\emph{assume} or postulate.

\begin{definition}
Suppose that for each positive integer $n$, there is associated in some way, a
number $x_n$. The numbers $x_1$, $x_2,\ldots, x_n,\ldots$ then form an infinite
sequence. The number $x_n$ is called the $n$th term of the sequence. We shall
indicate the sequence by $\{x_n\}$.
\end{definition}

For example, let the $n$th term of a sequence be given by $x_n = 2n$. The
positive integer 1 gives the first term $2 \cdot 1$. The positive integer 2
gives the second term $2 \cdot 2$. The third term is $2 \cdot 3$, and so on.
The sequence given by $x_n = 2n$, then, would look like this:

\begin{center}
\begin{tabular}{r@{\hspace{1.2em}}ccccccc}
Positive integers $\rightarrow$ & 1 & 2 & 3 & 4 & $\ldots$ & $n$ & $\ldots$\\[1pt]
 & $\updownarrow$ & $\updownarrow$ & $\updownarrow$ & $\updownarrow$ & &
   $\updownarrow$ & \\[1pt]
Terms of the sequence $\rightarrow$ & 2 & 4 & 6 & 8 & $\ldots$ & $2n$ &
   $\ldots$\,.
\end{tabular}
\end{center}

\exercisegroup{11--7}
\begin{multicols}{2}
\begin{exlist}
\item Write the first six terms of the sequence whose $n$th term is given by
      $x_n = 1/n$; by $x_n = \frac{1}{2}n$; by $x_n = \frac{1}{5}n$.
\item Write the 100th term for each of the sequences in Exercise 1. When $n$ is
      very large, what number are the terms close to?
\item Write the first 10 terms of the sequence whose $n$th term is
      $x_n = (n+1)/n$. Write the 100th term; the 1000th. What number are these
      terms close to when $n$ is extremely large?
\item A sequence is given by $x_n = 2/(n+3)$. Write the first six terms; the
      100th term; the millionth term. Is there a number that the terms of the
      sequence are close to for large values of $n$?
\item Consider a regular polygon of $n$ sides inscribed in a circle. Let $p_n$
      be the perimeter of the polygon. Does this define a sequence?
\item A sequence is given by $x_n = 2/\sqrt{n}$. Write the first six terms. Is
      there a number that the terms of the sequence are close to for large
      values of $n$?
\item Imagine $\sqrt{2}$ written as an infinite decimal,
      $\sqrt{2} = 1.4142\ldots$. Define a sequence $x_n$ as follows: $x_n =$
      the decimal obtained by cutting off the decimal expansion of $\sqrt{2}$
      after the $n$th decimal place. Then $x_1 = 1.4$, $x_2 = 1.41$, etc. What
      is $x_4$? $x_5$? To what number is $x_n$ close when $n$ is large?
\end{exlist}
\end{multicols}

We can now define the \emph{limit} of a sequence, an idea that should be clear
from the above exercises. We shall be content with expressing this idea in an
intuitive way, but it should be mentioned that a more complicated definition is
necessary in order to \emph{prove} theorems about limits. We \emph{assume} that
the theorems we need concerning limits are true.

\begin{definition}
The number $L$ is called the \emph{limit} of the sequence $\{x_n\}$ if, for
large values of $n$, the term $x_n$ is necessarily arbitrarily close to $L$. We
then say that $L$ is \emph{the limit of the sequence of terms $x_n$} and write
``$L = \lim\,\{x_n\}$.''
\end{definition}

\exercisegroup{11--8}
In each of the following problems, the $n$th term of a sequence is given. You
are to decide by inspection what the limit of the sequence is. Note that a
sequence does not have to possess a limit.
\begin{multicols}{2}
\begin{exlist}
\item $x_n = 1/n$.
\item $x_n = 1/10n$.
\item $x_n = (n+1)/n$.
\item $x_n = \sqrt{n}$.
\item $x_n = 5$.
\item $x_n = 2n/(n+1)$.
\item $x_n = n + 2$.
\item[\stex 8.] $x_n = (2n + 100)/(4n + 50)$.
\item[\stex 9.] $x_n = (n + 1)/n^2$.
\item[\stex 10.] $x_n = (n^2 + 5n + 6)/[2n(n + 2)]$.
\end{exlist}
\end{multicols}

When we discuss the circumference and area of circles, we will need the
following two theorems on limits. We assume that these theorems are true, and
thus treat them as postulates. But they are not postulates about our
\emph{geometry}; they are postulates about the \emph{number system}.

\begin{theorem}
Consider two sequences with $n$th terms $x_n$ and $y_n$, respectively, and
suppose that $\lim\,\{x_n\} = X$ and $\lim\,\{y_n\} = Y$. Then the sequence
whose $n$th term is $x_n y_n$ has a limit, and $\lim\,\{x_n y_n\} = XY$.
\end{theorem}

\begin{theorem}
With the same notation as in Theorem 11--10, and the additional hypothesis that
$y_n$ and $Y$ are not zero, then the sequence whose $n$th term is $x_n/y_n$ has
a limit, and $\lim\,\{x_n/y_n\} = X/Y$.
\end{theorem}

\exercisegroup{11--9}
\begin{exlist}
\item Verify Theorem 11--10 if $x_n = (2n + 1)/n$ and $y_n = n/(n + 1)$.
\item If $\lim\,\{x_n\} = 2$, what is $\lim\,\{(x_n)^2\}$?
\item If $\lim\,\{x_n\} = 3$ and $\lim\,\{y_n\} = 4$, what is
      $\lim\,\{x_n/y_n\}$? $\lim\,\{x_n y_n\}$? $\lim\,\{y_n/x_n\}$?
\item If $x_n$ is ``very close'' to 5 and $y_n$ is ``very close'' to 7, what is
      $x_n y_n$ close to? $x_n/y_n$? $y_n/x_n$?
\end{exlist}

%========================================================= 11-6
\section*{11--6\quad CIRCUMFERENCE}
\addcontentsline{toc}{section}{11--6\quad Circumference}

From your everyday experience, you already have an idea of the perimeter or
circumference of a circle. For example, if you were asked to find the
circumference of a wheel, you might wrap a string around the wheel, then pull
it out straight, and measure its length. But such a practical method is not
adequate for geometry, since we cannot describe this process in our geometry.
It is necessary that we first \emph{define} circumference. The definition must
satisfy our intuition about what we \emph{feel} the circumference should be,
and should be one that can readily be \emph{used} to prove theorems.

Observe that we had no difficulty in defining the perimeter of a polygon,
because the polygon consisted of a finite number of segments, each of which had
a length. A circle contains no segment of a line, so we cannot define its
perimeter so simply. To arrive at our definition, we need the following
theorem, the truth of which we assume.

\begin{theorem}
For a given circle, let $p_n$ be the perimeter of a regular polygon of $n$
sides inscribed in the circle. The numbers $p_n$ $(n = 3, 4, 5,\ldots)$ form a
sequence that has a limit.
\end{theorem}

\begin{XIremarks}
The proof of this theorem involves two steps. First, it is shown that the
sequence is increasing, that is, each term is less than the one that follows
it. Then it is shown that the sequence is \emph{bounded}, that is, the terms do
not become indefinitely large. It then follows, as a fundamental property of
real numbers, that such a sequence has a limit.
\end{XIremarks}

It is natural to ask how we know that regular polygons with $n$ sides exist for
all positive integers $n \XIgeqq 3$. We do know that there are regular polygons
with the number of sides equal to $2^2$, $2^3$, $2^4$, etc., for we can
inscribe a square in the circle, and, by bisecting the central angles, we get a
regular octagon, and then a regular 16-gon, and so on, as in Fig.~11--10.

\FIGXITEN

Therefore, the perimeters of these special regular polygons form a sequence
$p_4$, $p_8$, $p_{16}$, $p_{32},\ldots$, and by Theorem 11--12 and the
definition of limit, this sequence has a limit.

We can now give a definition of circumference. Observe that in order to
\emph{define} the circumference of a circle, we are \emph{necessarily forced}
to use limits.

\begin{definition}
The \emph{circumference $C$} of a circle is the limit of the sequence of
perimeters $p_n$ of regular inscribed polygons. $C = \lim\,\{p_n\}$.
\end{definition}

\exercisegroup{11--10}
\begin{multicols}{2}
\begin{exlist}
\item For a given circle, show that the perimeter of an inscribed square is
      less than that of an inscribed octagon.
\item Find the perimeter of a square inscribed in a circle of radius 1. Hence
      show that the circumference of the circle is more than $4\sqrt{2}$.
\item By inscribing a regular hexagon in a circle of radius 1, show that the
      circumference of the circle is more than 6.
\end{exlist}
\end{multicols}

The fundamental result on the circumferences of circles is the following:

\begin{theorem}
The circumferences of two circles are proportional to their radii (or to their
diameters). That is, if the circles have circumferences $C$ and $C'$, radii $r$
and $r'$, and diameters $d$ and $d'$, then $C/C' = r/r' = d/d'$.
\end{theorem}

\noindent\emph{Proof.}
\begin{steps}
\step{$C/C' = \lim\,\{p_n\}/\lim\,\{p'_n\}$, where $p_n$, $p'_n$, are the
      perimeters of the regular inscribed polygons of $n$ sides in the
      circles.}{Definition 11--7.}
\step{$C/C' = \lim\,\{p_n/p'_n\}$.}{Theorem 11--11 on limits.}
\step{$p_n/p'_n = r/r' = d/d'$.}{Theorem 11--7.}
\step{$\lim\,\{p_n/p'_n\} = \lim\,\{r/r'\}$.}{$p_n/p'_n = r/r'$ for all $n$.}
\step{$\lim\,\{r/r'\} = r/r'$.}{$r$ and $r'$ are the same for all $n$.}
\step{$\lim\,\{p_n/p'_n\} = r/r' = d/d'$.}{Statements 3, 4, and 5.}
\step{$C/C' = r/r' = d/d'$.}{Statements 2 and 6.}
\end{steps}

\begin{remark}
Observe that this theorem about circumferences is basically a theorem about
similar polygons.
\end{remark}

\begin{theorem}
If $C$ is the circumference of a circle, and $d$ is its diameter, the ratio
$C/d$ is the same for all circles.
\end{theorem}

\noindent\emph{Proof.} This theorem is really just a corollary to Theorem
11--13. But it is so important that it is given a place by itself. The proof is
as follows. From Theorem 11--13, we have $C/C' = d/d'$. This proportion implies
the proportion $C/d = C'/d'$. This last equality says that $C/d$ is the same
number for any two circles. This completes the proof.

\begin{definition}
The number $C/d$, which is the same for all circles, will be denoted by the
Greek letter $\pi$ (pi).
\end{definition}

\begin{corollary}
$C = \pi d = 2\pi r$.
\end{corollary}

\begin{remark}
The definition of $\pi$ does not immediately give its numerical
value. So a problem remains---namely that of calculating $\pi$ to any desired
degree of accuracy. We examine later a bit of the history of this number and
methods for calculating its value. At the moment, and for purposes of the
exercises, we mention two approximate values commonly used for $\pi$, namely
22/7 and 3.1416.
\end{remark}

\exercisegroup{11--11}
\begin{multicols}{2}
\begin{exlist}
\item Find, approximately, the circumference of a circle whose radius is 7 in.;
      2.131 in.; 49 ft.
\item Find, approximately, the radius of a circle whose circumference is 66 ft;
      12.125 in.; 100.00 ft.
\item Give an exact expression for the area of a square whose perimeter is
      equal to the circumference of a circle of radius 1.
\item Suppose that a circular steel band surrounds the earth at the equator
      (radius $=$ 4000 mi). If the band is made longer by inserting a piece
      1 ft long, about how far will the band be from its previous position?
\item A tire on a car has a diameter of 26 in. If it revolves at 10 revolutions
      per second, what is the approximate speed of the car in miles per hour?
\item About how many revolutions per minute does a car wheel make (diameter
      $=$ 28 in.) when the car is traveling at 60 mph?
\item The radius of one circle is three times the radius of a second. How do
      their circumferences compare?
\end{exlist}
\end{multicols}

%========================================================= 11-7
\section*{11--7\quad COMPUTATION OF PI}
\addcontentsline{toc}{section}{11--7\quad Computation of $\pi$}

There are many ways of computing $\pi$, some much more convenient than others.
The obvious way is to compute the perimeters $p_n$ of regular polygons with a
large number of sides. Then $\pi$ would be given approximately by $p_n/d$. This
is a rather tedious process, and we use it here to get a very rough estimate of
$\pi$.

\FIGXIELEVEN

Consider a regular inscribed hexagon in a circle of radius 1, as shown in
Fig.~11--11. The perimeter of the hexagon is 6. Increasing the number of sides
will increase the perimeter. Hence $C > 6$, so that $\pi = C/d > 6/2 = 3$, and
$\pi > 3$.

If we consider a regular circumscribed hexagon, we find that its perimeter is
$4\sqrt{3}$, and since $C$ is less than this perimeter,
\[
\pi = \frac{C}{d} < \frac{4\sqrt{3}}{2} = 2\sqrt{3} < 3.465,
\]
and $\pi < 3.465$.

In this way, we have obtained a very rough value for $\pi$. It lies somewhere
between 3 and 3.5. By these same methods, Archimedes (3rd century B.C.) showed
that $\pi$ is between 3 1/7 and 3 10/71. The Babylonians used 3 as the value of
$\pi$, and there is a verse in the Bible that implies that the writer believed
$\pi = 3$.

By far the most convenient and rapid methods for calculating $\pi$ are based on
formulas derived in calculus.\XIsymfoot{A formula from calculus that can be
used to calculate $\pi$ is
\[
\pi = 4\left(1 - \tfrac{1}{3} + \tfrac{1}{5} - \tfrac{1}{7} + \tfrac{1}{9}
- \tfrac{1}{11} + \cdots\right).
\]}
We cannot take up these formulas here, but it is interesting to know that in
the 17th century a Scot by the name of Shanks,\XIsymfoot{One of your authors
makes no claim to be a descendant of this misguided computer.} using one of
these formulas, computed $\pi$ to 707 decimal places. His reasons for doing
this are at best obscure. More recently, a group working on one of the large
electronic computers set up a program for the calculation of $\pi$ on the
computer and obtained, in a few hours, $\pi$ to more than 2000 decimal places!
This demonstrates the extreme speed of modern computing machines. The
corresponding job by hand would take hundreds of years. The correct value of
$\pi$ to 14 decimal places is 3.14159265358979.

The number $\pi$ is an irrational number. A proof of this is not easy, and was
given first by the German mathematician Johann Lambert in 1766.

\exercisegroup{11--12}
\begin{multicols}{2}
\begin{exlist}
\item By considering inscribed and circumscribed squares, show that
      $2\sqrt{2} < \pi < 4$.
\item[\stex 2.] By considering an inscribed regular octagon, show that
      $\pi > 3.06$.
\item[\stex 3.] Find the perimeter of a regular 12-gon in a circle of radius 1.
      Hence find an approximate value for $\pi$. Can you continue to find the
      perimeter of a regular 24-gon?
\item[4.] What does the irrationality of $\pi$ imply about its decimal
      representation?
\end{exlist}
\end{multicols}

%========================================================= 11-8
\section*{11--8\quad AREA OF CIRCLES}
\addcontentsline{toc}{section}{11--8\quad Area of circles}

In studying area we first obtained the area of a rectangle by counting squares.
Areas of triangles may be obtained by forming parallelograms and observing that
the triangle is half the area of the parallelograms. Areas of polygons were
obtained by cutting up the polygons into triangles. To define what we mean by
the area of a circle, it is necessary to use again the idea of limit, just as it
was necessary to use limits in order to define circumference. We all have a
fairly clear idea of what we \emph{feel} the area of a circle should be, and we
use this impression to define area. Since the area of the circle should be
close to the area of the inscribed polygon with a large number of sides, we
first need to know that the limit of the areas of inscribed polygons exists.
This fact is explicitly stated as a theorem below, which we accept without
proof.

\begin{theorem}
Given a circle $C$, let $S_n$ denote the area of an inscribed regular polygon
of $n$ sides; then $\lim\,\{S_n\}$ exists. In other words, there is a number
$S$, such that the term $S_n$ is arbitrarily close to $S$ for large $n$.
\end{theorem}

\begin{XIremarkn}{1}
We know, from Section 11--6, that regular polygons exist with the number of
sides equal to powers of 2. In using Theorem 11--15, we can consider only
polygons that we know to exist.
\end{XIremarkn}

\begin{XIremarkn}{2}
Theorem 11--15 is a theorem about numbers. It has little geometric content,
so we need not feel that it is another geometric postulate. If we consider only
polygons with 4, 8, 16, 32,\ldots\ sides, it is easy to show that each area in
our sequence $S_4$, $S_8$, $S_{16}$, $S_{32},\ldots$ is less than the one that
follows. On the other hand, we can see that each of these polygons has an area
less than the area of a circumscribed square. This indicates that our sequence
is not only increasing, but is ``bounded above.'' It is a theorem about numbers
that such a sequence has a limit.
\end{XIremarkn}

\begin{definition}
The area $S$ of a circle is the limit of the areas of inscribed regular
polygons. That is, $S = \lim\,\{S_n\}$.
\end{definition}

We can now prove the main theorem about the area of a circle. But because the
definition of area uses the idea of limit, we must use some theorems about
limits. In particular, we need to know the following fact about the apothems of
the inscribed polygons.

\begin{theorem}
Let $a_n$ be the apothem of an inscribed regular polygon of $n$ sides in a
circle of radius $r$, then $\lim\,\{a_n\} = r$.
\end{theorem}

\FIGXITWELVE

\noindent\emph{Proof.} In Fig.~11--12, $\sg{AC}$ is the length of a side of the
polygon, and is a small number when the number of sides of the polygon is
large. By the Pythagorean theorem, $a_n = \sqrt{r^2 - \sg{AB}^2}$. Thus, when
$n$ is very large, $\sg{AB}^2$ is very small, and $r^2 - \sg{AB}^2$ is very
near $r^2$. Hence $a_n$ will be near $r$. This is exactly what we mean when we
say that $\lim\,\{a_n\} = r$.

\begin{theorem}
The area of a circle is equal to one-half the circumference times the radius.
$S = \frac{1}{2}$ (circumference) $\times$ (radius).
\end{theorem}

\noindent\emph{Proof.}
\begin{steps}
\step{$S = \lim\,\{S_n\}$, where $S_n$ is the area of the regular inscribed
      polygon of $n$ sides.}{Definition.}
\step{$S_n = \frac{1}{2}p_n \cdot a_n$, where $p_n$ is the perimeter, and $a_n$
      is the apothem, of the polygon.}{Theorem 11--8.}
\step{$\lim\,\{S_n\} = \frac{1}{2}(\lim\,\{p_n\}) \cdot
      (\lim\,\{a_n\})$.}{Statement 2 and Theorem 11--10 on limits.}
\step{$\lim\,\{p_n\} = C =$ circumference.}{Definition of circumference.}
\step{$\lim\,\{a_n\} = r =$ radius.}{Theorem 11--16.}
\step{$\lim\,\{S_n\} = \frac{1}{2}C \cdot r$.}{Statements 3, 4, and 5.}
\step{$S = \frac{1}{2}C \cdot r$.}{Statements 1 and 6.}
\end{steps}

\begin{corollary}
The area of a circle of radius $r$ is $\pi r^2$.
\end{corollary}

\begin{corollary}
The areas of two circles are proportional to the squares of their radii (or
diameters).
\end{corollary}

\noindent The proofs are left to the student.

\exercisegroup{11--13}
\begin{multicols}{2}
\begin{exlist}
\item Two circles have radii of 5 in. and 2.5 in. Find their areas. What is the
      ratio of their areas?
\item An annular ring (Fig.~11--13) has inner radius 2 ft and outer radius
      4 ft. Find its area.

      \FIGXITHIRTEEN

\item Figure 11--14 is composed of a circle and semicircles. If
      $\sg{AB} = \sg{BC} = 2$, $B$ is the center of the circle, and $AC$ is a
      diameter, find the area of the shaded portion.

      \FIGXIFOURTEEN

\item Show that the area of the inscribed circle is half the area of the
      circumscribed circle of the square (Fig.~11--15).

      \FIGXIFIFTEEN

\item The area of a circle is 154 in\textsuperscript{2}. Find its radius.
      (Use $\pi = 22/7$.)
\item The area of a circle is 100 ft\textsuperscript{2}. Find its radius.
\item Semicircles are constructed on the sides of a right triangle, as in
      Fig.~11--16. Show that the area of the largest is the sum of the areas of
      the other two.

      \FIGXISIXTEEN

\item A circle has its circumference equal to its area. What is its radius?
\item A rectangular barn is 30 ft by 40 ft. A goat is tethered at one corner by
      a rope 40 ft long. Over what area can he graze? Over what area can he
      graze if the rope is 70 ft long?
\item The circumference of a circle is 50 ft. What is its area?
\item The area of a circle is 288 in\textsuperscript{2}. What is the area if
      the unit of length is a foot?
\item On Mars the standard unit of length is the \emph{pfutth}, which is 0.88
      as long as our foot. A Martian states that the area of a circle is 16
      square \emph{pfutths}. What is the area in square feet?
\item[\stex 13.] Show that the area of a circle is the limit of the areas of
      circumscribed regular polygons.
\item[14.] The area of a circle is 49/4 times the area of another circle. What
      is the ratio of their radii?
\end{exlist}
\end{multicols}

%========================================================= review
\section*{REVIEW OF CHAPTER 11}
\addcontentsline{toc}{section}{Review of Chapter 11}

The following definitions were given in this chapter. Be sure that you know
them.

\begin{center}
\begin{minipage}{0.66\linewidth}
inscribed polygon (circle)\\
circumscribed polygon (circle)\\
tangent\\
center, radius, apothem of regular polygons\\
sequence, limit\\
circumference\\
$\pi$\\
area of a circle
\end{minipage}
\end{center}

The following is a list of the more important theorems of the chapter. Give
complete statements of these theorems. Then try to give a proof for each. Do
not refer to the text until you need help.

\begin{exlist}
\item A regular polygon has a circumscribed circle.
\item A line tangent to a circle is perpendicular to a radius.
\item Converse of 2 (almost).
\item Regular polygons are similar if and only if they have the same number of
      sides.
\item The perimeters of similar polygons are proportional to their radii.
\item The area of a regular polygon is one half the product of apothem and
      perimeter.
\item The circumferences of two circles are proportional to their radii.
\item The area of a circle is equal to one half the circumference times the
      radius.
\end{exlist}

\begin{center}\bfseries Miscellaneous Questions\end{center}
\nopagebreak

\begin{exlist}
\item Why is it important in this chapter to know that regular polygons of the
      same number of sides are similar?
\item Why is it important that circumference is proportional to the radius?
      (Tell why this enables us to define $\pi$.) This theorem depends on what
      theorems on similar polygons?
\item Tell why we must \emph{define} circumference and area for circles.
\item In proving that area $= \pi r^2$, what theorems about polygons did we
      need? what theorems on limits?
\end{exlist}
\clearpage
% ==================== ch12 ====================
\setcounter{theorem}{0}\setcounter{definition}{0}\setcounter{construction}{0}\setcounter{exgroup}{0}\setcounter{footnote}{0}
\setcounter{chapter}{12}

\chapter*{\raggedright\Huge MEASUREMENT OF\\ANGLES AND ARCS}
\addcontentsline{toc}{chapter}{12.\ Measurement of Angles and Arcs}
\markboth{MEASUREMENT OF ANGLES AND ARCS}{}
\vspace{-1em}\noindent\rule{\linewidth}{0.6pt}\vspace{1em}
\figurekey

%========================================================= 12-1
\section*{12--1\quad INTRODUCTION}
\addcontentsline{toc}{section}{12--1\quad Introduction}

In the first chapter of this book you used a protractor to measure
representations of angles. Finally, at this late stage, we prove that the
process of angle measurement can be made precise in our geometry. In order to
discuss angle measurement, it is convenient first to define the measure of an
arc of a circle. This definition is quite like the definition of circumference
and really only makes precise the intuitive idea of \emph{length of arc} that
you already have. We needed limits to discuss circumference, and of course we
also need limits in order to discuss arc length.

Once we know what is meant by arc length, it will be relatively easy to discuss
measurement of angles. This discussion will go quite like Chapter 5 on
measurement of length of segments.

One reason that a precise description of measurement of angles in degrees is
difficult is that describing a unit angle is troublesome. (Is there an angle of
one degree?) Furthermore, how do we know that it is possible to divide a degree
into 60 congruent angles (minutes) and a minute into successively smaller
pieces? This chapter shows how to resolve these difficulties. Incidentally, it
should become quite clear that what we choose as the \emph{unit angle} is quite
arbitrary, and it is important for you to understand how measurements differ
when different units are used. The idea that ties all this together is that of
\emph{length of circular arc}.

%========================================================= 12-2
\section*{12--2\quad CIRCULAR ARCS}
\addcontentsline{toc}{section}{12--2\quad Circular arcs}

We have already described the interior and exterior of an angle (Chapter 3),
and we use Definition 3--10 to describe an arc of a circle.

\begin{definition}
If $\ang{AOB}$ is a central angle of a circle, \emph{the arc $\XIIarc{AB}$ of
the circle} subtended by $\ang{AOB}$ is defined to be the set of points of the
circle which are \emph{interior} to $\ang{AOB}$ or are on the angle. The set of
points of the circle which are \emph{exterior} to or on the angle is also called
\emph{an arc of the circle} (for example, the arc $\XIIarc{ABC}$ in
Fig.~12--1). If the central angle is a straight angle (like $\ang{COB}$ in
Fig.~12--1), then it has no interior, but the set of points of the circle which
are on one side of $CB$ or on $CB$ is called a \emph{semicircular arc} and is
said to be subtended by $\ang{COB}$. Finally, when we speak simply of an
\emph{arc of a circle}, we mean an arc subtended by a central angle or the arc
exterior to some central angle.
\end{definition}

\FIGXIIONE

\exercisegroup{12--1}
\begin{multicols}{2}
\begin{exlist}
\item In Fig.~12--2, show that if $\ang{AOB} > \ang{COB}$, with $C$ in the
      interior of $\ang{AOB}$, then every point of the arc $\XIIarc{CB}$ is a
      point of the arc $\XIIarc{AB}$.
\item If the arc $\XIIarc{AB}$ is subtended by the central angle $AOB$, define
      the \emph{midpoint of the arc}, and show that such a point exists.
\item Use the result of Exercise 2 to prove that there are infinitely many
      points on the arc.
\item In Fig.~12--3, the arc $\XIIarc{EF}$ is subtended by $\ang{EOF}$. Consider
      the segment $EF$ and any point $P$ on $EF$. Show that there is, on the ray
      $OP$, a point $Q$, which is on the arc $\XIIarc{EF}$. Show that if
      $P \neq P'$, then $Q \neq Q'$. How many points are there on the arc
      $\XIIarc{EF}$?
\item Referring to Fig.~12--4, suppose that an arc is exterior to a central
      angle $AOB$. Show how to define the midpoint of the arc. Why is
      $\ang{AOC} \cong \ang{BOC}$?
\end{exlist}
\end{multicols}

\FIGXIITWOTHREE

\FIGXIIFOUR

%========================================================= 12-3
\section*{12--3\quad ARC LENGTH}
\addcontentsline{toc}{section}{12--3\quad Arc length}

In defining the length of a segment, we first had to choose a unit of length.
So also in defining arc length, it will be understood that we have already
chosen a segment of unit length. We shall also suppose that we have defined the
midpoint of an arc as in Exercise Group 12--1.

We describe how to construct a regular inscribed polygonal path in an arc, and
indeed, polygonal paths of 2, 4, 8, 16,\ldots\ sides.

In Fig.~12--5, let $AB$ be a circular arc, and let $C$ be the midpoint of the
arc. (See Exercise Group 12--1.) The segments $AC$ and $CB$ form a regular
inscribed polygonal path from $A$ to $B$, with two sides. There are now formed
two congruent angles, $\ang{AOC}$ and $\ang{COB}$. Each of these angles subtends
an arc. Let $D$ and $E$ be the midpoints of these arcs. We can now consider the
regular inscribed polygonal path of four sides, $ADCEB$. We then consider each
of the arcs $\XIIarc{AD}$, $\XIIarc{DC}$, $\XIIarc{CE}$, and $\XIIarc{EB}$ and
their midpoints $F$, $G$, $H$, and $I$. Thus we obtain a regular inscribed path
$AFDGCHEIB$ from $A$ to $B$, having eight sides. Continuing in this manner, for
each positive integer we obtain a regular inscribed polygonal path from $A$ to
$B$. That is to say, first we have a path of two sides, next a path of four
sides, then a path of eight sides, and so on.

\FIGXIIFIVESIX

The following facts are clear from Fig.~12--6, and a complete proof is not
difficult to give, although several steps are needed. We shall use these facts
without proof. The \emph{length} of the circumscribed path in Fig.~12--6 will be
denoted by $L$. The length of a regular inscribed path of 2, 4, 8,
16,\ldots, sides will be denoted by $l_n$, where $n$ is used to denote the
number of segments in the path.

(a) \emph{For every $n$, $l_n < L$}; that is to say, the length of each regular
inscribed path is less than the length of any circumscribed path.

(b) The length of each regular inscribed path of 2, 4, 8, 16,\ldots\ sides is
less than the one that follows, that is, $l_n < l_{2n}$.

(c) From (a) and (b) and properties of real numbers, it follows that
$\lim \{l_n\}$ exists.

We are now able to define length of arc.

\begin{definition}
Given a circular arc, let $l_n$ be the length of a regular inscribed polygonal
path of 2, 4, 8, 16, 32,\ldots, sides inscribed in the arc. The \emph{length $l$
of the arc} is defined to be the following limit: $l = \lim \{l_n\}$.
\end{definition}

\exercisegroup{12--2}
\begin{multicols}{2}
\begin{exlist}
\item A regular hexagon is inscribed in a circle of radius 1, as in
      Fig.~12--7. Show that the length of the arc $\XIIarc{ACB}$ is greater
      than 2.
\item In a circle of radius 2, what is the length of arc subtended by a right
      angle? by a straight angle?
\item In Fig.~12--8, the central angles $AOB$ and $BOC$ are congruent, and
      $\ang{AOC}$ is a right angle. What is the length of the arc $\XIIarc{AB}$
      if $\sg{AO} = 5$?
\item[\stex 4.] Prove that the length of any arc is a positive number which is
      greater than the length of the line segment determined by the endpoints of
      the arc.
\end{exlist}
\end{multicols}

\FIGXIISEVENEIGHT

\begin{theorem}
If, in a circle, two central angles are congruent, they subtend arcs of equal
length.
\end{theorem}

\noindent\emph{Proof.} Because the angles are congruent, the inscribed polygons
at any stage are also congruent. Thus the lengths of the corresponding
inscribed polygonal paths are the same. Hence the arc lengths are the same.

\begin{remark}
In the same way it follows that the arcs exterior to the central angles also
have the same length.
\end{remark}

Having defined arc length, we now would like to prove some more theorems about
arc lengths. Indeed, it would be possible to proceed at once to this task, but
if we were to do so, our proofs would be quite troublesome. In constructing
later proofs it will be most convenient for us to use a theorem (12--2 below)
which tells us that in approximating to the circle by inscribed polygons, we
need \emph{not} restrict ourselves to regular polygons. The proof of this
``helping'' theorem is rather hard to make and involves greater knowledge of
limits than is presupposed here. Nevertheless, the geometric significance of the
theorem should be clear, and it should appeal to your intuition. We shall use
this theorem without proof, so that actually we treat it like a postulate.

\begin{theorem}
Let $\{P_n\}$ be any sequence of polygonal paths inscribed in the arc
$\XIIarc{AB}$ of a circle (Fig.~12--9), and let $\{l_n\}$ be the sequence of
lengths of these polygonal paths; let $s_n$ be the length of the longest side of
the path $P_n$. If $\lim \{s_n\} = 0$, then $\lim \{l_n\} = $ length
$\XIIarc{AB} = l$.
\end{theorem}

\begin{remark}
Evidently, when the length of the longest side of $P_n$ is small, the number of
sides, $n$, must be large. We can now use this theorem to prove some facts about
arc length.
\end{remark}

\begin{theorem}
If $\XIIarc{AB}$ and $\XIIarc{BC}$ (Fig.~12--10) are adjacent arcs (that is, if
they have only point $B$ in common), then length $\XIIarc{ABC} = $ length
$\XIIarc{AB} + $ length $\XIIarc{BC}$.
\end{theorem}

\FIGXIININETEN

\noindent\emph{Proof.} The length of $\XIIarc{ABC}$ is, according to
Theorem 12--2, the limit of lengths of inscribed polygonal paths from $A$ to
$C$. Let us always choose these paths so that $B$ is one of the vertices of the
path. The path from $A$ to $C$ can then be separated into a path from $A$ to $B$
and a path from $B$ to $C$. Let $l'_n$ be the length of the path from $A$ to
$B$, and $l''_n$ be the length of the path from $B$ to $C$. Then, if $l_n$ is
the length of the path from $A$ to $C$, we have $l_n = l'_n + l''_n$.

Now we consider sequences of paths $\{P'_n\}$ and $\{P''_n\}$, with the
sequences of lengths of longest sides approaching zero. Thus,
\begin{equation}
\left.
\begin{aligned}
\lim \{l_n\} &= \text{length } \XIIarc{ABC},\\
\lim \{l'_n\} &= \text{length } \XIIarc{AB},\\
\lim \{l''_n\} &= \text{length } \XIIarc{BC}.
\end{aligned}
\right\}
\end{equation}
But from a theorem on limits, we have
\begin{equation}
\lim \{l'_n + l''_n\} = \lim \{l'_n\} + \lim \{l''_n\}.
\end{equation}
From Eqs.~(12--1) and (12--2), we see that
\begin{equation}
\text{length } \XIIarc{ABC} = \text{length } \XIIarc{AB} + \text{length }
\XIIarc{BC}.
\end{equation}
This completes the proof. The theorem might be stated briefly as ``the length of
the `sum' is equal to the sum of the lengths.''

\FIGXIIELEVENTWELVE

\begin{theorem}
If, in a circle, one central angle is greater than another central angle, the
greater angle subtends the longer arc.
\end{theorem}

\noindent\emph{Proof.} Since congruent angles subtend arcs of equal length, we
may suppose that the central angles are as in Fig.~12--11, and that
$\ang{AOB} > \ang{COB}$. The remainder of the proof (using Theorem 12--3) is
left to the student.

\exercisegroup{12--3}
\begin{multicols}{2}
\begin{exlist}
\item Apply Theorem 12--3 to the arcs subtended by a straight angle and two
      adjacent right angles.
\item State the form Theorem 12--3 would have for three adjacent arcs; for four
      adjacent arcs; for $n$ adjacent arcs.
\item Two adjacent arcs in a circle of radius $r$ have lengths $\pi r/3$ and
      $2\pi r/3$. What angle subtends the two arcs together?
\item Two arcs on a circle of radius $r$ have lengths $r$ and $\pi r/3$. Which
      is subtended by the greater central angle?
\end{exlist}
\end{multicols}

\refstepcounter{theorem}
\par\medskip\noindent\textbf{Theorem \thetheorem} Converse to Theorem 12--4.\
\emph{If one arc of a circle is longer than a second arc of the circle, the
longer arc is subtended by the greater central angle.}\par\medskip

The proof is left to the student.

\begin{theorem}
Two central angles of a circle are congruent if and only if they subtend arcs of
equal length.
\end{theorem}

The proof is left to the student.

%========================================================= 12-4
\section*{12--4\quad ARCS ON DIFFERENT CIRCLES}
\addcontentsline{toc}{section}{12--4\quad Arcs on different circles}

So far we have compared arcs on the same circle. We now compare the lengths of
arcs on different circles. Before studying Theorem 12--7, try to discover the
relation between the lengths of two arcs in different circles if they are
subtended by congruent angles.

\begin{theorem}
If $\XIIarc{AB}$ and $\XIIarc{A'B'}$ are arcs on two circles of radii $r$ and
$r'$, respectively, and are subtended by congruent angles, then the arc lengths
are proportional to the radii (Fig.~12--12). That is, if
$\ang{AOB} \cong \ang{A'O'B'}$, then
\[
\frac{\text{length } \XIIarc{AB}}{\text{length } \XIIarc{A'B'}} =
\frac{r}{r'}, \qquad \text{or} \qquad
\frac{\text{length } \XIIarc{AB}}{r} =
\frac{\text{length } \XIIarc{A'B'}}{r'}.
\]
\end{theorem}

\noindent\emph{Proof.} The proof depends upon three things: (1) the definition
of arc length, (2) a theorem on limits, and (3) properties of similar triangles.

Suppose that $APQB$ is a polygonal path inscribed in the arc $AB$. Then there is
a corresponding polygonal path $A'P'Q'B'$ inscribed in the arc $A'B'$ such that
corresponding central angles are congruent, that is,
$\ang{AOP} \cong \ang{A'O'P'}$, $\ang{POQ} \cong \ang{P'O'Q'}$, etc. (Why?)
Hence, there is a one-to-one correspondence between polygonal paths inscribed in
the two arcs. Let $L_n$ be the length of a polygonal path of $n$ sides inscribed
in $\XIIarc{AB}$, and $L'_n$ be the length of the corresponding polygonal path
inscribed in $\XIIarc{A'B'}$. Because the corresponding central angles in the
two polygons are congruent, for example $\ang{AOP} \cong \ang{A'O'P'}$, the
triangles formed are similar. (Why?) Since $\sg{AP}/r = \sg{A'P'}/r'$, etc., we
have $L_n/r = L'_n/r'$ and $L_n/L'_n = r/r'$.

From the definition of length of arc, length $\XIIarc{AB} = \lim \{L_n\}$ and
length $\XIIarc{A'B'} = \lim \{L'_n\}$. From a theorem on limits,
\[
\frac{\text{length } \XIIarc{AB}}{\text{length } \XIIarc{A'B'}} =
\frac{\lim \{L_n\}}{\lim \{L'_n\}} = \lim \left\{\frac{L_n}{L'_n}\right\} =
\frac{r}{r'}.
\]
This completes the proof.

\exercisegroup{12--4}
\begin{multicols}{2}
\begin{exlist}
\item A central angle $B$ subtends an arc of length 5 in.\ on a circle of radius
      4 in. A central angle $B' \cong \ang{B}$ subtends an arc of length 10 in.\
      on another circle. What is the radius of the second circle?
\item A central angle subtends an arc of 2 in.\ on a circle of radius 5 in. What
      would be the length of an arc subtended by this same angle on a circle
      with the same center but with a radius of 2 ft?
\item The radii of two circles have a ratio of 3 to 5. If a central angle
      subtends an arc of length 2.454 on the first circle, what length arc would
      a central angle that is congruent to the first angle subtend on the second
      circle?
\end{exlist}
\end{multicols}

%========================================================= 12-5
\section*{12--5\quad MORE ABOUT ARC LENGTH}
\addcontentsline{toc}{section}{12--5\quad More about arc length}

In the preceding sections we learned that every arc has associated with it a
number called its length. We shall not consider an exhaustive proof of the
converse theorem: \emph{For every positive number, less than the circumference,
there is an arc having this number as its length.} The complete proof of this
theorem requires a rather difficult argument using limits, so we will limit our
proof to certain special numbers. It will then be clear that the theorem is
true.

In proving the theorem, we confine ourselves to arcs of length not more than one
quarter of the circumference, that is to arcs which are subtended by acute
angles or right angles.

\begin{theorem}
Suppose that $L$ is a positive number less than or equal to one fourth of the
circumference of a circle of radius $r$ $(L \leqq \pi r/2)$.\XIIsymfoot{Since
the circumference of a circle is $2\pi r$, it follows that one fourth of the
circumference is $\frac{1}{4} \times 2\pi r$, or $\pi r/2$.} Then there is an
arc of the circle whose length is $L$.
\end{theorem}

\noindent``\emph{Proof}.'' If $L = \pi r/2$, then the required arc is subtended
by a right angle. If $L = \frac{1}{2} \cdot \pi r/2$, then bisecting a central
right angle will give two arcs $\XIIarc{AC}$ and $\XIIarc{BC}$, as in
Fig.~12--13, of length $L$. If $L = \frac{3}{4} \cdot \pi r/2$, then the arc
$AD$ of Fig.~12--13 has length $L$, where $OD$ bisects $\ang{BOC}$. Similarly,
if $OE$ bisects $\ang{DOC}$, then length $\XIIarc{AE} = \frac{5}{8} \cdot
\pi r/2$.

Thus we see that if $L$ is a certain special fractional part of $\pi r/2$, then
there is an arc of length $L$. (Try to describe these fractions.) The fractions
for which arcs are easy to construct are 1/2, 1/4, 3/4, 1/8, 3/8, 5/8, 7/8,
1/16,\ldots. In other words, the arc is easy to construct if its length is given
simply in terms of halves, quarters, eighths, etc., of a quarter of the
circumference.

\FIGXIITHIRTEENFOURTEEN

You are familiar with measuring lengths with a ruler whose scale is divided into
halves, quarters, eighths, sixteenths, etc. Hence, if the number $L$ is not
exactly expressible in halves, quarters, eighths, etc., we can still construct
an arc whose length is very close to $L$.

In Fig.~12--14, $L$ appears to be a little greater than $\frac{5}{8} \cdot
\pi r/2$. It is plausible therefore, that for every $L < \pi r/2$, there is
exactly one arc $AP$, as in the figure, with length $\XIIarc{AP} = L$.

\exercisegroup{12--5}
\begin{multicols}{2}
\begin{exlist}
\item In this section we gave a list of fractions whose associated arcs are easy
      to construct. The list is computed according to a definite rule. Continue
      the list to 20 more terms.
\item On a circle of radius $r$, construct arcs of length
      $7/8 \cdot \pi r/2$; $9/16 \cdot \pi r/2$; $7/4 \cdot \pi r/2$.
\item Suppose that there is an arc of length 3 on a circle of radius 2 in. What
      part (as a decimal) of $\pi r/2$ is 3? Show how to approximate this
      decimal by a sum of halves, quarters, eighths, etc. Use the approximation
      3.14 for $\pi$, and then rework using 3.1416.
\item An arc of length $r$ on a circle of radius $r$ is what part of the
      circumference? of one quarter of the circumference?
\item On a circle of radius 32 in.\ is an arc of length 22 in. What part of a
      right angle does it subtend? Use $\pi = 22/7$. Give a construction
      approximating this arc.
\end{exlist}
\end{multicols}

%========================================================= 12-6
\section*{12--6\quad MEASUREMENT OF ANGLES}
\addcontentsline{toc}{section}{12--6\quad Measurement of angles}

In the previous section we have in effect been measuring angles by measuring
arcs. In fact we have expressed the arc length as a fraction of $\pi r/2$ that
is as a fraction of one quarter of the circumference. In doing this we have
measured the corresponding central angle in terms of the right angle as a unit
of angle. In this section we shall define the measure of an angle in terms of
other units.

In deciding how to measure angles, we must first define the unit. This involves
choosing some angle whose measure is 1. The choice is quite arbitrary, but there
are some restrictions based on convenience. If the unit angle is small, the
measures of common angles will be large numbers; whereas if the unit angle is
large, the measures of common angles will be small numbers. It turns out that
the familiar degree is a convenient size unit for ordinary measurement, because
the numbers involved are a ``good'' size. Obviously, other unit angles could be
selected. A natural one would be obtained by dividing the circle into 100 arcs
of equal length, or perhaps the right angle into 100 congruent angles.

How does it happen that we use that rather strange unit, the degree? We measure
angles in degrees because the ancient Sumerians\XIIsymfoot{The Sumerians had a
highly developed civilization in Mesopotamia, in the valleys of the Tigris and
Euphrates rivers, more than 5000 years ago. Roughly contemporary with ancient
Egypt, the Sumerian culture discovered much about arithmetic, geometry, and
astronomy. We inherit many things from these gifted people. Their method of
counting was based on 60 rather than our base of 10. Thus we divide the hour
into 60 minutes, and the minute into 60 seconds. They also invented place value
in number notation, although they had no ``zero.'' Many of our words go back to
Sumerian roots. For more Sumerian history, see \emph{A History of Science}, by
George Sarton, Harvard University Press, 1952, or your encyclopedia.} (2000
B.C.\ and earlier) did. They divided the circle into 360 ``equal parts,''
perhaps because they thought at one time that the year was 360 days long!

\begin{definition}
On a circle of radius $r$, an arc of length $1/360 \times 2\pi r$ is subtended
by a central angle of one degree $(1^\circ)$.
\end{definition}

\begin{XIIremarkn}{1}
The definition of degree is independent of the radius $r$, by
Theorem 12--7.
\end{XIIremarkn}

\begin{XIIremarkn}{2}
Any angle congruent to an angle of $1^\circ$ is also an angle of
$1^\circ$, by Theorem 12--6.
\end{XIIremarkn}

\exercisegroup{12--6}
\begin{multicols}{2}
\begin{exlist}
\item Write out the proof of Remark 1.
\item Write out the proof of Remark 2.
\item Define ``minute'' and ``second.''
\item Define a new unit of angle, which we will call the \emph{centile}, by
      dividing one quarter of the circumference into 100 equal arcs. Using this
      unit, what is the measure of a $30^\circ$ angle? a $45^\circ$ angle? a
      straight angle?
\item Divide the circumference of a circle into two arcs of equal length. What
      angle is subtended by each arc? Each arc is what part of the
      circumference? What is the measure in degrees of the central angle that is
      subtended by each arc?
\item Divide the circumference of a circle into four arcs of equal length. What
      kind of an angle is subtended by each of these arcs? Each arc is what part
      of the circumference? What is the measure in degrees of each central
      angle?
\item Consider an arc whose length is one sixth of the circumference. What is
      the measure in degrees of the central angle that is subtended by this arc?
\item Consider an arc whose length is one sixth of the length of a semicircular
      arc. What is the measure in degrees of the central angle that is subtended
      by this arc?
\item If a semicircular arc has length 8 in., and a second arc has length 2 in.,
      what part is the length of the second arc of the first? What is the
      measure in degrees of the central angle that is subtended by the second
      arc?
\item On a circle of radius 4, what is the length of an arc subtended by an
      angle of $1^\circ$? of $2^\circ$? of $30^\circ$?
\end{exlist}
\end{multicols}

\begin{definition}
If, on a circle of radius $r$, a central angle subtends an arc of length $L$,
then the measure of the angle in degrees is $L/\pi r \times 180$.
\end{definition}

\begin{XIIremarkn}{1}
This definition agrees with Definition 12--3 for angles of $1^\circ$.
\end{XIIremarkn}

\begin{XIIremarkn}{2}
By Theorem 12--7, the measure of an angle in degrees is independent of the
radius.
\end{XIIremarkn}

\begin{XIIremarkn}{3}
If $M$ is the measure of the central angle in degrees, we have from
Definition 12--4 that $L = M\pi r/180$.
\end{XIIremarkn}

\exercisegroup{12--7}
\begin{multicols}{2}
\begin{exlist}
\item If $\ang{AOB}$ and $\ang{BOC}$ are adjacent angles of $1^\circ$ each, show
      that the measure of $\ang{AOC}$ is $2^\circ$.
\item Show that a straight angle has measure $180^\circ$.
\item On a circle of radius 6, find the number of degrees in a central angle
      that subtends an arc of length 3; 6; $\pi$; $2\pi/3$.
\item On a circle of radius 9, find the length of an arc subtended by an angle
      of $30^\circ$; $180^\circ$; $60^\circ$; $(180/\pi)^\circ$.
\item Show that bisecting a right angle gives an angle whose measure is
      $45^\circ$.
\item Prove Remark 3 under Definition 12--4.
\end{exlist}
\end{multicols}

%========================================================= 12-7
\section*{12--7\quad COUNTING DEGREES AND\\MEASURING ANGLES}
\addcontentsline{toc}{section}{12--7\quad Counting degrees and measuring angles}

Measurement of angles is quite similar to the counting process which led to the
idea of length of a segment. In this section we give a description of this
process and find that it gives the same measure of an angle as is given by
Definition 12--4. Detailed proofs will be omitted. The following theorem shows
how counting applies.

\begin{theorem}
Suppose that $\ang{AOB}$ is an acute angle, and we consider successive arcs
subtending central angles of $1^\circ$ on the circle from $B$ to $A$ so that the
point $A$ is a point of the $(n + 1)$st arc (Fig.~12--15). Then the measure of
$\ang{AOB}$ in degrees is either equal to $n$, or $(n + 1)$, or is between $n$
and $(n + 1)$; $n \leqq$ measure of $\ang{AOB} \leqq n + 1$.
\end{theorem}

\FIGXIIFIFTEEN

\noindent\emph{Proof.} The length of the arc of $n^\circ$ is
$n \cdot \pi r/180$; the length of the arc of $(n + 1)^\circ$ is
$[(n + 1)\pi r]/180$. The length $L$ of the arc $\XIIarc{AB}$ is then between
these two numbers, by Theorem 12--5.
\begin{equation}
n \cdot \frac{\pi r}{180} \leqq L \leqq \frac{(n + 1)\pi r}{180}.
\end{equation}
By Definition 12--4, the measure of $\ang{AOB}$ in degrees is
$L \cdot 180/\pi r$. From Eq.~(12--4), we see that
\[
n \leqq L \cdot \frac{180}{\pi r} \leqq (n + 1),
\]
and this is what we wished to prove.

\begin{remark}
The measure to the nearest minute could be obtained by counting in a similar
way.
\end{remark}

Consider an angle $A'O'B'$ (Fig.~12--16) that is congruent to $\ang{AOB}$ of
Theorem 12--9. Consider also successive arcs subtending central angles of
$1^\circ$ on arc $B'A'$, as described for arc $BA$. If the point $A$ is on the
$(n + 1)$st arc, then the point $A'$ is on the $(n + 1)$st arc. This can be
shown using the addition and subtraction of angles postulate and the definition
of ``less than'' and ``greater than'' for angles. This proves the following
theorem:

\begin{theorem}
If two angles are congruent, they have equal measures.
\end{theorem}

The converse of Theorem 12--10 can be proved, using the fact that
$\ang{O} > \ang{O'}$, or $\ang{O} \cong \ang{O'}$, or $\ang{O} < \ang{O'}$.

\FIGXIISIXTEEN

\exercisegroup{12--8}
\begin{center}Give your answers in terms of $\pi$.\end{center}
\nopagebreak
\begin{multicols}{2}
\begin{exlist}
\item If the radius of a circle is 14, what is the length of the arc subtended
      by an angle of $30^\circ$?
\item On a circle of radius 5 in.\ is an arc of length 6 in. What is the measure
      in degrees of the central angle subtending the arc? What if the radius is
      20 in.? 6 in.?
\item A central angle of a circle subtends an arc of length equal to the radius.
      How many degrees in the angle?
\item An arc of length 8 in.\ on a circle of radius 6 in.\ subtends a central
      angle of how many degrees? an arc of length 4 in.? an arc of length
      6 in.?
\item On a circle of radius $\pi$, there is an arc of length $\pi^2/4$. How many
      degrees in the central angle?
\item An arc of length $\pi r/2$ is subtended on a circle of radius $r$ by what
      central angle?
\item Find the length of arc on a circle of radius $r$ subtended by an angle of
      $1^\circ$.
\end{exlist}
\end{multicols}

We have been measuring arcs by their lengths. It is convenient to measure them
by the measure in degrees of the central angles they subtend.

\begin{definition}
The measure in degrees of an arc subtended by a central angle is equal to the
measure of the central angle. The measure in degrees of an arc exterior to a
central angle is equal to 360 minus the measure of the central angle. (See
Fig.~12--17.)
\end{definition}

\FIGXIISEVENTEENEIGHTEEN

\exercisegroup{12--9}
\begin{multicols}{2}
\begin{exlist}
\item An arc of length $\pi$ in.\ on a circle of radius 3 in.\ is of how many
      degrees?
\item An arc of length 4 in.\ on a circle of radius 6 in.\ has how many degrees?
\item On a circle of radius 2 in., how long is an arc of $120^\circ$?
      $240^\circ$?
\item How many degrees has an arc of length 15 in.\ on a circle of radius 3?
\end{exlist}
\end{multicols}

%========================================================= 12-8
\section*{12--8\quad RADIAN MEASURE}
\addcontentsline{toc}{section}{12--8\quad Radian measure}

Besides the familiar unit of degree for measuring angles, there is another one
which is of great importance in advanced mathematics. The reasons for using it
cannot be given here, but the definition is quite simple, and you will meet it
again if you study trigonometry. The unit is called the \emph{radian}.

\begin{definition}
An angle of one radian is one which subtends an arc of length $r$ on a circle of
radius $r$ (Fig.~12--18).
\end{definition}

\begin{definition}
The measure of an angle in radians is equal to $L/r$, where $L$ is the length of
the arc subtended on a circle of radius $r$.
\end{definition}

\begin{XIIremarkn}{1}
Definition 12--7 agrees with Definition 12--6 for angles of 1 rad.
\end{XIIremarkn}

\begin{XIIremarkn}{2}
The measure in radians does not depend on the radius, by Theorem 12--7.
\end{XIIremarkn}

\exercisegroup{12--10}
\begin{multicols}{2}
\begin{exlist}
\item On a circle of radius 2 is an arc of length 4. How many radians are there
      in the central angle?
\item How many radians are there in a right angle? in a straight angle?
\item On a circle of radius 2, find the length of arc subtended by a central
      angle of $30^\circ$. How many radians are there in this angle?
\item How many degrees are there in one radian?
\item How many radians are there in one degree?
\item From Exercises 4 and 5, derive the formulas for changing measure in
      degrees to measure in radians, and vice versa.
\item What are the measures in radians of the following angles? $30^\circ$;
      $45^\circ$; $60^\circ$; $90^\circ$; $120^\circ$; $135^\circ$;
      $150^\circ$; $180^\circ$.
\item The following numbers are measures of angles in radians. What are their
      measures in degrees? $\pi/4$; $5\pi/4$; 3; 2; 1.5; $\pi/3$; $\pi/6$;
      $\pi/2$.
\item Which is the larger unit, radian or degree? For which unit will the
      measure of an angle give the larger number?
\item Consider a central angle of 1 radian on a circle of radius $r$. What is
      the length of the arc subtended by this angle? an angle of 2 radians? an
      angle of 3 radians? an angle of $n$ radians?
\end{exlist}
\end{multicols}

\par\medskip\noindent\emph{Consequence of Definition} 12--7.\ If $R$ is the
measure in radians of a central angle in a circle of radius $r$, then the angle
subtends an arc of length $r \cdot R$.\par\medskip

%========================================================= 12-9
\section*{12--9\quad OTHER UNITS}
\addcontentsline{toc}{section}{12--9\quad Other units}

Units other than degrees and radians could be chosen. One which has some use in
artillery is the \emph{mil}. It is approximately the angle which subtends an arc
of 1 yd at a distance of 1000 yd. Apart from its military use, it is of little
importance.

%========================================================= 12-10
\section*{12--10\quad INSCRIBED ANGLES}
\addcontentsline{toc}{section}{12--10\quad Inscribed angles}

Now that we have a measure for angles, it is easy to prove a theorem about the
measure of angles \emph{inscribed} in a circle. We recall the definition.

An angle is said to be \emph{inscribed} in a circle if its vertex and one point
on each side of the angle, other than the vertex, lie on the circle. The arc of
the circle interior to the angle is said to be subtended by the angle. (See
Fig.~12--19.)

\begin{theorem}
The measure of an inscribed angle in degrees is equal to one half of the degree
measure of the arc it subtends.
\end{theorem}

\FIGXIININETEENTWENTY

First, we prove a special case of this theorem, which will enable us to prove
the theorem easily.

\par\medskip\noindent\textbf{Lemma 12--11--1.}\ \emph{If $\ang{ABC}$ is an
inscribed angle and $BC$ is a diameter of the circle, then measure
$\ang{B} = \frac{1}{2}$ measure $\XIIarc{AC}$.} (See Fig.~12--20.)\par\medskip

\noindent\emph{Proof.}
\begin{steps}
\step{$\ang{OAB} \cong \ang{B}$.}{Why?}
\step{Measure $\ang{OAB} + $ measure $\ang{OBA} = $ measure $\ang{COA}$.}{Why?}
\step{2 (measure $\ang{B}$) $= $ measure $\ang{COA}$.}{Why?}
\step{Measure $\ang{COA} = $ measure $\XIIarc{AC}$.}{Why?}
\step{Measure $\ang{B} = \frac{1}{2}$ (measure $\XIIarc{AC}$).}{Why?}
\end{steps}

\noindent\emph{Proof of the theorem.} There are two cases to consider to
complete the proof of the theorem. They are: \emph{Case} 1. The center of the
circle is interior to $\ang{ABC}$ (Fig.~12--21). \emph{Case} 2. The center of
the circle is exterior to $\ang{ABC}$ (Fig.~12--22).

\FIGXIITWENTYONETWENTYTWO

In Case 1 the diameter through $B$ divides $\XIIarc{AC}$ into arcs
$\XIIarc{AP}$ and $\XIIarc{PC}$, and $\ang{ABC}$ into angles $x$ and $y$, as
shown. By Lemma 12--11--1, measure $\ang{x} = \frac{1}{2}$ (measure
$\XIIarc{AP}$), and measure $\ang{y} = \frac{1}{2}$ (measure $\XIIarc{PC}$).
Hence measure $\ang{ABC} = $ measure $\ang{x} + $ measure
$\ang{y} = \frac{1}{2}$ (measure $\XIIarc{AP}$) $+ \frac{1}{2}$ (measure
$\XIIarc{PC}$) $= \frac{1}{2}$ (measure $\XIIarc{AC}$).

The proof of Case 2 is left to the student.

\exercisegroup{12--11}
\begin{multicols}{2}
\begin{exlist}
\item If $\ang{ABC}$ is inscribed in a circle and measure
      $\XIIarc{AC} = 80^\circ$, how many degrees in $\ang{B}$?
\item An angle of $100^\circ$ is inscribed in a circle. How many degrees are
      there in the arc it subtends?
\item Prove that an angle ``inscribed in a semicircle'' is a right angle
      (Fig.~12--23).
\item How large can an angle be that is inscribed in a circle?
\item Prove that two inscribed angles in one circle which subtend arcs of equal
      measure are congruent.
\item Prove that a parallelogram inscribed in a circle must be a rectangle. More
      generally, prove that if a quadrilateral is inscribed in a circle, then the
      opposite angles are supplementary.
\item In Fig.~12--24 lines $AB$ and $CD$ are parallel. Prove that measure
      $\XIIarc{AC} = $ measure $\XIIarc{BD}$.
\item A trapezoid is inscribed in a circle. Prove that the base angles are
      congruent. [\emph{Hint:} Use the result of Exercise 7.]
\item[\stex 9.] In Fig.~12--25, $AC$ is tangent to the circle. Prove that
      measure $\ang{ABD} = \frac{1}{2}$ measure $\XIIarc{BD}$. [\emph{Hint:}
      Draw $DE \pll AC$. Show that arcs $\XIIarc{BD}$ and $\XIIarc{BE}$ have the
      same measure. Then use $\ang{ABD} \cong \ang{BDE} \cong \ang{BED}$.]
\item[\stex 10.] Prove, for Fig.~12--26, that measure
      $\ang{A} = \frac{1}{2}$ (measure $\XIIarc{CD} - $ measure $\XIIarc{BE}$).
      State this theorem in words.
\item[\stex 11.] State and prove theorems similar to that of Exercise 10 for
      cases in which either $AC$ or $AD$, or both, are tangent to the circle.
\item[\stex 12.] Prove, for Fig.~12--27, that measure
      $\ang{AEB} = \frac{1}{2}$ (measure $\XIIarc{AB} + $ measure
      $\XIIarc{DC}$). [\emph{Hint:} Draw $BC$.] State this theorem in words.
\item[13.] Using Fig.~12--27, prove that $\tri{AED} \simto \tri{BEC}$. Show that
      $\sg{AE} \cdot \sg{EC} = \sg{BE} \cdot \sg{ED}$. State this last property
      as a theorem.
\item[\stex 14.] Using Fig.~12--26, prove that $\tri{AEC} \simto \tri{ABD}$.
      Hence prove that $\sg{AB} \cdot \sg{AC} = \sg{AE} \cdot \sg{AD}$. State
      this last property as a theorem.
\item[15.] Two chords of a circle intersect as in Fig.~12--28. If
      $\sg{AO} = 12$, $\sg{OB} = 4$, and $\sg{CO} = 8$, find $\sg{OD}$.
      [\emph{Hint:} Use Exercise 13.]
\item[\stex 16.] From point $A$ outside a circle, secants are drawn as in
      Fig.~12--29. If $\sg{AD} = 6$, $\sg{DE} = 10$, and $\sg{AB} = 8$, find
      $\sg{BC}$.
\item[\stex 17.] In Fig.~12--30, $AB$ is tangent to the circle at $B$. Prove
      that $\tri{ACB} \simto \tri{ABD}$. You will need the result of Exercise 9
      for this. Prove also that $\sg{AB}^2 = \sg{AD} \cdot \sg{AC}$.
\item[18.] If, in Fig.~12--30, $\sg{AC} = 8$ and $\sg{CD} = 10$, determine
      $\sg{AB}$.
\end{exlist}
\end{multicols}

\FIGXIITWENTYTHREE

\FIGXIITWENTYFOURTWENTYFIVE

\FIGXIITWENTYSIXTWENTYSEVEN

\FIGXIITWENTYEIGHT

\FIGXIITWENTYNINETHIRTY

%========================================================= review
\par\bigskip
\begin{center}\bfseries\large REVIEW OF CHAPTER 12\end{center}
\addcontentsline{toc}{section}{Review of Chapter 12}
\nopagebreak

In this chapter definitions have been given of the following terms; be sure that
you can state them in your own words: arc, subtended arc, inscribed polygonal
path, arc length, degree, measure of angle in degrees, measure of arc in
degrees, radian, measure of angle in radians, inscribed angle.

The following is an outline of the important definitions and theorems of the
chapter. Study the outline so that you understand clearly how each step leads to
the next.

\par\smallskip\noindent\emph{Theorem:} The limit of the lengths of a sequence of
regular inscribed polygons exists.

\noindent\emph{Definition:} Arc length.

\noindent\emph{Theorem:} Can use inscribed polygons other than regular ones to
define arc length.

\noindent\emph{Theorem:} Length of ``sum'' of two arcs equals sum of lengths.

\noindent\emph{Theorem:} Greater angle subtends longer arc.

\noindent\emph{Theorem:} Angles congruent if and only if arc lengths equal.

\noindent\emph{Theorem:} Arc length is proportional to radius.

\noindent\emph{Theorem:} For every positive number less than the circumference,
there is an arc having length equal to the number.

\noindent\emph{Definition:} Measure of angles in degrees.

\noindent\emph{Definition:} Measure of angles in radians.

\noindent\emph{Theorems:} Angle measurement and counting.

\noindent\emph{Theorem:} Measure of inscribed angle is one half measure of
subtended arc.

\par\medskip
\begin{center}\bfseries Important Remark\end{center}
\nopagebreak

The purpose of the chapter has been to define measurement of angles. To make the
definition, we first defined \emph{arc length}, and then \emph{angle
measurement} in terms of an arc length.

\par\medskip
\begin{center}\bfseries Review Exercises\end{center}
\nopagebreak
\begin{exlist}
\item From the definition of arc length, show that the length of arc subtended
      by a right angle (Fig.~12--31) is greater than $\sqrt{2}\,r$. Hence,
      conclude that $\pi/2 > \sqrt{2}$.
\end{exlist}

\FIGXIITHIRTYONE

\begin{exlist}\setcounter{exlisti}{1}
\item Angles of $15^\circ$ subtend arcs of lengths 10 and 12 on two circles.
      What is the ratio of the radii of the circles? Find the radii.
\item Theorem 12--7 states that arc length is proportional to the radius. What
      theorem about triangles is used in the proof of Theorem 12--7?
\item State briefly what is involved in proving that the length of the ``sum''
      of two adjacent arcs is equal to the sum of their lengths.
\item Compute the length of an arc in a circle of radius 10 if it is subtended
      by an angle of (a) 40 degrees, (b) 2 radians, (c) $\pi/4$ radians,
      (d) $\frac{1}{2}$ degree.
\item[\stex 6.] A sequence of points $A_1$, $A_2$, $A_3$, \ldots\ is marked on a
      circle so that each pair of consecutive points determines a central angle
      of 1 rad. The direction from $A_1$ to $A_2$ to $A_3$, etc., is
      counterclockwise. Compute the approximate length of the arc from $A_{51}$
      to $A_1$ in a counterclockwise sense if the radius of the circle is 10.
      Use $\pi = 3.1416$.
\end{exlist}
\clearpage
% ==================== ch13 ====================
\setcounter{theorem}{0}\setcounter{definition}{0}\setcounter{construction}{0}\setcounter{exgroup}{0}\setcounter{footnote}{0}
\setcounter{chapter}{13}

\chapter*{\raggedright\Huge LOCI AND SETS}
\markboth{LOCI AND SETS}{}
\addcontentsline{toc}{chapter}{13.\ Loci and Sets}
\vspace{-1em}\noindent\rule{\linewidth}{0.6pt}\vspace{1em}
\figurekey

%========================================================= 13-1
\section*{13--1\quad TERMINOLOGY}
\addcontentsline{toc}{section}{13--1\quad Terminology}

In this chapter we use the term \emph{locus}*\XIIIlocusnote{} (plural,
\emph{loci}), which we inherit from the days when Latin was the universal
language of science. (Even in the first part of the 19th century, most
mathematics was published in Latin.) As will be clear from the definition given
below, a locus is just a set of points. The language may be met in later work,
so we devote this short chapter to it. We also collect a number of geometric
constructions in this chapter.

\begin{definition}
A locus of points is the set of all points, and only those points, which
satisfy some given condition.
\end{definition}

We may now use the language, ``The locus of points $P$ such that\ldots.'' This
is just another way of saying, ``The set of all points $P$ such that\ldots.''

Examples:

(a) The \emph{locus} of points at a given distance $r$ from a given point is
obviously a circle.

(b) The \emph{locus} of points at a given distance $d$ from a given line $l$ is
a pair of lines parallel to $l$, one on each side of $l$ (Fig.~13--1).

\FIGXIIIONE

Example (b) requires a bit of proof, which is omitted. We will discuss such
proofs in the next section. This locus could also be described as the \emph{set}
of all points at distance $d$ from $l$.

\FIGXIIITWO

(c) The \emph{locus} of points equidistant from two parallel lines is a line
between the given lines.

An equivalent way of describing (c) would be to say: the \emph{set} of all
points equidistant from two parallel lines (Fig.~13--2). Again, the fact that
locus (c) is a line requires proof.

\exercisegroup{13--1}

In the following exercises, state what you believe the locus to be. Also
describe this locus as a \emph{set}.

\begin{multicols}{2}
\begin{exlist}
\item What is the locus of points $P$ such that $\sg{PP_0} = 4$, where $P_0$ is
      a given point?
\item What is the locus of all points at a distance 1 from a line segment three
      units long (Fig.~13--3)? (By the distance from a point to a line segment,
      we mean the shortest distance from the point to any point of the
      segment.)
\item What is the locus of all points at a distance four units from the sides of
      a unit square? a square 12 units on a side? eight units?
\item A circle of radius five units is given. What is the locus of all points at
      a distance two units from the given circle?
\item A triangle has a fixed base four units long (Fig.~13--4). What is the
      locus of all possible positions of the vertex $V$ if the area is four
      units? if the area is $A$ units?
\end{exlist}
\end{multicols}

\FIGXIIITHREE
\FIGXIIIFOUR

%========================================================= 13-2
\section*{13--2\quad SOME LOCUS THEOREMS}
\addcontentsline{toc}{section}{13--2\quad Some Locus Theorems}

In many locus problems the set of points can be characterized in more than one
way, or we can say that two locus statements may describe the same set of
points. When this is the case, there is the following to be shown: (1)~Each
point of the first locus is also a point of the second locus. (The first set is
a \emph{subset} of the second set.) (2)~Each point of the second locus is also a
point of the first locus. (The second set is a subset of the first set.) If each
is a subset of the other, the two sets are the same.

\FIGXIIIFIVE

\begin{theorem}
The locus of points $P$ equidistant from two given points, $A$ and $B$, is the
set of points on the perpendicular bisector of the segment $AB$ (Fig.~13--5).
\end{theorem}

To make the proof we must show two things: (1)~that if a point $P$ is
equidistant from $A$ and $B$, then it is on the perpendicular bisector, and
(2)~that if a point is on the perpendicular bisector, then it is equidistant
from $A$ and $B$. These should be easy for you to prove now; in fact, both (1)
and (2) were proven previously as theorems or exercises.

A brief proof will be sketched here, although your teacher may require you to
write out the details. To prove (1), let $C$ denote the midpoint of $AB$. Then
if $P \neq C$, $\tri{PAC} \cong \tri{PBC}$. (Why?) Thus $\ang{PCA} \cong
\ang{PCB}$, and $P$ is on the perpendicular bisector. (How do we argue if
$P = C$?) To prove (2), suppose that $P$ is on the perpendicular bisector and
$P \neq C$. Then $\tri{PCA} \cong \tri{PCB}$ (why?) and $PA \cong PB$. (How do
we argue if $P = C$?)

A similar theorem concerns angle bisectors.

\begin{theorem}
The locus of points interior to an angle and equidistant from the sides of the
angle is the set of points on the bisector of the angle, excluding the vertex
(Fig.~13--6).
\end{theorem}

\FIGXIIISIX

\noindent\emph{Proof sketch.} Suppose that $P$ is interior to the angle and
equidistant from the sides. Let $A$ and $B$ be the feet of perpendiculars from
$P$ to the sides of the angle. (How do you know that perpendiculars from $P$
meet the sides of the angle?) Then $\tri{POA}$ and $\tri{POB}$ are right
triangles with two congruent sides and are congruent. Hence $\ang{POA} \cong
\ang{POB}$, and $OP$ is the angle bisector. Conversely, if $P$ is on the
bisector, then $P \neq O$, $PA$ and $PB$ are perpendicular to the sides of the
angle, and the right triangles $POA$ and $POB$ are congruent. Hence
$\sg{PA} = \sg{PB}$.

\begin{theorem}
Let $AB$ be a given segment. The set of all points $P$ such that $\tri{APB}$ is
right, with $AB$ the hypotenuse, is the circle on $AB$ as diameter, excluding
points $A$ and $B$ (Fig.~13--7).
\end{theorem}

\FIGXIIISEVEN

Let $O$ be the midpoint of $AB$. It is necessary to show that if $P$ is a point
such that $AP \perp PB$, then $P$ is on the circle whose center is $O$ and which
contains $A$ and $B$.

Conversely, if $P$ is on the circle, and different from $A$ and $B$, it is
necessary to show that $AP \perp PB$. But this follows from theorems on the
measure of inscribed angles. (There are other proofs, too.) The details are left
to the student.

\exercisegroup{13--2}

\begin{multicols}{2}
\begin{exlist}
\item What is the set of all points that are centers of circles tangent to both
      sides of an angle?
\item What is the locus of all points exterior to a given circle and at a
      distance $d$ from the circle?
\item A barn has dimensions $30 \times 50$ ft. A cow is tethered to a corner of
      the barn by a rope 60 ft long. What is the locus of points where the cow
      may graze?
\item[\stex 4.] In Fig.~13--8, points $A$ and $B$ are fixed. What is the locus
      of all points $P$, on a given side of the line $AB$, such that
      $\ang{APB}$ is $50^\circ$?
\item[5.] What is the locus of points equidistant from three noncollinear
      points?
\item[6.] As in Fig.~13--9, sketch the path of a point on the rim of a wheel
      rolling on a line. The curve obtained is called a \emph{cycloid}.
\item[\stex 7.] With the same wheel as in Exercise 6, sketch the locus if $P$ is
      on a spoke of the wheel; on an extension of a spoke.
\item[8.] A small circle of radius $R/4$ rolls on the inside of a circle of
      radius $R$ (Fig.~13--10). What is the path of a point on the small circle?
      (The curve is called a \emph{hypocycloid}.)
\item[\stex 9.] What are the various possibilities that arise from Exercise 8 if
      the radius of the small circle is $xR$? Investigate the locus if
      $x = \frac{1}{2}$; $x = \frac{1}{3}$; $x$ is a rational number; $x$ is an
      irrational number.
\end{exlist}
\end{multicols}

\FIGXIIIEIGHT
\FIGXIIININE
\FIGXIIITEN

%========================================================= 13-3
\section*{13--3\quad UNIONS AND INTERSECTIONS}
\addcontentsline{toc}{section}{13--3\quad Unions and Intersections}

Some loci are just unions or intersections of other simpler loci. To be precise,
we make a definition.

\begin{definition}
Given two sets of points $S_1$ and $S_2$, the union of the two sets is the set of
all points of $S_1$ and $S_2$. The \emph{intersection} of $S_1$ and $S_2$ is the
set of points that are common to $S_1$ and $S_2$.
\end{definition}

\exercisegroup{13--3}

\begin{multicols}{2}
\begin{exlist}
\item Describe an angle as a union of two sets.
\item Describe an angle and its interior as the intersection of two sets.
\item Describe the interior of a triangle as the intersection of three sets; as
      the intersection of two sets.
\item Describe a triangle as the union of three sets.
\item Describe a quadrilateral as the union of four sets. Describe the interior
      of a simple convex quadrilateral as the intersection of two sets; as the
      intersection of four sets.
\end{exlist}
\end{multicols}

\begin{theorem}
The bisectors of the angles of a triangle intersect in one point (called the
\emph{incenter}), which is equidistant from the three sides (Fig.~13--11).
\end{theorem}

\FIGXIIIELEVEN

\noindent\emph{Proof sketch.} The bisectors of angles $A$ and $B$ must intersect
in one point $O$. (Why?) Then this point $O$ is equidistant from $AB$ and from
$BC$. (Why?) Likewise, $O$ is equidistant from $AB$ and $AC$. Therefore, $O$ is
on the bisector of $\ang{C}$. (Why?) This proves the theorem.

The point $O$ is called the \emph{incenter} because it is the center of a circle
\emph{inscribed} in the triangle. One may also state the theorem by saying that
\emph{the bisectors of the angles are concurrent}.

\FIGXIIITWELVE

\begin{theorem}
The perpendicular bisectors of the sides of a triangle intersect in one point
(called the \emph{circumcenter}), which is equidistant from the three vertices
(Fig.~13--12).
\end{theorem}

The proof is left to the student.

\exercisegroup{13--4}

\begin{multicols}{2}
\begin{exlist}
\item Consider all lines that are parallel to the base of a triangle and that
      intersect the other sides in points $C$ and $D$ (Fig.~13--13). Show that
      the locus of midpoints of such segments $CD$ is the median to the base
      (except for one point!).
\item Show how to inscribe a circle in a triangle. State why your construction
      is correct. Why is there a unique inscribed circle?
\item Show how to circumscribe a circle about a triangle. State why there is a
      unique circumscribing circle.
\item In Fig.~13--14, $\ang{ABC}$, neither right nor straight, is given. The
      point $O$ is the intersection of the perpendicular bisector of $AB$ and
      the line perpendicular to $BC$ through $B$. Prove that there is a circle
      with center $O$ containing $A$ and $B$. If $Q$ is a point of this circle
      and is on the same side of $AB$ as $O$, prove that
      $\ang{AQB} \cong \ang{ABC}$.
\end{exlist}
\end{multicols}

\FIGXIIITHIRTEEN
\FIGXIIIFOURTEEN

\begin{theorem}
The medians of a triangle intersect in one point (called the \emph{centroid}),
which is two thirds of the distance from each vertex to the midpoint of the
opposite side.
\end{theorem}

\noindent\emph{Outline of proof.} (Details are left as an exercise.) (1)~Any two
medians, say $CD$ and $AE$ in Fig.~13--15, intersect. Call this intersection
point $O$. (2)~Choose $F$ and $G$ as the midpoints of $OC$ and $OA$,
respectively. (3)~Show that $D$, $E$, $F$, and $G$ are vertices of a
parallelogram. (4)~Thus $AG \cong GO \cong OE$, and $FC \cong FO \cong OD$.
(5)~The point $O$ has, therefore, the required distance properties. Since there
is but one such point on each median, the medians are concurrent.

The centroid is the ``balance point'' in the following sense. A sheet of metal
of uniform thickness, the shape of the triangle, would balance if supported at
$O$. Exercise 1 of Group 13--4 may help you to see why this is so.

\FIGXIIIFIFTEENSIXTEEN

\begin{theorem}
The altitudes of a triangle intersect in one point (called the
\emph{orthocenter}).
\end{theorem}

\noindent\emph{Outline of proof.} (Details are left as an exercise.)
(1)~Through each vertex, consider the line that is parallel to the opposite side
of the given $\tri{ABC}$. These lines will intersect to form $\tri{A'B'C'}$, as
in Fig.~13--16. (2)~Show that $AC'BC$, $ABCB'$, and $ABA'C$ are parallelograms,
and hence observe that $C'B \cong BA'$, $C'A \cong AB'$, and $B'C \cong CA'$.
(3)~Show that the circumcenter $O$ of $\tri{A'B'C'}$ is the point of
intersection of the altitudes of $\tri{ABC}$.

\exercisegroup{13--5}

\begin{multicols}{2}
\begin{exlist}
\item Construct a triangle, given the lengths of its base, altitude, and one
      side.
\item Construct a triangle, given the base, the length of the altitude to this
      base, and one of the base angles.
\item[\stex 3.] As in Fig.~13--17, construct a triangle, given the lengths of
      two sides and an angle other than the included angle. [\emph{Hint:} Use
      Exercise 4 in Group 13--4.]
\item[4.] Construct a triangle, given the lengths of an altitude to one side and
      the other two sides.
\item[5.] Construct a triangle, given the midpoints of the three sides.
\item[6.] Construct a triangle, given the length of the altitude to the base and
      the two base angles.
\item[7.] Given a circle and a line segment $AB$, construct the locus of all
      points $P$ such that tangents of length $\sg{AB}$ can be drawn to the
      circle from $P$.
\item[8.] Construct a right triangle with a given segment as hypotenuse and
      having a side of given length.
\item[9.] Construct a circle tangent to a given line at a given point and having
      a given radius.
\item[10.] Construct the locus of all points $P$ such that there is a circle,
      with given radius and center at $P$, which is tangent to a given circle.
\item[11.] Construct all circles with a given radius which are tangent to two
      given circles of unequal radii.
\item[12.] Construct a line tangent to a given circle and perpendicular to a
      given line.
\item[13.] Given a circle and a line segment $AB$, construct the set of all
      points $M$ such that $M$ is the midpoint of a chord (of the circle)
      congruent to $AB$.
\item[14.] Given an arc of a circle, construct the center of the circle.
\item[15.] Construct an isosceles triangle, given the base, and the radius of
      the circumscribed circle.
\item[16.] Construct an isosceles triangle, given the base, and the radius of
      the inscribed circle.
\item[17.] Construct an isosceles triangle, given the base and the vertex angle.
\item[18.] Construct a triangle, given two sides and the length of the median to
      one of these sides.
\item[\stex 19.] Let $P$ be a given point not on a given circle whose center is
      $O$. Consider all segments $PX$ such that $X$ is on the circle. Show that
      the locus of midpoints $M$ of these segments is a circle whose diameter is
      a segment on the line through $P$ and $O$. Construct this circle.
\item[\stex 20.] Given three segments, construct a triangle two of whose sides
      are congruent to two of these segments, and such that the median to the
      other side is congruent to the third segment.
\end{exlist}
\end{multicols}

\FIGXIIISEVENTEEN

%========================================================= Review
\section*{REVIEW OF CHAPTER 13}
\addcontentsline{toc}{section}{Review of Chapter 13}
\markboth{REVIEW OF CHAPTER 13}{}

In earlier chapters we were using the phrases \emph{set of points} and
\emph{geometric figure} interchangeably. In this chapter we introduced a third
term, \emph{locus}, which has exactly the same meaning as the earlier
expressions.

It is important that you understand what is involved in establishing that two
different ``rules'' determine the same locus. A statement like

\begin{quote}
\emph{the set of points having property $p$ is the same set as that having
property $q$}
\end{quote}

\noindent must be established by proving two things:

(1) Every point in the first set is also in the second. That is, every point
having property $p$ also has property $q$.

(2) Every point in the second set is also in the first. That is, every point
having property $q$ also has property $p$.

% The book heads this "Review Exercises", with no "Exercise"/"Exercise Group"
% prefix, so \exercise{} (which prepends "Exercise") is wrong here.
\par\medskip\begin{center}\bfseries Review Exercises\end{center}\par\nopagebreak

\begin{exlist}
\item If $AB$ is a fixed segment, are the properties
      \begin{quote}
      $p$: $\ang{AXB}$ is acute
      \end{quote}
      \noindent and
      \begin{quote}
      $q$: $X$ is in the exterior of the circle on $AB$ as diameter
      \end{quote}
      \noindent equivalent to each other?
\item Restate Exercise 1 in conventional locus terminology.
\item If $R$ is the set of points $X$ such that, for a fixed segment $AB$,
      % As printed: the book sets the right-hand side of this first relation
      % WITHOUT an overbar, though the two relations below it carry one.
      % Verified at 320 dpi (PDF 249); "AB," opens the line, so nothing is
      % clipped.  Reproduced as printed -- see ch13-UNCERTAIN.md.
      $\sg{AX} + \sg{XB} = AB$, and $S$ is the set of points between $A$
      and $B$, prove that $R = S$. How else can you describe the locus of points
      $X$ such that $\sg{AX} - \sg{XB} = \sg{AB}$? such that
      $\sg{BX} - \sg{XA} = \sg{AB}$?
\item Let $AB$ and $CD$ be two given segments. Describe in a second way the set
      of points $X$ such that
      \begin{quote}
      $\ang{AXB}$ is obtuse and $\ang{CXD}$ is right;\\
      $\ang{AXB}$ is obtuse or $\ang{CXD}$ is right;\\
      $\ang{AXB}$ is acute and $\ang{CXD}$ is obtuse;\\
      $\ang{AXB}$ is acute or $\ang{CXD}$ is obtuse.
      \end{quote}
\end{exlist}
\clearpage
% ==================== ch14 ====================
\setcounter{theorem}{0}\setcounter{definition}{0}\setcounter{construction}{0}\setcounter{exgroup}{0}\setcounter{footnote}{0}
\setcounter{chapter}{14}

\chapter*{\raggedright\Huge SPACE GEOMETRY\textsuperscript{*}}
\addcontentsline{toc}{chapter}{14.\ Space Geometry}
\markboth{SPACE GEOMETRY}{}
\vspace{-1em}\noindent\rule{\linewidth}{0.6pt}\vspace{1em}
\figurekey

\XIVnote{*\,This chapter is divided into three parts. For Part A (Section
14--3), Chapters 1 through 3 are prerequisites. For Part B (Sections 14--4 and
14--5), Chapters 4 through 8 are prerequisites. For Part C (Sections 14--6
through 14--9), Chapters 9 through 12 are prerequisites.}

%========================================================= 14-1
\section*{14--1\quad INTRODUCTION}
\addcontentsline{toc}{section}{14--1\quad Introduction}

In this chapter we outline the geometry of space, that is, three-dimensional
Euclidean space. Sometimes this study is called \emph{solid} geometry. Just as
plane geometry is an abstraction derived from our experience with flat surfaces,
so solid geometry is an abstraction derived from our ordinary perception of
objects in space.

In this chapter you will find many figures and many theorems and definitions,
but very few proofs. Proofs are omitted because this chapter is intended to be
only an introduction to the subject. Nevertheless, if you study the theorems in
the order they occur, you will find that their arrangement is logical because
they can be proved in that order. Furthermore, most of them have fairly simple
proofs. To get an understanding of the \emph{facts} of solid geometry, it will
be sufficient to read over the theorems and be sure that you know what they say.

%========================================================= 14-2
\section*{14--2\quad GEOMETRIC OBJECTS}
\addcontentsline{toc}{section}{14--2\quad Geometric Objects}

As in plane geometry, we shall be concerned mostly with simple sets of points
that you have already encountered, namely, planes, spheres, pyramids, cubes,
prisms, etc., as in Fig.~14--1. More complicated sets of points, such as those
in Fig.~14--2, will occur to you. As before, any set of points may be studied.
The drawings in Fig.~14--2 represent sets of points that do not lie in a plane.

\FIGXIVONE

\FIGXIVTWO

\exercisegroup{14--1}
\begin{multicols}{2}
\begin{exlist}
\item Draw figures of several three-dimensional sets of points.
\item Sketch a polyhedron.
\item What might be meant by parallel planes? perpendicular planes?
\end{exlist}
\end{multicols}

\XIVpart{A.\quad LINES, PLANES, AND BETWEENNESS IN SPACE}
\addcontentsline{toc}{section}{A.\ Lines, Planes, and Betweenness in Space}

%========================================================= 14-3
\section*{14--3\quad THE POSTULATES}
\addcontentsline{toc}{section}{14--3\quad The Postulates}

It is perhaps somewhat surprising that most of the essential difficulties in
geometry already occur in the plane, so that extension to space geometry
requires no fundamentally new concepts. As a matter of fact, very few more
postulates are needed than for the plane, indeed only four more.

\begin{postulate}{VII--1}
There is a set of points called \emph{space}. Certain subsets of these points
are called planes. There are at least four points that are not in any one plane.
\end{postulate}

\begin{postulate}{VII--2}
There is exactly one plane containing any three noncollinear points
(Fig.~14--3).
\end{postulate}

\FIGXIVTHREE

\begin{postulate}{VII--3}
Each plane satisfies all the postulates for plane geometry. In particular, the
line joining two points of a plane lies in that plane. Furthermore, the
postulates for congruence of segments and angles apply whether or not the
segments and angles are in the same plane.
\end{postulate}

\begin{postulate}{VII--4}
Each plane separates space. By this we mean that the set of points of space not
on the plane consists of two subsets, called the sides of the plane, with the
following properties. If two points are in the same subset, then the segment
joining them does not meet the plane, but if the two points are in different
subsets, the segment joining them does intersect the plane.
\end{postulate}

Postulate VII--1 asserts that no plane fills up all of space, so that we have a
geometry of more than two dimensions (Fig.~14--4).

Postulate VII--2 gives us a simple requirement for fixing a definite plane.

Postulate VII--3 states that our planes are as we want them to be from our study
of plane geometry. In addition, this postulate relates congruent lines and
angles in different planes.

Postulate VII--4 expresses the familiar fact that a wall between neighbors
separates them.

\FIGXIVFOURFIVE

\exercise{14--2}
Restate Postulate Group VII in your own words.

From these postulates and the first two postulate groups for plane geometry, we
can prove a number of theorems concerning intersections of lines and planes.

\begin{theorem}
If a line $l$ does not lie in a plane $\pi$, and has at least one point in
$\pi$, then $l$ and $\pi$ have exactly one point in common (Fig.~14--5).
\end{theorem}

For the rest of the chapter, proofs will be left out, but we include a proof of
Theorem 14--1 as a guide to start you out correctly.

\noindent\emph{Proof.} There is, by assumption, at least one point $P$ in both
$l$ and $\pi$. Suppose that there were a second point $Q$. Then, by Postulate
VII--3, the line $l = PQ$ would lie in the plane $\pi$. But we have assumed that
$l$ is not in $\pi$. Hence $P$ is the only point in both $l$ and $\pi$.

\FIGXIVSIXSEVEN

\begin{theorem}
There is a unique plane containing a given line and a point not on the line
(Fig.~14--6).
\end{theorem}

\begin{theorem}
There are infinitely many planes containing a given line.
\end{theorem}

\noindent Figure 14--7 may suggest a proof.

\begin{theorem}
There is a unique plane containing two intersecting lines (Fig.~14--8).
\end{theorem}

\begin{theorem}
If two distinct planes have at least one point in common, they have precisely a
line in common (Fig.~14--9).
\end{theorem}

\FIGXIVEIGHTNINE

\begin{definition}
If two planes $\pi_1$ and $\pi_2$ intersect in a line $l$ (Fig.~14--10), let
$\pi_1^+$ and $\pi_2^+$ be half planes of $\pi_1$ and $\pi_2$, determined by
$l$. The set of all points in $\pi_1^+$ and $\pi_2^+$ is called a
\emph{dihedral angle}. The line $l$ is called the \emph{edge}, and $\pi_1^+$ and
$\pi_2^+$ are called the \emph{faces}, of the dihedral angle.
\end{definition}

\FIGXIVTEN

\begin{definition}
The interior of the dihedral angle with faces $\pi_1^+$ and $\pi_2^+$ is the set
of all points of space on the same side of $\pi_2$ as $\pi_1^+$ and on the same
side of $\pi_1$ as $\pi_2^+$.
\end{definition}

\exercisegroup{14--3}
\begin{multicols}{2}
\begin{exlist}
\item Give a precise definition for a tetrahedron (Fig.~14--11). The
      tetrahedron is usually considered to be just the surface. What do you
      define the interior to be?
\item Give a definition for a polyhedron (surface). See Fig.~14--12.
\item Give a precise definition for a half plane.
\end{exlist}
\end{multicols}

\FIGXIVELEVENTWELVE

\XIVpart{B.\quad CONGRUENCE IN SPACE}
\addcontentsline{toc}{section}{B.\ Congruence in Space}

%========================================================= 14-4
\section*{14--4\quad PERPENDICULAR LINES\\AND PLANES}
\addcontentsline{toc}{section}{14--4\quad Perpendicular Lines and Planes}

Before giving some theorems, we observe that all right angles are congruent
regardless of whether they are in the same plane. We know this because our
postulates for angle congruence state that, given an $\ang{AOB}$ in a plane
$\pi$ and a ray $r'$ from $O'$ in a plane $\pi'$, there is exactly one ray $r''$
from $O'$ (on a given side of the line of $r'$ in the plane $\pi'$) that forms
with $r'$ an angle congruent to $\ang{AOB}$. See Fig.~14--13.

\FIGXIVTHIRTEENFOURTEEN

\begin{definition}
A line $l$ is perpendicular to a plane $\pi$ if $l$ intersects $\pi$ at a point
$P$ and $l$ is perpendicular to every line in $\pi$ that passes through $P$
(Fig.~14--14).
\end{definition}

\begin{theorem}
If a line $l$ is perpendicular to two intersecting lines, then it is
perpendicular to the plane containing the lines (Fig.~14--15).
\end{theorem}

\noindent[\emph{Hint:} Choose $A$ and $A'$ such that $AP \cong A'P$. Choose $B$
and $C$. Then prove that $\tri{ABC} \cong \tri{A'BC}$. Then prove that
$\tri{ABQ} \cong \tri{A'BQ}$, so that $AQ \cong A'Q$. This will imply that line
$m \perp l$.]

\FIGXIVFIFTEENSIXTEEN

\begin{theorem}
There is exactly one plane perpendicular to a line at a point of the line.
\end{theorem}

\noindent See Fig.~14--16(a).

\noindent[\emph{Hint:} Use the previous theorem to show that there is at least
one plane. To prove there is only one, it is sufficient to show that if $l_1$,
$l_2$, and $l_3$ in Fig.~14--16(b) are distinct lines perpendicular to $l$ at
$P$, then they are in a single plane. Thus, the plane of $l_2$ and $l_3$ meets
the plane of $l$ and $l_1$ in a line $l' \perp l$. Since there is only one line
in the plane of $l$ and $l_1$ perpendicular to $l$ at $P$, we have $l_1 = l'$.
Hence $l_1$, $l_2$, and $l_3$ are coplanar.]

\begin{theorem}
There is exactly one plane perpendicular to a line and containing a point not on
the line (Fig.~14--17).
\end{theorem}

\exercisegroup{14--4}
\begin{multicols}{2}
\begin{exlist}
\item State two theorems, analogous to Theorems 14--7 and 14--8, concerning
      existence of lines perpendicular to planes.
\item Prove the theorems that you stated in Exercise 1.
\end{exlist}
\end{multicols}

\begin{theorem}
If two lines are perpendicular to the same plane, they are parallel
(Fig.~14--18).
\end{theorem}

\noindent[\emph{Hint:} First one must show that the two lines are in one plane.]

\FIGXIVSEVENTEENEIGHTEEN

\begin{definition}
Two lines are called \emph{skew} if they do not lie in one plane. Two planes are
called \emph{parallel} if they do not intersect.
\end{definition}

\begin{XIVremarks}
Two lines in one plane either intersect or are parallel. Two lines in parallel
planes need not be parallel. See Fig.~14--19.
\end{XIVremarks}

\FIGXIVNINETEEN

\begin{theorem}
Through a point $P$ not in a plane $\pi$, there is exactly one plane parallel to
$\pi$ (Fig.~14--20).
\end{theorem}

\noindent[\emph{Hint:} Choose intersecting lines $l_1$ and $l_2$ in $\pi$. Then
$P$, with these lines, determines two planes $\pi_1$ and $\pi_2$, in which there
are unique lines $l_1'$ and $l_2'$ parallel to $l_1$ and $l_2$, respectively, as
in Fig.~14--21. These lines determine a plane parallel to $\pi$. The uniqueness
is yet to be proved.]

\begin{theorem}
Two planes perpendicular to the same line are parallel (Fig.~14--22).
\end{theorem}

\FIGXIVTWENTYTWENTYTWO

\begin{theorem}
If two lines are cut by three parallel planes, the intercepted segments have
lengths that are proportional. In Fig.~14--23, we would have
$\sg{AB}/\sg{A'B} = \sg{BC}/\sg{B'C'}$.
\end{theorem}

\FIGXIVTWENTYTHREE

\exercisegroup{14--5}
Prove the following theorems.
\begin{multicols}{2}
\begin{exlist}
\item Two planes parallel to one plane are parallel to each other.
\item If a line intersects one of two parallel planes, it intersects the other.
\item If a line is perpendicular to one of two parallel planes, it is
      perpendicular to the other.
\item If a plane intersects one of two parallel planes, it intersects the other.
\item If one of two parallel lines is perpendicular to a plane $\pi$, then the
      other line is perpendicular to $\pi$.
\end{exlist}
\end{multicols}

%========================================================= 14-5
\section*{14--5\quad POLYHEDRA AND\\POLYHEDRAL ANGLES}
\addcontentsline{toc}{section}{14--5\quad Polyhedra and Polyhedral Angles}

\begin{definition}
Let $P_1P_2\ldots P_n$ be a simple closed polygon in a plane $\pi$, and let $V$
be a point not in $\pi$, as in Fig.~14--24. The set of points $Q$ on rays from
$V$ through all the points $P$ of the polygon is called a polyhedral angle. $V$
is the \emph{vertex}, and $VP_1$, $VP_2, \ldots VP_n$ are called \emph{edges}.
The planes $VP_1P_2$, $VP_2P_3, \ldots$ contain the \emph{faces} of the angle.
The angles $\ang{P_1VP_2}$, $\ang{P_2VP_3}, \ldots$ are called \emph{face
angles}. The polyhedral angle is called convex if the polygon $P_1 \ldots P_n$
is convex.
\end{definition}

\exercisegroup{14--6}
\begin{multicols}{2}
\begin{exlist}
\item Draw a trihedral angle and a tetrahedral angle. These have three and four
      faces, respectively.
\item Can a trihedral angle have three right-angled face angles?
\end{exlist}
\end{multicols}

\begin{theorem}
The sum\textsuperscript{*} of two face angles of a trihedral angle is greater
than the third face angle (Fig.~14--25).
\end{theorem}

\XIVnote{*\,The word \emph{sum} is not defined, but the idea is clear.}

\noindent[\emph{Hint:} It suffices to prove the theorem in the case in which the
third angle is the largest. Suppose that that angle is $\ang{AVB}$, as in
Fig.~14--25. Choose $C$ such that $\ang{CAV} \cong \ang{BAV}$. Choose $D$ on
$AB$ such that $AC \cong AD$. Now show that $VC \cong VD$ and $CB > DB$. Hence
conclude that $\ang{CVB} > \ang{DVB}$.]

\FIGXIVTWENTYFOURTWENTYFIVE

\begin{theorem}
The sum of the face angles of a convex polyhedral angle is less than four right
angles.
\end{theorem}

\exercise{14--7}
Illustrate the truth of Theorem 14--14 by constructing a polyhedral angle
(Fig.~14--26) of paper and observing that when it is ``flattened out'' in a
plane, it does not ``encircle'' a point.

\FIGXIVTWENTYSIXTWENTYSEVEN

\begin{definition}
A \emph{polygonal cell} (Fig.~14--27) is the sets of points on, or interior to,
a simple plane polygon $P_1P_2\ldots P_{n-1}$. The edges of the cell are the
sides $P_1P_2$, $P_2P_3, \ldots$ of the polygon.
\end{definition}

A \emph{polyhedron} is the set of points on a finite number of polygonal cells
joined together in the following way:

\begin{enumerate}[label=(\arabic*), leftmargin=2.1em, itemsep=2pt, topsep=3pt]
\item Any two cells have either no points in common, exactly one vertex in
      common, or exactly one edge in common.
\item Every edge is on precisely two cells.
\end{enumerate}

\exercise{14--8}
With Scotch tape and cardboard, construct a number of polyhedra.

\begin{theorem}
A polyhedron separates space in the following sense. The points not on the
polyhedron consist of two sets (called the inside and the outside) such that a
polygonal path joining an inside point and an outside point must intersect the
polyhedron, and two inside (or two outside) points can be joined by a polygonal
path not intersecting the polyhedron.
\end{theorem}

\exercise{14--9}
Describe the interior of a tetrahedron.

\begin{definition}
A polyhedron is called \emph{convex} if the segment joining any two points of
the polyhedron lies either in the polyhedron or its interior. A polyhedron is
\emph{regular} if it is convex and all faces are congruent regular polygons.
\end{definition}

\exercise{14--10}
With Scotch tape and cardboard, construct regular polyhedra with faces that are
triangles; squares; pentagons.

\begin{theorem}
There are but five (except for size) regular polyhedra (Fig.~14--28).
\end{theorem}

\FIGXIVTWENTYEIGHT

\XIVpart{C.\quad MEASUREMENT IN SPACE}
\addcontentsline{toc}{section}{C.\ Measurement in Space}

%========================================================= 14-6
\section*{14--6\quad ANGLES}
\addcontentsline{toc}{section}{14--6\quad Angles}

\begin{theorem}
Parallel planes intercept a dihedral angle in congruent angles (Fig.~14--29).
\end{theorem}

\begin{definition}
The measure of a dihedral angle is equal to the measure of the angle intercepted
by a plane perpendicular to the edge of the dihedral angle (Fig.~14--30).
\end{definition}

\exercise{14--11}
Restate Theorems 14--13 and 14--14, using degrees for measuring angles.

\FIGXIVTWENTYNINETHIRTY

%========================================================= 14-7
\section*{14--7\quad AREA}
\addcontentsline{toc}{section}{14--7\quad Area}

\begin{definition}
The area of a polyhedron is the sum of the areas of its faces.
\end{definition}

\exercisegroup{14--12}
\begin{multicols}{2}
\begin{exlist}
\item A \emph{prism} is a polyhedron two of whose faces are congruent polygons
      in parallel planes, called bases, and whose remaining edges are all
      parallel to one another, as in Fig.~14--31. The lateral area is the total
      area of the faces less the area of the bases. Find a formula for the
      lateral area of a prism.
\item Using Fig.~14--32, define \emph{pyramid} and \emph{regular pyramid}. Find
      a formula for the lateral area of a regular pyramid.
\item A \emph{parallelepiped} is a special kind of prism. Give a definition that
      you think fits the drawing in Fig.~14--33. Prove that the opposite faces
      of a parallelepiped are congruent.
\end{exlist}
\end{multicols}

\FIGXIVTHIRTYONETHIRTYTHREE

%========================================================= 14-8
\section*{14--8\quad VOLUME}
\addcontentsline{toc}{section}{14--8\quad Volume}

Just as area in the plane offers special difficulties, so does volume in space.
We content ourselves with listing the theorems and sketching a few ``proofs,''
obtained by drawing figures. Proofs that involve the notion of limit are not
presented.

\begin{theorem}
The volume of a rectangular parallelepiped (box) is equal to the product of the
lengths of three concurrent edges. Volume $=lwh$.
\end{theorem}

\noindent See Fig.~14--34.

\FIGXIVTHIRTYFOURTHIRTYFIVE

\begin{theorem}
If a polyhedron is the union of two other polyhedra, with only a face and no
interior points in common, then its volume is the sum of the volumes of the two
polyhedra forming it (Fig.~14--35).
\end{theorem}

\exercise{14--13}
A parallelepiped is shown in Fig.~14--36. Show, by ``chopping up'' the
parallelepiped to form a box, that the following theorem is true.

\FIGXIVTHIRTYSIXTHIRTYSEVEN

\begin{theorem}
The volume of a parallelepiped is equal to the area of a base multiplied by the
corresponding altitude (Fig.~14--37).
\end{theorem}

The same formula holds for the volume of a prism. To see this, we first prove
the theorem for prisms with triangular bases.

\exercisegroup{14--14}
\begin{multicols}{2}
\begin{exlist}
\item By constructing a parallelepiped of which a triangular prism is half, show
      that the volume of the prism is equal to the area of a base times the
      altitude. See Fig.~14--38(a).
\item Show that the same result holds for an arbitrary prism. See
      Fig.~14--38(b).
\end{exlist}
\end{multicols}

\FIGXIVTHIRTYEIGHT

\begin{theorem}
The volume of a prism is equal to the area of a base times the altitude.
\end{theorem}

Volumes of pyramids are somewhat harder to discuss. The following theorem is to
be taken without proof.

\begin{theorem}
If $B$ is the area of the base of a pyramid, then the volume $=
\frac{1}{3}Bh$, where $h$ is the altitude (Fig.~14--39).
\end{theorem}

\FIGXIVTHIRTYNINE

\exercisegroup{14--15}
\begin{multicols}{2}
\begin{exlist}
\item Find the volume of a regular pyramid with a base 4 ft square and altitude
      equal to 4 ft.
\item Find the volume of a regular tetrahedron whose side has length $a$.
\item Show that if the formula for volume of a pyramid is correct for pyramids
      with triangular bases, then it is correct for all pyramids.
\item Define \emph{similar polyhedra}, and prove that the volumes of similar
      polyhedra are proportional to the cubes of corresponding edges. Thus,
      $V/V' = l^3/l'^3$, as in Fig.~14--40.
\end{exlist}
\end{multicols}

\FIGXIVFORTY

%========================================================= 14-9
\section*{14--9\quad CYLINDERS, CONES, AND SPHERES}
\addcontentsline{toc}{section}{14--9\quad Cylinders, Cones, and Spheres}

The familiar shapes in Fig.~14--41 are analogous to prisms, pyramids, and
regular polyhedra.

The lateral surface area of cylinders and cones is obtained as a limit of the
areas of prisms and pyramids.

\FIGXIVFORTYONE

\exercisegroup{14--16}
\begin{multicols}{2}
\begin{exlist}
\item Give precise definitions for cylinders, cones, and spheres.
\item Find a formula for the lateral surface of a ``regular'' cone, called a
      \emph{right circular cone} (Fig.~14--42).
\end{exlist}
\end{multicols}

The surface area of a sphere is harder to discuss. We merely state the result.

\FIGXIVFORTYTWOFORTYTHREE

\begin{theorem}
The area of the surface of a sphere of radius $r$ is $4\pi r^2$.
\end{theorem}

With the aid of Theorem 14--23, it is not too hard to arrive at the volume of a
sphere. Try to prove the following theorem by considering small ``pyramids,'' as
shown in Fig.~14--43.

\begin{theorem}
The volume of a sphere of radius $r$ is $\frac{4}{3}\pi r^3$.
\end{theorem}

\exercisegroup{14--17}
\begin{multicols}{2}
\begin{exlist}
\item Show that if the radius of a sphere is doubled, the volume is multiplied
      by a factor of 8.
\item Oranges 2 in.\ in diameter are 50\textcent{} per dozen. Neglecting skin
      thickness, how much should one pay for oranges 3 in.\ in diameter?
\item Prove that if a plane intersects a sphere, then it intersects the sphere
      in either a circle or one point.
\item What is the formula for the lateral area of a cylinder?
\end{exlist}
\end{multicols}
\clearpage
% ==================== ch15 ====================
\setcounter{theorem}{0}\setcounter{definition}{0}\setcounter{construction}{0}\setcounter{exgroup}{0}\setcounter{footnote}{0}
\setcounter{chapter}{15}

% Facing-page plate for the chapter opener (book p.250 = PDF 264, facing
% p.251). Halftone photograph, never redrawn -- same treatment and the same
% ownership rule as ch01's Pythagoras/Archimedes and ch02's Aristotle plates:
% the chapter whose opener the plate faces emits it. Previously orphaned:
% ch14-UNCERTAIN handed it to ch15, ch15-UNCERTAIN handed it back to ch14, so
% no unit emitted it at all.
\XVplate{Ren\'e Descartes (1596--1650)}{Courtesy of \emph{Scripta Mathematica}.}

\chapter*{\raggedright\Huge ANALYTIC GEOMETRY}
\addcontentsline{toc}{chapter}{15.\ Analytic Geometry}
\markboth{ANALYTIC GEOMETRY}{}
\vspace{-1em}\noindent\rule{\linewidth}{0.6pt}\vspace{1em}
\figurekey

%========================================================= 15-1
\section*{15--1\quad SOME HISTORY}
\addcontentsline{toc}{section}{15--1\quad Some history}

We have commented before on the fact that numbers, other than integers, did
not occur explicitly in Greek geometry. One reason for this was that numbers
were hard to work with at that time. Even after the introduction, during the
Middle Ages, of Arabic numerals and the place value of the decimal system,
geometry continued in the Greek spirit. As a result, it often happened that
problems that we would handle naturally with algebra were solved by geometric
constructions.

Of course, as time passed, there was a gradual tendency toward the use of
algebra in geometric problems, but there was nothing systematic in its use.
And until there was sufficient algebraic machinery available and people became
skilled in using numbers, there was little advantage in changing a geometric
problem into an algebraic problem.

By the beginning of the 17th century, considerable progress had been made in
algebra and in algebraic notation;* in particular, quadratic, cubic, and
quartic equations could be solved. Mathematicians were interested in studying
tangents to curves, which is more easily done with a little analytic geometry.
So it was that in the 17th century the times were right for something new.
Indeed, many men had actually been using some analytic geometry without making
it a definite system for attacking problems in geometry. Today, looking back
on those times, it is hard to imagine why some ideas were so slow in coming.
But we must remember that discovery is always much harder than understanding
what has been discovered.%
\XVnote{*The beginning student hardly realizes how much easier the shorthand
notation of algebra makes his problems. The interested student is urged to
read in the history of mathematics and to try working a problem in earlier
notations.}

It is universally agreed that the French philosopher Ren\'e Descartes
(1596--1650) was the discoverer (inventor?) of analytic geometry, and the date
is generally set as 1637, when his \emph{Geometry} was published. Descartes is
also noted for his work in physics and in philosophy. In the latter, he
aspired to write a philosophy that would be perfectly logical and free from
objectionable assumptions. His efforts in philosophy include the book
\emph{Discourse on Method}, which expounds what is referred to today as
Cartesian philosophy, and which has been a source of dispute ever since! But
his \emph{Geometry} is undisputed. Analytic geometry is today a part of a
liberal education as well as a necessity for pure and applied science.

Some brief remarks about Descartes' life may be of interest. Born in 1596, he
started school in a Jesuit institution at the age of eight. Never a strong,
healthy boy, he was allowed by his school to stay in bed mornings to do his
studies, and later in life he did much of his writing in bed. He left school
in 1612, studied law, and even spent some time in the Dutch army as a
gentleman volunteer, leaving in 1619. It appears that by 1620 most of his
great ideas had been established in outline. In other words, he had thought
them over and had a fairly clear picture of how they should go. Much of his
life thereafter was spent in Holland. A devout man, he nevertheless had
troubles with the church over his philosophical ideas. By the time his
\emph{Geometry} was published, he was a well-known figure. In 1649 he visited
the court of Queen Christina of Sweden, after much urging, to teach her his
system of philosophy. Assigned the outlandish hour of 5 a.m. three days a week
for the instructions, he fell ill in February 1650 and died a few days later.

\emph{Geometry} was his only book on mathematics and, in a sense, contained
nothing totally new. What Descartes did was to show how a systematic use of
coordinates helps to attack problems. He classified various kinds of curves
and suggested further work. The book was suggestive and tremendously fruitful.
We owe to him our convention of using the first letters of the alphabet
($a$, $b$, $c$) for knowns and the last letters ($x$, $y$, $z$) for unknowns.

It is not easy to give a concise definition of analytic geometry. In a sense,
this entire chapter is such a definition, for a complete definition should
include some of its uses. But for the moment, it may suffice to say that, in
analytic geometry, geometric figures are studied by means of equations. Then,
if we wish to know something concerning the geometric figure, it \emph{may} be
easy to learn this by examining an equation. The equations that are used
involve numbers called coordinates, which serve to locate points. But why the
word ``analytic''? Analytic geometry does enable one to analyze geometric
figures, but this is also done in regular Euclidean geometry. Perhaps a better
description would be to call this new way of looking at geometry, plane
\emph{coordinate} geometry, or plane \emph{Cartesian} geometry. Remember that
it is the same plane we are studying all the time, no matter how we look at it.

\exercisegroup{15--1}
\begin{multicols}{2}
\begin{exlist}
\item Read the history of mathematics during the 17th century.
\item Read a biography of Descartes in an encyclopedia or in \emph{Men of
  Mathematics}, by E. T. Bell.
\end{exlist}
\end{multicols}

%========================================================= 15-2
\section*{15--2\quad EXAMPLES}
\addcontentsline{toc}{section}{15--2\quad Examples}

Although analytic geometry is distinguished by use of coordinates, the idea
should already be familiar from some common situations. Some cities are laid
out in squares, with the east-west and north-south roads called avenues and
streets, respectively. The numbers of street and avenue fix any intersection.
In the same way, you unconsciously use coordinates in describing a nearby town
as, perhaps, 7 miles north and 5 miles west of ``here.''

\exercisegroup{15--2}
\begin{multicols}{2}
\begin{exlist}
\item Try to think of other examples of coordinates, that is situations in
  which position is determined by a pair of numbers.
\item Discuss the need to have streets straight in order to ``coordinatize'' a
  city.
\end{exlist}
\end{multicols}

By far the most important example of coordinates is our method (Chapter 5) of
putting coordinates on a line. We recall how this was done:

Given a line $l$ and a unit of length, choose a point $O$ on the line $l$, as
in Fig.~15--1. Now choose* a point $P_1$ on $l$ at a distance 1 from $O$. This
choice determines which half of the line we will call the \emph{positive
axis}, or \emph{positive direction}. For every positive real number $x$, there
is exactly one point $P_x$ of $l$ on the same side of $O$ as the point $P_1$
such that the length $\sg{OP_x} = x$. For every negative real number $x'$,
there is exactly one point $P_{x'}$ of $l$ on the opposite side of $O$ from
$P_1$ such that the length $\sg{OP_{x'}} = -x' = |x'|$.%
\XVnote{*There are two ways of choosing this point, namely on either side of
$O$. Which side is chosen is not important.}

Each of the numbers $x$ and $x'$ is called the \emph{coordinate} of the
corresponding point of the line.

\FIGXVONE

\exercisegroup{15--3}
\begin{multicols}{2}
\begin{exlist}
\item How are the points $P_{x_1}$, $P_{x_2}$, and $O$ related if $x_1 < 0$
  and $x_2 > 0$? if $0 < x_2 < x_1$? if $x_1 < x_2 < 0$?
\item If we agree that points with larger coordinates ``lie to the right'' of
  points with lesser coordinates, how are we thinking of the drawn line $l$
  and the chosen point $P_1$?
\item Two points on the line have coordinates $-3$ and 7. What is the distance
  between them? between 5 and 12? between $x_1$ and $x_2$?
\end{exlist}
\end{multicols}

\begin{definition}
We will refer to the above construction of coordinates as a coordinatization
of the line $l$, with the point $O$ as \emph{origin of coordinates}.
\end{definition}

%========================================================= 15-3
\section*{15--3\quad THE EXISTENCE OF COORDINATES}
\addcontentsline{toc}{section}{15--3\quad The existence of coordinates}

The construction of plane coordinates is best carried out in a number of
simple steps. We assume that a unit of length has been selected in the plane.

\begin{definition}
Choose a point $O$ in the plane. Call this point the \emph{origin of
coordinates}. Choose any two perpendicular lines through $O$. Call these lines
\emph{coordinate axes}. Coordinatize each of the axes, with the point $O$ as
origin of coordinates. Select one of these lines and call it the \emph{axis of
abscissas}.* The other axis is called the \emph{axis of ordinates}.\dag\ See
Fig.~15--2.%
\XVnote{*This is not a fortunately chosen word. It would be better to say
``first coordinates.'' In most drawings in textbooks, the axis of abscissas is
drawn horizontally as the reader views it. There is nothing sacred about this,
we just agree to consider one coordinate before we do the other (something
must be first) and call this first coordinate by the strange name
\emph{abscissa}.}%
\XVnote{\dag\ This word is not much better, but the literature of mathematics
is filled with this word too, so make the best of it and learn it. Of course
``second coordinate'' is much better.}
\end{definition}

\FIGXVTWO

\exercisegroup{15--4}
\begin{multicols}{2}
\begin{exlist}
\item Draw a figure, different from Fig.~15--2, illustrating Definition
  15--2.
\item Compare your figure with your neighbor's. Did you both choose the
  ``same'' axis for abscissas? Turn your figure sideways. Do you have a
  different opinion about which is the first coordinate?
\item With the same drawing as in Exercise 1, reverse the coordinates on each
  of the axes. Is this a possible way of satisfying Definition 15--2? Suppose
  that we reversed the coordinates on only one of the axes.
\end{exlist}
\end{multicols}

Perhaps you will have observed a certain arbitrariness in the definition, in
that we \emph{choose} a point $O$ and \emph{choose} lines for axes. This is
actually the way things operate in practice. In applying analytic geometry to
problems in physics or engineering, it is usually necessary to \emph{choose}
coordinates. \emph{There are no fixed coordinate systems inherent in nature.}
Coordinate systems are convenient inventions that facilitate mathematical
computations.

Having set up our axes, we can make certain statements concerning them. We
give some of these in a theorem.

\begin{theorem}
The coordinate axes separate the plane into four regions, called quadrants,
such that if $P$ and $Q$ are points not on either axis and in the same
quadrant, then the segment $PQ$ does not meet either axis, whereas if $P$ and
$Q$ are in different quadrants, then $PQ$ meets at least one of the axes. See
Fig.~15--3.
\end{theorem}

\FIGXVTHREE

\noindent\emph{Proof.}
Each axis separates the plane. (Why?) Consider any point $P$ not on any axis.
The point $P$ is on one side of the axis of abscissas and one side of the axis
of ordinates. Let $Q$ be any point on the same side of the axis of abscissas
as $P$ and on the same side of the axis of ordinates as $P$. Then $PQ$ meets
neither axis. (Why?) The set of all such points $Q$ is a quadrant. If $Q_1$
and $Q_2$ are two points of that quadrant, then $Q_1Q_2$ does not intersect
either axis. (Why?)

There are four ways of choosing the point $P$, according to which sides of the
axes it is on. This gives rise to the four quadrants.

\exercisegroup{15--5}
\begin{multicols}{2}
\begin{exlist}
\item Describe the four ways of choosing $P$ to get the four quadrants.
\item Does the above proof require the axes to be perpendicular?
\end{exlist}
\end{multicols}

\begin{definition}
The half plane that contains the positive axis of abscissas will be called the
\emph{right half plane}, and similarly for \emph{left half plane}.* The half
plane that contains the positive axis of ordinates will be called the
\emph{upper half plane}, and similarly for \emph{lower half plane}.%
\XVnote{*This agrees with the way axes are conventionally drawn in textbooks.}
\end{definition}

\FIGXVFOUR

The quadrant that is in the right half plane and in the upper half plane will
be called the \emph{first quadrant;} the \emph{second quadrant} lies in the
upper and left half planes; the \emph{third quadrant}, in the lower and left
half planes; and the \emph{fourth quadrant}, in the lower and right half
planes. See Fig.~15--4.

\exercisegroup{15--6}
\begin{multicols}{2}
\begin{exlist}
\item In Fig.~15--5, which is quadrant I? quadrant II? quadrant III? quadrant
  IV?
\item[\stex 2.] Explain why Definition 15--3 does not require any previous
  knowledge of what the words ``left'' and ``right'' mean.
\item[\stex 3.] Has there been any use made of the fact that the axes have
  been selected to be perpendicular? What changes would have to be made so far
  if the axes were not at right angles?
\end{exlist}
\end{multicols}

\FIGXVFIVE

\begin{definition}
Let $P$ be any point of the plane, and $l_1$ and $l_2$ the lines through $P$
parallel, respectively, to the axis of ordinates and axis of abscissas. The
line $l_1$ meets the first axis in a point $A$, with coordinate $x$. The line
$l_2$ meets the second axis in a point $B$, with coordinate $y$. See
Fig.~15--6. The two numbers $x$ and $y$ are called the \emph{coordinates} of
$P$. We write them in parentheses, $(x, y)$, and refer to $(x, y)$ as a
\emph{point}. The number $x$ is the \emph{abscissa} of $P$, and the number $y$
is the \emph{ordinate} of $P$.
\end{definition}

Because the letters $x$ and $y$ are so frequently used for abscissa and
ordinate, the axes are often called the \emph{x-axis} and the \emph{y-axis}.
The plane then, when coordinatized in this way, is often called the
\emph{xy-plane}.\dag%
\XVnote{\dag\ Of course there is nothing really special about the letters $x$
and $y$. Use of different letters, say $u$ and $v$, would give us a
\emph{uv-plane}.}

Given the coordinates of a point, if one marks that point in a figure, it is
then said that the point has been \emph{plotted}.

\begin{remark}
The point $O$ has coordinates $(0, 0)$. (Why?)
\end{remark}

If $P$ does not lie on either axis, neither coordinate of $P$ can be 0. If the
first coordinate of $P$ is 0, then $P$ is on the axis of ordinates. The
converse is also true.

\FIGXVSIX

\exercisegroup{15--7}
\begin{multicols}{2}
\begin{exlist}
\item Draw axes, and plot the points $(2, 4)$; $(-2, 4)$;
  $(-3, -\tfrac{1}{2})$; $(0, -5)$; $(-4, 0)$; and $(7\tfrac{1}{2}, -3)$.
\item A point $P$ lies in the first quadrant. The perpendicular distance from
  the $x$-axis is $4\tfrac{1}{2}$ units; the distance from the $y$-axis is 3
  units. What are the coordinates of $P$? What are the coordinates if
  distances are the same but $P$ is in the second quadrant?
\item[\stex 3.] Suppose that, using the same lines as axes, the unit of
  distance is chosen twice as large. What would be the effect upon the
  coordinates of points?
\item[\stex 4.] Where are all the points in the plane with abscissa 0?
  abscissa negative? ordinate positive?
\item[\stex 5.] Where are all the points in the plane whose ordinates are
  equal to their abscissas?
\item[6.] Draw the line that passes through the points $(-1, 5)$ and
  $(2, -1)$. Can you easily draw a line parallel to it, which passes through
  the origin?
\item[\stex 7.] Carry out the construction of coordinates when the axes are
  not perpendicular.
\item[8.] Both coordinates of the point $S$ are negative numbers. The points
  $S$ and $T$ are on the same side of the $x$-axis, but on opposite sides of
  the $y$-axis. What can you say about the coordinates of $T$?
% Manual \item[...] labels do not advance the list counter; resync so the
% plain items below number 9-17 instead of restarting at 3.
\setcounter{exlisti}{8}
\item The positive and negative $x$- and $y$-axes are rays from $O$. Describe
  the four quadrants as interiors of angles formed by these rays.
\item The abscissa of a point is 0. Where must it lie?
\item The ordinate of a point is 0. Where can it be?
\item What can you say about a point if the product of its coordinates is
  negative? 0? positive?
\item Plot five distinct points with first coordinate equal to 3. Describe the
  set of all points in the plane such that the first coordinate of every point
  of the set is 3.
\item Consider the set of all points in the plane such that the sum of the
  coordinates of each point is 10. Plot six such points, two in each of
  quadrants I, II, and IV. Why are there no such points in quadrant III?
\item Describe the set of all points $(x, y)$ such that $x + 1 < 0$.
\item What is the perpendicular distance from $(3, -4)$ to the $x$-axis? to
  the $y$-axis? from $(-5, -7)$ to each axis?
\item How far is $(3, 4)$ from $(0, 0)$?
\end{exlist}
\end{multicols}

%========================================================= 15-4
\section*{15--4\quad STRAIGHT LINES}
\addcontentsline{toc}{section}{15--4\quad Straight lines}

Having coordinates, we naturally ask how lines may be described by
coordinates. In exercises we have already examined a few lines.

For example, the set of all points $(x, y)$ with abscissa 3 is the line shown
in Fig.~15--7, and every point on the line parallel to the $y$-axis, at a
distance 3 to the right of the $y$-axis, has its $x$-coordinate equal to 3.
Hence, this line can be specified briefly by writing just $x = 3$. We say then
that ``$x = 3$'' is the \emph{equation} of this line.

\FIGXVSEVEN

Conversely, when we see the equation $x = 3$, we ask ourselves, ``What is the
set of all points $(x, y)$ for which $x = 3$?''

In the same way, all lines parallel to the $y$-axis have an equation $x = a$,
where $a$ is some real number, positive or negative. (If $a = 0$, what is the
line?)

Likewise, lines parallel to the $x$-axis will have equations $y = b$.

\exercisegroup{15--8}
\begin{multicols}{2}
\begin{exlist}
\item A line is parallel to the $y$-axis and passes through the point
  $(-3, 5)$. What is its equation? What if it is parallel to the $x$-axis?
\item What is the equation of the line that bisects the first and third
  quadrants? the second and fourth?
\item[\stex 3.] Plot several points $(x, y)$ whose coordinates satisfy the
  equation $x + y = 5$. Can you describe briefly the set of all points
  $(x, y)$ whose coordinates satisfy this equation?
\end{exlist}
\end{multicols}

We will prove shortly that every straight line has an equation of a simple
kind. But first we see that such simple equations turn out to be straight
lines.

\begin{XVexample}[ 1]
Consider the set of all points $(x, y)$ whose coordinates satisfy the equation
$2x - 3y + 6 = 0$. We first find a few of these points, by choosing a value
for $x$ (or $y$) and then solving for $y$ (or $x$). Thus, if $x = 0$, we have
$0 - 3y + 6 = 0$, or $y = 2$. If $x = 3$, then $y = 4$. If $y = 0$, then
$x = -3$.
\par\medskip
\begin{center}
\begin{tabular}{c|c}
\multicolumn{1}{c}{$x$} & \multicolumn{1}{c}{$y$} \\
0 & 2 \\
3 & 4 \\
$-3$ & 0 \\
\end{tabular}
\end{center}
\end{XVexample}

\FIGXVEIGHTNINE

When these points are plotted, we see (by looking at Fig.~15--8) that they
apparently lie on a line.

Find other points whose coordinates satisfy the equation, and plot them.

\begin{XVexample}[ 2]
Consider the set of all points $(x, y)$ whose coordinates satisfy the equation
$-4y - 8 = 0$. From the equation, we see that $y = -2$, regardless of what $x$
is. The set of all points $(x, -2)$ is the line parallel to the $x$-axis and a
distance 2 units below. See Fig.~15--9.
\end{XVexample}

\exercisegroup{15--9}
Plot several points $(x, y)$ whose coordinates satisfy the following
equations:
\begin{multicols}{2}
\begin{exlist}
\item $2x + 3y - 6 = 0$.
\item $x - 2y + 4 = 0$.
\item $3x - 9 = 0$.
\item $-x - y + 4 = 0$.
\item $2ax + 3ay - 6a = 0$; $a \neq 0$.
\item $-3x + 5y - 15 = 0$.
\item $y = 2x + 2$.
\item $x = 3y - 6$.
\end{exlist}
\end{multicols}

From the examples we observe that equations of the form $Ax + By + C = 0$,
where $A$, $B$, and $C$ are given numbers, are satisfied by the coordinates of
points on lines. We now prove this theorem.

\begin{XVstheorem}
If $A$, $B$, and $C$ are given numbers, with not both $A$ and $B$ zero, then
all points $(x, y)$ whose coordinates satisfy the equation $Ax + By + C = 0$
lie on a single line.
\end{XVstheorem}

\noindent\emph{Proof.}
We separate the proof into three cases.

\emph{Case} 1. $A = 0$. Then $B \neq 0$ (why?) and the equation is
$By + C = 0$, or $y = -C/B$. Then all points with this $y$-coordinate lie on
the line parallel to the $x$-axis and passing through $(0, -C/B)$.

\emph{Case} 2. $B = 0$. In this case, the points lie on a line parallel to the
$y$-axis.

\emph{Case} 3. $A \neq 0$ and $B \neq 0$. It suffices to show that any two
points $(x_1, y_1)$ and $(x_2, y_2)$ whose coordinates satisfy the equation
lie on one line with the point $(0, -C/B)$. This will be the case if and only
if the right triangles indicated in Fig.~15--10 are similar, where

\begin{equation}\tag{15--1}
y_1 = \frac{-Ax_1-C}{B}; \qquad y_2 = \frac{-Ax_2-C}{B}.
\end{equation}

\FIGXVTENELEVEN

But these triangles are similar if and only if

\begin{equation}\tag{15--2}
\frac{y_2-(-C/B)}{x_2} = \frac{y_1-(-C/B)}{x_1}.
\end{equation}

From Eq.~(15--1) we see that
\[
\frac{y_2+C/B}{x_2} = \frac{-A}{B} = \frac{y_1+C/B}{x_1}.
\]

Hence we conclude that the triangles are similar and the points
$[0, -(C/B)]$, $(x_1, y_1)$, and $(x_2, y_2)$ are collinear.

\begin{XVstheorem}
If $l$ is a line, then there are numbers $A$, $B$, and $C$, with not both $A$
and $B$ zero, such that every point $(x, y)$ on $l$ has coordinates that
satisfy the equation
\begin{equation}\tag{15--3}
Ax + By + C = 0.
\end{equation}
\end{XVstheorem}

\noindent\emph{Proof.}
As in the previous theorem, we consider three cases.

\emph{Case} 1. $l$ is parallel to the $x$-axis. Then there is a point
$(0, b)$ on $l$, and every point $(x, y)$ on $l$ satisfies $y = b$. This is of
the form of Eq.~(15--3), with $A = 0$, $B = 1$, and $C = -b$.

\emph{Case} 2. $l$ is parallel to the $y$-axis. This case is left to the
student.

\emph{Case} 3. $l$ is parallel to neither axis. Choose $(x_1, y_1)$ and
$(x_2, y_2)$, distinct points on $l$. Then $x_1 \neq x_2$ and
$y_1 \neq y_2$. (Why?) Let $(x, y)$ be any other point on $l$.

As shown in Fig.~15--11, the points $(x_2, y_1)$ and $(x, y_1)$ form, with
$(x_1, y_1)$, $(x_2, y_2)$, and $(x, y)$, two right triangles that are
similar. Because of similarity,\dag

\begin{equation}\tag{15--4}
\frac{y-y_1}{y_2-y_1} = \frac{x-x_1}{x_2-x_1}
\end{equation}
\XVnote{\dag\ Figure 15--11 is drawn so that $x_1 < x_2 < x$, and
$y_1 < y_2 < y$, but all orders can occur for different points and different
lines. In every case, however, Eq.~(15--4) will be satisfied. Thus it makes no
difference which point is considered as $(x_1, y_1)$, so long as one is
consistent in its use.}

\noindent or
\[
(y_2 - y_1)x - (x_2 - x_1)y + x_2y_1 - y_2x_1 = 0,
\]
which is of the form of Eq.~(15--3). This completes the proof.

\begin{XVremark}{1}
An equation equivalent to Eq.~(15--4) is
\begin{equation}\tag{15--5}
y - y_1 = \frac{y_2-y_1}{x_2-x_1}\,(x-x_1).
\end{equation}
This is a useful equation to remember because it enables one to write at once
an equation satisfied by the coordinates of any point on a line through two
given points.
\end{XVremark}

\begin{XVremark}{2}
Two different equations may correspond to the same line. Thus for example,
$x = 2$ and $2x - 4 = 0$ represent the same line. Such equations are said to
be \emph{equivalent} equations. In arriving at equations for lines, we usually
do not care which equivalent equation we obtain.
\end{XVremark}

\begin{definition}
If $Ax + By + C = 0$ is an equation satisfied by the coordinates $(x, y)$ of
all points on a line $l$, we say that $Ax + By + C = 0$ is the equation of
$l$. We also speak of the \emph{``line $Ax + By + C = 0$.''} Such an equation
of the first degree in $x$ and $y$ will be called a \emph{linear equation} in
$x$ and $y$.
\end{definition}

\FIGXVTWELVE

\begin{XVexample}
Find the equation of the line through $(-2, 3)$ and $(1, 5)$ in Fig.~15--12,
and draw the line. Using Eq.~(15--5), we have
\[
y - 3 = \frac{5-3}{1-(-2)}\,[x-(-2)]
\qquad\text{or}\qquad
y = \tfrac{2}{3}x + \tfrac{13}{3}.
\]
\end{XVexample}

\exercisegroup{15--10}
\begin{multicols}{2}
\begin{exlist}
\item Draw the lines
  \begin{enumerate}[label=(\alph*),leftmargin=1.5em,itemsep=0pt,topsep=2pt,
    parsep=0pt]
  \item $3x - 4y = 12$,
  \item $3x - 4y = 8$,
  \item $4x + 3y = 0$,
  \item $4x + 15 = 0$,
  \item $-2x - 3y + 5 = 0$,
  \item $y = \dfrac{-1}{2}\,x + 3$,
  \item $y = \dfrac{-1}{2}\,x$,
  \item $x = 3y + 3$,
  \item[(k)] $5y - 12 = 0$,
  \item[(l)] $2x + y = 7$.
  \end{enumerate}
\item Find the equations of the lines through the following pairs of points:
  \begin{enumerate}[label=(\alph*),leftmargin=1.5em,itemsep=0pt,topsep=2pt,
    parsep=0pt]
  \item $(-1, 1)$ and $(3, -1)$,
  \item $(2, 3)$ and $(-2, -3)$,
  \item $(a, 0)$ and $(0, b)$,
  \item $(-3, 0)$ and $(0, -5)$,
  \item $(-2, -4)$ and $(-5, -5)$,
  \item $(0, b)$ and $(1, m + b)$.
  \end{enumerate}
\item In Exercise 1, how many points did you plot to draw each line? Is there
  a reason for plotting more than two?
\item Find the coordinates of the point of intersection of the lines
  $3x - 5y = 7$ and $2x + 3y = -1$.
\item A line meets the $y$-axis in $(0, -3)$ and the $x$-axis in $(-4, 0)$.
  Where does this line intersect the line $y = 7x - 19$?
\item Show that the lines $2x + 3y = 0$ and $4x + 6y = 5$ are parallel.
\end{exlist}
\end{multicols}

%========================================================= 15-5
\section*{15--5\quad DISTANCE}
\addcontentsline{toc}{section}{15--5\quad Distance}

The distance between any two points of the plane is easy to calculate from the
coordinates of the points. It is particularly easy if the line through the
points is parallel to one of the axes. Thus, $\sg{AB} = 1 - (-3) = 4$, and
$\sg{CD} = |x_2 - x_1| = |x_1 - x_2|$ in Fig.~15--13.

\FIGXVTHIRTEENFOURTEEN

If the line through the points is not parallel to one of the axes, then with
the point $Q$ $(x_2, y_1)$ a right triangle $P_1P_2Q$ is formed, as in
Fig.~15--14, and the Pythagorean theorem asserts that
$\sg{P_1P_2}^2 = \sg{P_1Q}^2 + \sg{P_2Q}^2 = |x_2 - x_1|^2 +
|y_2 - y_1|^2$. Taking square roots, we obtain
$\sg{P_1P_2} = \sqrt{(x_2-x_1)^2 + (y_2-y_1)^2}$.

\begin{remark}
In using this formula it makes no difference which point is considered to be
$(x_1, y_1)$ because $(x_1 - x_2)^2 = (x_2 - x_1)^2$ and
$(y_1 - y_2)^2 = (y_2 - y_1)^2$.
\end{remark}

\exercisegroup{15--11}
\begin{multicols}{2}
\begin{exlist}
\item Show that the triangle with vertices $(2, 6)$, $(5, 7)$, and
  $(-8, -2)$ is a right triangle.
\item What kind of triangle is the one with vertices $(3, 2)$, $(-4, 4)$, and
  $(1, -5)$?
\item[\stex 3.] Write an equation that states that the point $(x, y)$ is
  equidistant from $(-3, 2)$ and $(2, 1)$. Show that this equation is
  equivalent to a linear equation. Hence conclude that the set of all points
  equidistant from $(-3, 2)$ and $(2, 1)$ is a straight line.
\item[4.] Show that $(-1, -18)$, $(-13, -2)$, $(11, -8)$, and $(-1, 8)$ are
  vertices of a parallelogram.
% Manual \item[...] labels do not advance the list counter; resync so the
% plain items below number 5-7 instead of restarting at 3.
\setcounter{exlisti}{4}
\item Show that if $(x, y)$ is a point on a circle whose center is $(-2, 1)$
  and radius is 3, then $x^2 + y^2 + 4x - 2y = 4$.
\item Show that the points $(2, 1)$, $(3, \tfrac{1}{2})$, and $(-2, 3)$ lie on
  a line. What geometric theorems do you use?
\item Show that the distance formula yields $|x_1 - x_2|$ if $y_1 = y_2$.
\end{exlist}
\end{multicols}

%========================================================= 15-6
\section*{15--6\quad MIDPOINT}
\addcontentsline{toc}{section}{15--6\quad Midpoint}

The midpoint of the segment $P_1P_2$ may be shown to have coordinates
\[
\left(\frac{x_1+x_2}{2},\ \frac{y_1+y_2}{2}\right).
\]

The proof is left to the student. If the line through the points is parallel
to one of the axes, the proof is very simple. If $x_1 \neq x_2$ and
$y_1 \neq y_2$, a right triangle may be formed as in Fig.~15--15, and
parallels to the coordinate axes through the midpoint will bisect the sides in
points whose coordinates are easily computed.

\begin{XVexample}
Show in Fig.~15--16 that the midpoint formula actually gives a point
equidistant from $(-1, 3)$ and $(1, -1)$.
\end{XVexample}

\FIGXVFIFTEENSIXTEEN

The point $M$ is
$\left(\dfrac{-1+1}{2},\ \dfrac{3+(-1)}{2}\right) = (0, 1)$. Then
\[
\sg{P_1P_2} = \sqrt{2^2 + 4^2} = \sqrt{20} = 2\sqrt{5},
\]
\[
\sg{P_1M} = \sqrt{1^2 + 2^2} = \sqrt{5},
\]
and
\[
\sg{MP_2} = \sqrt{1^2 + 2^2} = \sqrt{5}.
\]
Hence
\[
\sg{P_1P_2} = 2\sqrt{5} = \sqrt{5} + \sqrt{5} = \sg{P_1M} + \sg{MP_2}.
\]

\exercisegroup{15--12}
\begin{multicols}{2}
\begin{exlist}
\item Find the midpoints of the sides, and the lengths of the medians, of the
  triangle whose vertices are $(-4, 10)$, $(12, 10)$, and $(4, -2)$.
\item Find the coordinates of the points that divide the segment $(-3, -7)$
  $(11, 17)$ into fourths.
\item[\stex 3.] Find the points that trisect the segment $(-3, -7)$
  $(6, 11)$.
\item[\stex 4.] Derive a formula for the coordinates of the points that
  trisect the segment $(x_1, y_1)$ $(x_2, y_2)$.
\item[\stex 5.] Show that the medians of the triangle with vertices
  $(x_1, y_1)$, $(x_2, y_2)$, and $(x_3, y_3)$ intersect in the point
  $[(x_1 + x_2 + x_3)/3, (y_1 + y_2 + y_3)/3]$.
\item[6.] The point $(-3, 4)$ is the midpoint of a segment, one end of which
  is $(5, 7)$. Find the other end.
\item[\stex 7.] In Fig.~15--17, show that the line joining the midpoints of
  the sides of a triangle is half as long as the third side, and parallel to
  it. [\emph{Hint:} Choose your coordinate system so that one vertex is at the
  origin and another is on the positive $x$-axis.]
\item[\stex 8.] Prove that the diagonals of a parallelogram bisect each
  other.
\end{exlist}
\end{multicols}

\FIGXVSEVENTEEN

%========================================================= 15-7
\section*{15--7\quad CURVES AND EQUATIONS}
\addcontentsline{toc}{section}{15--7\quad Curves and equations}

We have seen that certain sets of points can be described by equations. Thus a
straight line has a linear or first degree equation, in the sense that the
coordinates of every point on the line satisfy an equation of first degree in
$x$ and $y$ of the form $Ax + By + C = 0$, and conversely. In the same way,
other sets of points can be described by different equations. Perhaps the
simplest example of this is furnished by circles.

\FIGXVEIGHTEENNINETEEN

Suppose that $(x, y)$ in Fig.~15--18 is a point on the circle whose center is
$(h, k)$ and whose radius is $r$, then
\[
\sqrt{(x-h)^2 + (y-k)^2} = r
\]
or, equivalently,

\begin{equation}\tag{15--6}
(x-h)^2 + (y-k)^2 = r^2.
\end{equation}

Conversely, a point $(x, y)$ whose coordinates satisfy Eq.~(15--6) will be at
a distance $r$ from $(h, k)$ and therefore lies on the circle.

We speak of Eq.~(15--6) as ``the'' equation of a circle, although any
equivalent equation would also be called ``the'' equation. For example, the
equation of the circle whose center is at $(-2, 3)$, and whose radius is 4, is
\[
(x + 2)^2 + (y - 3)^2 = 16
\]
or
\[
x^2 + y^2 + 4x - 6y = 3.
\]

Given an equation in the form of Eq.~(15--6), we can recognize it at once for
the equation of a circle. Any equation equivalent to Eq.~(15--6) is the
equation of a circle. For example, consider the equation
\[
x^2 + y^2 + 6x - 8y + 19 = 0.
\]

Rewrite the equation as
\[
(x^2 + 6x + \,?) + (y^2 - 8y + \,?) = -19 + \,? + \,?.
\]

By completing the squares of the expressions in parentheses, we have
\[
(x^2 + 6x + 9) + (y^2 - 8y + 16) = -19 + 9 + 16
\]
or
\[
(x + 3)^2 + (y - 4)^2 = 6,
\]
which is the equation of a circle whose center is at $(-3, 4)$ and whose
radius is $\sqrt{6}$. See Fig.~15--19.

Other curves (sets of points) will have other kinds of equations. The student
is referred to books on analytic geometry for more detail.

\exercisegroup{15--13}
\begin{multicols}{2}
\begin{exlist}
\item Write the equation of the circle whose center is $(3, 5)$ and whose
  radius is 6.
\item Write the equation of the circle whose center is at $(-4, -3)$ and which
  (a) is tangent to the $y$-axis, (b) passes through the origin.
\item Find center and radius of the following circles:
  \begin{enumerate}[label=(\alph*),leftmargin=1.5em,itemsep=0pt,topsep=2pt,
    parsep=0pt]
  \item $x^2 + y^2 + 8x - 9 = 0$,
  \item $x^2 + y^2 - 6x - 10y = 2$,
  \item $x^2 + y^2 - 7 = 0$,
  \item $2x^2 + 2y^2 + 5x + 4y - 5 = 0$,
  \item $x^2 + y^2 - 6x - 8y + 25 = 0$.
  \end{enumerate}
\item[\stex 4.] Find an equation satisfied by the coordinates of every point
  $(x, y)$ whose distance from $(2, 0)$ is equal to its distance from the line
  whose equation is $x + 2 = 0$.
\item[\stex 5.] If the sum of the distances from $(x, y)$ to $(-3, 0)$ and
  from $(x, y)$ to $(3, 0)$ is equal to 10, show that
  $16x^2 + 25y^2 = 400$.
\end{exlist}
\end{multicols}
\clearpage
% ==================== ch16 ====================
\setcounter{theorem}{0}\setcounter{definition}{0}\setcounter{construction}{0}\setcounter{exgroup}{0}\setcounter{footnote}{0}
\setcounter{chapter}{16}

\chapter*{\raggedright\Huge PHILOSOPHY OF\\MATHEMATICS}
\addcontentsline{toc}{chapter}{16.\ Philosophy of Mathematics}
\markboth{PHILOSOPHY OF MATHEMATICS}{}
\vspace{-1em}\noindent\rule{\linewidth}{0.6pt}\vspace{1em}

%========================================================= 16-1
\section*{16--1\quad WHAT IS PHILOSOPHY?}
\addcontentsline{toc}{section}{16--1\quad What is philosophy?}

The word ``philosophy'' is a combination of two Greek words and means
literally \emph{love of wisdom}. In the early Greek civilization, philosophy
was the study of all branches of knowledge. As man's learning increased
through the ages, certain disciplines grew until they split away from
philosophy and became separate areas for study. We no longer think of
medicine, law, mathematics, physics, chemistry, biology, economics, etc., as
parts of philosophy, although philosophy was the father of all these sciences.
That is, the first men who called attention to important concepts in
mathematics, physics, chemistry, economics, political science, etc., were
philosophers. It was philosophy that suggested how knowledge about these new
areas should be collected and organized.

Philosophy has always concerned itself with the unknown---with those
situations in which man is not sure of himself. Because of this, many of the
conclusions that philosophers reach cannot be checked by experiment, as we
check physical laws, and consequently many people have the mistaken belief
that time devoted to philosophical thought is largely wasted.

It is impossible to describe briefly and clearly just what the subject matter
of philosophy consists of today. Some of the branches into which philosophy is
subdivided will be listed here, but we make no pretense of completeness.

\emph{Moral philosophy}, or \emph{ethics}, investigates the nature of right
and wrong. It is concerned with distinguishing between good and bad conduct.
It raises questions like, ``Is an action to be judged good or bad according to
its effect or its intent?'' ``What are \emph{good} and \emph{evil}?'' ``How
shall we determine whether one action is better or worse than another?''

\emph{Religious philosophy} is the study of concepts common to religions. On
the other hand, \emph{theology} is concerned with the logical structure of a
single religion.

% NOTE (source oddity): the book jumps straight from "are constructed." into
% the two quoted questions, with no "It raises questions like," lead-in of the
% kind the moral-philosophy paragraph uses above. Verified at 300 dpi; it is
% printed that way. Do not "repair" this.
\emph{Political philosophy} does not consider details of government, but
rather the principles upon which various systems of government are
constructed. ``How much authority may the state exert over an individual?''
``Are there fundamental human rights, which no state can usurp? If so, what
are they?''

\emph{Metaphysics} (``beyond'' physics) is sometimes described as the study of
final or ultimate meanings. Much that is nonsense has been written in this
area, but it includes many deep and fascinating insights.

\emph{Epistemology} can be described briefly as the theory of knowledge. It
considers the general questions, ``How do we learn?'' and, ``How shall we
organize our knowledge?'' As a matter of fact, the philosophy of
\emph{mathematics} is just a branch of epistemology; so you may think of this
chapter as an introduction to a fragment of epistemology, which is itself a
fragment of philosophy.

It is impossible to discuss the philosophy of mathematics without considering
the philosophy of science, for mathematics is used as a tool in every science
to organize and interpret scientific knowledge. We shall mention how
mathematics thus often leads to the discovery of new scientific facts.

%========================================================= 16-2
\section*{16--2\quad THE TASK OF SCIENCE}
\addcontentsline{toc}{section}{16--2\quad The task of science}

All information about what we call the \emph{real} or \emph{external} world is
obtained from the five senses: sight, sound, taste, smell, and touch. Our
sense organs, for example the eye, send messages to the brain. A newborn baby
cannot do much with these messages, but as he grows older, he learns to
interpret them. We are always amazed at the accomplishments of a person who is
both blind and deaf (for example, Helen Keller), for it is difficult for us to
grasp how such a person can have much understanding of the world. What we are
inclined to forget is that all of us have learned how to interpret information
from our sense organs so as to make useful responses. As you read this page,
your eye receives impressions and sends messages to the brain. You have
learned to interpret and respond to these messages.

% NOTE (source oddity): "that is a touch" is printed without a comma after
% "is" (modern usage would set "that is, a touch"). Verified at 300 dpi.
% Reproduce as printed.
\emph{Every scientific observation is, at first, a sense impression, that is a
touch, smell, taste, sight, or sound, for some experimenter. It is the task of
scientists to understand, explain, and predict happenings in the world of our
sense impressions.}

Some philosophers have suggested that because our only knowledge of the
external world comes from sense perception, \emph{we cannot even prove that
the external world has any existence apart from our sense perception of it;
that perhaps the world that we believe we see, touch, and smell exists only as
things exist in a dream.}

% NOTE (source error, reproduce as printed): Berkeley is called a
% "17th-century" philosopher. He is 18th-century (1685--1753). The error is
% the book's --- see ch16-UNCERTAIN.md item 4. Do not silently correct.
The story is told that the great 17th-century philosopher Bishop Berkeley was
so convinced by this line of thought that one day, when he was walking and
meditating upon philosophical ideas, he encountered a large stone in his path
and so much doubted the actual existence of the stone that he drew back his
foot and gave it a vigorous kick. We have no record of the good bishop's
thoughts as he hobbled home.

%========================================================= 16-3
\section*{16--3\quad KINDS OF SCIENCE}
\addcontentsline{toc}{section}{16--3\quad Kinds of science}

The world is very complicated, and we have little understanding of it. To
simplify the task of understanding our environment, certain topics that seem
somewhat related are studied by themselves as separate sciences. For example,
biology is the study of living organisms, and biology is itself separated into
several subsciences, like zoology, the study of animal life, and botany, the
study of plant life. To understand all aspects of biology, it is necessary to
know some chemistry and physics. For this reason chemistry and physics are
regarded as more basic, or more fundamental, than biology. Of the sciences
concerned with the real world, the most fundamental is physics, which deals
with the properties of matter.

All sciences become \emph{quantitative} when they have progressed far enough.
This means that measurements are made and numbers are introduced into the
sciences. Those things which we measure, and to which we assign numbers, are
called \emph{quantities}. As soon as numbers are introduced into a science,
mathematics becomes necessary for its study. As a matter of fact, one needs
geometry in order to discuss the physics of motion of bodies, even before one
makes measurements. Even psychology and sociology, two of the newer sciences
that have only recently split away from philosophy, need statistics in order
to organize and interpret the huge number of facts with which they deal. Thus
all the sciences depend upon mathematics, so that it has been termed ``the
handmaiden of the sciences.''\XVIstar
% NOTE (source oddity): the footnote names Baltimore twice. Verified at
% 300 dpi --- the duplication is genuinely printed. Do not correct.
\XVIstarnote{E.~T. \textsc{Bell}, \emph{The Handmaiden of the Sciences}.
Baltimore: The Williams and Wilkins Co., Baltimore, 1937.}

Sometimes mathematics itself is classified as one of the sciences; and
sometimes one hears the phrase ``science \emph{and} mathematics.'' In any
event, call it what you may, mathematics is different from the other sciences
in that it is \emph{independent of experiments}. \emph{Mathematics is a
creation of the mind alone.}

%========================================================= 16-4
\section*{16--4\quad IS PHYSICS TRUE?}
\addcontentsline{toc}{section}{16--4\quad Is physics true?}

We cannot give an acceptable definition of physics in a few words, but it will
not take long to make the question of this section meaningful. You have
studied a little physics previously, so you should know that physics is
concerned with forces, motion, heat, sound, electricity, and much else. For
example, in that part of physics called mechanics, it is learned that if at
some instant we know the position and velocity of a ball, and if in addition
we know the forces acting upon it, then we can \emph{calculate} where it will
be at any later time.

You need not concern yourself here with the exact method used to make this
calculation. What is important is that \emph{given certain information, we can
predict what will happen}. Physics, and indeed every science, is just as
``true'' as its predictions are \emph{reliable}. As soon as a physical theory
fails to predict what will happen, then it becomes necessary to construct a
new theory that will predict more accurately. In physics, the final test of
``truth'' is whether the predictions we get from our theory actually agree
with our observations. As an example of this, Aristotle had written in his
\emph{Physics} that heavier bodies would fall more rapidly than lighter ones.
This prediction is a logical consequence of Aristotle's physical
\emph{theories}, or \emph{laws}. Every school child knows how Galileo
\emph{proved} these theories \emph{false} by dropping balls of various sizes
and weights from the leaning tower of Pisa. Every observer who was open-minded
enough to watch the experiment (and there were indeed some who were not)
became convinced that Aristotle's laws of motion were false and that it was
necessary to replace them by new laws that would come closer to predicting the
events that were actually observed.

The physical laws used by the physicist to make his predictions are actually
\emph{assumptions} and are expressed in terms of mathematical formulas. If the
mathematical calculations employed do not lead to correct predictions, this
does not mean that the mathematics is incorrect.\XVIstar\ Instead it means
that the physical laws do not agree well enough with reality, that the
\emph{physical assumptions} are at fault.
\XVIstarnote{Of course someone could make a computational mistake, but this is
the fault of the computer, not of the mathematics he is using.}

When the physicist defines physical concepts like force, weight, etc., and
then gives physical laws that tell how these concepts are related, he is
creating what we call a \emph{mathematical model}. What he is saying is that
\emph{if} there are things in the real world that correspond to these
concepts, and \emph{if} these things are related to one another according to
the rules called physical laws, \emph{then} mathematics tells us what will
happen in the real world.

If the predictions come out wrong, then the physicist must construct a better
mathematical model. Often it happens that the physicist cannot solve the
mathematical problems that arise within the mathematical model. Then he calls
upon the mathematician for help. In this way physics supplies mathematics with
interesting problems.

As an example of a simple mathematical model, consider our solar system. The
physicist thinks of the sun and planets as points that are tracing out curves
as they move through space. By using calculus and analytic geometry, these
curves are described by mathematical equations that predict how the planets
ought to move through the sky if the law of gravitation is valid.

% NOTE (source error, reproduce as printed): the book dates the
% Uranus/Neptune episode to "early in the 18th century". Adams and LeVerrier
% worked in the 1840s, and Uranus was not discovered until 1781. The error is
% the book's --- see ch16-UNCERTAIN.md item 4. Do not silently correct.
It is most interesting that early in the 18th century, when astronomers
calculated how the planets \emph{should} move according to the law of
gravitation, it was noticed that one of the planets, Uranus, did not follow
the exact path predicted for it by the mathematical model. The astronomers had
so much faith in the mathematical model that they became sure that there must
be an ``invisible'' planet, which was pulling Uranus off its course with its
gravitational attraction. The Englishman John C. Adams and the Frenchman
LeVerrier set about attempting to solve the problem of determining the path of
the ``invisible'' planet.\XVIstar\ Working with the mathematical model, they
computed where it \emph{ought} to be to affect the course of Uranus as it did.
When observers focused their telescopes upon that part of the heavens, they
saw the planet Neptune precisely where the mathematical model predicted that
it would be found. This is but one of the many dramatic illustrations of how
an \emph{abstract} mathematical model can lead us to the discovery of new
facts about our universe.
\XVIstarnote{A fascinating account of this work can be found in
J.~R.~Newman's \emph{The World of Mathematics}. Vol.~2, pp.~822--839.}

%========================================================= 16-5
\section*{16--5\quad IS MATHEMATICS TRUE?}
\addcontentsline{toc}{section}{16--5\quad Is mathematics true?}

If by this question we mean, ``Does the real world behave like the
mathematical models that we construct?'' then we are asking a nonmathematical
question. The \emph{truth} of mathematics is independent of the real world.
The \emph{usefulness} of one kind of mathematics or another does depend upon
the real world.

\emph{Pure mathematics is concerned with the logical consequences of postulate
systems.} In pure mathematics, postulate systems are studied primarily because
they are interesting, and often without regard to any future use. In the
history of mathematics, there are many postulate systems that have been
studied for their own sakes and only much later found to be applicable to the
external world. For example, the non-Euclidean geometry needed in the theory
of relativity was first studied many years before any physical application of
it was made.

\emph{Applied mathematics is concerned with those parts of pure mathematics
that are of greatest use in the science of today.} Scientific applications
often suggest the constructions of new postulate systems. Mathematics and
science thus serve as stimuli to each other.

We have not yet given a direct answer to the question, ``Is mathematics
true?'' What is ``truth,'' anyway? Each of us has some sort of feeling about
what ``truth'' means. Truth, for the physicist, is \emph{that which works},
and this attitude is characteristic of science. Truth, in practical affairs,
usually refers to \emph{whether some event actually took place}. Here, we
often take the word of others whom we trust. Truth, in spiritual matters,
depends upon \emph{faith}. In general, agreement among people regarding
spiritual beliefs is difficult to obtain, for there is no test by experience
that we can all agree upon.

When we say that a system of mathematical postulates is ``true,'' we mean
only, \emph{these postulates are free from contradiction}. That is, if it is
impossible to prove from our postulates that a statement is true, and then
also prove that the same statement is false, in this case and only in this
case do we call our postulates true. Remember that we do not assert in
mathematics that anything actually exists in the physical world. Rather, our
theorems have the form: \emph{If} a set of things has so-and-so properties,
\emph{then} it has such-and-such properties.

We have developed our geometry this year from statements that dealt with a set
of points that we called our plane. A mathematician is naturally pleased to
encounter in the physical world a set of objects that actually has the
properties he has postulated for one of his mathematical systems, for then he
can be sure that the set he finds will have other properties that he has
established by studying his abstract system. A natural question to ask at the
close of this year's study of geometry is whether there actually exists in the
real world a set of things that satisfies the postulates of our geometry. Does
the space we live in actually have the properties of the Euclidean space we
study in geometry?

The answer will probably not satisfy you. The answer is, \emph{we do not know.
Furthermore, no mathematician ever expects to know.} Our geometric plane
consists of \emph{infinitely} many points. \emph{It will never be possible to
determine whether our universe is finite or infinite in extent.} Although we
have no assurance that any part of the real world actually fits our geometric
postulates, experience has shown us that much of our geometric theory is
applicable to the world around us; and we can indeed use our geometry to make
very accurate predictions.

%========================================================= 16-6
\section*{16--6\quad FREEDOM FROM CONTRADICTION}
\addcontentsline{toc}{section}{16--6\quad Freedom from contradiction}

In this section we discuss a very deep question. When we have a postulate
system before us, how do we know that these postulates are
\emph{consistent}? How do we know that we shall never deduce from these
postulates, by correct logical reasoning, two contradictory conclusions? Let
us give an example, first given by Bertrand Russell,\XVIstar\ to illustrate
what we mean.
\XVIstarnote{An English philosopher who is interested in the foundations of
mathematics. At about the turn of the century, he wrote \emph{Principia
Mathematica} with A.~N. Whitehead. In the book they attempt to derive all
mathematics from logic alone.}

\emph{Russell's barber.} Let us assume that, in a certain small town, there is
but one barber. We assume that this barber shaves everyone in the town who
does not shave himself, and no one else.

We now ask who shaves the barber.

\emph{If} the barber does not shave himself, \emph{then} he is a person who
does not shave himself, and therefore he is shaved by the barber---himself.
Since this is impossible, we have proved that the statement \emph{``The barber
does not shave himself''} is false. Hence, the barber shaves himself. But, by
our assumptions, \emph{if} the barber shaves himself, \emph{then} he is one of
the men who is not shaved by the barber---again himself. This proves that the
statement \emph{``The barber shaves himself''} is false, and hence,
\emph{``The barber does not shave himself,''} is true. But now we have arrived
at a contradiction, having shown from our postulates that the statement
\emph{``The barber does not shave himself'' is both true and false.}

This example shows that in forming a set of postulates, we must be careful
that our postulates are consistent; that is, that they do not lead to
contradictory conclusions. When our postulate sets are very simple, it is
sometimes possible to prove that no contradiction can ever be reached by
logical reasoning. For more complicated postulate sets, this may not be
possible.

David Hilbert spent much of the latter portion of his life trying to prove the
consistency of certain interesting mathematical systems. At about the same
time that he wrote his \emph{Foundations of Geometry}, he showed that geometry
is consistent if arithmetic is consistent. He accomplished this proof by using
the ideas of Chapter 15. The details need not concern us here. All we need
know is that our geometry is consistent if arithmetic is. We should remark
that non-Euclidean geometries have also been proved to be consistent,
providing that arithmetic is. Of course this has nothing to do with the
question of whether the space we live in is Euclidean or non-Euclidean.

% NOTE: "0 = 1" is set inline; the exclamation mark after it is rhetorical,
% NOT a factorial. Do not typeset it as $0 = 1!$ in a way that reads as one.
Now, how do we know that arithmetic is consistent? The answer may again seem
disappointing, for the answer is, \emph{we do not know!} Perhaps we never will
know! There is no proof that we can never be led to an impossible conclusion
in arithmetic. For example, we do \emph{not} know that it is impossible to
\emph{prove}, by a valid chain of logical arguments, that $0 = 1$!

All mathematicians believe that arithmetic is consistent. This belief is based
upon faith, not proof. Of course, arithmetic is useful. It works, and helps us
reach conclusions that turn out to be correct. But here is the position in
which the scientist finds himself. If an inconsistency is discovered in
arithmetic, the postulates for arithmetic will have to be changed. The
possibility seems to be remote, so it does not worry working mathematicians
too much.

%========================================================= 16-7
\section*{16--7\quad MATHEMATICAL RESEARCH}
\addcontentsline{toc}{section}{16--7\quad Mathematical research}

Euclid's geometry is more than 2000 years old. The average person holds the
erroneous opinion that all of mathematics was discovered long ago, and that
today we study only old stuff. There are, however, many journals published
throughout the world that are devoted primarily to \emph{new} mathematics. In
fact, during the past fifty years, more new mathematics has been
invented\XVIstar\ than in all previous history! Mathematical progress is
stimulated by both intellectual curiosity and the needs of modern science.
Because of the mathematical nature of modern science and engineering,
mathematicians are in greater demand than ever before and command excellent
salaries.
\XVIstarnote{Is mathematics \emph{discovered} or \emph{invented}? The first
word suggests that mathematics is ``sitting around waiting to be found.'' The
second suggests that mathematics is created by the mind, that human activity
is essential.}

It is impossible to describe here the breadth and depth of current
mathematical research. As yet, we have not undertaken enough elementary
mathematics to appreciate such a description. However, in your study of plane
geometry, you have made a fine beginning by studying a single mathematical
system in some detail. The logical features of geometry are common to all
mathematics, so that you have before you one example of a postulate system and
some of its consequences.

%========================================================= 16-8
\section*{16--8\quad MATHEMATICS AND LIFE}
\addcontentsline{toc}{section}{16--8\quad Mathematics and life}

How may an understanding of the philosophy of mathematics affect you as an
individual in your everyday life? Mathematicians are much like any other
group. If you were to attend the International Congress of Mathematicians,
which convenes once every four years, you would see a group of people
differing widely in dress, speech, religious and political beliefs, and
physical appearance. But you would, if you were observant enough, see many
common characteristics. Everyone there would have a way to communicate with
everyone else---not through English, French, German, or Russian, but through
the universal symbolism of mathematics. You would soon realize that these
folk, while actively interested in the real world, were still far more
interested in the world of ideas.

You would see that each one present had a tremendous respect for the
personalities and intellectual capabilities of his neighbors.

The philosophy of mathematics speaks to those who will listen and assures them
that even more important than wresting food, shelter, and material riches from
the earth is the acquisition of knowledge and understanding.

%========================================================= Suggested Reading
\par\medskip
\begin{center}\bfseries SUGGESTED READING\end{center}
\par\nopagebreak\medskip

\noindent From \emph{The World of Mathematics}, by J.~R.~Newman.
\begingroup\raggedright
\begin{exlist}
\item \textsc{Philip} E.~B. \textsc{Jourdain}, ``The Nature of Mathematics.''
      Vol.~1, pp.~4--71.
\item J.~W.~N. \textsc{Sullivan}, ``Mathematics as an Art.''
      Vol.~3, pp.~2015--2021.
\item \textsc{Raymond} L. \textsc{Wilder}, ``The Axiomatic Method.''
      Vol.~3, pp.~1647--1667.
\end{exlist}
\endgroup
\clearpage
% ==================== appendix ====================
\setcounter{theorem}{0}\setcounter{definition}{0}\setcounter{construction}{0}\setcounter{exgroup}{0}\setcounter{footnote}{0}
\setcounter{chapter}{17}
% The Appendix text is verbatim, so line breaking must bend, not the wording:
% give TeX a little extra stretch rather than reword to clear overfull boxes.
\emergencystretch=2em

% The book sets this opener flush RIGHT: "APPENDIX" occupies the slot that
% "CHAPTER 10" occupies on a numbered chapter opener (cf. book p.180).
\chapter*{\raggedleft\Huge APPENDIX}
\addcontentsline{toc}{chapter}{Appendix}
\markboth{APPENDIX}{}
\vspace{-1em}\noindent\rule{\linewidth}{0.6pt}\vspace{1em}
\figurekey

This appendix lists all the postulates for plane geometry, along with
some of the definitions and theorems. In addition to being convenient
for reference purposes, the appendix merits study, because reading it
over and thinking about the order of the theorems will help you fix the
structure of Euclidean geometry in your mind.

%========================================================= the postulates
\APPhead{THE POSTULATES}{The Postulates}

There are three groups of postulates, and three postulates listed singly.
They are:

\begin{center}
\begin{minipage}{0.86\linewidth}
Postulates of incidence (3 postulates)\\
Postulates of order or betweenness (5 postulates)\\
Postulates of congruence (7 postulates)\\
The postulate of Archimedes\\
The completeness postulate\\
The parallel postulate.
\end{minipage}
\end{center}

\APPsub{The Incidence Postulates}

\begin{applist}
\item There are at least three points not all on a line.
\item For any two different points, there is exactly one line containing
      these points.
\item Every line contains at least two points.
\end{applist}

\APPsub{The Betweenness Postulates}

\begin{applist}[resume]
\item If $B$ is between $A$ and $C$, then $A$, $B$, and $C$ are three
      different, or distinct, points on a line.
\item For every three points on a line, exactly one of them is between the
      other two.
\item Any four points on a line may be named $A_1$, $A_2$, $A_3$, $A_4$, so
      that the only betweenness relations are the same as the order of the
      subscript numbers. That is, the betweenness relations are: $A_2$
      between $A_1$ and $A_3$; $A_2$ between $A_1$ and $A_4$; etc.
\item If $A$ and $B$ are two points, then there is at least one point $C$
      such that $B$ is between $A$ and $C$, and at least one point $D$ such
      that $D$ is between $A$ and $B$.
\item Every line separates the plane. By this we mean that all points of
      the plane not on the line are divided into two sets, called the two
      sides of the line, having the following properties:
      \begin{applist}
      \item If $P$ and $Q$ belong to one of these sets, then no point of
            $l$ is between $P$ and $Q$. We say then that $P$ and $Q$ are on
            the same side of $l$.
      \item If $P$ and $Q$ are in different sets, then there is a point on
            $l$ which is between $P$ and $Q$.
      \end{applist}
\end{applist}

\APPsub{Linear Congruence Postulates}

\begin{applist}[resume]
\item Given points $A$ and $B$ on a line $l$ and a point $A'$ on a line
      $l'$, then on a given side of $A'$ on $l'$ there is exactly one point
      $B'$ such that $AB$ is congruent to $A'B'$. To indicate that $AB$ is
      congruent to $A'B'$, we write ``$AB \cong A'B'$.''

      In congruence the order of the points does not matter, that is, if
      $AB \cong A'B'$, then also $AB \cong B'A'$, $BA \cong A'B'$, and
      $BA \cong B'A'$.
\item Congruence also satisfies the following:
      \begin{applist}
      \item $AB \cong AB$.
      \item If $AB \cong A'B'$, then $A'B' \cong AB$.
      \item If $AB \cong A'B'$ and $A'B' \cong A''B''$, then
            $AB \cong A''B''$.
      \end{applist}
\item (This postulate tells how to ``add'' and ``subtract'' segments.)
      Suppose that $B$ is between $A$ and $C$ on a line $l$, and that $B'$
      is between $A'$ and $C'$ on a line $l'$.
      \begin{applist}
      \item If $AB \cong A'B'$ and $BC \cong B'C'$, then $AC \cong A'C'$.
      \item If $AC \cong A'C'$ and $BC \cong B'C'$, then $AB \cong A'B'$.
      \end{applist}
\end{applist}

\APPsub{Archimedes' Postulate}

\begin{applist}[resume]
\item If $AB$ is any segment on a line $l$ and $CD$ any other segment, then
      there is a finite number $n$ of points $A_1, A_2, \ldots, A_n$ on $l$
      such that the points $A, A_1, A_2, \ldots, A_n$ are arranged in this
      order, and all the segments $AA_1, A_1A_2, \ldots, A_{n-1}A_n$ are
      congruent to $CD$, and either $B = A_n$ or $B$ lies between $A_{n-1}$
      and $A_n$.
\end{applist}

\APPsub{The Completeness Postulate}

\begin{applist}[resume]
\item For every line $l$ and any given point $A$ on $l$, and for any
      positive real number $x$, there is, on a given side of $A$, a point
      $B$ such that $\sg{AB} = x$.
\end{applist}

\APPsub{Congruence Postulates for Angles}

\begin{applist}[resume]
\item If $\ang{A}$ is not a straight angle and if $r'$ is a ray from $A'$ on
      a line $l'$, then on a given side of $l'$ there is exactly one ray
      $s'$ from $A'$ such that the angle $A'$, with sides $r'$ and $s'$, is
      congruent to angle $A$. We write this as ``$\ang{A'} \cong \ang{A}$.''
      If $\ang{A}$ is a straight angle, then the straight angle at $A'$ with
      one side $r'$ is the only angle with one side $r'$ congruent to angle
      $A$.
\item ~
      \begin{applist}
      \item $\ang{A} \cong \ang{A}$.
      \item If $\ang{A} \cong \ang{B}$, then $\ang{B} \cong \ang{A}$.
      \item If $\ang{A} \cong \ang{B}$ and $\ang{B} \cong \ang{C}$, then
            $\ang{A} \cong \ang{C}$.
      \end{applist}
\item If in the figure any two of the pairs $\ang{A}$ and $\ang{A'}$,
      $\ang{B}$ and $\ang{B'}$, $\ang{C}$ and $\ang{C'}$ are congruent, then
      the third pair of angles are congruent. In other words:
      \begin{applist}
      \item If $\ang{A} \cong \ang{A'}$ and $\ang{B} \cong \ang{B'}$, then
            $\ang{C} \cong \ang{C'}$.
      \item If $\ang{B} \cong \ang{B'}$ and $\ang{C} \cong \ang{C'}$, then
            $\ang{A} \cong \ang{A'}$.
      \item If $\ang{C} \cong \ang{C'}$ and $\ang{A} \cong \ang{A'}$, then
            $\ang{B} \cong \ang{B'}$.
      \end{applist}
      The angle $A$ may be a straight angle.
\end{applist}

\FIGAPPONE

\begin{applist}[resume]
\item If for $\tri{ABC}$ and $\tri{A'B'C'}$, $AB \cong A'B'$,
      $AC \cong A'C'$, and $\ang{A} \cong \ang{A'}$, then
      $\ang{B} \cong \ang{B'}$ and $\ang{C} \cong \ang{C'}$.
\end{applist}

\APPsub{The Parallel Postulate}

\begin{applist}[resume]
\item Through a point not on a line there is no more than one line parallel
      to the given line.
\end{applist}

%========================================================= definitions
\APPhead{SELECTED DEFINITIONS}{Selected Definitions}

Although every definition has a purpose, certain ones are more necessary
than others, and most of these essential definitions are listed here. In
some cases an abbreviated form of the definition is given when no
misunderstanding can arise.

\begin{applist}
\item Each set of points in the plane is called a \emph{geometric figure}.
\item If $A$ and $B$ are any two points on a line, then the set of points
      consisting of $A$ and $B$ and all the points between $A$ and $B$ will
      be called the \emph{line segment} (or \emph{segment}) determined by
      $A$ and $B$. The segment will be indicated by ``$AB$.'' The points
      $A$ and $B$ are called the \emph{ends} of the segment.
\item If $A$, $B$, and $C$ are three points not all on a line, the set of
      points of the segments $AB$, $BC$, $AC$ is called a \emph{triangle}.
      We call the points $A$, $B$, and $C$ \emph{vertices} and the segments
      $AB$, $AC$, $BC$ \emph{sides} of the triangle. To indicate ``triangle
      $ABC$,'' we write ``$\tri{ABC}$.''
\item If $O$ is a point of a line $l$ and $A$ and $B$ are two other points
      of $l$, then we say that $A$ and $B$ are on the \emph{same side of
      $O$} if $O$ is not between $A$ and $B$. We say that $A$ and $B$ are on
      \emph{opposite sides of $O$} if $O$ is between $A$ and $B$.
\item A \emph{ray} from $O$, or a ray with end $O$, is a set of
      points consisting of $O$ and all the points on one side of $O$ of a
      line $l$ through $O$.
\item The set of all points on two rays from a point $O$ is called an
      \emph{angle}. The point $O$ is called the \emph{vertex} of the angle,
      and the two rays are called the \emph{sides} of the angle. If the two
      rays are on the same line, the angle is called a \emph{straight
      angle}.
\item Consider an angle, not a straight angle, with vertex $O$ and whose
      sides are the rays $r_1$ and $r_2$ on the lines $l_1$ and $l_2$,
      respectively. The \emph{interior of the angle} is the set of all
      points on the same side of $l_1$ as $r_2$ and on the same side of
      $l_2$ as $r_1$. The \emph{exterior of the angle} is the set of all
      points in the plane which are not points of the angle and are not in
      the interior of the angle.
\item Two angles are \emph{adjacent} to each other if they have a common
      side and vertex and have no common interior points.
\item Let $AB$ be any segment and $A'B'$ any other segment on a line $l'$.
      On the same side of $A'$ as $B'$, let $B''$ be the point such that
      $AB \cong A'B''$. If $B'$ is between $A'$ and $B''$, we say that
      $AB > A'B'$. If $B''$ is between $A'$ and $B'$, we say that
      $AB < A'B'$.
\item For the definition of \emph{length of a segment} read Chapter 5,
      Section 4.
\item A \emph{right} angle is one which is congruent to its supplement.
\item Two triangles are said to be \emph{congruent} if they can be labeled
      $\tri{ABC}$ and $\tri{A'B'C'}$ such that $AB \cong A'B'$,
      $AC \cong A'C'$, $BC \cong B'C'$, $\ang{A} \cong \ang{A'}$,
      $\ang{B} \cong \ang{B'}$, and $\ang{C} \cong \ang{C'}$. To indicate
      that ``triangle $ABC$ is congruent to triangle $A'B'C'$,'' we write
      ``$\tri{ABC} \cong \tri{A'B'C'}$.''
\item Let $A$ and $B$ be any angles, not straight angles. Let $\ang{B'}$ be
      an angle congruent to $\ang{B}$ such that one side of $\ang{B'}$ is
      the same as the first side of $\ang{A}$, and the other side of
      $\ang{B'}$ and the second side of $\ang{A}$ lie on the same side of
      the first side of $\ang{A}$. If the second side of $\ang{B'}$ is
      interior to $\ang{A}$, we say $\ang{B} < \ang{A}$. If the second side
      of $\ang{B'}$ is exterior to $\ang{A}$, we say $\ang{B} > \ang{A}$. If
      $\ang{A}$ is a straight angle and $\ang{B}$ is not a straight angle,
      we say $\ang{A} > \ang{B}$ and $\ang{B} < \ang{A}$.
\item Let $P_1, P_2, \ldots, P_n$ be any set of $n$ points. The set of
      points on all the segments $P_1P_2, P_2P_3, \ldots, P_{n-1}P_n$ is
      called a \emph{polygonal path}. The path is said to join, or connect,
      the endpoints $P_1$ and $P_n$. The points $P_1, \ldots, P_n$ are
      called the \emph{vertices}. In particular, if $P_1 = P_n$, the
      polygonal path is called a \emph{polygon} and designated as polygon
      $P_1P_2 \ldots P_{n-1}$.

      A polygon is \emph{convex} if for every two points on the polygon each
      point between them is either on the polygon or in the interior. A
      convex, equiangular, and equilateral polygon is called \emph{regular}.
\item Two lines are \emph{parallel} if they do not intersect.
\item If the ray $OB$ is in the interior of $\ang{AOC}$ (or $\ang{AOC}$ is
      straight), then $\ang{AOC}$ is the \emph{sum} of $\ang{AOB}$ and
      $\ang{BOC}$.
\item If $AB$ and $CD$ are segments, the \emph{ratio} of $AB$ to $CD$ is the
      number $\sg{AB} \div \sg{CD}$. The statement that the segments $AB$
      and $CD$ are \emph{proportional} to the segments $A'B'$ and $C'D'$
      means that $\sg{AB}/\sg{A'B'} = \sg{CD}/\sg{C'D'}$. We say that the
      sides of $\tri{ABC}$ are proportional to the sides of $\tri{A'B'C'}$
      if $\sg{AB}/\sg{A'B'} = \sg{BC}/\sg{B'C'} = \sg{CA}/\sg{C'A'}$.
\item A \emph{circle} is determined by a point $O$, called the center, and a
      positive number $r$. The circle is the set of all points $A$ of the
      plane such that the distance $\sg{OA} = r$. If $A$ is a point of the
      circle, then the segment $OA$ is called a \emph{radius} (plural,
      \emph{radii}). If $A$ and $B$ are on the circle and the segment $AB$
      contains $O$, then $AB$ is called a \emph{diameter}.
\item If a circle has center $O$ and radius $r$, the \emph{interior of the
      circle} is the set of all points $B$ such that $\sg{OB} < r$. The
      \emph{exterior of the circle} is the set of all points $C$ such that
      $\sg{OC} > r$. A \emph{chord} is a segment determined by two points
      $E$ and $F$ on the circle. The line through $E$ and $F$ is a
      \emph{secant}. A \emph{central angle} of a circle is an angle whose
      vertex is at the center of the circle. An \emph{inscribed angle} is an
      angle with vertex on the circle and a second point of the circle on
      each of the two sides of the angle.
\item Suppose that for each positive integer $n$ there is associated in some
      way, a number $x_n$. The numbers $x_1, x_2, \ldots, x_n$ then form an
      \emph{infinite sequence}. The number $x_n$ is called the $n$th term of
      the sequence. We designate the sequence by the symbol $\{x_n\}$.
\item The number $L$ is called the \emph{limit of the sequence} $x_n$ if,
      for large values of $n$, $x_n$ is necessarily arbitrarily close to
      $L$. We then say that $L$ is the limit of the sequence of terms $x_n$
      and write $L = \lim \{x_n\}$.
\item The \emph{circumference} $C$ of a circle is the limit of the
      perimeters $p_n$ of regular inscribed polygons. That is,
      $C = \lim \{p_n\}$.
\item The number $C/d$, which is the same for all circles, will be denoted
      by the Greek letter $\pi$ (pi).
\item The \emph{area $S$ of a circle} is the limit of the areas of inscribed
      regular polygons. That is, $S = \lim \{S_n\}$.
\item Given a circular arc, let $l_n$ be the length of a regular polygonal
      path of $2^n$ sides inscribed in the arc. The \emph{length $l$ of the
      arc} is defined to be the following limit:
      $l = \mathrm{limit}\, \{l_n\}$.
\item On a circle of radius $r$ an arc of length $\pi r/180$ is subtended by
      a central angle of \emph{one degree} ($1^\circ$).
\item If on a circle of radius $r$ a central angle subtends an arc of length
      $L$, then the \emph{measure of the angle in degrees} is
      $(L/\pi r) \cdot 180^\circ$.
\end{applist}

%========================================================= theorems
\APPhead{SELECTED THEOREMS}{Selected Theorems}

The theorems selected below have been chosen to help you remember how the
ideas have been developed and \emph{in what order}. There is no attempt at
completeness, and in some cases the statement is less precise than in the
text; when this is the case, you should try to fill in the details or refer
back to the text.

\begin{applist}
\item Two different lines intersect in at most one point.
\item Let $O$ be any point of a line $l$. The point $O$ separates the line
      into two sets, called the two sides of the point $O$ on the line,
      having the following properties: (1) If $P$ and $Q$ belong to the same
      side of $O$, then either $OPQ$ or $OQP$. (2) If $P$ and $Q$ belong to
      different sides of $O$, then $POQ$.
\item If $P$ and $Q$ are both in the interior of $\ang{O}$, then the segment
      $PQ$ is in the interior of $\ang{O}$.
\item If rays $r_1$ and $r_2$ form $\ang{O}$, a ray $r$ from $O$, other than
      $r_1$ or $r_2$ lies either completely in the interior or completely in
      the exterior of $\ang{O}$.
\item If $P$ is in the interior of $\ang{O}$ and $Q$ is in the exterior of
      $\ang{O}$, then the segment $PQ$ contains a point of the angle.
\item If $AB$ and $A'B'$ are any two segments, either $AB > A'B'$, or
      $AB \cong A'B'$, or $AB < A'B'$, and exactly one of these holds.
\item If $AB > CD$ and $CD \cong EF$, then $AB > EF$.
\item If $AB \cong A'B'$, then $\sg{AB} = \sg{A'B'}$.
\item If $\sg{AB} = \sg{A'B'}$, then $AB \cong A'B'$.
\item $AB > CD$ if and only if $\sg{AB} > \sg{CD}$.
\item The S.A.S. theorem for congruent triangles.
\item The A.S.A. theorem.
\item The S.S.S. theorem.
\item For any two angles $\ang{A}$ and $\ang{B}$ either
      $\ang{A} < \ang{B}$, or $\ang{A} \cong \ang{B}$, or
      $\ang{A} > \ang{B}$, and exactly one of these is true.
\item All right angles are congruent to one another.
\item An exterior angle of a triangle is greater than either opposite
      interior angle.
\item In $\tri{ABC}$, $AB > AC$ if and only if $\ang{C} > \ang{B}$.
\item In $\tri{ABC}$, $\sg{AB} + \sg{BC} > \sg{AC}$.
\item If $l$ is any line and $A$ any point not on $l$, then there exists at
      least one line on $A$ parallel to $l$.
\item If parallel lines are cut by a transversal, the alternate interior
      angles are congruent.
\item The sum of the angles of a triangle is a straight angle.
\item If a transversal is cut by three parallel lines in congruent segments,
      the same is true for every transversal of these lines.
\item If for $\tri{ABC}$ and $\tri{A'B'C'}$, $\ang{A} \cong \ang{A'}$,
      $\ang{B} \cong \ang{B'}$, and $\ang{C} \cong \ang{C'}$, then the
      triangles are similar.
\item If $\sg{AB}/\sg{A'B'} = \sg{AC}/\sg{A'C'} = \sg{BC}/\sg{B'C'}$, then
      $\tri{ABC}$ is similar to $\tri{A'B'C'}$.
\item In a right triangle the altitude on the hypotenuse forms two right
      triangles similar to the given triangle.
\item In $\tri{ABC}$ with $\ang{C}$ a right angle,
      $\sg{AC}^2 + \sg{BC}^2 = \sg{AB}^2$.
\item The area of a triangle is one half the base times the altitude.
\item Suppose that circles with centers $A$ and $B$ each have radius
      $r > \frac{1}{2}\sg{AB}$. Then these circles intersect in two points
      on opposite sides of the line $AB$.
\item Exactly one circle is circumscribed about a regular polygon.
\item If regular polygons of $n$ sides are inscribed in two circles, then
      the perimeters of the polygons are proportional to the radii of the
      circles.
\item For a given circle let $p_n$ be the perimeter of a regular polygon of
      $n$ sides inscribed in the circle. The numbers $p_n$,
      $n = 3, 4, 5, \ldots$, form a sequence which has a limit.
\item The ratio $C/d$ (circumference to diameter) is the same for all
      circles.
\item If $S_n$ is the area of a regular polygon of $n$ sides inscribed in a
      circle, then $\mathrm{limit}\, \{S_n\}$ exists.
\item The area of a circle is equal to one half the circumference times the
      radius.
\item If two central angles are congruent, they subtend arcs of equal
      length.
\item If $\widehat{AB}$ and $\widehat{BC}$ are adjacent arcs (that is, if
      they have only point $B$ in common), then length $\widehat{ABC}$
      ${}={}$ length $\widehat{AB}$ ${}+{}$ length $\widehat{BC}$.
\item The measure of an angle inscribed in a circle is one half the measure
      of the central angle which subtends the same arc.
\end{applist}
\clearpage
\end{document}
